\documentclass[11pt,a4paper]{article}
\usepackage[T1]{fontenc}
\usepackage[utf8]{inputenc}
\usepackage{lmodern}
\usepackage[margin=25mm]{geometry}
\usepackage{amsmath,amssymb,amsthm,mathtools,mathrsfs}
\usepackage{booktabs,longtable,array,enumitem}
\usepackage{microtype}
\usepackage{xurl}
\usepackage[colorlinks=true,linkcolor=blue,urlcolor=blue,citecolor=blue,hypertexnames=false]{hyperref}
\setlength{\emergencystretch}{3em}
\setlength{\parskip}{3pt}
\setcounter{tocdepth}{2}
\makeatletter
\renewcommand{\@pnumwidth}{2.5em}
\renewcommand{\@tocrmarg}{3.5em}
\renewcommand*\l@subsection{\@dottedtocline{2}{1.5em}{3.3em}}
\makeatother
\newtheorem{theorem}{Theorem}[section]
\newtheorem{proposition}[theorem]{Proposition}
\newtheorem{lemma}[theorem]{Lemma}
\newtheorem{corollary}[theorem]{Corollary}
\theoremstyle{definition}
\newtheorem{definition}[theorem]{Definition}
\newtheorem{remark}[theorem]{Remark}
\newtheorem{example}[theorem]{Example}
\renewcommand{\paragraph}[1]{\par\smallskip\noindent\textbf{#1}\quad\ignorespaces}
\newcommand{\Gal}{\operatorname{Gal}}
\newcommand{\Spec}{\operatorname{Spec}}
\newcommand{\Sha}{\operatorname{Sha}}
\providecommand{\longtwoheadrightarrow}{\relbar\joinrel\twoheadrightarrow}
\title{KTOWER\\The Galois action on the tower,\\the $K$-groups of $\mathbf Z$, and cote 149}
\author{Research continuation of K17, AN\_GRH, and GRH\_NIGHT\\GPT-6 Astra Pro; effort setting not exposed}
\date{21 November 2026 --- active research snapshot}
\begin{document}
\maketitle
\begin{abstract}
This report separates exact new results, carried theorems, and unresolved
arithmetic gates. It gives uniform compact-unitary and classical finite
models, independent trace and source audits, a stronger explicit bound
for higher $K$-groups, and newly completed finite vanishing certificates.
It also records counterexamples to finite-field section surjectivity
and unrestricted higher-dimensional prime-to-characteristic point
separation, and audits the marked tangential statements at cote149.
The all-weight vanishing target and the remaining Ihara and geometric
cover questions are not asserted solved.
\end{abstract}
\tableofcontents
\clearpage

\part{Checked findings and exact source comparisons}

% BEGIN sections/00_introduction.tex
\section{Scope, conventions, and present conclusions}

The principal target is $K_{4k}(\mathbf Z)=0$ for every positive
integer $k$, and hence Vandiver's conjecture at every prime.
That target remains open in this research snapshot.
The checked results include new complete finite certificates for
$K_{52}(\mathbf Z)=K_{56}(\mathbf Z)=K_{60}(\mathbf Z)=K_{64}(\mathbf Z)=K_{68}(\mathbf Z)=K_{72}(\mathbf Z)=K_{76}(\mathbf Z)=K_{80}(\mathbf Z)=0$, uniform dimension and
spectrum formulas, an independent computation of $157899$, a closed
source extraction of the archive's nine-hypothesis implication, and
\[
 \boxed{\ \log|K_{4k}(\mathbf Z)|<\frac14(2k)^4\quad(k\ge1).\ }
\]
This order bound is not a vanishing theorem. The completed finite
degrees, together with K17, give $K_{4k}(\mathbf Z)=0$ for
$1\le k\le20$. The all-weight arithmetic elimination remains
active. All finite computations started in this continuation are complete;
no higher degree is claimed.

The task's source convention is retained throughout: Alp\"oge's
\emph{Remark}, step~17, and the pinned Lean archive are cited as
theorems. This report does not constitute a new validation of the
archive's entire foundational or imported theorem chain. The fresh
source extraction is distinguished from a Lean kernel run: no Lean
executable was available and no build or new axioms report was run.
The source files actually retrieved have pinned Git-blob provenance.

\paragraph{Carried results.}
The supplied K17 paper and the intake reports settle the following
inputs, which are cited rather than re-derived:
K17-S17-UNIT; the arithmetic comparison VK-PT, (K1)--(K5), and
VK-EXCEPTIONS; K17-U3-HECKE's lattice class number and integral
harmonic model; K17-GL2-GENERAL's integral level-three basis;
K17-MINUS7-GENERAL-SOURCE; K17-U3-RESIDUAL-RING; and the finite
vanishings through $K_{48}$. Exact arithmetic comparison references
are Weibel, \emph{The K-book}, VI.8.2, VI.8.3, VI.10.9, and
Kurihara, \emph{Some remarks on conjectures about cyclotomic fields
and $K$-groups of $\mathbf Z$} (1992), Lemma~1.2, Corollary~1.5,
Proposition~3.7 and Remark~3.11, as reconciled in K17.
The modern comparison is required for the two-way zero-equivalence;
no unconditional 1992 proof of the then-conjectural comparison is
being attributed to Kurihara.

At odd primes below $2^{31}$ we use Hart--Harvey--Ong,
\emph{Irregular primes to two billion}, \S1, item~(2), as carried
by K17. Prime two is separate: Rognes--Weibel,
\emph{Two-primary algebraic $K$-theory of rings of integers in
number fields}, Theorem~0.7, proved as Theorem~7.11,
with K17's localization specialization, gives the required
$2$-primary vanishing. The large-modulus calculations never
assume a residual integer is prime.

For the geometric questions AN\_GRH A05, A33 and A46, and the
previous failed route-two constructions A36--A44, are carried at
their stated intake scope. K17-VANDIVER-KERNEL,
K17-GREENBERG-KERNEL, K17-P2-MOMENT, and K17-CUBIC-ODD retain
exactly their stated arithmetic hypotheses. No rational
mixed-Tate statement is substituted for an integral hypothesis.

\paragraph{Reading the paper and evidence.}
Part~I contains checked findings, including negative answers and
source comparisons. Part~II records unresolved questions and
proposals. A theorem that invokes a named background theorem is
checked as an implication from that background, rather than by
re-proving the entire literature. The numerical evidence includes
complete matrices, exact finite arithmetic, input hashes, and
scripts. Producer completion, a full independent replay, a
mathematical proof reading, and a formal kernel check are different
events in \path{CHECK_EVENTS.tsv}.

All file locators within the supplied K17, AN\_GRH and GRH\_NIGHT
returns refer to the exact attached versions. Scan numbers refer
to cote~149's indexing as preserved in those supplied transcripts;
no fresh palaeographic reading of every original sheet is claimed.
The registers \path{QUESTIONS.tsv}, \path{LEDGER.md}, and
\path{IMPLICATIONS.md} state the live boundaries of each conclusion.

% END sections/00_introduction.tex


% BEGIN evidence/lean_chain/onechar_section.tex
\section{The closed source proof of \texorpdfstring{$\mathrm{OneCharU\_s}$}{OneCharU\_s}}
\label{sec:kt-onechar-closed}

\paragraph{Finding KT-ONECHAR-CLOSED.}
Under the stipulated convention that the archive declarations are cited as
theorems, \path{UpperTheory.V.Z.J.OneCharU_s} has a closed proof term with
no additional mathematical hypotheses.  This is a direct source trace,
not a new Lean kernel build.  No Lean executable was available in this
session, and the accompanying file
\path{OneCharUClosed_UNCOMPILED.lean} is explicitly marked uncompiled.
The universal arithmetic conditions inside the definition of
\path{OneCharU_s}, including \path{GD}, remain part of the statement.

All locators in this section use the pinned commit
\path{7cbf7a0c32fca1badf05931879a37c3f8a2b03b4} of
\url{https://github.com/CocoSgt/riemann-hypothesis-lean}.
Abbreviate
\path{src/Development/AutomorphicArithmetic/AutomorphicVarieties/}
by AV/ and
\path{AV/ModularityArgument/ThresholdArguments/DenseThresholdCases/}
by DTC/.  The actual implication is
\path{UpperTheory.V.Z.J.oneCharU_of_katoTwoW_rig_admI_automorphy},
DTC/\path{FOneCharU.lean}:75--95.  Its nine hypotheses, in order, are
the following; all namespace prefixes shown are exact.

\begin{enumerate}
\item \path{UpperTheory.V.U3.CohInp}, supplied by
\path{UpperTheory.V.HV.Core0.cohInp_alt hH.litN hK},
DTC/\path{CohInp.lean}:16--17.
\item \path{UpperTheory.V.Alg.LemmaM}, supplied by
\path{UpperTheory.V.Alg.M.main_holds},
AV/\path{CommutativeAlgebra/GradedModulesLength/MainCotSpan.lean}:32--70.
\item \path{UpperTheory.V.Alg.LemmaM'}, supplied by
\path{UpperTheory.V.Alg.M.mprime_holds},
AV/\path{CommutativeAlgebra/GradedModulesLength/Mprime.lean}:38--80.
\item \path{UpperTheory.V.Alg.LemmaP}, supplied by
\path{UpperTheory.V.Alg.lemmaP},
AV/\path{CommutativeAlgebra/LemmaP.lean}:256--424.
\item \path{UpperTheory.V.Z.Kato.WeilSphericalRecords.KatoTwoW},
supplied by
\path{UpperTheory.V.Z.Kato.WeilSphericalRecords.katoTwoW_of_printsH hH},
DTC/\path{KatoTwoWOfPublishedH.lean}:48--49.
\item \path{UpperTheory.V.Th.RigidityFirstCase}, supplied by
\path{UpperTheory.V.Z.J.rigIUE_alt hC hH hL hE},
DTC/\path{ERigIU.lean}:80--146.
\item \path{UpperTheory.V.Th.RigiditySecondCase}, supplied by
\path{UpperTheory.V.Z.Win.RR.rigII_a_alt hC hH hL},
DTC/\path{RRRigII.lean}:25--102.
\item \path{UpperTheory.V.Z.J.AdmI_s}, supplied by
\path{UpperTheory.V.Z.J.admissible_typeI_even_level_alt hC hH hE},
DTC/\path{AdmissibleTypeIEvenLevel.lean}:76--136.
\item \path{UpperTheory.V.Aut.AutomorphyUnramifiedCoefficientsAlt},
supplied by
\path{UpperTheory.V.HV.Core0.automorphyUnramified_of_litN_of_gdEnd}
applied to \path{hH.litN} and
\path{(UpperTheory.V.Z.Win.Dn.ends0_alt hH).1};
DTC/\path{CohInp.lean}:18 and DTC/\path{DnEndsTop.lean}:15--16.
\end{enumerate}

Here \path{hC} is the term in item 1;
\path{hK := UpperTheory.V.ZK.lemmaK_litN_alt hH.litN}
uses DTC/\path{KFinal.lean}:11;
\path{hL := UpperTheory.V.HV.DescentVanishingReductions.adjoint_lifting_of_cohomology hH}
uses DTC/\path{AdjointLiftingOfCohomology.lean}:31; and
\path{hE := unlEnd3E (UpperTheory.V.Z.J.linkedUnlE_alt hH)}.
The latter application is DTC/\path{LinkedUnlE.lean}:51;
\path{linkedUnlE_alt} is declared there at lines 43--48 and
\path{unlEnd3E} at
AV/\path{ModularityArgument/LinkedUnramifiedChains/LogarithmicNormBounds/BlockValuationRep.lean}:234.

\paragraph{Proof of the closedness claim.}
The application with all nine proof terms is written literally in
DTC/\path{LinkedUnlE.lean}:49--54, in the proof of
\path{count_i_of_printed_bound_k_kato_automorphy}.
Its caller, DTC/\path{WindowBoundOfLiteratureLevel.lean}:25--29,
supplies \path{hK}, the Kato term in item 5, and the automorphy term in
item 9.  The remaining record can be closed by setting
\path{hH} equal to
\path{published_hypotheses_from_patching_alt}
applied to \path{UpperTheory.V.SmallClassNumbers.litN}
and \path{UpperTheory.V.Z.Tf.SR.descWeak}.
The record constructor is
DTC/\path{PublishedHypothesesFromPatching.lean}:98--103,
and \path{SmallClassNumbers.litN} is the closed theorem in
AV/\path{SmallClassNumbers/CorGKOf.lean}:20--23.
The other argument, \path{UpperTheory.V.Z.Tf.SR.descWeak}, is
the closed theorem at
AV/\path{ModularityArgument/TraceFormula/SplitRigidity/HasOfSplit.lean}:29--30,
from \path{descWeakAll} at lines 26--27.
Indeed, \path{CorGKOf.lean}:10--18 repeats the same inner
\path{OneCharU_s} proof term, and line 25 supplies precisely this closed
record to prove \path{corGK}.  Taking that inner term after the same
substitution proves the claim.  No inference from mere membership in
a syntactic dependency closure is needed. \hfill$\square$

\paragraph{The actual route from the named root.}
\path{src/Main.lean}:4 defines \path{riemannHypothesis := rhOfAll}.
\path{src/Supplements/Kernels/MainTheorems.lean}:7--9 takes
\path{rhOfAll} from \path{allChi}, itself
\path{UpperTheory.FinalTheorems.everyCharacter}.
\path{src/Development/AutomorphicArithmetic/CharacterBoxes/MainConclusion.lean}:7--11
applies \path{everyCharacter_of_row} to
\path{UpperTheory.Gal.Int.galoisLocalGlobalStrong_recParameter}
and uses \path{UpperTheory.One.vwin_of_row}.
DTC/\path{VwinOfRow.lean}:4--5 supplies the open-image record,
\path{UpperTheory.BaseChangeIdentities.baseChangeN}, and that same
Galois row to \path{automorphic_window_of_inputs}.
DTC/\path{TopMainZj.lean}:71--74 and 55--58 construct the required
\path{LiteratureRankNHypothesis} record.
DTC/\path{WindowBoundOfLiteratureLevel.lean}:30--34 and 25--29 then
reach the literal nine-term application already identified.
The Galois-row theorem itself is closed at
\path{src/Development/AutomorphicArithmetic/Foundations/GaloisRepresentations/IntegralStructures/RowSrc.lean}:6.
These declarations are cited archive inputs; their entire transitive
dependency tree was not re-proved in this session.

\paragraph{Exact indexing and multiplicity.}
DTC/\path{NumberU3US.lean}:57--59 makes $X$ a finite set of actual
three-dimensional Galois representations.  Besides admissibility it
requires
\[
\rho,\rho'\in X,\quad\operatorname{RepIso}(\rho,\rho')
\quad\Longrightarrow\quad\rho=\rho'.
\]
This is distinctness of isomorphism classes, not distinctness of one
Frobenius trace.  DTC/\path{FOneCharU.lean}:59--65 proves the
injectivity required to pass to $X$; lines 66--74 use
\path{Finset.sum_image}.  Thus a unit criterion needs pointwise coverage,
whereas a numerical product bound must keep enough eigenform indices
to cover all pairwise nonisomorphic systems, even if some traces coincide.

\paragraph{The range of the rational specialization.}
AV/\path{SmallClassNumbers/ClassNumberThree.lean}:64--86 proves
\path{GD} for the trivial character modulo $2$, $K=\mathbf Q(\sqrt{-3})$,
every even $e\ge1$, and every prime $\ell\ge5$ satisfying
$(\ell-1)\nmid2e$.  Its statement contains no inequality $e<\ell$.
The exclusion $\ell\ne7$ belongs to the subsequent Frobenius-$7$
evaluation, not to this \path{GD} theorem.

\paragraph{Interpretation and verification.}
This closes the source-standing gate in K17-S17-UNIT.  It supplies the
archive's bound on strict Galois-cohomology extension classes of Tate
twists over $\operatorname{Spec}\mathbf Z[1/\ell]$.
Vanishing still requires the arithmetic trace input and the carried
K17-VK-PT comparison.  The complete extraction proof and its source
applications were read and recorded; no new kernel run or global
axiom audit is asserted.

% END evidence/lean_chain/onechar_section.tex


% BEGIN evidence/unitary_audit/unitary_results.tex
\section{Uniform compact-unitary dimensions and an independent trace computation}
\label{sec:kt-unitary}

All statements here use the standard positive Hermitian form on
$E^3$, where $E=\mathbf Q(\omega)$ and $\omega^2+\omega+1=0$,
and the full stabilizer $K$ of $O_E^3$ at the finite places.
Write $V_d$ for the compact $U(3)$ representation of highest weight
$(d,0,-d)$ and $M_d=M(U(3),K,V_d)$.
The lattice class-number-one statement and
$\Gamma=U(3)(\mathbf Q)\cap K=\mu_6^3\rtimes S_3$ are carried
from K17-U3-HECKE, proof paragraphs ``The lattice class number'' and
``The invariant harmonic polynomial'', equations (UH.2)--(UH.4).
In particular, $M_d=V_d^\Gamma$ and $|\Gamma|=1296$.

\subsection{The necessary weight correction}
\begin{proposition}[KT-U3-WEIGHT]
The coefficient space used by K17 for the Selmer exponent $e$ is
$M_{e-1}$, whose highest weight is $(e-1,0,1-e)$.
Its Harish--Chandra parameter is $(e,0,-e)$.
Consequently the number $157899$ belongs to coefficient weight
$(5,0,-5)$ and parameter $(6,0,-6)$.
The literal coefficient space of highest weight $(6,0,-6)$ has
dimension three, not one.
\end{proposition}
\begin{proof}
The compact $U(3)$ half-sum of positive roots is $(1,0,-1)$.
Adding it to $(e-1,0,1-e)$ gives $(e,0,-e)$.
K17-U3-COVERAGE states this shift explicitly; K17-U3-HECKE
explicitly uses highest weight $(5,0,-5)$.
The dimension assertions follow independently from
KT-U3-DIM below. Thus the prompt's reference to weight $(6,0,-6)$
must be read as a parameter reference when discussing the original
$K_{12}$ certificate.
\end{proof}

\subsection{An exact dimension formula in every weight}

Put $c_j=[x^j](1+x+\cdots+x^5)^3$ and, for $d\geq0$, set
\begin{align*}
 I(d)&=\sum_{\substack{0\leq j\leq\min(15,d)\\j\equiv d\ (6)}}
      c_j\binom{(d-j)/6+2}{2}^{\!2},\\
 d_0+1&=6q+s,\qquad d_0=\lfloor d/2\rfloor,\quad0\leq s<6,\\
 T(d)&=6q^2+2qs+s,\qquad C(d)=\mathbf1_{3\mid d},\\
 A(d)&=\frac{I(d)+3T(d)+2C(d)}6,
 \qquad A(-1)=0.
\end{align*}

\begin{theorem}[KT-U3-DIM]
For every $d\geq0$,
\[
       \dim M_d=D(d):=A(d)-A(d-1).
\]
Equivalently, for $d=12n+r$, $n\geq0$, $0\leq r<12$, the
dimension is the following cubic quasipolynomial:
\[
\begin{array}{c|l@{\qquad}c|l}
r&D(12n+r)&r&D(12n+r)\\\hline
0&8n^3+2n^2+2n+1&6&(4n+3)(2n^2+2n+1)\\
1&2n(4n^2+2n+1)&7&(2n+1)(4n^2+6n+3)\\
2&8n^3+6n^2+4n+1&8&8n^3+18n^2+16n+5\\
3&(2n+1)(4n^2+2n+1)&9&2(n+1)(4n^2+6n+3)\\
4&(4n+1)(2n^2+2n+1)&10&8n^3+22n^2+22n+7\\
5&(2n+1)^3&11&8(n+1)^3
\end{array}
\]
In particular, $\dim M_5=1$, $\dim M_6=3$, and for every
positive $e\equiv0\pmod6$,
\[
                    \dim M_{e-1}=(e/6)^3.
\]
\end{theorem}
\begin{proof}
Use the harmonic realization of $V_d$ in bidegree $(d,d)$ in
three variables $z_i$ and their independent dual variables $w_i$.
The contraction map
$\Delta=\sum_i\partial_{z_i}\partial_{w_i}$ onto bidegree
$(d-1,d-1)$ is surjective, and its kernel is $V_d$.
This realization and its surjectivity in every degree were proved in
K17-HARMONIC, \path{harmonic_hecke_proof.md}, Section~1;
we use that result rather than reprove it.
Taking $\Gamma$-invariants is exact, since the characteristic is zero.
It therefore suffices to compute the invariant dimension in the full
bidegree space.

A monomial $z^a w^b$ is invariant under $\mu_6^3$ exactly when
$a_i\equiv b_i\pmod6$ for all $i$.
The remaining $S_3$ action permutes the three pairs $(a_i,b_i)$.
Burnside's lemma counts their orbits by averaging over the identity,
the three transpositions, and the two three-cycles.
For the identity, write $a_i=r_i+6u_i$ and $b_i=r_i+6v_i$ with
$0\leq r_i<6$. If $\sum r_i=j$, both $\sum u_i$ and $\sum v_i$
equal $(d-j)/6$, giving precisely $I(d)$ fixed monomials.
For a transposition the pairs in the two swapped positions agree.
Writing the exponents in these positions as $u,v$, the remaining
exponents are $d-2u,d-2v$. Thus $0\leq u,v\leq\lfloor d/2\rfloor$
and $u\equiv v\pmod6$; the congruence for the remaining position
then follows automatically. Among $d_0+1=6q+s$ integers, $s$
residue classes have size $q+1$ and $6-s$ have size $q$.
The count is $s(q+1)^2+(6-s)q^2=T(d)$.
A three-cycle fixes one monomial exactly when $3\mid d$, with
all six exponents equal to $d/3$. This gives $C(d)$.
The Burnside average is $A(d)$, and exactness gives $D(d)$.

For the table, substitute $d=12n+r$ in this proved finite formula.
One may sum over all $0\leq j\leq15$ with $j\equiv d\pmod6$:
the added terms have $(d-j)/6=-1$ or $-2$ and their polynomial
binomial coefficients vanish. Hence this substitution is a polynomial
identity, valid even at $n=0$, rather than an interpolation assumption.
Subtracting the adjacent $A$ polynomials gives the displayed table.
The last assertion is its $r=5,11$ cases.
\end{proof}

\paragraph{Scope.}
This is the total dimension. The endoscopic construction is supplied
separately in KT-U3-ENDO-INJECT; in particular KT-U3-ENDO-SIX
identifies the entire space at parameter $(6,0,-6)$ by its global
packet, in addition to the independent character computation below.

\subsection{A second method for the value at seven}

Write $\pi=3+\omega$ and $\bar\pi=2-\omega$.
For $r=1,2,3$, let $J_r$ be the $3$ by $3$ matrix whose first
$r$ by $r$ block consists of ones and whose other entries vanish.
Define
\[
 U_r=\pi I+\frac{\omega^{r-1}\bar\pi-\pi}{r}J_r,
 \qquad g_r=U_r/\pi.
\]
These are the three integral matrices
\[
\begin{aligned}
U_1&=\operatorname{diag}(2-\omega,3+\omega,3+\omega),\\[3pt]
U_2&=\begin{pmatrix}2+2\omega&-1+\omega&0\\
-1+\omega&2+2\omega&0\\0&0&3+\omega\end{pmatrix},
\\[3pt]
U_3&=\begin{pmatrix}1&-2-\omega&-2-\omega\\
-2-\omega&1&-2-\omega\\-2-\omega&-2-\omega&1\end{pmatrix}.
\end{aligned}
\]

\begin{theorem}[KT-U3-SECOND]
The scaled minuscule Hecke eigenvalue on $M_5$ is
\[
                         7^5T_7=157899.
\]
Equivalently, the normalized Satake trace is $157899/7^6$.
This is obtained by a character average on $\Gamma$, using the
three rank-one reflections above, independently of the harmonic
polynomial and neighbor data in the K17 computation.
\end{theorem}
\begin{proof}
The rank-one orthogonal projector onto the line spanned by
$(1,\ldots,1,0,\ldots,0)$ with $r$ ones is $J_r/r$.
Thus $U_r$ has eigenvalue $\omega^{r-1}\bar\pi$ on that line
and eigenvalue $\pi$ on its orthogonal complement.
Both eigenvalues have norm seven, so $U_r^*U_r=7I$.
The three displayed quotients are integral. Modulo $\pi$ they
have rank one, with image the indicated support line.
The elementary-divisor argument of K17-U3-HECKE, equations
(UH.7)--(UH.8), therefore identifies $g_r$ with the negative
minuscule correspondence at that line.

Reduction identifies $\mu_6$ with $\mathbf F_7^\times$.
Consequently $\Gamma$ acts on $\mathbf P^2(\mathbf F_7)$ with
exactly three orbits, determined by the number $r$ of nonzero
coordinates, of sizes
\[
             m_1=3,\qquad m_2=18,\qquad m_3=36.
\]
Every line has a representative with nonzero coordinates in
$\mu_6$, and conjugating $g_r$ by a corresponding monomial
unitary matrix gives its unitary reflection.
Distinct lines give distinct right cosets. Thus these independently
constructed matrices and their conjugates exhaust all $57$ cosets.

Let $P=|\Gamma|^{-1}\sum_{\gamma\in\Gamma}\rho_d(\gamma)$.
It is the projector onto $V_d^\Gamma$.
Because $P$ commutes with $\Gamma$, all conjugates of $g_r$ have
the same value of $\operatorname{tr}(P\rho_d(g_r))$.
The trace of the Hecke operator on $M_d$ is therefore
\[
 \operatorname{tr}(T_7\mid M_d)=
 \frac1{1296}\sum_{r=1}^3m_r
       \sum_{\gamma\in\Gamma}\chi_d(g_r\gamma).
\]
For a matrix $A$ let $h_j(A)$ denote the complete homogeneous
symmetric polynomial of degree $j$ in its eigenvalues, and put
$h_{-1}=0$.
The exact contraction sequence for the harmonic representation gives,
for a unitary $g$,
\[
 \chi_d(g)=h_d(g)\overline{h_d(g)}
                 -h_{d-1}(g)\overline{h_{d-1}(g)}.
\]
For $A=U_r\gamma$, multiplication by $7^d$ consequently gives
the integral expression
\[
       7^d\chi_d(g_r\gamma)
           =N(h_d(A))-7N(h_{d-1}(A)).                 \tag{KT-U.1}
\]
This computation needs no invariant harmonic polynomial and no
interpolation of its values.

Here is the complete compressed arithmetic certificate for $d=5$:
\[
\begin{array}{c|r|r|r}
r&m_r&\displaystyle\sum_{\gamma\in\Gamma}
       (N(h_5(U_r\gamma))-7N(h_4(U_r\gamma)))
       &\text{sum}/1296\\\hline
1&3&21781872&16807\\
2&18&5270832&4067\\
3&36&1233792&952
\end{array}
\]
Every entry is reproduced by the following fully specified recurrence.
For each of the six permutations $\sigma$ and each
$(u_1,u_2,u_3)\in\mu_6^3$, form $A_{ij}=(U_r)_{i,\sigma(j)}u_j$.
Set $s_1=\operatorname{tr}A$,
$s_2=\sum_{i<j}(A_{ii}A_{jj}-A_{ij}A_{ji})$, and $s_3=\det A$.
Starting at $h_0=1$ and $h_j=0$ for $j<0$, compute
\[
             h_j=s_1h_{j-1}-s_2h_{j-2}+s_3h_{j-3}
             \quad(1\leq j\leq5).
\]
Use $N(a+b\omega)=a^2-ab+b^2$ in (KT-U.1) and sum.
The standalone \path{character_trace.py} implements exactly
these integer operations and no imported K17 computation or data.
This is a finite exact computation, not a Lean kernel certificate.
Since $\dim M_5=1$, its trace is its sole eigenvalue, namely
\[
           3(16807)+18(4067)+36(952)=157899.
\]
Complex conjugation exchanges the two minuscule correspondences
and preserves this one-dimensional coefficient space, so they have
the same eigenvalue. Gross, \emph{On the Satake isomorphism},
Section~3, equations (3.13)--(3.14), gives $T_7^+=7t_7$ under
the maximal-compact-volume-one convention. This proves the stated
Satake normalization.
\end{proof}

\subsection{An independent reflection calculation at thirteen}

\begin{theorem}[KT-U3-SECOND13]
The independent character-projector method at thirteen gives
\[
          13^3T_{13}|_{M_3}=54291,
          \qquad 13^5T_{13}|_{M_5}=6110859.
\]
The same calculation determines the trace of $13^dT_{13}$ in every
coefficient degree $d$ by the finite formula below.
\end{theorem}
\begin{proof}
Put $\pi=4+\omega$, $\bar\pi=3-\omega$, and $\delta=1-\omega$.
For each row of the following table set
\[
      U=\pi I+\frac{\eta\bar\pi-\pi}{n}vv^*,
      \qquad n=v^*v,\qquad g=U/\pi.
\]
\[
\begin{array}{c|c|r|r|r|r}
v&\eta&n&\text{orbit size}&b_3&b_5\\\hline
(1,0,0)&1&1&3&2197&371293\\
(1,1,0)&-\omega^2&2&18&1417&129233\\
(1,\delta,0)&\omega^2&4&18&-143&38233\\
(1,1,1)&\omega&3&36&532&7048\\
(1,1,\delta)&-\omega&5&108&52&16008
\end{array}
\]
The coefficient $(\eta\bar\pi-\pi)/n$ is respectively
$-1-2\omega$, $\omega$, $-2-\omega$, $-1+\omega$, and
$-1-\omega$, so all five numerator matrices are integral.
The rank-one projector $vv^*/n$ shows that the eigenvalues of
$U$ are $\eta\bar\pi,\pi,\pi$. Consequently
$U^*U=13I$ and $\det U=\eta\pi^2\bar\pi$.
Its reduction modulo $\pi$ has rank one and image the reduction
of $v$. The same local elementary-divisor argument as in
KT-U3-SECOND therefore places each $g$ in the required minuscule
double coset.

Here is the complete orbit classification. Modulo $\pi$ we have
$\omega=-4$ in $\mathbf F_{13}$, and the six Eisenstein units
are exactly the six nonzero squares. The residue of $\delta$ is
$5$, a nonsquare. Support-one lines form one orbit of size three.
For support two, the square class of the ratio of the two nonzero
coordinates gives two orbits, each of size $3\cdot6=18$;
interchanging the two coordinates preserves that square class.
For support three, consider the three square classes up to common
scaling. They are either all equal or have one different class.
The all-equal class gives an orbit of size $6^2=36$.
Permuting which coordinate has the different class gives one orbit
of size $3\cdot6^2=108$. Within each pattern the independent unit
multiplications are transitive. These five disjoint orbits have
total size $183=13^2+13+1$, so they exhaust the minuscule cosets.
The program also enumerates every projective line and checks this
partition exactly.

For each numerator $U$ put
\[
 b_d(U)=\frac1{1296}\sum_{\gamma\in\Gamma}
       \bigl(N(h_d(U\gamma))-13N(h_{d-1}(U\gamma))\bigr),
 \qquad h_{-1}=0.
\]
The character identity and scaling argument of KT-U3-SECOND
apply verbatim with $7$ replaced by $13$. Thus
$\operatorname{tr}(13^dT_{13}|M_d)$ is the sum of these five
$b_d$ values weighted by their orbit sizes. The $d=3,5$ values
from the complete-homogeneous recurrence are displayed in the
last two columns. Their weighted sums are respectively $54291$
and $6110859$. Since $D(3)=D(5)=1$, they are the eigenvalues.
The exact program \path{character_trace13.py} imports only
the new character arithmetic from this audit and no archived
neighbor or operator data. It saves the five numerator matrices,
all 183 image lines, and all character sums through degree thirty.
\end{proof}

\subsection{A finite uniform formula for the entire spectrum}

\begin{proposition}[KT-U3-SPECTRUM]
The three reflections give a uniform exact formula for the
characteristic polynomial of $7^dT_7$ on $M_d$ for every $d\geq0$.
It does not give a uniform upper bound on the prime divisors of the
corresponding comparison integer.
\end{proposition}
\begin{proof}
For each projective line $L\in\mathbf P^2(\mathbf F_7)$ choose its
representative $v_L$ with first nonzero coordinate one and all
nonzero coordinates in $\mu_6$. If it has $r$ nonzero entries, put
\[
 U_L=\pi I+\frac{\omega^{r-1}\bar\pi-\pi}{r}v_Lv_L^*.
\]
The preceding proof gives a complete fixed list of $57$ integral
matrices, independent of $d$, each satisfying $U_L^*U_L=7I$.
The list is saved in \path{minuscule7_cosets.json}; the independent
\path{generate_cosets.py} checks every Gram identity, determinant,
and distinct projective image line.
Set $D=D(d)$. For $j\geq1$ and a sequence $L_1,\ldots,L_j$,
write $A=U_{L_1}\cdots U_{L_j}\gamma$.
The projector trace identity, applied to $T_7^j$, gives the exact
power sum
\[
 S_j=\operatorname{tr}((7^dT_7)^j\mid M_d)
 =\frac1{1296}
   \sum_{L_1,\ldots,L_j}\sum_{\gamma\in\Gamma}
       \bigl(N(h_d(A))-7^jN(h_{d-1}(A))\bigr).       \tag{KT-U.2}
\]
Indeed the product matrix satisfies $A^*A=7^jI$ and the same
homogeneity calculation as in (KT-U.1) applies.
Each $h_i$ is calculated by the three-term recurrence above.
Starting with $c_0=1$, define
\[
           c_j=-\frac1j\sum_{i=1}^j c_{j-i}S_i
           \quad(1\leq j\leq D).
\]
Newton's identities show that
$P_d(X)=X^D+c_1X^{D-1}+\cdots+c_D$ is the required
characteristic polynomial. Thus its roots, divided by $7^{d+1}$,
give all normalized Satake traces, with multiplicity, uniformly
in $d$. This is an exact finite formula, though its direct evaluation
is inefficient as $D$ grows.
The scaled action preserves the integral harmonic-polynomial
lattice over $O_E$, so $P_d\in O_E[X]$.
Formula (KT-U.2) puts every $S_j$ in $\mathbf Q$, and Newton's
identities therefore put $P_d$ in $\mathbf Q[X]$.
Since $O_E\cap\mathbf Q=\mathbf Z$, one has
$P_d\in\mathbf Z[X]$.
Complex conjugation exchanges the two minuscule correspondences
and conjugates their matrix coefficients. Their characteristic
polynomials are consequently equal. Thus the same polynomial
works for the positive operator required by the Satake formula.
No step in this construction bounds the prime factors of
$P_{e-1}(7^{2e}+7^e+1)$ independently of $e$.
\end{proof}

\paragraph{Check record and boundaries.}
The character program additionally verifies every reflection Gram
identity and rank-one reduction. Through degree $30$, its full
group-character averages agree with KT-U3-DIM, and independent
enumeration of exponent-pair orbits agrees with $A(d)$.
These finite comparisons check the implementation; the all-degree
dimension statement rests on the Burnside proof.
All-degree arithmetic prime elimination remains a separate question.
The identities $157899=7^6+40250$ and
$6110859=13^6+1284050$ have a global explanation in
KT-U3-ENDO-SIX: its proof first constructs the endoscopic packet
from the non-CM level-three weight-thirteen pair and then uses
the exact dimension to identify the compact form.

% END evidence/unitary_audit/unitary_results.tex


% BEGIN evidence/gl2_uniform_section.tex
\section{The classical level-three family in every even degree}
\label{sec:gl2-uniform}

Put $\chi=\chi_{-3}$ and let $e\geq2$ be even. The complete
classical space is $S_e=S_{2e+1}(\Gamma_0(3),\chi)$.
Its dimension and integral basis are already supplied by
K17-GL2-GENERAL, \path{gl2_general_section.tex}, (G1)--(G5).
They are cited here, with the empty-space convention at $e=2$:
\[
 r_e=\dim S_e=\left\lfloor\frac{2e-2}{3}\right\rfloor,\qquad
 f_j=D a^{\,2e-5-3j}b^j\quad(0\leq j<r_e),
 \tag{GL.1}
\]
where
\[
 a=\sum_{\alpha\in\mathbf Z[\omega]}q^{N\alpha},\qquad
 b=\frac{\eta(3\tau)^9}{\eta(\tau)^3},\qquad
 D=\eta(\tau)^6\eta(3\tau)^6 .
\]
The primary ring statement is Cl\'ery--van der Geer,
\emph{Generating Picard modular forms by means of invariant theory},
arXiv:2110.00849, paragraph \emph{Modular forms on $\Gamma_1(3)$},
printed p.~12; the same work, Remark~15.2, printed p.~33,
gives the weight-six cusp form. The basis begins
$f_j=q^{j+1}+\cdots$ and is integral and unitriangular.

\begin{proposition}[KT-GL2-CM]
The CM subspace of $S_e$ has dimension
\[
 c_e=\begin{cases}1,&3\mid e,\\0,&3\nmid e.\end{cases}
 \qquad
 \dim S_e^{\mathrm{nonCM}}=r_e-c_e.
 \tag{GL.2}
\]
When $3\mid e$, its unique normalized eigenform is the theta series
of the conductor-one Hecke character
$\psi_e((\alpha))=\alpha^{2e}$ of $\mathbf Q(\omega)$.
At a split prime $p=\pi\bar\pi$ its eigenvalue is
\[
             a_p(\psi_e)=\pi^{2e}+\bar\pi^{2e}.
 \tag{GL.3}
\]
At an inert prime its eigenvalue is zero.
\end{proposition}

\begin{proof}
The nebentypus has primitive conductor three. Therefore the space
has no oldforms from a proper divisor of the level, and every
normalized eigenform has primitive level three.
The CM classification says that such a form comes from an
algebraic Hecke character of an imaginary quadratic field, with
level the absolute field discriminant times the norm of its
conductor. We use exactly Sch\"utt,
\emph{CM newforms with rational coefficients},
arXiv:math/0511228v5, Theorem~1.1 (Ribet's classification),
printed p.~2, and Theorem~1.4 (Hecke--Shimura),
printed p.~3. The cited classification is not restricted to
rational coefficient fields.

The level formula forces the field to be $\mathbf Q(\omega)$
and the conductor to be the unit ideal. Its class number is one.
A character with infinity type $2e$ and conductor one exists
precisely when $\alpha^{2e}$ is unchanged on replacing a generator
of a principal ideal by a unit multiple. Since its units are
$\mu_6$, this means $6\mid2e$, equivalently $3\mid e$.
When this holds, class number one makes the character unique.
The character with conjugate infinity type gives the same theta
series, by reindexing ideals by conjugation, so it gives no second
eigenform. The two prime ideals over a split prime give (GL.3).
There is no ideal of norm an inert rational prime, giving zero
there. Hecke--Shimura supplies a cuspidal eigenform in exactly
the stated level and nebentypus. This proves both existence and
exhaustion. Subtracting its dimension from (GL.1) proves (GL.2).
\end{proof}

\subsection{A uniform exact determinant formula}

For $p\ne3$ let
\[
 L_{ij}=[q^i]f_j,\quad
 B_{ij}=[q^{pi}]f_j+
       {\bf1}_{p\mid i}\chi(p)p^{2e}[q^{i/p}]f_j
 \quad(1\leq i\leq r_e,\ 0\leq j<r_e).
\]
Then, by the already proved integral basis and the classical
Hecke formula,
\[
       Q_{e,p}(X)=\det(XI-L^{-1}B)\in\mathbf Z[X].
 \tag{GL.4}
\]
This gives every eigenvalue, with multiplicity, uniformly in $e$.
It is the matrix form of K17-GL2-GENERAL, not a new modularity
theorem. The required Fourier coefficients extend only through
degree $pr_e$.

The fresh computation \path{evidence/gl2_uniform.py}
constructs these matrices for every even $2\leq e\leq30$
and $p=7,13$. It expands the eta products by their logarithmic
derivatives. Explicitly, if
$F(q)=\prod_{n\geq1}(1-q^n)^s(1-q^{3n})^t
=\sum_{n\geq0}c_nq^n$, then
\[
 c_0=1,\qquad
 n c_n=\sum_{j=1}^n
   \bigl(-s\sigma_1(j)-3t\,{\bf1}_{3\mid j}\sigma_1(j/3)\bigr)c_{n-j}.
 \tag{GL.5}
\]
This follows by multiplying $qF'/F$ by $F$ and comparing
coefficients. Every division by $n$ is checked to be exact.
This implementation uses (GL.5), rather than the predecessor's
binomial expansion of the eta products.
It checks further Hecke coefficients after the identifying
ones and compares the characteristic-polynomial evaluation
with a direct determinant. The complete higher4 matrices at
$e=26,28,30$ agree entry by entry with the attached K17 data.
The known $e=6$ CM eigenvalue is independently recovered as
$-153502$ at seven; the other two eigenvalues are $40250$.

These classical calculations settle the classical numerical half
of the three requested higher degrees. Arithmetic vanishing
also uses the full unitary comparison and source coverage.

\paragraph{Geometric interpretation.}
The eigenforms give rank-two compatible systems of pure weight
$2e$; their twists occurring in the Selmer comparison have
weight zero. The CM summand is an induced rank-one system on
$\operatorname{Spec}\mathbf Q(\omega)$, and its Euler factor
at a split prime has roots $\pi^{2e},\bar\pi^{2e}$.
The decomposition (GL.2) is a statement about these classical
automorphic systems. It does not identify the stable and
endoscopic summands of the compact unitary space.

% END evidence/gl2_uniform_section.tex


% BEGIN evidence/unitary_audit/coverage_results.tex
\section{Coverage rechecked, and a stronger use of multiplicities}
\label{sec:kt-coverage}

\subsection{The source classes and normalization}
The source commit is
\path{7cbf7a0c32fca1badf05931879a37c3f8a2b03b4}.
Paths in this paragraph are relative to
\path{src/Development/AutomorphicArithmetic/AutomorphicVarieties/}.
The exact predicate is
\path{ModularityArgument/ThresholdArguments/DenseThresholdCases/NumberU3US.lean},
lines 27--29: \path{Adm3U_s} extends \texttt{Adm3Alt} and adds
\texttt{Th.Unl}. The latter base predicate, in \path{AfStmt.lean},
lines 29--32, supplies irreducibility, the automorphic condition for
$r=\rho\chi_\ell^{-1}$, and
$\operatorname{Level}(\ell,2,(e,0,-e),\rho)$.
\path{ModularityArgument/AutomorphyLifting/IsAutCsdA.lean},
lines 11--12, retains both \texttt{AutCsd0} and the implication
from continuity to \texttt{WitA}.
In the parallel \path{Foundations/GaloisRepresentations/AdmissibleResidues/}
directory, \path{IsAutomorphicOverImagQuad.lean}, lines 35--36,
defines \texttt{AutCsd0} by a regular algebraic conjugate-self-dual
\emph{cuspidal} $\mathrm{GL}_3$ representation with attachment.

\path{AutomorphicForms/BlochKatoSelmerVanishes.lean},
lines 66--67 and 89--114, defines the root-of-unity and local
level conditions. Because $\mu_2\subset\mathbf Q$, every
\texttt{KillsMu 2} character is trivial; the same level witnesses
therefore work at modulus one.
\path{BuildingsResidueChars/LevelTwistCyclotomicPow.lean},
lines 19--61, changes the exponent vector to
$(e-1,-1,-e-1)$ after the common twist.
Continuity is part of \texttt{Level}, so \texttt{WitA} is available.
\path{ModularityArgument/AutomorphyLifting/AdmissibleLevelRestriction/UnlAllOf.lean},
lines 36--38, gives the required above-$\ell$ nonadjacent flag;
the common twist preserves its inequalities.
\path{KeepLevel.lean}, lines 30--82 and 95--101, then supplies
an everywhere-finite-unramified cusp representation, its infinitesimal
character $(e,0,-e)$, and Frobenius--Satake compatibility at every
place away from $\ell$.
At modulus one the group \texttt{KLev} is the full integral adelic
maximal compact: this follows directly from
\path{AutomorphicForms/KLevTop.lean}, \path{mem_KLev},
lines 17--48. All these relevant proofs and definitions were read
directly in the pinned files; their imported foundations were not
independently rebuilt in Lean.

The rank-two predicate is
\path{ModularityArgument/AutomorphyLifting/AutomorphyDescent/AutNewformW.lean},
\texttt{Adm2A}, lines 14--20.
It supplies an irreducible continuous $W$, a weight-$2e+1$
newform attached to $W\chi_\ell^{-e}$, the actual
$\operatorname{Level}(\ell,6,(e,-e),W)$, and
$\rho\simeq1\oplus W|_{G_E}$, as well as its extra
\texttt{WitAQ} and determinant conditions.
The definition of \texttt{Newform} and its attachment in
\path{AutomorphicForms/Newform.lean}, lines 21--33 and 43--47,
includes normalization, a minimal-level clause, semisimplicity,
and cofinite equality of full Frobenius characteristic polynomials.

\subsection{Published input for the compact transfer}

\paragraph{Rechecked inherited theorem: K17-U3-COVERAGE and
K17-U3-GENERAL-COVERAGE.}
For every even $e\geq2$ and odd $\ell$, every source
$\operatorname{Adm3U\_s}(\iota,e,2)$ representation is covered
by the full compact space $M_{e-1}$, and its split-prime trace is
$\alpha/p^e$, where $\alpha$ is an eigenvalue of
$p^{e-1}T_p$ on that space.
The following checks explain why the inherited proof uses the
published theorems within their stated scope.

KMSW, \emph{Endoscopic classification of representations: inner forms
of unitary groups}, arXiv:1409.3731v3, Section~0.3.3, pp.~17--18,
identifies the Hermitian-space unitary groups as pure inner forms;
the standard positive form is one of them, and in odd rank its finite
local groups are quasi-split.
Propositions~1.3.1 and 1.3.3, pp.~67 and 70--71, attach and localize
the simple generic parameter of the cusp representation.
The parity criterion of Section~1.2.3 and Lemma~1.2.5, pp.~59--61,
applies because the distinct infinity characters
$(z/\bar z)^e,1,(z/\bar z)^{-e}$ all have positive conjugate parity.
The reduced global component group is trivial by the definitions
and formula (1.3.6), Section~1.3.4, pp.~69--70.
The compact member at infinity has highest weight $(e-1,0,1-e)$.
The real-packet relevance assertion is Bellaiche--Chenevier,
\emph{Families of Galois representations and Selmer groups},
Appendix~A.11, Remark~A.11.8 and its footnote~142, p.~236;
the entire remark and its explicit local-parameter argument were
read here. Only this local statement is used, not that appendix's
conditional global-parameter formalism.

The global multiplicity formula used is precisely KMSW
Theorem~5.0.5 and its preceding paragraph, p.~206.
That paragraph verifies its two hypotheses for generic parameters
and pure inner twists. Its proof, read in full here, gives the
direct sum over the packet selected by the global character.
Section~1.7, pp.~88--89, shows that the product of the local central
characters descends to the reduced component group. Since that
group is trivial here, independently chosen relevant local members
have multiplicity one.
No unrestricted version of the starred Theorem~1.7.1 is assumed.

Shin, \emph{Galois representations arising from some compact Shimura
varieties}, Annals of Mathematics 173 (2011), Theorem~1.2 and
Corollary~1.3, pp.~1647--1648, applies to regular algebraic
conjugate-self-dual cuspidal representations of odd rank.
It makes every finite local parameter tempered; its additional
slight-regularity condition applies only in even rank.
Corollary~7.9 and its printed proof, p.~1728, were also read.
At a split prime the required vector is the ordinary
$\mathrm{GL}_3(O_{E_w})$-fixed vector.
At an inert prime the parameter is unramified and the hyperspecial
member exists by KMSW Proposition~1.5.2(5), p.~79.

At the ramified prime three a full-stabilizer check is necessary.
The central character of the global $\Pi$ is unramified at every
finite place and trivial at infinity, hence trivial because $E$
has class number one. The roots of its local unramified parameter
are consequently $\{a,1,a^{-1}\}$, with $|a|=1$.
For $F=\mathbf Q_3$ and $L=E_3$, choose the unramified character
$\chi:L^\times\to\mathbf C^\times$ with Frobenius value $a$.
On the diagonal torus of $U(2,1)$ take the character
$\lambda(\operatorname{diag}(\alpha,\beta,\bar\alpha^{-1}))
=\chi(\alpha)$.
Its stable torus lift is $(\chi,1,\chi^{-1})$.
The induced $\mathrm{GL}_3(L)$ representation is irreducible
because these unramified characters are unitary.

Adler--Lansky, \emph{Depth-zero base change for ramified $U(2,1)$},
arXiv:0807.1528v2, Section~4.4, pp.~14--15, lists the reducibility
possibilities. The nonunitary case is excluded, and the case whose
restriction to $F^\times$ is the ramified quadratic character is
excluded because $\chi$ is unramified. The remaining reducible
case is the nontrivial unramified quadratic $\chi$, whose two
constituents form one packet.
Proposition~6.1(1),(4), p.~21, and Proposition~7.1(i),(iv), p.~28,
identify the stable base change of the irreducible or two-member
packet as the required $\Pi_L$.
The proof of Proposition~7.1 refers to the earlier unramified paper:
arXiv:math/0601695, Proposition~5.1, pp.~19--20 (published
numbering Proposition~4.1). Its entire proof, read here, derives
the used cases directly from Rogawski, Proposition~4.10.2 and
the paragraph before Theorem~13.2.1.
Thus neither the separate supercuspidal restriction
$q>3$, $q\ne7$, nor the depth-preservation Hypothesis~1.5 is an
input to these principal-series cases.

For completeness the standard lattice is the full vertex-$y$
stabilizer in Adler--Lansky Sections~4.1 and 4.2.1, pp.~10--12.
Choose $b\in\mathbf Z_3$ with $b^2=-2$ by Hensel's lemma.
The vector $(1,1,b)$ is primitive isotropic; pairing it with an
integral vector of pairing one and subtracting half the latter's
norm gives an integral hyperbolic plane. If $d$ is the norm of
the complementary unit line and $B$ is the basis matrix,
$-d=N(\det B)$; scaling the last vector by $(\det B)^{-1}$
makes its norm $-1$. This identifies the form with their matrix
$\left(\begin{smallmatrix}0&0&1\\0&-1&0\\1&0&0\end{smallmatrix}\right)$.
Their full reductive quotient is $O_3(k)$, including its central
component, and the vertex is special. Iwasawa decomposition gives
one full-stabilizer-fixed vector in the principal series because
both $\lambda$ and the modulus are trivial on $B(F)\cap K_3$.
Exactness of compact-open invariants in characteristic zero places
such a vector in some constituent of the required packet.
Global assembly by KMSW now proves the inherited coverage.

Finally \texttt{FrobEig} uses geometric Frobenius, whereas
$\chi_\ell(\sigma)=p$ for arithmetic Frobenius.
Therefore $r(\sigma^{-1})=p\rho(\sigma^{-1})$.
The source's \texttt{algSat},
\path{BlochKatoSelmerVanishes.lean}, lines 69--82, multiplies
the unitary Satake roots in rank three by $p$.
The two factors cancel, and Gross's minuscule normalization
$T_p=p\sum\beta$ gives the asserted $\alpha/p^e$.

\subsection{Published input for the rank-two conductor}

\paragraph{Rechecked inherited theorem: K17-GL2-CONDUCTOR.}
For $e\geq1$ and $\ell>3$, every
$\operatorname{Adm2A}(\iota,e,6)$ is represented by a normalized
newform $f\in S_{2e+1}(\Gamma_0(3),\chi_{-3})$, of primitive
conductor exactly three, and has split-prime trace
$1+a_p(f)/p^e$.

Here is a reconstructed proof with the exact source scopes.
The source attachment identifies $R=W\chi_\ell^{-e}$ with
the semisimple Galois representation of the primitive eigenform,
by Chebotarev and Brauer--Nesbitt.
Every \texttt{KillsMu 6} character factors through
$\operatorname{Gal}(\mathbf Q(\mu_3)/\mathbf Q)$, so it is
$1$ or $\chi_{-3}$.
At $q\ne\ell$ the actual inertia action is a direct sum of these
characters. It has finite image and zero monodromy.
Thus the local parameter is unramified at $q\ne3,\ell$, and at
three it is tame with these two characters and zero monodromy.
Carayol, \emph{Sur les representations $\ell$-adiques associees
aux formes modulaires de Hilbert}, Section~0.7, Theorem~A,
pp.~410--411, gives the full away-from-$\ell$ compatibility.
Its extra discrete-series hypothesis, Section~0.3, p.~409,
is imposed only when the totally real degree is even, so it
does not apply over $\mathbf Q$.

At $\ell$, the local flag after twisting has characters
$\delta_1$ and $\delta_2\chi_\ell^{-2e}$ with unramified
continuous $\delta_i$.
Their Frobenius values are $\ell$-adic units because the image
of a compact group under a continuous character is compact.
Thus the two crystalline characters have slopes differing by $2e$.
Saito, \emph{Galois representations in arithmetic geometry},
Section~2c, Theorem~3, p.~9, with the definitions on pp.~7--8,
gives potential semistability and full Frobenius-semisimplified
Weil--Deligne compatibility for classical modular forms over
$\mathbf Q$, including monodromy. This author-stated theorem
refers explicitly to Saito's 1997 \emph{Modular forms and $p$-adic
Hodge theory}; the original 1997 proof was not reread here.

After a finite extension over which $R$ is semistable, exactness
of $D_{\rm st}$ shows that its two slopes are those of the two
flag quotients. The relevant exactness statement is
Colmez--Fontaine, \emph{Construction des representations $p$-adiques
semi-stables}, author version, Section~4.1, Proposition~4.2,
p.~23. The rank-one period calculation is Section~4.2,
Proposition~4.4 and its proof, p.~24.
Section~3.3, p.~21, gives $N\varphi=\ell\varphi N$ and says
that $N$ lowers slopes by one. Since $2e\geq2$, the monodromy
operator in this two-dimensional module is zero.
The Weil--Deligne inertia has finite image and its two graded
characters are trivial; averaging over that finite group makes
its entire action trivial. Saito's theorem therefore makes $f$
unramified at $\ell$ as well.

The nebentypus is now unramified outside three and has conductor
dividing three. Its odd weight forces it to be odd, hence equal
to $\chi_{-3}$. Exactly one of the two inertia characters at
three is consequently nontrivial. Their tame Artin conductor
is one, and zero monodromy adds nothing, proving primitive
conductor three. Finally
$R(\sigma^{-1})=p^eW(\sigma^{-1})$ and attachment at the split
prime give the trace formula. All extra source predicates are
retained; this proof simply does not need to weaken them.

\subsection{Multiplicity is available when the whole systems are retained}

\begin{lemma}[KT-TRACE-MULTIPLICITY]
Fix an even $e\geq2$, $\ell>3$, $\ell\ne7$, and an embedding
$\iota:\overline{\mathbf Q}_\ell\simeq\mathbf C$.
Let $X$ be a finite set of pairwise nonisomorphic source
representations, each satisfying either
$\operatorname{Adm3U\_s}(\iota,e,2)$ or
$\operatorname{Adm2A}(\iota,e,6)$.
Let $P_{e,7}$ be the characteristic polynomial of
$7^{e-1}T_7$ on the full compact space $M_{e-1}$, and
$Q_{e,7}$ the characteristic polynomial of $T_7$ on the full
$S_{2e+1}(\Gamma_0(3),\chi_{-3})$.
Put
\[
 N_{e,7}=P_{e,7}(7^{2e}+7^e+1)\,Q_{e,7}(7^{2e}+1).
\]
With $s$ a geometric Frobenius at a selected prime above seven,
$t=7^{-e}+1+7^e$, and $v_\ell(\ell)=1$, one has
\[
 \sum_{\rho\in X}v_\ell(\operatorname{tr}\rho(s)-t)
                       \leq v_\ell(N_{e,7}).           \tag{KT-M.1}
\]
The polynomial evaluations are nonzero. Under the source
$\mathrm{GD}$ hypothesis and $\mathrm{OneCharU\_s}$,
$V_n\leq2v_\ell(N_{e,7})$ follows.
\end{lemma}
\begin{proof}
Partition $X$ into its irreducible rank-three branch and its
reducible $1\oplus W|_{G_E}$ branch. These branches are disjoint.
For each rank-three member choose its attached cusp representation
$\Pi$. If two members had the same $\Pi$, their twists would be
continuous semisimple representations with equal full Frobenius
characteristic polynomials outside a finite set, by
\texttt{AttachedSys},
\path{BlochKatoSelmerVanishes.lean}, lines 77--82.
Chebotarev and Brauer--Nesbitt would make the two members
isomorphic, contrary to the hypothesis.
Thus the chosen $\Pi$ are distinct.
The generic multiplicity formula puts their compact transfers
in distinct automorphic parameter summands. Selecting one
$K$-fixed Hecke eigenline in each summand gives linearly
independent eigenvectors in $M_{e-1}$.
Accordingly their eigenvalues at seven occur in $P_{e,7}$ with
at least the required multiplicities, even when several values
at seven coincide.

For the second branch choose the normalized primitive newform
supplied by the conductor theorem. A fixed form determines $W$
up to isomorphism by its semisimple attachment and hence determines
$1\oplus W|_{G_E}$. Distinct elements of this branch therefore
require distinct normalized primitive eigenforms. These belong
to the standard simultaneous eigenbasis of the full cusp space,
so their eigenvalues occur in $Q_{e,7}$ with sufficient
multiplicity.

After applying $\iota^{-1}$, the two lists of trace differences
are respectively
\[
 \frac{\alpha-(7^{2e}+7^e+1)}{7^e},\qquad
 \frac{a_7-(7^{2e}+1)}{7^e}.
\]
Both characteristic polynomials are monic integral, so all their
roots are integral at $\ell$. Since $\ell\ne7$, the displayed
denominators have valuation zero. Each selected factor has
nonnegative valuation; adding all remaining characteristic
factors only increases the sum. Their product is $N_{e,7}$,
which proves (KT-M.1).

For nonvanishing, equip $V_{e-1}$ with a unitary norm.
The compact Hecke operator is a sum of $57$ unitary operators,
so every scaled eigenvalue has complex absolute value at most
$57\cdot7^{e-1}<7^{2e}+7^e+1$ for $e\geq2$.
The classical bound $|a_7|\leq2\cdot7^e<7^{2e}+1$ gives the
other nonvanishing. Alternatively its weaker elementary Hecke
operator norm also suffices.
The source comparison trace is therefore different from every
admissible trace, and \path{OneCharU_s} applies.
Its finite set is pairwise nonisomorphic by
\path{NumberU3US.lean}, lines 57--59, so (KT-M.1) bounds its
valuation sum as asserted.
\end{proof}

\paragraph{A bound for the comparison integer.}
If $D_e=\dim M_{e-1}$ and
$r_e=\dim S_{2e+1}(\Gamma_0(3),\chi_{-3})$, the same norms give
\[
 \log|N_{e,7}|\leq(D_e+r_e)(2e\log7+1)\qquad(e\geq2).
\]
Indeed
$7^{2e}+7^e+1+57\cdot7^{e-1}
 =7^{2e}(1+(64/7)7^{-e}+7^{-2e})$
and the parenthesis is less than $\exp(1)$ for $e\geq2$;
the classical factor satisfies the same estimate.
This is a bound on size, not on smoothness or on the prime
divisors individually. Turning it into an all-prime bound on
$K_{2e}(\mathbf Z)$ additionally requires the exact source
$\mathrm{GD}$ conditions and the inherited Selmer--$K$ comparison
for each prime; those arithmetic conditions must be checked
separately.

\begin{theorem}[KT-QUARTIC-BOUND]
Under the report's carried-source convention, for every even
$e\geq2$,
\[
      |K_{2e}(\mathbf Z)|\mid N_{e,7}^{\,2},\qquad
      \log|K_{2e}(\mathbf Z)|<\frac25 e^4
      \quad(e\geq2\text{ even}).                     \tag{KT-M.2}
\]
The first divisibility is understood with $|N_{e,7}|$ if necessary.
This is an explicit alternative bound; it does not evaluate the
particular existential constant selected inside VK-EXP7.
\end{theorem}
\paragraph{Dependencies.}
The source theorem used is KT-ONECHAR-CLOSED, the closed
proof-term extracted from
\path{SmallClassNumbers/CorGKOf.lean}, lines 10--25,
with the \path{OneCharU_s} implication at lines 15--18 and
the published-hypotheses instantiation at lines 20--25.
Its exact implication is in \path{DenseThresholdCases/FOneCharU.lean},
lines 75--95. The source chain is read and audited in this report;
it is not a new compilation of the entire Lean project.
The arithmetic inputs are the inherited Selmer--$K$ comparison
VK-PT and exceptional-character vanishing VK-EXCEPTIONS
(\path{vandiver_k.md}, Sections 9.2--9.3), the stipulated
Vandiver verification for odd primes below $2^{31}$ with
Kurihara's zero-equivalence, and the carried Rognes--Weibel
vanishing of the $2$-primary part of $K_{4k}(\mathbf Z)$.
The needed \texttt{GD} assertion is a proved source statement,
\path{SmallClassNumbers/ClassNumberThree.lean}, lines 64--86;
its full displayed proof was also read here.
\begin{proof}
For every prime $\ell\geq2^{31}$ with $\ell-1\nmid2e$,
the quoted \texttt{gd} theorem supplies $\mathrm{GD}$ for
$E=\mathbf Q(\sqrt{-3})$, modulus two and trivial character.
Its hypotheses are precisely $e\geq1$, $e$ even,
$\ell\geq5$, and $\ell-1\nmid2e$; no upper bound on $e$
occurs. Apply
KT-TRACE-MULTIPLICITY to the finite set supplied by
$\mathrm{OneCharU\_s}$ and use VK-PT to identify its normalized
length with $v_\ell(|K_{2e}(\mathbf Z)|)$.
If $\ell-1\mid2e$, VK-EXCEPTIONS makes the primary group zero.
All odd primes below $2^{31}$, including seven, contribute zero
by the stated Vandiver theorem and Kurihara's equivalence.
The $2$-primary part is zero by the separate Rognes--Weibel input.
Thus every prime exponent in the group order is bounded by twice
its exponent in $N_{e,7}$, proving the divisibility and
$\log|K_{2e}(\mathbf Z)|\leq2\log|N_{e,7}|$.

The exact dimension formula specializes for positive even $e$ to
\[
D_e=\begin{cases}
e^3/216,&e\equiv0\pmod6,\\
e^3/216+e/9-7/27,&e\equiv2\pmod6,\\
e^3/216+e/9+7/27,&e\equiv4\pmod6.
\end{cases}
\]
The full classical dimension is
$r_e=\lfloor(2e-2)/3\rfloor$.
Hence for $e\geq4$ even,
\[
 S_e:=D_e+r_e\leq e^3/216+7e/9-11/27.
\]
Dividing this right side by $e^3$ gives
$1/216+7/(9e^2)-11/(27e^3)$.
Its derivative is $(33-42e)/(27e^4)<0$ for $e\geq4$,
and its value at four is $3/64$. Therefore $S_e\leq3e^3/64$.
Using the preceding bound for $N_{e,7}$ and $\log7<2$,
\begin{align*}
 \log|K_{2e}(\mathbf Z)|
 &\leq 2S_e(2e\log7+1)\\
 &<\frac3{32}e^3(4e+1)
 \leq\frac{51}{128}e^4<\frac25e^4.
\end{align*}
For $e=2$ both spaces have dimension zero, so the divisibility
already says $|K_4(\mathbf Z)|=1$ and the strict bound follows.
All arithmetic interface conditions have now been supplied by
the listed inputs. Pointwise trace coverage alone would not
give this conclusion: the whole-system injection in
KT-TRACE-MULTIPLICITY is what preserves the necessary
valuation multiplicities.
\end{proof}

% END evidence/unitary_audit/coverage_results.tex


% BEGIN evidence/unitary_audit/endoscopic_results.tex
\section{A proved endoscopic injection and the origin of the small eigenvalue}
\label{sec:kt-endoscopic}

This section uses the same standard compact unitary group and full
lattice stabilizer as KT-U3-DIM. Write $\epsilon=\chi_{-3}$,
$E=\mathbf Q(\omega)$, and let $e\geq2$ be even.
For a normalized primitive form
$f\in S_{2e+1}(\Gamma_0(3),\epsilon)$, write $\pi_f$ for its
unitary automorphic representation of $\mathrm{GL}_2(\mathbf A_{\mathbf Q})$.
Thus its central character is $\epsilon$, and at a good prime the
sum of its unitary Satake roots is $a_p(f)/p^e$.

\begin{theorem}[KT-U3-ENDO-INJECT]
Primitive automorphic twisting by $\epsilon$ is an involution on
the normalized newforms of $S_{2e+1}(\Gamma_0(3),\epsilon)$.
Its fixed points are exactly the CM forms. For each two-element
orbit $\{f,f\otimes\epsilon\}$ there is a Hecke eigenline in
$M_{e-1}$ with generic endoscopic parameter
\[
             \psi_f=1\boxplus\operatorname{BC}_{E/\mathbf Q}(\pi_f).
\]
Different two-element orbits give independent eigenlines.
At every split prime $p\ne3$ their scaled eigenvalue is
\[
                  \alpha_p=p^e+a_p(f).                 \tag{KT-E.1}
\]
Consequently the dimension of the endoscopic part of $M_{e-1}$
is at least
\[
 \frac{r_e-c_e}{2}=\left\lfloor\frac{e-1}{3}\right\rfloor,
 \qquad
 r_e=\left\lfloor\frac{2e-2}{3}\right\rfloor,
 \quad c_e=\mathbf1_{3\mid e}.                         \tag{KT-E.2}
\]
For $e=4$ and $e=6$, this is the entire one-dimensional compact
space. In particular its stable part is zero in those two degrees.
\end{theorem}

\begin{proof}
\paragraph{Twists, base change, and the local level.}
The exact cyclic base-change inputs are Langlands,
\emph{Base change for $\mathrm{GL}(2)$}, author retypeset version,
Section~2, properties (A)--(F), p.~13, especially (C), and
Lemma~11.3(b), p.~133, Proposition~11.4, p.~136, and
Lemma~11.7, p.~142. They give existence with local compatibility,
cuspidality outside the dihedral exception, the two-element
quadratic-twist fibres, and the characterization of a nontrivial
quadratic self-twist by automorphic induction from that quadratic
field. These statements were read directly; the entire trace-formula
proof of base change was not independently reconstructed here.

The level and primitive nebentypus force the local conductor of
$\pi_f$ at three to be one and its determinant on inertia to be
the ramified quadratic character. Its Weil--Deligne representation
there has zero monodromy and inertial characters $1,\epsilon$.
Here is the elementary conductor check: the nontrivial determinant
forces the Artin conductor of the Weil representation to be at
least one. The Weil--Deligne conductor is this Artin conductor
plus the nonnegative monodromy term, so conductor one forces
Artin conductor one, one inertia invariant, and no Swan
contribution. Semisimple tame inertia is therefore exactly
$1\oplus\epsilon$. Monodromy commutes with inertia; on these
two distinct one-dimensional inertia lines a nilpotent commuting
operator is zero. Frobenius preserves the two distinct inertia lines,
so the local parameter is a principal-series pair
$\mu\oplus\nu\epsilon$ with $\mu,\nu$ unramified.
The conductor and local-parameter identification use the same
classical newform/local-compatibility input as K17-GL2-CONDUCTOR.

Twisting this pair by $\epsilon$ again gives conductor one.
At other finite primes the twist is unramified, and its weight
and nebentypus are unchanged. Hence the associated primitive
normalized twist has level three. This is an assertion about
the primitive automorphic twist; the coefficientwise Dirichlet
twist at multiples of three need not itself be that primitive
newform.
Langlands's self-twist criterion and KT-GL2-CM now show that
the fixed points are precisely the CM forms, all of which here
are induced from $E$.
For a nonfixed orbit the base change $\Pi=\operatorname{BC}(\pi_f)$
is cuspidal. Its finite local parameters are unramified:
the only ramified inertial character just exhibited is killed
on restriction to $W_{E_3}$, and there is no monodromy.
Its central character is $\epsilon\circ N_{E/\mathbf Q}=1$.
Also $\Pi^c\simeq\Pi$ by base change and
$\Pi^\vee\simeq\Pi$ by the rank-two trivial determinant,
so $\Pi$ is conjugate self-dual.
Its infinity characters are $(z/\bar z)^e$ and $(z/\bar z)^{-e}$,
distinct positive-parity characters. Here the global sign deduction
uses both KMSW seed theorems: Proposition~1.3.1, p.~67,
attaches $\Pi$ to the unique embedding $\eta_{\chi_\kappa}$,
and Proposition~1.3.3, p.~71, localizes this same embedding
at infinity. The local image has parity $(-1)^{2-1}\kappa$
by the bilinear-form definition in Section~1.2.3, p.~60,
and the construction in Section~1.2.4, Lemma~1.2.5, p.~61.
The two distinct one-dimensional positive-parity summands
cannot support a negative-parity form: such a form must be
diagonal on these distinct conjugate-self-dual characters, and
an alternating form on a one-dimensional space is zero.
Thus $-\kappa=1$. Proposition~1.3.1 identifies this as
global conjugate-orthogonal parity, so $\Pi$ is a permitted
positive-parity simple constituent of the rank-three parameter.
No even-rank application of Shin is needed.

The finite local parameters of $\Pi$ are tempered. At good
primes this follows from the classical Ramanujan bound.
At three the two principal-series characters of $\pi_f$ have
distinct restrictions to inertia. A unitary principal series
with such characters is tempered: the non-tempered unitary
principal series has inducing characters with the same unitary
part, and hence the same restriction to inertia. Base change
preserves this temperedness. Thus at every finite nonsplit
place the roots of $\Pi$ are $\{b,b^{-1}\}$, with $|b|=1$.

\paragraph{The compact packet member and its multiplicity.}
The global parameter $\psi=1\boxplus\Pi$ is discrete and generic.
Its reduced component group is of order two, generated by the
class of $a$, where $a$ acts as $-I_2$ on $\Pi$ and as $1$ on
the one-dimensional summand. At split finite places choose the
usual spherical representation. At inert places choose the
hyperspecial member of the unramified packet. At the ramified
place three, the full-stabilizer principal-series argument in
the coverage section supplies a $K_3$-fixed packet member for
the parameter $\{b,1,b^{-1}\}$. More explicitly, take the
unramified $E_3^\times$ character $\lambda_1$ with value $b$
and take $\lambda_2=1$ in Adler--Lansky's notation (4.4.1).
Section~4.4, pp.~14--15, excludes (4.4PS--1) by unitarity
and (4.4PS--2) because its restriction to $\mathbf Q_3^\times$
is ramified. The only reducible possibility is (4.4PS--3),
the nontrivial unramified quadratic character $b=-1$.
For $b=1$ the principal series is irreducible.
Proposition~6.1(1),(4), p.~21, gives the one-member or
two-member packet, and Proposition~7.1(i),(iv), p.~28,
gives its base change. The induced $\mathrm{GL}_3$ representation
is irreducible even when roots repeat: a ratio of its unitary
characters cannot equal $|\cdot|^{\pm1}$.
Iwasawa decomposition for the full stabilizer and the triviality
of $\lambda$ on $B\cap K_3$ give one fixed vector in the
principal series; exactness puts it in a constituent of this
packet. These are local statements, without a global
cuspidality requirement on a rank-three base change.

The value of every finite local packet character on $a$ is one.
At a nonsplit place this can be checked directly in the
centralizer, without selecting a sign for the ramified
constituent. If $b\ne b^{-1}$, the pair $b,b^{-1}$ contributes
the connected torus $\{\operatorname{diag}(t,t^{-1})\}$,
which contains $-I_2$. If $b=b^{-1}$, the multiplicity-two
space has an orthogonal or symplectic form. The element
$-I_2$ lies respectively in $\mathrm{SO}_2(\mathbf C)$ or
$\mathrm{Sp}_2(\mathbf C)$, both connected. If it merges with
the trivial third character, $a$ has determinant one in
$O_3(\mathbf C)$ and lies in $\mathrm{SO}_3(\mathbf C)$.
All these connected subgroups lie in $\mathrm{SL}_3$.
The exact definition is KMSW Section~0.4.4, p.~27,
repeated in Section~1.6.1, p.~80:
\[
 S_\phi^{\mathrm{rad}}=
       \operatorname{Cent}(\phi,\widehat G_{\mathrm{der}})^\circ,
 \qquad S_\phi^{\natural}=S_\phi/S_\phi^{\mathrm{rad}}.
\]
Here $\widehat G_{\mathrm{der}}=\mathrm{SL}_3$.
Consequently $a$ belongs to $S_{\psi_v}^{\mathrm{rad}}$, so
it is trivial even in the group $S_{\psi_v}^{\natural}$ on
which Theorem~1.6.1(2), p.~80, places the packet characters.
At a split place the
spherical packet character has the same trivial value.

At infinity the ordered weights are $(e,0,-e)$ and the compact
member has coefficient highest weight $(e-1,0,1-e)$.
The real discrete-series packet rule says that for a compact
root between positions $i<j$, the packet character takes product
$(-1)^{i+j}$ on the two corresponding sign generators.
This is Gan--Gross--Prasad, \emph{Restrictions of representations
of classical groups: examples}, Ast\'erisque 346 (2012),
Section~2, p.~115, the compact-root formulas following
Corollary~2.3. In the compact group all roots are compact.
Only the known real local-packet description is used here,
not the restriction conjecture studied elsewhere in that paper.
The two characters from $\Pi$ occupy positions one and three,
so the value on $a$ is $(-1)^{1+3}=1$.
This product rule does not require choosing the individual signs
of the two generators.

The product of all local characters is therefore trivial on
the reduced global component group. KMSW Section~1.7,
p.~89, states that $\epsilon_\psi=1$ for a generic parameter.
The proved generic/pure-inner case of Theorem~5.0.5, p.~206,
therefore places the chosen tensor product in the discrete
automorphic spectrum with multiplicity one. Its finite
$K$-fixed vector and its compact coefficient type give a
Hecke eigenline in $M_{e-1}$.

\paragraph{Independence, counting, and the eigenvalues.}
Langlands's fibre assertion (C) identifies two cuspidal base
changes exactly when their original forms differ by the
quadratic twist. Distinct two-element orbits therefore give
distinct global parameters. KMSW puts these in distinct
direct summands, so the selected eigenlines are independent.
The count of orbits is $(r_e-c_e)/2$ by KT-GL2-CM;
substitution in the three even residue classes modulo six
gives the floor in (KT-E.2).

At a split prime the three unitary roots consist of $1$ and
the two unitary roots of $\pi_f$, so their sum is
$1+a_p(f)/p^e$. Multiplication by $p^e$ in the already checked
compact Hecke normalization gives (KT-E.1).
Finally KT-U3-DIM gives $D(3)=D(5)=1$, while (KT-E.2)
gives one line at $e=4,6$. Hence these full spaces are
endoscopic, and neither contains a stable summand.
\end{proof}

\begin{corollary}[KT-U3-ENDO-SIX]
The unique compact form with Harish--Chandra parameter
$(6,0,-6)$, equivalently coefficient highest weight
$(5,0,-5)$, is the endoscopic lift of the non-CM quadratic-twist
pair in $S_{13}(\Gamma_0(3),\chi_{-3})$.
It has
\[
  7^5T_7=7^6+40250=157899,
  \qquad
  13^5T_{13}=13^6+1284050=6110859.
\]
\end{corollary}
\begin{proof}
The classical space has dimension three and CM dimension one.
The exact classical calculation K17-GL2-MODEL, independently
recomputed in the report's classical section, gives $a_7=40250$
on its two non-CM eigenforms; the CM eigenvalue is $-153502$.
The corresponding non-CM value at thirteen is $1284050$.
KT-U3-ENDO-INJECT and the one-dimensional compact dimension
identify the complete global parameter and give both displayed
values. This argument does not infer endoscopy from agreement
at finitely many primes; it first constructs the global
endoscopic eigenline and then uses the exact dimension.
\end{proof}

\paragraph{The all-weight converse.}
The injection proves the lower bound (KT-E.2), and already
proves the entire stable/endoscopic split at $e=4,6$.
The following section supplies the converse classification,
including the outer-character and other parameter cases.
KT-U3-ENDO-EXACT proves that this family exhausts the endoscopic
spectrum for every even $e\ge2$ and identifies
$D(e-1)-\lfloor(e-1)/3\rfloor$ as the stable dimension.

% END evidence/unitary_audit/endoscopic_results.tex


% BEGIN evidence/unitary_audit/endoscopic_converse.tex
\section{Exhausting the endoscopic spectrum in every even degree}
\label{sec:kt-endoscopic-converse}

This section extends KT-U3-ENDO-INJECT to an equality. Its additional
published inputs are the rank-three spectral classification of Rogawski
and the spherical representation theorem for a special maximal parahoric.
The latter is used to prove the converse local-level assertion, rather
than to assume that a compact fixed vector automatically gives an
unramified base change.

\begin{lemma}[KT-U3-LEVEL-CONVERSE]
Every irreducible smooth representation of $G(\mathbf Q_3)$ with a
vector fixed by the full standard lattice stabilizer has base change to
$\mathrm{GL}_3(E_3)$ with unramified Weil--Deligne parameter and zero
monodromy. For a fixed such local parameter, at most one member of the
packet has a full-stabilizer fixed vector, and its fixed space is one
dimensional.
\end{lemma}
\begin{proof}
Let $K_3^\circ$ be the connected special maximal parahoric inside $K_3$.
Adler--Lansky Sections~4.1--4.2.1, pp.~10--12, identify their reductive
quotients with $\mathrm{SO}_3$ and $O_3$, respectively.
Haines, \emph{On Satake parameters for representations with parahoric
fixed vectors}, author version~4, Section~3, Proposition~3.1, p.~4,
shows that every $K_3^\circ$-spherical irreducible representation is the
spherical constituent of a principal series induced from a weakly
unramified character of the minimal Levi. This is the ordinary
representation-theoretic assertion of that proposition; no conjectural
compatibility between Haines's Satake parameter and local Langlands is
being used.

Let $I(\lambda)$ be such a principal series containing our representation.
Taking invariants under the compact open subgroup $K_3$ is exact in
characteristic zero, so $I(\lambda)^{K_3}\ne0$. The Iwasawa decomposition
$G(\mathbf Q_3)=B(\mathbf Q_3)K_3$ identifies this space with the
functions on one coset satisfying the inducing character relation.
It is therefore one dimensional, and forces $\lambda$ to be trivial on
$B\cap K_3$. In the coordinates of Adler--Lansky (4.4.1), this says
that $\lambda_1$ is unramified on $E_3^\times$ and $\lambda_2=1$ on
$E_3^1$.

Adler--Lansky Section~4.4 now leaves only the irreducible principal
series, case (4.4PS--1), and case (4.4PS--3). Case (4.4PS--2) is
impossible because its restriction to $\mathbf Q_3^\times$ has ramified
quadratic part, whereas the restriction of $\lambda_1$ is unramified.
In (4.4PS--1), the constituents are the trivial representation and the
Steinberg representation. The trivial constituent already has a fixed
vector; exactness and $\dim I(\lambda)^{K_3}=1$ show that the Steinberg
constituent has none. In (4.4PS--3), exactly one of the two constituents
has a fixed vector, again by exactness.

The full list of base changes is Proposition~7.1(i)--(iv), p.~28.
In case (i), it is the irreducible unramified principal series of
$\mathrm{GL}_3(E_3)$. In case (ii), where the unitary-group principal
series is irreducible but the general-linear principal series is
reducible, the displayed lift is induced from a one-dimensional
$\mathrm{GL}_2$ representation and an unramified character; its
parameter is again a direct sum of unramified characters with zero
monodromy. In case (iii), the full-stabilizer constituent is the trivial
representation and its lift is trivial, also with zero monodromy.
Case (iv) is the irreducible unramified general-linear principal series.
These cover every case above. They are the explicit principal-series
cases, whose proof refers to the earlier unramified Proposition~4.1
(Proposition~5.1 in that author preprint); neither the depth-preservation
hypothesis nor restrictions on the size of the residue field are used.

For a fixed base-change parameter the unordered unramified inducing
triple determines $\lambda_1$ up to inversion, hence determines
$I(\lambda)$ up to the Weyl action. The same fixed-space calculation
then shows that at most one constituent in its packet is $K_3$-spherical,
with a one-dimensional fixed space. This includes the reducible
quadratic case, where the two members share one principal series.
\end{proof}

\begin{theorem}[KT-U3-ENDO-EXACT]
For every even $e\ge2$, the full compact space $M_{e-1}$ is a direct
sum of a stable part and the endoscopic lines constructed in
KT-U3-ENDO-INJECT. Its endoscopic dimension is exactly
\[
             E_e=\left\lfloor\frac{e-1}{3}\right\rfloor,
\]
and its stable dimension is
\[
 S_e=D(e-1)-E_e=
 \begin{cases}
 e^3/216-e/3+1,&e\equiv0\pmod6,\\
 e^3/216-2e/9+11/27,&e\equiv2\pmod6,\\
 e^3/216-2e/9+16/27,&e\equiv4\pmod6.
 \end{cases}                                                   \tag{KT-E.3}
\]
Here stable means that the generic parameter is a single cuspidal
$\mathrm{GL}_3/E$ constituent. Each endoscopic parameter contributes
exactly one line, and is uniquely of the form
$1\boxplus\operatorname{BC}_{E/\mathbf Q}(\pi_f)$ for a non-CM
quadratic-twist orbit in $S_{2e+1}(\Gamma_0(3),\epsilon)$.
\end{theorem}
\begin{proof}
\paragraph{The rank-three spectral input and regularity.}
We use Rogawski's classification for unitary groups in three variables,
including their inner forms: the discrete spectrum is partitioned by
discrete global Arthur parameters, with the expected infinitesimal
characters and local base changes. This rank-three assertion is a
proved published input, not the unrestricted general-rank Arthur
formalism. An explicit source locator is Bella\"iche--Chenevier,
\emph{Families of Galois representations and Selmer groups},
Appendix~A.9, p.~244 in the final author version
(pp.~228--229 in the supplied preprint): the paragraph
preceding Remark~A.9.1 states that this spectral description was proved
by Rogawski for $m\le3$, referring to his monograph, Chapters~13--14.
Remark~A.9.1 records the compact regularity obstruction. Section~6.7,
p.~148 (p.~140 in the supplied preprint), defines regularity as strict inequalities between coefficient
highest weights.

In our case the coefficient weights are $(e-1,0,1-e)$, strictly
decreasing for $e\ge2$. One can check the required obstruction directly
in rank three. A nontrivial Arthur $\mathrm{SL}_2$ block must have size
two or three. A size-two block places a pair of weights differing by
one in the infinitesimal character; a size-three block places three
consecutive weights there. The required multiset $(e,0,-e)$ has no
such pair. Therefore every parameter contributing to this space is
generic. The only remaining partitions of the rank are $3$, $2+1$,
and $1+1+1$. For these generic parameters we use the already checked
KMSW Theorem~5.0.5, including its multiplicity formula for this pure
inner form.

At each finite place, the standard compact subgroup forces the
base-change parameter to be unramified with zero monodromy:
this is the usual spherical local correspondence at split and
unramified nonsplit places, and KT-U3-LEVEL-CONVERSE at three.
Consequently every simple global constituent is unramified at every
finite place of $E$. The determinant of the rank-three parameter is
an everywhere unramified Hecke character with trivial infinity type.
Since $E$ has class number one, this determinant is trivial.

\paragraph{All possible one-dimensional constituents.}
For an integer $s$, an everywhere unramified Hecke character of $E$
with infinity character $(z/\bar z)^s$ exists exactly when $3\mid s$,
and is then unique; denote it by $\theta_s$.
Indeed, the class number is one, so its values on ideals are forced
by a generator. The rule is well defined on different generators
exactly when $(u/\bar u)^s=1$ for every $u\in\mu_6$, which is the
condition $3\mid s$. This also proves uniqueness; $\theta_0=1$.
Changing the reciprocity convention inverts all these characters
and leaves the assertions below unchanged.

At an inert rational prime the local value on a uniformizer is one,
because that uniformizer can be chosen rational. At the ramified
prime three, use a generator proportional to $\sqrt{-3}$; the ratio
of it to its conjugate is $-1$, so
\[
                  \theta_s(\varpi_3)=(-1)^s.
\]
In particular $\theta_e$ and $\theta_{-e}$ are trivial at every
finite nonsplit place whenever they exist, since $e$ is even.

\paragraph{The middle character gives precisely the classical lifts.}
Suppose the generic parameter has partition $2+1$ and its
one-dimensional summand occupies the middle infinity weight zero.
That summand is $\theta_0=1$, so write the parameter as $1\boxplus\Pi$.
The cuspidal $\mathrm{GL}_2/E$ representation $\Pi$ has determinant
one. The rank-two duality identity therefore gives
$\Pi^\vee\simeq\Pi$. Since it is conjugate self-dual, also
$\Pi^c\simeq\Pi$. Langlands, \emph{Base change for $\mathrm{GL}(2)$},
Section~2(B),(C),(E), p.~13, and Lemma~11.6(b), p.~142,
then give exactly two cuspidal descents to $\mathrm{GL}_2/\mathbf Q$,
related by twisting with $\epsilon$.

Their real parameter restricts to $(z/\bar z)^{\pm e}$, so their
classical weight is $2e+1$. Their central characters have trivial
base change, hence are either $1$ or $\epsilon$. The odd weight
forces the odd character $\epsilon$.
At every finite prime other than three, local base change is
unramified only if the original inertia and monodromy are trivial:
at a nonsplit unramified prime the inertia groups are identical,
and at a split prime the assertion is immediate.
At three the restriction to $W_{E_3}$ is unramified with zero
monodromy, so inertia factors through the quadratic quotient.
Its determinant is $\epsilon$, hence its characters are exactly
$1$ and $\epsilon$. Thus the conductor is one at three and zero
elsewhere. The descents are precisely the level-three primitive
forms in KT-U3-ENDO-INJECT. They are non-CM because their base
change is cuspidal.

\paragraph{Outer characters and three-character parameters are excluded.}
If a $2+1$ parameter has its one-dimensional constituent at an
outer weight, that constituent is $\theta_e$ or $\theta_{-e}$,
so $3\mid e$. Let $a$ act by $-I_2$ on its cuspidal rank-two
constituent and by $1$ on the character. At every finite nonsplit
place the character is trivial, and the determinant of the
rank-two constituent is one. Its unramified roots are therefore
$b,b^{-1}$. Exactly the connected determinant-one centralizer
argument in KT-U3-ENDO-INJECT puts $a$ in
$\operatorname{Cent}(\psi_v,\mathrm{SL}_3)^\circ$.
This argument does not require a bound on $|b|$.
At split places the same element belongs to the connected
determinant-one centralizer of the semisimple parameter.
Every finite local packet character consequently has value one
on $a$.

At infinity the rank-two constituent now occupies adjacent
positions, either $(1,2)$ or $(2,3)$, in the ordered weights
$(e,0,-e)$. The compact-root rule from Gan--Gross--Prasad,
Section~2, p.~115, gives value $(-1)^{1+2}=-1$ or
$(-1)^{2+3}=-1$. The global product is nontrivial, whereas
$\epsilon_\psi=1$ for a generic parameter. The multiplicity
formula excludes this tensor product.

For partition $1+1+1$, the only possible parameter is
$\theta_e\boxplus1\boxplus\theta_{-e}$, again requiring
$3\mid e$. At every finite nonsplit place all three characters
are trivial. Choose the global sign element acting as $-1$ on
two adjacent infinity positions and as $1$ on the third.
It lies in $\mathrm{SO}_3(\mathbf C)$, hence in the connected
determinant-one centralizer, at every such finite place. At
split places it belongs to the connected determinant-one torus
centralizer. All finite character values are one, and the
compact real value is minus one. The same generic multiplicity
formula excludes this parameter from the compact space.

\paragraph{Multiplicity and dimensions.}
The preceding cases exhaust the rank-three classification.
Every surviving endoscopic parameter is therefore one of the
middle-character parameters already constructed. At the real
place it has only one member on the compact group with this
coefficient type. At finite split and unramified nonsplit places
the spherical member and its fixed line are unique; at three
this is KT-U3-LEVEL-CONVERSE. KMSW gives multiplicity one for
the resulting global tensor product. Thus each non-CM twist
orbit contributes exactly one line and no further endoscopic
lines occur. KT-U3-ENDO-INJECT gives their number as
$\lfloor(e-1)/3\rfloor$. Subtracting this from KT-U3-DIM in
the three even residue classes modulo six gives (KT-E.3).
\end{proof}

\begin{corollary}[KT-U3-SPECTRAL-SPLIT]
Let $p\ne3$ be split in $E$, and write $Q_{e,p}(Y)$ for the
characteristic polynomial of the classical $T_p$ on
$S_{2e+1}(\Gamma_0(3),\epsilon)$. Put
\[
 C_{e,p}(Y)=
 \begin{cases}
 1,&3\nmid e,\\
 Y-(\varpi^{2e}+\bar\varpi^{2e}),&3\mid e,
 \end{cases}
 \qquad \varpi\bar\varpi=p.
\]
The endoscopic polynomial $P^{\rm end}_{e,p}(X)$ is the unique
monic polynomial satisfying
\[
 \bigl(P^{\rm end}_{e,p}(X)\bigr)^2
       =\frac{Q_{e,p}(X-p^e)}{C_{e,p}(X-p^e)}.             \tag{KT-E.4}
\]
It belongs to $\mathbf Z[X]$. If $P^{\rm full}_{e,p}$ is the
characteristic polynomial of the scaled compact operator
$p^{e-1}T_p$, then
\[
 P^{\rm st}_{e,p}=P^{\rm full}_{e,p}/P^{\rm end}_{e,p}
       \in\mathbf Z[X]                                  \tag{KT-E.5}
\]
is the polynomial on the stable part. In particular the exact
uniform character algorithm at $p=7$, together with the classical
algorithm, computes the two factors separately in every even degree.
\end{corollary}
\begin{proof}
KT-GL2-CM gives exactly the displayed CM factor. When it occurs,
$2e$ is divisible by six, so its value is independent of the unit
chosen in a generator $\varpi$ of the split prime ideal.
Every other normalized classical newform occurs in a two-element
quadratic-twist orbit. At a split prime $\epsilon(p)=1$, so the
two eigenvalues $a_p$ coincide. Removing the CM factor therefore
leaves the square of a monic polynomial, and KT-E.1 shifts each
orbit eigenvalue by $p^e$. The quotient in (KT-E.4) is in
$\mathbf Z[X]$; comparing coefficients in its monic square root
puts that root in $\mathbf Q[X]$, while its roots are algebraic
integers. Hence the square root lies in $\mathbf Z[X]$.
KT-U3-ENDO-EXACT identifies it with the entire endoscopic factor.
The full polynomial is integral by KT-U3-SPECTRUM (or the same
scaled-lattice argument at the chosen split prime). Polynomial
division by the monic endoscopic factor gives the stable quotient
in $\mathbf Q[X]$; its integral roots again put it in
$\mathbf Z[X]$.
\end{proof}

\paragraph{Refinement of rank-three source coverage.}
An admissible rank-three source in K17-U3-GENERAL-COVERAGE
is attached to a cuspidal $\mathrm{GL}_3/E$ representation.
Its compact transfer therefore lies in the stable summand.
Consequently KT-TRACE-MULTIPLICITY remains valid with the
rank-three factor $P^{\rm full}_{e,7}(7^{2e}+7^e+1)$ replaced
by $P^{\rm st}_{e,7}(7^{2e}+7^e+1)$. The classical factor
is retained unchanged. This removes the identified endoscopic
factors from that comparison but makes no assertion that the
remaining stable comparison values have only small prime divisors.

\paragraph{Scope of the new source checks.}
The complete arguments above were checked in this run. The
published rank-three classification and the local/global
Langlands correspondences remain imported theorems; no claim
is made to have reconstructed Rogawski's trace formula or
compiled a new formal proof of it. The generic multiplicity
formula, its pure-inner scope, and the exact ramified
principal-series cases were checked at their printed
statements and used proofs as recorded in the check log.

% END evidence/unitary_audit/endoscopic_converse.tex


% BEGIN evidence/unitary_audit/symmetric_square_results.tex
\section{A uniform symmetric-square subspace of the stable spectrum}
\label{sec:kt-symmetric-square}

The stable spectrum itself contains a uniform classical family.  The
construction in this section uses classical forms of weight $e+1$,
whereas the endoscopic construction used weight $2e+1$.  We identify
the resulting subspace and its exact polynomial factor.  The word
``constructed'' below specifies its scope: we do not classify every
possible adjoint origin over $E$.

\begin{lemma}[KT-GL2-LOWER-SCOPE]
The integral basis, Hecke determinant formula and CM classification
in KT-GL2-CM hold for the parameter $n$ in every integer degree
$n\geq1$, with the empty-space convention when $n=1,2$.
In particular, for even $e\geq2$,
\[
 L_e=S_{e+1}(\Gamma_0(3),\epsilon),\qquad
 \dim L_e=\left\lfloor\frac{e-2}{3}\right\rfloor,
 \qquad
 \dim L_e^{\rm CM}={\bf1}_{6\mid e}.                 \tag{KT-S.1}
\]
Twisting by $\epsilon=\chi_{-3}$ pairs the non-CM normalized
eigenforms without fixed points.
\end{lemma}
\begin{proof}
Use the same graded-ring statement on $\Gamma_1(3)$ as in (GL.1).
The weight and character of $Da^{2n-5-3j}b^j$ give an integral
unitriangular basis with $0\leq j<\lfloor(2n-2)/3\rfloor$;
no part of that graded-ring argument requires $n$ to be even.
For $n=1,2$ the permitted index set is empty.  The Fourier
coefficient formula for $T_p$ has factor $\epsilon(p)p^{2n}$
and is unchanged.  The CM proof uses only the level-three
conductor formula and the unit condition $6\mid2n$, so also
does not require even $n$.  Substituting $n=e/2$ proves (KT-S.1).

For completeness, the local twisting argument is likewise
independent of that parity.  Primitivity of the nebentypus
excludes oldforms.  At three the conductor-one Weil--Deligne
parameter has inertia characters $1,\epsilon$ and monodromy
zero, as in KT-U3-ENDO-INJECT.  Twisting interchanges these
two inertia characters and preserves the conductor.  There
is no change of conductor at the other finite primes.
The twist therefore stays in the same primitive level-three
space.  A fixed point is precisely a CM form for $E$, by the
quadratic self-twist criterion already used there.  Hence the
non-CM forms occur in two-element orbits.
\end{proof}

\begin{theorem}[KT-U3-SYM2-INJECT]
For every even $e\geq2$, the non-CM quadratic-twist orbits
in $L_e$ give distinct stable automorphic eigenlines in
$M_{e-1}$.  Their span $M_{e-1}^{\rm sym}$ has dimension
\[
                 H_e=\left\lfloor\frac{e-2}{6}\right\rfloor.
                                                               \tag{KT-S.2}
\]
The corresponding simple global parameter is
\[
 \Pi_f=\operatorname{BC}_{E/\mathbf Q}
             \bigl(\operatorname{Ad}(\pi_f)\bigr),
 \qquad
 \operatorname{Ad}(\pi_f)
       =\operatorname{Sym}^2(\pi_f)\otimes\epsilon^{-1}.
                                                               \tag{KT-S.3}
\]
At every split prime $p\ne3$, the eigenvalue of the scaled
compact operator $p^{e-1}T_p$ is
\[
                   a_p(f)^2-p^e.                              \tag{KT-S.4}
\]
The same value occurs for the opposite minuscule direction.
\end{theorem}
\begin{proof}
\paragraph{The global lift and its local parameters.}
We use unitary normalizations for the automorphic representations.
The Gelbart--Jacquet theorem, \emph{A relation between automorphic
representations of $\mathrm{GL}(2)$ and $\mathrm{GL}(3)$},
Theorem~9.3, printed p.~534, gives the adjoint lift with its local
parameters at every place and makes it cuspidal when the original
representation has no nontrivial character self-twist.  Its
normalization is the determinant-one adjoint representation:
Definition~3.1.3 and Proposition~3.2, pp.~485--486, identify it
at infinity, and Proposition~3.3, in particular (3.3.6), p.~487,
gives the principal-series local formula used at every finite
place here.  The self-twist case is treated in Section~3.7.
A non-CM holomorphic form here has no such self-twist:
any nontrivial self-twist is quadratic by the central character,
and the quadratic self-twist criterion makes the form CM.
Thus $A_f=\operatorname{Ad}(\pi_f)$ is cuspidal on
$\mathrm{GL}_3/\mathbf Q$, self-dual, and has central character one.

Arthur--Clozel, \emph{Simple algebras, base change, and the
advanced theory of the trace formula}, Chapter~3,
Theorem~4.2(a), p.~202, gives
$\Pi_f=\operatorname{BC}_{E/\mathbf Q}(A_f)$.
Theorem~5.1, pp.~212--213, makes this strong base change at
every finite and infinite place, in the sense of Definition~1.2,
p.~199.  The cuspidality criterion says that failure of
cuspidality would give the
nontrivial self-twist $A_f\simeq A_f\otimes\epsilon$.
The central characters of these two representations are $1$
and $\epsilon^3=\epsilon$, respectively, so this is impossible.
Only quadratic base change and its fibers are used here;
no statement about a general tower of solvable descent is needed.

At every prime other than three, the classical representation
is unramified, so its adjoint and base change are unramified
with zero monodromy.  At three its two-dimensional inertia
representation is $1\oplus\epsilon$ and $N=0$.  The adjoint
inertia is therefore $1\oplus\epsilon\oplus\epsilon$.
Restriction to $W_{E_3}$ is trivial, and monodromy stays zero.
The base change is thus unramified at every finite place of $E$.

At the complex place its characters have exponents
$(e,0,-e)$ in the notation $(z/\bar z)^s$.  This follows either
by applying the adjoint representation to the two real
discrete-series weights, or from the ratios of their complex
restrictions $(z/\bar z)^{\pm e/2}$.
The global representation $\Pi_f$ is both self-dual and invariant
under conjugation.  Its conjugate self-dual parity is positive.
To check the latter with precisely the source convention used
above, apply the global seed theorem and its localization in
KMSW Propositions~1.3.1 and~1.3.3.  At infinity the three
distinct characters $(z/\bar z)^e,1,(z/\bar z)^{-e}$
are each of positive conjugate-dual parity.  Their multiplicities
are one, so a negative conjugate-dual form is impossible there.
The global parity must consequently be positive.

\paragraph{The compact member and multiplicity.}
The representation $\Pi_f$ is a simple generic parameter for
the rank-three unitary group.  Its real packet contains the
compact member of coefficient highest weight $(e-1,0,1-e)$.
At every finite place it has a member fixed by the prescribed
standard compact subgroup: the spherical local correspondence
at split and unramified nonsplit primes, and the checked
ramified principal-series cases of Adler--Lansky at three,
give these members.  Each fixed line is one dimensional,
by KT-U3-LEVEL-CONVERSE and the usual spherical uniqueness.

For a simple generic parameter the reduced global component
group is trivial.  Hence there is no remaining sign constraint
in the generic pure-inner multiplicity formula of KMSW
Theorem~5.0.5.  The tensor product of the compact real member
and these finite members occurs with multiplicity one.  It
is stable because its base change is a single cuspidal
$\mathrm{GL}_3/E$ representation.

\paragraph{Precisely the asserted fibers.}
The adjoint is unchanged under twisting a two-dimensional
representation, so $f$ and $f\otimes\epsilon$ give the same
parameter.  Conversely, suppose $\Pi_f\simeq\Pi_g$.
The fibers of prime-cyclic base change, Arthur--Clozel,
Chapter~3, Theorem~3.1, p.~201 (equivalently the fiber assertion
in Theorem~4.2(d), p.~202), give
$A_g\simeq A_f\otimes\epsilon^j$ for $j=0$ or $1$.
The determinant-one normalization excludes $j=1$, exactly
as in the cuspidality argument.  Therefore $A_f\simeq A_g$.

Ramakrishnan, \emph{Modularity of the Rankin--Selberg $L$-series,
and multiplicity one for $\mathrm{SL}(2)$},
Theorem~4.1.2, printed p.~90, says that equality of the adjoint
representations of two cuspidal $\mathrm{GL}_2$ representations
forces one to be a character twist of the other.  Since the
central characters here agree, the twisting character $\eta$
is quadratic.  At each prime other than three both classical
representations are unramified.  Twisting an unramified
two-dimensional Weil--Deligne representation by a ramified
character makes its inertia nontrivial, so $\eta$ is unramified
outside three.  The only quadratic characters of $\mathbf Q$
unramified outside three are $1$ and $\epsilon$, as follows
from the classification of fundamental quadratic discriminants.
Thus $g=f$ or $f\otimes\epsilon$.

Distinct orbits consequently have distinct simple global
parameters and independent eigenlines.  Their number is
\[
 \frac12\left(
        \left\lfloor\frac{e-2}{3}\right\rfloor
          -{\bf1}_{6\mid e}\right)
        =\left\lfloor\frac{e-2}{6}\right\rfloor.
\]
Uniqueness of the finite fixed lines and global multiplicity
one show that their span has exactly this dimension.

\paragraph{Eigenvalue normalization.}
At a split prime let $\alpha,\beta$ be the classical roots,
so $\alpha+\beta=a_p(f)$ and $\alpha\beta=p^e$.
The adjoint roots are $\alpha/\beta,1,\beta/\alpha$.
The normalization already checked from Gross's Satake formula
multiplies their trace by $p^e$ for $p^{e-1}T_p$.  Thus
\[
 p^e\left(\frac\alpha\beta+1+\frac\beta\alpha\right)
       =\alpha^2+\alpha\beta+\beta^2
       =a_p(f)^2-p^e.
\]
The multiset of roots is invariant under inversion; the
opposite minuscule direction therefore has the same scaled
eigenvalue.  This proves (KT-S.4) for both directions.
\end{proof}

\begin{corollary}[KT-U3-SYM2-POLYNOMIAL]
Let $Q_{e/2,p}(Y)$ denote the classical characteristic polynomial
on $L_e$ for a split $p\ne3$.  Set
\[
 C^{\rm low}_{e,p}(Y)=
 \begin{cases}
 1,&6\nmid e,\\
 Y-(\varpi^e+\bar\varpi^e),&6\mid e,
 \end{cases}
 \qquad\varpi\bar\varpi=p.
\]
Let $H_{e,p}(Y)$ be the unique monic square root of
$Q_{e/2,p}(Y)/C^{\rm low}_{e,p}(Y)$.  Then
\[
 H_{e,p}(Y)\in\mathbf Z[Y],\qquad
 P^{\rm sym}_{e,p}(X)
   =\operatorname{Res}_{Y}
          \bigl(H_{e,p}(Y),X+p^e-Y^2\bigr)\in\mathbf Z[X].
                                                               \tag{KT-S.5}
\]
This is precisely the characteristic polynomial on the
constructed symmetric-square subspace, of degree $H_e$.
There is an exact monic integral factorization
\[
 P^{\rm full}_{e,p}
     =P^{\rm end}_{e,p}P^{\rm sym}_{e,p}P^{\rm rem}_{e,p},
 \qquad
 \deg P^{\rm rem}_{e,p}=D(e-1)-E_e-H_e.              \tag{KT-S.6}
\]
In particular this last degree is
\[
 \begin{cases}
 e^3/216-e/2+2,&e\equiv0\pmod6,\\
 e^3/216-7e/18+20/27,&e\equiv2\pmod6,\\
 e^3/216-7e/18+34/27,&e\equiv4\pmod6.
 \end{cases}
\]
\end{corollary}
\begin{proof}
At a split prime the two members of a quadratic-twist pair
have the same $a_p$.  Removing the unique CM factor therefore
leaves a monic square.  Its monic square root is rational by
coefficient recursion and integral because its roots are
algebraic integers.  The resultant is the product over these
twist orbits of $X+p^e-a_p(f)^2$, so KT-U3-SYM2-INJECT proves
its assertion and shows that it divides the stable polynomial.
The remaining quotient is rational, monic, and has only
algebraic integer roots.  It lies in $\mathbf Z[X]$.
The endoscopic factorization was proved in KT-U3-SPECTRAL-SPLIT.
\end{proof}

\paragraph{Exact independent polynomial consistency.}
The program \path{symmetric_square_check.py} freshly computes
the lower classical matrices from the uniform integral basis,
for every integer $1\leq n\leq15$.  It checks (KT-S.5)--(KT-S.6)
for all thirty pairs consisting of an even $2\leq e\leq30$
and $p=7,13$.  No archived factorization lists are used.
After the CM removal it takes an exact monic square root,
constructs the companion matrix of $H_{e,p}$, and computes
the characteristic polynomial of $C_{H_{e,p}}^2-p^eI$.
Exact monic division by this polynomial succeeds in every
stable polynomial saved by the earlier audit.  The full
coefficient arrays, fresh lower classical matrices, and
comparison evaluations are saved in
\path{symmetric_square_check.json}.

At the three new degrees $e=26,28,30$, the symmetric-square
dimension is four, and the remaining stable dimensions are
$72,92,112$, respectively.  For the smaller residual moduli
\[
 697446897311,\qquad323026426255399,\qquad
 236017358500033404991,
\]
the greatest common divisor with the symmetric-square
comparison value at $t_p=1+p^e+p^{2e}$ is the entire modulus,
at both seven and thirteen.  The greatest common divisor with
the remaining stable value is one, at both primes.  Each of
the three larger residual moduli is coprime to both stable
factors at both primes.  These are exact integer gcd assertions;
they require no primality assumption about the listed moduli.

\paragraph{Meaning for a source-aware residual comparison.}
The constructed symmetric-square subspace consists of actual
stable cuspidal parameters.  It cannot be discarded from the
rank-three source comparison on the ground that it is classical
in origin.  If a lower classical form has
\[
               a_p(f)\equiv1+p^e\pmod\lambda
\]
at all split good primes, its two compact minuscule eigenvalues
both satisfy the full Eisenstein trace congruence
\[
       a_p(f)^2-p^e\equiv1+p^e+p^{2e}\pmod\lambda.
\]
More intrinsically, if its residual two-dimensional Galois
representation restricted to $G_E$ has semisimplification
$1\oplus\overline\chi_{\rm cyc}^{\,e}$, then the associated
adjoint restricted to $G_E$ has semisimplification
$\overline\chi_{\rm cyc}^{-e}\oplus1\oplus
 \overline\chi_{\rm cyc}^{e}$.
Thus the all-prime residual condition itself is compatible
with this stable family.  Any elimination argument must
retain its arithmetic source conditions.  The factorization
separates the endoscopic, constructed symmetric-square and
remaining stable obligations, while the original classical
rank-two branch of the comparison is also retained.

\begin{corollary}[KT-U3-SYM2-EISENSTEIN]
Let $e\geq2$ be even and put
\[
      \beta_e=\frac{B_{e+1,\epsilon}}{e+1}.
\]
Suppose $\ell>3$ is a prime for which
$\operatorname{ord}_\ell(\beta_e)>0$.  Suppose also that
either $6\nmid e$ or $\ell\nmid1-2^e$.  Then there is a stable
compact automorphic eigenline in $M_{e-1}^{\rm sym}$ and a
coefficient prime above $\ell$ such that its weight-zero
rank-three Galois representation has residual
semisimplification
\[
 \overline\chi_{\rm cyc}^{-e}\oplus1\oplus
       \overline\chi_{\rm cyc}^{e}
       \quad\hbox{on the whole group }G_E.                     \tag{KT-S.7}
\]
At every split $p\ne3,\ell$, both scaled minuscule eigenvalues
are congruent to $1+p^e+p^{2e}$ at that same coefficient prime.

In particular, this conclusion holds for every prime divisor
$\ell$ of each of the following moduli:
\[
 \begin{array}{c|r}
 e&M_e\\\hline
 26&697446897311\\
 28&323026426255399\\
 30&236017358500033404991.
 \end{array}                                                  \tag{KT-S.8}
\]
No primality assertion about $M_e$ is needed.
\end{corollary}
\begin{proof}
Fretwell--Roberts, \emph{Newform Eisenstein congruences of
local origin}, Theorem~3.6, printed pp.~517--518, applies
with $N=3$, $M=1$, $k=e+1$, $\psi=1$, and
$\phi=\epsilon$.  Its sufficient direction requires only
$k>2$, a coefficient prime away from $6NM$, and positive
valuation of the displayed special-value product.
Here that product is $L(-e,\epsilon)=-\beta_e$.
There is no condition $\ell>k+1$ in this theorem.
It produces a normalized eigenform $f\in L_e$ and one
coefficient prime $\lambda\mid\ell$ with
\[
     a_q(f)\equiv1+\epsilon(q)q^e\pmod\lambda
                 \qquad(q\ne3).
\]
The paragraph immediately after its proof makes the form new
because the nebentypus has full conductor three.

If $6\nmid e$, KT-GL2-LOWER-SCOPE shows that no CM forms
occur in this space.  Otherwise, at the inert prime two a CM
form has $a_2=0$, whereas the constructed form has
$a_2\equiv1-2^e\not\equiv0\pmod\lambda$.
Thus $f$ is non-CM in either case.

The classical Galois representation of $f$, its determinant,
the good-prime congruences, and Chebotarev give
\[
     \overline\rho_{f,\lambda}^{\rm ss}
       \simeq1\oplus\epsilon\overline\chi_{\rm cyc}^{e}.
\]
Restriction to $G_E$ removes $\epsilon$.
Taking the adjoint gives the three displayed composition
factors in (KT-S.7).  Equivalently, use the filtration of the
endomorphisms of a two-dimensional triangular representation
and remove its scalar line; $\ell>3$ makes this removal
unambiguous.  KT-U3-SYM2-INJECT constructs the corresponding
stable compact line.  Compatibility at the good places and
Chebotarev identify its normalized rank-three system with this
adjoint restricted to $G_E$.  Formula (KT-S.4) proves the
claim about both minuscule directions at every split good
prime.

For the numerical cases, the exact generalized Bernoulli
values are
\[
 \begin{aligned}
 \beta_{26}&=\frac{80907883885186609178}{81},\\
 \beta_{28}&=-\frac{516454004244279410402}{3},\\
 \beta_{30}&=\frac{102431533589014497766094}{3}.
 \end{aligned}                                                \tag{KT-S.9}
\]
For each row of (KT-S.8), exact integer arithmetic gives
\[
 M_e\mid\operatorname{num}(\beta_e),\qquad
 \gcd\bigl(M_e,6\operatorname{den}(\beta_e)(1-2^e)\bigr)=1.
\]
Every prime divisor of $M_e$ therefore satisfies all the
preceding hypotheses.  The values in (KT-S.9) can be computed
directly from
\[
 \beta_e=\frac{2\,3^e}{e+1}
                 B_{e+1}\!\left(\frac13\right),
\]
which follows from the defining generalized Bernoulli sum
and $B_{e+1}(1-x)=-B_{e+1}(x)$.
The exact rational values and gcds are recorded in
\path{symmetric_square_check.json}, with independent
standard-library rational-arithmetic replay in
\path{evidence/main_gate/lower_bernoulli_certificate.py}
and its JSON output.  The source-audit sibling independently
read and replayed that certificate.
\end{proof}

\paragraph{Consequence for the higher-degree obstruction.}
The smaller moduli at $e=26,28,30$ are therefore supported by
actual stable automorphic systems satisfying the whole
rank-three residual condition, not just by isolated
determinant coincidences.  No finite or infinite collection
of split spherical Hecke generators can exclude those
systems from the full stable space.  The exact coprimality
of the remaining stable factor locates these particular
obstructions inside the constructed symmetric-square
subspace.  The three larger moduli remain in the original
classical comparison branch.  A successful arithmetic
vanishing proof must use the more specific deformation or
Selmer conditions; the existence of these automorphic
congruences by itself asserts neither realization of an
admissible deformation source nor nonvanishing of a
$K$-group.

\paragraph{Source scope.}
The new construction combines the cited Gelbart--Jacquet,
prime-cyclic Arthur--Clozel, and Ramakrishnan theorems with
the local and pure-inner packet arguments checked in the
preceding sections.  Those published lifting theorems are
imported inputs.  The polynomial computations are consistency
checks of their application and are not a proof of global
functoriality.  No classification of all adjoint origins over
$E$, or vanishing of the symmetric-square arithmetic source
branch, is asserted.

% END evidence/unitary_audit/symmetric_square_results.tex


% BEGIN evidence/unitary_audit/dual_hecke_results.tex
\section{Testing both minuscule directions}
\label{sec:kt-dual-hecke}

A trace comparison in rank three does not by itself impose the
inverse-trace condition.  The stored harmonic model permits both
conditions to be tested at the same coefficient prime.  The following
identity specifies that test without identifying a coefficient prime
with its complex conjugate.

\begin{lemma}[KT-U3-DUAL-MATRIX]
Let $C$ interchange the two sets of polynomial variables in the
harmonic model, and choose the saved basis with
$C=\operatorname{diag}(s_i)$, $s_i\in\{1,-1\}$.
Write $M_p=M_{p,0}+\omega M_{p,1}$ for the scaled negative
minuscule operator.  Put $\delta=1+2\omega$, $\delta^2=-3$,
and let $B$ be diagonal with entry $1$ on the plus basis
vectors and $\delta$ on the minus basis vectors.  In the
rational coefficient basis used in the data, the two scaled
minuscule matrices are
\[
       R_p=B^{-1}M_pB,\qquad R_p^{\rm opp}=CR_pC.
                                                               \tag{KT-D.1}
\]
At any modulus for which all recorded matrix denominators and
$6$ are units, the simultaneous kernel at a rational target $t_p$
is the direct sum of its two $C$-parity kernels.  On a vector of
one parity, it is enough to impose $R_p-t_pI$ for each $p$.
\end{lemma}
\begin{proof}
The full harmonic action is
$\rho(g)f(z,w)=f(g^{-1}z,\overline g^{-1}w)$.
Consequently $C\rho(g)C=\rho(\overline g)$.
Conjugation interchanges the positive and negative minuscule
coset sets, so the opposite matrix is $CM_pC$ on this same
space.  Since the saved basis is rational, it also equals
$\overline{M_p}$, giving
\[
 CM_{p,0}C=M_{p,0}-M_{p,1},\qquad
 CM_{p,1}C=-M_{p,1}.
\]
For equal-parity indices the $\omega$ coefficient is zero.
For unequal-parity indices it is twice the rational coefficient,
so that entry of $M_p$ is $\delta(M_{p,0})_{ij}$.
Conjugating by $B$ makes such an entry
$-3(M_{p,0})_{ij}$ in a plus row and
$(M_{p,0})_{ij}$ in a minus row.  This is exactly the saved
rational matrix.  Since $B$ commutes with $C$, it also proves
the second formula in (KT-D.1).

These are equalities before reduction.  At a coefficient
prime away from three, $B$ is invertible; both rational
matrices reduce in one and the same residue field.  No action
of complex conjugation on that prime is assumed.
The common kernel of $R_p-t_pI$ and $CR_pC-t_pI$, over all
chosen $p$, is invariant under $C$.  Since two is a unit,
the projectors $(1\pm C)/2$ split it into its parity parts.
For $Cv=sv$, the identity
\[
        (CR_pC-t_pI)v=sC(R_p-t_pI)v
\]
shows that the opposite condition is automatic after the
first one.  Restricting the columns of the stacked matrices
$R_p-t_pI$ to each parity therefore computes the exact
simultaneous kernel.
\end{proof}

\paragraph{Exact computation at the three smaller residual moduli.}
The program \path{check_dual_hecke.py} verifies all coefficient
identities above against the complete stored harmonic matrices,
and checks every denominator for invertibility modulo the
specified modulus.  Its elimination divides only by units.
The target at $p$ is $t_p=1+p^e+p^{2e}$.
For $p=7$ alone, for $p=13$ alone, and for the two primes
together, the parity ranks are the same:
\[
\begin{array}{c|r|c|c}
e&\text{modulus}&\text{plus rank/dimension}&
                         \text{minus rank/dimension}\\\hline
26&697446897311&44/44&39/40\\
28&323026426255399&55/55&49/50\\
30&236017358500033404991&65/65&59/60
\end{array}
\]
Every reduced matrix has zero residual rows after the listed
unit pivots.  In each minus block the program constructs a
kernel vector with a coordinate equal to one and verifies it
directly in all four full matrices.  Thus the vector remains
nonzero at every prime divisor of the listed modulus.
The complete pivot data and primitive kernel vectors are in
\path{dual_hecke_check.json}.  No primality test is needed
for these assertions.

The result is a failure of this finite Hecke exclusion test:
adding the opposite minuscule generators does not eliminate
these three residual obstructions.  It does not assert that
every resulting finite-prime congruence is an admissible
deformation-theoretic source.  KT-U3-SYM2-INJECT and the exact
factor checks locate the smaller-modulus spectral support in
the constructed symmetric-square family; both directions in
that family necessarily have equal eigenvalues.

% END evidence/unitary_audit/dual_hecke_results.tex


% BEGIN evidence/unitary_audit/quartic_refinement.tex
\section{A sharper explicit quartic size bound}
\label{sec:kt-quartic-refinement}

The exact stable/endoscopic decomposition improves the coefficient
in KT-QUARTIC-BOUND.  This refinement retains the full stable
rank-three space, including the constructed symmetric-square
subspace, and the full classical rank-two space.  It uses the
already proved stable-only source coverage and the same arithmetic
comparison as the preceding bound.

\begin{theorem}[KT-QUARTIC-QUARTER]
Under the report-wide carried-source convention of
KT-TRACE-MULTIPLICITY and KT-QUARTIC-BOUND, every even
$e\geq2$ satisfies
\[
             \log|K_{2e}(\mathbf Z)|<\frac14e^4.          \tag{KT-Q.1}
\]
More precisely, for every even $e\geq4$ the same argument gives
\[
             \log|K_{2e}(\mathbf Z)|<\frac{239}{960}e^4.
                                                               \tag{KT-Q.2}
\]
This is a new explicit coefficient.  It is not an evaluation
of the existentially selected constant in the archived
VK-EXP7 theorem.
\end{theorem}
\begin{proof}
\paragraph{The comparison integer and its exact degree.}
Set
\[
 N^{\rm st}_{e,7}
  =P^{\rm st}_{e,7}(7^{2e}+7^e+1)\,
                    Q_{e,7}(7^{2e}+1),
\]
where the first polynomial is the full stable compact factor
and the second is the full classical characteristic polynomial
at weight $2e+1$.  The stable-only refinement of
KT-TRACE-MULTIPLICITY and the arithmetic proof of
KT-QUARTIC-BOUND give
\[
             |K_{2e}(\mathbf Z)|\mid(N^{\rm st}_{e,7})^2.
                                                               \tag{KT-Q.3}
\]
The number of complex linear factors in this product is
\[
 s_e=D(e-1)-\left\lfloor\frac{e-1}{3}\right\rfloor
                  +\left\lfloor\frac{2e-2}{3}\right\rfloor.
\]
The dimension formulas give its exact values in the three
even residue classes:
\[
 s_e=\begin{cases}
 e^3/216+e/3,&e\equiv0\pmod6,\\
 e^3/216+4e/9-25/27,&e\equiv2\pmod6,\\
 e^3/216+4e/9-2/27,&e\equiv4\pmod6.
 \end{cases}                                                  \tag{KT-Q.4}
\]
In particular, $s_2=0$ and $s_4=2$.
For $e\geq6$, each line of (KT-Q.4) gives
\[
 \frac{s_e}{e^3}\leq\frac1{216}+\frac4{9e^2}
       \leq\frac1{216}+\frac1{81}
       =\frac{11}{648}<\frac1{32}.
\]
Since $s_4=2=4^3/32$, this proves the uniform inequality
\[
                       s_e\leq\frac{e^3}{32}
                   \qquad(e\geq4\hbox{ even}).                \tag{KT-Q.5}
\]

\paragraph{A common bound for every linear factor.}
The compact operator is a sum of $57$ operators unitary for
an invariant positive-definite Hermitian norm on its
coefficient representation.  Therefore every eigenvalue
$\alpha$ of its scaled operator obeys
$|\alpha|\leq57\,7^{e-1}$.  Restriction to the stable
summand preserves this bound.  The classical bound already
used above is $|a_7|\leq2\,7^e$.  Consequently every linear
factor in the absolute value of $N^{\rm st}_{e,7}$ is at most
\[
 7^{2e}+\frac{64}{7}\,7^e+1
   =7^{2e}\left(1+\frac{64}{7^{e+1}}+\frac1{7^{2e}}\right).
\]
For $e\geq2$ the last parenthesis is bounded by
\[
       1+\frac{64}{343}+\frac1{2401}
       =1+\frac{449}{2401}<\frac43.
\]
The targets also strictly exceed the corresponding eigenvalue
bounds, as in KT-TRACE-MULTIPLICITY, so the comparison integer
is nonzero.  Multiplying the bounds and taking logarithms gives
\[
      \log|N^{\rm st}_{e,7}|
       \leq s_e\left(2e\log7+\log\frac43\right).       \tag{KT-Q.6}
\]
This estimate uses the universal compact operator norm, with
no additional Ramanujan assertion about its stable spectrum.

\paragraph{The numerical coefficient and the two small cases.}
The elementary inequalities
$\log(4/3)<1/3$ and $\log7<39/20$ suffice.  The second can
be checked without decimal approximations: the positive
exponential series gives
\[
 \exp(39/20)>\sum_{j=0}^{8}\frac{(39/20)^j}{j!}
    =7+\frac{3130848254169}{114688000000000}>7.
\]
For every even $e\geq4$, combine (KT-Q.3), (KT-Q.5), and
(KT-Q.6), and use $1/(3e)\leq1/12$:
\[
 \begin{aligned}
 \log|K_{2e}(\mathbf Z)|
 &\leq2s_e\left(2e\log7+\log\frac43\right)\\
 &<\frac{e^4}{16}\left(\frac{39}{10}+\frac1{3e}\right)\\
 &\leq\frac{e^4}{16}\left(\frac{39}{10}+\frac1{12}\right)
   =\frac{239}{960}e^4<\frac14e^4.
 \end{aligned}
\]
This includes $e=4$, whose exact degree was treated above.
For $e=2$, both polynomial spaces are empty, so
$N^{\rm st}_{2,7}=1$.  Equation (KT-Q.3) gives
$|K_4(\mathbf Z)|=1$, and its logarithm is zero, proving
(KT-Q.1) in this remaining case.
\end{proof}

\paragraph{Dependencies retained.}
The stable-only source injection is the refinement following
KT-U3-SPECTRAL-SPLIT.  The arithmetic divisibility in
(KT-Q.3) retains the closed OneChar extraction, the all-even-degree
GD statement, the inherited exceptional-prime argument, the
carried small-prime Vandiver input, and the separate
Rognes--Weibel two-primary input already stated in
KT-QUARTIC-BOUND.  This section changes the dimension and
archimedean estimates; it adds no new arithmetic vanishing
claim and does not omit the stable symmetric-square branch.

% END evidence/unitary_audit/quartic_refinement.tex


% BEGIN evidence/higher4/higher4_section.tex
\section{Complete residual-ring certificates in the next weights}
\label{sec:kt-higher4}

This section uses K17-MINUS7-GENERAL-SOURCE and
K17-U3-RESIDUAL-RING as carried results.  Their entire written
proofs and the complete verifier were read again for the present
audit.  The direct source proof KT-ONECHAR-CLOSED supplies the
\path{OneCharU_s} input.  All conclusions below use actual complete
reports; progress logs and producer output alone do not certify a row.

\subsection{The native determinant preserves the exact criterion}

Put
\[
 B_M=(\mathbf Z/M\mathbf Z)[w]/(w^2-w+2),\qquad
 \overline{a+bw}=(a+b)-bw.
\]
The norm is $N(a+bw)=a^2+ab+2b^2$, and
\[
 x\in B_M^\times\quad\Longleftrightarrow\quad
                  \gcd(N(x),M)=1.                 \tag{KT-H.1}
\]
No assumption that $M$ is prime or squarefree is involved.

\begin{lemma}[KT-NATIVE-DETERMINANT]
Whenever the included native determinant routine terminates
successfully, its output is the exact determinant over $B_M$.
A failure to find an individual unit pivot is inconclusive.
Replacing the two determinant routines in the full residual-ring
verifier by this routine preserves every successful unit conclusion.
\end{lemma}
\begin{proof}
The source \path{ring_determinant.cpp}, lines 12--22,
implements multiplication using $w^2=w-2$, conjugation, the norm,
and inversion by $\bar xN(x)^{-1}$.  Formula (KT-H.1) justifies
every inverse.  The loop at lines 31--42 uses row interchanges and
subtracts pivot-row multiples.  A row interchange changes the
accumulated sign.  Earlier column entries in both swapped rows
are zero, so swapping only the active suffix is equivalent to
the full row interchange.

Inductively the original determinant equals the signed product
of previous pivots times the determinant of the active block.
For a unit pivot, the usual elementary block elimination proves
the next instance of this identity over any commutative ring.
The final entry is multiplied without inversion.  Thus a returned
value is the exact determinant, even if that last entry is a
nonunit.  If no unit pivot is available, the program exits without
a determinant certificate.  The distinction is necessary:
$\left(\begin{smallmatrix}2&3\\3&2\end{smallmatrix}\right)$
has determinant one over $\mathbf Z/6\mathbf Z$ but no unit
entry in its first column.

The wrapper requires successful exit, the expected number of
pivot inversions, residue representatives in the correct range,
and norm-unit determinants.  All full-model assertions precede
the substituted routine and remain unchanged.  When a producer
report is present, the comparison point, dimensions, and all
three determinant values must also agree with it.  Otherwise the
report is explicitly a new complete exact computation, rather
than an independently repeated numerical result.
\end{proof}

The independent controls in \path{audit_native_controls.py}
compare both the existing executable and a fresh compilation with
literal Leibniz determinants.  All 34 matrices are saved with
their seed and outputs.  There were 28 exact returns and six
correctly inconclusive pivot searches; every returned determinant
agreed.  Three malformed inputs were rejected.  These controls
test ring arithmetic and are not substitutes for large-model
computations.

\paragraph{Completeness of the coefficient space.}
The preserved verifier reconstructs every signed-permutation
orbit in the source and target bidegrees, checks every saved
column against the full contraction, rejects duplicate sparse
indices, and verifies a distinct-coordinate identity minor.
The split-contraction argument in K17-U3-RESIDUAL-RING, followed
by reduction at every maximal ideal and Nakayama, proves that
these columns form the entire harmonic kernel over $B_M$.
All six norm-two and seven norm-four frames are independently
enumerated and compared.  Every identifying-point evaluation
for every frame is recomputed.  Parallel completion order is
removed by restoring the original row indices and rejecting
missing or repeated rows.  The two additional stored points are
not checked by this residual-ring verifier.

For $k=e-1$ and $E_e=1+2^e+2^{2e}$, the matrix thus obtained is
\[
 S_e=E_e^2A_0-E_eB_2-B_4.
\]
The denominator cancellation in the carried proof identifies it
with the evaluation matrix of
\[
 D_e=E_e^2I-E_e2^kS_{11}-2^{2k}\operatorname{inc}S_{21}.
                                                        \tag{KT-H.2}
\]
Since $\det A_0$ is a unit, $\det S_e/\det A_0=\det D_e$.
When $\gcd(M,336(2e-1)!)=\gcd(M,E_e)=1$, a unit value in
(KT-H.2) proves the full two-class comparison is a unit at every
coefficient place over every prime divisor of $M$.  The inherited
integral conjugation involution also covers the opposite
minuscule direction.

\subsection{Completed rows and the arithmetic conclusion}

Use the following exact residual integers:
\[
\begin{array}{c|r}
 M_{26,1}&697446897311\\
 M_{26,2}&1346178024590512195284390892511906772585800385915601\\
 M_{28,1}&323026426255399\\
 M_{28,2}&1850161514202910077636429426241454885347275557889859421603\\
 M_{30,1}&236017358500033404991\\
 M_{30,2}&17825126775144482230877986030193435571179905679430331.
\end{array}                                             \tag{KT-H.3}
\]

\begin{proposition}[KT-HIGHER4-CERTIFIED]
The full alternate-field unitary comparisons are units at every
coefficient place above every prime divisor of the six
integers in (KT-H.3), in their corresponding weights.
The full classical cubic comparisons are units for all six
integers in (KT-H.3).
\end{proposition}
\begin{proof}
The completed unitary reports give the following residues;
all displayed elements have zero $w$ coefficient:
\[
\begin{array}{c|c|r}
(e,i)&\text{quantity}&\text{residue}\\\hline
(26,1)&\det A_0&65303065346\\
(26,1)&\det D_e&243170653113\\
(26,2)&\det A_0&1151061946436479205314087272743639226699780220614009\\
(26,2)&\det D_e&268141211335182978418440922413180866515460840367849\\
(28,1)&\det A_0&222297174863416\\
(28,1)&\det D_e&110403897672730\\
(28,2)&\det A_0&52319251349733998033648500176324776672876486412159569342\\
(28,2)&\det D_e&1469637909832672408423505694160941756010299985599249215332\\
(30,1)&\det A_0&129139130004479025515\\
(30,1)&\det D_e&166602791931417452321\\
(30,2)&\det A_0&11603147216938395361272501808831427722560693296986105\\
(30,2)&\det D_e&8004246273119905385907729021266082827377552171226936.
\end{array}                                             \tag{KT-H.4}
\]
The full invariant source, target and kernel dimensions are
$(5197,4465,732)$ for $e=26$, $(6942,6027,915)$ for $e=28$,
and $(9094,7969,1125)$ for $e=30$.
The reports recompute respectively all 732, 915 or 1125 identifying
rows.  Both norm gcds, the denominator gcd, and the scalar-pivot
gcd are one.  The full-model argument above proves the unitary
assertions.  The report names are
\path{e26_m697446897311/verify_ring_schur.json} and the
\path{verify_ring_schur_parallel.json} reports in the other
two corresponding case directories, together with
\path{ktower_complete_ring.json} in the large $e=28$ case directory
and both $e=30$ case directories.  These last three are new
complete exact computations; no prior producer determinant is
claimed to have been independently replayed.

For
\[
 J_e=\det\bigl((1+2^{2e})^3I-T_2^3\mid
                         S_{2e+1}(\Gamma_1(7))\bigr),
\]
the completed full classical calculations give
\[
\begin{array}{c|r}
(e,i)&J_e\bmod M_{e,i}\\\hline
(26,1)&564487428902\\
(26,2)&540110375711722242001383842413974682809788702229075\\
(28,1)&291047394117369\\
(28,2)&1361695747498278489585646473526121447210917280902628497429\\
(30,1)&106624733533355298729\\
(30,2)&11520444014513425541585786124609238893503523917343666.
\end{array}                                             \tag{KT-H.5}
\]
Their full dimensions are respectively 101, 109 and 117.
Every residue and every inverted pivot is coprime to its modulus.
The carried full integral cusp-model proof includes the diamond
correction in $T_2$ and all cubic twists.  It therefore gives the
asserted classical units.  These files are
\path{gl2_e26_residual_certificate.json},
\path{gl2_e28_residual_certificate.json}, and
\path{gl2_e30_residual_certificate.json}.
\end{proof}

\begin{theorem}[KT-K52-K56-K60-VANISH]
\label{thm:kt-k52-k56-k60-vanish}
Under the stipulated archive-theorem convention and the carried
arithmetic comparisons and small-prime results,
\[
       K_{52}(\mathbf Z)=K_{56}(\mathbf Z)=K_{60}(\mathbf Z)=0.
\]
\end{theorem}
\begin{proof}
The full initial-field $E=\mathbf Q(\sqrt{-3})$ computation at
$e=26$, at the test primes seven and thirteen, has completed
carried original, sparse, and integer-Faddeev replay reports.
Their common hash is
\path{d68b88a445f2c3726ce84aee3465c290a03e184b2eed8f521d71c010642c2515}.
Fresh reconciliation verifies their complete data wrapper, its
entire degree object, and its neighbor hash against the present
support calculation.  This is a carried complete computation
with a new hash reconciliation, not a new execution of those
three replays.

The corresponding full $e=28$ integer-Faddeev replay has also
completed, in \path{ktower_full_replay_e28.json}.  Its report hash is
\path{93a510c1136be190c45333eac43e62598b7b275d2bb28f8466489dbb190b4c84};
the current data hash is
\path{e5f4785307ebaec12f692c36e08e84cca2e9a8264204d11c5f9c879c3697df85}.
Fresh reconciliation checks that data hash, the full degree
object, and the common neighbor hash.  The report verifies every
neighbor evaluation, exact characteristic polynomials, and both
extra points for its full 105-dimensional harmonic space.

The full $e=30$ integer-Faddeev replay is complete in
\path{ktower_full_replay_e30.json}, with report hash
\path{b1ddcc47506e3fa0613a39f83f82eb14a8bbc37a0d3960bb88bfc7f7bae8255a}
and data hash
\path{629d1a34dbeaae989d471f10178c4813d3b056b98530f7a9474b8fcce01d1cc0}.
Its complete input object and neighbor hash have also been
reconciled.  The invariant source and target dimensions are
1057 and 932, and the entire 125-dimensional harmonic space
passes all neighbor evaluations, exact characteristic
polynomials, and both extra points at each test prime.

Fix $e\in\{26,28,30\}$.
The exact support certificate and its independent check prove
that every prime at least $2^{31}$ dividing both full
initial-field comparison integers divides $M_{e,1}$ or
$M_{e,2}$.  Repeated exact gcd division removes the entire
remaining integer.  The other factors are proved to have
prime divisors below the cutoff.  No factorization or primality
claim about either large residual integer is needed.

Suppose a prime $\ell$ divides the order of $K_{2e}(\mathbf Z)$.
The carried Rognes--Weibel input excludes $\ell=2$; the
stipulated Vandiver verification and Kurihara equivalence
exclude all odd $\ell<2^{31}$.  For the remaining primes the
source datum over $\mathbf Q(\sqrt{-3})$ holds, and $\ell$ differs
from both test primes.  The closed OneChar theorem and full
initial-field coverage force $\ell$ into the preceding residual
support.  Both alternate unitary comparisons and both classical
cubic comparisons in this weight are units by
KT-HIGHER4-CERTIFIED.  Since $\ell>2e+1$, the hypotheses of
K17-MINUS7-GENERAL-SOURCE hold.  That theorem makes the strict
negative Tate Selmer group zero, and VK-PT makes the same
$\ell$-primary $K$-group zero, a contradiction.  Every primary
part vanishes; finiteness proves the result.
\end{proof}

\paragraph{Completed-report boundary.}
At the updated reconciliation cutoff, 16:26 UTC on
21 November 2026, all six alternate unitary comparisons,
all six classical comparisons, and all three initial-field
full computations have complete reports.  The exact hashes
and unit checks are recorded in
\path{certificate_audit_snapshot.json}.  This reconciliation
does not claim another execution of the large-model computations
by the reviewer.  The earlier identifiers KT-K52-VANISH and
KT-K52-K56-VANISH remain valid restrictions of the present theorem.

\subsection{The remaining uniform problem}

The source argument is uniform in $e$.  In the full proof of
K17-MINUS7-GENERAL-SOURCE the inequality $\ell>2e+1$ is used
only to imply $(\ell-1)\nmid2e$.  The same proof therefore works
for every even $e\ge2$ and prime $\ell\ge5$ with $\ell\ne7$
and $(\ell-1)\nmid2e$, provided its two full determinant-unit
hypotheses hold.  The local slope argument imposes no additional
inequality between $e$ and $\ell$.

Let $R_e$ record the remaining prime support of the initial-field
gcd after independently known vanishing primes are removed.
Let $F_e$ be the full integral two-class unitary comparison over
$\mathbf Q(\sqrt{-7})$, and let $J_e$ be the full classical cubic
comparison.  A sufficient uniform assertion is
\[
               \gcd(R_e,N(F_e)J_e)=1               \tag{KT-H.6}
\]
in every weight, with exceptional source characters handled
separately.  No such all-weight coprimality theorem has been
proved here.  A failure of (KT-H.6) would identify a common
automorphic congruence, without by itself producing a Selmer
class.  The quartic order bound controls size in each weight
and supplies no uniform bound on the prime support.

The modular harmonic and Schur shortcuts impose two further
\emph{computational} conditions: invertibility of the contraction
denominators and of $E_e$.  They can be bypassed using a saturated
integral characteristic-zero harmonic lattice and the full
two-class determinant.  They are not intrinsic obstructions
to the global unit criterion.

\begin{lemma}[KT-SCHUR-PIVOT-SUPPORT]
For even positive $e$, the union of the prime divisors of
$E_e=1+2^e+2^{2e}$ is unbounded.  Every prime $\ell\ne3$
dividing $E_e$ satisfies $(\ell-1)\nmid2e$.
\end{lemma}
\begin{proof}
Given $B\ge7$, choose a positive even $e$ divisible by $p-1$
for every prime $5\le p\le B$.  Then $E_e\equiv3\pmod p$
for each such $p$.  It is odd and has exactly one factor three:
if $x=2^e\equiv1\pmod3$, then $x^2+x+1\equiv3\pmod9$.
Since $E_e>3$, it has a prime divisor greater than $B$.
If a prime $\ell\ne3$ dividing $E_e$ also satisfied
$(\ell-1)\mid2e$, then $x=2^e$ would satisfy $x^2=1$ and
$x^2+x+1=0$ modulo $\ell$.  Hence $x=-2$ and $4=1$,
forcing $\ell=3$, a contradiction.
\end{proof}

Thus a fixed small-prime cutoff alone cannot justify the scalar
Schur pivot in every weight, even in the generic source range.
This lemma does not assert that its pivot primes divide $R_e$;
the unresolved global assertion remains (KT-H.6).

% END evidence/higher4/higher4_section.tex


% BEGIN evidence/higher4/native_full_rows_section.tex
\subsection{Exact native evaluation of every identifying row}
\label{sec:kt-native-full-rows}

The new implementation consists of
\path{ring_full_evaluator.cpp} and
\path{ktower_complete_native_ring.py}.  The original
\path{verify_ring_schur_parallel.py} is not edited.  Its complete
structural prefix is executed before any native evaluation; it
reconstructs the thirteen neighbor frames, the full signed-permutation
source and target spaces, the contraction, the harmonic basis and
its identity minor, and compares every integer of the native input
\path{model.txt} with the verified model.  We use the coefficient-space
and Schur-complement conclusions of K17-U3-RESIDUAL-RING.

\begin{proposition}[KT-NATIVE-FULL-ROWS]
Suppose the new wrapper returns a complete unit certificate for a
model with $k=e-1$, modulus $M>1$, and positive harmonic dimension
$d$.  Then all $d$ identifying rows and both determinants in that
certificate are exact over
\[
 B_M=(\mathbf Z/M\mathbf Z)[w]/(w^2-w+2).
\]
The certificate proves the same full two-class comparison unit
criterion as K17-U3-RESIDUAL-RING, at every coefficient place above
every prime divisor of $M$.  Primality and squarefreeness of $M$
are not hypotheses.  The conclusion requires successful completion
of the structural prefix and the full native calculation, not a
sample of rows or a progress log.
\end{proposition}
\begin{proof}
Write $o(P)$ for the row of all signed-permutation orbit sums
evaluated at a point $P=(z,w)$, and let $C$ be the checked sparse
matrix whose columns are the harmonic basis.  Thus the harmonic
evaluation row is $o(P)C$.  For a frame $F$ the transformed coordinates
used by both implementations are
\[
 z'_j=\sum_i\overline{F_{ij}}z_i,
 \qquad w'_j=\sum_iF_{ij}w_i.
\]
The stored convention lists columns first; the native loops use
that same convention.  The two verified frame lists have six and
seven members.  Put $E=1+2^e+2^{2e}$ modulo $M$.  For each identifying
point, the native routine computes the two rows
\[
 \begin{split}
 A(P)&=o(P)C,\\
 S(P)&=\left(E^2o(P)-E\sum_{F\in\mathcal F_2}o(FP)
                         -\sum_{F\in\mathcal F_4}o(FP)\right)C.
 \end{split}                                                    \tag{KT-NF.1}
\]
Linearity of multiplication by $C$ shows that (KT-NF.1) equals
the original verifier's expression, which projects separately
after each frame evaluation.  Each orbit sum uses every distinct
permutation of its sorted triple of exponent pairs exactly once.
Precomputing powers and pair products only reuses equal ring
elements; it neither deletes a monomial nor changes its multiplicity.
All operations use arbitrary-precision integers reduced modulo $M$.
Rows are written to their original indices in disjoint storage;
each worker has its own temporary arithmetic objects.  Every one
of the $d$ rows is evaluated using all thirteen frames.  The native
completion counter and the wrapper require exactly $d$ rows.

Consequently the matrices are precisely the original $A_0$ and
$S_e=E^2A_0-EB_2-B_4$.  The preserved structural proof identifies
the columns with the whole harmonic kernel.  Under the checked
conditions
\[
 \gcd\bigl(M,336(2e-1)!\bigr)=\gcd(M,E)=1,
\]
the carried denominator cancellation identifies $S_e=A_0D_e$,
where
\[
 D_e=E^2I-E2^kS_{11}-2^{2k}\operatorname{inc}S_{21}.
\]

For $x=a+bw$, multiplication and inversion use
\[
 (a+bw)(c+dw)=(ac-2bd)+(ad+bc+bd)w,
 \quad \bar x=(a+b)-bw,
 \quad N(x)=a^2+ab+2b^2.
\]
Here $x$ is a unit exactly when $\gcd(N(x),M)=1$: the forward
implication follows by taking the determinant of multiplication
by $x$ on the free rank-two module $B_M$, and the reverse follows
from $x\bar x=N(x)$.  The scalar determinant routine works in
$\mathbf Z/M\mathbf Z$ because every entry of $A_0$ is checked
to have zero $w$ coordinate.  The second routine works in $B_M$.
Both use only row interchanges and subtraction of multiples of
a row whose pivot is a unit.  At each stage the original
determinant is the signed product of the previous pivots times
the determinant of the remaining block.  The last entry is
multiplied without inversion.  Thus successful outputs are the
exact determinants, including when the final entry is a nonunit.
Failure to find a unit pivot is reported as inconclusive and
cannot yield a unit certificate.

The wrapper checks that $\det A_0$ is a unit and independently
checks the ring identity
\[
 \det D_e=\det S_e\,(\det A_0)^{-1}
\]
and $\gcd(N(\det D_e),M)=1$.  Therefore $D_e$ is invertible
over $B_M$ by its adjugate.  Base change preserves this inverse
at every coefficient place over a prime divisor of $M$.
The Schur-complement argument of K17-U3-RESIDUAL-RING gives the
full two-class comparison there; its integral conjugation
involution supplies the opposite minuscule direction as well.
\end{proof}

\paragraph{Completed control and provenance.}
For $e=26$, $M=697446897311$, all 732 native identifying rows
give
\[
 \det A_0=65303065346,\qquad
 \det S_e=212516823907,\qquad
 \det D_e=243170653113,
\]
with zero $w$ coordinates and determinant norm gcd one.
The unchanged completed report
\path{native_full_e26_small_control.json} records this fresh
native calculation.  A separate reconciliation
\path{native_full_e26_reconciliation.json} matches all three
residues, the dimensions, parameters, input hashes, frame hash
and full-row count to the carried complete original-Python
\path{e26_m697446897311/verify_ring_schur.json} report.
Fresh original-Python evaluation also agrees on identifying
rows $0,366,731$ and on the two stored extra points.  These
last two checks compare pointwise evaluation formulas; they
are not additional operator-identity tests.  No fresh full
Python rerun is asserted.

The independent integer controls in
\path{audit_full_ring_native.py} compare 24 small random
determinants with the literal permutation formula, over six
moduli including nonsquarefree composite moduli.  Every
returned determinant and unit status agrees.  The invertible
matrix $\left(\begin{smallmatrix}3&1\\5&2\end{smallmatrix}\right)$
over $\mathbf Z/15\mathbf Z$ correctly yields an inconclusive
pivot search, and a one-by-one nonunit is correctly detected.
Four malformed determinant inputs and four malformed or
nonunit-denominator model inputs are rejected.  These controls
supplement the complete mathematical argument and the full
$e=26$ control; they do not certify any unfinished higher-weight
calculation.

% END evidence/higher4/native_full_rows_section.tex


% BEGIN evidence/higher4/exact_cyclic_charpoly_section.tex
\subsection{An exact cyclic certificate for the whole characteristic polynomial}
\begin{proposition}[KT-CYCLIC-CHARPOLY]
Let $A\in M_n(\mathbf Q)$, $n>0$, and let
$f(X)=X^n+c_1X^{n-1}+\cdots+c_n\in\mathbf Q[X]$.
Choose a positive integer $D$ with $B=DA\in M_n(\mathbf Z)$,
and suppose $F(X)=D^nf(X/D)$ has integer coefficients.
If an integer vector $v$ and a prime $q$ satisfy
\[
 \det(v,Bv,\ldots,B^{n-1}v)\not\equiv0\pmod q,
 \qquad F(B)v=0\quad\hbox{over }\mathbf Z,
\]
then $f$ is the characteristic polynomial of $A$.
\end{proposition}
\begin{proof}
The nonzero determinant modulo $q$ proves that the same integer
determinant is nonzero. Thus $v,Bv,\ldots,B^{n-1}v$ form a basis
over $\mathbf Q$. Since $F(B)$ commutes with $B$, the asserted
Horner identity implies that $F(B)$ kills every vector of this
basis. No nonzero polynomial of degree less than $n$ can
annihilate $v$, by linear independence. Therefore the minimal
polynomial of $B$ has degree $n$ and equals the monic polynomial
$F$; by Cayley--Hamilton it also equals the characteristic
polynomial of $B$. Rescaling gives
$\det(XI-A)=D^{-n}F(DX)=f(X)$.
Neither squarefreeness nor irreducibility of $f$ is required.
\end{proof}

The implementation \path{cyclic_charpoly_certificate.py} clears
the exact common denominator, checks the prime by trial division,
finds a cyclic vector using a full modular Krylov determinant,
and checks $F(B)v=0$ by exact integer Horner evaluation.
It calls no characteristic-polynomial routine. Failure to find
a cyclic vector is inconclusive and cannot produce a certificate.
The wrapper \path{replay_cyclic_charpoly.py} retains every
independent orbit, contraction, harmonic-kernel, identifying-point,
all-neighbor, operator-identity, conjugation, factor-product and
extra-point check of the frozen verifier; only its characteristic
polynomial calculation is replaced by the proved certificate.
The claimed coefficients are an explicit argument to that
certificate, and a nonzero final Horner vector rejects them.

% END evidence/higher4/exact_cyclic_charpoly_section.tex


% BEGIN evidence/higher4/exact_sum_first_replay_section.tex
\subsection{Summing complete orbit evaluations before projection}
\label{sec:kt-sum-first-replay}

\begin{proposition}[KT-SUM-FIRST-REPLAY]
The separate wrapper \path{replay_cyclic_sum_first.py}
preserves the full finite-model assertions of the frozen
\path{harmonic_hecke_certificate.py}.  A completed report
certifies every checked operator entry and the whole
characteristic polynomial, with the same mathematical
meaning as a successful full frozen replay.
\end{proposition}
\begin{proof}
Work over $\mathbf Z[\omega]$, with $\omega^2+\omega+1=0$.
Let $o(P)$ be the row of all invariant orbit sums evaluated
at $P$, and let $H$ be the full checked integer harmonic
basis matrix.  The original transformed evaluation is
\[
                    \sum_{g\in\mathcal N} o(gP)H.
\]
The new routine evaluates the same complete row $o(gP)$
for every member of the same checked neighbor list and
returns
\[
                    \left(\sum_{g\in\mathcal N}o(gP)\right)H.
                                                               \tag{KT-SF.1}
\]
These expressions are equal by finite distributivity.
The coordinate transformations, conjugation convention,
orbit enumeration and integer coefficients are unchanged.
Both coordinates in $1,\omega$ are summed exactly.  The
cached sparse columns of $H$ delete only zero coefficients;
the cache retains its matrix object and is replaced whenever
the degree uses another matrix.

The wrapper retains complete source and target orbit
enumeration, the exact contraction, the full harmonic-kernel
rank and containment checks, and every identifying-point
evaluation.  It additionally checks the relevant matrix
shapes and that the number of identifying points equals
the harmonic dimension.  The original nonzero evaluation
determinant and all operator identities are retained, as
are both extra-point operator identities at every prime.
Equation (KT-SF.1) therefore changes neither an entry nor
the scope of an entry check.

The conversion to the rational real form, conjugation
identities, factor-product reconstruction and comparison
values are also retained.  Only the characteristic-polynomial
algorithm is replaced by KT-CYCLIC-CHARPOLY: its explicit
monic candidate is certified by a full modular cyclic basis
and exact integer Horner evaluation.  Failure to find that
basis is inconclusive and cannot give a completed report.

The wrapper executes byte snapshots of the frozen source,
setup driver and cyclic module.  It records source and input
hashes at the beginning and end and requires equality.
The frozen verifier first writes a distinct pending file;
only after these checks does the wrapper atomically write
the named final report, with its effective source hash and
all cyclic certificates.  Thus an unfinished computation
or a changed source/input does not produce a final report
from this run.
\end{proof}

\paragraph{Complete carried control.}
The full degree $k=13$, $e=14$, of harmonic dimension $14$
was extracted without modifying any of its contents from
the carried higher-two dataset.  The new replay completed
both $p=7$ and $p=13$, including every neighbor, the two
extra-point identities at each prime, and both entire
characteristic polynomials.  The complete numerical result
object equals the $k=13$ result of the carried full report.
The old dataset hash, old report hash, extracted dataset
hash, new report hash, and source-hash equality are recorded
in \path{sum_first_k13_reconciliation.json}.
This is a fresh full replay of that carried degree.  Higher
weights require their own completed reports.

% END evidence/higher4/exact_sum_first_replay_section.tex


% BEGIN sections/component_extension.tex
\section{The new finite degrees in the cyclotomic tower}
\label{sec:component-extension}

\begin{corollary}[KT-COMPONENT-EXTENSION]
Use the notation of K17-COMPONENT-CONTROL. For every odd prime
$p$ and positive integer $k$ whose residue modulo $(p-1)/2$
has a representative $j$ with
\[
                         -16\le j\le15,
\]
one has
\[
 K_{4k}(\mathbf Z)\{p\}=0,\qquad
 A_n^{[-2k]}=0\ (n\ge0),\qquad X^{[-2k]}=0.
\]
For $p\ge5$, the actual ramified component $M^{[2k+1]}$ is
integrally free of rank one over the Iwasawa algebra $\Lambda$.
For every odd $p$, the full-domain Soul\'e characters in weights
$27$, $29$ and $31$ are onto in the precise cyclotomic-product
normalization of K17, (K19).
\end{corollary}
\begin{proof}
The carried K17-TWELVE-DEGREES-PERIOD supplies representatives
$-16\le j\le12$. Theorem KT-K52-K56-K60-VANISH supplies $j=13,14,15$.
K17's exact comparison (K2) gives $A_0^{[-2j]}=0$ in these
additional degrees. Congruent exponents select the same
$\mathbf F_p[\operatorname{Gal}(\mathbf Q(\mu_p)/\mathbf Q)]$
eigenspace, so (K2) gives the primary vanishing for the stated
residue classes. K17-COMPONENT-CONTROL, beginning at line~10975
of the supplied paper, then gives every finite-layer and
inverse-limit vanishing and the integral rank-one freeness.
For the exceptional odd characters $\omega^{\pm1}$, its cited
K17-ODD-COMPANION-EXCLUSION, beginning at line~11084,
supplies the same conclusion. These are carried control
statements, not a new Iwasawa-control proof.

Finally the full-domain index formula K17 (K18), in the section
beginning before line~8390, gives
\[
 [\mathbf Z_p(m):\kappa_m(N_p)]
       =|K_{2m-2}(\mathbf Z)\{p\}|.
\]
For $m=27,29,31$ the right sides are one. This proves the three
surjectivity claims with exactly its domain and normalization.
No assertion about the image of a deeper filtration subgroup
follows from this formula.
\end{proof}

The result enlarges selected character ranges in LM6 A.44--A.45
and the bridges Q06; it does not exhaust the characters at a
large prime or establish the full Ihara kernel assertion.
The periodic statement is about vanishing, not equality of the
orders or an isomorphism of all $K$-groups in congruent degrees.

% END sections/component_extension.tex


% BEGIN evidence/higher4/next_vanishing_section.tex
\section{The completed degree sixty-four and its tower consequence}
\label{sec:kt-next-vanishing}

This section uses the same stipulated archive-theorem convention,
the closed source theorem KT-ONECHAR-CLOSED, and the carried
arithmetic comparisons as KT-K52-K56-K60-VANISH.  Its new finite
input is restricted to actually completed full reports.  At the
independent reconciliation on 21 November 2026 at 18:08 UTC,
the complete-degree list is $\{32\}$; the weights $e=34,36$
remain pending.  The frozen snapshot is
\path{next_certificate_independent_snapshot.json}, and
\path{next_review_checks.json} also checks the exact factor
identity below.

\subsection{The full comparisons at parameter 32}

For the initial field $\mathbf Q(\sqrt{-3})$, let
$P_{32,p}$ and $Q_{32,p}$ be the characteristic polynomials
of the full unitary and classical operators in the carried
normalizations, for $p=7,13$.  Set
\[
 C_p=P_{32,p}(1+p^{32}+p^{64})\,
                         Q_{32,p}(1+p^{64}).                 \tag{KN.1}
\]
The invariant source and target dimensions are $1355$ and
$1200$, so the entire harmonic coefficient space has dimension
$155$.  The full classical space has dimension $20$.

Both complete initial-field reports
\path{ktower_full_cyclic_replay_e32.json} and
\path{ktower_full_native_replay_e32.json} have passed.
Their respective SHA256 hashes are
\begin{quote}\small
\path{eb4777af4d01ebed96d13bae591c84281aba671bf76979f107fba260b7c008cb},\\
\path{d5273d5fbd09856ab0b98fc013dcfeeb0dbe6ddcace991d41156591b6ce728a0}.
\end{quote}
They use the same complete input, of SHA256
\begin{quote}\small
\path{992b4de93f23ef658081a0b13ff57d948b9bc340697ae12616d1c34569db435a}.
\end{quote}
The checks include the whole invariant harmonic kernel, every
neighbor evaluation at both test primes, identifying matrices,
both exact characteristic polynomials, and two additional
points at each prime.  The cyclic method uses
KT-CYCLIC-CHARPOLY: the full Krylov determinants modulo the
proved prime $1000003$ are $6457$ at seven and $430478$ at
thirteen, and the exact integer Horner identities are zero.
The second report uses explicit integer Faddeev recursion.
The original model checks are retained in both methods.

Define the two residual integers by
\[
\begin{array}{c|r}
 M_1&25078104767246023\\
 M_2&297281331984585601081730755448375077567651428334282032418227870624001.
\end{array}                                                   \tag{KN.2}
\]
The full comparison evaluations are nonzero, and exact integer
arithmetic gives
\[
\begin{split}
 \gcd(C_7,C_{13})={}&
 2^{667}3^{798}5^{22}11^{11}19^4\,23\,97^4\,36293\,38177
                         \;M_1M_2^3.
\end{split}                                                 \tag{KN.3}
\]
The primality of each displayed small factor was checked by
trial division through its square root.  No primality or
squarefreeness assertion about $M_1$ or $M_2$ is used.
The exact support report has SHA256
\begin{quote}\small
\path{f2473d0d970bcd9f5b726a5108388788374d253bdb3ac935f3d342248e366a7d}.
\end{quote}
Thus every prime at least $2^{31}$ dividing both initial
comparisons divides $M_1M_2$.

For the alternate field $\mathbf Q(\sqrt{-7})$, use the
quadratic coefficient ring and full two-class comparison
$D_{32}$ from (KT-H.1)--(KT-H.2).  The complete signed-orbit
source, target, and harmonic dimensions are respectively
$11712$, $10347$, and $1365$.  Both native full-row reports
have verified all $1365$ identifying rows and all thirteen
frames per row, following the unchanged structural verifier
that reconstructs the whole harmonic kernel and every entry
of the saved integer model.  KT-NATIVE-FULL-ROWS justifies
the exact evaluation and determinant computation.

The resulting residues are as follows.  Every displayed
quadratic-ring determinant has zero $w$ coefficient.
\begin{center}\small\setlength{\tabcolsep}{5pt}
\begin{tabular}{c|c|r}
modulus&quantity&residue\\\hline
$M_1$&$\det A_0$&16143349325810270\\
$M_1$&$\det D_{32}$&22433259464780866\\
$M_2$&$\det A_0$&90816647413234740296735891833530247285152510937657893588581667927649\\
$M_2$&$\det D_{32}$&194069810455851942622854423335689544104099551086410455684338376369326
\end{tabular}
\end{center}
The two completed native-report hashes, in the same order,
are
\begin{quote}\small
\path{01339d0f8a0a7e14b70dc5761ac853df4b1bd725d73197a372c45791bd715a70},\\
\path{e56b077438d147c558c31a6fb7ce7ced658b018821f60bde3609017d613a5597}.
\end{quote}
Both determinant norms are units.  Direct integer checks give
\[
 \gcd(M_i,336\cdot63!)=
 \gcd(M_i,1+2^{32}+2^{64})=1\qquad(i=1,2).                 \tag{KN.4}
\]
Consequently K17-U3-RESIDUAL-RING identifies these with units
in the full two-class comparison at every coefficient place
above every prime dividing the corresponding $M_i$.

The full classical cubic comparison
\[
 J_{32}=\det\bigl((1+2^{64})^3I-T_2^3\mid
                                  S_{65}(\Gamma_1(7))\bigr)
                                                                  \tag{KN.5}
\]
is computed on the entire $125$-dimensional space.  Its residues
are
\[
\begin{array}{c|r}
 M_1&8605330201981558\\
 M_2&211994204480983529461394011887090065119566629384444459141606880381499.
\end{array}                                                   \tag{KN.6}
\]
Each residue, and every inverted pivot, is coprime to its
modulus.  The full cusp-model comparison includes its diamond
correction and the cubic twists.  These two complete numerical
objects, and the initial-field classical object, were freshly
recomputed in \path{verify_next_classical.json}; its hash is
\begin{quote}\small
\path{aeaf27c0f365a8334c4eef5337ef7cbf05dcaace4a34eb6cd160026625ac90cc}.
\end{quote}
That report contains nine complete classical replays across
the three candidate weights.  Only its three $e=32$ rows are
needed here.  The independent reconciliation checks the
current source and data hashes and the entire numerical
objects.  These are fresh executions of the previously audited
classical algorithms; no new classical construction is asserted.

\subsection{Arithmetic and cyclotomic consequences}

\begin{theorem}[KT-K64-VANISH]
\label{thm:kt-k64-vanish}
Under the stipulated archive-theorem convention and the carried
arithmetic comparisons and small-prime results,
\[
                              K_{64}(\mathbf Z)=0.
\]
\end{theorem}
\begin{proof}
Suppose that a prime $\ell$ divides the order of the finite
group $K_{64}(\mathbf Z)$.  The carried Rognes--Weibel result
excludes $\ell=2$.  The stipulated Vandiver verification and
the carried Kurihara equivalence exclude every odd
$\ell<2^{31}$.  Thus $\ell\geq2^{31}$.

For $e=32$, such a prime satisfies $\ell\geq5$,
$(\ell-1)\nmid64$, and $\ell\ne7,13$.  The exact source
datum over $\mathbf Q(\sqrt{-3})$ therefore holds by
\path{ClassNumberThree.lean:gd}.  KT-ONECHAR-CLOSED and
the full initial-field coverage imply that $\ell$ divides
both integers in (KN.1).  Equation (KN.3) forces
$\ell\mid M_1M_2$.

For either residual modulus containing this prime, the
alternate full unitary comparison is a unit by (KN.4)
and the complete native certificates.  The full classical
cubic comparison is a unit by (KN.6).  In particular
$\ell>65=2e+1$, so all the hypotheses of
K17-MINUS7-GENERAL-SOURCE hold.  That carried theorem makes
the rational strict negative Tate Selmer group zero.
K17-VK-PT then makes $K_{64}(\mathbf Z)\{\ell\}$ zero,
contradicting the choice of $\ell$.  Every primary part
vanishes, and finiteness proves the assertion.
\end{proof}

\begin{corollary}[KT-COMPONENT-EXTENSION-16]
For every odd prime $p$ and positive integer $k$ whose residue
modulo $(p-1)/2$ has a representative $j$ satisfying
\[
                              -16\leq j\leq16,
\]
the notation and normalization of K17-COMPONENT-CONTROL give
\[
 K_{4k}(\mathbf Z)\{p\}=0,\qquad
 A_n^{[-2k]}=0\ (n\geq0),\qquad X^{[-2k]}=0.
\]
For $p\geq5$ the actual ramified component $M^{[2k+1]}$ is
integrally free of rank one over $\Lambda$.  For every odd
$p$, the full-domain Soul\'e character $\kappa_{33}$ is onto
in the precise cyclotomic-product normalization of K17 (K19).
\end{corollary}
\begin{proof}
The carried K17-TWELVE-DEGREES-PERIOD supplies representatives
$-16\leq j\leq12$, and KT-K52-K56-K60-VANISH supplies
$j=13,14,15$.  Theorem KT-K64-VANISH supplies $j=16$.
The carried comparison (K2) gives $A_0^{[-32]}=0$ and makes
the vanishing periodic in $k$ modulo $(p-1)/2$.  Apply
K17-COMPONENT-CONTROL and its exceptional-character
K17-ODD-COMPANION-EXCLUSION exactly as in the earlier
component extension.  This proves the finite-layer,
inverse-limit, and integral rank-one conclusions.

Finally the full-domain index formula K17 (K18) is
\[
 [\mathbf Z_p(m):\kappa_m(N_p)]
                      =|K_{2m-2}(\mathbf Z)\{p\}|.
\]
At $m=33$ its right side is one.  This proves the asserted
surjectivity on that exact domain, with no assertion about
the image of a deeper filtration subgroup.  The negative
endpoint remains $-16$; no additional Bernoulli-component
vanishing is inferred from the new positive representative.
\end{proof}

\subsection{The two remaining finite checks}

For $e=34$ or $36$, the following is a conditional criterion
only.  If the saved full initial-field operators at seven and
thirteen are completely verified, the exact support reduction
is reconciled to those operators, and both full alternate
unitary comparisons are certified units at that weight's
two residual moduli, then $K_{2e}(\mathbf Z)=0$.  The full
classical comparisons in those weights have already been
replayed.  The proof is the preceding prime-by-prime argument:
the residual-support certificate confines every prime at least
$2^{31}$ to the two listed moduli, all such primes exceed
$2e+1$, and the two alternate unit conditions invoke the
same carried source theorem.

At the reconciliation cutoff above, those combined hypotheses
have not all been completed for $e=34,36$.  This section
therefore asserts neither $K_{68}(\mathbf Z)=0$ nor
$K_{72}(\mathbf Z)=0$.  A progress line, a producer
characteristic polynomial, or a completed comparison in only
one residual modulus is not substituted for a full report.

% END evidence/higher4/next_vanishing_section.tex


% BEGIN evidence/higher4/later_vanishing_section.tex
\section{The subsequent completed degree sixty-eight}
\label{sec:kt-later-vanishing}

This section records the next completed finite gate separately
from the frozen degree-sixty-four proof.  The independent
reconciliation at 18:17 UTC on 21 November 2026 gives complete
weights $e=32,34$ and leaves $e=36$ pending.  Its frozen file
\path{later_certificate_snapshot_e34.json} has SHA256
\begin{quote}\small
\path{bd30578e286d1eefb223718f05f4e2be2967678ebc958feaa8d3e99dd44d5fb8}.
\end{quote}

\subsection{Completed finite input for parameter 34}

The full initial-field invariant source, target, and harmonic
dimensions are $(1714,1528,186)$.  Both test primes seven and
thirteen have complete original-transform cyclic and
sum-first cyclic replays.  Their entire numerical model
results agree.  Each verifies the full harmonic kernel,
every neighbor evaluation, both exact characteristic
polynomials, and two additional points at each test prime.
The report hashes are
\begin{quote}\small
\path{b1b044cefdef2820b62d725b9a6a30beaeebd9e636dc559c5b7fab1398d9ef71}
\quad(original-transform cyclic),\\
\path{dd92f79563e1642353ff0d0730fcab007ec197043d107c882e81f4624f189480}
\quad(sum-first cyclic).
\end{quote}
Both reconcile to the full input hash
\begin{quote}\small
\path{c1d76928cc874ad8ef83c013e9120576629c2a1c39ddc986c224482cba375646}.
\end{quote}
The sum-first report also records identical starting and ending
hashes for every executed source and both input files.
KT-CYCLIC-CHARPOLY and KT-SUM-FIRST-REPLAY supply the
checked algebraic justification.  The two full Krylov
determinants modulo $1000003$ are $783829$ and $815725$;
both exact integer Horner identities are zero.

Put
\[
\begin{array}{c|r}
 L_1&2371981599141956245883063\\
 L_2&905931677830145683660433370785395963800605403339683999361486128103361399603.
\end{array}                                                    \tag{KL.1}
\]
For the product $C_p$ of the full initial-field unitary and
classical comparisons, evaluated respectively at
$1+p^{34}+p^{68}$ and $1+p^{68}$, exact integer arithmetic gives
\[
 \gcd(C_7,C_{13})
       =2^{956}3^{1005}5^{90}19^{14}23^4 31^3\,1249
                                                    \;L_1L_2^3.
                                                               \tag{KL.2}
\]
The classical initial-field dimension is $22$.  All displayed
small factors are proved primes below $2^{31}$; no primality
or squarefreeness assumption is made about $L_1,L_2$.
The exact support report has SHA256
\begin{quote}\small
\path{4623a966fe77df99943621e17ce8584cac848c42aebb13368375121611cb27e7}.
\end{quote}

The alternate-field source, target, and harmonic dimensions
are $(14860,13222,1638)$.  Both full native reports have
completed all $1638$ identifying rows, each using all thirteen
frames, after the unchanged structural prefix reconstructs
the entire harmonic kernel and integer model.  Their
determinant values are below; every $w$ coefficient is zero.
\begin{center}\footnotesize
\begin{tabular}{c|c|r}
modulus&quantity&residue\\\hline
$L_1$&$\det A_0$&890445239151571061622227\\
$L_1$&$\det D_{34}$&1502530045303762432482728\\
$L_2$&$\det A_0$&579980560248564379336863309163990105435794656506833824085404889575597989344\\
$L_2$&$\det D_{34}$&577158008572250025673753936842704259290368727141758842908187240473442064211
\end{tabular}
\end{center}
The corresponding report hashes are
\begin{quote}\small
\path{683159f0f99e47f56d3ef7af8d8cfbc3477f25b3faaaa03b3b03e44aeb732b1e},\\
\path{91dfc6c0d064845abb12274b25ae2afb69312ef44fdf50dbcf029369a3581fd4}.
\end{quote}
Every determinant norm is a unit, and the exact gcd checks give
\[
 \gcd(L_i,336\cdot67!)=
 \gcd(L_i,1+2^{34}+2^{68})=1\qquad(i=1,2).                  \tag{KL.3}
\]
Thus KT-NATIVE-FULL-ROWS and K17-U3-RESIDUAL-RING prove
the full two-class unitary comparison is a unit at every
coefficient place above every prime divisor of $L_1L_2$.

The full alternate classical cubic comparison on
$S_{69}(\Gamma_1(7))$, of dimension $133$, has residues
\[
\begin{array}{c|r}
 L_1&1818462064779772315365618\\
 L_2&4600576442480975660019765412403247352055986629143676003304820577585327697.
\end{array}                                                    \tag{KL.4}
\]
Each is coprime to its modulus and every inverted pivot is a
unit.  The full diamond-corrected cubic model and the
initial-field classical model have freshly completed replays
in \path{verify_next_classical.json}, already recorded with
its exact hash in the degree-sixty-four section.  The
reconciliation checks the current sources and complete
numerical objects for all three classical rows at $e=34$.

\begin{theorem}[KT-K68-VANISH]
\label{thm:kt-k68-vanish}
Under the same stipulated archive-theorem convention and
carried arithmetic inputs as KT-K64-VANISH,
\[
                              K_{68}(\mathbf Z)=0.
\]
\end{theorem}
\begin{proof}
The carried small-prime results exclude two and every odd
prime below $2^{31}$.  If a remaining prime $\ell$ divided
the order of $K_{68}(\mathbf Z)$, then
$\ell\geq2^{31}>69$, so the exact initial-field GD
hypotheses hold and $\ell\ne7,13$.  KT-ONECHAR-CLOSED
and full initial-field coverage force $\ell$ to divide
both comparison integers in (KL.2), hence $L_1L_2$.

The complete alternate unitary and classical comparisons
are units above this prime by (KL.3)--(KL.4).  Because
$\ell>2\cdot34+1$, K17-MINUS7-GENERAL-SOURCE applies
and gives zero rational strict negative Tate Selmer group.
K17-VK-PT contradicts the proposed nonzero primary part.
Finiteness of $K_{68}(\mathbf Z)$ concludes the proof.
\end{proof}

\begin{corollary}[KT-COMPONENT-EXTENSION-17]
The conclusions of KT-COMPONENT-EXTENSION-16 hold with
the representative interval enlarged to
\[
                             -16\leq j\leq17.
\]
For every odd prime $p$, the full-domain Soul\'e character
$\kappa_{35}$ is onto in the normalization of K17 (K19).
\end{corollary}
\begin{proof}
Theorem KT-K68-VANISH adds the positive representative
$j=17$.  Apply the same exact (K2) comparison,
K17-COMPONENT-CONTROL, and exceptional-character
K17-ODD-COMPANION-EXCLUSION.  The full-domain index
formula at $m=35$ has right side
$|K_{2\cdot35-2}(\mathbf Z)\{p\}|=1$.
The negative endpoint is unchanged, and no assertion about
a deeper-filtration Soul\'e image is made.
\end{proof}

The interrupted redundant initial-field $e=34$
integer-Faddeev job is not counted as a completed replay.
The two full cyclic reports above are complete.  The
parameter $e=36$ still has an unfinished combined gate at
this section's cutoff, so this section makes no assertion
that $K_{72}(\mathbf Z)$ vanishes.

% END evidence/higher4/later_vanishing_section.tex


% BEGIN evidence/higher4/last_vanishing_section.tex
\section{The completed degree seventy-two}
\label{sec:kt-last-vanishing}

This section records the last completed gate among parameters
$32,34,36$, separately from the frozen proofs for degrees
sixty-four and sixty-eight.  The complete-report reconciliation
at 18:40 UTC on 21 November 2026 confirms all three parameters.
The frozen file \path{last_certificate_snapshot_e36.json}
has SHA256
\begin{quote}\small
\path{67e3c068ee2b5c017696745a1d2fd5d65800df528b6ccbec8d1406ba25c28f85}.
\end{quote}
No progress log or interrupted calculation is used in this
completion claim.

\subsection{Completed finite input for parameter 36}

The full initial-field invariant source, target, and harmonic
dimensions are $(2133,1917,216)$.  The sum-first cyclic replay
has completed the full harmonic kernel, every neighbor
evaluation at seven and thirteen, both exact characteristic
polynomials, and two additional points at each test prime.
Its complete report has SHA256
\begin{quote}\small
\path{cae636f52eef4f0335ad8d70e779005d9b0fd753ea397ec39bb84259de71d9b1},
\end{quote}
and reconciles to the full input hash
\begin{quote}\small
\path{3155883f47539624c718e506391f28cca9c04f15cb60c0f930c5d097991ed974}.
\end{quote}
The report records identical starting and ending hashes for
all executed sources and both inputs.  The full Krylov
determinants modulo $1000003$ are $654191$ at seven and
$473965$ at thirteen; both exact integer Horner identities
are zero.  KT-CYCLIC-CHARPOLY and KT-SUM-FIRST-REPLAY
justify these complete characteristic-polynomial checks.

Put
\[
\begin{array}{c|r}
 M_1&1269913439616071\\
 M_2&15161168936170074565138296075385332429054971641697794594125219820150853.
\end{array}                                                     \tag{KM.1}
\]
For the product $C_p$ of the full initial-field unitary and
classical comparison integers, evaluated respectively at
$1+p^{36}+p^{72}$ and $1+p^{72}$, exact arithmetic gives
\[
\begin{aligned}
 \gcd(C_7,C_{13})
   ={}&2^{932}3^{1132}5^{22}7\,11^8 13^{10}19^2 29^4 37^2\\
     &\quad\cdot73^2 421^3\,3343\,18541\;M_1M_2^3.
\end{aligned}                                                    \tag{KM.2}
\]
The initial classical dimension is $23$.  All displayed
small factors are proved primes below $2^{31}$; no
primality or squarefreeness assertion is made about
$M_1,M_2$.  The exact support report has SHA256
\begin{quote}\small
\path{564a0ab4d38a8dc2669db0129994b29ec349655a674d23afc26f03d23a8df093}.
\end{quote}

The alternate-field source, target, and harmonic dimensions
are $(18603,16659,1944)$.  Each complete native report
verifies all $1944$ identifying rows, each using all thirteen
frames, after the unchanged structural prefix reconstructs
the entire harmonic kernel and all integers of the model.
The determinant residues are as follows; every $w$
coefficient is zero.
\begin{center}\footnotesize
\begin{tabular}{c|c|r}
modulus&quantity&residue\\\hline
$M_1$&$\det A_0$&146346680122672\\
$M_1$&$\det D_{36}$&1227124847574193\\
$M_2$&$\det A_0$&6766244652232482632897056974519737998248211708082697277035180883759652\\
$M_2$&$\det D_{36}$&7951138340699857056120169457228565191079127276248573439707562030772663
\end{tabular}
\end{center}
The corresponding complete report hashes are
\begin{quote}\small
\path{16a114520ea7b70f445a989b603de1c2226c5456dbbf25a9b01fa811b761b514},\\
\path{9a98e10449a1c35913256dbe55ad2245f57c1ecf24d3729b93031d368dcef554}.
\end{quote}
All determinant norms are units, and fresh exact gcd
checks give
\[
 \gcd(M_i,336\cdot71!)=
 \gcd(M_i,1+2^{36}+2^{72})=1\qquad(i=1,2).                  \tag{KM.3}
\]
KT-NATIVE-FULL-ROWS and K17-U3-RESIDUAL-RING therefore
give a unit full two-class unitary comparison at every
coefficient place above every prime dividing $M_1M_2$.
The native calculation makes no assumption that its
modulus is prime, and its complete rows are distinguished
from an independent full Python-row replay.

The alternate classical cubic comparison on
$S_{73}(\Gamma_1(7))$, of dimension $141$, has residues
\[
\begin{array}{c|r}
 M_1&40849963779002\\
 M_2&12002152405727334510087168008797785444498771439491866852279373158352988.
\end{array}                                                     \tag{KM.4}
\]
Each is coprime to its modulus, and all inverted pivots
are units.  The complete diamond-corrected cubic model
and the initial-field classical model have fresh full
replays in \path{verify_next_classical.json}, whose hash
is recorded in the degree-sixty-four section.  The current
reconciliation checks the source hashes and all numerical
objects for exactly one initial-field and two alternate-field
classical rows at parameter $36$.

\begin{theorem}[KT-K72-VANISH]
\label{thm:kt-k72-vanish}
Under the stipulated archive-theorem convention and
carried arithmetic inputs of KT-K64-VANISH,
\[
                              K_{72}(\mathbf Z)=0.
\]
\end{theorem}
\begin{proof}
The carried small-prime results exclude two and every odd
prime below $2^{31}$.  Suppose that a remaining prime
$\ell$ divides $|K_{72}(\mathbf Z)|$.  Then
$\ell\geq2^{31}>73$, so $\ell\geq5$, $\ell\ne7,13$,
and $(\ell-1)\nmid72$.  These are the required initial-field
GD conditions.  KT-ONECHAR-CLOSED and the full initial-field
coverage force $\ell$ to divide both comparison integers
in (KM.2), and hence to divide $M_1M_2$.

By (KM.3)--(KM.4), the complete alternate unitary and
classical comparisons are units above this prime.  Since
$\ell>2\cdot36+1$, K17-MINUS7-GENERAL-SOURCE applies
and gives zero rational strict negative Tate Selmer group.
K17-VK-PT then contradicts the assumed nonzero primary
part.  Finiteness of $K_{72}(\mathbf Z)$ concludes the proof.
\end{proof}

\begin{corollary}[KT-COMPONENT-EXTENSION-18]
The conclusions of KT-COMPONENT-EXTENSION-17 hold with
the representative interval enlarged to
\[
                             -16\leq j\leq18.
\]
For every odd prime $p$, the full-domain Soul\'e character
$\kappa_{37}$ is onto in the normalization of K17 (K19).
\end{corollary}
\begin{proof}
Theorem KT-K72-VANISH adds the positive representative
$j=18$.  Apply the same exact (K2) comparison,
K17-COMPONENT-CONTROL, and exceptional-character
K17-ODD-COMPANION-EXCLUSION used in the preceding
corollary.  The full-domain index formula at $m=37$ has
right side $|K_{2\cdot37-2}(\mathbf Z)\{p\}|=1$.
The negative endpoint is unchanged.  The conclusion
concerns the full domain, with no assertion about a
deeper-filtration Soul\'e image.
\end{proof}

The redundant interrupted initial-field integer-Faddeev
calculation is not a completed replay.  The completed
sum-first report above suffices for the finite input;
no unfinished further-degree computation is used here.

% END evidence/higher4/last_vanishing_section.tex


% BEGIN evidence/higher4/further_vanishing_e38.tex
\section{The completed degree seventy-six}
\label{sec:kt-further-e38}

The following proof uses the stipulated archive-theorem convention
and the carried arithmetic inputs of KT-K64-VANISH.  The complete
initial-field replay at parameter $38$ finished at 20:14:03 UTC on
21 November 2026.  All required finite reports are reconciled in
the frozen \path{further_certificate_lean_snapshot_e38.json},
of SHA256
\begin{quote}\small
\path{0d036f4bb18ead1388dfb6e683e3a49be07a1785ba98bb0e44921a182cd5edd5}.
\end{quote}
The separate \path{further_e38_review_checks.json} checks every
displayed factor, determinant residue, and unit condition below.
No unfinished parameter-$40$ computation is used.

\subsection{Complete finite comparisons}

Over $\mathbf Q(\sqrt{-3})$, write $P_{38,p}$ and $Q_{38,p}$ for
the characteristic polynomials of the full unitary and classical
operators in the carried normalizations, and put
\[
 C_p=P_{38,p}(1+p^{38}+p^{76})\,Q_{38,p}(1+p^{76}),
                  \qquad p=7,13.                         \tag{KF38.1}
\]
The invariant source, target, and harmonic dimensions are
$(2632,2374,258)$; the full classical dimension is $24$.
The completed \path{ktower_full_sum_first_replay_e38.json}
has SHA256
\begin{quote}\small
\path{4e4519817db13b7661618f31428b0db521637c6c7b7200f9e1b0cfee152395dc},
\end{quote}
and its complete input has SHA256
\begin{quote}\small
\path{80ef1296f3baea6292a9decb7bc274d9c2954852d4d4f1645722543e39b06bd7}.
\end{quote}
The report verifies the entire harmonic kernel, every neighbor
evaluation, both exact characteristic polynomials, and two
additional points at each test prime.  It records identical
starting and ending hashes for every executed source and both
inputs.  The full Krylov determinants modulo the proved prime
$1000003$ are $418843$ at seven and $575934$ at thirteen;
both exact integer Horner identities vanish.
KT-CYCLIC-CHARPOLY and KT-SUM-FIRST-REPLAY justify these checks.

Define the two residual integers by
{\footnotesize
\[
\begin{array}{c|r}
 M_1&35578577811063982960672862780147\\
 M_2&15251674497332588499311578723828894246360050688455945759751526497655448146123593.
\end{array}                                                \tag{KF38.2}
\]
}
Both integers in (KF38.1) are nonzero, and exact arithmetic gives
\[
 \gcd(C_7,C_{13})=
 2^{1262}3^{1392}5^{127}11^4\,23\,31\,2557^3\,15901^3
                                  \;M_1M_2^3.             \tag{KF38.3}
\]
Every displayed small factor is a proved prime below $2^{31}$.
No primality or squarefreeness assertion about $M_1$ or $M_2$
is used.  The exact support report has SHA256
\begin{quote}\small
\path{33d58cbd24beff9e24850be3f0b5c46a83bdb176858f6fe064a8d0b4a6b0185d}.
\end{quote}

For $\mathbf Q(\sqrt{-7})$, the full signed-orbit source, target,
and harmonic dimensions are $(23014,20728,2286)$.
Each completed native report verifies all $2286$ identifying
rows, with all thirteen actual frames in every row.  Its
unchanged structural prefix reconstructs the entire harmonic
kernel and verifies every integer in the saved model.
Use the notation of (KT-H.2): $A_0$ is the identifying
evaluation matrix, $S_{38}$ is the evaluation matrix of the
Schur operator $D_{38}$, and
$\det D_{38}=\det S_{38}/\det A_0$.
The exact residues are as follows; every $w$ coefficient is zero.
\begin{center}\scriptsize\setlength{\tabcolsep}{4pt}
\begin{tabular}{c|c|r}
modulus&quantity&residue\\\hline
$M_1$&$\det A_0$&13359736957765420668996360679866\\
$M_1$&$\det S_{38}$&12227698356557419998998820128830\\
$M_1$&$\det D_{38}$&25961938570910564303816764703704\\
$M_2$&$\det A_0$&14348380500075505276532267286279171076181263192472552281939481855405800823012461\\
$M_2$&$\det S_{38}$&11806158317438633481119271535689622100249752432683044850797048802985457362450094\\
$M_2$&$\det D_{38}$&12832984267076485504924734210768266922232258447209384964867423882967717792266714
\end{tabular}
\end{center}
The complete report hashes, in the order $M_1,M_2$, are
\begin{quote}\small
\path{f394945b425c33a6aef1e9a823c547c8a8410252ef4bf03580ffedeca85be6a2},\\
\path{6a6bc6711974bb2eeb7289065ba41f7367d1ca559deb3a724eeb0f5c90a9407e}.
\end{quote}
Every determinant norm is a unit, and direct integer checks give
\[
 \gcd(M_i,336\cdot75!)=
 \gcd(M_i,1+2^{38}+2^{76})=1\qquad(i=1,2).                \tag{KF38.4}
\]
KT-NATIVE-FULL-ROWS and K17-U3-RESIDUAL-RING therefore give
a unit full two-class unitary comparison at every coefficient
place above every prime dividing $M_1M_2$.  The native
computation uses unit-norm pivots over the quadratic coefficient
ring, without assuming the modulus prime.  It is a complete
computation; no additional independent Python-row replay is
claimed.

The full alternate classical cubic comparison on
$S_{77}(\Gamma_1(7))$, of dimension $149$, is
\[
 J_{38}=\det\bigl((1+2^{76})^3I-T_2^3
                              \mid S_{77}(\Gamma_1(7))\bigr).
                                                               \tag{KF38.5}
\]
Its exact residues are
{\footnotesize
\[
\begin{array}{c|r}
 M_1&28120119304299729255421071437422\\
 M_2&13636901229281977226565975129352822785735618961606779966242920847684659072051446.
\end{array}                                                \tag{KF38.6}
\]
}
Each residue and every inverted pivot is coprime to its modulus.
The complete diamond-corrected cubic models, together with the
initial-field classical model, have fresh full replays in
\path{verify_further_classical_e38.json}, of SHA256
\begin{quote}\small
\path{fd8d0a6492118bcf276e09733c98e62f7d6aeaac956c02ab370e34581e44c8da}.
\end{quote}
This report checks exactly one initial-field and two
alternate-field numerical objects, and records unchanged
starting and ending source and input hashes.

\begin{theorem}[KT-K76-VANISH]
\label{thm:kt-k76-vanish}
Under the stipulated archive-theorem convention and carried
arithmetic inputs of KT-K64-VANISH,
\[
                         K_{76}(\mathbf Z)=0.
\]
Consequently $K_{4k}(\mathbf Z)=0$ for every $1\leq k\leq19$.
\end{theorem}
\begin{proof}
The carried small-prime results exclude two and every odd
prime below $2^{31}$.  Suppose a remaining prime $\ell$
divides the order of the finite group $K_{76}(\mathbf Z)$.
Then $\ell\geq2^{31}>77$, so $\ell\geq5$, $\ell\ne7,13$,
and $(\ell-1)\nmid76$.  The exact initial-field source datum
therefore holds by \path{ClassNumberThree.lean:gd}.
KT-ONECHAR-CLOSED and the full initial-field coverage force
$\ell$ to divide both integers in (KF38.1).  Equation (KF38.3)
then forces $\ell\mid M_1M_2$.

For a modulus containing this prime, the full alternate
unitary comparison is a unit by (KF38.4) and the complete
native report.  The full classical comparison is a unit
by (KF38.6).  Since $\ell>2\cdot38+1$, all the hypotheses
of K17-MINUS7-GENERAL-SOURCE hold, including the integral
Hecke lattices supplied by K17-U3-RESIDUAL-RING and the
classical construction.  That theorem gives zero rational
strict negative Tate Selmer group.  K17-VK-PT contradicts
the proposed nonzero primary part.  Finiteness concludes
$K_{76}(\mathbf Z)=0$.  Combining this with the carried
first twelve degrees and KT-K52-K56-K60-VANISH,
KT-K64-VANISH, KT-K68-VANISH, and KT-K72-VANISH gives
the assertion for $1\leq k\leq19$.
\end{proof}

\begin{corollary}[KT-COMPONENT-EXTENSION-19]
For every odd prime $p$ and positive integer $k$ whose residue
modulo $(p-1)/2$ has a representative $j$ with
\[
                              -16\leq j\leq19,
\]
the notation of K17-COMPONENT-CONTROL gives
\[
 K_{4k}(\mathbf Z)\{p\}=0,\qquad
 A_n^{[-2k]}=0\quad(n\geq0),\qquad X^{[-2k]}=0.
\]
For $p\geq5$, the actual ramified component $M^{[2k+1]}$
is integrally free of rank one over $\Lambda$.
For every odd $p$, the full-domain Soul\'e character
$\kappa_{39}$ is onto in the normalization of K17 (K19).
\end{corollary}
\begin{proof}
KT-COMPONENT-EXTENSION-18 supplies the interval $[-16,18]$,
and KT-K76-VANISH supplies $j=19$.
The exact comparison (K2) makes this vanishing periodic
modulo $(p-1)/2$.  Apply K17-COMPONENT-CONTROL and its
exceptional-character K17-ODD-COMPANION-EXCLUSION exactly
as in the preceding component extension.  The full-domain
index formula K17 (K18), at $m=39$, gives
\[
 [\mathbf Z_p(39):\kappa_{39}(N_p)]
                  =|K_{76}(\mathbf Z)\{p\}|=1.
\]
The negative endpoint remains $-16$.  The Soul\'e assertion
uses the full domain $N_p$; it asserts no deeper-filtration
surjectivity.
\end{proof}

% END evidence/higher4/further_vanishing_e38.tex


% BEGIN evidence/higher4/further_vanishing_e40.tex
\section{The completed degree eighty}
\label{sec:kt-further-e40}

The following proof uses the stipulated archive-theorem convention
and the carried arithmetic inputs of KT-K64-VANISH.
The complete initial-field replay at parameter $40$ finished at
21:14:43 UTC on 21 November 2026.  All required finite reports
are reconciled in the frozen
\path{further_certificate_lean_snapshot_e40.json}, of SHA256
\begin{quote}\small
\path{e361ae1f564b607dd6abe5b12f84abaf8038eccd99ebab625018198c23492200}.
\end{quote}
The separate \path{further_e40_review_checks.json} checks the
displayed support, determinant residues, and unit conditions.
All computations used below have completed final reports.

\subsection{Complete finite comparisons}

Over $\mathbf Q(\sqrt{-3})$, write $P_{40,p}$ and $Q_{40,p}$ for
the characteristic polynomials of the full unitary and classical
operators in the carried normalizations, and put
\[
 C_p=P_{40,p}(1+p^{40}+p^{80})\,Q_{40,p}(1+p^{80}),
                     \qquad p=7,13.                      \tag{KF40.1}
\]
The invariant source, target, and harmonic dimensions are
$(3216,2915,301)$; the full classical dimension is $26$.
The completed \path{ktower_full_sum_first_replay_e40.json}
has SHA256
\begin{quote}\small
\path{d90be25e89d8b416da3ab3e701ee75a39b9cd97c3807ef23852f6756b0d81214},
\end{quote}
and its complete input has SHA256
\begin{quote}\small
\path{15e9d86e9489fabdf6823b4c069a76e86e14b4e769f0f2240ef61f1b4297d764}.
\end{quote}
The report verifies the entire harmonic kernel, every neighbor
evaluation, both exact characteristic polynomials, and two
additional points at each test prime.  It records identical
starting and ending hashes for every executed source and both
inputs.  The full Krylov determinants modulo the proved prime
$1000003$ are $186012$ at seven and $203153$ at thirteen;
both exact integer Horner identities vanish.
KT-CYCLIC-CHARPOLY and KT-SUM-FIRST-REPLAY justify these checks.

Define the two residual integers by
{\scriptsize
\[
\begin{array}{c|r}
 M_1&97006879910137295142114566339844001\\
 M_2&7108763171297761022563973098474159328699317449544482052462302898913655347936016491916627179.
\end{array}                                                \tag{KF40.2}
\]
}
Both integers in (KF40.1) are nonzero, and exact arithmetic gives
\[
 \gcd(C_7,C_{13})=
 2^{1237}3^{1481}5^{30}\,17\,19^3\,23\,53^3\,61^2\,263^3
                                      \;M_1M_2^3.          \tag{KF40.3}
\]
Every displayed small factor is a proved prime below $2^{31}$.
No primality or squarefreeness assertion about $M_1$ or $M_2$
is used.  The exact support report has SHA256
\begin{quote}\small
\path{bfc4fdcd5276394b20b3b84e7a06c8f7e4fa101fead4a0b5a527b5352a0e2686}.
\end{quote}

For $\mathbf Q(\sqrt{-7})$, the full signed-orbit source, target,
and harmonic dimensions are $(28167,25500,2667)$.
Each completed native report verifies all $2667$ identifying
rows, with all thirteen actual frames in every row.
Its unchanged structural prefix reconstructs the entire harmonic
kernel and verifies every integer in the saved model.
Use the notation of (KT-H.2): $A_0$ is the identifying evaluation
matrix, $S_{40}$ is the evaluation matrix of the Schur operator
$D_{40}$, and $\det D_{40}=\det S_{40}/\det A_0$.
The exact residues are as follows; every $w$ coefficient is zero.
\begin{center}\scriptsize\setlength{\tabcolsep}{0.5pt}
\begin{tabular}{c|c|r}
modulus&quantity&residue\\\hline
$M_1$&$\det A_0$&24508767132937698280274696417479690\\
$M_1$&$\det S_{40}$&57701312716688370785284167967093223\\
$M_1$&$\det D_{40}$&18475600431218418646318122285471634\\
$M_2$&$\det A_0$&6602630443705463115559079483133840822758782262835980542851052045757427148377946169622677766\\
$M_2$&$\det S_{40}$&2826085683880373070743038779071061534477046118495422200970314778769155029063802009223721753\\
$M_2$&$\det D_{40}$&4127418073500754648834354537349014264644056349444060619639548562532345933174679396418989537
\end{tabular}
\end{center}
The complete report hashes, in the order $M_1,M_2$, are
\begin{quote}\small
\path{c7f90cdd49d13f573d6d12ef8dc1e70d258911a2d5c76f2488d538c7c7d6688d},\\
\path{4d2a876b8bbc6bb04ec2bebf8a8811c7d953c5cff00805c29b263d90f4c1dbfc}.
\end{quote}
Every determinant norm is a unit, and direct integer checks give
\[
 \gcd(M_i,336\cdot79!)=
 \gcd(M_i,1+2^{40}+2^{80})=1\qquad(i=1,2).                \tag{KF40.4}
\]
KT-NATIVE-FULL-ROWS and K17-U3-RESIDUAL-RING therefore give
a unit full two-class unitary comparison at every coefficient
place above every prime dividing $M_1M_2$.
The native computation uses unit-norm pivots over the quadratic
coefficient ring, without assuming the modulus prime.
It is a complete computation; no additional independent
Python-row replay is claimed.

The full alternate classical cubic comparison on
$S_{81}(\Gamma_1(7))$, of dimension $157$, is
\[
 J_{40}=\det\bigl((1+2^{80})^3I-T_2^3
                              \mid S_{81}(\Gamma_1(7))\bigr).
                                                               \tag{KF40.5}
\]
Its exact residues are
{\scriptsize
\[
\begin{array}{c|r}
 M_1&16333300733145557954211546506408385\\
 M_2&4806930230144634588779829558730356227270450680470411532734969121338126545951002472836095182.
\end{array}                                                \tag{KF40.6}
\]
}
Each residue and every inverted pivot is coprime to its modulus.
The complete diamond-corrected cubic models, together with the
initial-field classical model, have fresh full replays in
\path{verify_further_classical_e40.json}, of SHA256
\begin{quote}\small
\path{5349dd81a15907e1f253a7bcaa74bbfc2eefcec5e9b3c932a94ac46eaa308717}.
\end{quote}
This report checks exactly one initial-field and two
alternate-field numerical objects, and records unchanged
starting and ending source and input hashes.

\begin{theorem}[KT-K80-VANISH]
\label{thm:kt-k80-vanish}
Under the stipulated archive-theorem convention and carried
arithmetic inputs of KT-K64-VANISH,
\[
                         K_{80}(\mathbf Z)=0.
\]
Consequently $K_{4k}(\mathbf Z)=0$ for every $1\leq k\leq20$.
\end{theorem}
\begin{proof}
The carried small-prime results exclude two and every odd prime
below $2^{31}$.  Suppose a remaining prime $\ell$ divides
the order of the finite group $K_{80}(\mathbf Z)$.
Then $\ell\geq2^{31}>81$, so $\ell\geq5$, $\ell\ne7,13$,
and $(\ell-1)\nmid80$.  The exact initial-field source datum
therefore holds by \path{ClassNumberThree.lean:gd}.
KT-ONECHAR-CLOSED and the full initial-field coverage force
$\ell$ to divide both integers in (KF40.1).  Equation (KF40.3)
then forces $\ell\mid M_1M_2$.

For a modulus containing this prime, the full alternate
unitary comparison is a unit by (KF40.4) and the complete
native report.  The full classical comparison is a unit
by (KF40.6).  Since $\ell>2\cdot40+1$, all hypotheses of
K17-MINUS7-GENERAL-SOURCE hold, including the integral Hecke
lattices supplied by K17-U3-RESIDUAL-RING and the classical
construction.  That theorem gives zero rational strict
negative Tate Selmer group.  K17-VK-PT contradicts the
proposed nonzero primary part.  Finiteness concludes
$K_{80}(\mathbf Z)=0$.
The range through $k=19$ is KT-K76-VANISH, so adjoining this
case gives the assertion for $1\leq k\leq20$.
\end{proof}

\begin{corollary}[KT-COMPONENT-EXTENSION-20]
For every odd prime $p$ and positive integer $k$ whose residue
modulo $(p-1)/2$ has a representative $j$ with
\[
                              -16\leq j\leq20,
\]
the notation of K17-COMPONENT-CONTROL gives
\[
 K_{4k}(\mathbf Z)\{p\}=0,\qquad
 A_n^{[-2k]}=0\quad(n\geq0),\qquad X^{[-2k]}=0.
\]
For $p\geq5$, the actual ramified component $M^{[2k+1]}$
is integrally free of rank one over $\Lambda$.
For every odd $p$, the full-domain Soul\'e character
$\kappa_{41}$ is onto in the normalization of K17 (K19).
\end{corollary}
\begin{proof}
KT-COMPONENT-EXTENSION-19 supplies the interval $[-16,19]$,
and KT-K80-VANISH supplies $j=20$.
The exact comparison (K2) makes this vanishing periodic
modulo $(p-1)/2$.  Apply K17-COMPONENT-CONTROL and its
exceptional-character K17-ODD-COMPANION-EXCLUSION exactly
as in the preceding component extension.  The full-domain
index formula K17 (K18), at $m=41$, gives
\[
 [\mathbf Z_p(41):\kappa_{41}(N_p)]
                  =|K_{80}(\mathbf Z)\{p\}|=1.
\]
The negative endpoint remains $-16$.  The Soul\'e assertion
uses the full domain $N_p$; it asserts no deeper-filtration
surjectivity.
\end{proof}

% END evidence/higher4/further_vanishing_e40.tex


% BEGIN sections/uniform_gate_heading.tex
\section{The uniform arithmetic gate and residual congruences}

The finite certificates above settle their stated weights.
The following results address uniform elimination of the remaining
residual primes and the precise limitations of proposed extensions.

% END sections/uniform_gate_heading.tex


% BEGIN evidence/main_gate/residual_global_gate.tex
\subsection{KT-JOINT-RECHECK: the carried global-congruence reduction}
\label{sec:kt-global-residual}

The supplied checkpoint already proves K17-S17-JOINT,
K17-S17-GROUP-RING, and the joint operator test in
\path{evidence/work/hecke_uniform/joint_trace_section.tex}.
The source audit below independently rechecks that carried
argument and makes its semisimplification consequence explicit;
it is not counted as a new arithmetic theorem. The argument
uses the lower-level model theorem underlying
\path{OneCharU_s}, not merely its final quantified statement.
All archive inputs have the source-standing convention and the
pinned commit of KT-ONECHAR-CLOSED. No new Lean kernel build is
claimed. Write AV/ for the archive directory
\path{src/Development/AutomorphicArithmetic/AutomorphicVarieties/},
U3/ for AV/\path{UnitaryThreeGroups/}, and DTC/ for
AV/\path{ModularityArgument/ThresholdArguments/DenseThresholdCases/}.

Fix data $(\Omega,K,\ell,c,N,e,\psi,\iota)$ satisfying the
archive's \texttt{GD}, with unramified coefficient ring
$\mathcal O=c.O$; in particular $e(\mathcal O)=1$.
This includes the rational specialization $\mathcal O=\mathbf Z_\ell$.
Let $\mathscr A$ be the union of the two \emph{actual} admissible
classes in DTC/\path{NumberU3US.lean}:49--59:
\[
 \operatorname{Adm3U\_s}(\iota,e,N,\rho)
 \quad\hbox{or}\quad
 \operatorname{Adm2A}(\iota,e,N|\operatorname{disc}K|,\rho).
                                                               \tag{GR.1}
\]
Put $\Psi=\operatorname{Chr.Psi}(\ell,e,\operatorname{dirG}\psi)$
and write
\[
 \overline T=\overline\Psi|_{G_K}\oplus1
                    \oplus\overline\Psi^{-1}|_{G_K}.
                                                               \tag{GR.2}
\]
Residual representations are over $\overline{\mathbf F}_\ell$.
An integral model means matrices over the valuation ring of
$\overline{\mathbf Q}_\ell$ with inverse matrices also integral.

\begin{theorem}[K17-S17-JOINT, rechecked unramified specialization]
Under the nine closed archive inputs recorded in
KT-ONECHAR-CLOSED, if no $\rho\in\mathscr A$ has residual
semisimplification isomorphic to $\overline T$, then
\[
                   \operatorname{Vlen}(\Omega,c,e,\psi)=0.
                                                               \tag{GR.3}
\]
Equivalently, nonzero strict Selmer length forces one actual
admissible source representation with the global residual
semisimplification (GR.2). In fact the construction supplies
an integral model whose reduction is already the displayed
diagonal representation.
\end{theorem}

\begin{proof}
Use precisely the construction in
DTC/\path{FOneCharU.lean}:31--38: the closed inputs supply
$D,g_1,U,S,M$, where $S.R$ is the deformation ring and $M$
its graded model. They also supply finiteness of its
weight-zero subalgebra and the point hypothesis used in
\path{model_upper}; the identical applications are written
at lines 55--57 of that file. These data are obtained before
using either the chosen Galois element or its trace hypothesis.

For a local point $f:S.R\to\overline{\mathbf Q}_\ell$ define
$\rho_f=D.\operatorname{rhoK}(\operatorname{push}f\,S.r)$.
Theorem \path{isPt_push}, U3/\path{IsPtOf.lean}:124--166,
shows that the pushed representation has integral matrix
entries and, for every $g\in G_{\mathbf Q}$ and every pair of
indices $a,b$, satisfies
\[
 \left| (\operatorname{push}f\,S.r)(g)_{ab}
              -r_0(\eta,\Psi)(g)_{ab}\right|_\ell<1.
                                                               \tag{GR.4}
\]
Its entire proof of this clause is at lines 136--145:
$S.r.\operatorname{red}$ makes the difference lie in the
maximal ideal of $S.R$, and the defining local-point property
maps that ideal into elements of norm strictly less than one.
The integrality clause follows from \path{ipo_le_one},
lines 95--112. Applying it to $g^{-1}$ also gives integral
inverse matrices. These are the exact \texttt{int} and
\texttt{red} clauses of \texttt{Dat.IsPt},
U3/\path{RhoKApply.lean}:76--83.

On restricting to $G_K=\ker\eta$, the definition of
$D.\operatorname{rhoK}$ and its application lemma
(U3/\path{RhoKApply.lean}:70--74) identify these matrices
with those of $\rho_f$. Reduction of (GR.4) is therefore
exactly (GR.2). For every \texttt{TwoPt} point, the
\texttt{hadm} argument in DTC/\path{FOneCharU.lean}:47--54
places $\rho_f$ in (GR.1). Here \texttt{TwoPt} includes
\texttt{IsLocPt} by definition,
U3/\path{ConjStable.lean}:94--95.
Consequently the hypothesis of the theorem implies that
there are no \texttt{TwoPt} points at all.

Apply \path{model_upper} with the identity Galois element.
Its sign condition is $\eta(1)=1$. Its remaining trace
condition says that the trace difference is nonzero at
every \texttt{TwoPt}; this is now vacuous. The theorem
(U3/\path{ModelUpper.lean}:104--112) yields finite model
length and a finite set $T$ satisfying
\[
 \operatorname{CoversPt}(T),\qquad
 \operatorname{nrm}(M.\operatorname{len})
       \leq2\sum_{f\in T}v_\ell(\operatorname{trDiff}(1,f)).
                                                               \tag{GR.5}
\]
By the first clause of \texttt{CoversPt}, every member of
$T$ is a \texttt{TwoPt}; see U3/\path{ConjStable.lean}:97--99.
Thus $T=\varnothing$ and the right side of (GR.5) is zero.
The normalized length is nonnegative; it is the finite
length divided by $e(\mathcal O)=1$, by
AV/\path{CharacterLevelBounds/ShaFamMono.lean}:71--83.
Hence $M.\operatorname{len}=0$. Finally
\path{vlen_le}, U3/\path{VlenLe.lean}:52--67,
gives $\operatorname{Vlen}\leq M.\operatorname{len}=0$.
\end{proof}

\paragraph{Scope and use.}
The theorem supplies a necessary global congruence, not an
equivalence between a congruence and nonzero Selmer length.
The example KT-TWOFIELD-661 has genuine congruences while
the rational component vanishes. No ambient counting
formula alone rules out (GR.2). The reduction removes the
need to find a single Galois element separating every
ambient source, and it does not assume that the finite set
$X$ in \path{OneCharU_s} is independent of that element.
It also keeps the \texttt{GD} and admissibility hypotheses:
it is not an assertion for arbitrary representations or
for an auxiliary field outside the established coverage.

\begin{proposition}[Horizontal form of the carried joint operator test]
Let $k$ be a field, $M$ a finite-dimensional $k$-space of
dimension $d$, and $A_1,\ldots,A_s$ commuting endomorphisms.
For specified scalars $t_1,\ldots,t_s\in k$, put
\[
 B=[\,A_1-t_1I\mid\cdots\mid A_s-t_sI\,].
                                                               \tag{GR.6}
\]
Then $B$ has row rank $d$ if and only if there is no
joint eigencharacter of the algebra $k[A_1,\ldots,A_s]$
on $M\otimes_k\overline k$ taking every $A_j$ to $t_j$.
Equivalently, the ideal generated by the $A_j-t_jI$ is
the unit ideal in that algebra.
\end{proposition}

\begin{proof}
Let $A=k[A_1,\ldots,A_s]\subseteq\operatorname{End}_k(M)$
and $I=(A_j-t_jI)_j$. The column space of $B$ is $IM$,
so full row rank is equivalent to $IM=M$.
The finite $A$-module $M$ is faithful by the definition of
$A$, and its support is therefore all of $\operatorname{Spec}A$.
For completeness, if $M_{\mathfrak m}=0$ at a maximal ideal,
finite generation supplies an element outside $\mathfrak m$
annihilating $M$, contradicting faithfulness. If $I$ were
proper, choose $\mathfrak m\supset I$. Nakayama applied to
the nonzero finite module $M_{\mathfrak m}$ gives
$IM_{\mathfrak m}\ne M_{\mathfrak m}$, a contradiction.
Thus $IM=M$ implies $I=A$; the converse is immediate.

If $I$ is proper, a maximal quotient of the finite
commutative algebra $A/I$ is a finite field extension of
$k$ and embeds in $\overline k$. This yields the required
character. Its localized module is nonzero by the same
faithfulness argument, so it occurs as a composition
factor of $M\otimes\overline k$ and, equivalently, as a
joint eigencharacter. Conversely such a character kills
$I$ and prevents $I=A$.
\end{proof}

In an established integral Hecke model covering all the
sources in (GR.1), apply this proposition to reductions
of commuting trace operators and the residual Eisenstein
traces. Full row rank for a finite collection of operators
rules out the global congruence (GR.2), even if every
individual determinant $\det(A_j-t_jI)$ vanishes.
Combining that certificate for every required source
branch with K17-S17-JOINT proves the desired vanishing.
The coverage and normalization remain explicit hypotheses;
a deficient block for finitely many primes does not by
itself establish a global congruence at all primes.


% END evidence/main_gate/residual_global_gate.tex


% BEGIN evidence/main_gate/residual_trace_avoidance.tex
\subsection{KT-TRACE-AVOID: the exact single-element obstruction}
\label{sec:kt-trace-avoid}

Let $\Gamma$ be a profinite group and let
$\overline\rho_1,\ldots,\overline\rho_m,\overline T$
be continuous finite-dimensional representations over
$\overline{\mathbf F}_\ell$, with finite image. They
factor through their common finite image $G$. In the
application $\overline T=\overline\Psi\oplus1\oplus
\overline\Psi^{-1}$ and all dimensions are three.
Write $K_0(G)$ for the Grothendieck ring of finite-dimensional
$\overline{\mathbf F}_\ell[G]$-modules, using tensor product
as multiplication, and form
\[
        V=\prod_{i=1}^m([\overline\rho_i]-[\overline T])
                          \quad\hbox{in }K_0(G).
                                                               \tag{TA.1}
\]

\begin{proposition}[KT-TRACE-AVOID]
There is an element $\sigma\in\Gamma$ such that
\[
 \operatorname{tr}\overline\rho_i(\sigma)
       \ne\operatorname{tr}\overline T(\sigma)
                     \qquad\hbox{for every }i
                                                               \tag{TA.2}
\]
if and only if $V\notin\ell K_0(G)$. Thus the obstruction
is an exact virtual tensor-product relation modulo $\ell$;
it is not in general equivalent to one factor vanishing.
\end{proposition}

\begin{proof}
Let $S_1,\ldots,S_r$ be the simple modules over the finite
group algebra $k[G]$, where $k=\overline{\mathbf F}_\ell$.
Their trace functions are linearly independent over $k$.
Indeed a relation on all group elements extends linearly
to the whole group algebra. Its semisimple quotient is
$\prod_j\operatorname{End}_k(S_j)$, and a lift of a single
diagonal matrix unit in the $j$th factor forces the $j$th
coefficient of the relation to be zero. This argument
does not require $\ell\nmid|G|$.

Trace is additive in short exact sequences and multiplicative
under tensor product. Therefore the trace of $V$ at $g$ is
\[
 \operatorname{tr}V(g)=
 \prod_i\bigl(\operatorname{tr}\overline\rho_i(g)
                   -\operatorname{tr}\overline T(g)\bigr).
                                                               \tag{TA.3}
\]
Since $K_0(G)$ is free on the classes $[S_j]$, the trace
function of $V$ vanishes identically exactly when every
integer coefficient of $V$ in this basis is divisible by
$\ell$, by the independence just proved. Equation (TA.3)
is nonzero at some $g$ exactly when every factor there is
nonzero. Lift $g$ to $\Gamma$ to obtain (TA.2).
\end{proof}

\paragraph{Individual traces versus simultaneous avoidance.}
For two semisimple representations of the same dimension
$n<\ell$, equality of traces on every group element implies
isomorphism. The preceding independence proves this directly:
each difference of simple multiplicities lies in $[-n,n]$,
so divisibility by $\ell$ forces it to be zero. Thus for
$n=3$, $\ell>3$, a single factor in (TA.1) vanishes modulo
$\ell$ exactly when its residual semisimplification is
$\overline T$. Nevertheless a product of nonzero factors
can vanish, as the following explicit example shows.

\begin{example}[KT-TRACE-COVER]
For every prime $\ell>3$, take
$G=C_3\times C_2\times C_2$ and choose a faithful character
$\Psi:C_3\to\overline{\mathbf F}_\ell^\times$.
Let $\chi_1,\chi_2,\chi_3$ be the three nontrivial
characters of $C_2\times C_2$, extended trivially on $C_3$.
Set
\[
 T=\Psi\oplus1\oplus\Psi^{-1},\qquad
 R_i=(\Psi\chi_i)\oplus1\oplus(\Psi^{-1}\chi_i).
                                                               \tag{TA.4}
\]
The three $R_i$ are pairwise nonisomorphic semisimple
three-dimensional representations; none is isomorphic to
$T$, and $\Psi^2\ne1$. Each has determinant one and is
self-dual. For $g=(h,v)$,
\[
 \operatorname{tr}R_i(g)-\operatorname{tr}T(g)
       =(\chi_i(v)-1)(\Psi(h)+\Psi(h)^{-1}).
                                                               \tag{TA.5}
\]
Every $v\in C_2\times C_2$ lies in the kernel of at least
one nontrivial $\chi_i$. Hence every $g$ makes at least
one difference in (TA.5) zero; no single element separates
all three systems. These statements hold for arbitrarily
large $\ell$ while the number and dimensions of the
systems stay fixed. This is an abstract residual example,
not an assertion that these representations are members
of the automorphic admissible class (GR.1).
\end{example}

\begin{proof}
Choose independent quadratic characters $\chi_1,\chi_2$
and put $\chi_3=\chi_1\chi_2$. The three nonzero elements
of $C_2\times C_2$ each have one of these characters equal
to $1$ and the other two equal to $-1$; at the identity
all three are $1$. Their kernels therefore cover the
group. The characters occurring in (TA.4) are distinct
because $\Psi$ has order three and the $\chi_i$ have
order two. These observations prove all the assertions,
including the displayed trace formula.
\end{proof}

This example also illustrates the stronger multi-operator
criterion rechecked alongside K17-S17-JOINT. At $h=1$, choose nonzero
$v_1,v_2$ with respectively $\chi_1(v_1)=1$ and
$\chi_2(v_2)=1$. On the three-system eigenvalue space,
the two trace-difference operators are
\[
 A_1=\operatorname{diag}(0,-4,-4),\qquad
 A_2=\operatorname{diag}(-4,0,-4).
                                                               \tag{TA.6}
\]
Each determinant is zero, but $[A_1\mid A_2]$ has full
row rank three for every odd $\ell$. Thus a multi-operator
certificate can succeed even when every individual group
element fails. A linear combination of group elements is
not itself a Galois element, so this does not change the
literal single-element assertion of \path{OneCharU_s}.


% END evidence/main_gate/residual_trace_avoidance.tex


% BEGIN evidence/main_gate/two_field_eisenstein_obstruction.tex
\subsection{KT-TWOFIELD-661: simultaneous genuine Eisenstein congruences}
\label{sec:kt-twofield-661}

Put $\chi_D=\chi_{-D}$ for $D=3,7$, and use the standard
normalization
\[
 \sum_{a=1}^{D}\frac{\chi_D(a)t e^{at}}{e^{Dt}-1}
       =\sum_{j\geq0}B_{j,\chi_D}\frac{t^j}{j!}.
                                                               \tag{TF.1}
\]
The following example gives a concrete obstruction to proving
an all-weight unit theorem merely by switching between the
two classical auxiliary levels and varying split Frobenius
primes. It is not an obstruction to the rational $K$-group
vanishing, which the separate reflected test actually proves
at this prime.

\begin{proposition}[KT-TWOFIELD-661]
Let $e=268$, $w=2e+1=537$, and $\ell=661$. Then
\[
 \ell\mid\operatorname{num}\frac{B_{w,\chi_3}}w,
 \qquad
 \ell\mid\operatorname{num}\frac{B_{w,\chi_7}}w,
 \qquad
 B_{\ell-1-e}=B_{392}\equiv53\pmod\ell.
                                                               \tag{TF.2}
\]
For each $D=3,7$ there are a normalized primitive newform
$f_D\in S_{537}(\Gamma_0(D),\chi_D)$ and a coefficient
prime $\lambda_D\mid661$ for which
\[
 a_q(f_D)\equiv1+\chi_D(q)q^{536}\pmod{\lambda_D}
                  \qquad(q\text{ prime},\ q\nmid D).
                                                               \tag{TF.3}
\]
Consequently every split-prime classical characteristic
polynomial comparison at $1+q^{536}$ is divisible by $661$
at both auxiliary levels. In the full level-seven model,
the twist-inclusive determinant
$\det((1+2^{536})^3I-T_2^3)$ is also nonunit at $661$.
Nevertheless the carried K17-REFLECTOR gives
\[
                         K_{536}(\mathbf Z)\{661\}=0.
                                                               \tag{TF.4}
\]
\end{proposition}

\begin{proof}
The prime $661$ is proved prime by trial division through
$25$. It satisfies $661>538$ and is coprime to $42\cdot537$.
The complete standard-library certificate
\path{bernoulli_661_certificate.py} uses two independent
finite identities to verify the generalized residues.

First, the defining recurrence
\[
 B_0=1,\qquad
 B_m=-\frac1{m+1}\sum_{j=0}^{m-1}\binom{m+1}{j}B_j
                                                               \tag{TF.5}
\]
is computed modulo $661$ through $m=537$. Every divisor
$m+1$ used is smaller than $661$, so it is a unit.
Bernoulli polynomials and (TF.1) give the finite identity
\[
 B_{w,\chi_D}
   =\sum_{a=1}^{D}\chi_D(a)
     \sum_{j=0}^{w}\binom wj B_j a^{w-j}D^{j-1}.
                                                               \tag{TF.6}
\]
Its residues are zero for both $D$; division by $w$ is
permitted. The same recurrence gives $B_{392}\equiv53$.

As an independent check of the generalized values, set
$S_D=\sum_{a=1}^{D\ell}\chi_D(a)a^w$.
Multiplication of (TF.1) by $(e^{D\ell t}-1)/t$ gives
\[
 S_D=\sum_{j=0}^{w}\binom wj B_{j,\chi_D}
             \frac{(D\ell)^{w+1-j}}{w+1-j}
       \equiv D\ell B_{w,\chi_D}\pmod{\ell^2}.
                                                               \tag{TF.7}
\]
All denominators are units, and the $B_{j,\chi_D}$ are
$\ell$-integral by (TF.5)--(TF.6). Direct modular exponentiation
gives $S_3\equiv S_7\equiv0\pmod{661^2}$, again proving
the first two congruences. A separate FLINT exact-rational
calculation, stored with complete numerators, gives
valuation one in both cases, with respective denominators
$9$ and $7$ for $B_{537,\chi_D}/537$. Exact valuation is
not needed for the conclusion.

Now apply Fretwell--Roberts, \emph{Newform Eisenstein
congruences of local origin}, Ramanujan J.~64 (2024),
505--527, Theorem~3.6, pp.~517--518, with
\[
 N=D,\quad M=1,\quad k=537,\quad
 \psi=1,\quad\phi=\chi_D,\quad\lambda'=(661).
\]
The conductor is squarefree, $k>2$, and $661\nmid6D$.
The required positive valuation is precisely that of
$L(1-k,\chi_D)=-B_{k,\chi_D}/k$, established above.
The theorem supplies (TF.3); its paragraph after the proof
states that $M=1$ and primitive character make the form new.
The theorem's entire displayed proof was read in the original
published source. Its modular reduction and lifting inputs
are cited, not independently re-proved here.

For a split prime, $\chi_D(q)=1$, so (TF.3) makes
$1+q^{536}-a_q(f_D)$ nonunit. The characteristic polynomial
of the full classical space contains this eigenvalue,
and its evaluation is therefore divisible by $661$.
For level seven the quadratic-character space is contained
in the full $\Gamma_1(7)$ space. Since $\chi_7(2)=1$,
the eigenvalue $a_2(f_7)$ makes the twist-inclusive cube
determinant nonunit as asserted.

Finally $2\leq e=268\leq658=\ell-3$ and the nonzero
residue $B_{392}\equiv53$ meets exactly the hypotheses
of K17-REFLECTOR, (RB.1), K17.tex lines~4533--4619.
That carried theorem kills the cyclotomic even component
$A_{661}^{[-268]}=A_{661}^{[392]}$, hence (TF.4).
\end{proof}

\paragraph{The precise all-weight gate exposed by the example.}
For even $e\geq2$ define
\[
 \mathcal E_D(e)=\left\{\ell>2e+2:\ \ell\nmid6D,\quad
       \ell\mid\operatorname{num}
                 \frac{B_{2e+1,\chi_D}}{2e+1}\right\}.
                                                               \tag{TF.8}
\]
Theorems~3.5--3.6 of Fretwell--Roberts, in their $M=1$
case, identify this with the prime support of the indicated
classical Eisenstein eigensystem congruences. These are
simultaneous congruences at all good primes, stronger than
an accidental common divisor of two determinant evaluations.
Equation (TF.2) disproves the hoped-for assertion
$\mathcal E_3(e)\cap\mathcal E_7(e)=\varnothing$.

A sufficient arithmetic statement for eliminating this
particular common support would instead be
\[
 \ell\in\mathcal E_3(e)\cap\mathcal E_7(e)
          \quad\Longrightarrow\quad
                   \ell\nmid B_{\ell-1-e}.
                                                               \tag{TF.9}
\]
It holds in the displayed example. No theorem establishing
(TF.9) uniformly has been obtained in this run. Moreover,
(TF.9) alone would not complete the full Selmer argument:
the independent rank-three unitary comparison can also be
nonunit, even when a classical branch is a unit. The
complete all-weight problem requires its actual joint
exceptional support to avoid the rational component
$A_\ell^{[-e]}$. A generalized Bernoulli numerator attached
to $\chi_D$ cannot simply be identified with this component.
The reflected ordinary Bernoulli number in (TF.9) records
the exact distinct character required by the supplied
class-field argument.

One can also see the distinction directly in the residual
representations. The eigensystem in (TF.3) has semisimplified
characters $1$ and $\chi_D\omega^{2e}$; after the weight-zero
twist these become $\omega^{-e}$ and $\chi_D\omega^e$.
An extension between them has ratio character
$\chi_D\omega^{2e}$ or $\chi_D\omega^{-2e}$, whereas
the rational negative Selmer problem uses $\omega^{-e}$.
Restriction to $G_{\mathbf Q(\sqrt{-D})}$ removes
$\chi_D$ but does not remove the doubling of the exponent.
For odd $\ell$, restriction of a rational cohomology class
to this quadratic field is injective by restriction--corestriction,
whose composite is multiplication by two. Thus field
variation supplies a new automorphic comparison, not a
formal annihilation of the original rational class.


% END evidence/main_gate/two_field_eisenstein_obstruction.tex


% BEGIN evidence/main_gate/two_field_global_obstruction.tex
\subsection{KT-TWOFIELD-GLOBAL: arbitrary Galois elements do not remove the example}
\label{sec:kt-twofield-global}

\begin{corollary}[KT-TWOFIELD-GLOBAL]
For the two forms in KT-TWOFIELD-661, let
$r_D:G_{\mathbf Q}\to\operatorname{GL}_2(
\overline{\mathbf Q}_{661})$ be the standard Galois
representation attached to $f_D$, with arithmetic-Frobenius
trace $a_q(f_D)$ and determinant
$\chi_D\epsilon^{536}$, where $\epsilon$ is the
$661$-adic cyclotomic character, and put
$\omega=\epsilon\bmod661$. Define
\[
 R_D=1\oplus(r_D\otimes\epsilon^{-268})|_{G_{E_D}},
 \qquad E_D=\mathbf Q(\sqrt{-D}).
                                                               \tag{TG.1}
\]
At the coefficient prime supplied by KT-TWOFIELD-661,
\[
 \overline R_D^{\mathrm{ss}}\simeq
          1\oplus\omega^{-268}\oplus\omega^{268}
                    \qquad(D=3,7).
                                                               \tag{TG.2}
\]
Hence the difference between $\operatorname{tr}R_D(\sigma)$
and $1+\epsilon(\sigma)^{-268}+\epsilon(\sigma)^{268}$
is nonunit for \emph{every} $\sigma\in G_{E_D}$.
Both $f_3$ and $f_7$ are non-CM.
\end{corollary}

\begin{proof}
The attached integral Galois representation and its trace
and determinant normalization are recalled in
Fretwell--Roberts, Section 2.1. The first paragraph of
the proof of their Theorem 3.5, equation (5), explicitly
applies Chebotarev and Brauer--Nesbitt to an Eisenstein
congruence and obtains its residual semisimplification.
Here it gives
\[
              \overline r_D^{\mathrm{ss}}
                     \simeq1\oplus\chi_D\omega^{536}.
                                                               \tag{TG.3}
\]
Equivalently, this follows by continuity from (TF.3) on
Frobenius classes and linear independence of simple
residual trace functions; $661>2$ prevents a multiplicity
ambiguity. Restrict (TG.3) to $G_{E_D}$, where $\chi_D=1$,
twist by $\omega^{-268}$, and add the trivial character.
This proves (TG.2) and the assertion for every $\sigma$.

At level three, KT-GL2-CM gives zero CM dimension because
$3\nmid268$. At level seven the same CM level formula
from Sch\"utt, Theorems 1.1 and 1.4, forces any CM newform
of this level to be induced from an unramified Hecke
character of $\mathbf Q(\sqrt{-7})$. The prime three is
inert in that field, so its Hecke eigenvalue would be
$a_3=0$. But (TF.3) gives
\[
                     a_3(f_7)\equiv1-3^{536}
                                      \equiv50\pmod{661},
                                                               \tag{TG.4}
\]
a contradiction. The last residue is a direct modular
exponentiation. Thus neither congruence can be removed
by deleting only CM newforms from the classical spaces.
\end{proof}

This assertion concerns the indicated classical residual
systems. It does not assert a new automorphic admissibility
or packet-coverage theorem. Whenever an established ambient
model includes these systems, its trace comparison has
this global obstruction. KT-TWOFIELD-661 still proves
$K_{536}(\mathbf Z)\{661\}=0$ independently: a genuine
global congruence is a necessary source condition in
the carried K17-S17-JOINT, not a sufficient condition for the
rational Selmer group to be nonzero.


% END evidence/main_gate/two_field_global_obstruction.tex


% BEGIN evidence/main_gate/two_field_infinite_family.tex
\subsection{KT-TWOFIELD-FAMILY: an unbounded family of common congruences}
\label{sec:kt-twofield-family}

\begin{proposition}[KT-TWOFIELD-FAMILY]
For every integer $t\geq0$, set
\[
            e_t=268+330t,\qquad w_t=2e_t+1=537+660t.
                                                               \tag{TI.1}
\]
Then, for each $D\in\{3,7\}$,
\[
        \frac{B_{w_t,\chi_D}}{w_t}\in661\mathbf Z_{661}.
                                                               \tag{TI.2}
\]
There is a normalized non-CM primitive newform
$f_{D,t}\in S_{w_t}(\Gamma_0(D),\chi_D)$, with a
coefficient prime $\lambda_{D,t}\mid661$, satisfying
\[
 a_q(f_{D,t})\equiv1+\chi_D(q)q^{w_t-1}
          \pmod{\lambda_{D,t}}
                 \qquad(q\text{ prime},\ q\nmid D).
                                                               \tag{TI.3}
\]
Its normalized three-dimensional classical residual system
$1\oplus(r_{f_{D,t}}\otimes\epsilon^{-e_t})|_{G_{E_D}}$
has semisimplification
$1\oplus\omega^{-e_t}\oplus\omega^{e_t}$.
Consequently no Galois element separates that system from
the Eisenstein triple modulo $661$.
Nevertheless, for every $t\geq0$,
\[
                    K_{536+660t}(\mathbf Z)\{661\}=0.
                                                               \tag{TI.5}
\]
\end{proposition}

\begin{proof}
Use Zhi-Hong Sun, \emph{Congruences concerning Bernoulli
numbers and Bernoulli polynomials}, Discrete Applied
Mathematics 105 (2000), 193--223, Lemma 8.1(a),
pp.~219--220, with $p=661$, character $\chi_D$, and
$n=1$. For any positive $b$ not divisible by $660$ it gives
\[
 (1-\chi_D(661)661^{b-1})\frac{B_{b,\chi_D}}b
 \equiv
 (1-\chi_D(661)661^{b+659})
           \frac{B_{b+660,\chi_D}}{b+660}
                                  \pmod{661}.
                                                               \tag{TI.4}
\]
The conductor is coprime to $661$. Every $b=w_t$ is
positive, exceeds one, and is congruent to $537$ modulo
$660$, so the lemma applies. Both Euler factors in
(TI.4) are $661$-adic units congruent to one. The base
case (TI.2) is KT-TWOFIELD-661. Induction using (TI.4)
proves (TI.2) for all $t$. This argument remains valid
when $661\mid w_t$: the cited congruence concerns the
normalized rational number, so no division by a residue
class of $w_t$ is being performed.

Apply Fretwell--Roberts Theorem 3.6 with
$N=D,M=1,k=w_t,\psi=1,\phi=\chi_D$, and the prime
$(661)$. Its sufficient-condition theorem requires
$k>2$, $661\nmid6D$, and the positive valuation of
$L(1-k,\chi_D)=-B_{k,\chi_D}/k$. These conditions hold
for every $t$. In contrast to their necessary-condition
Theorem 3.5, Theorem 3.6 does \emph{not} impose
$661>k+1$. Primitive nebentypus again makes the resulting
eigenform new. This proves (TI.3).

Since $330\equiv0\pmod3$, every $e_t$ is congruent to
one modulo three. KT-GL2-CM therefore excludes CM at
level three. At level seven, a CM newform would have
zero eigenvalue at the inert prime three. Fermat's theorem
and (TI.3) instead give
\[
 a_3(f_{7,t})\equiv1-3^{536+660t}
                  \equiv1-3^{536}\equiv50\pmod{661}.
\]
The CM classification and inert-prime rule used here
are the same as in KT-TWOFIELD-GLOBAL. Thus both forms
are non-CM. Finally the residual semisimplification follows
from the all-good-primes congruence exactly as in
KT-TWOFIELD-GLOBAL, followed by restriction and twisting.

To prove (TI.5), the already certified residue
$B_{392}\equiv53\pmod{661}$ and the independent short
recurrence certificate $B_{62}\equiv511\pmod{661}$
apply K17-REFLECTOR at $e=268$ and $e=598$, respectively.
Both exponents are even and belong to $[2,658]$.
Thus $A_{661}^{[-268]}=A_{661}^{[-598]}=0$.
The exponent $e_t$ modulo $660$ is $268$ for even $t$
and $598$ for odd $t$. The carried comparison (K2),
which holds for every positive $k$, now gives
$K_{2e_t}(\mathbf Z)\{661\}=0$. This uses periodicity
only of the class-group character index and of the
vanishing criterion, not an isomorphism or order formula
for $K$-groups in different degrees.
\end{proof}

The statement is an infinite-weight obstruction at the
fixed prime $661$. Only its initial member $t=0$ lies in
the large-prime range $661>2e_t+2$ used to define
$\mathcal E_D(e)$ in (TF.8). It does not claim infinitely
many distinct common primes in that range, nor does it
contradict the proposed reflected exclusion (TF.9).
It proves that discarding finitely many weights, changing
the test Galois element, and deleting CM forms cannot
remove this fixed common classical congruence support.



% END evidence/main_gate/two_field_infinite_family.tex


% BEGIN evidence/main_gate/stable_residual_propagation.tex
\section{KT-STABLE-PROPAGATION: genuine stable congruences in unbounded weights}
\label{sec:kt-stable-propagation}

We use the compact unitary group, standard full lattice stabilizer,
and class-number-one realization already fixed in KT-U3-DIM and
KT-U3-ENDO-EXACT. Thus, for $a\geq0$, the coefficient representation
$V_a$ has highest weight $(a,0,-a)$ and dimension $(a+1)^3$, and
the algebraic modular-form space is
\[
                     M_a=V_a^{\Gamma},\qquad |\Gamma|=1296.
                                                               \tag{SP.1}
\]
The scalar center acts trivially on $V_a$. In even degrees
$e=a+1\geq2$, the current-run theorem KT-U3-ENDO-EXACT states that
the nonstable part of $M_a$ has dimension $\lfloor a/3\rfloor$;
its complement consists of stable parameters with a single
cuspidal $\mathrm{GL}_3/E$ constituent. Here
$E=\mathbf Q(\sqrt{-3})$. Those spectral and local-level results
are cited, not reconstructed in this section.

The theorem below concerns actual stable automorphic eigensystems
in this full compact space. It does not assert the converse to
source coverage, namely membership in \path{Adm3U_s} with its
additional \texttt{WitA} and \texttt{Unl} predicates.

\begin{lemma}[Finite-group multiplicity growth]
Fix a prime $\ell>3$ and put $G=\operatorname{PGL}_3(\mathbf F_\ell)$.
Let $\overline V_a$ be the reduction of an algebraic integral
lattice in the representation with highest weight $(a,0,-a)$.
For every simple $\overline{\mathbf F}_\ell[G]$-module $S$,
\[
 [\overline V_a:S]
   =\frac{[\overline{\mathbf F}_\ell[G]:S]}{|G|}(a+1)^3+O_{\ell,S}(a+1).
                                                               \tag{SP.2}
\]
Every displayed leading coefficient is strictly positive.
\end{lemma}

\begin{proof}
Fix compatible lifts of prime-to-$\ell$ roots of unity to complex
roots of unity, and use Brauer characters on $\ell$-regular
elements. The reduction has the same formal weight character
as its integral algebraic lattice. If a lift of such an element
to $\mathrm{GL}_3(\mathbf F_\ell)$ has lifted eigenvalues
$x_1,x_2,x_3$, its Brauer character is therefore the Weyl
character
\[
 \beta_a(g)=
 \frac{\det\bigl(x_i^{a+2},x_i,x_i^{-a}\bigr)_{i=1}^3}
      {\det\bigl(x_i^2,x_i,1\bigr)_{i=1}^3}.
                                                               \tag{SP.3}
\]
The expression at coinciding eigenvalues is its continuous
confluent-determinant value. A lift of an $\ell$-regular
projective element can be chosen $\ell$-regular: its order
divides the projective order times $\ell-1$.

If the three eigenvalues are distinct, the denominator in
(SP.3) is fixed and nonzero while each of the six terms in
the numerator has complex absolute value one. Thus the
character is bounded independently of $a$. If exactly two
eigenvalues coincide, cancel their common Vandermonde factor
by differentiating one row. The derivative of $x^{a+2}$ or
$x^{-a}$ contributes at most a linear factor in $a$, while
the resulting denominator is fixed and nonzero. Hence
$\beta_a(g)=O_g(a+1)$ in this case. If all three eigenvalues
coincide, the projective element is the identity and the
central weight zero gives
\[
                      \beta_a(1)=(a+1)^3.
                                                               \tag{SP.4}
\]
There are only finitely many classes, so the constants are
uniform for fixed $\ell$ on all nonidentity $\ell$-regular
elements.

Here are sufficient algebraic details for extracting
multiplicities, without assuming semisimplicity of
$\overline{\mathbf F}_\ell[G]$. The ordinary modular trace
functions of its simple modules are linearly independent:
extend a relation to the group algebra and use a lift of a
diagonal matrix unit in one factor of its split semisimple
quotient. Traces on an element equal traces on its
$\ell$-regular part, since the commuting $\ell$-part acts
unipotently on each eigenspace of the regular part. Thus a
matrix of simple trace values on regular classes has a
nonzero maximal minor over $\overline{\mathbf F}_\ell$.
Lifting that minor entrywise by Brauer characters gives a
nonzero complex determinant: its reduction is the original
nonzero minor. The simple Brauer characters are therefore
linearly independent over $\mathbf C$ as well.

Let $r_S=[\overline{\mathbf F}_\ell[G]:S]$. The regular
module's Brauer character is $|G|$ at the identity and
zero at every other regular element. Indeed its matrices
are permutation matrices with cycles of length the order
of the element, and the roots of unity in a nontrivial
cycle sum to zero. Subtract from $\beta_a$ the regular
character times $(a+1)^3/|G|$. The resulting class function
has every value $O_\ell(a+1)$. Inverting the fixed nonzero
minor just described expresses every simple multiplicity
difference as $O_{\ell,S}(a+1)$, proving (SP.2).
Finally $r_S>0$, since every simple module is a quotient
of the regular module.
\end{proof}

\paragraph{A crude effective version.}
Set
\[
                  N_G=|G|=\ell^3(\ell^2-1)(\ell^3-1).
\]
The preceding proof gives the explicit lower bound
\[
 [\overline V_a:S]\ \geq
 \frac{(a+1)^3}{N_G}
       -N_G^{\,2N_G^2+4}(a+2).
                                                               \tag{SP.2a}
\]
Here are the estimates, to make its effectiveness auditable.
Since the character has central weight zero, rescale an
eigenvalue triple by its first entry. The resulting ratios
have orders dividing the projective element's order, hence
dividing $N_G$. Distinct such roots of unity have distance
at least $4/N_G$. The ordinary or confluent Weyl determinant
therefore gives
$|\beta_a(g)|\leq N_G^3(a+2)$ for every nonidentity regular
element. In the fixed invertible minor of simple Brauer
characters, every entry and every Galois conjugate has
absolute value at most $N_G$. If the minor has size $r$,
then $r\leq N_G$, and its determinant and cofactors have
absolute values at most
$H=N_G^{2N_G}$. Its nonzero determinant is an algebraic
integer in a cyclotomic field of degree at most $N_G$:
the common order of the relevant prime-to-$\ell$ roots
divides the prime-to-$\ell$ part of $N_G$.
The absolute norm is a nonzero integer, so the absolute
value of our chosen determinant is at least
$H^{-(N_G-1)}$. Every inverse entry is at most
$H^{N_G}=N_G^{2N_G^2}$.
Multiplying the residual character-value bound by the at
most $N_G$ entries in a row gives the error in (SP.2a).

In particular,
\[
                  a\geq N_G^{\,N_G^2+3}
        \quad\Longrightarrow\quad
                       [\overline V_a:S]>a/3
                           \quad\hbox{for every }S.
                                                               \tag{SP.2b}
\]
Indeed $a\geq2$, so $a+2\leq2a$; if
$C=N_G^{2N_G^2+4}$, the lower bound is at least
$a(a^2/N_G-2C)$, and the stipulated $a$ gives
$a^2/N_G\geq N_G C$. Since $N_G>3$, this exceeds $a/3$.
This bound is deliberately far larger than a useful
computational threshold.

\begin{lemma}[Propagation through the fixed Hecke action]
Assume $\ell>3$ splits in $E$. Fix a residue eigensystem
$\lambda$ occurring in some $M_{a_0}$ modulo $\ell$ for
all good split-prime Hecke operators, including both
minuscule directions. Then its composition multiplicity
in the reductions of $M_a$ is at least
\[
                            c_\lambda(a+1)^3-C_\lambda(a+1)
                                                               \tag{SP.5}
\]
for constants $c_\lambda>0$ and $C_\lambda$ independent
of $a$.
\end{lemma}

\begin{proof}
At the split place $\ell$, the integral unitary group is
$\mathrm{GL}_3(\mathbf Z_\ell)$. Choose the integral
algebraic lattice $L_a$ in $V_a$. Central weight zero
means its reduction factors through $G$ above. The image
$\overline\Gamma$ has order prime to $\ell$; consequently
the averaging idempotent
\[
       u=|\overline\Gamma|^{-1}
                      \sum_{\gamma\in\overline\Gamma}\gamma
                   \in\overline{\mathbf F}_\ell[G]
                                                               \tag{SP.6}
\]
makes invariants exact and gives
$(L_a^\Gamma)\otimes\overline{\mathbf F}_\ell=u\overline V_a$.

For each rational split prime $q\ne\ell$, choose once and
for all the rational neighbor representatives defining
the unscaled compact Hecke operator in either direction.
Their denominators are supported at $q$, so they are
integral units at $\ell$. The same representatives work
in every coefficient representation. Reduction therefore
gives a fixed element
\[
                h_{q,\pm}=u\left(\sum_j\overline g_{q,\pm,j}\right)u
                          \in u\overline{\mathbf F}_\ell[G]u
                                                               \tag{SP.7}
\]
whose action on $u\overline V_a$ is the reduced Hecke
operator. Let $H$ be the finite-dimensional algebra
generated by all these fixed elements and the identity
$u$. No weight-dependent conjugation of a computed matrix
is being used.

One may realize $u\overline{\mathbf F}_\ell[G]$ as the
weight-zero algebraic modular forms at the auxiliary level
whose subgroup at $\ell$ reduces to the scalar matrices.
The class-number-one decomposition identifies its finite
set with $\overline\Gamma\backslash G$. The operators
(SP.7) are the usual spherical Hecke operators away from
$\ell$ on that space, so they commute. Alternatively,
commutativity is unnecessary for the following exact-functor
argument; on each actual $M_a$ it is already the usual
commuting Hecke action.

The functor $W\mapsto uW$ from finite $G$-modules to
finite $H$-modules is exact. Since $\lambda$ occurs in
$u\overline V_{a_0}$, at least one simple $G$-module $S$
has $[uS:\lambda]>0$. Applying the functor to a composition
series and using additivity of Jordan--H\"older multiplicities
gives
\[
 [u\overline V_a:\lambda]
       =\sum_S[\overline V_a:S][uS:\lambda]
       \geq[\overline V_a:S_0]
                                                               \tag{SP.8}
\]
for one such $S_0$. Equation (SP.2) proves (SP.5).
\end{proof}

\begin{lemma}[Integral eigenflags]
Let $L$ be a free lattice over a complete DVR $O$, and
let a commuting family of integral endomorphisms of $L$
be simultaneously diagonalizable after a finite extension
of its fraction field. Over a sufficiently large finite
coefficient extension, the composition multiplicity of a
specified residue eigencharacter equals the number of
characteristic-zero eigenlines, counted with their vector-space
multiplicities, whose characters reduce to it.
\end{lemma}

\begin{proof}
There are finitely many independent endomorphisms and
eigencharacters, since they lie in $\operatorname{End}_O(L)$;
choose a finite extension splitting them all. Intersect $L$
with a common eigenline. This is a saturated rank-one
sublattice, invariant under every endomorphism, and the
quotient lattice is free. Repeat on the quotient. This
constructs an invariant flag with free rank-one successive
quotients and the characteristic-zero eigencharacters on
them. Reduction is exact at each step because the
quotients are free. Jordan--H\"older multiplicities of the
resulting residue characters give the assertion. The
flag may begin with all eigenlines in any specified
characteristic-zero invariant direct summand.
\end{proof}

\begin{theorem}[KT-STABLE-PROPAGATION]
Fix a split prime $\ell>3$, an even $e_0\geq2$, and a
simultaneous residual Eisenstein eigensystem in $M_{e_0-1}$.
Use the normalization in which the raw compact minuscule
operator is $q$ times the normalized standard Satake
trace, and the opposite direction is $q$ times its dual
trace. Suppose the seed character satisfies, in both
directions and at every split prime $q\ne\ell$,
\[
                 \lambda(h_{q,\pm})
                       =q^{1-e_0}+q+q^{1+e_0}\pmod\ell.
                                                               \tag{SP.9}
\]
Then, for every sufficiently large even $e$ with
$e\equiv e_0\pmod{\ell-1}$, the stable summand of
$M_{e-1}$ contains a characteristic-zero eigenline
whose normalized standard and dual traces satisfy
\[
                \operatorname{tr}_{q,\pm}
                   \equiv q^{-e}+1+q^e\pmod\lambda_e
                                                               \tag{SP.10}
\]
simultaneously at all these primes, for one coefficient
prime $\lambda_e\mid\ell$.
For the normalized attached Galois representation, this
is the global residual semisimplification
\[
                     1\oplus\omega^e\oplus\omega^{-e}
                              \quad\hbox{over }G_E.
                                                               \tag{SP.11}
\]
The explicit sufficient inequality
$e-1\geq N_G^{N_G^2+3}$ is available from (SP.2b).
\end{theorem}

\begin{proof}
Set $a=e-1$. By (SP.5), the multiplicity of the fixed
residue character $\lambda$ in $M_a$ grows at least
cubically in $a$. The entire endoscopic dimension is
$\lfloor a/3\rfloor$ by KT-U3-ENDO-EXACT. The integral
eigenflag lemma applied to the commuting Hecke operators
therefore forces at least one stable characteristic-zero
eigenline reducing to $\lambda$, once $a$ is sufficiently
large. In fact the number of such stable eigenlines,
counting eigenspace multiplicities after splitting
coefficients, is bounded below by (SP.5) minus
$\lfloor a/3\rfloor$. Semisimplicity here is the usual
characteristic-zero automorphic spectral decomposition
of this finite compact modular-form space.

For $e\equiv e_0\pmod{\ell-1}$, the raw Eisenstein targets
in (SP.9) equal $q^{1-e}+q+q^{1+e}$ modulo $\ell$.
Dividing the raw operator eigenvalue by the unit $q$
gives (SP.10). The single residue character is defined
on the algebra generated by \emph{all} these operators,
so the resulting coefficient prime and eigenline work
simultaneously at all of them. They are not selected
separately for different Frobenius primes.

For the Galois attachment use Shin, \emph{Galois representations
arising from some compact Shimura varieties}, Annals of
Mathematics 173 (2011), Theorem~1.2 and its following
compatible-system paragraph, pp.~1647--1648. It applies to
the regular algebraic conjugate-self-dual cuspidal
rank-three constituents in the stable spectrum and supplies
a continuous semisimple system over a finite coefficient
field with the needed local compatibility. Its extra
slight-regularity hypothesis is confined to even rank.
Use the same unitary normalization as in the established
compact source coverage. No characteristic-zero
irreducibility assertion is needed here. The two
directions account for both prime ideals of $E$ over
each split rational prime. Those degree-one prime ideals,
away from any fixed finite exceptional set, have Frobenius
classes dense in $G_E$ for finite quotients: in a normal
closure over $\mathbf Q$, any desired element fixing $E$
is realized by Chebotarev at a rational prime splitting
in $E$, and a prime above it realizes the desired class.
Consequently (SP.10) gives the same residual trace function
as (SP.11) on all of $G_E$. The representations factor
through a finite common residual quotient. Linear
independence of simple modular trace functions, with
dimension $3<\ell$, then identifies their semisimplifications,
as in KT-TRACE-AVOID. This proves (SP.11).
\end{proof}

\begin{corollary}[KT-PROPAGATION-SMALL-SUMMANDS]
Under the hypotheses of KT-STABLE-PROPAGATION, let
$D_a\subset M_a$ be any family of characteristic-zero
Hecke-invariant subspaces along the stipulated progression,
and suppose
\[
                         \dim D_a=o((a+1)^3).
                                                               \tag{SP.11a}
\]
For all sufficiently large $a$ in that progression, a
Hecke-invariant complement of $D_a$ contains a common
eigenline with the same all-prime residual Eisenstein
character. More precisely the number of such lines in
the complement, counted with vector-space multiplicity,
is at least
\[
             c_\lambda(a+1)^3-C_\lambda(a+1)-\dim D_a.
                                                               \tag{SP.11b}
\]
The constants are those of (SP.5). If $D_a$ contains the
entire endoscopic subspace, these surviving lines are stable.
\end{corollary}

\begin{proof}
The characteristic-zero Hecke operators are simultaneously
diagonalizable. Thus $D_a$ is a sum of subspaces of common
eigenspaces and has a Hecke-invariant complement. Choose
an integral eigenflag to count the total dimensions of
the common eigenspaces whose characters reduce to
$\lambda$. Their intersections with $D_a$ have total
dimension at most $\dim D_a$. Subtract that number from
(SP.5), giving (SP.11b). This dimension argument also
allows $D_a$ itself to have been specified only after
extending scalars: a common eigenspace has an algebraic
basis, so if it is not contained in $D_a$, some algebraic
eigenvector lies outside $D_a$. The bound is positive for
all sufficiently large $a$, since $c_\lambda>0$ and
(SP.11a) holds. All-prime normalization and the Galois
conclusion are unchanged. No compatibility between the
subspaces $D_a$ at different weights is assumed.
\end{proof}

\begin{corollary}[KT-STABLE-LARGE-PRIME]
Let
\[
                      \ell=8361790201>2^{31}.
                                                               \tag{SP.12}
\]
For every sufficiently large $t\geq0$, there is a stable
cuspidal compact eigensystem with
\[
              e=18+(\ell-1)t,
\]
whose normalized Galois representation has residual
semisimplification $1\oplus\omega^{18}\oplus\omega^{-18}$
at a prime above $\ell$. Nevertheless, for \emph{every}
$t\geq0$,
\[
                K_{36+2(\ell-1)t}(\mathbf Z)\{\ell\}=0.
                                                               \tag{SP.13}
\]
Thus removing the endoscopic factors and changing the
split Frobenius prime cannot make the level-one compact
stable comparison a unit in these sufficiently large
weights at this fixed prime above $2^{31}$.
\end{corollary}

\begin{proof}
The standard-library certificate
\path{large_prime_seed_certificate.py} proves primality
by trial division, verifies $\ell\equiv1\pmod3$, and
computes exactly
\[
 \frac{B_{37,\chi_{-3}}}{37}
       =-\frac{1701990248643001964415638362402}{3},
 \qquad
 \ell\mid1701990248643001964415638362402.
                                                               \tag{SP.14}
\]
It uses the ordinary Bernoulli recurrence through index
$37$ and the exact Bernoulli-polynomial formula for
the character of conductor three. Fretwell--Roberts
Theorem 3.6, with $N=3,M=1,k=37$, gives a primitive
newform $f\in S_{37}(\Gamma_0(3),\chi_{-3})$ satisfying
\[
              a_q(f)\equiv1+\chi_{-3}(q)q^{36}
                            \pmod\lambda
                         \qquad(q\ne3)
                                                               \tag{SP.15}
\]
for one coefficient prime $\lambda\mid\ell$.
The congruence is not CM: at the inert prime two a
level-three CM form would have eigenvalue zero, whereas
the same certificate gives
\[
                    1-2^{36}\equiv6536635074\ne0\pmod\ell.
                                                               \tag{SP.16}
\]

KT-U3-ENDO-INJECT therefore supplies a genuine endoscopic
eigenline in $M_{17}$. Its unitary parameter is
$1\boxplus\operatorname{BC}(\pi_f)$. At each split prime
the raw eigenvalue is $q(1+a_q(f)/q^{18})$; the dual
trace has the same residual value. Equation (SP.15)
gives exactly (SP.9) with $e_0=18$. An integral primitive
vector on this eigenline reduces to the required seed
character, simultaneously for all Hecke operators.
KT-STABLE-PROPAGATION now gives the stable conclusion.

For the rational component, cite the carried complete
certificate K17-BERNOULLI-8361790201 and K17-REFLECTOR:
\[
             B_{\ell-19}\equiv1827838999\ne0\pmod\ell,
             \qquad A_\ell^{[-18]}=0.
\]
Every $e=18+(\ell-1)t$ has the same class-group character
index. The carried comparison (K2), valid in all positive
degrees $K_{4k}$, proves (SP.13). No large-index Bernoulli
enumeration is claimed to have been rerun in this section.
\end{proof}

\paragraph{Limits of the conclusion.}
The theorem gives an explicit, enormous sufficient weight
bound; no smallest weight threshold is computed. It gives actual
stable cuspidal automorphic systems with global residual
congruences, not just an eigenvalue coincidence in an
endoscopic summand. It gives neither a nonzero rational
Selmer class nor automatic membership in the stricter
archive admissible-source predicate. It leaves open the
possibility of eliminating the prime by a different
auxiliary field or by the already successful reflected
class-group test. For this fixed field and the entire
stable compact space, however, an all-weight unit claim
is false even when all split test primes are available.


% END evidence/main_gate/stable_residual_propagation.tex


% BEGIN evidence/main_gate/ordinary_symmetric_square_admissibility.tex
\section{Ordinary symmetric-square congruences in every weight of a progression}
\label{sec:kt-ordinary-symmetric-square}

Put $E=\mathbf Q(\sqrt{-3})$ and $\epsilon=\chi_{-3}$.
The cyclotomic character $\varepsilon_\ell$ below is normalized
by $\varepsilon_\ell(\operatorname{Frob}^{\rm arith}_q)=q$.
This section strengthens KT-U3-SYM2-EISENSTEIN by verifying
the ordinary local flags and characteristic-zero irreducibility.
Its last source-predicate implication is stated separately so
that it can be checked against the literal definitions.

\begin{theorem}[KT-ORDINARY-SYM2-FAMILY]
Let $e_0\geq2$ be even and let $\ell>3$ be prime. Suppose
\[
 v_\ell\!\left(\frac{B_{e_0+1,\epsilon}}{e_0+1}\right)>0,
 \qquad \ell\nmid1-2^{e_0}.                         \tag{OS.1}
\]
For each $t\geq0$ put $e=e_0+(\ell-1)t$. There is a
non-CM primitive normalized eigenform
\[
 f_t\in S_{e+1}(\Gamma_0(3),\epsilon)
\]
and a coefficient prime $\lambda\mid\ell$ such that
\[
 a_q(f_t)\equiv1+\epsilon(q)q^e\pmod\lambda
                       \quad\text{for every prime }q\ne3.
                                                               \tag{OS.2}
\]
In particular $a_\ell(f_t)$ is a $\lambda$-adic unit.
Let $r_f$ be its arithmetic-normalized two-dimensional
$\ell$-adic representation and set
\[
 \rho_t=\left(\operatorname{Sym}^2r_f\otimes
                         (\det r_f)^{-1}\right)|_{G_E}.
                                                               \tag{OS.3}
\]
Then $\rho_t$ has all of the following properties.
\begin{enumerate}
\item It is continuous and absolutely irreducible over
      $\overline{\mathbf Q}_\ell$, with determinant one.
\item It is unramified at every finite place of $E$ away from
      $\ell$. At every place above $\ell$ it has an actual
      upper triangular flag, in the displayed order, with
      inertia characters
      $\varepsilon_\ell^e,1,\varepsilon_\ell^{-e}$.
\item For the three diagonal characters $\kappa_i$ on any
      such decomposition group and all $i,j$ and $m\geq1$,
      $\kappa_i^m\ne(\kappa_j\varepsilon_\ell)^m$.
      The same assertion holds after twisting the whole
      representation by $\varepsilon_\ell^{-1}$.
\item It is the weight-zero system of the everywhere
      finite-unramified, regular algebraic, conjugate
      self-dual cuspidal representation
      $\Pi_f=\operatorname{BC}_{E/\mathbf Q}
                    (\operatorname{Ad}(\pi_f))$.
      The normalized system used in \texttt{AttC} is
      $\rho_t\varepsilon_\ell^{-1}$.
\item For a stable lattice its residual semisimplification is
\[
       \overline\rho_t^{\,\rm ss}
           \simeq\overline\varepsilon_\ell^e\oplus1\oplus
                           \overline\varepsilon_\ell^{-e}
                      \quad\text{on all of }G_E.              \tag{OS.4}
\]
\end{enumerate}
Consequently $\rho_t$ satisfies the literal predicate
$\mathrm{Level}(\ell,2,(e,0,-e))$, including its ordinary
flag and nonadjacency alternative. The representation is
also a stable compact $U_3$ eigenline in the constructed
symmetric-square subspace of KT-U3-SYM2-INJECT.
\end{theorem}

\begin{proof}
\paragraph{The infinite progression and ordinarity.}
Write $w=e+1$ and
$\beta_w=B_{w,\epsilon}/w$. Since $w$ is odd and $\ell-1$
is even, $\ell-1\nmid w$. Sun, \emph{Congruences concerning
Bernoulli numbers and Bernoulli polynomials}, Lemma~8.1(a),
printed pp.~219--220, applied with $n=1$, gives
\[
 (1-\epsilon(\ell)\ell^{w-1})\beta_w
 \equiv
 (1-\epsilon(\ell)\ell^{w+\ell-2})\beta_{w+\ell-1}
                         \pmod\ell.                         \tag{OS.5}
\]
The conductor is three, prime to $\ell$. Each Euler factor
is an $\ell$-adic unit since $w\geq3$. Starting at (OS.1)
and applying (OS.5) successively proves $v_\ell(\beta_w)>0$
for every weight in the progression. The lemma concerns
the normalized Bernoulli numbers themselves and remains
valid when $w$ is divisible by $\ell$.

Fretwell--Roberts, \emph{Newform Eisenstein congruences of
local origin}, Theorem~3.6, printed pp.~517--518, now applies
with $N=3$, $M=1$, $k=w$, $\psi=1$, $\phi=\epsilon$.
Its special value is $L(1-w,\epsilon)=-\beta_w$;
its coefficient-prime exclusion is just $\ell\nmid6$.
There is no hypothesis $\ell>k+1$ in this sufficient theorem.
It gives one eigenform and one coefficient prime satisfying
(OS.2) simultaneously for every prime $q\nmid NM$.
In particular $q=\ell$ is included, so $a_\ell\equiv1$.
Primitivity of the nebentypus makes the eigenform new, as
explained immediately after that theorem's proof.

All CM forms in this level-three space are CM by $E$, by
KT-GL2-LOWER-SCOPE. At the inert prime two they have $a_2=0$.
But Fermat's theorem and (OS.1) give
$a_2\equiv1-2^e\equiv1-2^{e_0}\ne0\pmod\lambda$.
Thus every chosen $f_t$ is non-CM.

\paragraph{The ordinary flag, with its orientation.}
The ordinary local theorem of Mazur--Wiles and Wiles says
that $a_\ell$ a unit gives an upper triangular realization
\[
        r_f|_{G_{\mathbf Q_\ell}}\sim
           \begin{pmatrix}\delta&*\\0&\eta\end{pmatrix},
       \qquad \eta\text{ unramified},\qquad
       \delta\eta=\epsilon\varepsilon_\ell^e.                \tag{OS.6}
\]
An exact published statement of the general theorem, with
this orientation, is Ghate--Vatsal, \emph{On the local
behaviour of ordinary $\Lambda$-adic representations},
printed p.~2144, equation~(1.1) and its preceding paragraph.
No residual irreducibility assumption is imposed on this
ordinary assertion. Their later splitting theorems impose
additional assumptions; none is used here.

If $v_1,v_2$ is the basis in (OS.6), the ordered basis
$v_1^2,v_1v_2,v_2^2$ gives an upper triangular symmetric
square with diagonal $\delta^2,\delta\eta,\eta^2$.
Dividing by $\delta\eta$ yields the actual flag of (OS.3)
with diagonal
\[
                    \delta/\eta,\quad1,\quad\eta/\delta.
                                                               \tag{OS.7}
\]
Because $\ell\ne3$, the character $\epsilon$ is unramified
at $\ell$; in fact its restriction to $G_E$ is trivial.
On inertia (OS.7) is exactly
$(\varepsilon_\ell^e,1,\varepsilon_\ell^{-e})$.
Restriction to each completion of $E$ preserves the flag.
Conjugation transports it to every choice of a place in the
algebraic closure; no splitting of the flag is asserted.

Let $h=(e,0,-e)$. If
$\kappa_i^m=(\kappa_j\varepsilon_\ell)^m$ for $m\geq1$,
restriction to inertia gives
$\varepsilon_\ell^{m(h_i-h_j-1)}=1$.
The cyclotomic image of inertia in a finite extension of
$\mathbf Q_\ell$ is infinite. Hence $h_i-h_j=1$ would
follow. This is impossible since
$h_i-h_j\in\{0,\pm e,\pm2e\}$ and $e\geq2$.
Twisting all characters by $\varepsilon_\ell^{-1}$ leaves
their differences unchanged. This proves every asserted
positive-power nonlink, not just the case $m=1$.

\paragraph{Characteristic-zero irreducibility at every prime.}
Ribet, \emph{On $\ell$-adic representations attached to modular
forms II}, printed p.~191, immediately before Theorem~3.1,
recalls Momose's theorem that, for each $\ell$, the image
on the fixed inner-twist-kernel subgroup is open in the
reduced-norm group of the specified quaternion algebra.
The original reference is F.~Momose, \emph{On $\ell$-adic
representations attached to modular forms}, J.~Fac.~Sci.\
Univ.~Tokyo Sect.~IA Math.~28 (1981), 89--109.
This is the all-prime input used here; Ribet's Theorem~3.1
itself is an almost-all-primes statement and is not used
as an all-prime theorem.
After extension of coefficients to $\overline{\mathbf Q}_\ell$,
the projection to the chosen coefficient embedding contains
$\mathfrak{sl}_2$. This follows as well by taking the
commutator of the central simple algebra in that theorem.
Restriction to the finite-index subgroup $G_E$ does not
change the Lie algebra of the image. The representation
$\operatorname{Sym}^2\otimes\det^{-1}$ of
$\mathfrak{sl}_2$ is its irreducible three-dimensional
adjoint representation in characteristic zero. Therefore
any invariant subspace of (OS.3) is zero or the whole space.
This proves absolute irreducibility even when the residual
representation is reducible. Ribet's separate almost-all-
primes integral-image theorem is not substituted for this
all-prime Lie algebra theorem.

\paragraph{The global and residual assertions.}
KT-U3-SYM2-INJECT applies to every even $e\geq2$ and every
non-CM $f_t$. It proves the cuspidal base change, its regular
algebraic conjugate self-duality, all finite unramified local
parameters, and the stable compact member. Its local check
at three uses inertia $1\oplus\epsilon$ and zero monodromy
for the classical form; restriction to $E_3$ kills
$\epsilon$. Classical local compatibility therefore makes
(OS.3) unramified at three as well as at every other place
away from $\ell$.

The weight-zero adjoint is matched to $\Pi_f$ by its good
Frobenius polynomials; its three Satake roots are invariant
as a multiset under inversion. The definition \texttt{AttC} uses
geometric Frobenius and the factor $q^{(3-1)/2}=q$ in
\texttt{algSat}; twisting the weight-zero system by
$\varepsilon_\ell^{-1}$ supplies exactly this factor.
Thus this is the automorphic normalization required by
\texttt{Adm3Alt}.

Continuity and a finite coefficient field give a stable
lattice. The determinant of $r_f$ is
$\epsilon\varepsilon_\ell^e$. The simultaneous congruence
(OS.2), Chebotarev, and Brauer--Nesbitt give
\[
       \overline r_f^{\,\rm ss}
              \simeq1\oplus\epsilon\overline\varepsilon_\ell^e.
\]
Restriction to $G_E$ kills $\epsilon$. Taking the symmetric
square and dividing by the determinant gives the three
composition factors in (OS.4); equivalently, use the flag
of the endomorphisms and remove the scalar line. Since
$\ell>3$, scalar and trace-zero endomorphisms split.

Finally, at level parameter $2$ choose all the global finite
characters in \texttt{LevelAway} and \texttt{LevelAt} to be
trivial. They satisfy \texttt{KillsMu 2}. Unramifiedness
gives the required diagonal inertia away from $\ell$;
(OS.7) gives \texttt{IsFlagOn} above $\ell$ and the required
inertia powers in the right order. The case $m=1$ of the
nonlink argument supplies the first disjunct of
\texttt{LevelAt}. No potentially crystalline alternative or
Fontaine--Laffaille range is needed.
\end{proof}

\begin{corollary}[KT-ORDINARY-SYM2-VANISHING-PROGRESSIONS]
Under the stipulated archive-theorem convention used in
KT-K52-K56-K60-VANISH, for every row
\[
\begin{array}{c|r}
e_0&M_{e_0}\\\hline
26&697446897311\\
28&323026426255399\\
30&236017358500033404991
\end{array}                                                   \tag{OS.8}
\]
and every prime $\ell\mid M_{e_0}$, all the conclusions of
KT-ORDINARY-SYM2-FAMILY hold for
$e=e_0+(\ell-1)t$, $t\geq0$. In addition,
\[
                  K_{2e}(\mathbf Z)\{\ell\}=0
                           \qquad(t\geq0).                   \tag{OS.9}
\]
\end{corollary}
\begin{proof}
The exact values in KT-U3-SYM2-EISENSTEIN, or the independent
rational-arithmetic certificate
\path{lower_bernoulli_certificate.json}, give
$M_{e_0}\mid\operatorname{num}(B_{e_0+1,\epsilon}/(e_0+1))$
and
$\gcd(M_{e_0},6\operatorname{den}(B_{e_0+1,\epsilon}/(e_0+1))
 (1-2^{e_0}))=1$.
Thus each prime divisor satisfies (OS.1) and is greater than
three. No primality assumption about the displayed moduli
is required.

KT-K52-K56-K60-VANISH gives $K_{2e_0}(\mathbf Z)=0$.
The carried comparison (K2), with its all-degree scope
K17-K2-PERIODIC, identifies the vanishing criterion with the
same cyclotomic class-group character indexed by $-e$.
Since $e\equiv e_0\pmod{\ell-1}$, that character and hence
the criterion are unchanged. This proves (OS.9). The
argument uses periodicity of the vanishing criterion, not
an isomorphism between $K$-groups in different degrees.
\end{proof}

\begin{corollary}[KT-ORDINARY-SYM2-ADM3]
Under the stipulated archive-theorem convention, the
representations in KT-ORDINARY-SYM2-FAMILY satisfy the actual
predicate
\[
                 \mathrm{Adm3U\_s}(\iota,e,2,\rho_t).          \tag{OS.10}
\]
Here $\iota:\overline{\mathbf Q}_\ell\simeq\mathbf C$ is
chosen compatibly with the classical eigenform and its
attached automorphic representation. This includes the
stronger \texttt{WitA} and the added \texttt{Th.Unl}
conditions, not just ordinary automorphy and level.
\end{corollary}
\begin{proof}
At the pinned source revision
\path{7cbf7a0c32fca1badf05931879a37c3f8a2b03b4},
\texttt{Adm3Alt} requires \texttt{IsIrr},
\texttt{IsAutCsdA} for $\rho\varepsilon_\ell^{-1}$, and
the precise \texttt{Level} proved above.
\path{Adm3U_s} adds \texttt{Th.Unl}.
Write $r=\rho_t\varepsilon_\ell^{-1}$ and let
$\pi_0=\Pi_f$. We have \texttt{AutCsd0}$(\iota,r)$ by the
regular algebraic conjugate self-dual cuspidal lift and
the normalized attachment already proved.

\paragraph{The first clause of \texttt{WitA}.}
The diagonal inertia exponents of $r$ are
$(e-1,-1,-e-1)$. The proved positive-power nonlink gives
\texttt{PotUnl}$(r,v)$ at every $v\mid\ell$, for every
choice of the valuation above $v$. It therefore proves,
in particular, the implication
\texttt{Unl1}$(r,v)\Rightarrow$\texttt{PotUnl}$(r,v)$
required by the first clause of \texttt{WitA}. The closed
source theorem \texttt{potOfFlag} gives exactly the same
deduction: take its exponent tuple to be
$(e-1,-1,-e-1)$ and its auxiliary power to be one.

\paragraph{The universally quantified window clause.}
Every finite local component of $\pi_0$ is spherical.
The closed theorem \texttt{sphUnl} consequently supplies
toral support and unlinked characters at every such place.
The closed theorem \texttt{winOf}, applied once to $\pi_0$,
gives a single real number $u$ such that
\[
                  \mathrm{WinAt}(\pi_0,u,v)
                        \quad\text{for every finite }v.     \tag{OS.11}
\]
The same $u$ works at every place: the existential real
parameter in \texttt{WinOf} precedes its place quantifier.
This uses the exact strict width-one window of the source.
It does not require an additional Ramanujan assertion.

Now let $\pi$ be any regular algebraic cusp representation
with \texttt{AttC}$(\iota,r,\pi)$. Its attachment and that of
$\pi_0$ give the same characteristic polynomial at all but
finitely many places. Both use the same nonzero $q$-factor
in \texttt{algSat}; uniqueness of the roots with
multiplicity gives the same Satake multiset there.
The closed theorem \texttt{suppSmo}, whose proof applies
\texttt{strongMultOneC}, transports (OS.11) to $\pi$.
Thus the same window holds at every place for this $\pi$,
and in particular at every place selected by the disjunction
in the second clause of \texttt{WitA}. This proves its
universal quantifier over attached cusp representations.

Since \texttt{IsAutCsdA} is precisely \texttt{AutCsd0}
together with the implication from continuity to
\texttt{WitA}, the preceding paragraphs prove that
predicate for $r$.

\paragraph{The extra unlinked flag and conclusion.}
The definition \texttt{FlagAt} asks for the flag, the
inertia powers and global finite characters killed by
$\mu_2$. Choose those characters to be one and use (OS.7).
The case $m=1$ of the proved nonlink is the remaining
inequality in \texttt{Th.Unl}. Thus \texttt{Th.Unl}
holds for $\rho_t$ and $(e,0,-e)$.
Combining this with absolute irreducibility,
\texttt{IsAutCsdA} for $r$, and the proved \texttt{Level}
constructs \texttt{Adm3Alt} and then
\path{Adm3U_s}, exactly as defined.
\end{proof}

\begin{corollary}[KT-ORDINARY-SYM2-GD]
For every row of (OS.8), every prime $\ell\mid M_{e_0}$
and every $t\geq0$, the parameters $E$, $N=2$, trivial
Dirichlet character $\psi$, and $e=e_0+(\ell-1)t$ satisfy
the actual predicate \texttt{GD} in \path{OneCharU_s}.
\end{corollary}
\begin{proof}
The already checked rational specialization in
\path{SmallClassNumbers/ClassNumberThree.lean},
lines 64--86, proves \texttt{GD} for every positive even
$e$ and prime $\ell\geq5$ with $(\ell-1)\nmid2e$.
It has no additional inequality comparing $e$ with $\ell$.
For the three rows of (OS.8), exact integer arithmetic gives
\[
                   \gcd\bigl(M_{e_0},(2e_0+1)!\bigr)=1.
\]
Thus each prime divisor satisfies $\ell-1>2e_0>0$.
Since $2e\equiv2e_0\pmod{\ell-1}$, its nondivisibility
condition persists for every $t$. The exact gcd outputs
are saved in \path{ordinary_family_gd_certificate.json}.
\end{proof}

\begin{corollary}[KT-ORDINARY-SYM2-EXPLICIT-LARGE]
The integer $\ell=697446897311$ is prime and exceeds
$2^{31}$. For every $t\geq0$, put
\[
                  e=26+697446897310\,t.
\]
Then the actual \texttt{GD} parameters of
KT-ORDINARY-SYM2-GD admit a representation satisfying
\path{Adm3U_s}, the global residual triple (OS.4),
and $K_{2e}(\mathbf Z)\{\ell\}=0$.
Thus full-field unit exclusion fails at a fixed prime
above $2^{31}$ in unbounded weights within the literal
source hypotheses of \path{OneCharU_s}.
\end{corollary}
\begin{proof}
All assertions except primality follow from the preceding
three corollaries. Here is an exact primality certificate
independent of trial division of the target $\ell$.
Write $n=697446897311$. Then
\[
              n-1=2\cdot5\cdot31\cdot1019\cdot2207879.
\]
Each displayed factor is prime: exhaustive division up to
its integer square root suffices, and the largest such
bound is $1485$. For the following witnesses $a_q$, exact
modular exponentiation gives $a_q^{n-1}\equiv1\pmod n$
and the displayed $b_q=a_q^{(n-1)/q}\bmod n$:
\[
\begin{array}{r|r|r}
q&a_q&b_q\\\hline
2&13&697446897310\\
5&2&682051080219\\
31&2&34767638009\\
1019&2&273082850090\\
2207879&2&201208097879
\end{array}
\]
In each row $\gcd(b_q-1,n)=1$.
For any prime $p\mid n$, these conditions imply that the
order of $a_q$ modulo $p$ contains the full $q$-power of
$n-1$. Since this holds for every prime $q\mid n-1$,
we obtain $n-1\mid p-1$. Thus $p\geq n$, forcing $p=n$.
This proves primality. The inequality $n>2^{31}$ is
immediate (the difference is $695299413663$).
The complete standard-library verification is
\path{ordinary_large_prime_certificate.py}, with
exact output saved beside it.
\end{proof}

\paragraph{Exact source locators for (OS.10).}
All paths here are relative to
\path{src/Development/AutomorphicArithmetic/}.
The predicates \texttt{IsFlagOn}, \texttt{LevelAway},
\texttt{LevelAt}, and \texttt{Level} are in
\path{AutomorphicVarieties/AutomorphicForms/%
BlochKatoSelmerVanishes.lean}, lines 37--39,
89--92, 102--107, and 111--114;
\texttt{FlagAt} is in the same directory's
\path{LocIsotypic.lean}, lines 31--34.
\texttt{AutCsd0} is in
\path{Foundations/GaloisRepresentations/AdmissibleResidues/%
IsAutomorphicOverImagQuad.lean}, lines 35--36.
Inside \path{AutomorphicVarieties/ModularityArgument/},
\path{AutomorphyLifting/WeakOfAtt.lean}, lines 19--30,
defines \texttt{Unl1}, \texttt{PotUnl}, and \texttt{WitA};
\path{AutomorphyLifting/IsAutCsdA.lean}, lines 11--12,
defines \texttt{IsAutCsdA}.
The definitions \texttt{WinAt}, \texttt{SuppSmo},
\texttt{WinOf}, and \texttt{SphUnl} are respectively at
lines 38--41, 142--145, 158--164 of
\path{LeviSmoothLocalization/WinUpDown.lean}.
In its \path{PrincipalSeriesMatching/} subdirectory,
\path{PotOfFlag.lean}, lines 88--99,
\path{SphUnl.lean}, lines 25--45,
\path{NormAbsPowUnif.lean}, lines 124--168, and
\path{SuppSmo.lean}, lines 6--12, contain the closed
theorems just used, with their complete proofs.
Finally \path{AutomorphicVarieties/TheoremsMuLinked/%
AutP.lean}, lines 10--13,
defines \texttt{Th.Unl}; the
\path{ModularityArgument/ThresholdArguments/%
DenseThresholdCases/} files \path{AfStmt.lean},
lines 29--32, and \path{NumberU3US.lean}, lines 27--29,
define \texttt{Adm3Alt} and \path{Adm3U_s}.

\paragraph{Consequence and scope.}
For every representation in (OS.10), (OS.4) makes
\[
 \operatorname{tr}\rho_t(\sigma)
  -\bigl(\varepsilon_\ell(\sigma)^e+1+
                         \varepsilon_\ell(\sigma)^{-e}\bigr)
\]
belong to the coefficient maximal ideal for every
$\sigma\in G_E$. Thus varying a single group element,
or any number of good-prime Hecke operators, cannot make
this trace difference a unit for all members of the
actual admissible rank-three comparison class.
The families (OS.8)--(OS.9) show that this obstruction
coexists with the vanishing of the indicated rational
$K$-group component.
Membership in \path{Adm3U_s} is not by itself a converse
to the source's map from deformation-ring points.
The next section supplies the additional argument:
KT-ADJOINT-ACTUAL-HOST-POINT realizes these representations
as actual Type-I and TwoPt points, so every literal
\texttt{CoversPt} cover contains their isomorphism class.
This stronger conclusion still does not assert that every
possible existential witness set in \path{OneCharU_s}
must contain them.

\paragraph{Primary inputs and check scope.}
Sun's Lemma~8.1(a), including its displayed proof, and
Fretwell--Roberts Theorem~3.6, including its complete proof,
were already read and independently checked in the preceding
family audit. The present ordinary input is checked at the
exact published statement of Ghate--Vatsal, p.~2144,
equation~(1.1); the original Wiles 1988 proof is an imported
theorem and is not claimed to have been reread. The exact
Momose open-image statement recalled by Ribet on p.~191
is the imported all-prime theorem; no reread of Momose's
original proof is claimed. An earlier draft incorrectly
called this Ribet Theorem~0.1 and was corrected after the
independent visual source check. The global lift is the independently
checked current-run KT-U3-SYM2-INJECT. Source PDFs and
hashes for the two new local/irreducibility inputs are
saved under \path{evidence/main_gate/literature/}.

% END evidence/main_gate/ordinary_symmetric_square_admissibility.tex


% BEGIN evidence/main_gate/level_variation_scope.tex
\subsection{Changing the allowed level retains the actual ordinary obstruction}

\begin{corollary}[KT-ORDINARY-SYM2-ALL-LEVELS]
Let $\rho_t$, $e$ and $\iota$ be as in
KT-ORDINARY-SYM2-ADM3. For every positive integer $N$,
\[
                  \mathrm{Adm3U\_s}(\iota,e,N,\rho_t)
\]
holds with the same representation. Its trace difference from
$\varepsilon_\ell^e+1+\varepsilon_\ell^{-e}$ is nonunit
at every element of $G_E$, exactly as at $N=2$.
This statement concerns admissibility; it asserts neither
the predicate GD nor membership in every possible existential
comparison set for the changed parameter $N$.
\end{corollary}
\begin{proof}
In the pinned source,
\path{BlochKatoSelmerVanishes.lean}, lines 66--67,
\texttt{KillsMu}$(N,\xi)$ means that $\xi$ is trivial on
every automorphism fixing all $N$th roots of unity. It is
not defined by saying that the order of $\xi$ divides $N$.
The constant character $1$ satisfies this exact predicate
for every $N$.

All global finite characters used in the proof of
KT-ORDINARY-SYM2-ADM3 are constant $1$. The same choices
therefore verify \texttt{LevelAway} and \texttt{LevelAt}
for every $N$: their other clauses, namely continuity,
unramified inertia, the actual ordinary flag, the ordered
cyclotomic inertia powers, and nonadjacency, are unchanged.
The source locators are lines 89--92, 102--107 and
111--114 of that file. The extra \texttt{Th.Unl}
predicate, \path{TheoremsMuLinked/AutP.lean}, lines
10--13, uses \texttt{FlagAt}; in
\path{AutomorphicForms/LocIsotypic.lean}, lines 31--34,
its only occurrence of $N$ is again \texttt{KillsMu}$(N,\xi)$.
Thus the same flag and trivial characters verify it.

The irreducibility and automorphy clauses in
\path{DenseThresholdCases/AfStmt.lean}, lines 29--32,
do not contain $N$. The definition of \path{Adm3U_s}
in \path{NumberU3US.lean}, lines 27--29, only adds the
unlinked predicate already checked. This proves the stated
admissibility for every positive $N$. Finally the global
residual semisimplification of $\rho_t$ is independent
of $N$, so (OS.4) gives the asserted nonunit trace difference.
\end{proof}

This rules out a specific level-variation shortcut: allowing
a different level while demanding a unit trace difference
for every representation in the entire admissible class
does not remove these systems. Changing the field, imposing
a genuinely smaller comparison class by an additional
source theorem, or controlling extension modules is a
different arithmetic operation. No conclusion about those
operations follows from this corollary.

% END evidence/main_gate/level_variation_scope.tex


% BEGIN evidence/unitary_audit/host_sign_results.tex
\section{The Selmer sign and actual ordinary adjoint points}
\label{sec:kt-host-sign}

This section uses the same pinned-source convention as
KT-ONECHAR-CLOSED.  The new local argument below is a mathematical
proof, not a claimed Lean compilation.  Put
$E=\mathbf Q(\sqrt{-3})$, $\eta=\chi_{-3}$, and
$\Psi=\varepsilon_\ell^e$, with the arithmetic cyclotomic
normalization.  Assume $\ell>3$, $e\geq2$ is even, and
$(\ell-1)\nmid2e$.  These are the relevant rational trivial-character
\texttt{GD} conditions.  Write $J$ for the anti-diagonal matrix with
all three anti-diagonal entries equal to one.  The source uses
\[
 \theta(g)=Jg^{-t}J,\qquad
 GG_3=\mathrm{GL}_3\rtimes\{\pm1\},\qquad
 r_0(g)=(\operatorname{diag}(\Psi(g),1,\Psi(g)^{-1}),\eta(g)).
                                                               \tag{HS.1}
\]
Here $-1$ acts by $\theta$.  Matrix units in this section use
indices $1,2,3$; the Lean files use $0,1,2$.

\begin{lemma}[KT-HOST-TANGENT-SIGN]
The weight $-1$ space for the adjoint action of $r_0$ splits as
\[
 \langle E_{21},E_{32}\rangle
  =\langle E_{21}-E_{32}\rangle
       \oplus\langle E_{21}+E_{32}\rangle,
                                                               \tag{HS.2}
\]
with global characters $\Psi^{-1}$ and $\eta\Psi^{-1}$,
respectively.  The rational strict Selmer injection used by
\texttt{ShaInjGen} takes its values in the first line.  The lower
root of the normalized principal adjoint embedding takes its values
in the second line.
\end{lemma}
\begin{proof}
The differential of the involution is
$\theta_{\rm Lie}(X)=-JX^tJ$, so it interchanges
$E_{21}$ and $-E_{32}$.  Conjugation by
$\operatorname{diag}(z,1,z^{-1})$ multiplies both matrix units
by $z^{-1}$.  This proves the two characters in (HS.2).
In particular complex conjugation acts by $+1$ and $-1$ on
the two lines, because $\Psi(c)=1$ and $\eta(c)=-1$.

The literal definition \path{shaInj_mat} is
$t\mapsto t(E_{21}-E_{32})$; the following theorem
\path{shaInj_ad} proves its equivariance with character
$\Psi^{-1}$ for every global group element.
For the second assertion, use the ordered basis
\[
 E=\begin{pmatrix}0&1\\0&0\end{pmatrix},\qquad
 H=\begin{pmatrix}1&0\\0&-1\end{pmatrix},\qquad
 F=\begin{pmatrix}0&0\\1&0\end{pmatrix}
\]
of $\mathfrak{sl}_2$ and replace $H$ by $sH$, where
$s^2=-1/2$.  Then
\[
 \operatorname{ad}(F)=-s^{-1}E_{21}+2sE_{32}
                  =-s^{-1}(E_{21}+E_{32}).                    \tag{HS.3}
\]
The preserved trace form in this basis is
\[
 K=\begin{pmatrix}0&0&1\\0&-1&0\\1&0&0\end{pmatrix}.
                                                               \tag{HS.4}
\]
The construction below identifies this basis with the involution
in (HS.1), proving the asserted normalization.
\end{proof}

The distinction in (HS.2) is a distinction between tangent
directions.  To decide whether it excludes a characteristic-zero
point, one must check the actual point conditions.  In particular,
\texttt{TypeI} means irreducibility on $G_E$; it does not require
that a point arise by exponentiating the first line in (HS.2).
The next lemma addresses the local condition that is missing from
an argument using residual semisimplification alone.

\begin{lemma}[KT-ORDINARY-ON-BASE-COMPONENT]
Let $O$ be the ring of integers of any finite extension of
$\mathbf Q_\ell$,
and let $D_{\rm loc}$ be the source's local data with unramified
$\eta$, trivial $\tau$, and $\Psi|_I=\varepsilon_\ell^e$.
Let $U$ be its unrestricted framed universal ring.  Assume the
carried local component theorem
\[
 \mathrm{primesAt0}(D_{\rm loc},U)=\{I_t\}.                  \tag{HS.5}
\]
Let $A$ be a complete Noetherian local $O$-algebra with the same
residue field.  Suppose an integral lift of $r_0\bmod\mathfrak m_O$
to $GG_3(A)$ is conjugate, by a matrix congruent to the identity,
to a lift valued in $B\rtimes\{\pm1\}$, where $B$ is the standard
upper Borel, and its diagonal character on inertia is exactly
$\operatorname{diag}(\varepsilon_\ell^e,1,
\varepsilon_\ell^{-e})$.  Then its map $U.R\to A$ kills $I_t$.
This holds whether the quadratic local extension is split or
unramified quadratic.
\end{lemma}

\begin{proof}
\paragraph{A smooth integral family with the prescribed inertia.}
First consider the framed functor of lifts valued in
$B\rtimes\{\pm1\}$, with the projection $\eta$ fixed and the
specified torus character on inertia fixed.  It is prorepresented
by a complete Noetherian local $O$-algebra: this is the usual
framed lifting construction, imposing the Borel and inertial
matrix equations.  Its finite tangent space follows from local
Galois cohomological finiteness.  The unipotent radical $N$,
its Lie algebra $\mathfrak n$, the Borel Lie algebra
$\mathfrak b$, and $\mathfrak t=\mathfrak b/\mathfrak n$
are all stable under $\theta$.

Here is the full lifting argument, including the unramified
quadratic case.  For a small surjection $A'\to A$ with kernel
$J$, the obstruction to lifting a homomorphism with fixed
projection $\eta$ is the usual two-cocycle in
$H^2(G_{\mathbf Q_\ell},\mathfrak b_{\bar r_0})\otimes J$.
The fixed outer action merely makes this the indicated twisted
adjoint module; it does not introduce an additional Frobenius
descent condition.  On inertia the composition factors of
$\mathfrak b$ have characters
\[
                   1,\quad\bar\varepsilon_\ell^e,
                         \quad\bar\varepsilon_\ell^{2e}.
                                                               \tag{HS.6}
\]
Every exponent is even.  The Tate dual of each factor has an
odd exponent, which is nonzero modulo the even number
$\ell-1$.  Local Tate duality therefore gives
\[
 H^2(G_{\mathbf Q_\ell},\mathfrak b_{\bar r_0})
   =H^2(G_{\mathbf Q_\ell},\mathfrak n_{\bar r_0})=0.
                                                               \tag{HS.7}
\]
This proves that a Borel lift exists over $A'$.  Its torus
restriction may differ from the required inertia character by
an element of $H^1(I,\mathfrak t)^{G/I}\otimes J$.
The correction map
\[
 H^1(G,\mathfrak b)\longrightarrow
 H^1(G,\mathfrak t)\longrightarrow
 H^1(I,\mathfrak t)^{G/I}                                  \tag{HS.8}
\]
is surjective: the first arrow is surjective by (HS.7), and
the second by $\operatorname{cd}_\ell(G/I)=1$.
Changing the lift by a one-cocycle consequently corrects its
inertia without changing the given reduction.  Here inertia acts
trivially on $\mathfrak t$, so equality of its cohomology classes
is equality of the actual torus correction cocycles; there is
no unaccounted inertia coboundary.  Thus the functor
is formally smooth and its representing ring $R_B$ is a power
series ring over $O$, in particular a domain.  No unramifiedness
of the coefficient DVR is used: (HS.6)--(HS.8) are statements
over its finite residue field, and formal smoothness is relative
to the given $O$.

For split $\eta$ this is precisely the lifting argument in
Patrikis, \emph{Deformations of Galois representations and
exceptional monodromy}, \S4.1, Proposition~4.4, preprint
pp.~12--13.  Its hypotheses \textnormal{(REG)} and
\textnormal{(REG*)} hold here: the negative-root inertia
characters are $\bar\varepsilon_\ell^{-e}$ and
$\bar\varepsilon_\ell^{-2e}$, nontrivial by
$(\ell-1)\nmid2e$, and their cyclotomic twists have odd
nonzero exponents.  The preceding obstruction calculation
provides the necessary adaptation for nonsplit unramified
$\eta$ directly.  In particular we have not substituted a
naive fixed-point condition for descent of a homomorphism.

Adjoin nine formal variables to $R_B$ and conjugate its
universal lift by the matrix $1+(X_{ij})$.  The resulting
domain $R'=R_B[[X_{ij}]]$ is a family containing every
congruently conjugate lift in the lemma.  It also contains
the exact base representation $r_0$.

\paragraph{The saturation and the chosen component.}
Recall the literal definitions in the source:
$I_Y=(I_{\rm ord}:\delta^\infty)$,
$I_{\rm cr}=\sqrt{I_Y}$, and $I_t$ is the intersection of
the minimal primes of $I_{\rm cr}$ lying under the base-point
kernel.  Here $I_{\rm ord}$ is generated by the prescribed
inertia characteristic polynomials and the entries of
\[
 (r(s_0)-\Psi(s_0))
 (r(s_1)-1)
 (r(s_2)-\Psi(s_2)^{-1}),\qquad s_i\in I.                   \tag{HS.9}
\]
Every expression in (HS.9) vanishes on an upper triangular
lift in the indicated order: the rightmost factor maps the
three-dimensional space to the second flag step, the next
maps that step to the first, and the leftmost kills the first.
The characteristic polynomial equations also vanish.  These
statements are preserved under conjugation, so the map
$\phi:U.R\to R'$ kills $I_{\rm ord}$.

The image of the saturation element $\delta$ in $R'$ is
nonzero.  To check this without assuming it is nonzero at
the base, specialize to the diagonal family
$r_0(\eta,\Psi u_T)$, where $u_T$ is unramified and sends
a chosen Frobenius to $1+T$.  If the source's chosen
$\mathrm{fr}$ has Frobenius exponent $m\geq1$, put
$y=\Psi(\mathrm{fr})(1+T)^m$ and
$v=\varepsilon_\ell(\mathrm{fr})$.  Writing $c_i$ for the
source's nonzero products of distinct inertia eigenvalue
differences, direct substitution gives
\[
 \delta=\ell\,c_0c_1^2c_2\,y^{-1}(y-v)^2\ne0
                                  \quad\text{in }O[[T]].     \tag{HS.10}
\]
The $c_i$ are nonzero because the chosen inertia element has
cyclotomic character of infinite order.  The function $y$ is
nonconstant and the coefficient ring has characteristic zero.
Formula (HS.10) is also the computation used in the archived
\path{del_unr_twist} proof.  Since $R'$ is a domain,
$\phi$ kills $I_Y$ and hence $I_{\rm cr}$.

Set $P=\ker\phi$.  It is prime.  The base specialization
$R'\to O$ shows $P\subseteq\ker x_0$, where $x_0$ is the
unique framed base-point map.  Choose a minimal prime $Q$ of
$I_{\rm cr}$ contained in $P$.  Such a prime exists since
$U.R$ is Noetherian.  Then $Q\subseteq\ker x_0$, so $Q$
belongs to \texttt{primesAt0}.  By (HS.5), $Q=I_t$.
Thus $I_t\subseteq P$.  Specializing the family to the lift
in the lemma proves the assertion.
\end{proof}

\begin{theorem}[KT-ADJOINT-ACTUAL-HOST-POINT]
Let $f\in S_{e+1}(\Gamma_0(3),\eta)$ be a non-CM primitive
eigenform and let $\lambda\mid\ell$ be a coefficient prime.
Suppose $a_\ell(f)$ is a $\lambda$-adic unit and
\[
 \bar r_f^{\,\mathrm{ss}}
       =1\oplus\eta\bar\varepsilon_\ell^e.                 \tag{HS.11}
\]
Keep the assumptions at the beginning of this section and
the carried local component theorem (HS.5).  For any of the
global data $D$, an element $g_1$ satisfying $D.\mathrm{Marked}(g_1)$,
local universal ring
$U$, and global ring $S=D.\mathrm{Ring}(U,g_1)$ used in the
rational trivial-character \path{OneCharU_s} construction,
there is a point
\[
                  z:S.R\longrightarrow\overline{\mathbf Q}_\ell
                                                               \tag{HS.12}
\]
satisfying the literal predicates \texttt{IsLocPt} and
\texttt{TwoPt}.  Its representation on $G_E$ is isomorphic
to $\operatorname{Ad}(r_f)|_{G_E}$ and is of \texttt{TypeI}.
Consequently it also satisfies the corresponding
$D.\mathrm{IsPt}$ conditions.  This conclusion concerns
actual points of the prescribed source ring, not only
membership in the ambient admissible class.
\end{theorem}

\begin{proof}
\paragraph{A split residual lattice, with its polarization.}
Use the cohomological dual $r_f^\vee$.  Its residual
semisimplification is $1\oplus\eta\bar\Psi^{-1}$; its
adjoint is isomorphic to that of $r_f$.  Choose a stable
lattice over a finite coefficient DVR.  Because two is a
unit and complex conjugation has distinct eigenvalues
$1,-1$, the lattice has a basis in which
$r_f^\vee(c)=\operatorname{diag}(1,-1)$.  The residual
representation is triangular in this basis: an invariant
residual line must be one of the two distinct $c$-eigenlines.
Thus one of the two off-diagonal matrix-entry ideals is
contained in the maximal ideal.

Enlarge the coefficient DVR by a square root of a uniformizer
$\pi$.  If the upper off-diagonal entries are arbitrary
integral and the lower ones lie in $(\pi)$, replace the
second basis vector by $\sqrt\pi$ times itself.  The new
upper entries are multiplied by $\sqrt\pi$ and the new
lower entries divided by it.  Both are now in the maximal
ideal.  In the opposite triangular orientation use the
opposite scaling.  This preserves the diagonal action of
$c$ and yields an integral representation whose reduction
is exactly diagonal.  This argument does not assume that
a split residual lattice existed over the original DVR.

Enlarge coefficients once more to contain a unit $s$ with
$s^2=-1/2$, and form the adjoint in the ordered basis
$(E,sH,F)$ of (HS.3).  Denote this representation by $A(g)$.
It preserves the perfect form $K$ of (HS.4), and
\[
 A(c)=D_c:=\operatorname{diag}(-1,1,-1),\qquad
 JK=-D_c,\qquad\theta(A(g))=D_cA(g)D_c.                    \tag{HS.13}
\]
If $j(g)\in\{0,1\}$ is defined by
$\eta(g)=(-1)^{j(g)}$, set
\[
                 r(g)=(A(g)D_c^{j(g)},\eta(g)).              \tag{HS.14}
\]
Equation (HS.13) proves directly that (HS.14) is a
homomorphism to $GG_3$.  It has $r(c).\mathrm{mat}=1$.
The residual adjoint matrix is
$\operatorname{diag}(\eta\bar\Psi,1,
\eta\bar\Psi^{-1})$, and multiplication by $D_c^{j(g)}$
cancels the two $\eta$ factors.  Hence the residual matrix
of $r$ is exactly $r_0$, entry by entry.

Let $O$ be the coefficient DVR of $D$, enlarge $O'$ if necessary
so that it contains $O$, and take
$A_0=O+\mathfrak m_{O'}$ inside the final coefficient DVR $O'$.
Every entry of $r$ and of its inverse lies in $A_0$, because
its reduction is in the residue field of $O$.
This order is finite over $O$, local with maximal ideal
$\mathfrak m_{O'}$, and complete; its residue field is exactly
that of $O$.  Indeed it is the inverse image of that residue
subfield, and its maximal-ideal topology agrees with its
$\ell$-adic topology.  Thus it is an admissible coefficient
algebra for the source's \texttt{IsCO} universal property.
The unit $s$ itself need not lie in $A_0$; the adjoint
matrices do, which is all that is required.

\paragraph{Compatibility of the ordinary flag and the involution.}
The ordinary local theorem gives for $r_f^\vee$ an actual
upper triangular flag ordered as an unramified character
followed by a character with inertia
$\varepsilon_\ell^{-e}$.  Its integral saturated line
reduces to the first standard line: the two residual
inertia characters are distinct by
$(\ell-1)\nmid2e$.  A matrix $q\equiv1\pmod{\mathfrak m_{O'}}$
therefore makes $q r_f^\vee q^{-1}$ upper triangular
on $G_{\mathbf Q_\ell}$.  Its adjoint $h=\operatorname{Ad}(q)$
in the basis above also satisfies $h\equiv1$, lies in
$\mathrm{GL}_3(A_0)$, preserves $K$, and obeys
$\theta(h)=D_chD_c$.  Consequently
\[
 (\operatorname{conj}_h r)(g)
       =(hA(g)h^{-1}D_c^{j(g)},\eta(g))                     \tag{HS.15}
\]
is valued in $B\rtimes\{\pm1\}$ locally at $\ell$.
Since $\eta$ is unramified there, its diagonal inertia
character is exactly
$(\varepsilon_\ell^e,1,\varepsilon_\ell^{-e})$.
This proves compatibility with the polarization, including
the inert quadratic case.  KT-ORDINARY-ON-BASE-COMPONENT
puts the local representation on $I_t$.

\paragraph{The global normalization and universal point.}
The representation $r$ kills the exact subgroup
$D.\mathrm{Kill}$: away from $\ell$, its restriction to
$G_E$ is unramified by the already checked local conductor
calculation for the adjoint lift, and membership in
$D.\mathrm{Kill}$ forces $\eta(g)=1$.  Continuity and all
integrality conditions hold over $A_0$.

Apply the source's \path{exists_normalised} theorem
over $A_0$ to the residual marked data.  It supplies a
matrix $v\equiv1$ such that $r'=\operatorname{conj}_v r$
is \texttt{Normalised} at $(c,g_1)$.  The local condition
proved above is preserved: locally $r'$ is still conjugate
by a matrix congruent to one to (HS.15), so the same local
lemma applies.  Killing $D.\mathrm{Kill}$ is also unchanged.
Thus $r'$ is an integral lift satisfying precisely
\texttt{NormCond(BoxCond)}.

The defining universal property of $S$ gives a local
$O$-algebra map $S.R\to A_0$ whose universal representation
is $r'$.  Composing with $A_0\hookrightarrow
\overline{\mathbf Q}_\ell$ gives (HS.12).
Every element of the maximal ideal of $S.R$ maps into
$\mathfrak m_{O'}$, so \texttt{IsLocPt} holds.
The representation on $G_E$ is isomorphic to
$\operatorname{Ad}(r_f)|_{G_E}$.  Its absolute
irreducibility is the all-prime non-CM adjoint
irreducibility proved in KT-ORDINARY-SYM2-FAMILY.
It is therefore of \texttt{TypeI} and can be neither
\texttt{Upper} nor \texttt{Lower}, proving \texttt{TwoPt}.
The source theorem \path{isPt_push} now gives
$D.\mathrm{IsPt}$ as well; alternatively each clause was
verified directly in the construction.
\end{proof}

\begin{corollary}[KT-ADJOINT-HOST-OBSTRUCTION]
Every ordinary symmetric-square family in
KT-ORDINARY-SYM2-GD gives actual \texttt{TypeI},
\texttt{TwoPt} points as in KT-ADJOINT-ACTUAL-HOST-POINT.
In particular, under the same carried arithmetic inputs,
for the certified prime $\ell=697446897311$ and
\[
                   e=26+697446897310\,t,\qquad t\geq0,
\]
these points coexist with $K_{2e}(\mathbf Z)\{\ell\}=0$.
Their trace difference from the Eisenstein triple is a
nonunit at every element of $G_E$.  Any finite set satisfying
the literal \texttt{CoversPt} condition contains a
representative of their isomorphism class.
\end{corollary}
\begin{proof}
KT-ORDINARY-SYM2-FAMILY supplies the ordinary forms and
(HS.11); KT-ORDINARY-SYM2-GD supplies
$(\ell-1)\nmid2e$.  Apply the preceding theorem.
The trace assertion follows from its entrywise residual
triple.  The asserted $K$-group vanishing and the certified
primality are KT-ORDINARY-SYM2-EXPLICIT-LARGE.
Finally \texttt{CoversPt} covers every \texttt{TwoPt}
up to torus conjugation, which preserves the isomorphism
class on $G_E$.
\end{proof}

\paragraph{What the sign distinction does and does not prove.}
The actual source points above rule out an exclusion based
only on the opposite signs in (HS.2).  The archive injects
the rational strict Selmer group into a larger weight
$-1$ deformation module and bounds that module using all
relevant points.  The presence of a point is not the
converse of a nonzero rational strict Selmer class.
A sharper argument would have to retain the first
eigensummand throughout the deformation-length comparison.
Neither a projector on its tangent space alone nor removal
of adjoint characteristic-zero points has been proved to
give such a comparison here.  Also, the existential
\path{OneCharU_s} statement may admit finite sets other
than the particular \texttt{CoversPt} set constructed in
its proof; the corollary makes no assertion about every
possible existential witness.

\paragraph{Exact source and proof-check scope.}
All Lean paths below are relative to the pinned
\path{AutomorphicVarieties/} directory.
\path{UnitaryThreeGroups/UnitaryThreeReflectionRight.lean}
defines $J$, $\theta$, $\theta_{\rm Lie}$, and $r_0$.
\path{ShaInjGen.lean}, lines 101--129, defines and checks
the Selmer matrix in (HS.2).
\path{BaseRepGL3NormalisedOfPsiEqOne.lean}, lines 15--24,
defines \texttt{TypeI}, \texttt{Upper}, and \texttt{Lower}.
\path{NormEx.lean}, lines 155 onward, proves universal
normalizability by a matrix congruent to one.
\path{GRepresents.lean}, lines 37--105, and
\path{ConjStable.lean}, lines 49 and 94--99, give the
universal ring and the literal point-cover definitions.
\path{RhoKApply.lean}, lines 61--83, defines
$D.\mathrm{Kill}$ and $D.\mathrm{IsPt}$; the one-way
\path{isPt_push} proof is in \path{IsPtOf.lean},
lines 124--166.

The local ideals are defined in
\path{PadicHodgeTheory/IsInertiaInfiniteCyclotomic.lean}.
The complete proof of (HS.5) is the carried
\path{component_basic} theorem in
\path{ModularityArgument/CrystallineOrdinaryIdeals/}
\path{OrdinaryComponentUniquePrimeDimThirteenNearBase.lean};
its local published inputs remain explicit dependencies.
The saturation calculation is checked against that
directory's \path{SaturatingElementUnramifiedTwistEqPowMulUnit.lean}.
The seemingly relevant \path{OnOrdinaryComponentIffContinuousCrystallineLift}
theorem is restricted to Artinian characteristic-zero
deformations of the exact base representation and is not
used as an arbitrary ordinary-point converse.
The separately named \path{UnitaryThreeDeformations/IsHostFrame}
fixes the Gaussian character and is not substituted for
the present Eisenstein-field data.

Patrikis's Proposition~4.4 and its entire proof were read
in the primary preprint, arXiv:1507.01294, pp.~12--13;
the elementary extension of its obstruction proof to
$B\rtimes\eta$ is written out above.  The original local
ordinary theorem for modular forms and the all-prime
non-CM image theorem remain the precisely named primary
inputs already recorded in KT-ORDINARY-SYM2-FAMILY.
The separate \path{host_sign_matrix_check.py/json}
certificate verifies sixteen matrix and scalar identities
over the rational function field with $2s^2+1=0$, including
both $\eta$ components of (HS.14), the residual cancellation,
the root signs, and (HS.10).  This finite algebra check is
not presented as a verification of the deformation-theoretic
or automorphic inputs.

% END evidence/unitary_audit/host_sign_results.tex


% BEGIN evidence/unitary_audit/projector_results.tex
\section{An integral involution projector and the remaining normal-length gate}
\label{sec:kt-projector}

This section separates a new integral conormal criterion from a
source-located limitation of the existing point-counting argument.
The subsequent KT-ODD-CONORMAL-SELMER-EQUIVALENCE identifies
that conormal with the dual of the original strict Selmer
module.  Its minor criterion is therefore an equivalent
formulation, with no arithmetic gate reduction claimed.
All source references use the pinned commit of KT-ONECHAR-CLOSED.
These are mathematical proofs and source inspections; no Lean
compilation is claimed.  The hypotheses and notation of
KT-ADJOINT-ACTUAL-HOST-POINT are retained.  In particular $O$ is
any finite coefficient DVR over $\mathbf Z_\ell$, $\ell>3$,
$e\geq2$ is even, $\eta=\chi_{-3}$, $\Psi=\varepsilon_\ell^e$,
and $(\ell-1)\nmid2e$.

\subsection{Part I: what the present length proof actually detects}

Write $\mathcal A_j$ for the source's torus weight spaces and
$x:A\to O$ for its base augmentation.  The literal definitions in
\path{CommutativeAlgebra/ResMapBijective.lean:245--268} are
\[
 \begin{split}
 \operatorname{satSq}(x)
   &=\{f:\text{ some }0\ne c\in O\text{ has }cf\in(\ker x)^2\},\\
 L_{-1}&=\operatorname{im}\bigl(\mathcal A_{-1}\cap\ker x
                     \longrightarrow A/\operatorname{satSq}(x)\bigr),\\
 C_A&=L_{-1}/\langle F_i\bmod\operatorname{satSq}(x)\rangle_O.
 \end{split}                                                    \tag{KP.1}
\]
Here $F_i$ are the negative weight-one equations.  The quantity
\texttt{cotLen} is $\operatorname{length}_O C_A$.  Saturation
makes the ambient lattice $L_{-1}$ torsion-free; it does not make
its quotient $C_A$ torsion-free.  The latter integral torsion is
precisely the quantity bounded.

The following is the exact chain used in the archived proof.
In the middle four entries, filenames are relative to
\path{CommutativeAlgebra/GradedModulesLength/}.
\begin{center}
\begin{tabular}{p{0.34\textwidth}p{0.59\textwidth}}
\textbf{Source locator} & \textbf{Operation relevant to the projector}\\\hline
\path{UnitaryThreeGroups/ModelLower.lean:296--369,374--398}
 & Constructs the normalized dual-number point from the strict
 Selmer cocycle, then lifts its derivation to the ambient model.\\
\path{Mprime.lean:38--80}
 & Injects the torsion input into
 $\operatorname{Hom}_O(C_A,O/(a))$.  It does not inject $C_A$
 into a generic point algebra.\\
\path{MainCotSpan.lean:32--70}
 & Takes the Smith normal form of the integral linearized
 negative equations, cuts the zero-weight variables, and invokes
 the Rees/Euler-characteristic estimate.\\
\path{ToDegInjective.lean:187--219}
 & Introduces $F_i^2$ and $E_i^2$ in the Rees sections and the
 auxiliary ideal $I_2$ in (KP.2).\\
\path{TwoSigned.lean:17--95}
 & Bounds localized lengths of the auxiliary quotient by
 $2^{r+1+r'}$ times generic orbit valuations.\\
\path{AlgMemCoef.lean:92--143}
 & Compares the lower contribution
 $2^{r+r'}\sum_i c_i$ with that upper bound, yielding the factor
 $2$ in the final estimate.\\
\path{EtHolds.lean:71--103,123--137}
 & Detects elements by characteristic-zero points only after
 every nonzero element of $O$ is assumed invertible.\\
\end{tabular}
\end{center}
The equations $F_i,G_i,E_i,E'_i,Z_i$ have weights
$-1,-2,1,2,0$, respectively.  The exact auxiliary ideal is
\[
        I_2=(F_i^2,\ G_i,\ E_i^2,\ E'_i,\ Z_i).              \tag{KP.2}
\]
Thus even if a chosen $F_i$ is odd under an involution,
its square is even, and the class of $F_i$ can remain an odd
nilpotent in the auxiliary quotient.  The localized-length
bound explicitly allows this thickening.  The generic reduced
orbit bound by one occurs before multiplying by the powers of
two which account for (KP.2).

More precisely, \path{EtHolds.et_holds} assumes
$\forall c\in O\setminus\{0\}$, \texttt{IsUnit (algebraMap O C c)}.
Its proof constructs an algebra structure over $\operatorname{Frac}O$
at lines 84--93 and obtains finiteness and reducedness over that
field at lines 99--103.  This is not a flatness or reducedness
assertion for the integral host ring.  None of the displayed
length statements has an involution as an argument.  The audited
chain therefore supplies no injection of the integral odd
conormal into a torsion-free reduced trace algebra.  This is a
limitation of the proposed refinement, not a defect in the
stated source inequality.

\begin{lemma}[KT-FIXED-GENERIC-POINTS]
Let $B$ be a finite torsion-free $O$-algebra with reduced generic
fiber, and let $\kappa$ be an $O$-algebra involution.  If every
geometric generic point is fixed by $\kappa$, then $\kappa=1$ on
$B$.  The analogous conclusion for the integral odd cotangent
of an arbitrary augmented complete local $O$-algebra does not
follow from its reduced generic point set.
\end{lemma}
\begin{proof}
The finite reduced algebra $B\otimes_O\operatorname{Frac}O$
is a product of finite separable fields, since the coefficient
field has characteristic zero.  Its geometric points separate
elements.  Hence $b-\kappa(b)$ vanishes in the generic fiber
for every $b\in B$, and torsion-freeness gives $b=\kappa(b)$.

For the second assertion let $\varpi$ be a uniformizer and, for
any $n\geq1$, set
\[
 R_n=O[[u]]/(\varpi^n u,u^2),\qquad
 x(u)=0,\qquad \kappa(u)=-u.                                \tag{KP.3}
\]
Give $u$ weight $-1$.  As an $O$-module,
$R_n=O\oplus (O/\varpi^n)u$.  Its generic fiber and its maximal
torsion-free reduced quotient are both $\operatorname{Frac}O$
and $O$, respectively; the unique generic point is fixed.
Nevertheless
\[
       (\ker x/(\ker x)^2)_{-1}^{-}
                    =(O/\varpi^n)u                           \tag{KP.4}
\]
has length $n$.  The example is a counterexample to the
generic-reduction inference alone; it is not claimed to satisfy
all the source hypotheses \texttt{M1--M4}.
\end{proof}

Even in a flat reduced algebra, an individual fixed point
can have nonzero integral odd cotangent when other generic
points are present.  For example
\[
 R'_n=O[[u]]/(u(u^2-\varpi^{2n})),\qquad \kappa(u)=-u          \tag{KP.5}
\]
is finite free of rank three over $O$, and reduced.  The generic
points $u=\varpi^n$ and $u=-\varpi^n$ are interchanged, while
$u=0$ is fixed.  At the augmentation $u=0$ its odd cotangent is
$O/(\varpi^{2n})$.  Indeed the only linear relation is
$-\varpi^{2n}u$, so its length is $2n$.  Reducedness follows
from the embedding into its reduced generic fiber, using
torsion-freeness.  This second example distinguishes a fixed
point from a branch with trivial integral normal conormal.
It does not refute a possible length bound counting only
nonfixed generic points: the two interchanged points could
account for the length $2n$ in this example.

\subsection{Part II: an involution on the actual normalized ring}

Keep $J$ and $K$ from (HS.1) and (HS.4), and put
$D=\operatorname{diag}(-1,1,-1)$.  The matrix automorphism
\[
 \kappa(M)=D\theta(M)D=K M^{-t}K                             \tag{KP.6}
\]
extends to $GG_3$ by leaving its sign component unchanged.
It commutes with $\theta$, fixes $r_0$, and commutes with
conjugation by $T(t)=\operatorname{diag}(t,1,t^{-1})$.
We first verify the local condition needed to use (KP.6) on
the source's integral functor.

\begin{lemma}[KT-ORDINARY-INTEGRAL-COMPONENT]
Under the hypotheses above and the carried local component
theorem, the source's local quotient $U.R/I_t$ represents the
framed ordinary functor with prescribed ordered inertial
characters $(\varepsilon_\ell^e,1,\varepsilon_\ell^{-e})$.
It is a power series ring in twelve variables over $O$.
Consequently $I_t$ is stable under $\kappa$.
\end{lemma}
\begin{proof}
The proof of KT-ORDINARY-ON-BASE-COMPONENT already establishes
formal smoothness of the fixed-inertia Borel lifting functor,
including unramified nonsplit $\eta$.  Its relative tangent
dimension is nine, as follows.  The twisted adjoint Borel Lie
algebra $\mathfrak b$ has dimension six and $H^2(G,\mathfrak b)=0$.
Local Euler characteristic over $G=G_{\mathbf Q_\ell}$ gives
$h^1(\mathfrak b)-h^0(\mathfrak b)=6$.  Therefore the framed
cocycle space $Z^1(G,\mathfrak b)$ has dimension twelve.
The map to $H^1(I,\mathfrak t)^{G/I}$ is onto by the proof of
(HS.8).  Since $\dim\mathfrak t=3$, $H^2(G,\mathfrak t)=0$,
and $\dim H^1(G/I,\mathfrak t)=h^0(\mathfrak t)$, the latter
space has dimension three.  Prescribing inertia leaves nine
parameters.  In the nonsplit case $h^0(\mathfrak t)=1$, whereas
in the split case it is three; the cancellation just used works
in both cases.  Thus $R_B\simeq O[[z_1,\ldots,z_9]]$.

Choose an inertia element $\tau$ for which the three residual
eigenvalues $\bar\Psi(\tau),1,\bar\Psi(\tau)^{-1}$ are distinct;
the stated GD condition ensures such an element.  The
characteristic polynomial of the unrestricted universal
matrix $r(\tau)$ has three simple residual roots.  Hensel
factorization in the unrestricted complete local ring itself
therefore gives three ordered roots and their spectral
idempotents.  The resulting ordered eigenflag lies in the big
cell around the standard flag.  For an ordinary lift it is the
unique invariant flag in that neighborhood: an invariant flag
is stable under $r(\tau)$ and its residual ordering specifies
which spectral lines it contains.

The big cell is uniquely parametrized by a lower unipotent
matrix congruent to one, with three coordinates.  Consequently
every framed ordinary lift is uniquely a conjugate of a Borel
lift by such a matrix.  This gives a representing ring
\[
              R_{\rm ord}=R_B[[z_{10},z_{11},z_{12}]].        \tag{KP.7}
\]
It also shows that this is a closed subfunctor of the
unrestricted framed functor.  Indeed use its canonical Hensel
eigenflag, require every $r(g)$ to preserve it, and require the
three diagonal inertia characters to have the prescribed
values.  These are closed matrix equations in the universal
ring.  For the sign $-1$ action, preservation means preservation
under the semidirect action on flags; conjugation into
$B\rtimes\{\pm1\}$ is the same condition because $\theta$
preserves the upper Borel.  Thus $U.R\to R_{\rm ord}$ is
surjective, with prime kernel $P$.

KT-ORDINARY-ON-BASE-COMPONENT gives $I_t\subset P$.  The carried
theorem \path{component_basic} states that $I_t$ is prime and
$\dim U.R/I_t=13$; see
\path{OrdinaryComponentUniquePrimeDimThirteenNearBase.lean:11--35}.
Equation (KP.7) gives $\dim U.R/P=13$.  A strict inclusion of
these primes would append one more prime to a chain of length
thirteen in $U.R/P$, a contradiction.  Hence $P=I_t$.

Finally $\kappa$ preserves $B$, reverses and inverts its three
diagonal characters, and therefore preserves their prescribed
ordered list.  It carries a congruent conjugating matrix to
another such matrix.  It thus preserves the represented
ordinary functor, proving the asserted stability of $I_t$.
\end{proof}

\begin{lemma}[KT-EQUIVARIANT-MARKED-SLICE]
The actual source ring
\[
 S=\operatorname{Univ}\bigl(\operatorname{NormCond}(c,g_1)
                        (\operatorname{BoxCond}(\mathrm{Kill},I_t))\bigr)
\]
admits an $O$-algebra involution $\kappa_S$, commuting with the
source torus action and fixing its base augmentation.  On
unmarked representations it is (KP.6); on the marked slice it
is this operation followed by a functorial normalization.
The ordinary adjoint points constructed above admit
$\kappa_S$-fixed representatives.  On the rational strict
Selmer dual-number directions, $\kappa_S$ acts by $-1$.
\end{lemma}
\begin{proof}
The local lemma, $\kappa(r_0)=r_0$, and $\kappa(1)=1$ show that
the unnormalized box conditions are stable.  Congruent
conjugation preserves them by the source's local conjugacy
theorem and the definition of the killed subgroup.

It remains to treat the actual marking, rather than assume
that its linearization suffices.  Work first with
$r(c).\mathrm{mat}=1$, as in the source's normalized functor,
and set $M=r(g_1).\mathrm{mat}$.  Since $\eta(g_1)=1$, the two
remaining marking equations are
\[
 Z_1(M)=M_{12}-M_{23}=0,\qquad
 Z_2(M)=M_{21}-M_{32}=0.                                    \tag{KP.8}
\]
This is the literal marking after
\texttt{MatGstar} and \path{NormEx.anti_shape}; indices here
are one-based.  Define genuinely odd functions
\[
           F_i(M)=\frac{Z_i(M)-Z_i(\kappa M)}2.              \tag{KP.9}
\]
They have torus weights $1$ and $-1$ and satisfy
$F_i(\kappa M)=-F_i(M)$.

Put $X=E_{12}-E_{23}$, $Y=E_{21}-E_{32}$, and use the Cayley
cross-section
\[
 W=xX+yY,\qquad
 H(x,y)=(1-W/2)^{-1}(1+W/2),\qquad x,y\in\mathfrak m_R.
                                                                  \tag{KP.10}
\]
Both $X$ and $Y$ lie in $\mathfrak{so}(J)$.
Consequently $H J H^t=J$, so conjugation by $H$ preserves
$r(c).\mathrm{mat}=1$.  The cross-section has the identities
\[
 \begin{split}
 H(x,y)^{-1}&=H(-x,-y),\\
 \kappa H(x,y)&=H(-x,-y),\\
 T(t)H(x,y)T(t)^{-1}&=H(tx,t^{-1}y).
 \end{split}                                                    \tag{KP.11}
\]
The first identity, in particular, stays within the same
cross-section; an arbitrary ordered product of root groups
would not have this property.

At $M_0=\operatorname{diag}(q,1,q^{-1})$, the Jacobian in
$(x,y)$ of either pair
$Z_i(HMH^{-1})$ or $F_i(HMH^{-1})$ is
\[
               \begin{pmatrix}q^{-1}-q&0\\0&q-q^{-1}\end{pmatrix}.
                                                                  \tag{KP.12}
\]
For $Z$ this follows by differentiating $[xX+yY,M_0]$.
For $F$ the same derivative follows because these gauge
directions are $\kappa$-odd and $M_0$ is fixed.  The residual
value $q=\bar\Psi(g_1)$ satisfies $q^2-1\in k(O)^\times$
by \texttt{D.Marked(g1)}.  Thus the formal implicit function
theorem gives a unique $(x,y)\in\mathfrak m_R^2$ making either
pair of marking functions zero.  Equivalently one may solve
successively modulo every power of $\mathfrak m_R$ using the
unit Jacobian, and take the unique limit.  This also proves
functoriality in every complete coefficient algebra.

Let $N_F$ send the old $Z$-slice to the new $F$-slice using
this gauge, and let $N_Z$ send the new slice to the old.
They are mutual inverses: if $r'=H\cdot r$, then
$H^{-1}=H(-x,-y)$ is a gauge in the same cross-section taking
$r'$ back to its required old marking; uniqueness forces it
to be the inverse normalization.  The same argument applies
in the other direction.  Formula (KP.11) and uniqueness show
that $N_F$ is torus-equivariant, and that the $F$-normalization
on the $c$-normalized functor is $\kappa$-equivariant.
Transport the genuine involution of the $F$-slice through
$N_Z\circ\kappa\circ N_F$.  This proves the functorial
involution on the old slice, and representability gives
$\kappa_S$.  The base lift needs no gauge, so its augmentation
is fixed.

The original global adjoint lift in (HS.14) is literally fixed by
(KP.6): its matrices are $A(g)D^{j(g)}$, where $A(g)$
preserves $K$, and both $A(g)$ and $D$ are fixed by $\kappa$.
Its matrix at $c$ is already one.  The local conjugation in
(HS.15) was used to verify its ordinary local property;
we do not replace the global representative by that locally
conjugated one.  The original lift is therefore on the
$F$-slice.  Applying $N_Z$ gives a $\kappa_S$-fixed point on
the actual old slice, with all the host conditions unchanged.

Finally the source's rational strict Selmer cocycles have
the literal shape $\gamma(g)=t(g)(E_{21}-E_{32})$ and satisfy
$t(c)=t(g_1)=0$; see
\path{ShaInjGen.lean:101--129,153--187,216--288}.
The differential of $\kappa$ is $X\mapsto-KX^tK$, which sends
$E_{21}-E_{32}$ to its negative.  The corresponding
dual-number lift has $M=M_0$ exactly, so both marking gauges
are one.  The transported involution therefore sends its
derivation to its negative, as claimed.
\end{proof}

\subsection{The exact unsquared conormal inequality}

Let $x_0:S\to O$ denote the base augmentation, put
$J_0=\ker x_0$, and write
\[
       C^-_{-1}=\bigl(J_0/J_0^2\bigr)_{\text{weight }-1}^{\kappa_S=-1}.
                                                                  \tag{KP.13}
\]
Unlike (KP.1), this is the ordinary integral cotangent of the
actual ring.  No saturation or reduced quotient is taken.
It is a finite $O$-module, although finite \emph{length} is an
additional assertion when needed below.  The finite module
decomposes by torus weights and by the two idempotents
$(1\pm\kappa_S)/2$.

\begin{theorem}[KT-ODD-CONORMAL-GATE]
For every $0\ne b\in O$, the rational strict Selmer torsion
chunk $\mathcal X[b]$ used in \path{sha_inj} injects into
\[
               \operatorname{Hom}_O(C^-_{-1},O/(b)).         \tag{KP.14}
\]
Here $\mathcal X$ is exactly the rational strict Selmer module
for the character $(\Psi)^{-1}$ in the source's convention.
In particular, if $C^-_{-1}$ has finite length, then
\[
           \operatorname{length}_O\mathcal X
                   \leq\operatorname{length}_O C^-_{-1}.    \tag{KP.15}
\]
Choose a finite torus- and involution-equivariant presentation
of $S$ over $O$.  Let $u_1,\ldots,u_n$ be its odd variables of
weight $-1$, and let $f_1,\ldots,f_m$ be its odd defining
equations of weight $-1$.  Write
\[
 f_j=\sum_{i=1}^n b_{ij}u_i+\text{terms of degree at least two},
       \qquad B_-=(b_{ij})\in M_{n\times m}(O).              \tag{KP.16}
\]
Then $C^-_{-1}=\operatorname{coker} B_-$.  If $\Delta$ is any
nonzero $n$-rowed minor, then
\[
       \operatorname{length}_O\mathcal X
                  \leq\operatorname{length}_O C^-_{-1}
                  \leq v_O(\Delta).                         \tag{KP.17}
\]
A unit maximal minor therefore proves vanishing.  If the
matrix has full rank, the best bound in (KP.17) is exactly
the valuation of its maximal-minor ideal.  For $n=0$ the
empty minor is one and the module is zero.
\end{theorem}
\begin{proof}
The explicit proof of \texttt{ShaInjGen} gives an injective
map from $\mathcal X[b]$ to the normalized cocycles just used.
Although the abstract \path{sha_inj} statement permits an
annihilator $a$, its proof takes $a=b$; see
\path{ShaInj.lean:29--38}.  The construction
\path{ModelLower.main_der} gives, before passing to the
ambient model, a derivation
$D:S\to O/(b)$ at $x_0$ with
\[
 D\bigl(r(g)_{ij}\bigr)=\Psi(g)^{\mathrm{wt}_1(j)}\gamma(g)_{ij}.
                                                                  \tag{KP.18}
\]
The unit factor in (KP.18) and the injectivity of the cocycle
map make the map to derivations injective.  Equivalently, two
derivations agreeing on all universal matrix entries agree
by uniqueness of the representing dual-number lift.  Each
derivation factors through $J_0/J_0^2$.  The cocycle has only
weight $-1$, so torus-equivariance, or (KP.18) applied to
torus conjugates, makes it zero on every other weight.
KT-EQUIVARIANT-MARKED-SLICE gives $D\circ\kappa_S=-D$.
Since two is a unit, it is zero on the even summand.  This
proves (KP.14).  The construction is additive; no unproved
compatibility with the saturated ambient quotient is used.

For a finite-length $O$-module $C$, the elementary divisor
decomposition shows
\[
\operatorname{length}_O\operatorname{Hom}_O(C,O/(b))
\leq\operatorname{length}_O C.
\]
The finite residue field
converts the injection into the same length inequality for
each torsion chunk, even if only its additive injectivity is
used.  The source's $\mathcal X$ is torsion, as in
\path{VlenLe.lean:13--42}; taking the increasing union of
its bounded torsion chunks proves (KP.15).

For completeness, an equivariant finite presentation is
obtained by choosing weight and involution eigenvectors in
$J_0/(J_0^2+\mathfrak m_OJ_0)$, lifting them to $J_0$, and
using the corresponding power series variables.  The torus
is diagonalizable and two is invertible, so the lifts can be
projected to the specified weight and sign.  The resulting
map from a finite power series ring is onto by complete
Nakayama.  Its ideal is closed and finitely generated;
apply the same weight/sign projections and Nakayama to
choose finitely many homogeneous eigen-equations.  The
formal torus action here is the one induced by the entry
weights; the Cayley change of slice is equivariant for it.

Modulo the square of the augmentation ideal, an equation
can affect the odd weight $-1$ variables only if it has that
same weight and sign.  Thus (KP.16) is exactly a presentation
of (KP.13).  Selecting a nonzero maximal minor gives a
square full-rank relation submatrix; its cokernel surjects
onto $\operatorname{coker}B_-$.  Smith normal form proves
the bound by $v_O(\Delta)$ and the equality for the
maximal-minor ideal.
\end{proof}

This is an exact integral conormal criterion, but no
matrix $B_-$ or unit minor has been evaluated in this run.
KT-ODD-CONORMAL-SELMER-EQUIVALENCE further proves that
it is equivalent to the original strict Selmer vanishing
problem.  In particular (KP.17) deletes no automorphic
branch from the existing trace-product bound and is not
counted as a stricter arithmetic gate.

\subsection{Unsquared normal equations along the adjoint branch}

The following exact algebra describes what a proposed
replacement by an arithmetic Jacobian would have to measure.
For any augmented ring with an involution as above, let
\[
 Q=(s-\kappa_S(s):s\in S),\quad B=S/Q,\quad N=Q/Q^2.         \tag{KP.19}
\]
Thus $B$ is the fixed-locus quotient, and $N$ its integral
normal conormal.  The involution is the identity on $B$
and is minus the identity on $N$.

\begin{lemma}[KT-UNSQUARED-NORMAL-JACOBIAN]
There is a canonical isomorphism
\[
       (J_0/J_0^2)^-\simeq N\otimes_{B,x_0}O.               \tag{KP.20}
\]
For an equivariant presentation
$S=O[[v_1,\ldots,v_a,u_1,\ldots,u_b]]/(f^+_j,f^-_j)$,
where $v$ is even and $u$ is odd, the $B$-module $N$ has the
presentation
\[
 B^{\{f^-_j\}}\xrightarrow{\,\mathcal J^-\,}B^b
              \longrightarrow N\longrightarrow0,
 \qquad
 \mathcal J^-_{ij}=\left.\frac{\partial f^-_j}{\partial u_i}\right|_{u=0}.
                                                                  \tag{KP.21}
\]
Consequently, at any integral fixed point
$\xi:B\to O_\xi$ with $O_\xi$ a finite coefficient DVR,
if $\mathcal J^-(\xi)$ has full row rank over its fraction
field, then
\[
 \operatorname{length}_{O_\xi}(N\otimes_{B,\xi}O_\xi)
       =v_{O_\xi}\!\left(\operatorname{Fitt}_0
                  (N\otimes_{B,\xi}O_\xi)\right)
       \leq v_{O_\xi}(\Delta_\xi)                            \tag{KP.22}
\]
for every nonzero maximal minor $\Delta_\xi$ of the evaluated
unsquared matrix.  The ideal in (KP.22) is independent of
the chosen presentation.
\end{lemma}
\begin{proof}
Write $J_0=J_0^+\oplus J_0^-$.  Commutativity and the two
projectors give
\[
              (J_0^2)^-=J_0^+J_0^- .                        \tag{KP.23}
\]
The ideal $Q$ is generated by $J_0^-$, its odd part is
$J_0^-$, and its even part consists of sums of products of
two odd elements, hence lies in $Q^2$.  Therefore $N$ is
entirely odd.  Also
$Q\cap J_0^2=J_0Q$: its odd part is (KP.23), and its even
part is contained in $Q^2\subset J_0Q$.  It follows that
\[
 Q/Q^2\otimes_B O=Q/J_0Q
       \simeq\operatorname{im}(Q\to J_0/J_0^2)
       =(J_0/J_0^2)^-,
\]
proving (KP.20).  This is an integral identity, not a
comparison with the reduced generic fixed locus.

In the displayed presentation, $Q$ is generated by the
classes of the $u_i$.  An even equation has no term linear
in the odd variables.  An odd equation, modulo $(u)^2$,
is exactly the column in (KP.21).  The even equations
specialized at $u=0$ define $B$.  These observations give
the right-exact presentation (KP.21), including all its
relations; multiples of the even equations vanish in $B$
and add none.  Base change to $O_\xi$ is right exact, and
Smith normal form proves (KP.22).  Presentation independence
is the defining invariance of the zeroth Fitting ideal.
\end{proof}

At an actual adjoint point, the coefficient module underlying
these odd normal cocycles has a particularly precise form.
Use left multiplicative cocycles
$r_\epsilon(g)=(1+\epsilon\gamma(g))r(g)$, exactly as in
\texttt{IsDeform}; thus their coefficient representation is
$\operatorname{ad}_{GG}r$, not its Tate dual.

\begin{lemma}[KT-ADJOINT-ODD-NORMAL-MODULE]
Let $\rho$ be the two-dimensional representation used in the
adjoint lift of KT-ADJOINT-ACTUAL-HOST-POINT, over a
characteristic-zero coefficient field.  For that literal
$\kappa$-fixed lift, the even and odd summands of
$\mathfrak{gl}_3$ as $G_{\mathbf Q}$-modules are
\[
 \begin{split}
   (\mathfrak{gl}_3)^+&\simeq\operatorname{Ad}^0\rho,\\
   (\mathfrak{gl}_3)^-&\simeq
       \eta\otimes\left(\mathbf1\oplus
                    \operatorname{Sym}^4\rho\otimes(\det\rho)^{-2}\right).
 \end{split}                                                     \tag{KP.24}
\]
Replacing $\rho$ by its dual leaves the normalized symmetric
fourth in this formula unchanged.
\end{lemma}
\begin{proof}
The differential of $\kappa$ is $X\mapsto-KX^tK$.
Its even space is $\mathfrak{so}(K)$ of dimension three,
and its odd space is the $K$-self-adjoint space of dimension
six.  Put $V=\operatorname{Ad}^0\rho$.  The invariant trace
form identifies $V$ with $V^*$, and the elementary
$\mathrm{GL}_2$ decomposition is
\[
 \bigwedge^2 V\simeq V,\qquad
 \operatorname{Sym}^2 V\simeq\mathbf1\oplus
                      \operatorname{Sym}^4\rho(\det\rho)^{-2}.
                                                                  \tag{KP.25}
\]
One way to verify the second identity is to write
$V=\operatorname{Sym}^2\rho\otimes(\det\rho)^{-1}$.
Multiplication of binary quadratics gives the surjection
$\operatorname{Sym}^2(\operatorname{Sym}^2\rho)
\to\operatorname{Sym}^4\rho$; its one-dimensional kernel is
$(\det\rho)^2$.  Over characteristic zero this sequence
splits.  The alternating identity follows from the
invariant volume and trace form of the determinant-one
three-dimensional representation.

The actual GG matrix is $A(g)D^{j(g)}$, with
$\eta(g)=(-1)^{j(g)}$ and $A(g)=\operatorname{Ad}^0\rho(g)$.
For $j(g)=1$ its action on the Lie algebra is
\[
   X\longmapsto A(g)D\theta_{\rm Lie}(X)D A(g)^{-1}
              =A(g)\kappa_{\rm Lie}(X)A(g)^{-1}.             \tag{KP.26}
\]
For $j(g)=0$ it is ordinary conjugation by $A(g)$.  Thus
the sign $\eta$ occurs on precisely the odd summand of
(KP.25), proving (KP.24).  Finally
$\rho^*\simeq\rho\otimes(\det\rho)^{-1}$ implies
\[
 \operatorname{Sym}^4(\rho^*)\otimes\det(\rho^*)^{-2}
       \simeq\operatorname{Sym}^4\rho\otimes(\det\rho)^{-2}.
\]
No Tate twist is introduced here.  A later use of local or
global duality would have to dualize and Tate-twist the
coefficient module separately.
\end{proof}

At a fixed two-sided point the carried characteristic-zero
rigidity statement \texttt{M4} leaves only the torus tangent.
That tangent is nonzero (some nonzero-weight function is
nonzero at the point) and is even because the point is fixed
and $\kappa_S$ commutes with the torus.  Hence every odd
normal tangent is zero.  Here the passage from the actual
ring to the finite graded model can be made explicitly.
A functional on $N\otimes_{B,\xi}\operatorname{Frac}O_\xi$
gives the odd derivation
$D(s)=\phi((s-\kappa_S(s))/2\bmod Q^2)$.
Since $N$ is a finite $B$-module and $\xi(B)\subset O_\xi$,
the values of $D$ are bounded in the fraction field.
After multiplication by some nonzero $c\in O$, the map
$f+\epsilon cD$ takes values in the complete coefficient
order
\[
 (O+\mathfrak m_{O_\xi})\oplus
                  \epsilon\mathfrak m_{O_\xi},\qquad\epsilon^2=0.
                                                                  \tag{KP.27}
\]
This order has the same residue field as $O$ and is
\texttt{IsCO}, by the same finite-order argument as in
KT-ADJOINT-ACTUAL-HOST-POINT.  If the restriction of $D$
to the model is zero, its value on every universal matrix
entry is zero.  Uniqueness in \texttt{S.rep}, applied to
(KP.27), makes the two homomorphisms $f+\epsilon cD$ and
$f$ equal; hence $D=0$.  The torus derivative is bounded
as well, by its functorial action over (KP.27).  Restrictions
therefore preserve linear independence between an odd
normal tangent and the even torus tangent, so \texttt{M4}
indeed excludes the former.  This avoids any assertion
about unrestricted algebraic differentials of a completed
power series ring.

It follows that (KP.21) has full generic row rank at such
a point, so (KP.22) is a finite local normal-length
invariant.  It does not make that invariant a unit:
characteristic-zero rigidity says only that the
corresponding integral normal module is torsion.

\paragraph{Precisely what remains for a branchwise replacement.}
The base inequality (KP.17) and the branch inequalities
(KP.22) are proved separately.  To replace an adjoint
trace-difference factor in the current global bound by
\emph{branchwise} normal Jacobian values, one still needs
an integral comparison bounding
$\operatorname{length}_O C^-_{-1}$ by those normal modules,
with appropriate residue degrees, ramification normalizations,
and a term controlling any vertical or embedded contribution.
The existing squared Rees construction (KP.2) does not give
that comparison.  A passage to the finite flat reduced
trace algebra does not give it either, by (KP.3)--(KP.4).
The exact missing assertion can also be stated as a kernel
condition: if a quotient $q:S\to S^{\rm flat,red}$ has trivial
involution, the image of the odd part of $\ker q$ in
$(J_0/J_0^2)_{-1}$ must be proved zero before that quotient
can discard the relevant Selmer conormal.  Point detection
alone proves no such containment in $J_0^2$.

The following conditional transfer makes the desired
numerical comparison concrete.  Its extra hypotheses are
not being asserted for the archived deformation model.

\begin{lemma}[KT-NORMAL-FLAT-TRANSFER, conditional criterion]
Let $B_0$ be a finite flat $O$-algebra whose generic fiber
is a product of fields $K_1,\ldots,K_s$.  Let $L$ be a
finite $B_0$-module equipped with a surjective $O$-linear map
$L\twoheadrightarrow C^-_{-1}$.  Suppose
\[
            B_0^m\xrightarrow{\mathcal J} B_0^n
                       \longrightarrow L\longrightarrow0
                                                                  \tag{KP.28}
\]
is a presentation, and $H$ is any $n$ by $n$ matrix obtained
by $B_0$-linear combinations of its relation columns such
that $\det H$ is nonzero in every $K_i$.  Let $O_i$ be the
integral closure of $O$ in $K_i$, let
$f_i=[k(O_i):k(O)]$, and let $H_i$ be the image of $H$ in
$M_n(O_i)$.  Then
\[
 \operatorname{length}_O\mathcal X
 \leq \operatorname{length}_O C^-_{-1}
 \leq \operatorname{length}_O L
 \leq v_O\!\left(\operatorname{Norm}_{B_0/O}(\det H)\right)
 =\sum_{i=1}^s f_i\,v_{O_i}(\det H_i).                      \tag{KP.29}
\]
In particular this replaces a trace comparison by normal
relation determinants whenever the stated integral
comparison is supplied.  Full row rank of $\mathcal J$
over every $K_i$ suffices for the existence of such an $H$.
\end{lemma}
\begin{proof}
The image of $H$ is contained in the image of
$\mathcal J$, so $\operatorname{coker}H$ surjects onto $L$.
Since $O$ is local, $B_0$ is finite free as an $O$-module.
After restricting scalars, $H$ is a square matrix over $O$
which is invertible over its fraction field.  Smith normal
form gives
\[
 \operatorname{length}_O\operatorname{coker}H
     =v_O(\det_O H)
     =v_O\!\left(\operatorname{Norm}_{B_0/O}(\det_{B_0}H)\right).
\]
The last equality can be checked after tensoring with the
fraction field and decomposing into the $K_i$; determinants
on a field extension restrict to field norms.  The usual
valuation of a field norm is $f_i$ times the normalized
valuation in $O_i$.  This proves the last equality of
(KP.29), and the quotient maps prove the two intermediate
inequalities.  Finite length now permits applying
KT-ODD-CONORMAL-GATE for the first one.

For the final assertion take $H=\mathcal J P$ with an
$m$ by $n$ matrix $P$ of indeterminates.  Full row rank in
$K_i$ makes $\det(\mathcal J_iP)$ a nonzero polynomial
for every $i$.  A finite union of their zero loci cannot
contain all matrices over the infinite field
$\operatorname{Frac}O$; equivalently expand coefficients
in bases of the finite extensions.  Choosing such a
matrix and clearing denominators supplies $P$ over $O$,
and hence the required $H$.
\end{proof}

For the actual application the full fixed-locus ring $B$
in (KP.19) retains torus directions and is not the finite
algebra required in this lemma.  One must descend the
relevant normal module to an appropriate finite
weight-zero algebra and control its integral torsion
before taking a flat reduced quotient.  A presentation
over that algebra, the surjection onto the base conormal,
and full generic rank on every retained component are
substantive hypotheses.  They are not consequences of
having a finite set of reduced automorphic isotypes.

Formula (KP.24) identifies the candidate arithmetic normal
deformation problem: the $\eta$-twisted normalized symmetric
fourth, together with the $\eta$ scalar line and the actual
ordinary/away-prime conditions.  An equality of its integral
Fitting ideal with an automorphic congruence ideal or a
special $L$-value would be an additional theorem.  None is
asserted here.  The present structural results are the
actual integral projector, its unsquared conormal presentation,
and the precise duality with the original strict Selmer
module.  No all-weight vanishing or removal of the adjoint
source family has been inferred from them.

% END evidence/unitary_audit/projector_results.tex


% BEGIN evidence/unitary_audit/projector_selmer_equivalence.tex
\section{The odd conormal is the original strict Selmer dual}
\label{sec:kt-projector-equivalence}

This is a scope audit of the new projector criterion, not an
additional arithmetic vanishing argument.  We retain the
hypotheses and the actual normalized ring $S$ of
KT-EQUIVARIANT-MARKED-SLICE.  The source's rational trivial
character is the character of modulus $N=2$:
\path{SmallClassNumbers/ClassNumberThree.lean:13--32}
identifies its module with
\[
 \mathcal X=\operatorname{Sha}_{2\ell}
              (G_{\mathbf Q},(K_O/O)(\alpha)),\qquad
 K_O=\operatorname{Frac}O,\quad
 \alpha=\Psi^{-1}=\varepsilon_\ell^{-e}.                    \tag{KE.1}
\]
Here $\operatorname{Sha}_{2\ell}$ means exactly the source's
\texttt{ShaU}: the cohomology class restricts to zero on
every inertia group and on every decomposition group over
primes dividing $2\ell$.  The literal definition is
\path{GaloisLibrary/SelmerTransport/ConjStableShaFam.lean:94--96,167--169}.
The identification of the coefficient module
$\texttt{QO}=O\otimes_{\mathbf Z_\ell}(\mathbf Q_\ell/\mathbf Z_\ell)$
with $K_O/O$ follows by tensoring with the finite free
$\mathbf Z_\ell$-module $O$.

\begin{theorem}[KT-ODD-CONORMAL-SELMER-EQUIVALENCE]
Put
\[
 C=\bigl(\ker(x_0:S\to O)/(\ker x_0)^2\bigr)_{-1}^{\kappa_S=-1}.
                                                                  \tag{KE.2}
\]
For every $0\ne b\in O$ there is an isomorphism
\[
        \operatorname{Hom}_O(C,O/(b))\simeq\mathcal X[b].    \tag{KE.3}
\]
Use $O/(b)\simeq(K_O/O)[b]$, $u\mapsto u/b$, to specify
the coefficient identification.  The isomorphisms are
compatible with the usual inclusions of torsion chunks;
in particular
\[
             \mathcal X\simeq\operatorname{Hom}_O(C,K_O/O).
                                                                  \tag{KE.4}
\]
Thus $C=0$ if and only if $\mathcal X=0$.  Under the
established finiteness of $\mathcal X$ in the carried
source chain, $C$ has finite length and
\[
                  \operatorname{length}_O C
                       =\operatorname{length}_O\mathcal X. \tag{KE.5}
\]
Consequently the unit-minor criterion of
KT-ODD-CONORMAL-GATE is equivalent to the original strict
Selmer vanishing problem.  Without an independently
evaluated minor or a new arithmetic comparison, it does
not constitute a stricter arithmetic gate or the removal
of an automorphic branch.
\end{theorem}

\begin{proof}
Write $M=(K_O/O)(\alpha)$ and $A_b=(O/(b))(\alpha)$.
We give both directions, including the coefficient and
local-condition checks.

\paragraph{Global coefficient identification and normalization.}
The marked element $g_1$ satisfies
$\Psi(g_1)^2-1\in O^\times$.  Therefore
$\alpha(g_1)-1\in O^\times$, so $M^{G_{\mathbf Q}}=0$.
The exact sequence
\[
       0\longrightarrow A_b\longrightarrow M
         \xrightarrow{\ b\ }M\longrightarrow0              \tag{KE.6}
\]
gives
$H^1(G_{\mathbf Q},A_b)\simeq H^1(G_{\mathbf Q},M)[b]$.
Every class on the left has a unique cocycle representative
$t$ with $t(g_1)=0$: subtract its uniquely determined
coboundary using the unit $\alpha(g_1)-1$.  This is the
normalization performed in
\path{ShaInjGen.shaInj_norm:33--99}.  Also $t(c)=0$,
because $c^2=1$, $\alpha(c)=1$, and two is invertible.

\paragraph{From an odd conormal functional to a scalar cocycle.}
A functional in $\operatorname{Hom}_O(C,O/(b))$ extends by
zero on every other weight and sign to an $O$-derivation
$D:S\to O/(b)$ at $x_0$.  Conversely such a derivation
factors uniquely through $C$.  The homomorphism
\[
           F(s)=x_0(s)+\epsilon D(s),\qquad
           \epsilon^2=0,\quad b\epsilon=0                  \tag{KE.7}
\]
takes values in the complete local ring
$O\oplus\epsilon O/(b)$, which is \texttt{IsCO}.
Continuity introduces no extra condition: if
$\mathfrak m_S=\mathfrak m_O S+\ker x_0$, Leibniz gives
$D(\mathfrak m_S^r)\subset\mathfrak m_O^{r-1}O/(b)$,
which is zero for all sufficiently large $r$.
The universal representation evaluated at (KE.7) is a
continuous dual-number lift.

Pass through the isomorphism to the genuinely equivariant
$F$-marked slice of KT-EQUIVARIANT-MARKED-SLICE.  On this
slice $\kappa$ is the literal matrix operation (KP.6).
The left multiplicative cocycle of the lift therefore has
weight $-1$ and is odd for
$X\mapsto-KX^tK$.  Hence it has exactly the form
\[
         \gamma(g)=t(g)Y_-,\qquad Y_-=E_{21}-E_{32},        \tag{KE.8}
\]
and the scalar cocycle relation is
$t(gh)=t(g)+\alpha(g)t(h)$, by (HS.2).
At $g_1$, put $q=\Psi(g_1)$.  The linearized second odd
marking (KP.9), applied to (KE.8), is
\[
                         (q+1)t(g_1)=0.                    \tag{KE.9}
\]
Since $q+1$ is a unit, $t(g_1)=0$.  The old marking has
the same linearized value on this odd direction, so both
Cayley changes of slice are the identity for this
dual-number lift.  Thus (KE.8) and its normalization
describe the functional on the actual original ring as
well; we have not confused a transported tangent sign
with a literal matrix-entry equation.

\paragraph{The local condition at $\ell$ is strict.}
KT-ORDINARY-INTEGRAL-COMPONENT identifies the actual local
quotient $U.R/I_t$ with the integral ordered ordinary
functor.  Therefore the local lift just constructed has
its unique ordinary flag.  After setting $\epsilon=0$
this is the standard flag of $r_0$: its ordered inertial
eigenvalues are distinct.  Its lower-unipotent flag
coordinate consequently has the form $1+\epsilon W$,
with $W$ strictly lower triangular over $O/(b)$.
Writing the lift as the conjugate of an upper Borel
lift and comparing lower entries gives
\[
       \gamma|_{G_{\mathbf Q_\ell}}
            =\operatorname{cob}(W)
                       \quad\text{in the lower-weight space}.
                                                                  \tag{KE.10}
\]
The base adjoint action preserves weights and commutes
with $\kappa_{\rm Lie}$.  Project (KE.10) to weight $-1$
and sign $-1$.  The projected $W$ is a scalar multiple
of $Y_-$, so $t|_{G_{\mathbf Q_\ell}}$ is a coboundary
already with coefficients in $A_b$.  This argument works
for both split and unramified nonsplit $\eta$.

\paragraph{All away-prime inertia conditions.}
For the trivial Dirichlet character, the exact
\path{DatFacts.dat_facts} description of
\texttt{Kill} is the union, over all places away from
$\ell$, of
\[
                        I_q\cap\ker\eta .                  \tag{KE.11}
\]
See \path{RhoKApply.lean:61--62} and
\path{DatFacts.lean:70--100}.
The cyclotomic character is unramified at each such
place, so $\alpha|_{I_q}=1$ and $t|_{I_q}$ is a
continuous additive homomorphism.  The box condition
makes it zero on (KE.11).  It therefore factors through
a group of order at most two.  Two is a unit on $A_b$,
so this homomorphism is zero on all of $I_q$.  This
includes the ramified quadratic place $q=3$; no
unramifiedness of $\eta$ at that place is being assumed.
Together with the local strict condition at $\ell$,
all inertia groups in the literal \texttt{shaFam} have
been handled.

\paragraph{The extra decomposition group at $2$.}
The parameter is $2\ell$, so the source also requires
triviality on a full decomposition group at $2$.
After embedding $A_b$ in $M$, the class just obtained
is unramified there.  Its unramified cohomology is
\[
 H^1_{\rm unr}(G_{\mathbf Q_2},M)
       =M/(2^{-e}-1)M=0.                                   \tag{KE.12}
\]
Indeed $e>0$, so $2^{-e}-1\ne0$ in $K_O$, and
multiplication by any nonzero element of $O$ is
surjective on $K_O/O$.  Formula (KE.12) uses divisible
coefficients.  The corresponding group with coefficients
in $O/(b)$ need not vanish, and no such assertion is
used.  The usual conjugacy of decomposition groups
handles every extension of the place to the algebraic
closure.  The definitions of \texttt{shaFam} and
\texttt{ShaU} now show that the image lies in
$\mathcal X[b]$.

\paragraph{The inverse construction.}
Given $\xi\in\mathcal X[b]$, use (KE.6) to obtain its
unique finite-coefficient class and normalize its
cocycle $t$ at $g_1$.  On every inertia group away from
$\ell$ the action on $A_b$ and $M$ is trivial.  Since
$\xi$ restricts to zero in $H^1(I_q,M)$ and the
coefficient inclusion is injective, $t|_{I_q}=0$.

At $\ell$, choose an inertia element $\tau$ for which
$\bar\Psi(\tau)\ne1$, as supplied by GD.  Then
$\alpha(\tau)-1$ is a unit and
$M^{G_{\mathbf Q_\ell}}=0$.  Applying (KE.6) locally
shows that
\[
 H^1(G_{\mathbf Q_\ell},A_b)
                   \longrightarrow H^1(G_{\mathbf Q_\ell},M)
                    \quad\text{is injective}.              \tag{KE.13}
\]
The strict local condition on $\xi$ therefore makes
$t$ a local coboundary with finite coefficients.  Form
the cocycle (KE.8) and the corresponding dual-number
GG lift.  It kills \texttt{Kill}.  Its local restriction
is conjugate by $1+\epsilon wY_-$ to the exact base
representation, hence is integrally ordinary with the
specified inertia; the local component lemma gives
the required $I_t$ condition.  The equations
$t(c)=t(g_1)=0$ give the source normalization.  The
universal property of $S$ now supplies (KE.7).
It is odd and of weight $-1$, so its derivation factors
through $C$.

Uniqueness of the $g_1$-normalized scalar cocycle and
uniqueness in \texttt{S.rep} show that these two
constructions are inverses.  They respect addition
and scalar multiplication.  The choice of cyclic
generator in the archived \path{sha_inj} may change
its displayed cocycle map by a coefficient unit;
the explicit identification $u\mapsto u/b$ fixes that
choice in (KE.3).

Finally use $b=\varpi^n$ and take the direct limit.
Every homomorphism from the finite $O$-module $C$ to
$K_O/O$ has bounded torsion image, giving (KE.4).
This dual functor detects every nonzero finite
$O$-module.  Its elementary-divisor decomposition
shows that it has a divisible summand for every free
summand of $C$, while a summand $O/(\varpi^a)$ is
dual to another summand of the same length.  Finiteness
of $\mathcal X$ therefore forces $C$ to be torsion
and proves (KE.5).
\end{proof}

% END evidence/unitary_audit/projector_selmer_equivalence.tex


% BEGIN evidence/unitary_audit/projector_flatness_audit.tex
\section{The finite normal module exists; flatness is a separate gate}
\label{sec:kt-projector-flatness}

The conditional tool KT-NORMAL-FLAT-TRANSFER asks for more
than generic automorphic point detection.  Two of its
integral hypotheses can nevertheless be checked for the
actual source presentation: a finite weight-zero algebra
and a finite normal module surjecting onto the base odd
conormal.  The following statement records exactly this
partial verification, without a flatness assertion.

\begin{lemma}[KT-FINITE-GRADED-NORMAL-MODULE]
Use the weighted presentation of the actual normalized
source ring $S$ in \texttt{ModelOfPres}, together with the
same carried finiteness of \texttt{weightZeroSubalgebra}
used in \texttt{ModelUpper}.  Let $S_j$ be its
weight spaces and let
\[
 A=\bigoplus_{j\in\mathbf Z}S_j\subset S,
 \qquad Q_A=(a-\kappa_S(a):a\in A),
 \qquad B=A/Q_A.
                                                                  \tag{KF.1}
\]
Then $B_0$ is a finite $O$-algebra, and
\[
                 L=(Q_A/Q_A^2)_{-1}                        \tag{KF.2}
\]
is a finite $B_0$-module.  There is a natural surjective
$O$-linear map
\[
                            L\twoheadrightarrow C,         \tag{KF.3}
\]
where $C$ is the odd weight $-1$ cotangent in (KE.2).
The involution acts by minus one on $Q_A/Q_A^2$.
\end{lemma}
\begin{proof}
The source construction supplies a weighted power series
presentation of $S$ with finitely many variables of
weights in $\{-2,-1,0,1,2\}$ and homogeneous kernel.
Precisely, \path{ModelOfPres.mop_cover} constructs the
surjection $q_h:R\to S$ and its finitely generated
homogeneous kernel; \path{mop_main} identifies the
variables with differences of universal matrix entries
and checks the weight bounds.  These are
\path{ModelOfPres.lean:34--121,168--238,243--284}.
Composing its weighted presentation of $R$ with $q_h$
gives the asserted presentation of $S$.

The degree-zero image of this presentation is the image
of the ambient model's degree-zero piece.  The complete
proof of
\path{ModelIdealFinitenessOfWeightZeroFinite.lean:101--189}
places it inside the already finite
\texttt{weightZeroSubalgebra} by compactness and proves
it finite over $O$.  Hence $S_0$ is finite over $O$.
The standard weighted-presentation argument
\path{CommutativeAlgebra/IsWRing.lean:8--56}, theorem
\texttt{isWRing}, makes $A$ a finite-type algebra over
$S_0$.  Its invoked generation theorem is exactly
\path{GrGen.lean:72--109}, \path{gr_gen}, whose
homogeneous-series decomposition is proved at lines
15--57: finitely many
nonzero-weight monomials generate all homogeneous
series, with their remaining zero-weight factors in
$S_0$.  Thus $A$ is Noetherian.  The involution commutes
with the torus, so it preserves every $S_j$ and $A$.
Its ideal $Q_A$ is homogeneous and finitely generated
by homogeneous odd elements.  Therefore $B_0$ is the
quotient of $S_0$ by $Q_A\cap S_0$ and is finite over $O$.

The conormal $N_A=Q_A/Q_A^2$ is a finite graded
$B$-module.  Each of its weight pieces is finite over
$B_0$.  Here is a proof that also applies with negative
weights.  For a fixed $j$, form the $B$-submodule of
$N_A$ spanned by $(N_A)_j$.  Since $N_A$ is Noetherian,
finitely many elements $z_1,\ldots,z_r$ of that weight
generate this submodule.  If
$z=\sum_i b_i z_i$ has weight $j$, its weight-$j$
projection is $z=\sum_i(b_i)_0z_i$.  Thus the same
finite list generates $(N_A)_j$ over $B_0$.  This is
the module version of the exact argument in
\path{GradedModulesLength/FgHolds.lean:5--45}.
Taking $j=-1$ proves the finiteness of $L$.

The augmentation kills every odd element, so the image
of $Q_A$ lies in $J_0=\ker x_0$ and that of $Q_A^2$
lies in $J_0^2$.  Every element of the odd weight $-1$
piece of $J_0/J_0^2$ has an odd homogeneous lift in
$S_{-1}\subset A$: apply the weight projector and
$(1-\kappa_S)/2$ to any lift.  That lift belongs to
$Q_A$, proving (KF.3).  Finally the proof of
KT-UNSQUARED-NORMAL-JACOBIAN, applied to $A$, shows
that $N_A$ is entirely odd.  This argument needs no
identification of an algebraic conormal with its
completed counterpart beyond the displayed
surjection.
\end{proof}

What is still absent is a proof that $B_0$ is
torsion-free, together with the appropriate reduced
generic structure and full generic normal rank on all
of its components.  The carried rigidity statement
handles the two-sided characteristic-zero orbits; it
does not by itself control vertical contributions or
every fiber of the degree-zero quotient.  Replacing
$B_0$ by its maximal flat reduced quotient may kill
part of $L$ and its image in $C$.  The reverse-map
theorem identifies the latter image with the original
strict Selmer dual, so this loss cannot be ignored.

There are two concrete source boundaries.  First,
\path{ModelIdealFinitenessOfWeightZeroFinite} proves
module finiteness only; its displayed conclusions
contain no torsion-freeness or flatness statement.
Second, \path{ModelOfPres.mop_cover} forms the
global ring by imposing the global and marking
relations on a power series extension of the local
ring.  Even the new integral smoothness of the local
ring does not prove that these global relations have
torsion-free quotient.  The characteristic-zero
point detector \texttt{EtHolds} does not fill either
gap, for the reasons in Part I of the projector audit.

\begin{lemma}[KT-FIXED-QUOTIENT-FLATNESS-COUNTEREXAMPLE]
A finite flat reduced augmented $O$-algebra with an
involution can have a fixed-locus quotient with
nonzero $O$-torsion, even when two is a unit.
\end{lemma}
\begin{proof}
For $a=\varpi^n$, $n\geq1$, let $R$ be the subalgebra
of $O^3$ with $O$-basis
\[
  1=(1,1,1),\quad u=(0,a,-a),\quad v=(0,a,a).
                                                                  \tag{KF.4}
\]
These vectors are linearly independent, since $a\ne0$
and two is a unit.  They are closed under
multiplication because
\[
                u^2=av,\qquad uv=au,\qquad v^2=av.          \tag{KF.5}
\]
Thus $R$ is finite free of rank three and reduced,
with generic fiber $K_O^3$.  Projection to its first
coordinate is an augmentation.  The involution
swapping the last two coordinates sends $u$ to
$-u$ and fixes $v$.  Its fixed-locus ideal is $(u)$,
whose $O$-span is generated by $u$ and $av$.
Consequently
\[
                  R/(u)\simeq O[v]/(av,v^2).               \tag{KF.6}
\]
The nonzero class of $v$ is killed by $a$.  In
particular, flatness of an ambient algebra does not
pass to the fixed-locus quotient used in (KF.1).
This is an algebraic scope counterexample, not a
counterexample to the complete source hypotheses.
\end{proof}

No determinant over the actual $B_0$, saturation
identity for $L$, or arithmetic normal Jacobian has
been computed or proved to be a unit.  The verified
finite algebra and module in (KF.1)--(KF.3) identify
where a further integral theorem would apply; they
do not discharge its flatness and support hypotheses.

% END evidence/unitary_audit/projector_flatness_audit.tex


% BEGIN evidence/unitary_audit/hecke_normal_results.tex
\section{An evaluated Hecke separator and the missing integral extension data}
\label{sec:kt-hecke-normal}

The complementary stable factors computed earlier do give an
integral, all-prime statement about the localized Hecke algebra.
This is an arithmetic consequence of an evaluated unit.  The
following audit then explains precisely why it does not yet give
a smaller strict Selmer bound.

Write $A_q=q^{e-1}T_q$ and $A_q^{\rm opp}=q^{e-1}T_q^{\rm opp}$
for the two integral scaled split operators, and put
$t_q=1+q^e+q^{2e}$.  Let $M^{\rm st}_e$ be the stable compact
space, with its proved decomposition into the constructed
symmetric-square subspace and its complementary stable subspace.
Write $R_{e,7}(X)$ for the characteristic polynomial of $A_7$
on the latter.  Thus its degrees at $e=26,28,30$ are
$72,92,112$.

\begin{proposition}[KT-STABLE-HECKE-SELFDUAL-LOCALIZATION]
For $e\in\{26,28,30\}$ put, respectively,
\[
 M_e=697446897311,\quad323026426255399,\quad
      236017358500033404991.
                                                               \tag{KH.1}
\]
Let $\ell$ be any prime divisor of $M_e$.  On the entire stable
space one has, for every split $q\ne\ell$,
\[
 R_{e,7}(A_7)(A_q-A_q^{\rm opp})=0.                         \tag{KH.2}
\]
Choose a common integral $\mathbf Z_\ell$-lattice and let
$\mathbb T^{\rm st}$ be the algebra generated by all the displayed
split operators.  At the Eisenstein residual maximal ideal
\[
 \mathfrak m=(\ell,A_q-t_q,A_q^{\rm opp}-t_q:q\ne\ell
                   \text{ split}),                        \tag{KH.3}
\]
one has the integral identities
\[
         A_q=A_q^{\rm opp}\quad\hbox{in }
                       \mathbb T^{\rm st}_{\mathfrak m}.  \tag{KH.4}
\]
The involution interchanging the two operators is the identity
on this localized order.  Its characteristic-zero cuspidal
systems all belong to the constructed symmetric-square subspace.

The same assertion holds for the augmented order generated by
the pairs $(t_q,A_q)$ and $(t_q,A_q^{\rm opp})$ inside
$\mathbf Z_\ell\oplus\mathbb T^{\rm st}$, localized at its
augmentation modulo $\ell$.  Here the first factor is a formal
Eisenstein character; it is not asserted to be a compact
cuspidal representation.
\end{proposition}
\begin{proof}
All the split Hecke operators commute.  On the complementary
subspace Cayley--Hamilton makes $R_{e,7}(A_7)$ zero.  On the
constructed symmetric-square subspace the associated adjoint
representation is self-dual, and the two normalized minuscule
eigenvalues agree at every split prime.  Semisimplicity therefore
makes $A_q-A_q^{\rm opp}$ zero there.  These facts prove (KH.2)
over characteristic zero, and hence on the integral lattice.
No integral direct-sum decomposition of that lattice is needed.

The integral harmonic lattice at $\ell>3$ supplies a common
Hecke-stable lattice; intersect it with the rational stable
subspace.  The involution given by interchanging the harmonic
variables preserves it after replacing the lattice by its sum
with its image, if necessary.  KT-U3-DUAL-MATRIX identifies its
conjugation action on the two split operators.  The algebra
$\mathbb T^{\rm st}$ is a finite torsion-free
$\mathbf Z_\ell$-module, being a submodule of the finite matrix
algebra.  The earlier Fretwell--Roberts construction provides a
cuspidal eigencharacter with the values (KH.3), so this ideal is
proper and maximal.

The exact calculations give
$\gcd(M_e,R_{e,7}(t_7))=1$.  Consequently the image of
$R_{e,7}(A_7)$ modulo $\mathfrak m$ is nonzero and this element
is a unit in the localized algebra.  Multiplying (KH.2) by its
inverse proves (KH.4).  The generators are all fixed, so the
involution is the identity.  A characteristic-zero factor from
the complement cannot survive this localization, since on that
factor $R_{e,7}(A_7)=0$.  This also proves the generic-support
assertion without an assumption of distinct eigenvalues at seven.

On the added Eisenstein coordinate the second factor in (KH.2)
is zero.  The same relation and unit argument therefore apply
to the augmented order, which is again finite and torsion-free.
Its first-coordinate projection gives the specified augmentation.
\end{proof}

The new small certificate \path{hecke_normal_audit.py}
evaluates the previously reviewed exact coefficient arrays
modulo $M_e$ and saves explicit Bezout identities.  At seven
the reduced values and their inverses are
\[
\begin{array}{c|r|r}
e&R_{e,7}(t_7)\bmod M_e&R_{e,7}(t_7)^{-1}\bmod M_e\\\hline
26&429636065482&455711584154\\
28&242701611239023&44010470192816\\
30&53393982565377757302&202700793921978609871
\end{array}                                                   \tag{KH.5}
\]
The analogous three certificates at thirteen are also saved.
No primality assumption about any $M_e$ is used.  These six unit
certificates strengthen the finite-prime comparison to (KH.4),
but they are not a new Selmer vanishing certificate.

\subsection{Why a trace comparison does not detect the required tangent}

\begin{lemma}[KT-TRACE-INVISIBLE-SELMER-TANGENT]
Let $S,x_0,C$ be the actual normalized source ring, its diagonal
augmentation, and the odd weight $-1$ cotangent from
KT-ODD-CONORMAL-SELMER-EQUIVALENCE.  For every nonzero $b\in O$,
each element of $\operatorname{Hom}_O(C,O/b)$ induces zero
first-order variation of every characteristic coefficient of
the universal representation restricted to $G_E$.  In particular
it is zero on every trace coordinate and on every polynomial
in such coordinates.
\end{lemma}
\begin{proof}
The proved inverse map identifies the corresponding lift over
$O\oplus\epsilon O/b$, $\epsilon^2=0$, as
\[
 r_\epsilon(g)=(1+\epsilon t(g)Y)r_0(g),\qquad
 Y=E_{21}-E_{32},\qquad g\in G_E,                           \tag{KH.6}
\]
with the prescribed normalized strict Selmer cocycle $t$.
The matrix $r_0(g)$ is diagonal, so $r_\epsilon(g)$ is lower
triangular with exactly the same diagonal entries.  Its
characteristic polynomial is therefore equal to that of
$r_0(g)$, as an exact equality over this square-zero coefficient
ring.  This includes traces of every group element, hence also
traces of all products of group matrices.  Equivalently the
displayed first derivative is
\[
             \operatorname{tr}(t(g)Y r_0(g))=0.            \tag{KH.7}
\]
The equality is integral and uses no semisimplification of the
first-order representation.  The product rule proves the claim
for polynomial trace coordinates.  The weight description gives
the same conclusion: such coordinates have torus weight zero,
while this derivation is supported on weight $-1$.
\end{proof}

Consequently restricting these derivations to the trace
subalgebra gives the zero map.  Showing this restriction to be
injective would already prove that the original strict Selmer
module vanishes.  The identity involution in (KH.4) supplies no
such injectivity.  Even an identification of an integral
pseudorepresentation algebra with the displayed Hecke order
would need an additional theorem controlling the off-diagonal
extension modules.  A generalized matrix algebra comparison,
or an equivalent comparison of integral lattices and their
reducibility pairings, could in principle provide those data;
neither the characteristic-zero source coverage nor the actual
point construction is that comparison.  The original classical
rank-two branch is also absent from the stable Hecke algebra and
cannot be removed by (KH.4).

\subsection{The trace pairing and Euler characteristic do not supply a unit}

There is a further concrete obstruction to treating an
irreducible adjoint point as integrally nondegenerate.  It
explains where denominators enter an attempt to recover matrix
entries from traces.

\begin{lemma}[KT-RESIDUAL-TRACE-DISCRIMINANT]
Let $r:G_E\to\mathrm{GL}_3(O')$ be an actual adjoint lift with
the split diagonal residual lattice constructed in
KT-ADJOINT-ACTUAL-HOST-POINT.  Suppose, as holds for that
non-CM lift, that $r$ is absolutely irreducible over the fraction
field $K'$ of the coefficient DVR $O'$.  Let $\mathcal A$ be
the $O'$-span of the matrices $r(g)$.  Then $\mathcal A$ is a
full lattice in $M_3(K')$, and the discriminant of its trace
pairing satisfies
\[
 v_{O'}\operatorname{disc}(\mathcal A)
       =2\operatorname{length}_{O'}(M_3(O')/\mathcal A)
       \geq12.                                             \tag{KH.8}
\]
In particular this trace discriminant is not a unit.
\end{lemma}
\begin{proof}
Absolute irreducibility gives the full matrix algebra as the
$K'$-span of the group, so $\mathcal A$ has rank nine.  Its
reduction inside $M_3(k')$ is contained in the three-dimensional
diagonal subspace.  At least six elementary divisors of its
inclusion in $M_3(O')$ are therefore divisible by the uniformizer,
giving the lower bound on the length.  The matrix-unit basis of
$M_3(O')$ has unimodular pairing $(X,Y)\mapsto\operatorname{tr}(XY)$.
Changing from that basis to a basis of $\mathcal A$ multiplies
the Gram determinant by the square of the inclusion determinant.
Smith normal form proves (KH.8).
\end{proof}

This is a statement about the actual split residual adjoint
lattice, not an evaluation of the odd normal Jacobian.  It
does not imply that the strict Selmer group is nonzero; it
only prevents replacing integral recovery of matrix entries by
the characteristic-zero density theorem without its denominators.

The local and global cohomological dimension formulas used
earlier also do not determine those denominators.  For example,
the two-term complex
$[O\xrightarrow{\varpi^n}O]$ has zero rational cohomology and
the same module ranks for every $n\geq1$, but its integral
cokernel has length $n$.  A perfect finite pairing alone leaves
the same ambiguity: $(x,y)\mapsto xy/\varpi^n\bmod O$ is a
perfect pairing on $O/\varpi^n$ for every $n$.  These are algebraic
scope examples, not counterexamples to the source hypotheses.
For the rational character $\chi^{-e}$ the Tate-dual coefficient
is $\chi^{e+1}$; the duality operation does not identify the
required integral group with a group already proved zero here.

At an adjoint branch the already proved odd normal coefficient
is $\eta\otimes(\mathbf1\oplus\operatorname{Sym}^4\rho
\otimes(\det\rho)^{-2})$.  Characteristic-zero rigidity proves
full generic rank of its normal Jacobian, but no normalized
arithmetic determinant or integral special-value comparison for
that Jacobian has been supplied.  Nor has its base change to the
Eisenstein augmentation been bounded independently of the
original strict Selmer module.  Thus the evaluated unit (KH.5)
does prove the new Hecke identity, while a strict Selmer upper
bound would still require the missing integral extension-module
comparison and an evaluated annihilator or normal determinant.
No stronger arithmetic vanishing conclusion is claimed.

\paragraph{A published comparison route and its failed hypothesis here.}
A targeted primary-source check finds a direct model for the
missing comparison in Hida--Tilouine, \emph{Symmetric Power
Congruence Ideals and Selmer Groups}, author manuscript dated
7 October 2018, pp.~2--4: residual irreducibility is imposed
on p.~2, hypothesis $(*)$ on p.~3 requires residual image
containing $\mathrm{SL}_2(k')$, and Theorems~1.1 and~1.3 use
these hypotheses for the deformation/congruence-ideal comparison.
The source is \url{https://www.math.ucla.edu/~hida/321.pdf}.
Zhang's \emph{Selmer Groups of Symmetric Powers of Ordinary
Modular Galois Representations}, arXiv:1802.08329v1, printed
p.~8, imposes the residual condition
$\operatorname{Sym}^{n-1}\mathrm{SL}_2(\mathbf F')
\subset\operatorname{im}\overline\rho$ in Condition~2.7;
Theorem~2.11 assumes this condition in its $R=T$ conclusion.
The exact version is
\url{https://arxiv.org/pdf/1802.08329v1}.

For the present classical forms,
$\overline\rho_f^{\rm ss}=\mathbf1\oplus\eta\overline\chi^e$.
Every residual lattice is therefore reducible, including a
possible nonsplit extension; its symmetric powers preserve
a nontrivial flag.  Scalar twisting or restricting the Galois
group preserves that flag.  Thus the stated residual-image
hypotheses fail for these branches.  These papers supply no
applicable integral comparison here on their stated hypotheses.
Only those hypotheses and relevant theorem statements were
used in this exclusion audit; their full comparison proofs
have not been imported as new inputs.

% END evidence/unitary_audit/hecke_normal_results.tex


% BEGIN evidence/main_gate/field_variation_nonvanishing.tex
\subsection{Field variation: an exact finite-character reduction}
\label{sec:kt-field-variation}

The field-selection question in this subsection is deliberately narrower
than the full automorphic exclusion question.  It asks for an imaginary
quadratic field whose class number and two specified normalized
generalized Bernoulli numbers are units.  Even a positive answer would
need a separate theorem to exclude every stable rank-three congruence.

Let $\ell>3$ be prime, let $e\geq2$ be even, and assume the actual
\texttt{GD} congruence condition
\[
                         (\ell-1)\nmid 2e.                 \tag{FV.1}
\]
Let $D<0$ be a fundamental discriminant, $f=|D|$, with $\ell\nmid f$,
and put $\chi=\chi_D$.  Fix the Teichm\"uller character
$\omega:(\mathbf Z/\ell\mathbf Z)^\times\to\mathbf Z_\ell^\times$.
Congruences involving its values use this fixed $\ell$-adic embedding.
Write $\beta_{n,\chi}=B_{n,\chi}/n$ for $n\geq1$.

\begin{proposition}[KT-FIELD-CHARACTER-REDUCTION]
For every odd $n\geq3$ with $n\not\equiv1\pmod{\ell-1}$, let $k$ be
the unique odd integer in $[3,\ell-2]$ congruent to $n$ modulo
$\ell-1$.  Then
\[
 \beta_{n,\chi}\equiv\beta_{k,\chi}
       \equiv B_{1,\chi\omega^{k-1}}\pmod\ell.              \tag{FV.2}
\]
All three quantities are $\ell$-integral.  The character in the last
term is primitive of conductor $f\ell$.  In particular,
\[
 \beta_{e+1,\chi_D}\equiv B_{1,\chi_D\omega^e},\qquad
 \beta_{2e+1,\chi_D}\equiv B_{1,\chi_D\omega^{2e}}
                                      \pmod\ell.           \tag{FV.3}
\]
There is no hypothesis comparing $\ell$ with the original weight $n$.
\end{proposition}

\begin{proof}
We use Sun's Lemma~8.1(a), with its difference order equal to one:
if $b\geq1$ and $(\ell-1)\nmid b$, then
\[
 (1-\chi(\ell)\ell^{b-1})\beta_{b,\chi}
 \equiv
 (1-\chi(\ell)\ell^{b+\ell-2})\beta_{b+\ell-1,\chi}
                                  \pmod\ell.               \tag{FV.4}
\]
The conductor hypothesis is exactly $\ell\nmid f$.  The published
statement is in Z.-H.~Sun, \emph{Congruences concerning Bernoulli
numbers and Bernoulli polynomials}, Discrete Applied Mathematics
105 (2000), 193--223, pp.~219--220.  Its full proof was already
checked for KT-TWOFIELD-FAMILY; the exact statement and proof are
reread for the present application.

First all $\beta_{m,\chi}$ for odd $m\geq1$ are $\ell$-integral.
For odd $1\leq b\leq\ell-2$, this follows from
\[
 B_{b,\chi}=f^{b-1}\sum_{a=1}^{f}\chi(a)B_b(a/f).
\]
Indeed the ordinary Bernoulli numbers $B_j$ for $0\leq j\leq\ell-2$
are $\ell$-integral: their elementary recurrence only divides by
$j+1<\ell$.  Also $b$ and $f$ are units.  Induction along each odd
residue class using (FV.4) proves the assertion for every odd $m$.
At the first step from $b=1$, the factor $1-\chi(\ell)$ can be zero;
this causes no difficulty because its product with $\beta_{1,\chi}$
is integral, while the Euler factor at $b+\ell-1$ is a unit.
At every later step both Euler factors are units congruent to one.
Since $\chi$ is odd, $B_{m,\chi}=0$ for even $m$, including $m=0$.
Consequently every $B_{m,\chi}$ is $\ell$-integral.

For $n$ and $k$ in the proposition, (FV.4) can be applied from $k$
to $n$ using only indices at least three.  Both Euler factors are
then congruent to one; hence $\beta_{n,\chi}\equiv\beta_{k,\chi}$.
This includes indices divisible by $\ell$: no division by their
residue classes is performed.

Put
\[
                        r=1+\ell(k-1),\qquad M=f\ell.
\]
Then $r\equiv k\pmod{\ell-1}$, $r\equiv1\pmod\ell$, and $r\geq3$.
The same congruence gives $\beta_{r,\chi}\equiv\beta_{k,\chi}$.
The generalized Bernoulli generating series, summed over $\ell$
complete periods of $\chi$, gives the exact identity
\[
 \frac1M\sum_{a=1}^{M}\chi(a)a^r
   =\sum_{j=0}^{r}\binom rj B_{r-j,\chi}\frac{M^j}{j+1}.
                                                               \tag{FV.5}
\]
For $j\geq1$ one has $j-v_\ell(j+1)\geq1$.  The integrality just
proved therefore implies
\[
 \frac1M\sum_{a=1}^{M}\chi(a)a^r\equiv B_{r,\chi}\pmod\ell.
                                                               \tag{FV.6}
\]
Removing multiples of $\ell$ from this sum changes its divided
value by
$\chi(\ell)\ell^{r-1}f^{-1}\sum_{b=1}^{f}\chi(b)b^r$,
which is divisible by $\ell$.

For every $\ell$-adic unit $a$,
$\omega(a)\equiv a^\ell\pmod{\ell^2}$.  Indeed write
$a=\omega(a)(1+\ell z)$ and raise to the $\ell$th power.
The primitive character $\theta=\chi\omega^{k-1}$ has conductor
$f\ell$: the two factors have coprime conductors, and $k-1$ is
nonzero modulo $\ell-1$.  Its defining formula is
\[
 B_{1,\theta}=\frac1{f\ell}
       \sum_{\substack{1\leq a\leq f\ell\\\ell\nmid a}}
               \chi(a)\omega(a)^{k-1}a.
\]
Replacing $\omega(a)$ by $a^\ell$ changes each summand by a
multiple of $\ell^2$, hence changes the divided sum by a multiple
of $\ell$.  Combining this observation with (FV.6) proves
\[
 B_{1,\theta}\equiv B_{r,\chi}
       =r\beta_{r,\chi}\equiv\beta_{k,\chi}\pmod\ell.
\]
It also proves the asserted integrality of $B_{1,\theta}$.
Finally (FV.1) excludes both $e\equiv0$ and $2e\equiv0$
modulo $\ell-1$, so the proposition applies to $e+1$ and $2e+1$.
\end{proof}

\begin{corollary}[The exact simultaneous field-selection requirement]
Let $\mathcal D$ be any prescribed family of negative fundamental
discriminants prime to $\ell$, with any additional local splitting
or parity conditions.  A member $D\in\mathcal D$ satisfies
\[
 \ell\nmid h(D),\qquad
 \beta_{e+1,\chi_D},\ \beta_{2e+1,\chi_D}\in\mathbf Z_\ell^\times
                                                               \tag{FV.7}
\]
if and only if
\[
 B_{1,\chi_D}\,B_{1,\chi_D\omega^e}\,
 B_{1,\chi_D\omega^{2e}}\not\equiv0\pmod\ell.               \tag{FV.8}
\]
The first twisted character $\chi_D\omega^e$ is never quadratic
under (FV.1).
\end{corollary}
\begin{proof}
The imaginary quadratic class-number formula is
$B_{1,\chi_D}=-2h(D)/w(D)$, with $w(D)\in\{2,4,6\}$.
Since $\ell>3$, its unit condition is exactly $\ell\nmid h(D)$.
Now apply (FV.3).  On the $\ell$-part of the conductor the order
of $\chi_D\omega^e$ is
$(\ell-1)/\gcd(e,\ell-1)>2$, by (FV.1), proving the last assertion.
\end{proof}

\paragraph{What the checked field-variation theorems supply.}
Wiles, \emph{On class groups of imaginary quadratic fields}, JLMS
92 (2015), 411--426, Theorem~A (accepted manuscript
Theorem~0.0.1, pp.~1--2), supplies the first factor in (FV.8),
with specified finite sets of odd split, inert, and ramified primes
subject to its three explicit residue restrictions.  Its original
statement allows $\ell$ in a local set; in particular one can ask
that $\ell$ split.  It is a quadratic-character theorem and does
not provide either nonquadratic twist in (FV.8).  Beckwith's
quantitative refinement, Theorem~1.1 and Corollary~1.2, concerns
the same class-number factor; its restatement excludes $\ell$
from the specified local sets.

Bruinier, Duke Math.~J.~98 (1999), 595--611, Theorem~3
(author PDF p.~11), treats one integral primitive half-integral
weight Hecke eigenform at a time.  For a Cohen series of weight
$k+\tfrac12$ and an odd prescribed local prime $p$, its exceptional
set condition is
\[
 \ell\nmid p(p-1)(p^{2k}-1)(p^{2k-2}-1).                    \tag{FV.9}
\]
This is obtained by substituting $1+p^{2k-1}$ for the Hecke
eigenvalue and factoring the two possible eigenvalue differences.
The theorem also requires the given series to be integral and
primitive at $\ell$; rescaling by a nonunit cannot establish
unit values in the original normalization.  Its Dirichlet
specialization Theorem~6 is stated for even $k$ and real fields,
with an explicit unit seed.  The odd-weight analogue is mentioned
afterwards, rather than supplied as an unconditional imaginary
field theorem.  These precise conditions are recorded in the
separate source audit.

The limitation (FV.9) is genuine for this application of the
theorem: $\ell=17$, $e=4$ satisfies (FV.1), but for the second
index $k=2e+1=9$, every $p\ne17$ has $p^{2k-2}=p^{16}\equiv1$.
Thus this version of the prescribed-local-data theorem cannot
be applied, even though both weights are smaller than $\ell$.
This does not assert that the desired fields fail to exist.

Burungale--Sun, Math.~Ann.~378 (2020), 317--358,
Theorem~5.2 and Corollary~5.4, give quantitative nonvanishing as
\emph{all} characters of a growing conductor vary.  Their family
does not restrict to quadratic characters; for prime conductor
there is only one nontrivial quadratic character, so their lower
bound cannot force that character to contribute.  Their
fixed-base-character prime-power tower results likewise vary
characters of growing nonquadratic order.  Neither statement is
a simultaneous assertion for the three factors in (FV.8).
There is also a direct conductor obstruction to applying
Theorem~5.2 to the fixed character $\omega^e$: that theorem
requires $\ell\nmid N$ for the fixed period $N$, whereas
$\omega^e$ has conductor $\ell$ under (FV.1).

\paragraph{The remaining precise implication.}
The reviewed sources do not establish existence in (FV.8) for
each admissible residue class $e\bmod(\ell-1)$ with the required
local conditions.  Separate infinite sets of discriminants with
individual unit values do not imply that their intersection is
nonempty.  Multiplying two Cohen $q$-series forms a convolution
of coefficients, not the same-discriminant product required in
(FV.8).  This is the exact additional field-selection theorem
needed by this Bernoulli strategy.  It is a scope statement about
the checked results, not a claim that no stronger theorem exists.

Fretwell--Roberts Theorem~3.6 supplies the direction from
Bernoulli divisibility to an Eisenstein congruence. Their
necessary-condition Theorem~3.5 has a large-prime hypothesis
relative to the actual weight; Kummer reduction alone does not
transport a high-weight system into that range. The missing
ordinary implication has now been proved by a separate argument:
KT-ORDINARY-RANK2-UNIT-EXCLUSION uses the exact auxiliary strict
group and an opposite stable lattice, with no weight bound.
KT-AUXILIARY-UNIT-NO-TYPEII matches it to actual Type-II source
points, under GD, $N=2$ and $\psi=1$, using the $a=2e$ unit.
The $a=e$ version also excludes the specified ordinary rank-two
seeds of the symmetric-square construction. Uniform field
selection with the actual GD restrictions, and exclusion of
the remaining stable rank-three sources, are still required.

\begin{lemma}[KT-FIELD-SPECIALIZATION-LIMIT]
Avoidance of common height-one primes by two power series does
not imply that one of their values is a unit at every integral
specialization, even if the two series use independent variables.
\end{lemma}
\begin{proof}
In $A=\mathbf Z_\ell[[X,Y]]$ take $f=\ell+X$ and $g=\ell+Y$.
The quotient $A/(f)$ is $\mathbf Z_\ell[[Y]]$ by the continuous
substitution $X=-\ell$; hence $(f)$ is a nonzero height-one
prime. Any height-one prime containing $f$ must equal $(f)$,
whereas the image $\ell+Y$ of $g$ in that quotient is nonzero.
Thus no height-one prime contains both. At the continuous
integral specialization $X=Y=0$ both values equal $\ell$,
so both are nonunits.
\end{proof}

% END evidence/main_gate/field_variation_nonvanishing.tex


% BEGIN evidence/main_gate/explicit_two_unit_family.tex
\subsection{An explicit two-unit family with either local behavior at the residue prime}
\label{sec:kt-explicit-two-unit}

The following is a partial solution of the simultaneous Bernoulli-unit
question.  It gives two of its three units, at a fixed prime and in an
unbounded weight progression, together with either splitting behavior at
that prime.  The third unit and any additional local conditions are not
asserted.

\begin{theorem}[KT-EXPLICIT-TWO-UNIT-FAMILY]
Let $\ell\geq7$ be a prime with $\ell\equiv3\pmod4$.  For each sign
$s\in\{1,-1\}$, write
\[
 \ell+2s=m_su_s^2,\qquad m_s>0\text{ square-free},\qquad
 D_s=-4m_s,\qquad E_s=\mathbf Q(\sqrt{D_s}).                 \tag{EU.1}
\]
Then $D_s$ is a negative fundamental discriminant and $\ell\nmid D_s$.
For every integer $t\geq0$, put
\[
 e=\frac{\ell-3}{2}+(\ell-1)t.                              \tag{EU.2}
\]
This is a positive even integer satisfying $(\ell-1)\nmid2e$.
The discriminants here are even; as explained below, this arithmetic
family does not satisfy the source's full GD condition when $N$ is even.
With the fixed $\ell$-adic Teichm\"uller character $\omega$, both
\[
 B_{1,\chi_{D_s}}\quad\text{and}\quad
 B_{1,\chi_{D_s}\omega^e}
                                                               \tag{EU.3}
\]
are $\ell$-adic units.  The normalized Bernoulli number
$B_{e+1,\chi_{D_s}}/(e+1)$ is therefore also a unit.
Moreover,
\[
 \chi_{D_+}(\ell)=-\left(\frac2\ell\right),\qquad
 \chi_{D_-}(\ell)= \left(\frac2\ell\right).                 \tag{EU.4}
\]
Thus one of the two fields is split at $\ell$ and the other is inert.
The two nontrivial Teichm\"uller twists $\omega^e$ and $\omega^{2e}$
have order $(\ell-1)/2>2$.

For either sign, the remaining third-unit condition is exactly
\[
 B_{1,\chi_{D_s}\omega^{2e}}\in\mathbf Z_\ell^\times
 \quad\Longleftrightarrow\quad
 \frac{B_{\ell-2,\chi_{D_s}}}{\ell-2}\in\mathbf Z_\ell^\times.
                                                               \tag{EU.5}
\]
Neither side of (EU.5) is claimed by this theorem.
\end{theorem}

\begin{proof}
Both $\ell+2$ and $\ell-2$ are positive odd integers congruent to $1$
modulo $4$.  Their square-free parts $m_s$ therefore also satisfy
$m_s\equiv1\pmod4$.  The integer $-m_s$ is square-free and congruent
to $3$ modulo $4$, so its quadratic field has fundamental discriminant
$D_s=-4m_s$.  This includes the square cases $m_s=1$, when $D_s=-4$.
Since $\ell$ divides neither $\ell+2$ nor $\ell-2$, it does not divide
$D_s$.  Reduction of (EU.1) modulo $\ell$ gives (EU.4), because
$(-1/\ell)=-1$ and $u_s^2$ is a nonzero square modulo $\ell$.

We first check the ordinary class-number factor, including all small
cases.  Theorem~4.1 on PDF p.~14 of Benmerieme--Movahhedi gives,
since $\chi_{D_s}(2)=0$,
\[
 h(D_s)=\frac{w(E_s)}4\sum_{r=1}^{2m_s}\chi_{D_s}(r).
                                                               \tag{EU.6}
\]
Only the $m_s$ odd terms can be nonzero, and each is at most $1$.
If $m_s>1$, then $w(E_s)=2$ and hence
$h(D_s)\leq m_s/2\leq(\ell+2)/2<\ell$.
If $m_s=1$, the same formula gives $h(-4)=1<\ell$.
Therefore $\ell\nmid h(D_s)$, including both signs at $\ell=7$.
The imaginary quadratic class-number formula
$B_{1,\chi_{D_s}}=-2h(D_s)/w(E_s)$, where $w(E_s)=2$ or $4$,
proves the first assertion in (EU.3).

Here is the imported arithmetic input for the second assertion.
Benmerieme--Movahhedi prove that
\[
             F_s=\mathbf Q(\sqrt{\ell(\ell+2s)})
                                                               \tag{EU.7}
\]
is $\ell$-rational for every prime $\ell>3$: their Proposition~3.2
gives the plus sign, and the proof of Proposition~4.4 gives the minus
sign.  Their Corollary~2.2 identifies this with
$L_\ell(1,\chi_{F_s})\in\mathbf Z_\ell^\times$.
Equation~(3.2) then gives, for $n_0=(\ell-1)/2$,
\[
 L_\ell(1,\chi_{F_s})
       \equiv L(1-n_0,\chi_{D_s})
       =-\frac{B_{n_0,\chi_{D_s}}}{n_0}\pmod\ell.          \tag{EU.8}
\]
Indeed the quadratic character of $\mathbf Q(\sqrt{-\ell})$
is $\omega^{n_0}$, and the primitive character underlying
$\chi_{F_s}\omega^{n_0}$ is $\chi_{D_s}$.  This is also equality
as Galois characters.  As an imprimitive Dirichlet function the
product on the left still omits the $\ell$ Euler factor; the
interpolation value therefore includes
$1-\chi_{D_s}(\ell)\ell^{n_0-1}$, which is $1$ modulo $\ell$
because $n_0\geq3$.  Thus it does not change (EU.8).
The hypotheses of that
displayed congruence hold for the square-free negative integer $-m_s$:
it has the prescribed sign, is prime to $\ell$, and is different from
$-\ell$.  Thus $B_{n_0,\chi_{D_s}}/n_0$ is a unit.

Now $n_0$ is odd, $3\leq n_0\leq\ell-2$, and
$e+1\equiv n_0\pmod{\ell-1}$.  Applying the already proved
KT-FIELD-CHARACTER-REDUCTION yields
\[
 \frac{B_{e+1,\chi_{D_s}}}{e+1}
 \equiv\frac{B_{n_0,\chi_{D_s}}}{n_0}
 \equiv B_{1,\chi_{D_s}\omega^e}\pmod\ell.                 \tag{EU.9}
\]
Its integrality statement applies even when the original weight
$e+1$ is divisible by $\ell$ or exceeds $\ell$.  It also proves
(EU.5), since $2e+1\equiv\ell-2\pmod{\ell-1}$.

Finally, $2e\equiv-2\pmod{\ell-1}$ proves the required GD
nondivisibility.  Writing $\ell=4a+3$ gives
$\gcd(\ell-1,e)=\gcd(4a+2,2a)=2$, so $\omega^e$ has order
$(\ell-1)/2$, an odd integer at least $3$.  Squaring does not change
this order, which proves the final assertion concerning the twists.
\end{proof}

\paragraph{Exact primary locators and check scope.}
Y.~Benmerieme and A.~Movahhedi, \emph{Multi-quadratic $p$-rational
number fields}, Journal of Pure and Applied Algebra 225 (2021),
106657, DOI \path{10.1016/j.jpaa.2020.106657}.
The version used here is the authors' arXiv:2007.04864v1,
9 July 2020, \url{https://arxiv.org/pdf/2007.04864v1}:
Corollary~2.2 on PDF p.~5 and its preceding derivation on pp.~4--5;
Proposition~3.2 and its entire proof on pp.~11--12;
equation~(3.2) and its explanatory paragraph on p.~12;
Theorem~4.1 on p.~14; Proposition~4.4 and its entire proof on p.~15.
The present proof imports those arithmetic statements, not fresh
proofs of their cited Coates or $p$-adic $L$-function results.
The first class-number unit follows directly from the finite character
sum (EU.6), with both square cases included.
This theorem does not prescribe local behavior at any prime other than
$\ell$, does not prove (EU.5), and does not exclude every stable
automorphic congruence even if that remaining unit is later verified.

\begin{proposition}[KT-EXPLICIT-TWO-UNIT-GD-SCOPE]
Neither discriminant in (EU.1) is an admissible auxiliary discriminant
for the actual source GD condition with even $N$, in particular with
$N=2$.
\end{proposition}
\begin{proof}
The actual source hypotheses \path{GD.N_even} and
\texttt{GD.cop} state that $N$ is even and coprime to the absolute
discriminant of the auxiliary quadratic field.  Thus $2\mid N$
and $\gcd(N,|D|)=1$ imply $2\nmid D$.  This exact deduction is
already written in the pinned source
\path{UnitaryThreeGroups/AdmissibleOfTypeII.lean:119} as
\path{Nat.coprime_two_right.1
(D.gd.cop.coprime_dvd_right D.gd.N_even.two_dvd)}.
But (EU.1) has $D_s=-4m_s$, divisible by four for both signs,
including $m_s=1$.  Therefore the family proves the two asserted
Bernoulli units and the exponent nondivisibility, but supplies no
field satisfying the full GD hypotheses.  The bounded third-unit
certificate for this family has the same parity limitation.
\end{proof}

% END evidence/main_gate/explicit_two_unit_family.tex


% BEGIN evidence/main_gate/odd_discriminant_two_unit_family.tex
\subsection{An odd-discriminant two-unit family satisfying the actual GD conditions}
\label{sec:kt-odd-discriminant-two-unit}

The even discriminants in KT-EXPLICIT-TWO-UNIT-FAMILY cannot satisfy
the source's coprimality condition at $N=2$.  The following related
construction repairs that specific obstruction.  Its elementary unit
calculation is given in full; the real class-number estimate and the
$p$-rationality criterion are imported at the exact locators below.

\begin{theorem}[KT-ODD-DISCRIMINANT-TWO-UNIT-FAMILY]
Let $p\geq7$ be a prime with $p\equiv3\pmod4$.  For each
$s\in\{1,-1\}$ write
\[
 p+4s=m_su_s^2,\qquad m_s>0\text{ square-free},\quad u_s>0,
 \qquad D_s=-m_s,\qquad E_s=\mathbf Q(\sqrt{D_s}).       \tag{OD.1}
\]
Then $m_s\equiv3\pmod4$, $m_s\geq3$, and $D_s$ is an odd
negative fundamental discriminant prime to $p$.  Moreover
\[
 0<h(E_s)<p,\qquad
 \chi_{D_+}(p)=-1,\qquad \chi_{D_-}(p)=1.                \tag{OD.2}
\]
For every integer $t\geq0$ put
\[
 e=\frac{p-3}{2}+(p-1)t.                                \tag{OD.3}
\]
Then $e$ is positive and even, $(p-1)\nmid2e$, and both
\[
 B_{1,\chi_{D_s}},\qquad B_{1,\chi_{D_s}\omega^e}
                                                            \tag{OD.4}
\]
are $p$-adic units.  Equivalently for the second term,
$B_{e+1,\chi_{D_s}}/(e+1)$ is a $p$-adic unit.
The twists $\omega^e$ and $\omega^{2e}$ both have order
$(p-1)/2>2$.  The remaining third-unit condition is precisely
\[
 B_{1,\chi_{D_s}\omega^{2e}}\in\mathbf Z_p^\times
 \quad\Longleftrightarrow\quad
 \frac{B_{p-2,\chi_{D_s}}}{p-2}\in\mathbf Z_p^\times.     \tag{OD.5}
\]
No assertion of (OD.5) is part of this theorem.
\end{theorem}

\begin{proof}
Both integers $p+4s$ are positive, odd, and congruent to $3$ modulo
$4$.  Thus $u_s$ is odd and $m_s\equiv3\pmod4$, which also gives
$m_s\geq3$.  Since $p\nmid p+4s$, neither $m_s$ nor $u_s$ is
divisible by $p$.  The discriminant of $E_s$ is consequently the
odd fundamental discriminant $D_s=-m_s$.  Modulo $p$ we have
$D_s\equiv-4s/u_s^2$, so the two local symbols in (OD.2) follow
from $(-1/p)=-1$.

We first bound the imaginary class number.  For $m_s>3$, the field
$E_s$ has exactly two roots of unity.  The finite class-number
formula of Benmerieme--Movahhedi, Theorem~4.1, gives
\[
 h(E_s)=\frac{1}{2-\chi_{D_s}(2)}
       \sum_{r=1}^{(m_s-1)/2}\chi_{D_s}(r)
 \leq\frac{m_s-1}{2}\leq\frac{p+3}{2}<p.                \tag{OD.6}
\]
Here $2-\chi_{D_s}(2)$ is either $1$ or $3$, and each summand
is at most $1$.  For $m_s=3$ the same formula, with six roots of
unity, gives $h(-3)=1$.  Positivity of class numbers proves (OD.2).
The imaginary quadratic formula
$B_{1,\chi_{D_s}}=-2h(E_s)/w(E_s)$, with $w(E_s)\in\{2,6\}$,
therefore proves the first unit in (OD.4), since $p\geq7$.

For the second unit consider the real quadratic field and its unit
\[
 d_s=pm_s,\qquad F_s=\mathbf Q(\sqrt{d_s}),\qquad
 \alpha_s=\frac{p+2s+u_s\sqrt{d_s}}2
          =\frac{p+2s+\sqrt{p(p+4s)}}2.                 \tag{OD.7}
\]
The integer $d_s$ is square-free and congruent to $1$ modulo $4$,
so it is the fundamental discriminant of $F_s$.  The numerator
coefficients in (OD.7) are odd, hence $\alpha_s$ is integral.
Its norm is $((p+2s)^2-p(p+4s))/4=1$, and $\alpha_s>1$.

In fact $\alpha_s$ is the fundamental unit of $F_s$.
No integral unit can have norm $-1$: writing it as
$(x+y\sqrt{d_s})/2$ would give $x^2-d_sy^2=-4$, which is
impossible modulo $p$ because $(-1/p)=-1$.  Let $\varepsilon>1$
be the fundamental unit, of norm $1$.  Write
$\varepsilon=(x+y\sqrt{d_s})/2$, with integers $x>0$ and
$y\geq1$.  If $\alpha_s=\varepsilon^n$ for $n\geq2$, the
strict monotonicity of $z+z^{-1}$ for $z>1$ gives
\[
 p+2s=\operatorname{Tr}(\alpha_s)
 \geq\operatorname{Tr}(\varepsilon^2)
 =x^2-2=d_sy^2+2\geq3p+2>p+2,                         \tag{OD.8}
\]
a contradiction.  Thus $n=1$.

There is a unique prime $v$ of $F_s$ above $p$, with ramification
index two and residue field $\mathbf F_p$.  In the completion
$F_{s,v}$ put $\pi=\sqrt{d_s}$ and normalize $v_\pi(\pi)=1$;
then $v_\pi(p)=2$ and $u_s$ is a local unit.  Directly from
(OD.7),
\[
 s\alpha_s-1=\frac{s(p+u_s\pi)}2,
 \qquad v_\pi(s\alpha_s-1)=1.                           \tag{OD.9}
\]
Suppose $\alpha_s=z^p$ in $F_{s,v}$.  The element $z$ is a unit
and its residue is $s$, because its $p$-th power has that residue
in $\mathbf F_p$.  Thus $sz=1+w$ with $v_\pi(w)\geq1$, and
$(1+w)^p=s\alpha_s$.  For $1\leq j<p$,
$v_\pi(\binom pj w^j)\geq2+j\geq3$, while
$v_\pi(w^p)\geq p\geq7$.  The binomial expansion therefore
gives $v_\pi((1+w)^p-1)\geq3$, contradicting (OD.9).
This proves that the fundamental unit is not a local $p$-th power,
without a logarithm convention or an additional weight hypothesis.

Le's Theorem~(a) applies to every real quadratic field and gives
\[
 h(F_s)\leq\left\lfloor\frac{\sqrt{d_s}}2\right\rfloor
 \leq\frac{\sqrt{p(p+4)}}2<p.                            \tag{OD.10}
\]
Its published correction changes only the attribution of one
reference, not the theorem.  Consequently $p\nmid h(F_s)$.
Benmerieme--Movahhedi, Corollary~2.6, now proves that $F_s$ is
$p$-rational; its hypotheses require only $p\geq5$ in the present
real quadratic setting.  Their Corollary~2.2 consequently gives
\[
                    L_p(1,\chi_{F_s})\in\mathbf Z_p^\times.
                                                            \tag{OD.11}
\]

Set $n_0=(p-1)/2$.  The primitive product of the quadratic
characters $\chi_{F_s}$ and $\omega^{n_0}=\chi_{-p}$ is
$\chi_{D_s}$: indeed $(pm_s)(-p)=-m_sp^2$.
This is equality of primitive characters, equivalently of Galois
characters.  Equation~(3.2) of Benmerieme--Movahhedi therefore gives
\[
 L_p(1,\chi_{F_s})\equiv
       -\frac{B_{n_0,\chi_{D_s}}}{n_0}\pmod p.          \tag{OD.12}
\]
As an imprimitive Dirichlet function the product still omits the
$p$ Euler factor.  That factor is
$1-\chi_{D_s}(p)p^{n_0-1}\equiv1\pmod p$, since $n_0\geq3$,
so (OD.12) has exactly the displayed normalization.  The square-free
integer $-m_s$ satisfies the hypotheses preceding that equation:
it is negative, prime to $p$, and different from $-p$.

Finally apply KT-FIELD-CHARACTER-REDUCTION, since
$e+1\equiv n_0\pmod{p-1}$ and $2e+1\equiv p-2\pmod{p-1}$.
It gives the second unit in (OD.4) and the equivalence (OD.5),
including weights divisible by $p$ and weights larger than $p$.
Writing $p=4a+3$ gives $e\equiv2a\pmod{p-1}$ and
$\gcd(4a+2,2a)=2$.  Hence the two twists have the asserted common
odd order $(p-1)/2$, and $2e\equiv-2\pmod{p-1}$ proves the
nondivisibility assertion.
\end{proof}

\begin{corollary}[KT-ODD-DISCRIMINANT-GD]
For every field $E_s$, weight $e$, and prime $p$ of the theorem,
the actual source conditions $\mathrm{GD}$ hold at $N=2$ and
$\psi=1$, after choosing the usual embeddings and coefficient
data.  Either prescribed splitting behavior at $p$ is available.
\end{corollary}
\begin{proof}
We check all twelve fields of
\path{HeightAndVanishing/Counts.lean:9--22}.
The prime is odd, the weight is positive, and $N=2$ is even.
The character $\psi=1$ has order one, its prime-to-$p$ and
$p$-power parts are trivial, and its parity agrees with even $e$.
The field is imaginary quadratic.  By (OD.1)--(OD.2),
$p\nmid |D_s|h(E_s)$ and $\gcd(|D_s|,2)=1$.
The final genericity condition is $(p-1)\nmid2e$, already proved.

For the remaining field \texttt{eta}, the source requires that
the intersection of the kernels of the residual cyclotomic
character and the trivial Dirichlet character is not contained
in $G_{E_s}$.  Choose any prime $q\mid m_s$; such a prime exists
because $m_s\geq3$, and $q\ne p$.  The quadratic character
$\chi_{D_s}$ is nontrivial on inertia at $q$, since $q$ divides
the fundamental discriminant.  The residual cyclotomic character
at $p$ is unramified there, and the Dirichlet character is trivial.
An inertia element on which $\chi_{D_s}$ is $-1$ therefore
belongs to the specified intersection but not to $G_{E_s}$.
This is precisely the remaining noncontainment condition.
\end{proof}

\paragraph{Arithmetic scope.}
For these actual GD fields, the already proved
KT-ORDINARY-RANK2-UNIT-EXCLUSION applies with $a=e$ to every
representation satisfying that theorem's explicit residual,
ordinary-quotient, and away-$p$ conditions.  Thus the constructed
ordinary lower-weight $\mathrm{GL}_2$ source branch is excluded
where those conditions hold.  The actual Type~II exclusion uses
$a=2e$ and still requires the third unit (OD.5).
Neither the third unit nor an exclusion of all genuinely rank-three
sources follows from this two-unit theorem.  No additional local
conditions away from $p$ are prescribed.

\paragraph{Exact primary locators and check scope.}
Maohua Le, \emph{Upper bounds for class numbers of real quadratic
fields}, Acta Arithmetica 68 (1994), 141--144,
DOI \path{10.4064/aa-68-2-141-144}: Theorem~(a) on printed
p.~141; its complete proof is in Section~3, pp.~142--143.
The complete four-page paper was read.  The numerical initial
range and the analytic lemmas cited there are imported, not
independently recomputed.  Le's one-page correction, Acta
Arithmetica 72 (1995), p.~399,
DOI \path{10.4064/aa-72-4-399}, was read completely; it corrects
the author of reference~[2] to B.~Oriat and makes no mathematical
change.  Publisher PDFs are saved with retrieval metadata under
\path{evidence/main_gate/literature/odd_discriminant/}.

Y.~Benmerieme and A.~Movahhedi, \emph{Multi-quadratic $p$-rational
number fields}, Journal of Pure and Applied Algebra 225 (2021),
106657, DOI \path{10.1016/j.jpaa.2020.106657}; version
arXiv:2007.04864v1, 9 July 2020,
\url{https://arxiv.org/pdf/2007.04864v1}:
Corollary~2.2 and preceding derivation, PDF pp.~4--5;
Proposition~2.3, Proposition~2.4, and Corollary~2.6 with their
intervening derivation, pp.~5--6; equation~(3.2) and its full
context, p.~12; Theorem~4.1, p.~14.
The authors mention the field $\mathbf Q(\sqrt{p(p-4)})$ on
p.~14, without supplying the proof used here.  The apparently
missing display after the last ``same method'' sentence on p.~15
is also blank in the rendered PDF; no unstated theorem is inferred
from it.  The present proof supplies both signs directly.

% END evidence/main_gate/odd_discriminant_two_unit_family.tex


% BEGIN evidence/main_gate/odd_family_reflection_scope.tex
\begin{proposition}[KT-ODD-FAMILY-RATIONAL-MIRROR]
The weight progression in KT-ODD-DISCRIMINANT-TWO-UNIT-FAMILY
already satisfies the original rational class-component vanishing
by the reflection argument of K17-REFLECTOR.  More precisely, if
$p\geq7$, $p\equiv3\pmod4$, and
$e=(p-3)/2+(p-1)t$ with $t\geq0$, then
\[
 A_p^{[1+e]}=A_p^{[-e]}=0,
 \qquad K_{2e}(\mathbf Z)\{p\}=0.                     \tag{ODR.1}
\]
This uses neither auxiliary field $E_s$ nor the third unit (OD.5).
\end{proposition}
\begin{proof}
Put $q=(p-1)/2$, $F=\mathbf Q(\mu_p)$, and
$K=\mathbf Q(\sqrt{-p})$.  The character $\omega^q$ is the
quadratic character of $K$, and $1+e\equiv q\pmod{p-1}$.
The imaginary quadratic class-number formula used in (OD.6), now
with discriminant $-p$, gives $0<h(K)\leq(p-1)/2<p$.
Thus the $p$-primary class group of $K$ is zero.

Let $H=\operatorname{Gal}(F/K)$, of order $q$ prime to $p$.
Extension of ideals and norm satisfy
$j\circ N=\sum_{\sigma\in H}\sigma$ on the class group of $F$.
For an $H$-invariant $p$-primary class $x$ this gives
$q x=jN(x)=0$, since $N(x)$ is a $p$-primary class of $K$.
Multiplication by $q$ is invertible on a finite $p$-group, so
$A_p^H=0$.  In particular $A_p^{[q]}=0$.

Apply the reflection half of the already proved K17-REFLECTOR
at $e_0=(p-3)/2$.  This is even and lies in $[2,p-3]$;
its mirror character is $\omega^{1+e_0}=\omega^q$.
It follows that $A_p^{[-e_0]}=0$.  The character is unchanged
on replacing $e_0$ by $e$, and the carried K17-K2-PERIODIC
identification gives the last assertion of (ODR.1).
\end{proof}

\paragraph{Scope of the new auxiliary construction.}
The odd-discriminant theorem supplies actual GD fields and two
specified unit conditions, with both local splitting choices.
It does not remove a previously uncontrolled rational component
in this particular progression: (ODR.1) is a direct consequence
of the already accepted quadratic class-number and reflection
inputs.  A proof of the third unit would still be a separate
arithmetic statement about the auxiliary fields.  No full
third-unit assertion, no uniform result for arbitrary weight
residues, and no new general rank-three exclusion follows here.

% END evidence/main_gate/odd_family_reflection_scope.tex


% BEGIN evidence/main_gate/jacobi_two_unit_fields.tex
\subsection{Two auxiliary units from a Jacobi heat lift}
\label{sec:kt-jacobi-fields}

The following construction varies the imaginary quadratic field.  Its
extra rational Bernoulli condition is $B_{2b}/(2b)$, whereas the
original rational reflection component at $e\equiv b-1$ is controlled
by $B_{p-b}/(p-b)$.  No unit condition on the latter number is assumed.
The construction does not assert simultaneous units for both twists
$\omega^e$ and $\omega^{2e}$ over the same field.

\begin{lemma}[KT-SMALL-DISCRIMINANT-CLASS-BOUND]
Let $D<0$ be an odd fundamental discriminant, put $d=|D|$ and
$R=\lfloor\sqrt{d/3}\rfloor$.  Then
\[
 h(D)\leq\frac{R(R+3)}2
       \leq\frac d6+\frac{\sqrt{3d}}2.                 \tag{JF.1}
\]
Consequently, if $p\geq29$ and $d<4p$, then $h(D)<p$.
\end{lemma}
\begin{proof}
Here is the elementary form-counting argument, including the
representation of ideal classes that the count uses.  For an integral
ideal $I$ and an oriented integral basis $\alpha,\beta$, the form
\[
 Q_I(x,y)=N(x\alpha+y\beta)/N(I)=ax^2+bxy+cy^2
\]
is integral and positive definite, and has discriminant $D$.
Integrality follows because $(u)\subseteq I$ implies $N(I)\mid N(u)$,
also applied to $u=\alpha+\beta$.  The discriminant follows from
the covolume $N(I)\sqrt d/2$.  An oriented basis change is a proper
equivalence.  Conversely, a proper equivalence between two such norm
forms gives an orientation-preserving similarity between their complex
ideal lattices, hence multiplication by some $\lambda\in\mathbf C$.
Since it maps a nonzero element of one ideal to an element of the other,
$\lambda\in\mathbf Q(\sqrt D)$; the two ideals have the same class.
Thus ideal classes inject into proper equivalence classes of these forms.

In each form class choose a primitive vector with smallest positive
value $a$ and extend it to an oriented integral basis.  The vector is
primitive because otherwise division by its common divisor would
decrease its positive value.  Changing the second basis vector by an
integer multiple of the first arranges $|b|\leq a$.  Its value $c$ is
at least $a$ by minimality.  Hence $d=4ac-b^2\geq3a^2$ and $a\leq R$.
For a fixed $a$, the odd integer $b$ lies in $[-a,a]$, giving at most
$a+1$ choices; $c=(b^2+d)/(4a)$ is then determined.  Counting all these
pairs, even the ones that are not integral forms, proves (JF.1).
If $d<4p$, its right side is less than $2p/3+\sqrt{3p}$, which is
less than $p$ when $p>27$.
\end{proof}

\begin{lemma}[KT-JACOBI-ODD-COEFFICIENT]
Let $p\geq5$, let $K$ be even, and let
$\Phi=\sum a(n,r)q^n\zeta^r\in J_{K,1}(\mathrm{SL}_2(\mathbf Z))$
have rational $p$-integral coefficients with bounded denominators.
If $\Phi\not\equiv0\pmod p$, then its odd-$r$ theta component is
nonzero modulo $p$.  Moreover, some coefficient $a(n,r)$ with $r$ odd
and $0<n\leq K/2$ is a $p$-adic unit.
\end{lemma}
\begin{proof}
Use $e(x)=\exp(2\pi i x)$ and the usual Jacobi slash operator.
Put $g(\tau,z)=\Phi(\tau,z+1/2)$ and $F=\Phi-g$.
The elliptic transformation law shows that $g$ has index one for
the full integral translation lattice.  For
$\gamma=\left(\begin{smallmatrix}a&b\\c&d\end{smallmatrix}\right)
\in\Gamma_0(4)$, the modular law followed by the integral translation
$(c/2,(d-1)/2)$ gives exactly
\[
                   g|_{K,1}\gamma=e(c(d-2)/4)g=g.       \tag{JF.2}
\]
Thus $F$ is a Jacobi form for $\Gamma_0(4)$, and in particular for
$\Gamma_1(2^2)\ltimes(2\mathbf Z\times\mathbf Z)$.
Holomorphy at its cusps follows directly by transforming to rational
elliptic translates: the resulting exponents are
$n+r\lambda+\lambda^2=(n-r^2/4)+(\lambda+r/2)^2\geq0$.

Suppose all odd-$r$ coefficients of $\Phi$ vanish modulo $p$.
Then $F\equiv0\pmod p$.  Schwagenscheidt, Theorem~1.2
(statement p.~2, complete proof pp.~11--12), applied with $N=2,m=1$,
implies $F|_{K,1}S\equiv0\pmod p$ for
$S=\left(\begin{smallmatrix}0&-1\\1&0\end{smallmatrix}\right)$.
Its integrality hypothesis is met by multiplying by a common
denominator prime to $p$.  An explicit calculation gives
\[
 F|_{K,1}S=\Phi(\tau,z)
       -q^{1/4}\zeta\Phi(\tau,z+\tau/2).                \tag{JF.3}
\]
Modulo $p$, the first summand has only even powers of $\zeta$ and
the second only odd powers.  Their vanishing difference forces
$\Phi\equiv0$, a contradiction.  This proves the first assertion.

Write the index-one theta decomposition as
\[
 \Phi=h_0\vartheta_0+h_1\vartheta_1,\qquad
 \vartheta_\mu(\tau,z)=
       \sum_{r\equiv\mu\ (2)}q^{r^2/4}\zeta^r.
\]
Then $h_1\not\equiv0\pmod p$, and
\[
 f(\tau)=\frac{\Phi(\tau,0)-\Phi(\tau,1/2)}2
        =h_1(\tau)\vartheta_1(\tau,0)
        \in M_K(\Gamma_0(4)).                           \tag{JF.4}
\]
The character is trivial, by (JF.2) at $z=0$; there is no unrecorded
multiplier.  The series $\vartheta_1(\tau,0)$ starts with $2q^{1/4}$,
so (JF.4) is nonzero modulo $p$; cancellation between the odd-$r$
coefficients cannot make it zero.  It is $p$-integral.
The ordinary Sturm bound gives a nonzero coefficient of $f$ at
$n\leq K[\mathrm{SL}_2(\mathbf Z):\Gamma_0(4)]/12=K/2$.
This is the usual elliptic Sturm theorem, stated in exactly this
coefficient-field generality on p.~129 of Choi--Choie--Kikuta (2013),
with the original reference Sturm (1987), pp.~275--280.
That coefficient is $\sum_{r\text{ odd}}a(n,r)$, so one of its
summands is a unit.  The Jacobi support condition rules out $n=0$.
\end{proof}

\begin{theorem}[KT-JACOBI-TWO-UNIT-FIELD]
Let $p\geq29$ be prime and let $b$ be odd with
\[
 3\leq b\leq p-2,\qquad b\ne(p-1)/2,
 \qquad B_{2b}/(2b)\in\mathbf Z_p^\times.               \tag{JF.5}
\]
There exists an odd negative fundamental discriminant $D$, prime to
$p$, such that
\[
 |D|\leq2(b+p+2)-1\leq4p-1,\qquad
 B_{1,\chi_D},\quad B_{1,\chi_D\omega^{b-1}}
                      \in\mathbf Z_p^\times.           \tag{JF.6}
\]
Equivalently, $p\nmid h(D)$ and
$B_{b,\chi_D}/b\in\mathbf Z_p^\times$.
The same field has $B_{n,\chi_D}/n$ a unit for every odd $n\geq3$
congruent to $b$ modulo $p-1$, with no bound on $n$.
\end{theorem}
\begin{proof}
Put $k=b+1$ and take the normalized Jacobi Eisenstein series
\[
                \phi=\zeta(1-2b)E_{b+1,1}\in J_{k,1}.
                                                               \tag{JF.7}
\]
The index-one coefficient formula of Eichler--Zagier, Theorem~2.1,
is also recalled for fundamental discriminants by Schwagenscheidt
on p.~3.  Its general-discriminant version is explicit in the
1983 MPIM preprint: the definition of $L_D(s)$ on printed p.~19a
(PDF p.~24), followed by Theorem~2.1 on printed p.~20 (PDF p.~25).
In full, if $r^2-4n=D_0u^2<0$ with $D_0$ fundamental,
then the coefficient of (JF.7) is
\[
 L(1-b,\chi_{D_0})
 \sum_{v\mid u}\mu(v)\chi_{D_0}(v)v^{b-1}
                      \sigma_{2b-1}(u/v).               \tag{JF.8}
\]
For zero discriminant it is $\zeta(1-2b)$.
The coefficient denominators are bounded; this property of (JF.7)
is also stated on p.~3 of Schwagenscheidt.

We verify $p$-integrality even when $p\mid D_0$, so this step does
not depend on an omitted condition on the discriminant.  For any
quadratic primitive character $\chi$ of conductor $f$, the Bernoulli
polynomial formula is
\[
 \frac{B_{b,\chi}}b
 =\frac1b\sum_{j=0}^{b}\binom bj B_j f^{j-1}
                     \sum_{a=1}^{f}\chi(a)a^{b-j}.       \tag{JF.9}
\]
Since $b\leq p-2$, every $B_j$ here is $p$-integral, as follows
from its elementary recurrence, and $b$ is a unit.  If $p\nmid f$
the assertion follows immediately.  If $p\mid f$, then $v_p(f)=1$
because $f$ is a fundamental discriminant and $p$ is odd.
Only the $j=0$ term can have a denominator $p$.  Decompose the
primitive quadratic character by the Chinese remainder theorem
as $\chi_p\chi_{f/p}$.  Modulo $p$, its numerator is a product
of the sums of $\chi_p(a)a^b$ over $a\in\mathbf F_p^\times$
and of $\chi_{f/p}$ over its complete period.  The latter sum is
zero unless $f/p=1$.  In that case the first sum is
$\sum_{a\in\mathbf F_p^\times}a^{b+(p-1)/2}$, which is zero
unless $b=(p-1)/2$.  That exception was excluded in (JF.5).
Thus (JF.9), and consequently every coefficient in (JF.8), is
$p$-integral.  The zero-discriminant coefficient is a unit by (JF.5).
In particular $\phi\not\equiv0\pmod p$.

Richter (2009), Proposition~1 on p.~872, gives the weight
congruence for nonzero Jacobi forms modulo $p$.
It implies that the filtration of $\phi$ is $k$, since
$4\leq k\leq p-1$.  At $k=p-1$ the only lower possible weight
is zero, and $J_{0,1}=0$; this also follows from the free
$E_{4,1},E_{6,1}$ decomposition recalled on p.~871 of Richter.
Proposition~2, including its proof on p.~873, now gives
\[
 \Omega(L_1\phi)=k+p+1=:K,\qquad L_1\phi\not\equiv0\pmod p,
                                                               \tag{JF.10}
\]
because $p\nmid 2k-1=2b+1$.  In the indicated range the only
failure of this last condition is the same excluded $b=(p-1)/2$.
For precision an actual lift is
\[
 \widehat\phi=L_1\phi-\frac{2k-1}{6}E_2\phi\in J_{k+2,1},
 \qquad
 \Phi=E_{p-1}\widehat\phi+
                   \frac{2k-1}{6}E_{p+1}\phi\in J_{K,1}.
                                                               \tag{JF.11}
\]
Here $\Phi\equiv L_1\phi\pmod p$, using
$E_{p-1}\equiv1$ and $E_{p+1}\equiv E_2$; all coefficients
are $p$-integral.  Formula (JF.11) is the explicit lift used in
Richter's proof, rather than merely a formal heat series.

Apply KT-JACOBI-ODD-COEFFICIENT to $\Phi$.  It gives a unit
coefficient at $(n,r)$ with $r$ odd and $n\leq K/2$.
Modulo $p$ that coefficient equals
$(4n-r^2)c_\phi(n,r)$.  Therefore $p\nmid4n-r^2$, and
$c_\phi(n,r)$ is itself a unit.  The integer
$D'=r^2-4n$ is negative and odd, and
\[
                    0<|D'|\leq4n-1\leq2K-1\leq4p-1.
                                                               \tag{JF.12}
\]
Write $D'=Du^2$ with $D$ fundamental.  Then $D$ is odd,
$p\nmid Du$, and $|D|\leq|D'|$.
The multiplier in (JF.8) is an integer.  Since
$L(1-b,\chi_D)=-B_{b,\chi_D}/b$ is $p$-integral,
the unit coefficient forces $B_{b,\chi_D}/b$ to be a unit.
KT-SMALL-DISCRIMINANT-CLASS-BOUND gives $h(D)<p$, hence the
first unit in (JF.6), by $B_{1,\chi_D}=-2h(D)/w(D)$.
Finally KT-FIELD-CHARACTER-REDUCTION, (FV.2), gives the second
unit in (JF.6) and the asserted transport to all weights $n$.
\end{proof}

\begin{corollary}[KT-JACOBI-TWO-UNIT-GD]
Under (JF.5), suppose also $(p-1)\nmid2(b-1)$ and let
$e=b-1+(p-1)t$, $t\geq0$.  The field supplied by (JF.6) gives
actual \texttt{GD} data for $N=2$, $\psi=1$, and this $e$.
It has the class-number unit and the unit
$B_{1,\chi_D\omega^e}$.

In particular, for every prime $p\geq31$ with $p\equiv3\pmod4$
and every
\[
                      e=(p+1)/2+(p-1)t,\qquad t\geq0,    \tag{JF.13}
\]
such an actual GD field exists without any additional rational
Bernoulli hypothesis.  Here $\omega^{1+e}$ has order $p-1$,
and $\omega^e$ has order $(p-1)/2$.
\end{corollary}
\begin{proof}
The twelve clauses of \texttt{GD} are the ones in
\path{Counts.lean}, lines 9--22.  The verification in
KT-ODD-DISCRIMINANT-GD applies: the discriminant is odd and
prime to $p$, its class number is prime to $p$, $e$ is positive
and even, and the required nondivisibility is explicitly assumed.
For $\psi=1$ the order and square-compatibility clauses are trivial;
$N=2$ is even and coprime to $D$.  The quadratic and imaginary
clauses hold by construction.  For the remaining \texttt{eta}
clause, $E$ cannot be contained in $\mathbf Q(\mu_p)$: some
prime dividing $D$ ramifies in $E$ and is different from $p$,
whereas $\mathbf Q(\mu_p)$ is unramified there.  This also makes
explicit that no splitting-at-$p$ assumption is used.

For (JF.13), set $b=(p+3)/2$.  It is odd, is in the required
range, and is not $(p-1)/2$.  The ordinary Kummer congruence
(the already checked Sun lemma with trivial character) gives
\[
 \frac{B_{2b}}{2b}=\frac{B_{p+3}}{p+3}
       \equiv\frac{B_4}{4}=-\frac1{120}\pmod p.          \tag{JF.14}
\]
This is a unit.  Also $2e\equiv2\pmod{p-1}$.
Writing $q=(p-1)/2$, which is odd, gives
$\gcd(2q,q+2)=1$ and $\gcd(2q,q+1)=2$.
These are exactly the stated character orders.  In particular
the mirror in (JF.13) is nonquadratic; no vanishing of its
rational class component was an input to the construction.
\end{proof}

\begin{corollary}[KT-JACOBI-TYPEII-FIELD]
Let $p\geq31$, let $e\geq2$ be even with $(p-1)\nmid2e$,
and let $b$ be the odd representative in $[3,p-2]$ of $2e+1$.
If $b\ne(p-1)/2$ and $B_{2b}/(2b)$ is a $p$-adic unit,
there is an odd imaginary quadratic GD field for $N=2$, $\psi=1$
over which the actual source has no TypeII points.
In particular this holds, without an extra Bernoulli hypothesis, for
\[
 p\equiv7\pmod8,\qquad
                e=(p+1)/4+(p-1)t,\qquad t\geq0.        \tag{JF.15}
\]
\end{corollary}
\begin{proof}
Apply KT-JACOBI-TWO-UNIT-FIELD to this $b$.  Its field has
$p\nmid h(D)$ and $B_{1,\chi_D\omega^{2e}}$ a unit.
The same twelve-clause verification just given supplies GD for
the specified $e$, whose genericity was assumed.  The previously
proved KT-AUXILIARY-UNIT-NO-TYPEII applies to these actual source
data, including its finite-coefficient descent and full rational
ordinary flag; it gives the conclusion about TypeII points.
For (JF.15), $e$ is even, $2e+1\equiv(p+3)/2$, and (JF.14)
supplies the seed.  Here $2e\equiv(p+1)/2\not\equiv0\pmod{p-1}$.
This corollary does not assert that $B_{1,\chi_D\omega^e}$ is a
unit for the same field, and does not exclude genuine rank-three
sources.
\end{proof}

\paragraph{Primary sources and exact scope.}
M.~Eichler and D.~Zagier, \emph{On the theory of Jacobi forms I},
MPIM Preprint 1983(8), printed pp.~19a--20 (PDF pp.~24--25),
gives the general-discriminant $L_D$ definition and Theorem~2.1
used in (JF.8).  The public primary scan is
\url{https://archive.mpim-bonn.mpg.de/id/eprint/1437/1/preprint_1983_8.pdf}.
It is a precursor of their 1985 monograph \emph{The Theory of Jacobi
Forms}, Progress in Mathematics~55; the displayed page locators here
refer to the 1983 preprint, not to the monograph.
O.~K.~Richter, \emph{The action of the heat operator on Jacobi forms},
Proc. Amer. Math. Soc.~137 (2009), 869--875,
doi:10.1090/S0002-9939-08-09566-X, provides the free decomposition
on p.~871, Proposition~1 on p.~872, and Proposition~2 with its
complete proof on p.~873.
M.~Schwagenscheidt, \emph{Nonvanishing modulo $\ell$ of Fourier
coefficients of Jacobi forms}, Research in Number Theory~2 (2016),
article~4, doi:10.1007/s40993-015-0035-1, supplies Theorem~1.2
(p.~2, full proof pp.~11--12) and the normalized Eisenstein and
bounded-denominator statements on p.~3.
D.~Choi, Y.~Choie and T.~Kikuta, \emph{Sturm type theorem for
Siegel modular forms of genus~2 modulo $p$}, Acta Arith.~158
(2013), 129--139, doi:10.4064/aa158-2-2, states the ordinary
elliptic Sturm theorem on p.~129, citing J.~Sturm,
\emph{On the congruence of modular forms}, Lecture Notes in
Mathematics~1240 (1987), 275--280.  The ordinary Sturm theorem is
imported through that exact published statement; no fresh reading
of Sturm's original proof is claimed.

% END evidence/main_gate/jacobi_two_unit_fields.tex


% BEGIN evidence/main_gate/jacobi_third_unit_counterexample.tex
\begin{proposition}[KT-JACOBI-THIRD-UNIT-NONIMPLICATION]
A field arising from a unit heat coefficient in
KT-JACOBI-TWO-UNIT-FIELD need not have the third required unit,
even in the automatic family (JF.13).
\end{proposition}
\begin{proof}
Take $p=31$, $e=16=(p+1)/2$ and $D=-79$.  This is an odd
fundamental discriminant prime to $p$.  The exact character sums give
\[
\begin{gathered}
 B_{1,\chi_D}=-5,\qquad
 \frac{B_{3,\chi_D}}3=992=31\cdot32,\\
 \frac{B_{17,\chi_D}}{17}
       =-23090968227437542810000061098240
       \equiv20\pmod{31}.
\end{gathered}
\]
The reproducible standard-library certificate
\path{jacobi_third_unit_counterexample.py} evaluates all three
numbers by the Bernoulli-polynomial formula and independently
reconciles the last residue using the defining conductor-$79\cdot31$
sum modulo $31^2$.  In particular $h(-79)=5$.
At $(n,r)=(20,1)$ the heat coefficient in (JF.11) is
$79(-20)\equiv1\pmod{31}$.  Also $K=50$, so $n\leq K/2$
and $|D|=79\leq3p+6=99$: this is inside the proved selection bound.
But $2e\equiv2\pmod{30}$, so (FV.2) identifies the third unit
with $B_{3,\chi_D}/3$, which is zero modulo $31$.
This disproves the automatic implication for an individually selected
field, not the existence of another field with all three units.
It is a scope discriminator and not a new rational vanishing result.
\end{proof}

% END evidence/main_gate/jacobi_third_unit_counterexample.tex


% BEGIN evidence/main_gate/jacobi_fixed_weight_reformulation.tex
\subsection{A fixed-weight formulation of the simultaneous field problem}

\begin{proposition}[KT-JACOBI-FIXED-WEIGHT-REFORMULATION]
Let $p\ge31$ be prime with $p\equiv3\pmod4$, put
$b=(p+3)/2$, and let $D<0$ be an odd fundamental discriminant
prime to $p$.  Set
\[
 d=|D|,\qquad F=\mathbf Q(\sqrt{pd}),\qquad
 \chi_F=\chi_D\chi_{-p}.
\]
The last equality is an equality of primitive quadratic characters,
and $pd$ is the fundamental discriminant of $F$.  Both sides of
\[
 \frac{B_{b,\chi_D}}b\equiv\frac{B_{2,\chi_F}}2\pmod p
                                                        \tag{JFW.1}
\]
are $p$-integral.  Thus, for every
$e=(p+1)/2+(p-1)t$, $t\ge0$, the three auxiliary units are
equivalent to the fixed-weight conditions
\[
 p\nmid h(D),\qquad
 \frac{B_{2,\chi_F}}2\in\mathbf Z_p^\times,\qquad
 \frac{B_{3,\chi_D}}3\in\mathbf Z_p^\times.              \tag{JFW.2}
\]
If $d\le3p+6$, the first condition follows from
KT-SMALL-DISCRIMINANT-CLASS-BOUND.  This identifies the two
remaining special values at the same $D$; it does not assert their
simultaneous nonvanishing.
\end{proposition}

\begin{proof}
Since $D$ and $-p$ are coprime odd fundamental discriminants,
their product $pd$ is a positive fundamental discriminant.  Euler's
criterion gives, for every $a$ prime to $pd$,
\[
 \chi_D(a)a^{b-1}
 =\chi_D(a)a^{(p-1)/2+1}
 \equiv\chi_F(a)a\pmod p.                              \tag{JFW.3}
\]
Choose, by the Chinese remainder theorem, an integer $c>1$ with
$c\equiv1\pmod p$, $(c,d)=1$, and $\chi_D(c)=-1$.
Then $\chi_F(c)=-1$ as well.  Apply Guerzhoy's Proposition~2
with $\Delta=d$, $(\xi_1,r_1)=(\chi_D,b)$ and
$(\xi_2,r_2)=(\chi_F,2)$.  Its hypotheses allow the second
character to have conductor divisible by $p$.
The congruence of the integrands is exactly (JFW.3).  The two
regularizing factors are
\[
 1-\chi_D(c)c^b\equiv
 1-\chi_F(c)c^2\equiv2\pmod p,
\]
and hence are units.  The first Euler factor is
$1-\chi_D(p)p^{b-1}\equiv1\pmod p$; the second is exactly
$1$, since $\chi_F(p)=0$.

Proposition~1 expresses each regularized value as an integral
against an integral measure.  Thus both unregularized values
are integral, and Proposition~2 permits division by the common
unit residue $2$.  Using $L(1-r,\xi)=-B_{r,\xi}/r$ proves
(JFW.1).  In particular no prime-to-$p$ conductor assumption has
been inserted for $\chi_F$.

By KT-FIELD-CHARACTER-REDUCTION, the first twisted unit is
equivalent to $B_{b,\chi_D}/b$ being a unit, while
$2e\equiv2\pmod{p-1}$ identifies the second twisted unit
with $B_{3,\chi_D}/3$.  The untwisted $B_1$ unit is
equivalent to $p\nmid h(D)$; the factors from the roots of
unity are units because $p\ge31$.  This proves (JFW.2) in
every indicated weight.  Finally $3p+6<4p$ for $p\ge31$,
so the cited class-number bound applies.
\end{proof}

\paragraph{Exact primary input and scope.}
P.~Guerzhoy, \emph{Quadratic congruences for Cohen--Eisenstein
series}, Glasgow Math.~J.~41 (1999), 141--144,
\url{https://doi.org/10.1017/S001708959997060X}:
Proposition~1 and formula~(2), printed p.~143, give the integral
Mazur measure; Proposition~2, printed pp.~143--144,
and formula~(4) give the change-of-character Kummer congruence.
The short deduction of Proposition~2 from the measure is included
there.  The construction of that measure is an imported published
input, not a newly checked proof here.  Only these propositions
are used, not the final type-a/type-b congruence argument.
The local PDF and retrieval metadata are
\path{literature/jacobi_simultaneous/Guerzhoy_1999.pdf}
and its corresponding retrieval JSON.
This proposition is a reformulation of the simultaneous target,
not an additional vanishing result for the original rational class
component or an exclusion of genuine rank-three sources.

% END evidence/main_gate/jacobi_fixed_weight_reformulation.tex


% BEGIN evidence/main_gate/kohnen_theta_discriminator.tex
\begin{proposition}[KT-KOHNEN-THETA-NONCONGRUENCE]
The normalized Kohnen plus projection of $\theta^{p+4}$ is not,
in general for primes $p\equiv3\pmod4$, a scalar multiple modulo $p$ of the Cohen Eisenstein
series $H_{(p+3)/2}$.  In particular, this fails at $p=31$.
Consequently the Frobenius identity
$\theta^{p+4}\equiv\theta(q^p)\theta^4\pmod p$ alone does
not supply the missing simultaneous unit in the automatic Jacobi
family, even after the genuine plus projection.
\end{proposition}

\begin{proof}
Write $\theta(\tau)=\sum_{a\in\mathbf Z}q^{a^2}$,
$\theta_o(\tau)=\sum_{a\ \mathrm{odd}}q^{a^2}$ and
$\theta^m=\sum_{n\ge0}r_m(n)q^n$.
For an odd integer $b\ge1$ put $m=2b+1$ and
$s=(-1)^{(b+1)/2}$.  We use the normalization
\[
 \mathcal P f=\frac12f+
 \frac{s}{2\sqrt2}\sum_{\nu=0}^3 f|_{m/2}\xi\widetilde A_\nu,
 \quad
 \xi=\left(\begin{pmatrix}4&1\\0&4\end{pmatrix},e^{\pi i/4}\right),
 \quad A_\nu=\begin{pmatrix}1&0\\4\nu&1\end{pmatrix}.
                                                        \tag{KTHT.1}
\]
Here $(f|_k(A,\phi))(\tau)=\phi(\tau)^{-2k}f(A\tau)$,
and $\widetilde A_\nu$ uses the usual theta automorphy factor.
This is the normalized plus operator in Matsusaka's thesis,
printed pp.~27--28.  It is $3/2$ times the idempotent
normalization; this nonzero rational scale has no effect on the
question of proportionality modulo $31$.

The four transformed theta functions in (KTHT.1), before taking
their $m$th powers, are
\[
 e^{-\pi i/4}\theta(\tau+1/4),\quad
 \sqrt2\theta_o(\tau),\quad
 e^{\pi i/4}\theta(\tau-1/4),\quad
 \sqrt2\theta(4\tau),                                  \tag{KTHT.2}
\]
in the order $\nu=0,1,2,3$.
Here are details controlling the phases.  For $\nu=2$, factor
\[
 \begin{pmatrix}3&1/4\\8&1\end{pmatrix}
 =\begin{pmatrix}3&1\\8&3\end{pmatrix}
  \begin{pmatrix}1&-1/4\\0&1\end{pmatrix};
\]
the theta multiplier of the integral matrix is
$(8/3)\epsilon_3^{-1}=i$.
For $A=\left(\begin{smallmatrix}1&0\\2&1\end{smallmatrix}\right)$,
Poisson summation gives
\[
 (2\tau+1)^{-1/2}\theta(A\tau)
     =\sum_{a\ \mathrm{odd}}q^{a^2/4}.
\]
Factoring the matrix for $\nu=1$ as
$A\left(\begin{smallmatrix}2&1/4\\0&1/2\end{smallmatrix}\right)$
and using $a^2\equiv1\pmod8$ gives its entry in (KTHT.2).
For $\nu=3$, factor it as
\[
 \begin{pmatrix}4&1/4\\12&1\end{pmatrix}
 =\begin{pmatrix}3&1\\8&3\end{pmatrix}
  \begin{pmatrix}0&-1/4\\4&1\end{pmatrix}
\]
and use $\theta(-1/(4z))=\sqrt{-2iz}\,\theta(z)$.
The square roots have their holomorphic branches fixed at
$\tau=iy$, $y>0$, so these identities extend to the upper
half-plane without a sign ambiguity.

For $n\equiv3\pmod4$, the sum of the first and third
terms of (KTHT.2), raised to the $m$th power, has coefficient
$s\sqrt2\,r_m(n)$.  The last term is supported at multiples
of $4$, and the second starts at $q^m$.  Therefore
\[
 [q^n]\mathcal P\theta^m=r_m(n)
 \quad(0<n<m,\ n\equiv3\pmod4).                        \tag{KTHT.3}
\]
For $p=31$ and $m=35$, Frobenius gives
$r_{35}(3)\equiv r_4(3)=32\equiv1$ and
$r_{35}(7)\equiv r_4(7)=64\equiv2\pmod{31}$.
On the other hand, the defining Bernoulli-polynomial sums give
\[
 H_{17}(3)=-\frac{B_{17,\chi_{-3}}}{17}\equiv26,
 \qquad
 H_{17}(7)=-\frac{B_{17,\chi_{-7}}}{17}\equiv16
 \pmod{31}.
\]
Hence the two coefficient vectors have determinant
$1\cdot16-2\cdot26\equiv26\ne0\pmod{31}$.
They cannot be scalar multiples.  The standalone certificate
\path{kohnen_theta_discriminator.py} computes both exact
Bernoulli values and independently checks them using conductor
$31d$ character sums modulo $31^2$, for $d=3,7$; it also
computes the theta coefficients by exact integer convolution.
This disproves the proposed uniform scalar identity, not the
existence of another modular or genus-average construction.
\end{proof}

\paragraph{Source and normalization audit.}
T.~Matsusaka, \emph{Fourier coefficients of polyharmonic weak
Maass forms}, Ph.D.~thesis, Kyushu University (2019), printed
pp.~27--28, gives (KTHT.1), its slash convention, and the
phase $e^{\pi i/4}$; Remark~2.3.2 explicitly warns against
inserting an additional factor $2\lambda+1$ into that exponent
under this slash convention.  The institutional PDF is
\url{https://catalog.lib.kyushu-u.ac.jp/opac_download_md/2236044/math0246.pdf}.
The operator's plus-space property is an imported statement;
the theta transformations and the coefficient calculation above
are proved here.  No whole-thesis review is claimed.
The author PDF of Bruinier--Jenkins--Ono,
\emph{Hilbert class polynomials and traces of singular moduli},
Math.~Ann.~334 (2006), 373--393, \S3.2, formula~(3.12),
printed pp.~14--15, uses the idempotent factors $1/3$ instead of
$1/2$ but prints that extra phase factor.  We do not silently
use that inconsistent phase with its stated slash convention.

% END evidence/main_gate/kohnen_theta_discriminator.tex


% BEGIN evidence/main_gate/same_discriminant_detectors.tex
\subsection{An exact quadratic detector for the same-discriminant product}

The result in this subsection is a reformulation with an explicit integral
normalization.  It does not assert that the detector has a unit value,
and it does not supply the remaining auxiliary field.

\begin{proposition}[KT-EISENSTEIN-QUADRATIC-DETECTOR]
Let $p\geq31$ be prime with $p\equiv3\pmod4$, and put
\[
 b=(p+3)/2,\qquad k=b+3=(p+9)/2.
\]
Let $D<0$ be an odd fundamental discriminant, $d=|D|$, $p\nmid d$,
and write $\beta_j(\chi_D)=B_{j,\chi_D}/j$.  Normalize the level-one
Eisenstein series by
\[
 G_k(z)=-\frac{B_k}{2k}+\sum_{n\geq1}\sigma_{k-1}(n)e^{2\pi inz}.
\]
For the standard Gauss sum $\tau(\chi_D)$ one has the exact identity
\begin{equation}
 \frac{\tau(\chi_D)L(G_k\otimes\chi_D,3)}{(2\pi i)^3}
       =\frac{\beta_3(\chi_D)\beta_b(\chi_D)}{4d^2}.
 \tag{QD.1}
\end{equation}
Both Bernoulli quotients on the right are $p$-integral.

Choose a prime $q\equiv1\pmod p$, and set
\[
 E_{k,q}(z)=G_k(z)-q^{k-1}G_k(qz).
\]
This is a normalized Eisenstein eigenform of level $q$ whose constant
terms at every cusp vanish modulo $p$.  Consequently it has the
Heumann--Vatsal modular symbol $M_{k,q}$ modulo $p$.  For every $D$ as
above with $q\nmid d$ and $\chi_D(q)=-1$, the following coefficient is
\begin{equation}
 [X^2Y^{k-4}](M_{k,q}|R_{\chi_D})(\{\infty\}-\{0\})
       =-\frac{15}{8}\,\beta_3(\chi_D)\beta_b(\chi_D)
          \quad\hbox{in }\mathbb F_p.
 \tag{QD.2}
\end{equation}
In particular, this specified coefficient is nonzero if and only if both
required Bernoulli quotients are units.
\end{proposition}

\begin{proof}
For $\Re s>k$, the divisor-sum expansion gives
\[
 L(G_k\otimes\chi_D,s)=L(s,\chi_D)L(s-k+1,\chi_D),
\]
and continuation gives the identity at $s=3$.  The ordinary Bernoulli
Fourier series, followed by the primitive Gauss-sum identity, gives
\[
 \begin{split}
 B_{3,\chi_D}
 &=d^2\sum_{a=1}^{d}\chi_D(a)B_3(a/d)\\
 &=-\frac{6d^2}{(2\pi i)^3}
       \sum_{n\ne0}\frac{1}{n^3}
           \sum_{a=1}^{d}\chi_D(a)e^{2\pi ina/d}\\
 &=-\frac{12d^2\tau(\chi_D)}{(2\pi i)^3}L(3,\chi_D).
 \end{split}
\]
The Fourier series is absolutely convergent, and the last equality uses
the odd parity of $\chi_D$.  Since $L(1-b,\chi_D)=-\beta_b(\chi_D)$,
this proves (QD.1), including its sign and conductor power.
The Bernoulli-polynomial expression for $B_{j,\chi_D}$ shows its
$p$-integrality for $j=3,b$: here $j\leq p-2$, all ordinary Bernoulli
numbers of indices at most $j$ are $p$-integral, $p\nmid d$, and
$p\nmid j$.

Put $c=-B_k/(2k)$.  The inequalities $k<p-1$ and $p>k$ hold in the
stated range, so $c$ is $p$-integral.  The constant of $E_{k,q}$ at
infinity is $c(1-q^{k-1})$.  With the determinant-normalized slash
operator and $W_q=\left(\begin{smallmatrix}0&-1\\q&0\end{smallmatrix}\right)$,
\[
 G_k|_kW_q=q^{k/2}G_k(qz),\qquad
 G_k(qz)|_kW_q=q^{-k/2}G_k(z).
\]
Thus its constant at zero is $c q^{k/2-1}(q-1)$.  Both constants
vanish modulo $p$.  These are the two $\Gamma_0(q)$ cusp classes;
refining to $\Gamma_1(q)$ changes only the equivalent cusp
normalizations, with $p$-unit width factors.  The positive Fourier
coefficients are those of the normalized Eisenstein series attached to
the trivial character and the principal character modulo $q$.
In particular $a_1=1$, $U_qE_{k,q}=E_{k,q}$, and the good-prime
eigenvalues are $1+r^{k-1}$.

Theorem 3.23 of Heumann--Vatsal therefore supplies the integral modular
symbol modulo $p$.  Its construction is used here as a published input.
Twisting the displayed stabilization gives
\[
 \frac{\tau(\chi_D)L(E_{k,q}\otimes\chi_D,3)}{(2\pi i)^3}
  =(1-\chi_D(q)q^{k-4})
        \frac{\beta_3(\chi_D)\beta_b(\chi_D)}{4d^2}.
\]
Their Corollary 4.7 says that the coefficient in (QD.2), before
reduction, is the preceding value times
$-\binom{k-2}{2}2!d^2=-(k-2)(k-3)d^2$.
Now $q\equiv1$, $\chi_D(q)=-1$, $k-2\equiv5/2$, and
$k-3\equiv3/2$ modulo $p$.  The resulting scalar is $-15/8$.
It is a unit because $p\geq31$, proving the final equivalence.
\end{proof}

\paragraph{Exact source scope.}
J.~Heumann and V.~Vatsal, \emph{Modular symbols, Eisenstein series,
and congruences}, J. Th\'eor. Nombres Bordeaux \textbf{26} (2014),
709--756, DOI 10.5802/jtnb.886: the definition and hypotheses are on
pp.~734--735, Theorem 3.23 and its proof on pp.~736--737, and
Theorem 4.6 with its complete proof and Corollary 4.7 are on
pp.~742--745.  The primary PDF is
\url{https://www.numdam.org/item/JTNB_2014__26_3_709_0.pdf}.
Theorem 5.1 and its full proof on pp.~747--748 show that this modular
symbol is nonzero.  Their proof obtains a unit using characters with
growing auxiliary odd-prime-power order, and chooses the $j=0$
coefficient.  It does not imply (QD.2) is nonzero for any quadratic
character.  Formula (QD.1) was derived directly above; no odd-$j$
constant from their Proposition 6.1 is needed.

\paragraph{What a quadratic detector would still have to prove.}
The necessary nonzero evaluation is the particular coefficient (QD.2),
with an odd negative fundamental $D$, $p q\nmid D$, and
$\chi_D(q)=-1$.  A merely nonzero value in a different polynomial
component does not imply this conclusion.  The third auxiliary
condition $p\nmid h(D)$ must also be supplied.  It follows from the
already proved small-discriminant class bound if $d<4p$; it is not
automatic for an unbounded discriminant.  Thus neither full-symbol
primitivity nor unrestricted complex quadratic-twist nonvanishing is
the missing assertion.

% END evidence/main_gate/same_discriminant_detectors.tex


% BEGIN evidence/main_gate/senary_trace_detector.tex
\subsection{A genus-six trace and the remaining integral detection problem}

The following gives an exact same-discriminant product.  Its trace is
integral, but a unit value is a further assertion.  In particular, the
result does not provide an auxiliary field or close the arithmetic gate.

\begin{proposition}[KT-SENARY-GENUS-TRACE-DETECTOR]
Let $p\geq31$ be prime, $p\equiv3\pmod4$, and put
$b=(p+3)/2$ and $k=b+3$.  Write $E_k^{(6)}$ for the degree-six,
level-one Siegel Eisenstein series with constant coefficient $1$, and set
\begin{equation}
 F_p=\kappa_k E_k^{(6)},\qquad
 \kappa_k=\frac{\zeta(1-k)}8
                  \prod_{j=1}^3\zeta(1-2k+2j).
 \tag{SG.1}
\end{equation}
Then $\kappa_k/B_k\in\mathbb Z_{(p)}^\times$.

Let $D<0$ be odd and fundamental, put $d=|D|$, suppose $p\nmid d$,
and let $t$ be the number of distinct prime divisors of $d$.
Let $\mathcal L_D$ be the genus of positive definite even lattices of
rank six whose local lattices are
$H\oplus H\oplus\mathcal O_{\mathbb Q(\sqrt D)}(-1)$.
This genus is nonempty and its lattices have determinant $d$.
For a lattice $L$ in this genus, let $T_L$ be half its Gram matrix.
Then
\begin{equation}
 a(F_p,T_L)=-\beta_b(\chi_D),\qquad
 \mathcal T_D(F_p):=
  \sum_{[L]\in\mathcal L_D/\cong}
             \frac{a(F_p,T_L)}{|\operatorname O(L)|}
 =-\frac{\beta_b(\chi_D)\beta_3(\chi_D)}
                    {2^{t-1}\,23040}.
 \tag{SG.2}
\end{equation}
Every summand in this finite trace is $p$-integral.  Its reduction is
nonzero if and only if both Bernoulli quotients are units.  If also
$d<4p$, the already proved small-discriminant bound supplies
$p\nmid h(D)$, so a nonzero trace in that range supplies all three
units for the automatic progression under consideration.
\end{proposition}

\begin{proof}
We use two explicitly located primary inputs.  First, Kohnen's
coefficient formula is Theorem 21 on p.~22 of his author-hosted
\emph{A short course on Siegel modular forms}.  The definitions of
$D_T$, $f_T$, and $\varphi(a;T)$ are on pp.~20--21.
For degree $2n$ it states
\[
 a(E_k^{(2n)},T)=
 \frac{2^n}{\zeta(1-k)\prod_{j=1}^n\zeta(1-2k+2j)}
 \sum_{a\mid f_T}a^{k-n-1}\varphi(a;T)
 H\bigl(k-n,|D_{T,0}|(f_T/a)^2\bigr).
\]
Here $(-1)^n\det(2T)=D_{T,0}f_T^2$.  In our case $n=3$,
$D_{T,0}=D$ and $f_T=1$.  The defining sum for $\varphi(1;T)$
has the single unimodular class and gives $\varphi(1;T)=1$.
Thus (SG.1) gives
$a(F_p,T_L)=H(b,d)=L(1-b,\chi_D)=-\beta_b(\chi_D)$.
The coefficient theorem is imported; its earlier lifting proof is not
being re-proved here.

Second, the genus and its local specification are given by
Ibukiyama--Williams, \emph{Hermitian modular forms and algebraic
modular forms on $\operatorname{SO}(6)$}, Section 5, author PDF
pp.~23--26.  We give the numerical mass calculation explicitly, to
fix the normalization as the sum of reciprocals of the \emph{full}
orthogonal automorphism groups.  Their Minkowski--Siegel formula is
\[
 \operatorname{mass}(\mathcal L_D)
   =\frac{3d^{7/2}}{2\pi^9}\prod_q\alpha_q^{-1}.
\]
For odd $q\nmid d$ the local factor is
$(1-q^{-2})(1-q^{-4})(1-\chi_D(q)q^{-3})$; for $q\mid d$ it is
$2q(1-q^{-2})(1-q^{-4})$.  At $2$ the displayed local quantities
$P=(1-2^{-2})(1-2^{-4})(1-2^{-6})$ and
$E=2/(1+\chi_D(2)2^{-3})$ give
\[
 \alpha_2=2^5PE
 =2^6(1-2^{-2})(1-2^{-4})(1-\chi_D(2)2^{-3}).
\]
Consequently
\[
 \prod_q\alpha_q^{-1}
   =\frac{\zeta(2)\zeta(4)L(3,\chi_D)}{2^{t+6}d},
 \qquad
 \operatorname{mass}(\mathcal L_D)
   =\frac{3d^{5/2}\zeta(2)\zeta(4)L(3,\chi_D)}
                         {2^{t+7}\pi^9}.
\]
The primitive odd quadratic functional equation, equivalently the
Bernoulli Fourier calculation used in (QD.1), yields
$d^{5/2}L(3,\chi_D)/\pi^3=2B_{3,\chi_D}/3$.
Substituting $\zeta(2)=\pi^2/6$ and
$\zeta(4)=\pi^4/90$ therefore gives
\begin{equation}
 \operatorname{mass}(\mathcal L_D)
   =\frac{B_{3,\chi_D}}{2^{t-1}2^9 3^3 5}
   =\frac{\beta_3(\chi_D)}{2^{t-1}23040}.
 \tag{SG.3}
\end{equation}
Combining this with the constant value of the Fourier coefficient
proves (SG.2).

We check the integral assertions separately.  A finite subgroup of
$\operatorname{GL}_6(\mathbb Z)$ has order prime to $p\geq31$:
otherwise Cauchy's theorem gives an element of order $p$, whose
nontrivial rational representation has dimension at least $p-1>6$.
Thus each automorphism denominator is a $p$-unit.  Moreover
$3,b\leq p-2$ and $p\nmid d$, so the Bernoulli-polynomial formula
shows directly that both $\beta_3(\chi_D)$ and $\beta_b(\chi_D)$
are $p$-integral.  Since $23040$ and $2$ are $p$-units, the trace
criterion follows.

Finally, the three indices $2k-2j$ are $p+7,p+5,p+3$.
Kummer congruences give
\[
 \frac{B_{p+7}}{p+7}\equiv\frac{B_8}{8}=-\frac1{240},\quad
 \frac{B_{p+5}}{p+5}\equiv\frac{B_6}{6}=\frac1{252},\quad
 \frac{B_{p+3}}{p+3}\equiv\frac{B_4}{4}=-\frac1{120}
 \pmod p.
\]
All are units in the stated range, and $k<p-1$.  Hence
$\kappa_k/B_k=-(8k)^{-1}\prod_{j=1}^3\zeta(1-2k+2j)$ is a
$p$-unit.  The assertion about $h(D)$ is exactly
KT-SMALL-DISCRIMINANT-CLASS-BOUND.
\end{proof}

\paragraph{A real cancellation in this trace.}
The previously certified example
KT-JACOBI-THIRD-UNIT-NONIMPLICATION has $p=31$, $b=17$, $D=-79$,
$\beta_{17}(\chi_{-79})\equiv20\pmod{31}$ and
$\beta_3(\chi_{-79})=992$.  Formula (SG.3) gives
$\operatorname{mass}(\mathcal L_{-79})=31/720$.
Thus every individual $a(F_{31},T_L)$ is a unit, while the integral
genus trace is zero modulo $31$.  Nonvanishing of an individual
full-rank Siegel coefficient therefore does not force this trace to
be a unit, even for the actual Eisenstein series used here.

\begin{lemma}[KT-ANALYTIC-RESIDUE-NONSATURATION]
An analytically computed unit residue of a Dirichlet series need not
detect a unit among its $p$-integral coefficients, even when every
index is prime to $p$.
\end{lemma}
\begin{proof}
The absolutely convergent series on $\Re s>1$
\[
 A_p(s)=p\sum_{p\nmid n}n^{-s}=p(1-p^{-s})\zeta(s)
\]
has every coefficient in $p\mathbb Z$, but its meromorphic residue
at $s=1$ is $p-1$, a $p$-unit.  A residue argument for the trace above
must therefore establish an integral coefficient functional in that
specific modular setting; complex continuation and the residue value
alone do not establish it.
\end{proof}

\paragraph{Primary locators and normalization caution.}
Kohnen's primary author notes are at
\url{https://www.mathi.uni-heidelberg.de/~winfried/siegel2.pdf},
pp.~20--22, especially Theorem 21.  Ibukiyama--Williams is published
in J. Number Theory \textbf{286} (2026), 230--279,
DOI 10.1016/j.jnt.2026.02.009; the checked author PDF is
\url{https://btw-47.github.io/hermitian_orthogonal.pdf}.
Its Proposition 12 statement on p.~24 prints $3^2$ in the mass
denominator, while the proof conclusions on pp.~25--26 have $3^3$.
The explicit calculation above confirms $3^3$; for $D=-3$ it gives
the $E_6$ mass $1/103680$.  We do not use the numerical prefactor of
the 1996 Ibukiyama--Katsurada Koecher--Maa\ss{} announcement, or
assert that an unexamined analytic residue supplies an integral unit
functional.  No determinant bound or unit trace is proved here.

% END evidence/main_gate/senary_trace_detector.tex


% BEGIN evidence/main_gate/witt_pullback_discriminator.tex
\subsection{What an integral Witt coefficient does not detect}

The following is a failure gate for a proposed same-discriminant selection
argument.  It neither proves nor disproves existence of a different
auxiliary field with all the required units.

\begin{proposition}[KT-WITT-PULLBACK-MASS-OBSTRUCTION]
Let $F_{31}$ be the normalized degree-six Eisenstein series of
KT-SENARY-GENUS-TRACE-DETECTOR, so $b=17$ and $k=20$.  Put
\[
 C=\begin{pmatrix}2&-1&0\\-1&2&-1\\0&-1&2\end{pmatrix},
 \qquad
 B=\begin{pmatrix}2&-1&0\\-1&2&-1\\0&-1&14\end{pmatrix}.
\]
For the Witt restriction
$W(F)(Z_1,Z_2)=F(\operatorname{diag}(Z_1,Z_2))$, define the
\emph{determinant-$79$ contribution} to its $(C/2,B/2)$ coefficient by
\[
 c_{79}=
 \sum_{\substack{R\in M_3(\mathbb Z)\\
 G_R=\left(\begin{smallmatrix}C&R\\R^t&B\end{smallmatrix}\right)>0\\
 \det G_R=79}}
 a(F_{31},G_R/2).
 \tag{WP.1}
\]
There are exactly $56$ summands, and
\[
 c_{79}\equiv27\pmod{31},\qquad
 \mathcal T_{-79}(F_{31})\equiv0\pmod{31}.
 \tag{WP.2}
\]
Thus even a unit determinant-filtered contribution to an ordinary Witt
coefficient does not force the product
$\beta_{17}(\chi_{-79})\beta_3(\chi_{-79})$ to be a unit.
The determinant filter in (WP.1) is just a finite coefficient selection;
no modularity of this filtering operation is asserted.
\end{proposition}

\begin{proof}
The coefficient formula for a Witt restriction follows immediately by
substituting a block diagonal argument in the Fourier expansion.  It sums
over the cross-block matrices $R$, including positive semidefinite matrices
in the unfiltered coefficient.  For the strictly positive matrices in
(WP.1), the two-by-two principal minors give
$|R_{ij}|^2<C_{ii}B_{jj}$.  The first two columns therefore have entries
in $\{-1,0,1\}$ and the third has entries in $\{-5,\ldots,5\}$.
This finite box has $3^6 11^3=970299$ elements.

Here $\det C=4$, $\det B=40$, and
\[
 M=4C^{-1}=\begin{pmatrix}3&2&1\\2&4&2\\1&2&3\end{pmatrix}.
\]
Writing the columns of $R$ as $x,y,v$, positivity requires
$x^tMx<8$, $y^tMy<8$, and $v^tMv<56$.  These leave respectively
$15,15,411$ vectors.  The scaled Schur complement is
\[
 S=4B-R^tMR=\begin{pmatrix}a&z&s\\z&b&t\\s&t&c\end{pmatrix},
 \qquad \det G_R=\det(S)/16.
\]
Sylvester's criterion says that $G_R>0$ exactly when
$a>0$, $ab-z^2>0$, and $\det S>0$.  The saved exact certificate enumerates
these vectors and checks
\[
 (ab-z^2)c-at^2-bs^2+2zst=16\cdot79.
\]
It obtains $56$ solutions.  A separate enumeration of the full coordinate
box uses integer Bareiss determinants of the four-, five-, and six-dimensional
principal matrices and obtains the identical list.  Thus both the finiteness
and completeness of the count are explicit.

Since $-79$ is fundamental, the coefficient formula already used in
KT-SENARY-GENUS-TRACE-DETECTOR gives
$a(F_{31},G_R/2)=-\beta_{17}(\chi_{-79})$ for every summand, independently
of its integral isometry class.  The previously proved exact values are
$\beta_{17}(\chi_{-79})\equiv20\pmod{31}$ and
$\beta_3(\chi_{-79})=992$; the new certificate also recomputes both
from the rational Bernoulli-polynomial formula.  Hence
$c_{79}\equiv-56\cdot20=27\pmod{31}$.
The genus trace is zero because its mass is $31/720$, by (SG.3).
\end{proof}

\begin{lemma}[KT-GENUS-PROJECTOR-INTEGRALITY]
Let $p\geq31$ and let $\mathcal L_D$ be a senary genus as in
KT-SENARY-GENUS-TRACE-DETECTOR.  On
$\Lambda=\bigoplus_{[L]}\mathbb Z_{(p)}\delta_L$, consider the canonical
mass-orthogonal projector onto constant functions:
\[
 P(f)=\frac{\sum_{[L]}f(L)/|\operatorname O(L)|}
 {\sum_{[L]}1/|\operatorname O(L)|}\,\mathbf1.
 \tag{WP.3}
\]
Then $P(\Lambda)\subseteq\Lambda$ if and only if
$\beta_3(\chi_D)$ is a $p$-unit.
\end{lemma}
\begin{proof}
All automorphism denominators are $p$-units, as proved in the cited
senary result.  Therefore the mass $M$ is $p$-integral and every numerator
$\langle\delta_L,\mathbf1\rangle=1/|\operatorname O(L)|$ is a unit.
Formula (WP.3) preserves the lattice exactly when $M$ is a unit.
Equation (SG.3) identifies this condition with the asserted Bernoulli unit.
For $p=31,D=-79$, its value on every $\delta_L$ has valuation $-1$.
This concerns the canonical mass-orthogonal projector; arbitrary
nonorthogonal projections onto the constant line do exist integrally.
\end{proof}

\paragraph{The remaining pullback requirement.}
An unnormalized integral coefficient functional could still prove a useful
simultaneous-unit result.  It must actually equal an integral combination
of the genus traces, and its unit value must be established.  Ordinary
cross-block summation weights lattices by their numbers of compatible
ordered bases; (WP.2) shows that these weights cannot simply be replaced
by the mass weights.  Normalizing a genus average by its total mass, as in
(WP.3), already requires the third unit.  Neither observation is counted
as movement on the remaining arithmetic gate.

\paragraph{Published integrality scope.}
Pitale--Saha--Schmidt, \emph{Integrality and cuspidality of pullbacks of
nearly holomorphic Siegel Eisenstein series}, Publ. Mat. \textbf{66}
(2022), 405--434, DOI 10.5565/PUBLMAT6612216, Theorem 1.1(ii) on p.~408
and Remark 1.2 on p.~409, give a general normalized coefficient-integrality
theorem with hypothesis $p\geq2k$.  It cannot be invoked directly at
$k=(p+9)/2$.  This is an applicability statement about that theorem,
not a claim that our holomorphic coefficients fail to be integral; the
coefficients in (WP.1) are integral by their explicit formula.  The
theorem statement and explanatory remark were read for this scope check;
no fresh review of its full proof is claimed.

% END evidence/main_gate/witt_pullback_discriminator.tex


% BEGIN evidence/main_gate/neighbor_pairing_attempt.tex
\subsection{What integral neighbor operators do to the same-genus trace}

This is an audit of a proposed method for detecting the same-discriminant
product.  It supplies no auxiliary field and no rational Selmer vanishing.
It uses full orthogonal isometry classes, as in the senary mass calculation.

\begin{proposition}[KT-NEIGHBOR-CONSTANT-LINE-LIMIT]
Let $\mathcal G$ be a genus of positive-definite even lattices of rank six
and determinant $d$.  Fix a prime $p\geq31$ with $p\nmid2d$, and write
\[
 \Lambda=\mathbf Z_p^{\mathcal G/\cong},\quad
 \langle f,g\rangle=\sum_{[L]}\frac{f(L)g(L)}{|O(L)|},\quad
 M=\langle1,1\rangle.
 \tag{NP.1}
\]
For an odd prime $q\nmid d$, let $T_q$ be the unscaled Kneser neighbor
operator on functions.  Put
\[
 N_q=q^4+q^3+(1+\epsilon_q)q^2+q+1,
 \qquad \epsilon_q=\left(\frac{-d}{q}\right).
 \tag{NP.2}
\]
Then $T_q1=N_q1$.  Consequently every polynomial $P$ with
$\mathbf Z_p$ coefficients in finitely many such neighbor operators satisfies
\[
 \langle P((T_q))\,a1,1\rangle
       =aP((N_q))M\qquad(a\in\mathbf Z_p).
 \tag{NP.3}
\]
If $p\mid M$, these integral neighbor operations followed by genus
averaging never give a unit on a constant input.  In particular
$N_p\equiv1\pmod p$ does not imply $M\in\mathbf Z_p^\times$.

More strongly, if $p\mid M$, there is a simultaneous
characteristic-zero eigenline for all these $T_q$, orthogonal to
every spinor character, whose eigenvalues satisfy
\[
                  \lambda(T_q)\equiv N_q\pmod{\mathfrak p}
                         \quad\hbox{for every odd }q\nmid d
 \tag{NP.4}
\]
at a prime of a finite coefficient extension above $p$.
Its $T_p$ eigenvalue is a unit.
\end{proposition}

\begin{proof}
All automorphism orders in (NP.1) are units at $p$.  Otherwise Cauchy's
theorem would give an order-$p$ element of $GL_6(\mathbf Q)$, whose
nontrivial cyclotomic constituent has dimension $p-1>6$.
Thus the displayed pairing is perfect over $\mathbf Z_p$.

The neighbor construction identifies the neighbors of $L$ with the
isotropic lines in $L/qL$.  Their number is (NP.2), by the rank-six
finite-quadric formula.  This proves the constant row sum and (NP.3).
Reversing an isometry from $L'$ onto a neighbor of $L$ gives the balance
identity
\[
 \frac{\#\{q\text{-neighbors of }L\text{ isometric to }L'\}}{|O(L)|}
 =\frac{\#\{q\text{-neighbors of }L'\text{ isometric to }L\}}{|O(L')|}.
 \tag{NP.5}
\]
Hence the operators are self-adjoint.  Their commutativity is the usual
commutativity of the unramified orthogonal Hecke algebras at each prime,
together with the independence of different primes.

Let $S$ be the group of spinor genera in $\mathcal G$ and put $s=|S|$.
It is a finite elementary abelian two-group; every spinor genus has
mass $M/s$.  For each $u\in S$ choose one lattice $L_u$ in that spinor
genus and set
\[
 \Lambda^0_{\mathrm{sp}}
  =\left\{f\in\Lambda:\sum_{[L]\in u}\frac{f(L)}{|O(L)|}=0
                       \text{ for every }u\in S\right\}.
 \tag{NP.6}
\]
The map of component means is onto $\mathbf Z_p^S$, so this submodule
is saturated.  It is the orthogonal complement of all spinor characters.
A $q$-neighbor shifts the spinor genus by a fixed element $\delta_q$.
Together with (NP.5) and the row sum, this shows that $T_q$ preserves
(NP.6): the mean of $T_qf$ on $u$ is $N_q$ times the mean of $f$ on
$u\delta_q$.  If $p\mid M$, the element
\[
 g=1-\frac Ms\sum_{u\in S}|O(L_u)|\,\mathbf1_{L_u}
       \in\Lambda^0_{\mathrm{sp}}
       \quad\hbox{satisfies}\quad g\equiv1\pmod p.
 \tag{NP.7}
\]
The factor $s$ is a unit.  In particular $g$ is primitive, and its
reduction spans a simultaneous eigenline with eigenvalues $N_q$.

For completeness, let $\mathcal A$ be the image of the algebra generated
by all $T_q$ in $\operatorname{End}_{\mathbf Z_p}(\Lambda^0_{\mathrm{sp}})$.
It is finite and torsion-free over $\mathbf Z_p$, since it is a
submodule of a finite free module, and is commutative.  The line just
constructed defines a character $\mathcal A\to\mathbf F_p$.
A maximal ideal defining that character contains a minimal prime that
avoids $p$: the element $p$ is a nonzerodivisor in $\mathcal A$.
The corresponding finite torsion-free domain has fraction field finite
over $\mathbf Q_p$.  Normalizing and completing at a prime above the
chosen maximal ideal gives a characteristic-zero character reducing to
the specified one.  The operators and the component-mean conditions
are defined over $\mathbf Q$, and their weighted pairing over $\mathbf R$
is positive definite.  Thus the commuting operators are simultaneously
diagonalizable over $\mathbf R$ and hence over an algebraic coefficient
field; scalar extension gives the same conclusion over $\overline{\mathbf Q}_p$.
Therefore this character
occurs on an eigenline of $\Lambda^0_{\mathrm{sp}}$ after finite scalar
extension.  This proves (NP.4).  Finally (NP.2) gives $N_p\equiv1$,
which proves the unit assertion.
\end{proof}

\paragraph{Exact application to the retained genus.}
In KT-SENARY-GENUS-TRACE-DETECTOR, take $p=31$ and $D=-79$.
The already checked values are $M=31/720$ and
$a=-\beta_{17}(\chi_{-79})\equiv11\pmod{31}$.
Here $\epsilon_{31}=1$, since $-79\equiv13^2\pmod{31}$, and
\[
                  N_{31}=962\cdot993=955266\equiv1\pmod{31}.
\]
Thus every expression (NP.3) is zero modulo $31$, although the
constant eigenvalue of $T_{31}$ is a unit.  There is also an actual
simultaneous orthogonal Hecke congruence (NP.4) outside all spinor
characters.  No assertion that this eigenline belongs to the compact
unitary source, or to an admissible deformation space, is made.

\paragraph{Primary inputs and remaining theta-lift issue.}
The primary reference is G.~Chenevier, \emph{Statistics for Kneser
$p$-neighbors}, arXiv:2104.06846v2, 3 May 2021: (1.2)--(1.4) on
p.~3; the full local construction and Lemmas 4.1--4.3 on pp.~13--14;
Definition 4.5 and the self-transpose Hecke algebra on p.~15;
the spinor-genus description and complete proof of Lemma 5.6 on
p.~20; the spinor shift in Remark (R2) on p.~21; and (6.1)--(6.2),
Lemma 6.1 and the self-adjointness paragraph on pp.~22--23.
Those precise passages, including the short local and equal-mass
proofs, were read for this audit.  The mass and Bernoulli values are
imported from the already checked senary trace.

A different theta-lift functional can have varying weights on the
classes of a genus.  It is not covered by the constant row-sum argument
unless it is followed by this same genus average.  Its value on the
constant input is then a different weighted sum, whose integrality,
factorization by the specified same-discriminant product, and unit
value must each be proved.  Existence of an integral theta map alone
proves none of those unit assertions.  This proposition rejects the
neighbor-polynomial and spectral-separation shortcuts above; it does
not rule out a suitable further theta correspondence.

% END evidence/main_gate/neighbor_pairing_attempt.tex


% BEGIN evidence/lean_chain/quadratic_evaluation_limits.tex
\subsection{What finite quadratic evaluations can and cannot detect}

This subsection records a bounded attempt to detect the explicit
Eisenstein class in KT-EISENSTEIN-QUADRATIC-DETECTOR.  It proves two
limitations of a proposed spanning argument and a sufficient short
discriminant window.  It does not prove that this window contains the
required field, and it gives no additional vanishing of a rational
class-group component.

Fix a prime $p\geq31$ with $p\equiv3\pmod4$, and put
\[
 b=(p+3)/2,\qquad k=(p+9)/2,\qquad e=b-1.
\]
For the actual auxiliary-field problem the three units are
$B_{1,\chi_D}$, $\beta_b(\chi_D)$ and $\beta_3(\chi_D)$:
the last index is obtained from $2e+1=p+2$ by the previously
checked Kummer transport.  Here $\beta_r(\chi)=B_{r,\chi}/r$.
The discriminant must remain odd and negative and prime to $p$.
The carried small-discriminant bound gives $p\nmid h(D)$ whenever
$|D|<4p$.

\begin{proposition}[KT-SMALL-QUADRATIC-SPAN-LIMIT]
Let $q$ be prime with $q\equiv1\pmod p$.  Consider the selected
$j=2$ quadratic evaluation functionals at the odd negative fundamental
discriminants $D$ satisfying
\[
 0<|D|<4p,\qquad \gcd(D,pq)=1,\qquad \chi_D(q)=-1.
\]
On a full integral lattice in either parity component of the
characteristic-zero cuspidal modular-symbol space of weight $k$ and
level $\Gamma_0(q)$, these functionals cannot span the entire dual
lattice, either rationally or after reduction modulo $p$ when the
functionals are integral on that lattice.
\end{proposition}

\begin{proof}
There are at most $p$ candidate integers: an odd negative fundamental
discriminant is $D=-d$ with $d\equiv3\pmod4$, and there are $p$ such
positive integers below $4p$.  Squarefreeness and the displayed local
conditions only decrease their number.

Either cuspidal parity component has rank
$s=\dim S_k(\Gamma_0(q))$.  In the standard dimension formula put
$e_2=1+(\frac{-4}{q})$ and $e_3=1+(\frac{-3}{q})$, so that
$0\leq e_2,e_3\leq2$.  Since there are two cusps, the formula is
\[
 s=\frac{(k-1)(q+1)}{12}-1
   +e_2\left(\left\lfloor\frac{k}{4}\right\rfloor
                         -\frac{k-1}{4}\right)
   +e_3\left(\left\lfloor\frac{k}{3}\right\rfloor
                         -\frac{k-1}{3}\right).
\]
For even $k$, the two parenthetical terms are at least $-1/4$
and $-1/3$, respectively.  Also $q\geq2p+1$, because both primes
are odd.  Hence
\[
 s\geq\frac{(p+7)(p+1)}{12}-\frac{13}{6}>p.
\]
The last difference is $(p^2-4p-19)/12$, positive for $p\geq31$.
At most $p$ vectors cannot span a space of dimension greater than
$p$.  Reduction of a free lattice still has its full rank as
dimension over $\mathbf F_p$, so the same conclusion holds there.
\end{proof}

This proposition does not say that the explicit Eisenstein class is
annihilated by the small-discriminant evaluations.  Nor does it bound
an Eisenstein-localized eigenspace by $s$.  It rules out the proposed
full-lattice spanning shortcut, not a proof of detection of that
particular class.

\begin{lemma}[KT-EISENSTEIN-HECKE-EVALUATION-IDEAL]
Let $R$ be a ring, let $M$ be an $R$-module with Hecke action, and
let $v\in M$ satisfy $Tv=a(T)v$ for every operator being used.
For any family of $R$-linear functionals $\lambda_D$, the ideal
\[
 (\lambda_D(Tv):D,T)\subset R
\]
is exactly $(\lambda_D(v):D)$ if the identity operator is included.
\end{lemma}

\begin{proof}
One has $\lambda_D(Tv)=a(T)\lambda_D(v)$, giving one inclusion.
The identity operator gives the reverse inclusion.
\end{proof}

For the Heumann--Vatsal eigenclass and the chosen local condition,
the generators in this lemma are the unit multiple
$-\frac{15}{8}\beta_3(\chi_D)\beta_b(\chi_D)$.
Thus taking the Hecke closure of the evaluation rows cannot itself
improve their scalar ideal on this class.  This observation does
not rule out proving that the original ideal is the unit ideal.
Varying $q$ likewise changes only the already displayed Euler
factor.  For this pair of Bernoulli indices, Corollary 4.7 has
$j=2$ or $j=k-4$; an unspecified nonzero critical coefficient
can concern different Bernoulli factors.

\paragraph{The exact finite quadratic orthogonality calculation.}
Let $m$ be odd and squarefree, prime to $pq$, with $m<4p$.
Put $G=(\mathbf Z/m\mathbf Z)^\times$ and
$X=\operatorname{Hom}(G,\{\pm1\})$.  For $r=3,b$ define the
$\mathbf F_p$-valued function
\[
 F_r(a)=\frac{m^{r-1}}r B_r(a/m)\pmod p,\qquad a\in G.
\]
All denominators are units: $r<p-1$ and $p\nmid m$.
The Bernoulli-polynomial definition gives
\[
 \beta_{r,m}(\chi)=\sum_{a\in G}\chi(a)F_r(a).
\]
The subscript $m$ allows the character to be imprimitive.
Writing $H=G^2$, ordinary finite-group character orthogonality gives
\[
 \sum_{\substack{\chi\in X\\\chi(-1)=\chi(q)=-1}}\chi(g)
 =\frac{|X|}{4}
  \bigl(1_H(g)-1_H(-g)-1_H(qg)+1_H(-qg)\bigr).
 \tag{QE.1}
\]
Indeed multiply the unrestricted sum by
$(1-\chi(-1))(1-\chi(q))/4$ before summing; this proof also
covers dependent local conditions or an empty selected family.
Consequently the sum of the products
$\beta_{3,m}(\chi)\beta_{b,m}(\chi)$ over the selected characters
is exactly
\[
 \frac{|X|}{4}\sum_{a,c\in G}F_3(a)F_b(c)
  \bigl(1_H(ac)-1_H(-ac)-1_H(qac)+1_H(-qac)\bigr).
 \tag{QE.2}
\]
This is a concrete finite expression, with the same character in
both factors.  A nonzero value proves that one product is a unit.
If $\chi^*$ is the primitive character inducing $\chi$, then
\[
 \beta_{r,m}(\chi)=\beta_r(\chi^*)
  \prod_{\substack{\ell\mid m\\\ell\nmid\operatorname{cond}(\chi^*)}}
       (1-\chi^*(\ell)\ell^{r-1}).
 \tag{QE.3}
\]
Thus a nonzero product in (QE.2) supplies the required two primitive
units and a discriminant of size at most $m<4p$; its class-number
unit follows from the carried bound.  The parity and local
conditions also pass to $\chi^*$.

No nonzero value of (QE.2) has been proved uniformly here.
Quadratic orthogonality detects only the sums of $F_r$ on the
squareclasses $G/H$, not its individual values.  In particular
$\delta_a-\delta_{ah}$ is in its kernel for $h\in H$ with
$ah\ne a$.  Even nonzero quadratic Fourier transforms of two
functions may have disjoint supports.  Therefore separate
nonvanishing of the two transforms does not prove that their
pointwise product is nonzero.  These are exact limitations of an
argument using only character orthogonality, not counterexamples
to nonvanishing for the specific Bernoulli functions $F_r$.
The factors in (QE.3) can also vanish modulo $p$, so a primitive
unit need not be visible in a chosen imprimitive conductor.

\begin{proposition}[KT-BETA3-SHORT-WINDOW]
Let $D=-d$ be an odd negative fundamental discriminant.
If $d\ne3,7$, then $\beta_3(\chi_D)$ is a positive integer and
\[
 0<\beta_3(\chi_D)<\frac{5}{216}d^{5/2}.
 \tag{QE.4}
\]
For $d=3,7$ the exact values are $2/9$ and $16/7$.
In particular, for every prime $p\geq31$, the condition
\[
 p\nmid d,\qquad d^5\leq1600p^2
 \tag{QE.5}
\]
implies that $\beta_3(\chi_D)$ is a $p$-adic unit.
\end{proposition}

\begin{proof}
Write $\chi=\chi_D=(\frac{\cdot}{d})$ and
$S_j=\sum_{a=1}^d\chi(a)a^j$.  Odd parity, paired by $a\mapsto d-a$,
gives $S_2=dS_1$, and the explicit polynomial
$B_3(x)=x^3-\frac32x^2+\frac12x$ gives
\[
 \beta_3(\chi)=\frac{S_3-d^2S_1}{3d}.
 \tag{QE.6}
\]
For each prime $\ell\mid d$, the Chinese remainder theorem shows
$S_3\equiv0\pmod\ell$ if $d/\ell>1$: the sum of the nontrivial
quadratic character modulo $d/\ell$ is zero.  If $d=\ell$ is prime,
Euler's criterion and the finite-field power sum give the same
conclusion unless $\ell=3$ or $7$.  Indeed the relevant exponent
is $(\ell-1)/2+3$, which is not a multiple of $\ell-1$ for
$\ell>7$.  Thus $d\mid S_3$ outside the two stated cases.

If $3\nmid d$, then $a^3\equiv a\pmod3$ and $d^2\equiv1\pmod3$,
so the numerator in (QE.6) is also divisible by $3$.
If $3\mid d$ and $d>3$, put $d=3m$.  For $3\nmid a$,
$(\frac a3)a^3\equiv1\pmod9$.  Hence
\[
 S_3\equiv
 \sum_{\substack{1\leq a\leq3m\\3\nmid a}}
          \left(\frac a m\right)=0\pmod9.
\]
The last equality follows by subtracting the sum over multiples
of $3$ from three complete periods of the nontrivial character
modulo $m$.  Also $9\mid d^2S_1$.  Since $d$ is squarefree, these
facts prove integrality in all the required cases.
Direct substitution in (QE.6) gives the two exceptional values.

The same Bernoulli Fourier calculation as in (QD.1), using
$\tau(\chi_D)=i\sqrt d$, gives
\[
 \beta_3(\chi_D)=\frac{d^{5/2}}{2\pi^3}L(3,\chi_D).
\]
The Euler product is positive and
$L(3,\chi_D)\leq\zeta(3)<5/4$; the latter follows from
$1+2^{-3}+\int_2^\infty x^{-3}\,dx=5/4$.
Together with $\pi>3$ this proves (QE.4).
Under (QE.5), its upper bound is at most
$200p/216=25p/27<p$.  A positive integer smaller than $p$
is a unit modulo $p$.  The exceptional values are also units
for $p\geq31$.
\end{proof}

Therefore a $\beta_b$-unit seed in the window (QE.5) would supply
the same-discriminant product and the class-number condition.
This window has size $O(p^{2/5})$; the established Jacobi
heat/Sturm construction gives a bound of order $p$ and has not
been improved to this window.  No existence assertion in (QE.5)
is made.

\paragraph{Source and verification scope.}
The modular-symbol and Bernoulli normalization is the already
whole-reviewed Heumann--Vatsal Corollary 4.7 and
KT-EISENSTEIN-QUADRATIC-DETECTOR.  The dimension input is
W.~Stein, \emph{Modular Forms: A Computational Approach},
Proposition 6.1, printed p.~93, with its definitions on p.~92,
and the parity pairing in Theorem 8.23, printed p.~141,
in the primary author PDF
\url{https://wstein.org/books/modform/stein-modform.pdf}.
Those precise statements were read; their foundational
Riemann--Roch and period-pairing inputs are imported,
not re-proved here.
The elementary character, scalar-ideal, and integrality
arguments above are given in full.  This is a limitation and
sufficient-window result, not arithmetic movement on the
remaining rational component.

% END evidence/lean_chain/quadratic_evaluation_limits.tex


% BEGIN evidence/lean_chain/cohen_heat_cycle_limit.tex
\subsection{The special Cohen weight does not lower the heat-cycle filtration}

The following calculation addresses the specific automatic family
$b=(p+3)/2$.  It does not rule out a stronger first-coefficient
theorem obtained by other arithmetic methods.  It does rule out
lowering the current Jacobi heat construction to bounded weight
by Hasse division or by moving around its heat cycle.

\begin{proposition}[KT-COHEN-SPECIAL-HEAT-FILTRATION]
Let $p\geq31$ be prime with $p\equiv3\pmod4$, put
$b=(p+3)/2$ and $\kappa=b+1=(p+5)/2$, and let
\[
 \phi=\zeta(1-2b)E_{\kappa,1}
\]
be the normalized Jacobi Eisenstein series used in
KT-JACOBI-TWO-UNIT-FIELDS.  Let $L=L_1$ be the heat operator,
which multiplies the coefficient at $q^n\zeta^r$ by $4n-r^2$.
Write $\Omega$ for Jacobi filtration modulo $p$ at index $1$.
Then
\[
 \Omega(\phi)=\kappa,\qquad
 \Omega(L^i\phi)=\kappa+i(p+1)\quad(1\leq i\leq p-2),
 \tag{HC.1}
\]
while $L^{p-1}\phi\ne0$ modulo $p$ and
\[
 \Omega(L^{p-1}\phi)\geq\kappa.
 \tag{HC.2}
\]
For $i\geq1$ the reductions of $L^i\phi$ are periodic in $i$
with period $p-1$.  Thus every nonzero iterate has filtration
at least $\kappa$, and the first heat lift has exact filtration
\[
 \Omega(L\phi)=\frac{3p+7}{2}.
 \tag{HC.3}
\]
\end{proposition}

\begin{proof}
The $p$-integrality of $\phi$ is part of the already checked
Jacobi construction.  Its constant is
\[
 \zeta(1-2b)=-\frac{B_{p+3}}{p+3}
             \equiv-\frac{B_4}{4}=\frac1{120}\pmod p,
 \tag{HC.4}
\]
by Kummer's congruence.  The Euler factors at $p$ in this
congruence are units with residue $1$.  Hence $\phi$ is nonzero.

The congruent-weight theorem for Jacobi forms of fixed index
implies that every weight representing its nonzero reduction
is congruent to $\kappa$ modulo $p-1$.  Since
$0<\kappa<p-1$ and Jacobi forms of negative weight vanish,
this proves $\Omega(\phi)=\kappa$.

For a nonzero reduction of index $1$, the heat-filtration
theorem states
\[
 \Omega(Lf)\leq\Omega(f)+p+1,
 \quad\hbox{with equality precisely when }
 p\nmid 2\Omega(f)-1.
 \tag{HC.5}
\]
It also supplies a Jacobi modular lift of the heat image.
For $0\leq i\leq p-3$, if the claimed value holds at $i$, then
\[
 2\bigl(\kappa+i(p+1)\bigr)-1
       \equiv4+2i\not\equiv0\pmod p.
 \tag{HC.6}
\]
The unique first index in $0,\ldots,p-1$ where this expression
vanishes is $i=p-2$.  Induction by (HC.5) therefore proves
(HC.1), including nonvanishing at every step.

Coefficientwise Fermat gives $L^p\phi=L\phi$ modulo $p$.
In particular $L^{p-1}\phi$ cannot vanish, since $L\phi$
does not vanish.  Apply the inequality in (HC.5) to that
nonzero reduction:
\[
 \kappa+p+1=\Omega(L^p\phi)
       \leq\Omega(L^{p-1}\phi)+p+1.
 \]
Subtracting $p+1$ proves (HC.2).
Finally $L^{i+p-1}\phi=L^i\phi$ for every $i\geq1$, again
coefficientwise.  This proves all the assertions.
\end{proof}

\paragraph{Meaning for the short-discriminant proposal.}
The first heat image is the one used to remove the
discriminant-zero coefficients before applying the odd-component
and Sturm argument.  Equation (HC.3) shows that its linear
weight cannot be lowered by dividing out a Hasse invariant
while remaining a holomorphic Jacobi form with the same
reduction.  Iterating the heat operator never reaches a
filtration smaller than $(p+5)/2$ either.  This excludes that
specific way of replacing the existing $O(p)$ bound by the
$O(p^{2/5})$ sufficient window in KT-BETA3-SHORT-WINDOW.
It does not give a lower bound on the actual first nonzero
fundamental coefficient.

The nonzero constant term alone is not a field coefficient.
In Jacobi notation, its translates have $4n-r^2=0$; the heat
operator kills all of them.  The proof needs a nonzero
coefficient with positive discriminant and with odd $r$.
Thus the constant in (HC.4) cannot simply be counted as the
desired seed.

The fixed-weight identity
$\beta_b(\chi_D)\equiv\beta_2(\chi_D\chi_{-p})$
does not contradict (HC.1).  Its real quadratic character has
conductor $p|D|$, and the corresponding weight-$5/2$ Cohen
coefficient is at index $p|D|$, not at index $|D|$.
A small ordinary Sturm bound for the fixed-weight form
does not bound its first unit coefficient in this
$p$-dependent subsequence.
For the short window itself the class-number requirement
would present no extra gap: $d^5\leq1600p^2$ implies
$d<4p$ for $p\geq31$, and the carried
KT-SMALL-DISCRIMINANT-CLASS-BOUND gives $h(-d)<p$.
Existence of a $\beta_b$-unit field in that window remains
unproved here.

\paragraph{Exact primary input and review scope.}
D.~Choi, Y.~Choie and O.~K.~Richter,
\emph{Congruences for Siegel modular forms of degree two},
Ann.~Inst.~Fourier \textbf{61} (2011), 1455--1466:
Section 2, Propositions 1 and 2 on printed p.~1457 give
the congruent-weight theorem and (HC.5).
Their hypotheses are even Jacobi weight, $p\geq5$, and
fixed index; here $\kappa$ and all lift weights are even
and the index is $1$.  The definitions and both full
statements, with the surrounding scope, were read directly
in the retained primary text
\path{main_gate/literature/character_unit_followup/Choi_Choie_Richter_2011}.
These published propositions are imported inputs; their
foundational Sofer and Richter proofs are not newly
reconstructed in this deduction.  The boundary case
$p=2\kappa-5$ is not inserted into either strict inequality
in their Theorem 2.

% END evidence/lean_chain/cohen_heat_cycle_limit.tex


% BEGIN evidence/lean_chain/cohen_weight_two_bound.tex
\subsection{A sharper first-unit bound in the special Cohen weight}

This argument improves the discriminant bound for the automatic family
$b=(p+3)/2$.  It supplies the class-number unit and the first twisted
Bernoulli unit.  It does not supply the second twisted unit
$\beta_3(\chi_D)$ required by the same-discriminant detector.

\begin{theorem}[KT-COHEN-WEIGHT-TWO-BOUND]
Let $p\geq31$ be prime with $p\equiv3\pmod4$, and set
$b=(p+3)/2$.  There exists an odd negative fundamental discriminant
$D=-d$ such that
\[
 3\leq d\leq p-4,\qquad
 p\nmid h(D),\qquad \beta_b(\chi_D)\in\mathbf Z_p^\times.
 \tag{CW.1}
\]
In particular $p\nmid D$.  By the already checked character transport,
$B_{1,\chi_D}$ and $B_{1,\chi_D\omega^{b-1}}$ are units, and
$\beta_n(\chi_D)$ is a unit for every odd $n\geq3$ congruent to $b$
modulo $p-1$.
\end{theorem}

\begin{proof}
Write $q=\exp(2\pi i\tau)$ and use the Cohen normalization
\[
 H_b(\tau)=\sum_{n\geq0}H(b,n)q^n
       \in M_{b+1/2}^{+}(\Gamma_0(4)).
 \tag{CW.2}
\]
The index-one Jacobi form corresponding to it is
$\zeta(1-2b)E_{b+1,1}$, as in KT-JACOBI-TWO-UNIT-FIELD.
Its full coefficient formula and the integrality proof there show
that every $H(b,n)$ is $p$-integral.  This includes indices divisible
by $p$: here $3\leq b\leq p-2$ and $b\ne(p-1)/2$.
The constant and support are
\[
 H(b,0)=-\frac{B_{p+3}}{p+3}\equiv\frac1{120}\pmod p,
 \qquad H(b,n)=0\quad(n\equiv1,2\pmod4).
 \tag{CW.3}
\]
The constant congruence is the same Kummer congruence used in
KT-COHEN-SPECIAL-HEAT-FILTRATION; in particular it is a unit.

Put
\[
 \theta(\tau)=\sum_{a\in\mathbf Z}q^{a^2},\qquad
 \Theta=\theta^4,
 \qquad G=2E_2(2\tau)-E_2(\tau).
 \tag{CW.4}
\]
Here $G\in M_2(\Gamma_0(2))\subset M_2(\Gamma_0(4))$ and
$\Theta\in M_2(\Gamma_0(4))$, both with trivial character and
integral Fourier coefficients.  The transformation law of $E_2$
shows the asserted modularity of $G$: its two nonmodular terms
cancel.  Their first coefficients are
\[
 \Theta=1+8q+24q^2+O(q^3),\qquad
 G=1+24q+24q^2+O(q^3).
 \tag{CW.5}
\]
The theta product
$\theta=\prod_{n\geq1}(1-q^{2n})(1+q^{2n-1})^2$
has no zeros in the upper half-plane.  The three cusps of
$\Gamma_0(4)$ are $\infty,0,1/2$.  At $\infty$ the theta
constant is $1$.  At $0$, Poisson summation gives
\[
 (\Theta|_2S)(\tau)=-\frac14\theta(\tau/4)^4,
 \qquad S=\begin{pmatrix}0&-1\\1&0\end{pmatrix},
 \tag{CW.6}
\]
so its constant is nonzero.  At $1/2$, with
$A=\left(\begin{smallmatrix}1&0\\2&1\end{smallmatrix}\right)$,
the same Poisson formula, with consistent holomorphic square roots,
is
\[
 (2\tau+1)^{-1/2}\theta(A\tau)
      =\sum_{a\text{ odd}}q^{a^2/4}
      =2q^{1/4}(1+O(q^2)).
 \tag{CW.7}
\]
Consequently $\Theta|_2A=16q+O(q^3)$; its sole zero is
the simple zero at this cusp.  The cusp has width $1$.
Also $G|_2A=G$, because $A\in\Gamma_0(2)$, so its constant
at that cusp is exactly $1$.

Now form the meromorphic modular form
\[
                       F=H_b/\theta^p.
 \tag{CW.8}
\]
Its weight is $b+1/2-p/2=2$.  More precisely the theta
automorphy factor of $H_b$ has exponent $2b+1=p+4$;
division by $\theta^p$ leaves the exponent $4$, which is
exactly the trivial-character weight-two factor on $\Gamma_0(4)$.
Thus there is no residual character in (CW.8).
This form is holomorphic away from the possible pole at $1/2$.
Its series at $\infty$ is $p$-integral, since $\theta$ has
constant $1$ and integral coefficients.  Frobenius gives
$\theta^p\equiv\theta(q^p)\pmod p$.  Therefore
\[
 [q^0]F\equiv1/120,
 \qquad [q]F\equiv[q^2]F\equiv0\pmod p.
 \tag{CW.9}
\]

We record explicitly which coefficients control its pole.
The two-component theta decomposition is
\[
 H_b(\tau)=h_0(4\tau)+h_1(4\tau),\qquad
 h_1(\tau)=\sum_{n\equiv3\ (4)}H(b,n)q^{n/4}.
 \tag{CW.10}
\]
For $v=\tau+1/2$ one has $4A\tau=2-1/v$.
The dual index-one theta representation has
$T^2=\operatorname{diag}(1,-1)$ and its $S$ matrix is an
eighth root of unity times
$2^{-1/2}\left(\begin{smallmatrix}1&1\\1&-1\end{smallmatrix}\right)$.
Multiplying its row $(1,-1)$ by that $S$ matrix leaves a
unit multiple of $(0,\sqrt2)$.  Hence, for an eighth root
of unity $u$ whose precise value is immaterial here,
\[
 (2\tau+1)^{-b-1/2}H_b(A\tau)
               =u\,2^{-b}h_1(\tau+1/2).
 \tag{CW.11}
\]
This derivation fixes both the component and the power of $2$;
all scalar factors are $p$-adic units.
Dividing (CW.11) by the $p$th power of (CW.7) shows that
$F|_2A$ has a $p$-integral Laurent expansion in integral
powers of $q$.  Its possible negative powers arise only
from $H(b,n)$ with $0<n<p$ and $n\equiv3\pmod4$.
Indeed its numerator terms have exponent $n/4$, whereas
its denominator is $2^p q^{p/4}$ times a power-series unit.
Since $n\equiv p\pmod4$, the resulting exponents are integral.
Its pole order is at most $(p-3)/4$.

Suppose, for contradiction, that all these odd-index coefficients
are zero modulo $p$.  Then every coefficient in the principal
part of $F|_2A$ is divisible by $p$.
The following explicit removal of that principal part justifies
the reduction step without an assumed base-change theorem for
half-integral weight forms or a special-fiber divisor assertion.
Work in a number field containing $u$ and the eighth roots of unity,
localized at a prime above $p$.  For every $j\geq1$ put
\[
                  W_j=G^{j+1}/\Theta^j.
 \tag{CW.12}
\]
These are weight-two meromorphic forms with trivial character.
They are holomorphic away from $1/2$, are integral at $\infty$,
and satisfy
\[
 W_j|_2A=16^{-j}q^{-j}+\text{higher powers of }q.
 \tag{CW.13}
\]
All Laurent coefficients here are integral at $p$, by
(CW.5) and (CW.7).  Start with the largest pole order of $F$.
Its coefficient, multiplied by $16^j$, is an element divisible
by $p$.  Subtract this multiple of $W_j$.  The leading pole
disappears; all remaining negative coefficients stay divisible
by $p$, because $W_j|_2A$ is integral.  Repeating finitely many
times yields a holomorphic characteristic-zero form
\[
                  F_0\in M_2(\Gamma_0(4))
 \tag{CW.14}
\]
whose expansion at $\infty$ is integral and congruent to that
of $F$ modulo $p$.  Every subtracted multiple is divisible
by $p$ and each $W_j$ is integral there.  Holomorphy at the
other cusps and in the upper half-plane was never changed.

The ordinary Sturm bound at weight $2$, level $4$ is
$2[\mathrm{SL}_2(\mathbf Z):\Gamma_0(4)]/12=1$.
Thus an integral form modulo $p$ is determined by its constant
and its $q$ coefficient.  The constant/$q$ coefficient matrix
of $\Theta,G$ has determinant $16$, a unit.  There are integral
scalars $a,c$ whose linear combination $a\Theta+cG$ has the
same first two coefficients as $F_0$.  By the Sturm theorem
they agree modulo $p$ at every coefficient.  Its $q$ and
$q^2$ coefficients would consequently satisfy
\[
 8a+24c=0,\qquad24a+24c=0\pmod p.
 \tag{CW.15}
\]
The determinant is $-384$, a unit for $p\geq31$, so
$a=c=0$ modulo $p$.  This contradicts the constant
$a+c\equiv1/120$ in (CW.9).  The contradiction proves that
some $H(b,n)$ with $0<n<p$, $n\equiv3\pmod4$, is a unit.

Choose the smallest such $n$.  Write $-n=D f^2$ with
$D<0$ fundamental.  It is odd, and $|D|\leq n<p$.
The full Cohen coefficient formula is
\[
 H(b,n)=-\beta_b(\chi_D)
       \sum_{v\mid f}\mu(v)\chi_D(v)v^{b-1}
                                  \sigma_{2b-1}(f/v).
 \tag{CW.16}
\]
The sum is an integer and the Bernoulli factor is $p$-integral.
A unit left side therefore forces $\beta_b(\chi_D)$ to be
a unit.  If $f>1$, the smaller fundamental coefficient
$H(b,|D|)=-\beta_b(\chi_D)$ would be a unit, contradicting
minimality.  Thus $n=d=|D|$ is fundamental.
Since $p$ and $d$ are both $3$ modulo $4$ and $d<p$,
one has $d\leq p-4$.
The carried elementary class-number estimate
KT-SMALL-DISCRIMINANT-CLASS-BOUND applies since $d<4p$;
it gives $h(D)<p$.  This proves (CW.1).
The stated character and all-weight consequences are exactly
the previously checked Kummer transport at conductor prime to $p$.
\end{proof}

\paragraph{Exact inputs and scope.}
The Cohen/Jacobi coefficient formula and its $p$-integrality,
including the conductor-divisible-by-$p$ case, are the already
whole-reviewed proof of KT-JACOBI-TWO-UNIT-FIELD, using
Eichler--Zagier's index-one theta decomposition and Theorem~2.1
(1983 primary preprint, printed pp.~19a--20 for the full
coefficient formula).  The two-component $S,T$ calculation is
displayed in (CW.10)--(CW.11); the theta Poisson normalization is
the same as the whole-reviewed KT-KOHNEN-THETA-NONCONGRUENCE.
The ordinary elliptic Sturm theorem is used with the exact
coefficient-field generality stated by Choi--Choie--Kikuta
(2013), printed p.~129.  These are imported modular-form inputs;
no new foundational proof of the Sturm theorem is claimed.
The finite principal-part removal proves the needed integral
reduction directly, so the argument does not require an
unverified identification of a special-fiber half-integral
weight space with reductions of characteristic-zero forms.

The new discriminant bound is linear in $p$.  It does not put
$d$ in the sufficient window $d^5\leq1600p^2$ of
KT-BETA3-SHORT-WINDOW.  No unit assertion for
$\beta_3(\chi_D)$, no full three-unit auxiliary field, and
no new vanishing of an original rational class-group component
is inferred from (CW.1).

% END evidence/lean_chain/cohen_weight_two_bound.tex


% BEGIN evidence/lean_chain/cohen_general_first_unit_bound.tex
\subsection{A general first-unit bound and its ramification boundary}

The same principal-part argument has a useful general form.
Its prime-to-$p$ scope must be retained: outside the special
weight, the first unit index could be $p$ itself.

\begin{proposition}[KT-COHEN-GENERAL-FIRST-UNIT-BOUND]
Let $p\geq29$ be prime, and let $b$ be odd with
\[
 3<b\leq p-2,\qquad b\ne(p-1)/2,
 \qquad C=-B_{2b}/(2b)\in\mathbf Z_p^\times.
 \tag{CG.1}
\]
Then there exists an odd negative fundamental discriminant
$D=-d$ such that
\[
 3\leq d\leq2b-7<2p,
 \qquad p\nmid h(D),\qquad
 \beta_b(\chi_D)\in\mathbf Z_p^\times.
 \tag{CG.2}
\]
If either $p\equiv1\pmod4$ or $2b-7<p$, the discriminant
in (CG.2) is automatically prime to $p$.  In these cases
the two generalized $B_1$ units and the all-weight transport
of KT-JACOBI-TWO-UNIT-FIELD apply to this same $D$.
For $b=3$, the prime-to-$p$ choice $D=-3$ has
$h(D)=1$ and $\beta_3(\chi_D)=2/9$.
\end{proposition}

\begin{proof}
Use the notation and normalizations of
KT-COHEN-WEIGHT-TWO-BOUND, but set
\[
                  m=2b-3,\qquad F_b=H_b/\theta^m.
 \tag{CG.3}
\]
The difference of theta-multiplier exponents is again
$2b+1-m=4$, so this is weight two with trivial character.
It is meromorphic only at the cusp $1/2$.
The complete coefficient-integrality proof in
KT-JACOBI-TWO-UNIT-FIELD applies under (CG.1), including
at indices divisible by $p$.

Since $H_b=C+O(q^3)$ and $\theta=1+2q+O(q^4)$,
the ordinary binomial expansion gives
\[
 F_b=C-2mCq+2m(m+1)Cq^2+O(q^3).
 \tag{CG.4}
\]
The cusp calculation (CW.10)--(CW.11) has not changed.
Its numerator terms are $p$-adic units times
$H(b,n)q^{n/4}$ for $n\equiv3\pmod4$;
the denominator in (CG.3) is $2^m q^{m/4}$ times
an integral power-series unit.  Here $m\equiv3\pmod4$.
Consequently, if every $H(b,n)$ with $0<n<m$ and
$n\equiv3\pmod4$ vanished modulo $p$, the principal
part of $F_b$ at $1/2$ would be divisible by $p$.

The explicit forms $W_j=G^{j+1}/\Theta^j$ in
(CW.12)--(CW.13) remove that principal part, using only
multiples divisible by $p$.  This would produce a
holomorphic characteristic-zero weight-two form, integral
at $p$, with the same three displayed coefficients modulo
$p$ as (CG.4).  The same coefficient-field Sturm bound
one applies.
Both basis forms $\Theta$ and $G$ have their $q^2$
coefficient equal to $24$ times their constant coefficient.
Their constant/$q$ coefficient determinant is $16$, so
the Sturm theorem implies this relation for every
holomorphic weight-two form modulo $p$ at level $4$.
For (CG.4) the required relation would be
\[
 \begin{aligned}
 0&=\bigl(2m(m+1)-24\bigr)C\\
  &=2(m+4)(m-3)C
    =4(2b+1)(b-3)C\pmod p.
 \end{aligned}
 \tag{CG.5}
\]
This is impossible.  The factors $4$ and $C$ are units;
$0<b-3<p$; and $0<2b+1<2p$, whose only possible
multiple of $p$ would require $b=(p-1)/2$, explicitly
excluded in (CG.1).

There is therefore a unit odd-index coefficient at some
$n<m$.  The minimal such index is fundamental, by exactly
the integral divisor-factor argument (CW.16).
Since both $m$ and $n$ are $3$ modulo $4$, it satisfies
$d=n\leq m-4=2b-7$.
The carried elementary class-number bound applies since
$d<2p<4p$, so $h(D)<p$.

The only positive multiple of $p$ that could lie below
$2p$ is $p$.  It cannot equal $d$ if $p\equiv1\pmod4$,
because $d\equiv3\pmod4$; it cannot occur if $d<p$
either.  This proves the asserted prime-to-$p$ cases.
There the existing character transport has its required
conductor hypothesis.  The separate $b=3$ assertion is
the stated exact elementary value and class number.
\end{proof}

\paragraph{What the general bound does not imply.}
For $p\equiv3\pmod4$ and $2b-7\geq p$, this proof
does not rule out $d=p$.  Integrality of a generalized
Bernoulli value at conductor $p$ does not imply its
divisibility by $p$.  No claim that (CG.2) always satisfies
the unramified auxiliary-field condition is made in that
range.  The earlier heat construction supplies its own
prime-to-$p$ field with its earlier bound; combining the
two arguments does not assign the sharper bound to that
field without a further step.

At the special weight $b=(p+3)/2$ with $p\equiv3\pmod4$,
one has $2b-7=p-4$, recovering (CW.1).
Neither the general nor the special bound proves a
same-discriminant unit for the additional twisted
Bernoulli value.  This is an auxiliary first-unit bound,
not a new vanishing statement for the original rational
class-group component.

% END evidence/lean_chain/cohen_general_first_unit_bound.tex


% BEGIN evidence/lean_chain/cohen_hecke_support.tex
\subsection{The special Cohen series and the limits of square-class Hecke recurrences}

This is a bounded attempt at the same-discriminant unit problem.
The first result strengthens the carried fixed-weight congruence to
the full Fourier series.  The second identifies precisely what the
coefficient recurrences fail to supply.  Neither result asserts a
simultaneous unit or removes a rational class component.

\begin{proposition}[KT-COHEN-UP-REFORMULATION]
Let $p\geq31$ be prime, $p\equiv3\pmod4$, and $b=(p+3)/2$.
Use the Cohen series
$H_r=\sum_{n\geq0}H(r,n)q^n$, and define coefficient extraction by
\[
 U_p\!\left(\sum_{n\geq0}a(n)q^n\right)
                  =\sum_{n\geq0}a(pn)q^n.
\]
Then every coefficient of $H_b$ and $H_2$ is $p$-integral, and
\[
                           H_b\equiv U_p H_2\pmod p
                  \quad\hbox{in }\mathbf F_p[[q]].       \tag{HS.1}
\]
This is a Fourier-series identity.  It does not assert that $U_p$
preserves the level-four, weight-$5/2$ space.
\end{proposition}

\begin{proof}
For a fundamental discriminant $D$, allowing $D=1$ with the
trivial character, put
\[
 P_r(D,m)=\sum_{v\mid m}\mu(v)\chi_D(v)v^{r-1}
                                    \sigma_{2r-1}(m/v).
                                                               \tag{HS.2}
\]
The full Cohen formula is
\[
 H(r,n)=-\beta_r(\chi_D)P_r(D,m),
                  \qquad (-1)^rn=Dm^2.               \tag{HS.3}
\]
At the constant coefficient,
$H(b,0)=-B_{p+3}/(p+3)\equiv-B_4/4=1/120=H(2,0)$.
If $n\equiv1,2\pmod4$, both sides of (HS.1) are zero,
since $pn\equiv3,2\pmod4$ respectively.

Now write $-n=Dm^2$ with $D<0$ fundamental.  The primitive
quadratic character associated with $\chi_D\chi_{-p}$ has positive
fundamental discriminant
\[
 D^\dagger=
 \begin{cases}
  -pD,&p\nmid D,\\
  -D/p,&p\mid D.
 \end{cases}
 \qquad
 m^\dagger=
 \begin{cases}
  m,&p\nmid D,\\
  pm,&p\mid D,
 \end{cases}
 \qquad pn=D^\dagger(m^\dagger)^2.                    \tag{HS.4}
\]
Here the second case includes $D=-p$, $D^\dagger=1$.
The factorization of fundamental discriminants is valid also
when $D$ is even, because $-p$ is odd and its prime occurs in
$D$ to exponent exactly one.

For every integer $a$ prime to the two conductors,
Euler's criterion gives
\[
 \chi_D(a)a^{b-1}
       \equiv\chi_D(a)\chi_{-p}(a)a
       =\chi_{D^\dagger}(a)a\pmod p.                 \tag{HS.5}
\]
Apply the change-of-character Kummer congruence in Guerzhoy's
Proposition~2 to $(\chi_D,b)$ and $(\chi_{D^\dagger},2)$.
Both conductors divide $\Delta p$ for the same $\Delta$ prime to
$p$.  A common regularizer can be chosen with unit residue.
If $\Delta>1$, its quadratic prime-to-$p$ factor is nontrivial;
choose $c\equiv1\pmod p$ and with that factor equal to $-1$.
Then both factors $1-\chi(c)c^r$ are $2$ modulo $p$.
If $\Delta=1$, necessarily $D=-p$, and $c=2$ gives residue
$1-c^2=-3$, again a unit.  In both cases the Euler factors
$1-\chi(p)p^{r-1}$ are $1$ modulo $p$.
The integral measure in Proposition~1 therefore gives
integrality and
\[
                    \beta_b(\chi_D)
                    \equiv\beta_2(\chi_{D^\dagger})\pmod p.
                                                               \tag{HS.6}
\]
This explicitly includes the conductor-divisible-by-$p$ branch,
rather than inferring it from a prime-to-$p$ formula.
The ordinary value $\beta_2(1)=B_2/2=1/12$ is integral as well.

It remains to compare the integer factors (HS.2).
They are multiplicative in $m$.  At a prime $\ell$ their factor is
\[
 P_r(D,\ell^s)=
 \sum_{i=0}^{s}\ell^{i(2r-1)}
   -\chi_D(\ell)\ell^{r-1}
                \sum_{i=0}^{s-1}\ell^{i(2r-1)},       \tag{HS.7}
\]
with an empty second sum when $s=0$.
For $\ell\ne p$ one has
\[
 \ell^{2b-1}\equiv\ell^3,\qquad
 \chi_D(\ell)\ell^{b-1}
      \equiv\chi_{D^\dagger}(\ell)\ell\pmod p.
                                                               \tag{HS.8}
\]
For $\ell=p$, (HS.7) is $1$ modulo $p$ for every $s\geq0$
and every $r\geq2$.  Thus the possible extra factor $p$ in
$m^\dagger$ has no effect, and
$P_b(D,m)\equiv P_2(D^\dagger,m^\dagger)\pmod p$.
Equations (HS.3) and (HS.6) prove (HS.1).

For completeness, the same measure argument gives the
$p$-integrality of every quadratic value $\beta_2$ involved:
its integrand $\chi(a)a$ has a regularizing unit since
$p\geq31$; the exceptional possibility of a residue character
$a^{-2}$ is impossible for a quadratic character at such $p$.
Alternatively, the degree-two Bernoulli-polynomial formula
gives it directly.  Together with (HS.2), this covers all
coefficients of $H_2$; the already checked Cohen integrality
also supplies the coefficients of $H_b$.
\end{proof}

\begin{proposition}[KT-HECKE-SQUARE-CLASS-SCOPE]
Let $r\geq3$ be odd.  Reduction of the full Cohen formula
has the following consequences.

First, if $\beta_r(\chi_D)=0$ modulo $p$, then
$H(r,|D|m^2)=0$ modulo $p$ for every $m$.
In particular, an existing first-unit seed at a discriminant
where $\beta_3$ vanishes cannot be repaired by replacing its
index by square multiples.

Second, the standard Fourier-coefficient equations
\[
 a(\ell^2n)+\left(\frac{-n}{\ell}\right)\ell^{r-1}a(n)
       +\ell^{2r-1}a(n/\ell^2)
       =(1+\ell^{2r-1})a(n),\qquad \ell\nmid2p,         \tag{HS.9}
\]
where $a(n/\ell^2)=0$ if $\ell^2\nmid n$,
do not connect distinct fundamental discriminants.
Even when imposed for every such $\ell$, they leave their
fundamental coefficients formally independent, and they leave
the constant coefficient independent of all positive indices.
\end{proposition}

\begin{proof}
The first assertion is immediate from (HS.2)--(HS.3).
For the second, prescribe arbitrary residues $A_D$ indexed by
negative fundamental discriminants, and an arbitrary constant $c$.
Define a formal series by
\[
 a(0)=c,\qquad
 a(|D|m^2)=A_D P_r(D,m),
 \qquad a(n)=0\ \hbox{at the forbidden plus indices}.   \tag{HS.10}
\]
Every positive allowed index has a unique such factorization.
At $\ell\nmid m$, the $s=1$ case of (HS.7) gives (HS.9),
including $\ell\mid D$ when the character value is zero.
At $\ell\mid m$, the character in (HS.9) is zero and
(HS.7) gives the elementary recurrence
\[
 P_r(D,\ell^{s+1})+\ell^{2r-1}P_r(D,\ell^{s-1})
       =(1+\ell^{2r-1})P_r(D,\ell^s),\qquad s\geq1.
                                                               \tag{HS.11}
\]
The factors at the other primes do not change.
The equation at $n=0$ is simply
$c+\ell^{2r-1}c=(1+\ell^{2r-1})c$.
This proves all the claimed formal equations without any
restriction on the $A_D$.
The Fourier derivative $\Theta=q\,d/dq$ also preserves every
zero coefficient: it multiplies $a(n)$ by $n$.
Iterating it, or combining it with these square-index
coefficient operations, does not create a coefficient in a
previously absent fundamental square class.

For an actual retained example, $p=83$, $D=-43$ satisfies
$\beta_3(\chi_D)=166\equiv0$ and
$\beta_{43}(\chi_D)\equiv76\ne0$.
In (HS.10) with $r=43$, choose $c=1/120$,
$A_{-43}=-76$, and all other $A_D=0$.
This series satisfies every equation (HS.9), has the required
nonzero constant, and has zero coefficientwise product with
$H_3$ at every positive index.  It is a formal example only:
no modularity is asserted.  It isolates the insufficiency of
the Hecke equations themselves, without claiming that the
actual Cohen series has this support.
\end{proof}

\paragraph{Remaining modular question.}
The simultaneous field problem would follow from a nonzero
coefficientwise product
$H(2,pd)H(3,d)$ at an odd fundamental $d<p$.
The first-unit theorem gives a nonzero factor $H(2,pd)$ in that
range, but does not give the other factor.
Neither ordinary multiplication of modular forms nor its
no-zero-divisor property applies to a coefficientwise product.
The modularity of (HS.10) is exactly the omitted global constraint;
the formal construction is not a counterexample to a possible
special modular-form theorem.
The zero set of $\beta_3$ has not been shown to contain only
finitely many fundamental square classes, so a theorem excluding
finite square-class support cannot be invoked here.
No fixed-level modular realization or bounded-weight Sturm
criterion for this coefficientwise product has been established
in this attempt.  Consequently (HS.1) and (HS.9) do not supply
the missing simultaneous unit, even though they give an exact
fixed-weight reformulation and identify the limits of
square-class Hecke amplification.

\paragraph{Inputs and check scope.}
The full Cohen formula is the already checked coefficient
formula, also stated in P.~Guerzhoy,
\emph{Quadratic congruences for Cohen--Eisenstein series},
Glasgow Math.~J.~41 (1999), 141--144, equation~(1), p.~142.
The square-index convention $D=1$ is explicit in
A.~Balog, W.~J.~McGraw and K.~Ono,
\emph{The arithmetic of the coefficients of half-integral weight
Eisenstein series}, Acta Arith.~107 (2003), 337--344,
definitions~(1.2)--(1.4), printed pp.~337--338.
Their formulas at $r=1$, $k=2$ give
$H(2,m^2)=-P_2(1,m)/12$, including $H(2,1)=-1/12$.
Only these definitions are used; their later main theorem
is not needed.
Guerzhoy's Propositions~1--2 and equations~(2)--(4), pp.~143--144,
give precisely the integral measure and character-changing
congruence used in (HS.5)--(HS.6), with conductor divisible
by $p$ permitted.
These statements and the short deduction of Proposition~2 were
read in full.  The published construction of the Mazur measure
remains an imported input.
The recurrence (HS.9) is proved directly from (HS.7) here,
so it requires no unverified converse from a formal eigenseries
to a modular form.
The $p=83$ values are carried exact controls; no prime census,
new discriminant search, or new finite $K$-degree computation
is involved.

% END evidence/lean_chain/cohen_hecke_support.tex


% BEGIN evidence/lean_chain/hadamard_modularity_scope.tex
\subsection{Global modularity tests for a same-discriminant unit}
\label{sec:hadamard-modularity-scope}

This is a bounded method audit, not a simultaneous-unit theorem.  Let
$p\ge31$ be prime with $p\equiv3\pmod4$ and put $b=(p+3)/2$.
The already proved full-series congruence
$H_b\equiv U_pH_2\pmod p$ turns the remaining coefficient target into
\[
 H(2,pd)H(3,d)\ne0\pmod p,
 \qquad -d<0\text{ odd fundamental},\quad d<p.
 \tag{HM.1}
\]
The first-unit theorem supplies such a $d$ for the first factor, and its
class number is already a unit.  It does not supply the second factor.
The operations below do not complete this step.

\paragraph{An exact same-index Dirichlet series is not yet an integral
modular detector.}
Mizuno's \emph{Rankin--Selberg convolutions of noncuspidal half integral
weight Maass forms in the plus space}, Nagoya Math.\ J.\ 237 (2020),
127--165, Section~1.2, equation~(6) and Theorem~1, printed pp.~133--134,
does treat a series whose summands contain
\[
 L_{-d}\bigl(\sigma-(k+1)/2\bigr)
 L_{-d}\bigl(\eta-(k+1)/2\bigr).
 \tag{HM.2}
\]
Here the $L_D$ notation includes the nonfundamental-discriminant divisor
factors.  For its required $k\equiv1\pmod4$, the theorem gives
meromorphic continuation in three complex
variables, an explicit holomorphic regularization, and a vector functional
equation.  The parameters exclude poles of the summands before continuation.
Proposition~3 and its complete proof, printed pp.~152--153, and the
complete proof of Theorem~1, pp.~153--154, establish these assertions by
regularized integrals and an invertible matrix of Whittaker integrals.
They do not identify the coefficients in (HM.2) with a holomorphic
modular $q$-series over a specified $p$-integral lattice.  In particular,
this cited theorem contains neither an integral residue formula forcing
one coefficient to be a unit, nor a bound $d<p$.  A nonzero complex
residue alone supplies neither conclusion.  The partial Euler product
in Section~3.4, p.~154, sums the square-discriminant subseries; it is not
an Euler product for the negative fundamental discriminants in (HM.1).

Mizuno's separate \emph{Petersson norms of Eisenstein series and
Kohnen--Zagier's formula}, J.\ Th\'eor.\ Nombres Bordeaux 32 (2020), 665--684,
Theorem~2.4, p.~670, gives the regularized Kohnen--Zagier formula for
\emph{one} Cohen Eisenstein coefficient squared.  For a fundamental $D$
with $(-1)^rD>0$, its proof in Section~5, pp.~680--681, uses
\[
 L(G_{2r},D,s)=L(s,\chi_D)L(s+1-2r,\chi_D)
 \tag{HM.3}
\]
and the functional equation to obtain the square of $L(1-r,\chi_D)$.
The complete proof and the norm-ratio normalization have been read.
This is not the mixed product of the weights $b$ and $3$ in (HM.1).
No $p$-unit conclusion follows just from the nonzero real regularized
norms.  The earlier computations of those norms are named inputs to
that proof, not newly rechecked in this audit.

\paragraph{A saturated Rankin--Cohen trial still fails the diagonal test.}
The standard Rankin--Cohen formula is a bilinear differential expression,
so its Fourier coefficients are additive convolutions.  One might try
order $p-1$ to select equal indices by Fermat's theorem, while keeping
the formal weight linear in $p$.  The following exact calculation rules
out this particular trial, even after its common factor $p$ is removed.
It is a statement about the displayed universal operator; it does not
exclude an accidental identity for the particular Cohen series.

\begin{proposition}[KT-SATURATED-RANKIN-COHEN-SCOPE]
Set $k=(p+4)/2$, $\ell=7/2$, and $N=p-1$, and define in
$\mathbb Z_{(p)}$
\[
 c_r=(-1)^r\binom{N+k-1}{N-r}\binom{N+\ell-1}{r}
     =(-1)^r\binom{3p/2}{p-1-r}\binom{p+3/2}{r},
 \qquad 0\le r\le p-1.
 \tag{HM.4}
\]
Every $c_r$ is divisible by $p$.  For $s_r=c_r/p$ one has
\[
 s_r\equiv0\pmod p\quad(b+1\le r\le p-2),
 \qquad
 s_r\not\equiv0\pmod p\quad(0\le r\le b\text{ or }r=p-1).
 \tag{HM.5}
\]
Put $S(X,Y)=\sum_{r=0}^{p-1}\bar s_rX^rY^{p-1-r}$.
There are no scalars $\alpha,\beta\in\mathbb F_p$ such that, for all
$x,y\in\mathbb F_p$ with $x+y\ne0$,
\[
 \alpha S(x,y)+\beta=\begin{cases}1,&x=y,\\0,&x\ne y.\end{cases}
 \tag{HM.6}
\]
Thus this saturated operator, even with an arbitrary scalar multiple
of the ordinary product added, is not a universal equal-index selector.
\end{proposition}

\begin{proof}
All factorial denominators in (HM.4) have indices below $p$.
If $s=p-1-r>0$, the numerator of $\binom{3p/2}{s}$, after multiplying
by the unit $2^s$, is
$\prod_{j=0}^{s-1}(3p-2j)$.  Exactly its $j=0$ factor is divisible by
$p$, and it has valuation one.  For $s=0$ the binomial is $1$.
The numerator of the other binomial, after multiplying by $2^r$, is
$\prod_{j=0}^{r-1}(2p+3-2j)$.  It contains a multiple of $p$ precisely
when $r\ge b+1$, and that multiple is its $j=b$ factor, equal to $p$.
Its valuation is then one.  This proves the exact valuations and (HM.5).

For $x+y\ne0$ the desired indicator is the value of the homogeneous
polynomial
\[
 D(X,Y)=(X+Y)^{p-1}-(X-Y)^{p-1}
       =\sum_{r=0}^{p-1}\bigl((-1)^r-1\bigr)X^rY^{p-1-r}
       \quad\text{over }\mathbb F_p.
 \tag{HM.7}
\]
If (HM.6) held, the homogeneous polynomial
$\alpha S+\beta(X+Y)^{p-1}-D$ would vanish on every point of the line
$X+Y=1$.  Its restriction to that line has degree at most $p-1$ and
has $p$ roots, so is the zero polynomial.  Homogenization recovers the
original homogeneous polynomial, which is therefore zero.
The coefficients of $X^{p-3}Y^2$ and $X^{p-2}Y$ now contradict each
other: both corresponding coefficients of $S$ vanish by (HM.5), while
those of $(X+Y)^{p-1}$ are $1,-1$ and those of $D$ are $0,-2$.
They would force $\beta=0$ and $\beta=2$.
\end{proof}

Writing $\Theta=q\,d/dq$, the formal operator just tested is
\[
 \mathcal R(f,g)=\sum_{r=0}^{p-1}s_r
          \Theta^rf\,\Theta^{p-1-r}g.
 \tag{HM.8}
\]
Its symbol has been checked on arbitrary residue indices.  The proposition
does not assert that a pair of genuine modular forms can have arbitrary
Fourier coefficients, and does not prove that no other bilinear operator
or global modular relation works.  If a compatible modular realization
is used, adding a multiple of $A^2fg$, with $A\equiv1\pmod p$ a Hasse
form of weight $p-1$, has exactly the additional symbol $\beta$ already
excluded.  Restricting to the actual Cohen forms or to a finite short
window could impose extra relations not addressed by this calculation.

The source formula for the trial is Zagier, \emph{Modular forms and
differential operators}, Proc.\ Indian Acad.\ Sci.\ (Math.\ Sci.)
104 (1994), 57--75, Section~1, equation~(1), p.~58.  The entire Section~1
and both modularity explanations on pp.~58--59 were read; the displayed
binomial formula was also visually checked.  That paper's opening
convention is integral-weight modular forms.  The proposition above is
pure coefficient algebra with the specified half-integral parameters,
and does not silently import a half-integral-weight theorem from that
convention.  No modularity assertion is needed for its obstruction.

\paragraph{Finite square-class support is a different hypothesis.}
Ahlgren--Choi--Rouse, \emph{Congruences for level four cusp forms},
Math.\ Res.\ Lett.\ 16 (2009), author Theorem~1.2, pp.~2--3, assumes a
nonzero cusp form in the level-four plus space, integral at a place
above a prime $p\ge5$, whose reduction is supported on \emph{finitely
many} square classes.  Under
$\lambda+1/2<p(p+1+1/2)$ it concludes that $\lambda$ is even and gives
a derivative-of-theta formula.  The complete stated theorem and its
final proof on author pp.~16--17 were read; the intermediate
propositions used in that proof are not newly certified here.
The weight bound is not the decisive missing premise in (HM.1).
$H_b$ is Eisenstein, and the assertion that its support lies in
$\{D:\beta_3(\chi_D)=0\bmod p\}$ does not imply that this set is finite.
Passing to a derivative would not establish that missing finiteness.
This theorem therefore does not convert the hypothetical disjoint
supports into a contradiction.

\paragraph{Evidence and boundary.}
The four primary PDFs and their exact retrieval hashes are saved in
\path{evidence/lean_chain/primary/hadamard_scope/}.
The separate standard-library script
\path{saturated_rankin_cohen_control.py} checks all coefficients
in (HM.4)--(HM.5) for the single diagnostic prime $31$, including the
two conflicting coefficients in (HM.7).  This is a finite algebra
control, not a prime census or a proof by computation of (HM.1).
No simultaneous seed, original class-component vanishing, or new
finite $K$-group degree is inferred in this section.

% END evidence/lean_chain/hadamard_modularity_scope.tex


% BEGIN evidence/main_gate/subp_third_unit_counterexample.tex
\begin{proposition}[KT-SUBP-SEED-THIRD-UNIT-NONIMPLICATION]
Requiring a two-unit auxiliary discriminant to satisfy $|D|\leq p-4$
does not itself force the remaining third unit, even in the automatic
family $p\equiv3\pmod4$, $e=(p+1)/2$.
\end{proposition}
\begin{proof}
Take $p=83$, $e=42$, $b=e+1=43$, and $D=-43$.  Both $83$ and $43$
are prime, $D$ is odd and fundamental, and $|D|=43\leq79=p-4$.
The exact defining Bernoulli sums give
\[
 B_{1,\chi_{-43}}=-1,\qquad
 \beta_3(\chi_{-43})=166=2\cdot83,
 \qquad \beta_{43}(\chi_{-43})\equiv76\pmod{83}.
\]
Thus $h(-43)=1$ and the first twisted value is a unit, whereas the
third is not.  The full exact high-weight quotient is retained in the
certificate report.
The standalone certificate
\path{subp_third_unit_counterexample.py} computes these rational
Bernoulli-polynomial values and independently obtains $76$ from the
defining conductor-$83\cdot43$ Teichm\"uller-character sum modulo $83^2$.
The field also satisfies the already checked source conditions at
$N=2,\psi=1$: its discriminant is odd, $83\nmid43h(-43)$, $e$ is
positive and even, and $82\nmid2e=84$.  The remaining quadratic and
character-intersection clauses are the same elementary verification
as in KT-JACOBI-TWO-UNIT-GD.  This is a single-field nonimplication;
it does not rule out choosing another field with all three units.
\end{proof}

% END evidence/main_gate/subp_third_unit_counterexample.tex


% BEGIN evidence/main_gate/weight_two_kernel_discriminator.tex
\subsection{The weight-two residue kernel does not factor through the third value}

\begin{proposition}[KT-WEIGHT-TWO-KERNEL-NONFACTORIZATION]
Write $\Theta=\theta^4$, $G=2E_2(2\tau)-E_2(\tau)$,
$T=G/\Theta$, and
\[
 R=256\frac{2-T}{(T-1)^2},\qquad
 A=\begin{pmatrix}1&0\\2&1\end{pmatrix}.
\]
At $p=83$, the coefficient of $q^{10}$ in $R(A\tau)$ is a
$p$-adic unit, whereas the matching third value for
$d=p-4\cdot10=43$ is $\beta_3(\chi_{-43})=166\equiv0\pmod p$.
Thus the reversed principal-part weights of this particular residue
functional do not vanish on the third-value zero set.  In particular
they cannot be factored coefficientwise by that third value over
$\mathbf F_{83}$.
\end{proposition}

\begin{proof}
The transformations already established in
KT-COHEN-WEIGHT-TWO-BOUND give
\[
 G|_2A=G,\qquad
 \Theta|_2A=16q\left(\sum_{j\geq0}q^{j(j+1)}\right)^4.
\]
Writing $\Theta_A=\Theta|_2A$ therefore gives the exact expansion
\[
 R(A\tau)=256\frac{\Theta_A(2\Theta_A-G)}{(G-\Theta_A)^2}.
 \tag{WK.1}
\]
Its denominator has constant $1$.  Equivalently the standard theta
identities give
\[
 R(A\tau)=-256\lambda(2\tau)(1-\lambda(2\tau))
   =-4096q\frac{\prod_{n\geq1}(1-q^{2n})^{12}}{\theta(\tau)^{12}}.
 \tag{WK.2}
\]
The attached finite certificate expands (WK.1) and (WK.2) separately
through $q^{10}$, using integer arithmetic, and finds equal complete
coefficient vectors.  In particular
\[
 [q^{10}]R(A\tau)=166601228288\equiv26\pmod{83}.
\]
The exact value $\beta_3(\chi_{-43})=166$ is the independently checked
certificate KT-SUBP-SEED-THIRD-UNIT-NONIMPLICATION.  The phase in the
theta decomposition for $d=p-4j$ contributes only a common unit times
$(-1)^j$; here $j=10$, so it does not change the displayed residue.
Any factorization by the zero third value would give zero.  This
proves the proposition.
\end{proof}

This is a failure of the specified functional to supply the required
coefficientwise factorization.  It does not rule out another integral
functional or another auxiliary field.  The first-value coefficient at
this field is itself nonzero, since the earlier certificate gives
$\beta_{43}(\chi_{-43})\equiv76\pmod{83}$; no new search over primes or
discriminants is involved.

% END evidence/main_gate/weight_two_kernel_discriminator.tex


% BEGIN evidence/lean_chain/weighted_character_attempt.tex
\subsection{A bounded character-sum attempt at the short unit window}
\noindent\textbf{KT-WEIGHTED-CHARACTER-ATTEMPT.}\quad

This is a check of a proposed route to the missing simultaneous
unit, not a new existence theorem.  Fix $p\geq31$ prime with
$p\equiv3\pmod4$, put $b=(p+3)/2$, and let $D=-d$ be odd,
negative, fundamental and prime to $p$.  Write
$\chi=\chi_D$, $\epsilon(r)=(\frac rp)$, and
$P(j)=\sum_{0\leq a<j}\chi(a)$ for $0\leq j<d$.
The sufficient simultaneous-unit window already proved is
$d^5\leq1600p^2$.  Within it the third value $\beta_3(\chi)$
and the class number are units; the missing assertion is a
unit of $\beta_b(\chi)$ for at least one such $D$.

\paragraph{The exact integer representative.}
KT-FIELD-CHARACTER-REDUCTION and the complete prefix formula give
\[
 \beta_b(\chi)\equiv I_D
 :=\chi(p)\sum_{r=1}^{p-1}r\epsilon(r)
                         P([p^{-1}r]_d)\pmod p.
 \tag{WC.1}
\]
Indeed $b-1=(p-1)/2+1$, so Euler's criterion changes
$r^{b-1}$ into $r\epsilon(r)$ in that formula.
This is a signed integer representative of the residue;
it is not an equality of ordinary real values with
$\beta_b(\chi)$.

For clarity, assume $d<p$, as holds in the proposed short
window.  For $0\leq j<d$ let $a_j\in\{1,\ldots,d\}$
represent $pj$ modulo $d$, set
\[
 N_j=1+\left\lfloor\frac{p-1-a_j}{d}\right\rfloor,
 \qquad u_j\equiv a_jd^{-1}\pmod p,
 \quad 0\leq u_j<p,
\]
and put
\[
 S_j=\sum_{t=0}^{N_j-1}\epsilon(u_j+t),\qquad
 T_j=\sum_{t=0}^{N_j-1}t\epsilon(u_j+t).
\]
Partitioning $1\leq r<p$ according to $r\equiv pj\pmod d$
gives the exact identity of integers
\[
 I_D=\chi(p)\epsilon(d)
        \sum_{j=0}^{d-1}P(j)(a_jS_j+dT_j).
 \tag{WC.2}
\]
The shifts $u_j$ can be anywhere modulo $p$; the character
is periodic, so the intervals may pass through $p$.
Their lengths are about $p/d$.  Formula (WC.2) has
$d$ intervals, their prefix coefficients, and the factors
$a_j,d$.  A bound for one unweighted character sum is
not itself a bound below $p$ for $I_D$.

\paragraph{What the direct Burgess estimate actually gives.}
For a shifted interval of length $X$ and weights
$t+\alpha$, $0<\alpha\leq1$, summation by parts and
the classical Burgess estimate give, for every integer
$R\geq1$,
\[
 \left|\sum_{0\leq t<X}(t+\alpha)\epsilon(u+t)\right|
 \ll_{R,\varepsilon}
       X^{2-1/R}p^{(R+1)/(4R^2)+\varepsilon}.
 \tag{WC.3}
\]
Primality of $p$ satisfies the cube-free modulus condition
for arbitrary $R$.  The normalized expression for a block
in (WC.2) is of this form with $\alpha=a_j/d$; recovering
that block multiplies by $d$.

Even discard this last factor, all prefix coefficients,
and the sum over $j$, which is more favorable than the
actual expression.  At the edge $d\asymp p^{2/5}$ of
the desired window one has $X\asymp p^{3/5}$, and the
power of $p$ in this one-block envelope is
\[
 E_R=\frac65-\frac7{20R}+\frac1{4R^2}.
 \tag{WC.4}
\]
Its minimum over positive integer $R$ is
$E_2=87/80>1$.  In fact $E_1=11/10$, and $E_R$
is increasing for real $R\geq2$ since
$E_R'=(7R-10)/(20R^3)>0$ there.
Thus this estimate does not bound even the favorable
normalized block by $p$ in the proposed range.
The constants and logarithmic factors cannot change
that asymptotic power obstruction.

More generally, for $d\asymp p^\alpha$ the same envelope
has exponent below $1$ only when
\[
 \alpha>
 \frac{4R^2-3R+1}{4R(2R-1)}.
 \tag{WC.5}
\]
The smallest right side over positive integer $R$ is
$11/24$, at $R=2$.  For $R\geq2$ its difference from
$11/24$ is
$(2R-3)(R-2)/(24R(2R-1))\geq0$, while $R=1$
gives $1/2$.  The short-unit exponent $2/5$ is smaller.
These calculations delimit this direct absolute-value
argument; they are not lower bounds on the actual sums
and do not exclude additional cancellation between blocks.

\paragraph{Why the positive fixed-weight value is not a small substitute.}
The already checked fixed-weight identity is
\[
 \beta_b(\chi_D)\equiv\beta_2(\chi_F)\pmod p,
 \qquad F=\mathbf Q(\sqrt{pd}),\quad
 \chi_F=\chi_D\chi_{-p}.
 \tag{WC.6}
\]
Put $f=pd$.  The elementary Bernoulli Fourier expansion,
with the positive quadratic Gauss sum $\sqrt f$, gives
\[
 \beta_2(\chi_F)=\frac{f^{3/2}}{2\pi^2}L(2,\chi_F)>0.
 \tag{WC.7}
\]
Every Euler factor of this positive $L$-value is at least
$(1+\ell^{-2})^{-1}$, including primes dividing $f$.
Consequently
\[
 L(2,\chi_F)\geq\frac{\zeta(4)}{\zeta(2)}
                  =\frac{\pi^2}{15},\qquad
 \beta_2(\chi_F)\geq\frac{(pd)^{3/2}}{30}>p
       \quad(p\geq43,\ d\geq3).
 \tag{WC.8}
\]
The final inequality follows from $27p>900$.
So positivity of this real value cannot prove that it
is a nonzero number smaller than $p$.  Passing instead
to the smaller residue representative (WC.1) loses
that positivity and requires its own nonvanishing argument.

The existence of a small quadratic nonresidue likewise
does not by itself supply a nonzero weighted moment.
For example, modulo $31$ the first three Legendre signs
are $1,1,-1$, but $1+2-3=0$.  More globally the
ordinary class-number formula gives
$\sum_{r=1}^{p-1}r\epsilon(r)=-p\,h(-p)$ for
$p>3$, $p\equiv3\pmod4$: a natural nonzero weighted
sum can be divisible by $p$.  These examples concern
the proposed inference from character signs, not a
counterexample to the target simultaneous field existence.

\paragraph{Precise remaining requirement and source scope.}
One would need a nonzero residue in (WC.1) for some
fundamental $d$ in the short window, with a method that
retains the same $D$ for all units.  Separate real
nonvanishing, a bound for the least nonresidue, or an
upper bound for one block in (WC.2) does not prove this.
No such additional cancellation or lower-bound statement
is established in this attempt.

The analytic input is the classical Burgess statement
as stated in E.~Trevi\~no, \emph{The Burgess inequality
and the least $k$-th power non-residue}, primary author
PDF \url{https://campus.lakeforest.edu/trevino/Burgess.pdf},
pp.~1--2.  The full statement, including the cube-free
qualification, and his explicit Theorem~1 and Corollary~1
on p.~2 were read.  Only these estimate scopes are
imported; the foundational Burgess proof was not newly
reviewed.  The Fourier calculation and exponent
optimization above are displayed deductions.
There is no broad numerical scan, no new main-gate
theorem, and no claim that every character-sum method
must fail.

% END evidence/lean_chain/weighted_character_attempt.tex


% BEGIN evidence/main_gate/high_order_twist_source_scope.tex
\paragraph{KT-HIGH-ORDER-TWIST-TARGET-SCOPE.}
The high-order twists used in Heumann--Vatsal Theorem~5.1
cannot simply replace the quadratic auxiliary field in the present
argument. This is a scope statement about the displayed source
specialization, not an impossibility theorem for a different construction.

The exact declaration \texttt{UpperTheory.V.HV.GD} in
\path{HeightAndVanishing/Counts.lean}, lines~9--22, requires
the field to have degree two and be totally complex. Its condition
\texttt{ine} says that the \emph{component at the coefficient prime}
of the finite character has square one. It does not say that the
entire finite character is quadratic: the definition of
\texttt{comp} is in
\path{HeightAndVanishing/DirichletPrimeToLPartMulLPowerPart.lean},
lines~75--86. High-order characters away from that prime are
therefore not excluded just by this particular condition.

They do, however, change the target. In the latter file,
\texttt{Vlen} at line~116 is
\texttt{Chr.V'} evaluated at the Galois character
\texttt{Chr.dirG} of the actual finite character. The orbit in
lines~140--145 consists of characters with one fixed kernel.
A nontrivial character does not belong to the trivial character's
orbit. The rational untwisted target used here is the
specialization with finite character one.

There is also no direct reduction-to-one substitution using a
nontrivial character of order prime to $p$. Reduction is injective
on roots of unity of order prime to $p$: if $\zeta\ne1$ has
order $m$ prime to $p$ and reduces to one, then
$1+\zeta+\cdots+\zeta^{m-1}=0$ would reduce to $m=0$,
a contradiction. The condition \texttt{GD.ord} expressly
requires prime-to-$p$ order. Thus the high-order character
nonvanishing theorem supplies neither a quadratic field nor a
vanishing result for the original finite-character-one Selmer
group through these declarations alone. An additional comparison
would have to be proved.

% END evidence/main_gate/high_order_twist_source_scope.tex


% BEGIN evidence/main_gate/bernoulli_prefix_formula.tex
\paragraph{KT-BERNOULLI-PREFIX-FORMULA.}
Here is an exact finite-sum identity used only to make the numerical
checks below efficient.  Let $\chi$ be a nontrivial quadratic Dirichlet
character of conductor $f$, let $p\nmid f$ be an odd prime, and let
$\psi$ be a nontrivial character modulo $p$.  For $0\leq j<f$, put
$P(j)=\sum_{b=0}^{j-1}\chi(b)$, and let $[x]_f$ be the representative
in $[0,f-1]$ of $x$ modulo $f$.  Then
\[
 B_{1,\chi\psi}
   =\chi(p)\sum_{r=1}^{p-1}\psi(r)P([p^{-1}r]_f).          \tag{EU.10}
\]
For the characters $\psi=\omega^k$ with
$k\not\equiv0\pmod{p-1}$, reduction modulo $p$ consequently gives
\[
 B_{1,\chi\omega^k}
 \equiv\chi(p)\sum_{r=1}^{p-1}r^kP([p^{-1}r]_f)\pmod p.  \tag{EU.11}
\]
The exact identity in particular proves the needed $p$-integrality.

\begin{proof}
In the defining sum
$B_{1,\chi\psi}=(fp)^{-1}\sum_{a=1}^{fp}a\chi(a)\psi(a)$,
write each term with $p\nmid a$ uniquely as $a=r+ph$, where
$1\leq r\leq p-1$ and $0\leq h\leq f-1$.  The sum of
$\chi(r+ph)$ over $h$ vanishes, so
\[
 B_{1,\chi\psi}=\frac1f\sum_{r=1}^{p-1}\psi(r)
                       \sum_{h=0}^{f-1}h\chi(r+ph).
\]
Set $A_j=\sum_{h=0}^{f-1}h\chi(h+j)$.  Shifting the summation
cyclically shows $A_{j+1}=A_j+f\chi(j)$, since $\sum\chi=0$.
Thus $A_j=A_0+fP(j)$.  Multiplicativity gives
$\sum_hh\chi(r+ph)=\chi(p)A_{[p^{-1}r]_f}$, while
$\sum_r\psi(r)=0$ cancels the $A_0$ term.  This proves (EU.10).
Reducing $\omega^k(r)$ to $r^k$ proves (EU.11).
\end{proof}

The formula is an elementary rearrangement of the defining Bernoulli
sum; no originality claim is made.  Its role is to reduce each test
from $O(fp)$ summands to $O(f+p)$.  In the two-unit family the weights
are $r^{(p-3)/2}=(r/p)r^{-1}$ and $r^{p-3}=r^{-2}$ modulo $p$.
An independent direct sum modulo $p^2$ and exact rational Bernoulli
polynomials check the normalization in the saved certificate.

\paragraph{KT-EXPLICIT-TWO-UNIT-FINITE-CHECK (bounded exact computation).}
The saved \path{explicit_two_unit_certificate.py} evaluates
(EU.11) for both signs and for every prime
$7\leq\ell\leq20000$ with $\ell\equiv3\pmod4$.
All $1135$ primes, hence $2270$ fields counted with their sign labels,
have a nonzero third residue.  Together with
KT-EXPLICIT-TWO-UNIT-FAMILY and KT-FIELD-CHARACTER-REDUCTION, these
exact values give all three units for both fields in this finite
prime range and every weight in (EU.2).
The complete rows are saved in the corresponding JSON file, with
row hash
\path{8328795e4544c175eaa37953962b681aee08ecb2b7aa9cf787aea3660994a472}.
There are $72$ independent comparisons with the defining conductor
$f\ell$ sum modulo $\ell^2$, and $12$ comparisons with exact rational
Bernoulli polynomials.  The finite computation supplies no conclusion
for a prime outside its stated range.

% END evidence/main_gate/bernoulli_prefix_formula.tex


% BEGIN evidence/main_gate/third_unit_character_obstruction.tex
\subsection{The remaining character and the limit of a Pell-unit argument}
\label{sec:kt-third-unit-character}

This subsection tests a specific proposed mechanism for the third unit in
KT-EXPLICIT-TWO-UNIT-FAMILY.  It identifies the character that must be
controlled, proves that units from multiquadratic fields have zero
projection to that character, and shows that requiring the entire
cyclotomic compositum to be $p$-rational is too strong.  None of these
statements asserts failure of the third unit itself.

\begin{lemma}[KT-PADIC-BRANCH-RESIDUE]
Let $p>2$ be prime and let $\psi$ be a nontrivial even Dirichlet
character whose order divides $p-1$, with values in $\mathbf Z_p$.
For all $s,t\in\mathbf Z_p$,
\[
 L_p(s,\psi)\in\mathbf Z_p,\qquad
 L_p(s,\psi)\equiv L_p(t,\psi)\pmod p.                 \tag{TU.1}
\]
There is no restriction that the conductor be prime to $p$.
\end{lemma}
\begin{proof}
Write $f$ for the conductor.  Choose an integer $c>1$, prime to $pf$,
with $\psi(c)\ne1$.  Such a residue class exists because $\psi$ is
nontrivial; imposing that it be prime to $p$ adds no obstruction.
Since the order of $\psi$ divides $p-1$, reduction is injective on its
values, so $1-\psi(c)$ is a unit.

Knospe--Washington, Corollary~2.5 on author-PDF p.~5, gives
\[
 -(1-\psi(c)\langle c\rangle^{1-s})L_p(s,\psi)
 =\lim_{n\to\infty}\sum_{\substack{1\le a\le fp^n\\p\nmid a}}
     \psi\omega^{-1}(a)\langle a\rangle^{-s}
     \epsilon_{a,c,fp^n},                                  \tag{TU.2}
\]
where
$\epsilon_{a,c,F}=(c-1)/2-\{-aF^{-1}\}_c\in\tfrac12\mathbf Z$.
Every summand is integral.  For $s,t\in\mathbf Z_p$,
$\langle a\rangle^{-s}\equiv\langle a\rangle^{-t}\equiv1\pmod p$.
Thus each finite sum is integral and its difference between $s$ and
$t$ is in $p\mathbf Z_p$.  Taking the two convergent limits preserves
these properties.  The factor on the left is a unit, congruent to
$1-\psi(c)$ independently of $s$.  Division proves (TU.1).
The cited domain $|s|<p^{1-1/(p-1)}$ contains $\mathbf Z_p$.
\end{proof}

\begin{theorem}[KT-THIRD-UNIT-CHARACTER-OBSTRUCTION]
Let $p>3$ be prime, let $D<0$ be a fundamental discriminant prime
to $p$, and set
\[
 L=\mathbf Q(\sqrt D,\mu_p),\qquad K=L^+,
 \qquad\theta=\chi_D\omega^{-1}.                          \tag{TU.3}
\]
Then $K/\mathbf Q$ is cyclic of degree $p-1$, and $\theta$ is a
faithful even character of its Galois group.  Moreover,
\[
 L_p(1,\theta)\equiv-B_{1,\chi_D\omega^{-2}}
       \equiv-\frac{B_{p-2,\chi_D}}{p-2}\pmod p.           \tag{TU.4}
\]
All terms are integral.  Consequently the remaining third unit of
KT-EXPLICIT-TWO-UNIT-FAMILY is equivalent to
$L_p(1,\theta)\in\mathbf Z_p^\times$.

Let $M/\mathbf Q$ be any finite multiquadratic extension and
$N=KM$.  Inflate $\theta$ to $G=\operatorname{Gal}(N/\mathbf Q)$.
With
\[
 e_\theta=\frac1{|G|}\sum_{g\in G}\theta(g)^{-1}g
          \in\mathbf Z_p[G],                              \tag{TU.5}
\]
the image of $\mathcal O_M^\times\otimes\mathbf Z_p$ in
$\mathcal O_N^\times\otimes\mathbf Z_p$ is killed by
$e_\theta$.  Thus this is also true of any collection of Pell
units from quadratic subfields, their Galois conjugates, and the
$\mathbf Z_p$-module they generate.  In particular, for $D=D_s$
in KT-EXPLICIT-TWO-UNIT-FAMILY, the explicit units used to
establish $p$-rationality of
$\mathbf Q(\sqrt{p(p+2s)})$ do not supply a nonzero vector in
the component controlling (TU.4).
\end{theorem}
\begin{proof}
The fields $\mathbf Q(\sqrt D)$ and $\mathbf Q(\mu_p)$ are
linearly disjoint: the unique quadratic subfield of the latter
is ramified at $p$, whereas the former is not.  Hence
\[
 \operatorname{Gal}(L/\mathbf Q)
       \simeq\{\pm1\}\times\mathbf F_p^\times.
\]
Under the Teichm\"uller identification, the character
$(\varepsilon,z)\mapsto\varepsilon\omega(z)^{-1}$ has kernel
$\{(1,1),(-1,-1)\}$.  The second element is complex
conjugation.  Its quotient is cyclic of order $p-1$, proving
the assertions about $K$ and $\theta$.

The character $\chi_D\omega^{-2}$ is primitive of conductor
$|D|p$, since $p>3$.  Interpolation at zero therefore gives the
exact equality
\[
                     L_p(0,\theta)=-B_{1,\chi_D\omega^{-2}}.
                                                               \tag{TU.6}
\]
There is no omitted $p$ Euler factor, since this character is
ramified at $p$.  For completeness, (TU.6) also follows directly
by substituting $s=0$ in formula~(1) on p.~2 of
Knospe--Washington.  The Bernoulli series then stops at $j=1$;
its first term is the negative defining $B_1$ sum, and its
second is one half the sum of the nontrivial character
$\chi_D\omega^{-2}$, hence zero.
Apply KT-PADIC-BRANCH-RESIDUE to $\theta$ and then the already
proved KT-FIELD-CHARACTER-REDUCTION at $n=p-2$ to obtain (TU.4).

For the projection assertion, every subextension of a
multiquadratic extension is multiquadratic.  Since $K/\mathbf Q$
is cyclic, $[K\cap M:\mathbf Q]\le2$.  Therefore
\[
 H=\operatorname{Gal}(N/M)\simeq
       \operatorname{Gal}(K/(K\cap M))
\]
has order at least $(p-1)/2>1$.  Faithfulness of $\theta$ on
$K$ implies that its restriction to $H$ is nontrivial.  Choose
$h\in H$ with $\theta(h)\ne1$.  If $u$ is the image of an
element of $\mathcal O_M^\times\otimes\mathbf Z_p$, then
$hu=u$, whereas $h(e_\theta u)=\theta(h)e_\theta u$.
The difference $\theta(h)-1$ is a $p$-adic unit, so
$e_\theta u=0$.  Also $p\nmid|G|$, so (TU.5) is integral.
The projection commutes with Galois action and all
$\mathbf Z_p$-linear operations.  This proves the claimed
scope, without an assertion about other, higher arithmetic
information that could conceivably be extracted from the same
fields.
\end{proof}

\begin{proposition}[KT-COMPOSITUM-PRATIONAL-OBSTRUCTION]
Suppose that $p>3$ is irregular, in the precise sense that for
some even $k$ with $2\le k\le p-3$ one has
$p\mid\operatorname{num}(B_k)$.  For every $D<0$ fundamental
and prime to $p$, the field
$K=(\mathbf Q(\sqrt D,\mu_p))^+$ is not $p$-rational.
Nevertheless, at $p=59$ both fields from
KT-EXPLICIT-TWO-UNIT-FAMILY satisfy the third-unit condition.
Thus full $p$-rationality of this compositum is not a necessary
condition for that third unit.
\end{proposition}
\begin{proof}
The even nontrivial character $\omega^k$ is a character of
$\mathbf Q(\mu_p)^+\subset K$.  Interpolation and (TU.1) give
\[
 L_p(1,\omega^k)\equiv L_p(1-k,\omega^k)
       =-(1-p^{k-1})\frac{B_k}{k}\equiv0\pmod p.          \tag{TU.7}
\]
Here $k$ is a unit and the ordinary Bernoulli denominator is
prime to $p$; the latter also follows from its recurrence up
to index $p-3$.  Every other nontrivial character of $K$ has
order dividing $p-1$, so its $L_p(1)$ value is integral by
KT-PADIC-BRANCH-RESIDUE.

Equation~(2.4) on pp.~4--5 of Benmerieme--Movahhedi states
that a real abelian field with character group $X$ is
$p$-rational precisely when
\[
             \prod_{1\ne\psi\in X}L_p(1,\psi)
                     \sim p^{1-m},                       \tag{TU.8}
\]
where $|\mu_{p^\infty}\cap K(\mu_p)|=p^m$ and $\sim$
means equality of valuations.  The abelian case satisfies
Leopoldt, as recalled there.  In the present case
$K(\mu_p)=L$ has degree $2(p-1)$ and contains $\mu_p$ but
not $\mu_{p^2}$, so $m=1$.  The right side of (TU.8) is a
unit, whereas the left side is integral with the nonunit
factor (TU.7).  This proves the obstruction.

For the concrete comparison, exact rational arithmetic gives
\[
 \frac{B_{44}}{44}
       =-\frac{2530297234481911294093}{2760},\qquad
       v_{59}(B_{44}/44)=1.                                \tag{TU.9}
\]
At $p=59$, the two discriminants are $D_+=-244$ and
$D_-=-228$.  The defining generalized Bernoulli sums, evaluated
modulo $59^2$ before dividing by the conductor, give
\[
 B_{1,\chi_{-244}\omega^{-2}}\equiv13,\qquad
 B_{1,\chi_{-228}\omega^{-2}}\equiv5\pmod{59}.             \tag{TU.10}
\]
The complete small certificate
\path{third_unit_59_certificate.py} computes (TU.9)
from the Bernoulli recurrence and (TU.10) directly from the
conductor sums; it does not use the prefix formula.  The
same residues occur in the independently reviewed bounded
certificate for KT-EXPLICIT-TWO-UNIT-FINITE-CHECK.  Therefore
both third units hold, although neither corresponding $K$
is $59$-rational.
\end{proof}

\paragraph{Primary input and status.}
H.~Knospe and L.~C.~Washington, \emph{Dirichlet series
expansions of $p$-adic $L$-functions}, Abhandlungen aus dem
Mathematischen Seminar der Universit\"at Hamburg 91 (2021),
325--334, DOI \path{10.1007/s12188-021-00244-0}; author
version arXiv:2102.02851v3, 9 September 2021,
\url{https://arxiv.org/pdf/2102.02851}:
formula~(1) on p.~2, Proposition~2.1 and its full proof on
pp.~2--3, Theorem~2.4 and its full proof on pp.~3--5, and
Corollary~2.5 on p.~5.  These items were read in full for
the present argument.  The analytic formula imported in
that paper is Washington's Theorem~5.11; its original book
proof is not claimed freshly checked here.
For (TU.8), use the same Benmerieme--Movahhedi version and
precise pp.~4--5 locators already checked for
KT-EXPLICIT-TWO-UNIT-FAMILY.  This subsection is an
obstruction to two proposed proof mechanisms, not a new
all-prime nonvanishing or Selmer-vanishing assertion.

% END evidence/main_gate/third_unit_character_obstruction.tex


% BEGIN evidence/main_gate/auxiliary_unit_selmer_bridge.tex
\subsection{What the three auxiliary units actually annihilate}
\label{sec:kt-auxiliary-unit-bridge}

\begin{theorem}[KT-AUXILIARY-UNIT-SELMER-BRIDGE]
Let $p>3$ be prime, let $D<0$ be a fundamental discriminant
prime to $p$, and put
\[
 E=\mathbf Q(\sqrt D),\quad L=E(\mu_p),\quad
 \chi=\chi_D,\quad A=\operatorname{Cl}(L)\{p\}.
                                                               \tag{AB.1}
\]
Let $a\ge0$ be even, with either $a=0$ or
$a\not\equiv0\pmod{p-1}$.  If
\[
                         B_{1,\chi\omega^a}\in\mathbf Z_p^\times,
                                                               \tag{AB.2}
\]
then
\[
             A^{\chi\omega^{-a}}=0,
             \qquad A^{\chi\omega^{a+1}}=0.               \tag{AB.3}
\]
Here a superscript denotes the eigenspace for the indicated
Galois character, with the fixed $p$-adic Teichm\"uller embedding.

Write $\varepsilon_p$ for the $p$-adic cyclotomic character and
$T_a=\mathbf Q_p/\mathbf Z_p(\chi\varepsilon_p^{-a})$.
Define the strict group by its actual local conditions:
\[
 \mathcal S_p(T_a)=\ker\left(H^1(\mathbf Q,T_a)\longrightarrow
 H^1(\mathbf Q_p,T_a)\oplus
       \prod_{q\ne p}H^1(I_q,T_a)\right),                 \tag{AB.4}
\]
where $q$ ranges over finite rational primes.  Then
\[
                            \mathcal S_p(T_a)=0.          \tag{AB.5}
\]
The same assertion holds with $T_a$ replaced by
$\mathbf F_p(\chi\omega^{-a})$.

For $a=0,e,2e$ under $e\ge2$ even and
$(p-1)\nmid2e$, the three units therefore give the following
exact character correspondence:
\[
\begin{array}{c|c|c|c}
\text{unit}&\text{odd class component}&\text{even mirror}
  &\text{strict coefficient module}\\ \hline
B_{1,\chi}&\chi&\chi\omega&\mathbf Q_p/\mathbf Z_p(\chi)\\
B_{1,\chi\omega^e}&\chi\omega^{-e}&\chi\omega^{e+1}
 &\mathbf Q_p/\mathbf Z_p(\chi\varepsilon_p^{-e})\\
B_{1,\chi\omega^{2e}}&\chi\omega^{-2e}&\chi\omega^{2e+1}
 &\mathbf Q_p/\mathbf Z_p(\chi\varepsilon_p^{-2e})
\end{array}                                                   \tag{AB.6}
\]
Every component in the second and third columns is negative
under $\operatorname{Gal}(L/\mathbf Q(\mu_p))$.
In contrast, the original rational component $\omega^{-e}$
is positive under that group.  Thus (AB.3)--(AB.6) do not,
by themselves, assert vanishing of the original rational
strict group with coefficient $\mathbf Q_p/\mathbf Z_p(-e)$.
\end{theorem}

\begin{proof}
Set $f=|D|$ and $G=\operatorname{Gal}(L/\mathbf Q)$.
The extensions $E$ and $\mathbf Q(\mu_p)$ are disjoint,
so $|G|=2(p-1)$ is prime to $p$.  All character values
under consideration lie in $\mathbf Z_p$, and integral
character projections are available.

First suppose $a\not\equiv0\pmod{p-1}$.  Then
$\eta=\chi\omega^a$ is primitive of conductor $fp$.
Use the Stickelberger element for the abelian field $L$,
with $\sigma_b$ induced by $\zeta_{fp}\mapsto\zeta_{fp}^b$:
\[
 \Theta_L=\frac1{fp}\sum_{\substack{1\le b\le fp\\(b,fp)=1}}
                     b\sigma_b^{-1}\in\mathbf Q[G].       \tag{AB.7}
\]
For every positive integer $c$ prime to $fp$,
$(c-\sigma_c)\Theta_L$ is integral and annihilates
$\operatorname{Cl}(L)$ by Stickelberger's theorem for an
abelian subfield, not merely for a full cyclotomic field.
Integrality can also be seen before restriction to $L$:
the coefficient at $\sigma_b^{-1}$ is
$(cb-[cb]_{fp})/(fp)\in\mathbf Z$.

On the $\eta^{-1}$ eigenspace, this annihilator acts by
\[
                  (c-\eta^{-1}(c))B_{1,\eta}.             \tag{AB.8}
\]
Indeed evaluation of $\sigma_b^{-1}$ at $\eta^{-1}$
gives $\eta(b)$, so the normalization and inverse in
(AB.8) follow directly from the defining Bernoulli sum.
Choose $c\equiv1\pmod p$ and $\chi(c)=-1$, with
$(c,f)=1$, by the Chinese remainder theorem.  Then
$c-\eta^{-1}(c)\equiv2\pmod p$.  Under (AB.2), (AB.8)
is a unit annihilator; hence $A^{\chi\omega^{-a}}=0$.
There is no assumption $p\nmid\varphi(f)$.

For $a=0$, the imaginary quadratic class-number formula
shows that $p\nmid h(E)$.  With
$H=\operatorname{Gal}(L/E)$, extension after norm is
$\sum_{h\in H}h$.  It acts as the unit $p-1$ on $A^H$,
while norm has target $\operatorname{Cl}(E)\{p\}=0$.
Thus $A^H=0$, proving $A^\chi=0$.  This treatment avoids
inserting the imprimitive Euler factor $1-\chi(p)$, which
would vanish when $p$ splits in $E$.

For the mirror assertion, write
$\alpha=\chi\omega^{a+1}$ and $\xi=\chi\omega^{-a}$.
Then $\alpha$ is even and $\xi=\omega\alpha^{-1}$ is
odd and different from $\omega$.  Apply the full
Kummer-radical argument of the carried K17-REFLECTOR to
$L$.  More explicitly, the quotient
$A^\alpha/pA^\alpha$ gives an everywhere unramified
elementary abelian $p$-extension of $L$, stable under $G$.
Its Kummer radical $V\subset L^\times/L^{\times p}$ has
character $\xi$ by the equivariant pairing into $\mu_p$.
Unramifiedness implies that $(x)=\mathfrak a^p$ for every
representative $x$ of $V$.  Sending $x$ to
$[\mathfrak a]\in A[p]$ is well defined and $G$-equivariant;
its kernel lies in the unit group modulo $p$th powers.
For a CM field, that unit group's odd part consists only
of roots of unity: $u/\overline u$ is a root of unity,
and $u^2=(u/\overline u)(u\overline u)$.
Here the $p$-primary roots of unity are exactly $\mu_p$
because $[L:\mathbf Q]=2(p-1)$ is prime to $p$.
Their character is $\omega$, different from $\xi$.
Consequently $V\hookrightarrow A[p]^\xi$ is injective.
The first part makes its target zero, so
$A^\alpha/pA^\alpha=0$.  Finiteness gives $A^\alpha=0$.

Put $W_a=T_a[p]=\mathbf F_p(\chi\omega^{-a})$.
Since $p\nmid|G|$, restriction identifies global
$H^1(\mathbf Q,W_a)$ with the $G$-invariant part of
$H^1(L,W_a)$.  A class in the residual strict group
restricts to an everywhere unramified homomorphism on
$G_L$: away from $p$ its inertia restriction is zero,
and at $p$ its entire local class is zero.  Global class
field theory therefore embeds that strict group into
\[
                 \operatorname{Hom}_G(A/pA,W_a).          \tag{AB.9}
\]
This is zero because the first component in (AB.3) is zero.
The equivariance convention is explicit: a homomorphism
$f$ in (AB.9) obeys $f(gx)=\chi(g)\omega(g)^{-a}f(x)$;
it sees $A^{\chi\omega^{-a}}$, not its inverse.

It remains to pass from residual to divisible coefficients
without changing the local conditions.  Complex conjugation
acts as $-1$ on $T_a$, so $H^0(\mathbf Q,T_a)=0$ and
\[
              H^1(\mathbf Q,W_a)\simeq
                       H^1(\mathbf Q,T_a)[p].            \tag{AB.10}
\]
For $q\ne p$, the cyclotomic character is unramified, and
inertia acts through the order-at-most-two character $\chi$.
Hence $T_a^{I_q}$ is either $0$ or $T_a$, and is
$p$-divisible.  At $p$, if $a\not\equiv0\pmod{p-1}$,
some inertia element acts nontrivially modulo $p$, so
$H^0(\mathbf Q_p,T_a)=0$.  If $a=0$, this group is
either $0$ or $T_a$, according as $\chi(p)=-1$ or $1$;
it is again $p$-divisible.  The long exact sequences for
$0\to W_a\to T_a\xrightarrow{p}T_a\to0$ consequently
make every residual local $H^1$ map in (AB.4) injective.
Thus (AB.10) identifies the $p$-torsion in
$\mathcal S_p(T_a)$ with the residual strict group,
which is zero by (AB.9).  A nonzero $p$-primary torsion
group has a nonzero element of order $p$, proving (AB.5).

Finally let $\tau$ be the nontrivial automorphism of
$L/\mathbf Q(\mu_p)$.  It has $\chi(\tau)=-1$ and
$\omega(\tau)=1$.  This proves every sign assertion
in (AB.6) and the stated distinction from the rational
character $\omega^{-e}$.
\end{proof}

\paragraph{Imported theorem and exact proof scope.}
The annihilation input is L.~C.~Washington,
\emph{Introduction to Cyclotomic Fields}, first edition,
Springer (1982), \S6.2: definition on p.~93,
Lemma~6.9 and Theorem~6.10 on p.~94, with descent from
the full cyclotomic field to the arbitrary abelian field
on pp.~99--100.  The accessible book OCR was read for
the definition, lemma, theorem, and concluding descent;
the preceding Gauss-sum factorization proof is imported,
not claimed freshly checked.  The hosting page labels
the upload as the 1997 edition, but its actual title
and copyright pages show the 1982 first edition.
The publisher's second-edition chapter is
\url{https://doi.org/10.1007/978-1-4612-1934-7_6};
its abstract independently states the arbitrary-abelian-field
scope, while its paywalled chapter text was not retrieved.
The mirror proof above is a fully spelled-out adaptation
of the already checked K17-REFLECTOR, carried
\path{K17.tex}, lines 4533--4612; that argument cites
Sharifi's Theorem~4.4.2 and Corollary~4.4.4, pp.~111--113.
No new claim identifies (AB.4) with a different deformation
problem's local conditions.  Such an identification must
be checked against that problem's actual local predicates.

\begin{lemma}[KT-RATIONAL-POSITIVE-ROOT-UNIT]
Let $p>3$, let $b\ge2$ be even, and suppose
$(p-1)\nmid b$.  If $B_b/b$ is a $p$-adic unit, then
\[
              H^1(G_{\mathbf Q,\{p\}},
                         \mathbf F_p(\omega^b))=0.       \tag{AB.11}
\]
Here $G_{\mathbf Q,\{p\}}$ denotes the Galois group of the
maximal extension unramified at finite primes other than
$p$; no strict condition at $p$ is imposed in (AB.11).
\end{lemma}
\begin{proof}
Let $b_0\in\{2,4,\ldots,p-3\}$ be congruent to $b$
modulo $p-1$.  The ordinary Kummer congruence, equivalently
Sun's Lemma~8.1(a) with trivial character and difference
order one, gives $B_b/b\equiv B_{b_0}/b_0\pmod p$;
all intervening Euler factors are congruent to one.
Herbrand's implication in the carried K17-REFLECTOR proof
therefore gives
\[
 \operatorname{Cl}(\mathbf Q(\mu_p))\{p\}^{\omega^{1-b}}=0.
                                                               \tag{AB.12}
\]
Set $F=\mathbf Q(\mu_p)$, $\Delta=\operatorname{Gal}(F/\mathbf Q)$,
and $S$ equal to its unique prime over $p$.  The $S$-unit
Kummer exact sequence is
\[
 0\longrightarrow\mathcal O_F[1/p]^\times/
                    \mathcal O_F[1/p]^{\times p}
 \longrightarrow H^1(G_{F,S},\mu_p)
 \longrightarrow\operatorname{Cl}(\mathcal O_F[1/p])[p]
 \longrightarrow0.                                      \tag{AB.13}
\]
The prime in $S$ is principal, so the last group equals
$\operatorname{Cl}(F)[p]$.  The only odd character in
ordinary units modulo $p$th powers is $\omega$.
Inverting the unique prime adds only a trivial-character
quotient, since its valuation is $\Delta$-invariant.
Thus the $\omega^{1-b}$ component of the first term is
zero: this character is odd and differs from $\omega$.
The corresponding component of the last term is zero
by (AB.12).  It follows that the middle component is zero.
Finally twist (AB.13) by $\omega^{b-1}$ and take exact
$\Delta$-invariants.  Restriction identifies the resulting
middle term with (AB.11), because $|\Delta|=p-1$ is a unit.
This proves the claim at every positive weight $b$ in the
stated residue class.
\end{proof}

This additional rational statement removes the positive
coefficient $\omega^b$ under its stated Bernoulli-unit
condition.  It does not remove $\omega^{-b}$ without
the corresponding reflected Bernoulli input.  Whether
these groups govern every arrow in a proposed deformation
graph remains a separate question.

% END evidence/main_gate/auxiliary_unit_selmer_bridge.tex


% BEGIN evidence/main_gate/small_character_order_reduction.tex
\section{Removing the small-order rational characters}
\label{sec:kt-small-character-order}

The low-degree real class-number theorems give a direct restriction
on the rational components still requiring the unit-trace method.
The class-number theorems in this section are known inputs, not new
results. The short argument records their precise implication for
the character and integral-cohomology conventions of this report.

\begin{proposition}[KT-SMALL-ORDER-COMPONENT-REDUCTION]
Let $p>3$ be prime and let $e\ge2$ be even. If the order of the
Teichm\"uller character $\omega^e$ is in $\{1,2,3,4\}$, then
\[
                 A_p^{[-e]}=0,
                 \qquad K_{2e}(\mathbf Z)\{p\}=0.       \tag{LO.1}
\]
Consequently a remaining rational component in the all-weight
problem has $\operatorname{ord}(\omega^e)\ge5$.
No bound on $e$ relative to $p$ is needed.
\end{proposition}

\begin{proof}
Put $L=\mathbf Q(\mu_p)$ and let $d$ be the order of
$\omega^{-e}$. Let $F=L^{\ker(\omega^{-e})}$, so
$[F:\mathbf Q]=d$. Since $e$ is even, $F$ is totally real.
We first recall the exact class-number inputs for these four
possible degrees.

For $d=1$, $F=\mathbf Q$ has class number one. For $d=2$,
$p\equiv1\pmod4$ and $F=\mathbf Q(\sqrt p)$; the real
quadratic bound in Le, Theorem~(a), printed p.~141, already
checked in KT-ODD-DISCRIMINANT-TWO-UNIT-FAMILY, gives
$h(F)\le\lfloor\sqrt p/2\rfloor<p$.

For $d=3$, $F$ is the unique cyclic cubic field of conductor
$p$. Moser--Payan, \emph{J. Math. Soc. Japan}~33 (1981),
701--706, \S4, unnumbered Theorem and Corollary on p.~705,
gives $h(F)<2p/3$, hence $p\nmid h(F)$.

For $d=4$, evenness of the character implies
$p\equiv1\pmod8$. Put $k=\mathbf Q(\sqrt p)\subset F$.
Moser, Grenoble seminar 1979--1980, expos\'e~3,
Theorem~B on printed p.~III.2, gives
\[
            h(k)\le\sqrt p/2,\qquad h(F)/h(k)\le p/4.  \tag{LO.2}
\]
The same expos\'e, Theorem~1(2) on p.~III.1, gives
$h(k)\mid h(F)$: the cyclic quadratic extension $F/k$
has exactly one ramified prime. Thus both $h(k)$ and
$h(F)/h(k)$ are positive integers less than $p$, and
$p\nmid h(F)$. These are applications of the cited known
theorems; their analytic proofs are not being reproved here.

In every case $\operatorname{Cl}(F)\{p\}=0$.
Write $H=\operatorname{Gal}(L/F)$, whose order is prime to
$p$. Let $j$ and $N$ be extension and norm on ideal class
groups. They satisfy
\[
                      jN(x)=\sum_{\sigma\in H}\sigma x.
                                                               \tag{LO.3}
\]
If $x\in A_p^H$, then $N(x)$ is a $p$-primary class of $F$,
so it is zero. Formula (LO.3) gives $|H|x=0$.
Multiplication by $|H|$ is invertible on the finite $p$-group
$A_p$, hence $A_p^H=0$. The character $\omega^{-e}$ is
trivial on $H$, so $A_p^{[-e]}\subseteq A_p^H=0$.
The carried comparison K17-(K2), together with its periodic
character convention, gives the second assertion of (LO.1).
\end{proof}

\paragraph{Scope and source boundary.}
The two primary sources are
\url{https://doi.org/10.2969/jmsj/03340701} and
\url{https://www.numdam.org/item/STNG_1979-1980__8__A3_0/}.
The cubic numerical bound used is the published $2p/3$ bound;
no stronger bound is needed. The exact quartic statements and
their hypotheses were checked against the scanned theorem page.
The supplied K17 comparison is cited as settled input.
This restriction removes the specified character orders from
the rational $H^2_{\acute et}(\operatorname{Spec}\mathbf Z[1/p],
\mathbf Z_p(e+1))$ problem in (FW.1). It neither computes every
class component of higher order nor proves the all-weight
vanishing theorem. In particular no degree-six assertion is
imported from a title or a secondary bibliographic description.

% END evidence/main_gate/small_character_order_reduction.tex


% BEGIN evidence/main_gate/sextic_class_product.tex
\subsection{The integral class-number ratio in prime-conductor degree six}

\begin{lemma}[KT-SEXTIC-CLASS-PRODUCT]
Let $p\equiv1\pmod{12}$ be prime, let $F$ be the real cyclic sextic
subfield of $\mathbf Q(\mu_p)$, and let $F_2,F_3$ be its quadratic
and cubic subfields.  For ordinary ideal class numbers one has
\[
                         h(F_2)h(F_3)\mid h(F).
 \tag{SC.1}
\]
Thus $h(F)/(h(F_2)h(F_3))$ is a positive integer in this setting.
\end{lemma}

\begin{proof}
For $i=2,3$, let $H_i$ be the ordinary Hilbert class field of $F_i$:
it is unramified at every finite place, every real place splits, and
$\operatorname{Gal}(H_i/F_i)$ is the ordinary ideal class group.
The field $H_i$ is Galois over $\mathbf Q$, since $F_i/\mathbf Q$ is
Galois and its defining maximal unramified abelian extension is
preserved by every automorphism of $F_i$.

The unique prime over $p$ is totally ramified in every nontrivial
subextension of $F/F_i$.  Since $H_i/F_i$ is unramified there,
$H_i\cap F=F_i$.  Consequently
\[
 K_i=FH_i,\qquad [K_i:F]=h(F_i),
 \qquad \operatorname{Gal}(K_i/F)\simeq\operatorname{Gal}(H_i/F_i).
 \tag{SC.2}
\]
These are abelian extensions of $F$, unramified at all finite places
and split at its real places, and are Galois over $\mathbf Q$.

Put $E=K_2\cap K_3$ and $A=\operatorname{Gal}(E/F)$.
The extension $E/\mathbf Q$ is Galois, and $A$ is abelian.  Therefore
conjugation defines an action of
$G=\operatorname{Gal}(F/\mathbf Q)\simeq C_6$ on $A$, independent
of the choice of lifts to $\operatorname{Gal}(E/\mathbf Q)$.
This action factors through both $\operatorname{Gal}(F_2/\mathbf Q)$
and $\operatorname{Gal}(F_3/\mathbf Q)$, as follows.
If $g\in\operatorname{Gal}(F/F_i)$, the disjointness
$F\cap H_i=F_i$ gives a lift to $K_i$ acting as $g$ on $F$ and as
the identity on $H_i$.  That lift centralizes
$\operatorname{Gal}(K_i/F)$.  Its action on the quotient $A$ is
therefore trivial.  The subgroups $\operatorname{Gal}(F/F_2)$ and
$\operatorname{Gal}(F/F_3)$ have orders three and two and generate
$G$, so its action on $A$ is trivial.

Thus $A$ is central in $\Gamma=\operatorname{Gal}(E/\mathbf Q)$,
and $\Gamma/A$ is cyclic.  The group $\Gamma$ is abelian: choose
a lift $x$ of a generator of the cyclic quotient, and write every
element as $ax^j$ with $a\in A$.
The only possible finite ramified prime of $E/\mathbf Q$ is $p$.
Its inertia subgroup $I$ has order exactly six, since $E/F$ is
unramified and $F/\mathbf Q$ is totally ramified of degree six at $p$.
Because $\Gamma$ is abelian, $I$ is normal and the field $E^I$ is
unramified at every finite rational prime.  A number field with that
property has absolute discriminant one and hence is $\mathbf Q$
by Minkowski's discriminant bound.  It follows that $\Gamma=I$ and
$[E:\mathbf Q]=6$, so $E=F$.

The two Galois extensions $K_2/F$ and $K_3/F$ are consequently
linearly disjoint.  Their compositum is an abelian extension of $F$
unramified at all finite places and split at every real place, of
degree $h(F_2)h(F_3)$.  It is contained in the ordinary Hilbert
class field of $F$.  This proves (SC.1).
\end{proof}

\begin{corollary}[KT-SEXTIC-REGULATOR-RATIO-FOUR]
Use the same fields.  Let $R_2,R_3,R_6$ be their regulators, and let
$R_*$ be the primitive relative regulator in Louboutin's normalization:
for a generating relative unit with absolute logarithms
$(a,a+c,c,-a,-a-c,-c)$ it is $R_*=a^2+ac+c^2$.
Let $\chi_6$ be either primitive even Dirichlet character of order six
defining $F$, and put
\[
                       C_6=\frac{p|L(1,\chi_6)|^2}{R_*}.
 \tag{SC.3}
\]
Then $C_6$ is a positive integer divisible by four.  More precisely,
if $Q=[U_F:U_{F_2}U_{F_3}U_F^*]$, where relative units are defined
by norms lying in $\{\pm1\}$, then
\[
 Q\mid12,\qquad
 C_6=\frac{48}{Q}\frac{h(F)}{h(F_2)h(F_3)}\in4\mathbf Z_{>0}.
 \tag{SC.4}
\]
There is no hypothesis on signatures of the relative units.
\end{corollary}

\begin{proof}
Louboutin's exact regulator determinant and unit-index lemma give
$R_6=12R_2R_3R_*/Q$ with $Q\mid12$.
For clarity, the logarithmic quadratic form in that reference is
$x^2+y^2-xy$, with $x=a$ and $y=a+c$, hence equals the $R_*$ here.
The conductor-discriminant formula gives
$d_{F_2}=p$, $d_{F_3}=p^2$, and $d_F=p^5$.
The ordinary analytic class-number formula therefore reads
\[
 \begin{aligned}
 \operatorname{Res}_{s=1}\zeta_{F_2}(s)&=\frac{2h(F_2)R_2}{\sqrt p},\\
 \operatorname{Res}_{s=1}\zeta_{F_3}(s)&=\frac{4h(F_3)R_3}{p},\\
 \operatorname{Res}_{s=1}\zeta_F(s)&=\frac{32h(F)R_6}{p^{5/2}}.
 \end{aligned}
 \tag{SC.5}
\]
The Dirichlet-character factorizations give
\[
 \frac{\zeta_F(s)\zeta(s)}{\zeta_{F_2}(s)\zeta_{F_3}(s)}
       =L(s,\chi_6)L(s,\overline{\chi_6}).
\]
Taking the limit at $s=1$ and inserting (SC.5) and the regulator
determinant yields (SC.4).  Its class-number ratio is an integer
by KT-SEXTIC-CLASS-PRODUCT.  Since $Q$ divides twelve, $48/Q$ is
an integer divisible by four.  This proves the assertion.
\end{proof}

\paragraph{Exact primary locator and scope.}
S.~R.~Louboutin, \emph{Class numbers of real cyclotomic fields},
Publ. Math. Debrecen \textbf{64} (2004), printed pp.~454--455:
the relative-unit definition and regulator normalization are in
Section 3.1 on p.~454; Lemma 3 and its entire proof on p.~455
give the determinant and $Q\mid12$.
The paragraph immediately following that proof explicitly states
$h(F_2)h(F_3)\mid h(F)$ and cites Cornell--Washington, Lemma 1.
The latter paper's complete PDF was not retrieved in this audit;
(SC.1) has instead been proved above directly in the needed
prime-conductor case.  Standard Hilbert class-field existence,
base change of unramified extensions, the analytic class-number
formula, and Minkowski's bound are the stated foundational inputs.
No upper bound for $C_6$, and no exclusion of $p$ from a class group,
is asserted in this lemma and corollary alone.

% END evidence/main_gate/sextic_class_product.tex


% BEGIN evidence/lean_chain/order_six_reduction.tex
\subsection{Removing the rational character of order six}
\label{sec:kt-order-six}

The class-number bounds used here are known results.  This section
records their consequence for the precise rational class component in
the report.  It uses the general relative-unit bound of Gras and the
explicit even-character bound of Louboutin; it does not import an
unread statement from the title of Moser's 1981 paper.

\begin{proposition}[KT-ORDER-SIX-COMPONENT-REDUCTION]
Let $p>3$ be prime and let $e\ge2$ be even.  If
$\operatorname{ord}(\omega^e)=6$, then
\[
                 A_p^{[-e]}=0,
                 \qquad K_{2e}(\mathbf Z)\{p\}=0.       \tag{OS6.1}
\]
There is no bound on $e$ relative to $p$.
Together with KT-SMALL-ORDER-COMPONENT-REDUCTION, a remaining
rational character has order $5$ or at least $7$.
\end{proposition}

\begin{proof}
Put $L=\mathbf Q(\mu_p)$ and
$F=L^{\ker(\omega^{-e})}$.  Then $F$ is the real cyclic
sextic field of conductor $p$; in particular $p\equiv1\pmod{12}$.
Write $F_2,F_3$ for its quadratic and cubic subfields,
and $h_i,R_i$ for their ordinary class numbers and regulators,
with $h_6=h(F)$ and $R_6=R(F)$.  All three conductors are $p$,
and the respective discriminants are $p,p^2,p^5$.

For $p<2^{31}$, the carried Hart--Harvey--Ong verification
gives $A_p^+=0$, so the assertion follows from the evenness
of $e$ and K17-(K2).  Hence assume $p\ge2^{31}$.

First,
\[
                         J:=\frac{h_6}{h_2h_3}
                            \quad\hbox{is a positive integer}. \tag{OS6.2}
\]
This is KT-SEXTIC-CLASS-PRODUCT, whose proof uses the ordinary
Hilbert class fields and total ramification at $p$.
The same divisibility is stated in Louboutin, printed p.~455,
immediately after Lemma~3.  The integrality in (OS6.2),
not just a real-valued class-number quotient, is used below.

Let $\sigma$ generate $\operatorname{Gal}(F/\mathbf Q)$,
and use the relative-unit conventions
\[
\begin{split}
 U^*&=\{u\in\mathcal O_F^\times:
       N_{F/F_2}(u),N_{F/F_3}(u)\in\{\pm1\}\},\\
 R^*&=(\log|\varepsilon_*|)^2+
       (\log|\varepsilon_*^\sigma|)^2
       -(\log|\varepsilon_*|)(\log|\varepsilon_*^\sigma|),
\end{split}                                                   \tag{OS6.3}
\]
where $\{-1,\varepsilon_*,\varepsilon_*^\sigma\}$ generates
$U^*$.  Louboutin, Lemma~3, gives
\[
                 R_6=\frac{12R_2R_3R^*}{Q},
                 \qquad Q\in\{1,3,4,12\}.             \tag{OS6.4}
\]
There is no extra signature index in $R^*$.
Indeed $N_{F/F_3}(u)=-1$ would give total norm $-1$,
whereas the norm through $F_2$ of $\pm1$ is $+1$.
Thus $N_{F/F_3}(u)=+1$, and replacing $u$ by $-u$
if necessary makes $N_{F/F_2}(u)=+1$ as well.
This identifies $U^*/\{\pm1\}$ with the exact norm-one
module $E_*$ in Gras, without changing absolute logarithms.

Gras, \S2(b), Theorem~(2.5), with its definition~(2.4),
therefore gives
\[
                    R^*\ge\frac3{16}
                    \log^2\!\left(\frac{p-20}{4}\right).
                                                               \tag{OS6.5}
\]
The bound applies to every real cyclic sextic field of
conductor at least $28$, not just to the special family
introduced later in her \S3.  Its normalization can also
be checked directly from her Lemma~(2.6): for the logarithms
$x,y,y-x,-x,-y,x-y$ of a relative unit, put
$m=\max(|x|,|y|,|x-y|)$.  The lemma gives
$2m\ge\log((p-20)/4)$, and
$x^2+y^2-xy\ge3m^2/4$.  Apply this to a generating unit
to obtain (OS6.5).

Let $\chi$ be either primitive even Dirichlet character of
order six associated with $F$.  The analytic class-number
formula and the factorization
$\zeta_F\zeta_{\mathbf Q}/(\zeta_{F_2}\zeta_{F_3})
=L(s,\chi)L(s,\overline\chi)$ give
\[
                  J=\frac{Qp}{48R^*}|L(1,\chi)|^2.
                                                               \tag{OS6.6}
\]
For example, the residues of the three zeta functions are
$2h_2R_2/\sqrt p$, $4h_3R_3/p$ and
$32h_6R_6/p^{5/2}$; substitution of (OS6.4) yields exactly
(OS6.6).  This also agrees with Louboutin's displayed
class-number identity on printed p.~456 after division
by the cubic class-number formula.

Louboutin, \emph{Acta Arith.}~101 (2002), Theorem~1(1),
equation~(3), gives, for every primitive even character,
\[
 |L(1,\chi)|\le\tfrac12(\log p+\kappa_0),
 \qquad\kappa_0=2+\gamma-\log(4\pi)<1.                \tag{OS6.7}
\]
Combining (OS6.5)--(OS6.7) and $Q\le12$ gives
\[
 J\le\frac p3
       \left(\frac{\log p+\kappa_0}
       {\log((p-20)/4)}\right)^2
   <\frac{49}{108}p<p.                                \tag{OS6.8}
\]
Here the last comparison has ample exact margins: set
$t=\log p>20$, using $\log2>2/3$.  Since $p>100$,
$(p-20)/4>p/5$, and $\log5<2$; hence the displayed
ratio is less than $(t+1)/(t-2)<21/18$.
The bound $\kappa_0<1$ follows from $\gamma<1$ and
$\log(4\pi)>2$.
This proves (OS6.8) without relying on a decimal
evaluation of Euler's constant.

The already imported quadratic and cubic bounds give
$h_2\le\sqrt p/2<p$ and $h_3<2p/3<p$.
Together with (OS6.2) and (OS6.8), they imply
$p\nmid h_6=h_2h_3J$.

Finally put $H=\operatorname{Gal}(L/F)$.
For extension and norm on ideal classes,
\[
                       jN(x)=\sum_{g\in H}g x.         \tag{OS6.9}
\]
If $x\in A_p^H$, then $N(x)=0$ because
$\operatorname{Cl}(F)\{p\}=0$.  Thus $|H|x=0$;
as $p\nmid|H|$, one has $A_p^H=0$.
The $\omega^{-e}$ component lies in these invariants,
so $A_p^{[-e]}=0$.  The carried comparison K17-(K2)
gives the stated $K$-group vanishing in every weight
with this periodic character.
\end{proof}

\paragraph{Exact primary inputs and read scope.}
M.-N.~Gras, \emph{Familles d'unit\'es dans les extensions
cycliques r\'eelles de degr\'e 6 de $\mathbf Q$},
Besan\c con 1984/85--1985/86, fascicule~2:
\S2(a), printed pp.~2--3, fixes the general-field and unit
conventions; \S2(b), definition~(2.4), Theorem~(2.5)
and Lemma~(2.6), printed p.~4, give the regulator bound,
with their complete proofs on printed pp.~4--6.
These are physical PDF pages~5--7 including the cover.
The general result precedes the special-family \S3.
The primary scan is
\url{https://www.numdam.org/article/PMB_1986___2_A2_0.pdf}.
The cited earlier resolvent-divisibility and general
lattice results remain named published inputs; no fresh
review of all of those earlier foundational proofs is claimed.

S.~Louboutin, \emph{Class numbers of real cyclotomic fields},
\emph{Publ. Math. Debrecen}~64 (2004), 451--461,
\S3.1, printed pp.~454--456, gives the relative-unit
definition, Lemma~3 and its proof, the class-number
divisibility, and the exact analytic identity used here:
\url{https://publi.math.unideb.hu/paper/1169/download/10_5486_PMD_2004_2976.pdf}.
His separate analytic lower bound in Theorem~5 is not used.

S.~Louboutin, \emph{Explicit upper bounds for $|L(1,\chi)|$
for primitive even Dirichlet characters},
\emph{Acta Arith.}~101 (2002), 1--18:
Theorem~1(1), equation~(3), printed p.~2, is (OS6.7).
The full \S3 proof, printed pp.~5--8, including Lemma~7
and the $r=0$ completion, was read; its hypotheses allow
nonreal characters of order six.  The official PDF is
\url{https://www.impan.pl/shop/publication/transaction/download/product/82711}.
Exact downloaded-source hashes and review scopes are
recorded in the accompanying evidence manifest.

The small-prime Hart--Harvey--Ong verification, the lower-degree
class-number bounds, and K17-(K2) remain the already accepted
inputs.  The degree-six class-number conclusion is known in the
literature; the contribution here is its explicit source-checked
application to the remaining rational character, not a claim to
a new degree-six class-number theorem.

% END evidence/lean_chain/order_six_reduction.tex


% BEGIN evidence/main_gate/order_five_length_bound.tex
\subsection{A length bound for the four rational characters of order five}

\begin{proposition}[KT-ORDER-FIVE-LENGTH-ONE]
Let $p$ be prime with $p\equiv1\pmod {10}$, let
$L=\mathbf Q(\mu_p)$, and let $F$ be its unique cyclic quintic
subfield.  Write $A_p=\operatorname{Cl}(L)\{p\}$ and let $\omega$
be the Teichm\"uller character.  Then
\[
 \sum_{j=1}^4
 \operatorname{length}_{\mathbf Z_p}
 A_p^{[j(p-1)/5]}\le1.                         \tag{O5.1}
\]
Equivalently, at most one of these four components can be nonzero,
and any nonzero component is isomorphic to $\mathbf Z/p\mathbf Z$.
For every $p\ge2^{31}$ one has the explicit ordinary class-number bound
\[
                         h(F)<\frac{p^2}{6}.   \tag{O5.2}
\]
The proposition gives a size bound, not the vanishing of a specified
order-five component.
\end{proposition}

\begin{proof}
For $p<2^{31}$, the carried Hart--Harvey--Ong verification gives
$A_p^+=0$.  All four characters in (O5.1) are even, because
$(p-1)/5$ is even, so (O5.1) follows in this range.
Assume henceforth $p\ge2^{31}$.

Let $R$ be the ordinary regulator of $F$ and put $h=h(F)$.
The conductor--discriminant formula gives $d_F=p^4$.
Schoof--Washington, Corollary~(2.7), proves for every cyclic
quintic field of conductor $f$ that
\[
                        R\ge\frac1{25}\log^4(f/2).
                                                        \tag{O5.3}
\]
This is a general-field theorem in their Section~2, preceding
and independent of the Lehmer family studied in Section~3.
In particular it applies here with $f=p$.

Let $\chi$ generate the group of Dirichlet characters defining $F$.
Each of $\chi,\chi^2,\chi^3,\chi^4$ is primitive, nontrivial,
even, and of conductor $p$.  The zeta factorization and ordinary
analytic class-number formula give
\[
       \frac{16hR}{p^2}
         =\prod_{j=1}^4 L(1,\chi^j)
         =|L(1,\chi)|^2|L(1,\chi^2)|^2.         \tag{O5.4}
\]
Louboutin's even-character bound, Theorem~1(1), equation~(3), gives
$|L(1,\chi^j)|\le(\log p+\kappa_0)/2$ with
$\kappa_0=2+\gamma-\log(4\pi)<1$.
Combining this with (O5.3)--(O5.4), and writing $t=\log p>20$, yields
\[
 \begin{aligned}
 h&\le\frac{25p^2}{256}
       \left(\frac{\log p+\kappa_0}{\log(p/2)}\right)^4\\
  &<\frac{25p^2}{256}
       \left(\frac{t+1}{t-1}\right)^4
   <\frac{25\cdot21^4}{256\cdot19^4}p^2
   <\frac{p^2}{6}.
 \end{aligned}                                            \tag{O5.5}
\]
Here $t>20$ follows from $p\ge2^{31}$ and $\log2>2/3$;
we used $\log2<1$, $\kappa_0<1$, and the decreasing function
$(t+1)/(t-1)$.  The final comparison is the exact integer inequality
$6\cdot25\cdot21^4=29172150<33362176=256\cdot19^4$.
Consequently $v_p(h)\le1$.

For completeness, this bounds exactly the components in (O5.1).
Put $H=\operatorname{Gal}(L/F)$, so $p\nmid |H|$.
Extension $j$ and norm $N$ on ideal classes satisfy
\[
 Nj=|H|\quad\hbox{on }\operatorname{Cl}(F)\{p\},
 \qquad
 jN=\sum_{g\in H}g\quad\hbox{on }A_p.
                                                        \tag{O5.6}
\]
Thus $j$ identifies $\operatorname{Cl}(F)\{p\}$ with $A_p^H$:
it is injective, and for $x\in A_p^H$ one has
$x=j(|H|^{-1}N(x))$.
The characters occurring in $A_p^H$ are precisely the trivial
character and the four characters in (O5.1).
The trivial component is zero by the same norm identity for
$L/\mathbf Q$, since $\operatorname{Cl}(\mathbf Q)=0$.
The Teichm\"uller idempotents are integral over $\mathbf Z_p$;
hence their four lengths add to
$\operatorname{length}_{\mathbf Z_p}\operatorname{Cl}(F)\{p\}
=v_p(h)\le1$.  This proves (O5.1).
\end{proof}

\paragraph{Exact primary input and limitation.}
R.~Schoof and L.~C.~Washington, \emph{Quintic polynomials and real
cyclotomic fields with large class numbers},
\emph{Math. Comp.}~50 (1988), 543--556:
Section~2 begins on printed p.~544; Theorem~(2.2) and its full proof
are on pp.~545--546, Corollary~(2.6) and its proof are on p.~546,
and Corollary~(2.7) with its full proof is on p.~547.
The last corollary is precisely (O5.3).
The author-hosted primary PDF is
\url{https://reneschoof.github.io/sw.pdf}.
The proof of Corollary~(2.7) imports Pohst's five-variable product
bound $M=16$; no fresh proof audit of that earlier input is claimed.
Louboutin's exact even-character theorem and the carried
Hart--Harvey--Ong verification have the same scopes as in
KT-ORDER-SIX-COMPONENT-REDUCTION.

The quintic source explicitly states on printed p.~544 that the
quintic field of prime conductor $641491$ has class number $1566401$.
Thus an analogue of the bound $h(F)<p$ is false in degree five.
That published numerical assertion is cited as source scope; no
independent class-group recomputation is claimed here.
Neither the bound (O5.2) nor this example decides whether
$p$ can divide $h(F)$.  The contribution is the exact length-one
restriction on a possible remaining rational obstruction.
No integral $K$-group or Selmer-group length comparison is asserted.

% END evidence/main_gate/order_five_length_bound.tex


% BEGIN evidence/unitary_audit/quintic_strict_cyclicity.tex
\subsection{The remaining order-five strict group is cyclic}

Write $\chi_p$ for the $p$-adic cyclotomic character and
$\omega$ for its Teichm\"uller reduction.  For an even integer
$e\ge2$ set
\[
 D=(\mathbf Q_p/\mathbf Z_p)(\chi_p^{-e}),\qquad
 \mathcal S_p(e)=
 \ker\!\left(H^1(G_{\mathbf Q,\{p\}},D)
       \longrightarrow H^1(G_{\mathbf Q_p},D)\right).
 \tag{O5S.1}
\]
Thus the classes are unramified away from $p$ and are zero
on the full decomposition group at $p$.
The carried VK-LOCAL comparison identifies this with the
source's $\operatorname{ShaU}(\chi_p^{-e};2p)$: its additional
decomposition condition at $2$ is automatic for the divisible
coefficients in (O5S.1).

\begin{corollary}[KT-ORDER-FIVE-STRICT-CYCLIC]
Let $p>3$ be prime, let $e\ge2$ be even, and suppose
$\operatorname{ord}(\omega^e)=5$.
Put $L=\mathbf Q(\mu_p)$, $\Delta=\operatorname{Gal}(L/\mathbf Q)$,
$A_p=\operatorname{Cl}(L)\{p\}$ and $\alpha=\omega^{-e}$.
There is a natural injection
\[
 \mathcal S_p(e)[p]\lhook\joinrel\longrightarrow
 \operatorname{Hom}_{\mathbf F_p[\Delta]}
       \bigl(A_p/pA_p,\mathbf F_p(\alpha)\bigr).
 \tag{O5S.2}
\]
Consequently
\[
            \dim_{\mathbf F_p}\mathcal S_p(e)[p]\le1.
 \tag{O5S.3}
\]
Under the carried finiteness and full-group VK-PT comparison,
there is an integer $a\ge0$ such that
\[
 \mathcal S_p(e)\simeq\mathbf Z/p^a\mathbf Z,
 \qquad K_{2e}(\mathbf Z)\{p\}\simeq\mathbf Z/p^a\mathbf Z.
 \tag{O5S.4}
\]
The argument asserts neither $a=0$ nor an upper bound on $a$.
There is no restriction $e<p$.
\end{corollary}

\begin{proof}
Let $M=\mathbf F_p(\alpha)$ and identify it with $D[p]$.
Since $\alpha$ has order five, the cyclotomic image of local
inertia at $p$ contains an element $g$ for which
$\chi_p^{-e}(g)-1\in\mathbf Z_p^\times$.
It follows that
\[
 H^0(G_{\mathbf Q},D)=H^0(G_{\mathbf Q_p},D)=0.
 \tag{O5S.5}
\]
The exact sequence
\[
          0\longrightarrow M\longrightarrow D
            \xrightarrow{\ p\ }D\longrightarrow0
 \tag{O5S.6}
\]
therefore identifies $H^1(G_{\mathbf Q},M)$ with
$H^1(G_{\mathbf Q},D)[p]$, and similarly gives an injection
$H^1(G_{\mathbf Q_p},M)\hookrightarrow
H^1(G_{\mathbf Q_p},D)$.
In particular, the finite-coefficient class corresponding to
$s\in\mathcal S_p(e)[p]$ is zero on $G_{\mathbf Q_p}$.
For any prime $q\ne p$, inertia acts trivially on $D$,
and $H^0(I_q,D)=D$ is $p$-divisible.
The same exact sequence gives an injection
\[
                    H^1(I_q,M)\hookrightarrow H^1(I_q,D).
 \tag{O5S.7}
\]
The finite-coefficient class is thus unramified at every
$q\ne p$ as well.  These arguments also justify using the
unramified-outside-$p$ quotient in (O5S.1).

The coefficient $M$ becomes trivial on $G_L$.
Since $|\Delta|=p-1$ is prime to $p$, inflation--restriction
and averaging give an injective restriction
\[
 H^1(G_{\mathbf Q},M)\lhook\joinrel\longrightarrow
 H^1(G_L,M)^\Delta.
 \tag{O5S.8}
\]
The restricted class is a continuous homomorphism $G_L\to M$.
It is unramified at every finite place of $L$: at the places
above $p$ this follows from the full local vanishing already
proved, and elsewhere from (O5S.7).
Global class field theory therefore makes it factor through
$A_p/pA_p$.
The equivariance convention is explicit: for a lift of
$\delta\in\Delta$ and $g\in G_L$, a restricted cocycle $c$
satisfies
\[
                  c(\delta g\delta^{-1})=\alpha(\delta)c(g).
 \tag{O5S.9}
\]
Functoriality of the Artin map gives precisely the Hom group
in (O5S.2), with character $\alpha$, not its inverse.
This proves the asserted injection.

Integral Teichm\"uller idempotents identify the dimension of
the Hom group with
\[
          \dim_{\mathbf F_p}
             A_p^{[-e]}/pA_p^{[-e]}.
 \tag{O5S.10}
\]
The character $\alpha$ is one of the four order-five characters
in KT-ORDER-FIVE-LENGTH-ONE.  That proposition bounds the sum
of their $\mathbf Z_p$-lengths by one, and proves (O5S.3).

The carried VK-PT theorem is a full group isomorphism
\[
              \bigl(K_{2e}(\mathbf Z)\{p\}\bigr)^\vee
                       \simeq\mathcal S_p(e),
 \tag{O5S.11}
\]
not merely a comparison of vanishing assertions.
The positive even-degree $K$-group here is finite by the
carried finiteness input.  Hence $\mathcal S_p(e)$ is a finite
abelian $p$-group.  For any such group $B$, the dimensions of
$B[p]$ and $B/pB$ are equal; dimension at most one means that
$B$ is cyclic.  Its Pontryagin dual is cyclic of the same
order.  This proves (O5S.4).
\end{proof}

\paragraph{Why the class-group length does not give the same exponent.}
The proof uses the fact that $D[p]$ is trivial on $G_L$.
This does not hold for all higher torsion chunks of $D$.
More precisely, put $r=v_p(e)$.  The image of $G_L$ under
$\chi_p$ is $1+p\mathbf Z_p$, and the lifting-the-exponent
identity gives
\[
                  v_p((1+p)^e-1)=r+1.
 \tag{O5S.12}
\]
Thus $D[p^{r+2}]$ has nontrivial $G_L$-action; when $p\nmid e$
this already happens for $D[p^2]$.
The same restriction argument therefore does not produce an
all-level replacement of (O5S.2) by homomorphisms from the same
ordinary ideal class group.
The finite groups $\mathbf Z/p^a\mathbf Z$ all have
one-dimensional $p$-torsion when $a\ge1$.
Therefore (O5S.3) alone does not bound the integer $a$.
This is the precise limitation of the comparison proved here;
no claim that an additional argument cannot bound $a$ is made.

\paragraph{The ordinary $R=T$ cyclicity hypothesis remains distinct.}
Suppose, in addition, that the hypotheses of
KT-RESIDUAL-CYCLICITY-SCOPE hold: the prime splits in
$E=\mathbf Q(\sqrt{-3})$, and the specified ordinary
irreducible rank-two congruence supplies its Urban class.
With the notation of that theorem, one has
\[
 \dim_k\mathcal H_E^{c=-1}\ge1,
 \qquad
 \dim_k\mathcal H_E^{c=+1}
        \ge\dim_{\mathbf F_p}\mathcal S_p(e)[p].
 \tag{O5S.13}
\]
If the rational group in (O5S.4) is nonzero, its $p$-torsion
has dimension exactly one, and the two eigenspaces in
(O5S.13) give $\dim_k\mathcal H_E\ge2$.
Thus the new one-generator conclusion for the rational
group does not supply the one-dimensional \emph{total}
ordinary Selmer hypothesis in the inspected Berger--Klosin
comparison.  Outside those Urban-branch hypotheses,
(O5S.3) still gives no upper bound for the additional
$c=-1$ eigenspace.  The previously checked weight-range,
ordered-flag, and divisible-Selmer-bound hypotheses of
those comparisons are unchanged.  No $R=T$ theorem is
being applied in this corollary.

\paragraph{Exact carried comparison and source scope.}
The full-group comparison (O5S.11) and the equality of
the local conditions are the accepted K17 theorem
\emph{VK-PT and VK-LOCAL}, equations (K5)--(K7), in
\path{RETURN_K17_checkpoint_08_20261121T0616Z/K17.tex},
lines 788--898.  Its entire proof was reread for this
application; its named foundational inputs remain carried.
The $\operatorname{ShaU}$ definition and its exact
finite-coefficient local transport were also recorded in
KT-ODD-CONORMAL-SELMER-EQUIVALENCE.
The order-five class bound is the separately checked
KT-ORDER-FIVE-LENGTH-ONE.  The additional ordinary
eigenspace in (O5S.13) is used only with all the hypotheses
of KT-RESIDUAL-CYCLICITY-SCOPE, not asserted for an arbitrary
prime or an arbitrary order-five character.

% END evidence/unitary_audit/quintic_strict_cyclicity.tex


% BEGIN evidence/main_gate/order_five_pdiv_weight.tex
\subsection{An exponent bound when the weight is divisible by the residue prime}

\begin{corollary}[KT-ORDER-FIVE-PDIV-WEIGHT]
Let $p>3$ be prime and let $e\ge2$ be even, with
$\operatorname{ord}(\omega^e)=5$ and $p\mid e$.
For the strict group $\mathcal S_p(e)$ of (O5S.1), one has
\[
 \mathcal S_p(e)\simeq K_{2e}(\mathbf Z)\{p\}
       \simeq0\quad\hbox{or}\quad\mathbf Z/p\mathbf Z.
                                                        \tag{O5P.1}
\]
The displayed isomorphisms are isomorphisms of abstract finite groups.
This is an exponent bound in the stated weights; it does not prove
vanishing of the remaining order-five class component.
\end{corollary}

\begin{proof}
Put $L=\mathbf Q(\mu_p)$, $\Delta=\operatorname{Gal}(L/\mathbf Q)$,
$D=(\mathbf Q_p/\mathbf Z_p)(\chi_p^{-e})$, and $M_2=D[p^2]$.
The cyclotomic image of $G_L$ is $1+p\mathbf Z_p$.
Since $p\mid e$, the binomial congruence
$(1+pa)^e\equiv1\pmod {p^2}$ shows that $M_2$ is trivial on
$G_L$.  Its action thus factors through $\Delta$ and is
\[
                M_2\simeq(\mathbf Z/p^2\mathbf Z)(\widetilde\alpha),
                \qquad\alpha=\omega^{-e},                \tag{O5P.2}
\]
where $\widetilde\alpha$ is the Teichm\"uller lift of $\alpha$.
This lift is unique because $|\Delta|=p-1$ is prime to $p$.

The residual inertia character has order five.  An inertia element
on which $\alpha-1$ is nonzero therefore acts on $D$ with
$\chi_p^{-e}-1$ a unit.  Thus the global and $p$-local invariant
groups of $D$ vanish.  Applying cohomology to
\[
 0\longrightarrow M_2\longrightarrow D
       \xrightarrow{\ p^2\ }D\longrightarrow0             \tag{O5P.3}
\]
identifies $p^2$-torsion global classes with finite-coefficient
classes and preserves their strict local vanishing at $p$.
For $q\ne p$, inertia acts trivially on $D$, and its invariant
group $D$ is $p^2$-divisible.  The same exact sequence makes
$H^1(I_q,M_2)\longrightarrow H^1(I_q,D)$ injective, so the
finite-coefficient class is unramified there as well.

Restriction to $G_L$ is injective by inflation--restriction,
since $|\Delta|$ is prime to $p$.
It now gives an unramified homomorphism to $M_2$, and global
class field theory yields an injection
\[
 \mathcal S_p(e)[p^2]\lhook\joinrel\longrightarrow
 \operatorname{Hom}_{\mathbf Z_p[\Delta]}
   \bigl(A_p, (\mathbf Z/p^2\mathbf Z)(\widetilde\alpha)\bigr),
 \qquad A_p=\operatorname{Cl}(L)\{p\}.                   \tag{O5P.4}
\]
The equivariance has the same convention as (O5S.9).
Integral Teichm\"uller idempotents show that the right side
only uses $A_p^{[-e]}$.  KT-ORDER-FIVE-LENGTH-ONE gives this
component length at most one, so the right side has order
at most $p$.

By KT-ORDER-FIVE-STRICT-CYCLIC and the carried finiteness theorem,
$\mathcal S_p(e)\simeq\mathbf Z/p^a\mathbf Z$ for some $a\ge0$.
Its $p^2$-torsion has order $p^{\min(a,2)}$.
The bound from (O5P.4) therefore gives $a\le1$.
The full-group VK-PT comparison identifies the Pontryagin dual
of $K_{2e}(\mathbf Z)\{p\}$ with $\mathcal S_p(e)$,
which proves (O5P.1).
\end{proof}

\paragraph{Scope.}
The exponent conclusion uses the finite-level class-field injection
(O5P.4) and the full-group VK-PT comparison.
For $p\nmid e$, the coefficient $D[p^2]$ already has nontrivial
$G_L$-action, so this argument does not extend the exponent bound
to those weights.  The corollary leaves the same possible nonzero
class component; it removes no additional character or prime.
All class-number and foundational comparison inputs are the
separately checked theorems cited above, with their existing
source and review records.

% END evidence/main_gate/order_five_pdiv_weight.tex


% BEGIN evidence/main_gate/order_five_greenberg_gate.tex
\subsection{The additional unit condition in the quintic capitulation criterion}

\begin{proposition}[KT-ORDER-FIVE-GREENBERG-UNIT-GATE]
Let $p\equiv1\pmod {10}$ be prime, let $F$ be the cyclic quintic
field of conductor $p$, and let $F_n/F$ be the degree-$p^n$ layer
of its cyclotomic $\mathbf Z_p$-extension.  Put
\[
 A_n=\operatorname{Cl}(F_n)\{p\},\qquad
 E_n=\mathcal O_{F_n}^{\times},\qquad
 G_n=\operatorname{Gal}(F_n/F),
 \qquad j_n:A_0\longrightarrow A_n.
                                                        \tag{O5G.1}
\]
For every $n\ge1$, all units in $E_0$ are norms from
$F_n^{\times}$, and there is an exact sequence
\[
 1\longrightarrow j_n(A_0)\longrightarrow A_n^{G_n}
 \longrightarrow E_0/N_{F_n/F}E_n\longrightarrow1,
 \qquad |A_n^{G_n}|=|A_0|.                              \tag{O5G.2}
\]
Consequently
\[
 [E_0:N_{F_n/F}E_n]=|\ker(j_n)|.                        \tag{O5G.3}
\]
If the remaining possibility in KT-ORDER-FIVE-LENGTH-ONE is
$A_0\simeq\mathbf Z/p\mathbf Z$, then the following are equivalent:
\begin{enumerate}
\item $[E_0:N_{F_n/F}E_n]=p$ for some $n$;
\item $A_0$ capitulates in some $F_n$;
\item $\varinjlim_n A_n=0$, where the transition maps extend ideals;
\item the orders $|A_n|$ are bounded, equivalently the cyclotomic
      Iwasawa invariants satisfy $\lambda=\mu=0$.
\end{enumerate}
In particular, the required extra condition concerns norms of
\emph{units}; the assertion that every base unit is a norm of a
field element already holds without capitulation.
\end{proposition}

\begin{proof}
The field $F_n$ is the unique degree-$5p^n$ subfield of
$\mathbf Q(\mu_{p^{n+1}})$.  It is totally real and has a unique
prime $\mathfrak p_n$ above $p$, totally ramified over $\mathbf Q$.
Moreover,
\[
 \pi_n=N_{\mathbf Q(\mu_{p^{n+1}})/F_n}
                       (1-\zeta_{p^{n+1}}),\qquad
                    (\pi_n)=\mathfrak p_n.              \tag{O5G.4}
\]
Indeed the ideal norm of the unique prime in the full cyclotomic
field has exponent its residue degree, which is one.  Thus the
only prime ramified in $F_n/F$ has principal ideal class.

At every finite place other than $p$, a base unit is a local norm:
the local extension is unramified and the norm on local units is
surjective.  At real places there is no complexification.
The product formula for the cyclic norm-residue symbols then
forces the symbol at the unique $p$-place to be trivial as well.
The Hasse norm theorem for the cyclic extension $F_n/F$ gives
\[
                    E_0\subset N_{F_n/F}(F_n^{\times}). \tag{O5G.5}
\]

Here is a direct class-field proof of the order assertion in
(O5G.2).  Let $H_n$ be the ordinary Hilbert $p$-class field of
$F_n$, and let $M_n$ be its maximal subfield abelian over $F$.
The field $H_n$ is Galois over $F$, and
\[
 \operatorname{Gal}(M_n/F_n)
       \simeq A_n/(\sigma-1)A_n,                       \tag{O5G.6}
\]
for a generator $\sigma$ of $G_n$.
This follows because the commutator subgroup is
$(\sigma-1)A_n$: after quotienting by this subgroup the kernel
is central and the quotient is cyclic, so the entire group is
abelian.  In the abelian extension $M_n/F$, the inertia group
$I$ at $p$ has order exactly $p^n$, since $M_n/F_n$ is
unramified and $F_n/F$ is totally ramified.  It maps
isomorphically to $G_n$.  Its fixed field $M_n^I$ is everywhere
unramified over $F$, including the real places, and hence is
contained in $H_0$.  Thus $M_n\subset F_nH_0$.
Conversely $F_nH_0/F_n$ is unramified abelian of $p$-power degree
and $F_nH_0/F$ is abelian, so $F_nH_0\subset M_n$.
Total ramification gives $F_n\cap H_0=F$, and therefore
$|A_n/(\sigma-1)A_n|=|A_0|$.
For a finite module the kernel and cokernel of $\sigma-1$
have equal orders, proving $|A_n^{G_n}|=|A_0|$.

For completeness, the ambiguous-class map sends an invariant
class represented by $\mathfrak a$, with
$\mathfrak a^{\sigma-1}=(x)$, to $N_{F_n/F}x$ modulo
$N_{F_n/F}E_n$.  This is well-defined.  Its kernel consists
of classes with an invariant ideal representative: if the
norm equals the norm of a unit, divide $x$ by that unit and
apply Hilbert's theorem 90.  Conversely an invariant ideal
representative gives the trivial unit class.
The map is onto $E_0\cap N(F_n^{\times})/N E_n$:
for a unit $N x$, the ideal $(x)$ has ideal norm one, and the
norm-one kernel on the permutation module of fractional ideals
is $(\sigma-1)I(F_n)$.
An invariant fractional ideal is a product of an extended ideal
of $F$ and powers of $\mathfrak p_n$.  Equation (O5G.4) removes
the latter classes.  Taking $p$-primary parts and using (O5G.5)
therefore gives exactly (O5G.2).  The unit quotient is a
$p$-group, since the $p^n$-th power of any base unit is its norm
when regarded as a unit upstairs.  Orders now prove (O5G.3).

Under $|A_0|=p$, (O5G.3) proves the equivalence of the first
two assertions.  The unique $p$-prime is totally ramified in
the whole tower, so Greenberg's one-prime capitulation criterion
applies.  In the notation of Ichimura's published version,
Theorem~1 gives $\varinjlim A_n=0$ if the image of $A_0$ is zero;
the converse is immediate from finiteness of $A_0$.
Under the same one-prime hypothesis, his Remark~1(1) identifies
the auxiliary extension $L_n$ with $H_nF_\infty$.
Total ramification gives $H_n\cap F_\infty=F_n$, so
$[L_n:F_\infty]=|A_n|$.  Remark~1(2) identifies vanishing of
the direct limit with finiteness of $\bigcup_nL_n/F_\infty$,
which is therefore equivalent to bounded $|A_n|$.
The standard Iwasawa class-number formula gives the stated
equivalence with $\lambda=\mu=0$.
\end{proof}

\paragraph{What the criterion does and does not supply.}
When $A_0\simeq\mathbf Z/p\mathbf Z$, the norm-unit index is
either $1$ or $p$.  If it is $1$ at every layer, $j_n$ is
injective at every layer.  If it becomes $p$, the base class
capitulates and the tower has bounded class numbers.
The order-five class bound does not evaluate this additional
norm-unit index.  The cited criterion assumes capitulation,
equivalently the displayed norm-unit condition; it does not
deduce it from cyclicity of $A_0$.
Even the bounded-growth conclusion asserts capitulation of
classes after extension, not their triviality in $F$.
No order-five component is removed by this proposition.

\paragraph{Exact primary scope.}
H.~Ichimura, \emph{A note on the ideal class group of the
cyclotomic $\mathbf Z_p$-extension of a totally real number
field}, Acta Arithmetica \textbf{105.1} (2002), 29--34:
Theorems~1--2 and Proposition~1 are on p.~29;
Remark~1(2) and the imported bounded-unit-index Lemma~1 are on
p.~30; the complete proofs of Lemma~2, Theorem~2 and
Proposition~1 are on p.~31.  These pages were read in full.
This is a direct published exposition of Greenberg's
1976 Theorem~1; no fresh read of the original 1976 proof is
claimed.  The norm-unit exact sequence is also stated as
equation (1.1), p.~4, of G.~Gras,
\emph{Algebraic norms and capitulation of $p$-class groups in
ramified cyclic $p$-extensions}, arXiv:2211.12279v3
(6 February 2023).  Its general setup on pp.~2--5 and
cyclotomic discussion on pp.~23--25 were read.
The restricted tame setting of that paper's Section~2.2 is
not substituted for the present cyclotomic extension.
The argument above proves the specialized exact sequence and
ambiguous-class order directly; the Hasse norm theorem,
Hilbert's theorem 90 and global class field theory are imported
foundational results, not freshly reproved here.

% END evidence/main_gate/order_five_greenberg_gate.tex


% BEGIN evidence/main_gate/order_five_greenberg_abstract_model.tex
\subsection{An abstract cyclic norm system with unbounded growth}

\begin{proposition}[KT-ORDER-FIVE-ABSTRACT-NORM-LIMIT]
Let $p\equiv1\pmod {10}$ be prime and identify
$\Lambda=\mathbf Z_p[[\Gamma]]=\mathbf Z_p[[T]]$ for
$\Gamma\simeq\mathbf Z_p$ with a topological generator $\gamma=1+T$.
There is an explicit continuous $\Lambda$-module $X$, with a commuting
action of $C_5$ through one primitive Teichm\"uller character, for which
the finite quotients $A_n=X/\omega_n(T)X$, where
$\omega_n(T)=(1+T)^{p^n}-1$, have the following properties:
\[
 A_n\simeq\mathbf Z/p^{n+1}\mathbf Z,\qquad
 |A_n^{\Gamma}|=p,\qquad
 j_{0,n}:A_0\lhook\joinrel\longrightarrow A_n,
 \qquad j_{0,n}(A_0)=A_n^{\Gamma}.                     \tag{O5M.1}
\]
The natural norm and transfer maps satisfy their degree and algebraic
norm identities.  Nevertheless $|A_n|$ is unbounded, and $X$ has
$\mu=0$, $\lambda=1$.
This is an abstract module example, not a class group of a number field.
\end{proposition}

\begin{proof}
Take
\[
             X=\Lambda/(T+p)\simeq\mathbf Z_p,
             \qquad \gamma=1+T\ \text{acts as }1-p.
                                                               \tag{O5M.2}
\]
For odd $p$, the lifting-the-exponent identity gives
\[
       v_p\bigl((1-p)^{p^n}-1\bigr)=n+1.               \tag{O5M.3}
\]
For clarity, this also follows inductively: if $v_p(z)=r\ge1$,
then the term $pz$ in $(1+z)^p-1$ has valuation $r+1$;
all middle terms have valuation at least $1+2r>r+1$, and
$z^p$ has valuation $pr>r+1$.  Start with $z=-p$.
Thus $A_n$ has the order and abstract structure in (O5M.1).
The action of $\gamma$ on it factors through a group of order $p^n$.

For $m\ge n$, put
\[
 \nu_{m,n}(T)=\frac{\omega_m(T)}{\omega_n(T)}
       =\sum_{i=0}^{p^{m-n}-1}(1+T)^{ip^n}.
                                                               \tag{O5M.4}
\]
Define $N_{m,n}:A_m\to A_n$ by the natural quotient and
$j_{n,m}:A_n\to A_m$ by multiplication by $\nu_{m,n}(-p)$.
These are well-defined because $\omega_n$ divides $\omega_m$.
Equation (O5M.3) gives
\[
              v_p(\nu_{m,n}(-p))=m-n.                 \tag{O5M.5}
\]
Every $N_{m,n}$ is onto and every $j_{n,m}$ is injective.
The maps compose correctly because the quotients $\nu_{m,n}$
multiply along a tower.  Moreover,
\[
 N_{m,n}j_{n,m}=p^{m-n}\quad\text{on }A_n,
 \qquad
 j_{n,m}N_{m,n}=\sum_{i=0}^{p^{m-n}-1}\gamma^{ip^n}
       \quad\text{on }A_m.                            \tag{O5M.6}
\]
For the first identity, $(1+T)^{p^n}=1$ on $A_n$;
the second is precisely the defining polynomial (O5M.4).

Since $\gamma-1$ acts as $-p$, its kernel on $A_n$ is the
subgroup of order $p$.  By (O5M.5), this is exactly the image
of $A_0$ under $j_{0,n}$, proving (O5M.1).
The inverse limit under the quotient norms is $\mathbf Z_p=X$.
Its distinguished characteristic polynomial is $T+p$, giving
$\mu=0$ and $\lambda=1$.  In the direct system of transfers,
the nonzero class of $A_0$ survives because every transfer is
injective.

Finally, choose a primitive fifth root of unity
$\xi\in\mathbf Z_p^{\times}$, which exists because $5\mid p-1$,
and let a generator of $C_5$ act by multiplication by $\xi$.
This commutes with $T$ and all the displayed maps.  Each $A_n$
therefore has just that one primitive order-five character.
\end{proof}

\paragraph{Exact limitation.}
The base component has length one, the ambiguous subgroup has the
same order at every level, and all the stated norm/transfer identities
hold, but bounded growth fails.  Thus those module data alone do not
force the extra unit-norm condition in
KT-ORDER-FIVE-GREENBERG-UNIT-GATE.
No realization by class fields, global units, or an actual
cyclotomic number-field tower is asserted.  In particular, this
example is not a counterexample to Greenberg's conjecture and does
not remove an arithmetic class component.

% END evidence/main_gate/order_five_greenberg_abstract_model.tex


% BEGIN evidence/lean_chain/class_pairing_scope.tex
\subsection{The class reciprocity pairing does not impose quintic parity}
\label{KT-QUINTIC-CLASS-PAIRING-SCOPE}
\noindent\textbf{KT-QUINTIC-CLASS-PAIRING-SCOPE.}\quad

This is an applicability check, not a class-group vanishing theorem.
Let $p$ be odd, let $F=\mathbb Q(\mu_p)$, and write
$\Delta=\operatorname{Gal}(F/\mathbb Q)$ and $\omega$ for the mod-$p$
cyclotomic character.  The surviving component in the quintic question has
character $\alpha$ of order five and $\alpha(-1)=1$.  The previously proved
length bound allows this component to have order $p$.

\paragraph{Exact primary construction.}
The source is A.~Morgan and A.~Smith, \emph{Field change for the
Cassels--Tate pairing and applications to class groups}, Research in
Number Theory \textbf{10}, 61 (2024),
\path{doi:10.1007/s40993-024-00545-2}.
Its notation is
\[
 \operatorname{Cl}^{*}L=
 \operatorname{Hom}(\operatorname{Gal}(H_L/L),\mathbb Q/\mathbb Z).
\]
The perfect reciprocity pairing in (1.5), printed pp.~3--4, is between
$\operatorname{Cl}^{*}L[p^\infty]$ and
$\operatorname{Cl}L[p^\infty]$.  It is not an asserted alternating
self-pairing of the class group.

Definition~5.2, Notation~5.3, and Lemma~5.4, printed pp.~20--21, give
\[
 1\longrightarrow\mathcal O_L^\times/(\mathcal O_L^\times)^p
 \longrightarrow\operatorname{Sel}_pL
 \longrightarrow\operatorname{Cl}L[p]\longrightarrow1,
 \qquad
 \operatorname{Sel}_pL\simeq\operatorname{Sel}\mu_p.
 \tag{CP.1}
\]
Thus the Kummer Selmer input is not simply a second copy of
$\operatorname{Cl}L[p]$; it has the indicated unit kernel.
For the unramified exact sequence $E(a,b)$ in (5.3), Proposition~5.5
identifies its Cassels--Tate pairing with minus the reciprocity pairing
after this class-group projection.  The general pairing theorem recalled
as Theorem~2.1, printed p.~8, specifies kernels coming from the middle
Selmer group; it does not assert perfection before those kernels are
removed.

\paragraph{What the roots of unity supply.}
In \S5.2 the chosen identification of cyclic groups with roots of unity
gives a morphism $g_m:\frac1m\mathbb Z/\mathbb Z\to\mu_m$ in the
category of Selmer objects when $\mu_m\subset L$.  The inclusion of local
conditions used to define this morphism is not asserted to be an equality.
Proposition~5.6, printed p.~22, assumes $\mu_{ab}\subset L$ so that the
whole exact sequence, including its middle term, admits the requisite
map to the dual sequence.  It proves
\[
 \operatorname{CTP}_{E(a,b)}(\phi,g_a\psi)
 =\operatorname{CTP}_{E(b,a)}(\psi,g_b\phi).
 \tag{CP.2}
\]
For $a=b$ this is explicitly \emph{symmetry}.  In particular the case
$a=b=p$ requires $\mu_{p^2}\subset L$, which fails for
$L=\mathbb Q(\mu_p)$ and for its totally real quintic subfield.
The proof of Proposition~1.3 on the same page is an exact reformulation
of (CP.2) as a reciprocity identity; it does not change symmetry into
alternation.  A one-dimensional group $\mathbb Z/p\mathbb Z$ admits the
nondegenerate symmetric pairing $(x,y)\mapsto xy/p\bmod\mathbb Z$, so
symmetry alone cannot exclude the allowed order-$p$ case.

\paragraph{The actual character obstruction.}
The source defines Cartier duality by
$M^\vee=\operatorname{Hom}(M,(\overline{\mathbb Q})^\times)$
(\S2.1, printed p.~7).  Consequently
\[
 \bigl(\mathbb F_p(\alpha)\bigr)^\vee
   =\mathbb F_p(\omega\alpha^{-1}),\qquad
 \operatorname{Hom}_{G_{\mathbb Q}}\!
 \left(\mathbb F_p(\alpha),
       \mathbb F_p(\omega\alpha^{-1})\right)=0.
 \tag{CP.3}
\]
Indeed complex conjugation is $+1$ on the first character and $-1$ on
the second.  Equivalently, the required identity $\alpha^2=\omega$
is impossible by parity.  Even a proposed untwisted equivariant
self-duality with values in the constant module would require
$\alpha=\alpha^{-1}$, contrary to its order five.  The ordinary
perfect evaluation pairing pairs the character-$\alpha$ class piece
with its dual character piece in $\operatorname{Cl}^{*}$; it supplies
no identification of that piece with itself.

Passing to a field containing $\mu_{p^2}$ would change the class group
under discussion, and the proposition there still supplies the symmetric
identity (CP.2).  It does not manufacture the missing equivariant
self-duality of (CP.3) on the original component.  Hence the cited
construction supplies no perfect alternating self-pairing that would
force the possible single order-$p$ quintic component to vanish.
This conclusion is limited to the exact pairing hypotheses just checked;
it does not exclude other arithmetic input for the component.

\paragraph{Checked source scope.}
The complete requested material in \S\S5.1--5.2, including the proofs
of Lemma~5.4, Propositions~5.5--5.6 and the deduction of Proposition~1.3,
was read, together with the defining conventions and Theorem~2.1 in
\S2.  The general Cassels--Tate machinery from the authors' preceding
paper remains a cited input.  The published correction
\path{doi:10.1007/s40993-024-00552-3} adds an author affiliation and
contains no mathematical correction.

% END evidence/lean_chain/class_pairing_scope.tex


% BEGIN evidence/unitary_audit/dimension2_comparison/dimension_two_comparison.tex
\subsection{What a two-generator ordinary comparison would require}

The order-five result bounds the original rational strict group by
one generator.  On the split branches of
KT-RESIDUAL-CYCLICITY-SCOPE, a nonzero rational group and the
separately constructed negative-conjugation class give
\[
 \dim_k\mathcal H_E\geq2.                                      \tag{D2.1}
\]
There is no upper bound of two for this entire ordinary group in
the results proved here.  In particular, ``dimension two'' below
is a proposed additional hypothesis, not a consequence of the
quintic class-number estimate.

\paragraph{Published comparisons do allow more than one extension.}
Berger--Klosin, \emph{On lifting and modularity of reducible residual
Galois representations over imaginary quadratic fields}, IMRN 2015,
use no dimension-one assumption for their extension space.
In the accepted author version, arXiv:1606.06535, Theorem~8.2
(pp.~27--28) requires a projectively unique modular basis and
finitely many characteristic-zero lifts for each member, together
with the preceding local hypotheses and a divisible Selmer upper
bound matched by a congruence lower bound.  Theorem~8.5 (p.~28)
replaces projective uniqueness by
\[
 \varpi H^1_\Sigma
       (E,\Psi^{-1}\otimes\operatorname{Frac}(\mathcal O)/\mathcal O)
       =0.                                                     \tag{D2.2}
\]
It retains the other hypotheses.  Assumptions~2.5 and~4.2
(pp.~6, 12) explicitly require
\[
 \#H^1_\Sigma(E,\Psi^{-1}\otimes
              \operatorname{Frac}(\mathcal O)/\mathcal O)
       \leq\#\mathcal O/(L)\leq\#T/J.                          \tag{D2.3}
\]
Lemma~4.4 (p.~13) uses Fitting ideals to compare the lattice
quotient with this Selmer group.  The final surjectivity in that
lemma uses both inequalities in (D2.3); it does not furnish the
left inequality from modularity.

These are two-constituent rank-two results.  The imaginary
quadratic principality argument, Proposition~7.1 and its proof
(p.~23), assumes an anticyclotomic residual ratio
$\bar\Psi^c=\bar\Psi^{-1}$.  For the Tate ratio of order five,
\[
 (\omega^{e}|_{G_E})^c=\omega^{e}|_{G_E},
 \qquad \omega^{e}|_{G_E}\ne\omega^{-e}|_{G_E},                 \tag{D2.4}
\]
provided $E$ is prime to $p$, as in the auxiliary-field setup.
Indeed $E\cap\mathbf Q(\mu_p)=\mathbf Q$ here, so restriction
preserves the order five.  A common twist does not change a
constituent ratio, and interchanging the two constituents only
inverts it.  Thus neither operation repairs (D2.4).
The rank-three residual representation also cannot be treated as
two absolutely irreducible constituents by combining two of its
three distinct characters into a block.

The published conclusion concerns the characteristic-zero
quotient $R_\tau^0$: the authors explicitly separate it from
possible characteristic-$p$ components, in Section~2.5
(pp.~8--9).  It supplies no injection of our integral odd
off-diagonal module into the resulting trace algebra.

\paragraph{The divisible-by-$p$ weight does not verify (D2.2).}
KT-ORDER-FIVE-PDIV-WEIGHT proves
\[
 p\,\mathcal S_p(e)=0\qquad(p\mid e),                         \tag{D2.5}
\]
for the rational strict group with coefficient
$(\mathbf Q_p/\mathbf Z_p)(\chi_p^{-e})$.
Restriction places these classes in the positive-conjugation
part of the corresponding group over $E$; it does not identify
them with its whole divisible Selmer group.  It bounds neither
the negative-conjugation part nor a relaxation at any allowed
ramified prime.  Moreover the character in (D2.2) is the
anticyclotomic character of the cited theorem, and its local
condition is the theorem's crystalline condition at every
prime over $p$ (Section~2.3, pp.~5--6).  These are not supplied
by the inclusion of the original rational strict group.
Even with the coefficient ring $\mathbf Z_p$, (D2.5) therefore
does not verify (D2.2).  With a ramified coefficient extension,
annihilation by $p$ would additionally be weaker than
annihilation by its uniformizer.

There is a separate local obstruction to applying that paper to
the actual large-weight Tate lifts: its ``crystalline'' deformation
condition is the short Fontaine--Laffaille condition
$\operatorname{Fil}^0=D$, $\operatorname{Fil}^{p-1}=0$
(Section~2.1, pp.~3--4).  If $p\mid e$ and $e>0$ is even, then
$e\geq2p$.  The Hodge--Tate weight difference $e$ of an adjacent
Tate pair cannot be put in that interval by a common twist.
Replacing $e$ by a congruent smaller integer changes the
characteristic-zero local condition and the divisible
coefficient group, so it does not prove the desired application.

\paragraph{Nonprincipal Eisenstein ideals are compatible with $R=T$.}
Wake--Wang-Erickson, \emph{The Eisenstein ideal with squarefree
level}, author version dated 19 January 2021,
Theorem~6.4.1 (p.~33), proves an actual $R=T$ theorem with
\[
 I/I^2\simeq
 \mathbf Z_p/(\ell_0-1)\oplus\mathbf Z_p/(\ell_1-1).            \tag{D2.6}
\]
Its hypotheses are $N=\ell_0\ell_1$, both
$\ell_i\equiv1\pmod p$, Atkin--Lehner signs $(-1,-1)$, and
neither prime being a $p$th power modulo the other.  Its
deformation problem is the rank-two residual
pseudorepresentation $1\oplus\omega$, with determinant
$\chi_p$ and the weight-two finite-flat condition at $p$.
Here $p\nmid N$ follows from $\ell_i\equiv1\pmod p$.
The proof computes the conormal module and verifies the
Wiles numerical criterion; Proposition~5.2.1 (p.~30) explains
that numerical comparison.  This supplies a genuine example
where one generator is unnecessary.  It supplies no such
conormal calculation for $1\oplus\omega^e\oplus\omega^{-e}$.
All pairwise ratios in our triple have order five, whereas the
ratio in this theorem is $\omega$, of even order $p-1$; twisting
or permuting the constituents does not identify the problems.

Berger--Klosin, \emph{Modularity of residual Galois extensions
and the Eisenstein ideal}, Trans. AMS 2019, accepted author
version, Proposition~3.9 (p.~7) and Theorem~4.8 (p.~10), also
allow a noncyclic extension space.  Under the stated matched
divisible-Selmer/congruence bound and principality of the
individual residual-branch ideals, nonprincipality of the
global Eisenstein ideal forces \emph{more} modular extension
lines than the dimension of the Selmer group.
Proposition~4.6 (p.~10) proves the needed branch principality
when the opposite extension space is one-dimensional.
This is not a vanishing criterion for either extension space.

\begin{lemma}[KT-GMA-MULTIGENERATOR-CONORMAL]
Let $A$ be a complete Noetherian local ring and let the
three-block generalized matrix algebra, trace $T$, and
anti-involution satisfy Assumption~2.7 of Berger--Klosin,
arXiv:2611.19007v1.  In particular the anti-involution
interchanges blocks one and three, fixes block two, and all
off-diagonal modules $A_{ij}$ are finite over $A$.
Write $A'_{ij}=A_{ik}A_{kj}$ for the remaining index $k$, and
let $I\subset\mathfrak m_A$ be the total reducibility ideal.
Assume
\[
 A_{12}=A'_{12},\qquad
 A_{13}/A'_{13}\text{ is cyclic},\qquad
 A_{32}/A'_{32}\text{ is generated by at most }d\text{ elements}.
                                                               \tag{D2.7}
\]
Put $B=A_{21}/A'_{21}$ and $C=A_{31}/A'_{31}$.
For any $A$-module $M$ for which the right side is finite,
\[
 \#\operatorname{Hom}_A(I/I^2,M)
 \leq\#\operatorname{Hom}_A(C,M)\,
       \#\operatorname{Hom}_A(B,M)^d.                         \tag{D2.8}
\]
If $A/I\simeq\mathcal O$ is a complete discrete valuation
ring with finite residue field and
$\operatorname{Hom}_A(B,\mathcal O)=
 \operatorname{Hom}_A(C,\mathcal O)=0$, then $I/I^2$ is finite.
For every $n$ annihilating it by $\varpi^n$, (D2.8) applies
with $M=\mathcal O/\varpi^n$ and with its left side replaced
by $\#I/I^2$.
\end{lemma}

\begin{proof}
This is a generator-count variant of the complete proof of
Proposition~2.12 and Corollary~2.16 in the cited preprint
(pp.~7--9).  Put
\[
 N=A_{31}/A_{32}A_{23}A_{31},\qquad
 K=A'_{31}/A_{32}A_{23}A_{31}.                                \tag{D2.9}
\]
The product identities of Lemma~2.9, Remark~2.10, and
Lemma~2.11, using
$A_{12}=A'_{12}$, give $I=T(A_{13}A_{31})$.
Multiplication gives the surjections
\[
 (A_{13}/A'_{13})\otimes_A N\longtwoheadrightarrow I/I^2,
 \qquad
 (A_{32}/A'_{32})\otimes_A B\longtwoheadrightarrow K.          \tag{D2.10}
\]
These are precisely the well-defined product maps (2.7) and
(2.10) of that proof; their construction does not require
cyclicity of $A_{32}/A'_{32}$.
Choose a cyclic generator for the first factor of the first
map, and at most $d$ generators for the first factor of the
second.  We obtain $N\twoheadrightarrow I/I^2$ and
$B^{\oplus d}\twoheadrightarrow K$.
The definition gives $0\to K\to N\to C\to0$.
Applying $\operatorname{Hom}_A(-,M)$ therefore gives an
injection from $\operatorname{Hom}_A(I/I^2,M)$ into
$\operatorname{Hom}_A(N,M)$ and an exact sequence
\[
 0\longrightarrow\operatorname{Hom}_A(C,M)
 \longrightarrow\operatorname{Hom}_A(N,M)
 \longrightarrow\operatorname{Hom}_A(B,M)^{\oplus d}.          \tag{D2.11}
\]
Exactness is asserted at the first two terms; the last arrow
need not be surjective.  This proves (D2.8).
Taking $M=\mathcal O$ proves that
$\operatorname{Hom}_{\mathcal O}(I/I^2,\mathcal O)=0$.
Since $I/I^2$ is finite over $A/I=\mathcal O$, it is torsion,
and hence finite.  The elementary divisor theorem identifies
its cardinality with that of its dual into
$\mathcal O/\varpi^n$ whenever $\varpi^n$ annihilates it.
\end{proof}

\paragraph{The exact extra arithmetic inequality.}
Suppose, conditionally, that the full deformation and Hecke
setup of the cited three-character preprint is satisfied,
except that the residual Selmer bound in (S2)(c) is replaced
by dimension at most $d$.  Retain (S2)(a), (S2)(b), the
finite divisible groups in (S1), the ordinary GMA injections
of Theorem~6.4, and all the Hecke and local hypotheses.
The preceding lemma replaces the conormal estimate in the
proof of Theorem~10.2 by
\[
 \operatorname{length}_{\mathcal O}(I/I^2)
 \leq\operatorname{length}_{\mathcal O}(H_{\rm long})
       +d\operatorname{length}_{\mathcal O}(H_{\rm adj}),       \tag{D2.12}
\]
where, in the notation of that theorem, and with
$V=\operatorname{Frac}(\mathcal O)/\mathcal O$,
\[
 \begin{aligned}
 H_{\rm adj}&=H^1_{\rm Gr}
   (G_\Sigma,V(\Psi_k^{-1}\chi_p)),\\
 H_{\rm long}&=H^1_{\rm Gr}
   (G_\Sigma,V(\Psi_k^{-1-c}\chi_p^2))^{c=-1}.
 \end{aligned}                                                \tag{D2.13}
\]
Consequently its Theorem~7.8 still yields $R=T^0\times_{T^0/J^0}
\mathcal O$ if one independently proves
\[
 \operatorname{length}_{\mathcal O}(H_{\rm long})
       +d\operatorname{length}_{\mathcal O}(H_{\rm adj})
 \leq\operatorname{length}_{\mathcal O}(T^0/J^0)<\infty.       \tag{D2.14}
\]
This is a conditional numerical criterion, not an application
to the actual KTOWER source.  For $d=2$ the adjacent Selmer
length is counted twice.  The original proof's two $L$-value
factors bound one copy of each divisible group; they do not
automatically establish (D2.14) with $d=2$.
Even a residual upper bound of two would not bound either
integral length or compute the congruence module.

\paragraph{Remaining boundary.}
For the present Tate triple we have not established the
source's ordinary-flag identification, its other residual
Selmer assumptions, a two-generator bound for the entire
adjacent ordinary group, or (D2.14).  Nor has a saturated
integral comparison been proved between the original odd
off-diagonal module and a trace algebra.  The order-five
class-component estimate and its $p\mid e$ consequence do
not supply these missing statements.  Thus this bounded
audit gives an explicit replacement inequality and records
genuine noncyclic comparison theorems, but removes no
remaining prime, character, or source point.

\paragraph{Source and review scope.}
The new algebraic proof above was checked against the whole
proof of Berger--Klosin 2026 Proposition~2.12, Corollary~2.16,
Theorem~7.8, and Theorem~10.2.  The cited 2015 proofs of
Lemma~4.4, Theorem~6.3, Proposition~7.1, and
Theorems~8.2/8.5, and the stated setup and assumptions, were
read; their imported modularity, lattice, and numerical
theorems remain published inputs.  The complete short
proofs of Wake--Wang-Erickson Proposition~5.2.1 and
Theorem~6.4.1 and Berger--Klosin 2019 Proposition~3.9 and
Theorem~4.8 were read for scope.  No claim is made to
reprove those papers or to have kernel-checked their
arguments.  The primary PDFs, exact retrieval hashes, and
audit records are retained separately.

% END evidence/unitary_audit/dimension2_comparison/dimension_two_comparison.tex


% BEGIN evidence/unitary_audit/dimension2_comparison/adjacent_factor_attempt.tex
\subsection{The adjacent congruence factor: a final bounded source check}

This note checks whether the polarization or the integral congruence
quotient in Berger--Klosin, arXiv:2611.19007v1, supplies the additional
adjacent length required by (D2.14) when $d=2$.  The answer supplied by
the stated results is negative: their inequalities count one adjacent
group and one long-root group.  This is a precise failure of the
attempted application, not an impossibility theorem for a stronger
comparison.  No remaining prime, character, or source point is removed.

\paragraph{The two displayed adjacent ideals are the same ideal.}\textbf{KT-ADJACENT-FACTOR-SCOPE.} 
In Assumption~2.7 and Lemma~2.9 (printed pp.~6--7), the involution
interchanges the outer blocks and fixes the middle block.  The complete
lemma and Remark~2.10 give
\[
 A_{12}\simeq A_{23},\qquad A_{21}\simeq A_{32},\qquad
 J_{\rm adj}:=T(A_{12}A_{21})=T(A_{23}A_{32}).                 \tag{AF.1}
\]
Consequently the total reducibility ideal is
\[
 I=J_{\rm adj}+T(A_{13}A_{31})+J_{\rm adj}
   =J_{\rm adj}+T(A_{13}A_{31}).                              \tag{AF.2}
\]
The addition here is ideal addition.  Repeating the ideal does not
square it, produce another congruence quotient, or double a length.
The multiplication identities in Remark~2.10 transport this same ideal.
Theorem~10.2 (p.~54) likewise uses complex conjugation to identify the
two adjacent Greenberg groups; it does not exhibit their independent
occurrence in a congruence module.

Nor does this anti-involution make the adjacent product alternating.
For an explicit algebraic check let $\mathcal O$ be a discrete valuation
ring with uniformizer $\pi$, and take the matrix order
\[
 \mathcal E=
 \begin{pmatrix}
 \mathcal O&\pi^2\mathcal O&\pi\mathcal O\\
 \pi\mathcal O&\mathcal O&\pi^2\mathcal O\\
 \pi^2\mathcal O&\pi\mathcal O&\mathcal O
 \end{pmatrix},\qquad \tau(X)=JX^tJ,\quad J_{i,4-i}=1.        \tag{AF.3}
\]
The entrywise valuation inequalities show that this is an algebra;
$\tau$ preserves it and satisfies Assumption~2.7.  Here
$A_{12}=A'_{12}$, $A_{13}/A'_{13}=\pi\mathcal O/\pi^4\mathcal O$
is cyclic, and $I=\pi^3\mathcal O$.  Nevertheless, for
$b=\pi E_{21}$,
\[
 \tau(b)b=\pi^2E_{31}\ne0
 \quad\hbox{in}\quad
 A'_{31}/(A_{32}A_{23}A_{31})
       =\pi^2\mathcal O/\pi^5\mathcal O.                    \tag{AF.4}
\]
Thus these polarization identities alone do not replace the tensor
product used in (D2.10) by an exterior square.  This example is only a
GMA calculation; it is not asserted to be a global deformation ring
or to satisfy the augmentation condition $A/I\simeq\mathcal O$.
In the paper's representation argument, Proposition~6.7 and
Lemma~6.8 (pp.~35--36) assert that $\tau$ acts identically on the
long block after the specified reduced base change.  Even such an
identity action is consistent with (AF.4); it supplies no
alternating-product identity on the integral conormal module.
There is also a separate sign issue in the displayed proof's
$M_2\times F$ case, recorded below rather than imported without
qualification.

\paragraph{The integral lower bound contains two factors, not three.}
Write, at an allowed weight $k$ of the cited paper,
\[
 A_k=L_p^\Sigma(0,\Psi_k),\qquad
 B_k=L_p^\Sigma(-1,(\Psi_+)_k\eta),\qquad
 \Psi_+=\Psi\circ\operatorname{Ver},                         \tag{AF.5}
\]
where $\eta$ is the quadratic character of the imaginary quadratic
field.  Equations~(8.3), Theorem~8.16 (p.~45), and (10.6) (p.~55)
give the actual integral surjection
\[
 T^0_k/J^0_k\longtwoheadrightarrow\mathcal O/(A_kB_k).        \tag{AF.6}
\]
Assumption (S3) and (10.7) separately give
\[
 \begin{aligned}
 a_k:=\operatorname{length}_{\mathcal O}H_{\rm adj}
       &\leq v_\pi(A_k),\\
 b_k:=\operatorname{length}_{\mathcal O}H_{\rm long}
       &\leq v_\pi(B_k),\\
 c_k:=\operatorname{length}_{\mathcal O}(T^0_k/J^0_k)
       &\geq v_\pi(A_k)+v_\pi(B_k).
 \end{aligned}                                               \tag{AF.7}
\]
Here finiteness and nonzero values are retained when writing finite
lengths.  These inequalities establish $c_k\geq a_k+b_k$.
They do not establish $c_k\geq2a_k+b_k$, the sufficient inequality
obtained from (D2.14) for two generators.  For example the numbers
$a_k=v_\pi(A_k)=c_k=1$ and $b_k=v_\pi(B_k)=0$ satisfy every
inequality in (AF.7) and fail the proposed conclusion.  This example
only checks the logical content of the bounds; it claims no arithmetic
realization.  A conjugate copy of (AF.6) does not yield a surjection
onto a quotient by $A_k^2B_k$ without an additional independence or
ideal-product statement.

Lemma~10.6 and its complete proof (pp.~57--58) do not supply that
statement.  Its second adjacent bound follows from its first by
conjugation.  Its Iwasawa specialization and local comparison give
the one-copy upper bounds, not another lower bound for the congruence
quotient.  Surplus valuation in either factor could in principle make
the stronger inequality true at some weight.  No such surplus for the
actual Tate ratios is calculated by these results.

\paragraph{The actual Tate substitution has separate scope barriers.}
Set $\chi=\chi_p$ and consider the actual characteristic-zero
triple $(\chi^e,1,\chi^{-e})$ over $G_E$, with $e>0$ even.
Matching the paper's characters
$(\Psi_k\chi^{-1},1,\Psi_k^{-c}\chi)$ gives
\[
 \Psi_k=\chi^{e+1},\qquad
 \Psi_k^{-1}\chi=\chi^{-e},\qquad
 \Psi_k^{-1-c}\chi^2=\chi^{-2e},\qquad
 (\Psi_+)_k=\chi^{2e+2}.                                    \tag{AF.8}
\]
The last equality is the transfer identity: the restriction to $G_E$
is the square of the parallel character, and at complex conjugation
both sides are one.  The long-root $c=-1$ group therefore descends
to a group over $\mathbf Q$ with coefficient
$\eta\chi^{-2e}$.  It is not another copy of the adjacent
$\chi^{-e}$ group.  For a residual Tate ratio of order five,
the powers $\omega^{-e}$ and $\omega^{-2e}$ remain distinct.
The order-five class bound and the $p\mid e$ rational strict exponent
bound do not identify these entire ordinary groups or their lengths.

Theorem~10.2 is set up in Section~10 (pp.~51--53) for infinity type
$(a+b,-b)$ with $a\geq3$ and $b\geq-1$.  The first substitution
in (AF.8) has parallel type $(e+1,e+1)$ and would require
\[
 (a,b)=(2e+2,-e-1).
                                                               \tag{AF.9}
\]
It fails $b\geq-1$.  Swapping the outer characters requires
$\Psi_k=\chi^{1-e}$ and $(a,b)=(2-2e,e-1)$, which fails
$a\geq3$.  Moreover the ordinary orders in (10.1) and (10.2)
are respectively $(\tau_1,\tau_3,1)$ and $(1,\tau_1,\tau_3)$
at the two primes over a split $p$.  Our existing source construction
does not supply these two flags.  A locally split representation may
have several flags, so no absolute impossibility is claimed.

The local conditions cannot simply be left implicit in changing the
weight.  Definition~8.22 and Example~8.23 (p.~48) impose the adjacent
inertia condition at one prime in the paper's weight cone, with the
conjugate condition at the other prime for the conjugate character.
For the parallel coefficient $\chi^{-e}$, both weights are negative,
so Definition~8.22 imposes the inertia condition at both primes.
Definition~8.24 treats the special boundary $b=-1$ and does not cover
(AF.9).  No equality of the relevant integral Selmer or conormal
modules after this change has been established.

Theorem~8.16 itself must be distinguished from the earlier unmodified
Hsieh statement.  Its proof explicitly uses Finis and He to remove
one additional finite-character ramification assumption and derives
the stated $\mu$ condition from $p\nmid\#\Delta_{\mathfrak c}$.
Thus the older Hsieh parity/ramification exclusion should not be
silently applied to this enlarged theorem.  The failed Section~10
weight and local identifications above, and the absent second
adjacent congruence bound, remain separate obstacles.  This audit
does not claim a cyclotomic factorization of the specialized Katz
value, whose interpolation, periods, and local factors would need
their own comparison.

\paragraph{A source normalization error and its ramification repair.}\textbf{KT-BK2026-NORMALIZED-DEPTH-REPAIR.} 
The PDF at printed p.~50, equation~(9.3), visibly writes
$\#(M\otimes_{\mathcal O}\mathcal O')=[F':F]\,\#M$ for a
finite $\mathcal O$-module $M$.  With $d=[F':F]=e'f'$ the correct
formula is
\[
 \#(M\otimes_{\mathcal O}\mathcal O')=(\#M)^d,\qquad
 \operatorname{length}_{\mathcal O'}(M\otimes\mathcal O')
       =e'\operatorname{length}_{\mathcal O}M.                \tag{AF.10}
\]
These follow either from the elementary divisors or from
$\mathcal O'\simeq\mathcal O^d$ and
$\pi\mathcal O'=(\pi')^{e'}\mathcal O'$.
The error also appears in the immediately following sentence.

Here is the uniform repair of the numerical depth in Proposition~9.3,
retaining that proposition's imported Galois, ordinary, sign, and
lattice inputs.  Their production and Selmer placement of primitive
classes are not proved by this numerical repair.  Its free Hecke
rank $r$ is fixed.  In Lemma~9.2 choose
$N=rn$.  Write $q=\#k$ and $q'=q^{f'}$.  The product inequality
in (9.3), together with $\#M\geq q^{rn}$, then gives
\[
 (q')^{\sum_i m_i}\geq q^{e'f'rn}
       =(q')^{e'rn},\qquad \max_i m_i\geq e'n.               \tag{AF.11}
\]
Changing the weight back to $k$ changes its characters only at depth
$e'rn$, so the chosen congruence retains depth $e'n$ over
$\mathcal O'$.  After any further field of definition $F''/F'$
used for the representation and its stable lattice, it retains depth
\[
 e''e'n=e(F''/F)n.                                          \tag{AF.12}
\]
Use the same lattice lemma at that retained exponent.  Its primitive
extension class then has order exponent $e(F''/F)n$ with respect to
the uniformizer of $\mathcal O''$.

This really contradicts a fixed finite divisible Selmer group, even
when $F''$ depends on $n$.  Let $C$ annihilate all the finitely many
base Selmer groups in Proposition~9.3 by $\pi^C$.  Finite free
coefficient extension commutes with continuous cochain complexes,
their cohomology, and the kernels defining the inertia conditions:
after choosing an $\mathcal O$-basis of $\mathcal O''$, these are
just finite direct sums.  Also
$(F/\mathcal O)\otimes_{\mathcal O}\mathcal O''
\simeq F''/\mathcal O''$; the $c=-1$ projector commutes with
this extension since $2$ is invertible.  Thus every extended group
is killed by $\pi^C$, or equivalently by a uniformizer power
$e(F''/F)C$.  The primitive class of depth (AF.12) contradicts this
as soon as $n>C$.  The paper's finite-coefficient classes map
injectively to the divisible groups because its residual character
distinctions give $H^0(G_\Sigma,F''/\mathcal O''(\tau))=0$.
These distinctions are imposed at one element of the decomposition
group in Definition~8.19 (p.~46).
This repairs the numerical normalization of the contradiction;
no independent second adjacent factor arises from coefficient
extension.

\paragraph{A separate sign gap in the displayed algebraic argument.}\textbf{KT-BK2026-LONG-ROOT-SIGN-SCOPE.} 
The inference in Lemma~6.8 that $\tau^2=1$ forces the implementing
matrix to be symmetric is valid for the full odd-dimensional matrix
algebra.  It is not valid for its stated third case, the outer
$M_2(F)$ block plus the middle $F$ block.  To see the precise issue,
over a field $F$ of characteristic different from two put
\[
 \mathcal B=
 \left\{\begin{pmatrix}a&0&b\\0&t&0\\c&0&d\end{pmatrix}
           :a,b,c,d,t\in F\right\},\qquad
 C=\begin{pmatrix}0&0&1\\0&1&0\\-1&0&0\end{pmatrix},
 \qquad \tau_C(X)=CX^tC^{-1}.                                \tag{AF.13}
\]
Then $\mathcal B\simeq M_2(F)\times F$, $\tau_C$ is an
anti-automorphism, and
\[
 \begin{gathered}
 C C^{-t}=\operatorname{diag}(-1,1,-1),\qquad
 \tau_C^2|_{\mathcal B}=1,\\
 \tau_C(E_{11})=E_{33},\qquad \tau_C(E_{22})=E_{22},\qquad
 \tau_C(E_{13})=-E_{13},\quad \tau_C(E_{31})=-E_{31}.
 \end{gathered}                                               \tag{AF.14}
\]
The diagonal matrix in (AF.14) lies in the centralizer of
$\mathcal B$ but is not scalar on $F^3$.  The determinant
pseudorepresentation is invariant under $\tau_C$.  Thus
determinant self-duality and outer-idempotent interchange alone
do not force the asserted positive action on the long block.
The definitions of self-duality in Section~4, including
Definitions~4.1/4.8 and Corollary~4.10, and the setup of
Section~6.1 impose those identities; no additional common
symmetric-polarization sign was found in those clauses.

This is a counterexample to the displayed algebraic inference,
not an arithmetic counterexample to Proposition~6.7 or the main
theorem.  The field-algebra model does not assert realization by
the paper's global Galois group, ordinary deformation ring, or
automorphic representations.  A prescribed symmetric ambient
polarization, an appropriate conjugate-orthogonal sign on every
irreducible block, or restriction to a full $M_3$ generic branch
would exclude it.  Supplying such a condition and proving it in
the intended application is a separate repair.  In particular,
the use of the ordinary GMA injections of Theorem~6.4 in (D2.12)--
(D2.14) remains an explicit hypothesis; this audit has not
discharged it for the actual Tate source.

\paragraph{Outcome and audit scope.}
The exact conditional source boundary in (D2.14) is unchanged.
The 2026 preprint supplies one adjacent factor and one long-root
factor in its own ordinary weight range.  Neither its polarization
identities nor its base-change cardinality argument supplies the
extra adjacent length for the actual Tate-ratio problem.  The
remaining order-five or order-at-least-seven gate stays open.
The complete cited polarization proofs, Proposition~6.7/Lemma~6.8,
Theorem~8.16, Section~8.5 definitions, Proposition~9.3,
Theorem~10.2, and Lemma~10.6 were read for this check.  Their imported
theorems are retained as inputs, not independently reproved here.
The PDF equation~(9.3) was checked visually.  The sign issue is
confined to the identified algebraic inference, with omitted
arithmetic hypotheses stated explicitly.  This is a source-scope
and normalization audit, not a Lean kernel verification or a new
automorphic vanishing result.

% END evidence/unitary_audit/dimension2_comparison/adjacent_factor_attempt.tex


% BEGIN evidence/unitary_audit/reducible_comparison/reducible_comparison_results.tex
\section{Residually reducible comparison theorems: exact applicability}
\label{sec:kt-reducible-comparison}

This audit concerns the integral extension module at the residual
point $\bar\rho=\bar\chi^e\oplus1\oplus\bar\chi^{-e}$.
Here $e>0$ is even, $\chi=\varepsilon_\ell$, $\ell>3$, and the
carried GD condition gives $(\ell-1)\nmid2e$.  The desired module
is the odd weight $-1$ cotangent identified with the original
rational strict Selmer dual in KT-ODD-CONORMAL-SELMER-EQUIVALENCE.
The conclusions below are an applicability audit, not a new
upper bound for that Selmer group.

\subsection{What the general matrix algebra results actually control}

Bella\"iche--Chenevier, \emph{$p$-adic families of Galois
representations and higher rank Selmer groups}, arXiv
\path{math/0602340v2}, 12 January 2007, is the exact author
version used here.  Its pagination is not that of the final
Ast\'erisque volume.  Proposition~1.5.1, pp.~32--34, applies
to any residually multiplicity-free pseudocharacter.  Thus
three distinct rank-one constituents cause no obstruction.
For a Cayley--Hamilton quotient with structural modules $A_{ij}$,
put
\[
 A'_{ij}=\sum_{k\ne i,j}A_{ik}A_{kj},\qquad
 I=\sum_{i\ne j}T(A_{ij}A_{ji}).                         \tag{RC.1}
\]
The total reducibility ideal $I$ is independent of that
Cayley--Hamilton quotient.  Theorem~1.5.5, p.~35, gives
\[
 \operatorname{Hom}_A(A_{ij}/A'_{ij},A/J)
       \lhook\joinrel\longrightarrow
       \operatorname{Ext}^1_{R/JR}(\rho_j,\rho_i),
       \qquad J\supset I.                                \tag{RC.2}
\]
Theorem~1.5.6, p.~36, identifies the image with extensions
factoring through the chosen Cayley--Hamilton quotient.
Remark~1.5.8, p.~38, explicitly distinguishes these extension
groups from those of the faithful trace quotient: their
inclusion can be strict even though the reducibility ideal
is unchanged.  Consequently a comparison of trace rings
alone supplies neither an identification of the required
off-diagonal module nor its integral saturation.

\begin{lemma}[KT-GMA-REDUCIBILITY-DOES-NOT-BOUND-EXTENSIONS]
For every integer $n\geq1$ there is a polarized
Cayley--Hamilton generalized matrix algebra of type $(1,1,1)$
over a complete discrete valuation ring $O$, with base
algebra $O$ and total reducibility ideal zero, whose two
adjacent off-diagonal modules have length $n$.
\end{lemma}
\begin{proof}
Put $T_n=O/(\varpi^n)$ and
$R_n=O\oplus\epsilon T_n$, with $\epsilon^2=0$.
Inside $M_3(R_n)$ take
\[
 \mathcal E_n=
 \left\{
 \begin{pmatrix}
 a&0&0\\ \epsilon u&b&0\\0&\epsilon v&c
 \end{pmatrix}:
 a,b,c\in O,\ u,v\in T_n\right\}.                       \tag{RC.3}
\]
It is closed under multiplication; every product of two
off-diagonal structural modules is zero.  Diagonal
evaluation defines a degree-three determinant law, with
characteristic polynomial $(X-a)(X-b)(X-c)$.
Cayley--Hamilton holds in the displayed matrix algebra,
also after every scalar extension, since the same direct
sum description embeds in the corresponding matrix
algebra after base change.  The anti-automorphism
$x\mapsto Jx^tJ$, with $J$ anti-diagonal, exchanges
$e_1,e_3$ and $A_{21},A_{32}$ and fixes $e_2$.
Thus the polarization has the required outer exchange.
Formula (RC.1) gives $I=0$, while
$A_{21}/A'_{21}\simeq A_{32}/A'_{32}\simeq T_n$.
The coefficient algebra $O$ is itself flat and a complete
intersection.  This is an algebraic scope example; it is
not asserted to satisfy the global automorphic source
conditions.
\end{proof}

\subsection{The earlier two-constituent comparisons}

Berger--Klosin, \emph{On deformation rings of residually
reducible Galois representations and $R=T$ theorems},
author version \texttt{hd32.pdf}, dated 10 February 2011,
uses two nonisomorphic absolutely irreducible residual
blocks and a nonsplit residual extension.  Assumption~5.1,
p.~13, imposes a one-dimensional extension Selmer group
and rigid diagonal deformation rings.  Theorem~8.5,
p.~23, requires an upper bound for the divisible extension
Selmer group in addition to a lower bound for the Hecke
congruence quotient.  The alternative Theorem~8.6 assumes
the needed conormal inequality.  Grouping two characters
of our residual triple into one block does not satisfy
absolute irreducibility.  Even for a two-character
subproblem these hypotheses do not supply the missing
upper bound: that bound is an input.

Their \emph{On modularity of reducible residual Galois
representations}, author version
\texttt{bergerklosinonlifting21.pdf}, dated 2 April 2013,
retains two absolutely irreducible blocks.  Assumption~2.7,
p.~6, is explicitly the required Bloch--Kato upper bound;
Assumption~4.2 and Lemma~4.4, p.~12, use it together with
a congruence lower bound to prove surjectivity onto the
Selmer group.  It cannot be removed when using that
surjectivity.  The construction also passes to the image
of the reduced ring in its localization at $p$; the
introduction, pp.~1--3, explains why characteristic-$p$
components can be lost.  This does not identify our full
integral cotangent before that passage.

Urban's \emph{Selmer Groups and the Eisenstein--Klingen
Ideal}, arXiv \path{math/9803171v1}, Theorem~1.1,
author p.~6, does provide a useful two-block stable
lattice.  Its nonsplit residual extension yields a
Selmer class.  This is a construction of classes, not
an upper bound for every class.  We use precisely this
positive consequence below.

Allen--Newton--Thorne, \emph{Automorphy lifting for
residually reducible $l$-adic Galois representations, II},
Theorem~1.1 (also Theorem~6.1), allows more constituents,
but its conditions (iv) and (vii) fail here.  After the
rank-three normalization by $\bar\chi^{-1}$, individual
polarization of either outer constituent requires
$\bar\chi^{2e}=1$, contrary to GD.  All three characters
become equal on $G_{E(\zeta_\ell)}$, contrary to (vii).
The full triple's polarization exchanging outer
constituents does not repair either condition.  The
independent source audit is recorded in
\path{lean_chain/thorne_scope/THORNE_SCOPE.md}.

\subsection{The three-character ordinary unitary preprint}

Berger--Klosin, \emph{On $p$-adic modularity of abelian
threefolds with imaginary multiplication and rational
$p$-isogeny}, arXiv \texttt{2611.19007v1}, 13 November
2026, is a recent preprint, not a claimed published
theorem.  The official arXiv PDF and the separately
retrieved mirror are byte-identical.  Its general
Proposition~2.12 and Corollary~2.16, pp.~7--9, bound
$I/I^2$ using off-diagonal modules under vanishing and
cyclicity hypotheses.  The direction of this estimate
does not bound those modules from $I/I^2$.

Section~10, pp.~51--56, is directly relevant but has
different ordinary data.  Write $\Psi$ for its
specialized character.  Its outer characters are
$\tau_1=\Psi\chi^{-1}$ and
$\tau_3=\Psi^{-c}\chi$, with middle constituent $1$.
Matching our characteristic-zero triple gives just
the following two possibilities:
\[
 \begin{array}{c|c|c}
 \tau_1&\Psi\text{ infinity type }(a+b,-b)&(a,b)\\ \hline
 \chi^e&(e+1,e+1)&(2e+2,-e-1)\\
 \chi^{-e}&(1-e,1-e)&(2-2e,e-1).
 \end{array}                                                \tag{RC.4}
\]
The first violates $b\geq-1$, the second $a\geq3$.
These are exact characteristic-zero weights, so reducing
the exponents modulo $\ell-1$ does not resolve the issue.

There is a separate local flag mismatch.  Equations
(10.1)--(10.2), p.~52, prescribe the upper-triangular
orders $(\tau_1,\tau_3,1)$ and $(1,\tau_1,\tau_3)$ at
the two split primes.  Our component has order
$(\chi^e,1,\chi^{-e})$ at both.  An identification
respecting polarization fixes the unique self-dual
constituent $1$ and can at most exchange the outer two;
it cannot identify these ordered functors.  A locally
split representation can have several flags, so this
observation does not assert that every individual point
fails the preprint's flag condition.  Our construction
does not supply that additional pair of flags.

For clarity, the following substitution is a
\emph{hypothesis audit outside the theorem's weight
range}, not an application.  Substituting
$\Psi=\chi^{e+1}$ into the Selmer conditions of
Theorem~10.2, p.~53, gives
\[
 \begin{array}{c|c}
 \text{condition}&\text{coefficient character}\\ \hline
 (S2)(a)&\chi^e\\
 (S2)(b),\ c=-1&\eta\chi^{2e}\text{ over }\mathbf Q\\
 (S2)(c)&\chi^{-e}\\
 (S3)&(K_O/O)(\chi^{-e}).
 \end{array}                                                \tag{RC.5}
\]
The class component in Theorem~10.4(5) similarly has
character $\eta\bar\chi^{2e}$; it is not literally the
untwisted target.  Condition~10.4(6), however, retains
the same one-dimensional $\chi^{-e}$ Selmer assumption.
Assumption (T), p.~53, independently requires reduced
specializations of the relevant ordinary Hecke family.
No such family comparison or reducedness hypothesis has
been established for our ring.

Definition~8.22, p.~48, places inertia conditions at
\emph{both} split primes on a parallel negative character
$\chi^{-e}$.  The one-sided examples in that section
belong to the preprint's different infinity-type cone.
In either interpretation the target strict group injects
into the group in (RC.5).  The proof of Theorem~10.2,
pp.~53--55, uses its Selmer upper bounds to bound
$I/I^2$ and then applies a numerical comparison.
It does not extract those upper bounds from $R=T$.
The fixed prime $697446897311$ in the $e=26$ calculation
is congruent to $2$ modulo $3$, hence is inert in
$\mathbf Q(\sqrt{-3})$; the preprint's split-prime
assumption already excludes that particular seed.

\subsection{The cyclicity condition already detects the target}

\begin{theorem}[KT-RESIDUAL-CYCLICITY-SCOPE]
Let $E=\mathbf Q(\sqrt{-3})$, $\eta=\chi_{-3}$,
and assume that $\ell$ splits in $E$.  Let $f$ be one
of the non-CM ordinary forms supplied by the checked
lower-weight Eisenstein congruence: its two-dimensional
representation $\rho_f$ is absolutely irreducible,
\[
 \bar\rho_f^{\rm ss}=1\oplus\eta\bar\chi^e,
 \qquad a_\ell(f)\equiv1,
                                                               \tag{RC.6}
\]
and $\rho_f|_{G_E}$ is unramified away from $\ell$.
Let $k$ be its residue field.  Let $\Sigma$ be any
conjugation-stable set containing the primes over
$3\ell$, and set
\[
 \mathcal H_E=
 \ker\left(H^1(G_{E,\Sigma},k(\bar\chi^{-e}))
       \longrightarrow\prod_{v\mid\ell}
                         H^1(I_v,k(\bar\chi^{-e}))\right).
                                                               \tag{RC.7}
\]
Let $\mathcal X$ be the original rational strict
Selmer group with coefficients
$(\mathbf Q_\ell/\mathbf Z_\ell)(\chi^{-e})$,
including the decomposition-group condition at $2$.
Then
\[
 \dim_k\mathcal H_E^{c=-1}\geq1,\qquad
 \dim_k\mathcal H_E^{c=+1}
          \geq\dim_{\mathbf F_\ell}\mathcal X[\ell].
                                                               \tag{RC.8}
\]
In particular, $\dim_k\mathcal H_E=1$ implies
$\mathcal X=0$.  The same implication holds with any
relaxation of the local conditions in (RC.7).
Thus an attempted use of the cyclicity hypothesis in
(RC.5) in these concrete branches would already require
a condition forcing the original vanishing problem.
\end{theorem}

\begin{proof}
Apply Urban's Theorem~1.1 with coefficient ring the
valuation ring of $f$, ideal its maximal ideal, and
ordered rank-one blocks $1$, $\eta\chi^e$.
Integrality and congruence of characteristic polynomials
follow from (RC.6).  Absolute irreducibility and
distinctness of the residual blocks are part of the
hypotheses, the latter also following from GD at
inertia over $\ell$.  The theorem supplies a stable
lattice whose reduction is a nonsplit extension
\[
 0\longrightarrow k\longrightarrow\bar L
       \longrightarrow k(\eta\bar\chi^e)
       \longrightarrow0.                                  \tag{RC.9}
\]
Urban's auxiliary module is a nonzero finitely generated
submodule of the coefficient field, hence is free of
rank one over the valuation ring; its reduction is
exactly one copy of $k$ as in (RC.9).
Indeed its assertion that the lattice has no quotient
isomorphic to the first residual block rules out a
split reduction.  This gives a nonzero class
$u\in H^1(G_{\mathbf Q},k(\eta\bar\chi^{-e}))$.

Ordinarity gives a characteristic-zero local stable
line whose character reduces to $\eta\bar\chi^e$:
the ordinary quotient reduces to $1$ because
$a_\ell(f)\equiv1$.  Its intersection with any stable
lattice is saturated.  In (RC.9) its reduction and the
displayed stable line $k$ are distinct, since their
inertial characters differ.  They therefore split
(RC.9) on the full decomposition group at $\ell$.
Thus $u$ is locally strict there, not merely
unramified.  Away from $3\ell$ it is unramified.
At three, restriction to $G_E$ is unramified by the
stated property of $\rho_f$, independently of lattice.
Consequently $u|_{G_E}$ belongs to (RC.7).

Restriction through the index-two subgroup $G_E$ is
injective because two is a unit.  Under the
identification of the restricted coefficient character
with $\bar\chi^{-e}$, the additional factor $\eta$
places its image in the $c=-1$ eigenspace.  Hence the
first inequality in (RC.8).

Put $M=(\mathbf Q_\ell/\mathbf Z_\ell)(\chi^{-e})$.
The global and local $H^0$ groups at $\ell$ vanish:
GD supplies an inertia element on which
$\bar\chi^{-e}-1$ is a unit.  The exact sequence
$0\to\mathbf F_\ell(\bar\chi^{-e})\to M
\xrightarrow{\ell}M\to0$ therefore identifies
global $\ell$-torsion classes with finite-coefficient
classes and preserves their strict local vanishing
at $\ell$.  Away from $\ell$, the action on inertia
is trivial, so $H^0(I_q,M)=M$ is $\ell$-divisible.
The same long exact sequence therefore makes
$H^1(I_q,\mathbf F_\ell(\bar\chi^{-e}))\to H^1(I_q,M)$
injective, preserving unramifiedness.  Restricting these finite classes to
$E$ and extending scalars to $k$ gives an injection
\[
 \mathcal X[\ell]\otimes_{\mathbf F_\ell}k
                  \lhook\joinrel\longrightarrow
                         \mathcal H_E^{c=+1}.              \tag{RC.10}
\]
No vanishing of finite-coefficient unramified cohomology
at $2$ is asserted: the extra condition at $2$ is
simply part of the domain $\mathcal X$.
This proves the second inequality.  Since the two
eigenspaces are disjoint and $\ell\ne2$, dimension
one for (RC.7) forces $\mathcal X[\ell]=0$.
Every nonzero element of the $\ell$-primary torsion
group $\mathcal X$ has a nonzero multiple of order
$\ell$, so $\mathcal X=0$.
\end{proof}

The theorem uses the already checked existence,
ordinarity and irreducibility of the classical form;
it does not supply the cyclicity assumption.
The inspected comparison results therefore leave a
specific missing input: an integral comparison of
the relevant extension module, or an independent
upper bound for it, with the actual ordinary flags
and conjugation sign.  Neither a semisimple $R=T$
statement nor a reducibility-ideal estimate reverses
the injections in (RC.2).

% END evidence/unitary_audit/reducible_comparison/reducible_comparison_results.tex


% BEGIN evidence/unitary_audit/reducible_comparison/ordinary_rank2_exclusion.tex
\section{A Bernoulli unit excludes the ordinary rank-two source}
\label{sec:kt-ordinary-rank-two-unit}

This section proves the necessary ordinary direction
that was absent from the earlier field-variation
discussion.  It uses the strict Selmer vanishing in
KT-AUXILIARY-UNIT-SELMER-BRIDGE and a stable lattice,
not a converse to an Eisenstein congruence theorem.

\begin{theorem}[KT-ORDINARY-RANK2-UNIT-EXCLUSION]
Let $p>3$, let $D<0$ be a fundamental discriminant
prime to $p$, put $E=\mathbf Q(\sqrt D)$ and
$\eta=\chi_D$, and let $a>0$ be even with
$a\not\equiv0\pmod{p-1}$.  Fix a finite coefficient
field $K/\mathbf Q_p$, with valuation ring $O$ and
residue field $k$.  There is no continuous
absolutely irreducible representation
\[
                 r:G_{\mathbf Q}\longrightarrow
                                  \mathrm{GL}_2(K)        \tag{OR.1}
\]
satisfying all the following conditions if
$B_{1,\eta\omega^a}$ is a $p$-adic unit:
\begin{enumerate}
\item Its residual semisimplification is
      $1\oplus\bar\eta\bar\varepsilon_p^a$.
\item Its restriction to $G_{\mathbf Q_p}$ has a
      stable line whose rank-one quotient is
      unramified.
\item Its restriction to $G_E$ is unramified at
      every finite place away from $p$.
\end{enumerate}
The conclusion holds whether $p$ splits or is inert
in $E$, and it places no upper bound on $a$.
\end{theorem}

\begin{proof}
Write $\delta=\bar\eta\bar\varepsilon_p^a$.
The residual characters $1,\delta$ are distinct
already on inertia at $p$, since $p\nmid D$ and
$a\not\equiv0\pmod{p-1}$.
Apply Urban's stable-lattice theorem in the order
$\rho_1=1$, $\rho_2=\eta\varepsilon_p^a$, with ideal
the coefficient maximal ideal.  More precisely use
Theorem~1.1 of \emph{Selmer Groups and the
Eisenstein--Klingen Ideal}, arXiv
\path{math/9803171v1}, author p.~6.
The image of the continuous representation is
compact and admits a stable lattice, so its
characteristic polynomials are integral; their
residual congruence is hypothesis~1.  The two
residual blocks are absolutely irreducible and
distinct, and (OR.1) is absolutely irreducible.
All hypotheses of that theorem hold.

It gives a stable lattice $L$ with a nonsplit
residual extension
\[
 0\longrightarrow k\longrightarrow\bar L
       \longrightarrow k(\delta)\longrightarrow0.        \tag{OR.2}
\]
Its auxiliary fractional module is free of rank
one over the valuation ring, so the left term is
exactly one copy of $k$.  The theorem's condition
that $L$ has no quotient isomorphic to the first
residual block makes (OR.2) nonsplit.  Its class is
therefore nonzero in
\[
             H^1(\mathbf Q,k(\bar\eta\omega^{-a})).       \tag{OR.3}
\]

Let $V^+\subset K^2$ be the ordinary stable line
in hypothesis~2.  Its intersection with $L$ is
saturated, and the quotient lattice is torsion-free.
Hence reduction gives a stable line and its
rank-one quotient in $\bar L|_{G_{\mathbf Q_p}}$.
The quotient is unramified.  Of the two residual
characters $1,\delta$, only $1$ is unramified on
inertia at $p$, so this quotient has character $1$
and the saturated line has character $\delta$.
That line is distinct from the $k$-line in (OR.2).
Its map to the quotient $k(\delta)$ is consequently
an isomorphism, splitting (OR.2) on the full local
decomposition group.  Thus (OR.3) is strict at $p$.

For every $q\ne p$, hypothesis~3 makes
$I_q\cap G_E$ act trivially on $r$, hence also
on (OR.2).  The latter representation on $I_q$
factors through
\[
                     I_q/(I_q\cap G_E),                 \tag{OR.4}
\]
which has order at most two.  Since $p>2$,
averaging splits every extension over this finite
quotient.  Therefore the restriction of (OR.3)
to $H^1(I_q,k(\bar\eta\omega^{-a}))$ is zero.
This argument includes the ramified primes $q\mid D$:
it does not incorrectly treat $\eta|_{I_q}$ as
trivial there.

The class (OR.3) is now in precisely the residual
strict group of KT-AUXILIARY-UNIT-SELMER-BRIDGE,
with vanishing on $G_{\mathbf Q_p}$ and on every
away-prime inertia group.  That theorem makes
the group with coefficients
$\mathbf F_p(\bar\eta\omega^{-a})$ zero under the
Bernoulli-unit hypothesis.  Extending scalars to
$k$ preserves this statement: these character
values lie in $\mathbf F_p$, and a choice of
$\mathbf F_p$-basis of $k$ identifies the coefficient
module and all local maps with a finite direct
sum of the original ones.  Thus (OR.3) must be
zero, contradicting its construction.
\end{proof}

\begin{corollary}[KT-ORDINARY-RANK2-RESTRICTION-VARIANT]
The first hypothesis in the theorem may be replaced
by the two conditions
\[
 \bigl(\bar r|_{G_E}\bigr)^{\rm ss}
                  =1\oplus\bar\varepsilon_p^a,
 \qquad \det\bar r=\bar\eta\bar\varepsilon_p^a.          \tag{OR.5}
\]
\end{corollary}
\begin{proof}
The residual image is finite.  For each simple
module for that finite image, its socle on the
normal subgroup coming from $G_E$ is nonzero and
stable under the full group, hence is the whole
module.  Thus restriction of the semisimplified
residual representation to $G_E$ is semisimple.
The two characters on $G_E$ are distinct because
$p\nmid D$ makes $E$ disjoint from
$\mathbf Q(\mu_p)$, and the exponent is nonzero
modulo $p-1$.  They are invariant under conjugation
by $G_{\mathbf Q}$.  In the semisimplified residual
representation their individual eigenspaces are
therefore stable under $G_{\mathbf Q}$.  The two
characters over $\mathbf Q$ must be
$\bar\eta^b$ and
$\bar\eta^{1-b}\bar\varepsilon_p^a$ for some
$b\in\{0,1\}$, by the determinant condition.
Replace $r$ by $r\otimes\eta^b$.
It remains absolutely irreducible; its restriction
to $G_E$ is unchanged, and its ordinary quotient
remains unramified at $p$ because $p\nmid D$.
Its residual pair now has the form in the theorem,
which gives the contradiction.
\end{proof}

\paragraph{Arithmetic use and remaining scope.}
For the original even parameter $e$ with GD,
both $a=e$ and $a=2e$ satisfy the nontrivial
inertia hypothesis.  The corresponding auxiliary
units therefore exclude the ordinary rank-two
sources with these residual pairs and with the
stated unramified restriction to $E$.
Equivalently, any such source forces the relevant
$B_{1,\eta\omega^a}$ to be a nonunit.
The integrality used in that formulation is also
contained in the Stickelberger calculation of
the auxiliary-unit bridge, whose auxiliary
multiplier is a $p$-adic unit.

When the source is described by classical forms
of weight $a+1$ and quadratic nebentypus $\eta$,
the determinant in (OR.5) is the usual residual
determinant $\bar\eta\omega^a$.  The ordinary and
unramified-after-$E$ hypotheses must still be
matched to the actual source's local predicates;
they are explicit hypotheses of the result above.
No ordinary condition is silently imposed on
an arbitrary cuspidal form.  No generic stable
rank-three source is excluded by this theorem.

\begin{corollary}[KT-AUXILIARY-UNIT-NO-TYPEII]
Retain the carried source convention and take actual
source data $\mathcal D$ for the imaginary quadratic
field $E$, with $N=2$, $\psi=1$, even $e\ge2$, and
GD.  Let $S$ be its normalized global source ring,
let $f:S.R\to\overline{\mathbf Q}_p$ be an actual
\texttt{IsLocPt} specialization, and let
$r=\operatorname{push}_f(S.r)$.
If
\[
                     B_{1,\eta\omega^{2e}}
                              \in\mathbf Z_p^\times,     \tag{OR.6}
\]
then $r$ is not of Type~II.
This is an exclusion of actual source points;
the bare \texttt{Adm2A} predicate alone is not
being asserted to include residual Eisenstein
data or the determinant $\eta$.
\end{corollary}

\begin{proof}
First the point has finite coefficients.
The literal private theorem \path{fin_quot}
in \path{UnitaryThreeGroups/TanDet.lean:212--216}
proves that $S.R/(\ker f)^2$ is finite over the
base coefficient valuation ring.  Its quotient
$S.R/\ker f$ is therefore finite as well.  This
quotient embeds in $\overline{\mathbf Q}_p$, so
its fraction field is a finite extension of the
base coefficient field and contains every matrix
entry of $r$.  The hypotheses are exactly
\texttt{IsCO} and \texttt{IsLocPt}; no
torsion-freeness of the entire source ring or
automorphic descent is used.

Suppose $r$ has Type~II.  The carried
\path{typeII_facts} theorem supplies a
continuous absolutely irreducible representation
$W$ with $\det W=\eta$ and matrices
\[
 W(g)=
 \begin{cases}
  \operatorname{corner}(r(g).\mathrm{mat}),&\eta(g)=1,\\
  \operatorname{corner}(r(g).\mathrm{mat})
                   \operatorname{diag}(-1,1),&\eta(g)=-1.
 \end{cases}                                               \tag{OR.7}
\]
These matrices belong to the same finite
coefficient field.  Enlarge it by the finitely
many entries of a local flag matrix when needed.
The exact source locators are
\path{CoboundaryLemmaNoColumn.lean:135--141}
for the properties of $W$ and
\path{BaseRepGL3NormalisedOfPsiEqOne.lean:28--34}
for its matrix formula.

The entrywise residual assertion
\path{IsPtOf.isPt_push:124--145}, or the
\texttt{red} field in
\path{RhoKApply.lean:76--83}, says that
$r(g).\mathrm{mat}$ reduces to
$\operatorname{diag}(\bar\chi^e,1,\bar\chi^{-e})$.
Formula (OR.7) therefore gives the actual global
residual matrix
\[
       \bar W(g)=
           \operatorname{diag}(\bar\eta\bar\chi^e,
                                      \bar\chi^{-e}).     \tag{OR.8}
\]
Consequently $R=W\otimes\chi^e$ has residual
semisimplification $\bar\eta\bar\chi^{2e}\oplus1$.
For this canonical $W$, no extra quadratic twist
is required.

The whole local argument in
\path{AdmissibleOfTypeII.lean:142--195}
constructs a flag on the full decomposition group
with ordered inertial characters
$(\chi^e,\chi^{-e})$, since $\psi=1$.
Thus $R$ has ordered inertia $(\chi^{2e},1)$
and an unramified quotient.  This statement
also holds at an inert prime; the proof uses
the full rational decomposition group.

Finally \path{PtInertiaAway.lean:54--69}
shows that on $I_q\cap G_E$ the source point is
conjugate to the base inertia representation.
For $q\ne p$ and $\psi=1$ the latter is the
identity.  Hence $W$, and also $R$, is
unramified away from $p$ after restriction to
$E$.  GD gives $(p-1)\nmid2e$ and $p\nmid D$.
All hypotheses of
KT-ORDINARY-RANK2-UNIT-EXCLUSION now hold for
$a=2e$, contradicting (OR.6).
\end{proof}

The actual GD hypotheses are retained here.  In
particular with $N=2$ they require odd auxiliary
discriminant, as explicitly derived in
\path{AdmissibleOfTypeII.lean:119}.
The generic rank-two theorem itself has no such
odd-discriminant restriction.  The unit for $a=e$
likewise excludes any ordinary lower-weight
rank-two seed satisfying its stated residual
and unramified-after-$E$ conditions; this includes
the construction underlying the ordinary
symmetric-square examples when those hypotheses
hold.  Neither assertion eliminates an arbitrary
irreducible stable rank-three source.

% END evidence/unitary_audit/reducible_comparison/ordinary_rank2_exclusion.tex


% BEGIN evidence/unitary_audit/reducible_comparison/adjoint_unit_exclusion.tex
\section{The same unit excludes the full normalized adjoint branch}
\label{sec:kt-adjoint-unit-exclusion}

The ordinary rank-two exclusion also applies when
only its adjoint residual representation is known.
The necessary scalar character can be recovered
from that residual adjoint and then removed.  This
is a new exclusion for the specified adjoint
branch; the general stable rank-three branch
remains outside its scope.

\begin{theorem}[KT-ADJOINT-BRANCH-UNIT-EXCLUSION]
Let $p>3$, let $D<0$ be a fundamental discriminant
prime to $p$, put $E=\mathbf Q(\sqrt D)$ and
$\eta=\chi_D$, and let $e\ge2$ be even with
$(p-1)\nmid2e$.  Let $V$ be a continuous absolutely
irreducible two-dimensional representation of
$G_{\mathbf Q}$ over a finite extension of
$\mathbf Q_p$, with
\[
                  \det V=\eta\varepsilon_p^e.             \tag{AJ.1}
\]
Put $A=\operatorname{Ad}^0(V)=\operatorname{End}^0(V)$,
with the conjugation action.  Suppose that:
\begin{enumerate}
\item $A|_{G_E}$ is unramified at all finite places
      away from $p$;
\item the residual semisimplification of $A|_{G_E}$ is
      $1\oplus\bar\varepsilon_p^e
                  \oplus\bar\varepsilon_p^{-e}$;
\item at a place $v\mid p$ of $E$, $A|_{G_{E_v}}$
      has a stable full flag with ordered inertial
      characters
      $(\varepsilon_p^e,1,\varepsilon_p^{-e})$.
\end{enumerate}
Then $B_{1,\eta\omega^e}$ is not a $p$-adic unit.
The assertion holds at split and inert $p$ and
has no upper bound on $e$.
\end{theorem}

\begin{proof}
We may extend the finite coefficient field whenever
necessary to define a stable lattice or a flag.
Write $\psi=\bar\varepsilon_p^e$ and
$d=\operatorname{ord}(\psi|_{G_E})$.  Since
$p\nmid D$, the fields $E$ and $\mathbf Q(\mu_p)$
are disjoint, and $(p-1)\nmid2e$ gives $d>2$.

\paragraph{The residual two-dimensional representation splits.}
Choose a stable lattice in $V$ and reduce it.
The induced lattice in $\operatorname{End}^0(V)$
has the same residual semisimplification as any
other stable adjoint lattice.  Hypothesis~2
therefore makes the finite image of
$\operatorname{Ad}^0(\bar V)|_{G_E}$ triangularizable:
all its composition factors are one-dimensional.
Conjugation on trace-zero matrices is faithful
on $\mathrm{PGL}_2$ when $p\ne2$.
Thus the finite projective image $P_E$ has a
normal $p$-subgroup $U$, namely the kernel of
the diagonal-character map of that triangular
representation, and $P_E/U$ is abelian of
order prime to $p$.

If $U\ne1$, it has a unique fixed point on
$\mathbf P^1$.  Here is the elementary argument.
A central element of order $p$ in $U$ is
unipotent and has a unique fixed point; centrality
forces all of $U$ to preserve that point.
After putting it at infinity, every element of
the finite $p$-group has trivial diagonal ratio,
so $U$ consists of translations and its common
fixed point is unique.  Normality then forces
$P_E$ to preserve that point.

If $U=1$, $P_E$ is abelian and has an element
of projective order greater than two: otherwise
all its adjoint characters would have order at
most two, contrary to $d>2$.
Such an element is semisimple and its centralizer
in $\mathrm{PGL}_2$ is a torus.  Indeed a commuting
element must preserve its two eigenlines;
exchanging them would invert the projective
element and force order at most two.
Consequently $P_E$ lies in that torus.
In both cases the residual representation on
$G_E$ is reducible.

Its semisimplified character ratio, in a suitable
order, is $\psi$.  Restriction of a semisimple
finite-image representation to a normal subgroup
is semisimple by the socle argument in
KT-ORDINARY-RANK2-RESTRICTION-VARIANT.
Conjugation by $G_{\mathbf Q}$ cannot exchange
these two distinct character lines: their ratio
$\psi$ is conjugation-invariant and differs from
$\psi^{-1}$.  Hence
\[
                  \bar V^{\rm ss}=\lambda\oplus\mu
                                                               \tag{AJ.2}
\]
for global characters, with
$\lambda\mu^{-1}=\bar\eta^t\psi$ for some
$t\in\{0,1\}$.

At complex conjugation, $\det\bar V=-1$ by
(AJ.1), whereas $\psi(c)=1$ because $e$ is
even.  Since both $\lambda(c)$ and $\mu(c)$
are signs, their ratio equals their product.
Thus $t=1$.  Combining the ratio and determinant
in (AJ.1) gives $\mu^2=1$.
Let $\nu:G_{\mathbf Q}\to\{\pm1\}$ be the
canonical characteristic-zero lift of $\mu$.
We have proved
\[
          \bar V^{\rm ss}=
                  \bar\nu\oplus\bar\nu\bar\eta\psi.
                                                               \tag{AJ.3}
\]

\paragraph{The adjoint ordinary flag gives the required quotient.}
Because $p\nmid D$, $E_v/\mathbf Q_p$ is trivial
or unramified quadratic; its inertia is the same
as that of $\mathbf Q_p$.
The highest line $L$ of the flag in hypothesis~3
is the unique inertia eigenline of character
$\varepsilon_p^e$: that character occurs once
among the three distinct composition factors.
It is consequently stable under the full
$G_{\mathbf Q_p}$, which normalizes inertia and
preserves the cyclotomic character.

Take $0\ne N\in L\subset\operatorname{End}^0(V)$.
Conjugation preserves $\det N$, while inertia
scales $N$ by $\varepsilon_p^e$.  The latter
character has infinite order, so $\det N=0$.
The trace is zero as well; thus $N$ is a
rank-one nilpotent and
$\ker N=\operatorname{im}N$ is a line in $V$
stable under $G_{\mathbf Q_p}$.
Write $\delta_+$ for its character and
$\delta_-$ for that of the quotient.
The action on $N$ is $\delta_+/\delta_-$, so
on inertia
\[
 \delta_+/\delta_-=\varepsilon_p^e,
 \qquad \delta_+\delta_-=\varepsilon_p^e,
 \qquad \delta_-^2=1.                                    \tag{AJ.4}
\]

The saturated local line in a stable lattice
shows that $\bar\delta_-|_{I_p}$ is one of
$\bar\nu|_{I_p}$ and $\bar\nu\psi|_{I_p}$.
The second is impossible: its square is
$\psi^2\ne1$ on inertia, whereas (AJ.4) has
square one.  Therefore
$\bar\delta_-|_{I_p}=\bar\nu|_{I_p}$.
Both $\delta_-|_{I_p}$ and $\nu|_{I_p}$ take
values in $\{\pm1\}$, on which reduction is
injective.  They are equal.  It follows that
\[
                            V'=V\otimes\nu              \tag{AJ.5}
\]
has an unramified rank-one quotient at $p$.
This argument does not assume an unproved
converse from adjoint crystallinity: it uses
the exact ordered inertia flag in hypothesis~3.
For classical forms unramified at $p$, their
usual crystalline realization is consistent
with this local conclusion but is not needed
as an extra input.

\paragraph{The same scalar twist removes away-prime inertia.}
For $q\ne p$ and $g\in I_q\cap G_E$,
hypothesis~1 gives $A(g)=1$, so $V(g)=z_g I$.
Its determinant is one by (AJ.1), since both
$\eta$ and $\varepsilon_p$ are trivial on
this inertia subgroup.  Hence $z_g\in\{\pm1\}$.
Formula (AJ.3) shows that its reduction is
$\bar\nu(g)$, and injectivity on signs gives
$z_g=\nu(g)$.  Thus $V'(g)=1$.
The twisted representation is therefore
unramified away from $p$ after restriction
to $G_E$.

Finally (AJ.3) gives
$(\bar V')^{\rm ss}=1\oplus\bar\eta\psi$.
The twist preserves absolute irreducibility.
All three hypotheses of
KT-ORDINARY-RANK2-UNIT-EXCLUSION hold for
$V'$ and $a=e$.  A unit $B_{1,\eta\omega^e}$
would contradict that theorem.
\end{proof}

\paragraph{Application to the source and its boundary.}
For any actual Type~I point whose restricted
rank-three system is of the form
$\operatorname{Ad}^0(V)|_{G_E}$ with (AJ.1),
the source's residual matrix, away-prime
unramifiedness and ordinary flag give exactly
the three hypotheses above.  For
$E=\mathbf Q(\sqrt{-3})$, the lower-weight
classical adjoint subspace constructed in
KT-U3-SYM2-INJECT has determinant normalization
(AJ.1) before taking the adjoint.  Thus the
$a=e$ auxiliary unit excludes every actual
source point in this normalized adjoint
subspace, including residual pairs that
initially carry an unknown quadratic scalar
character.  This is stronger than excluding
only the previously normalized Eisenstein
seeds.

The determinant and adjoint-origin hypotheses
are explicit.  For other auxiliary fields the
same exclusion applies whenever this stated
adjoint origin and determinant have been
identified; no additional spectral construction
over those fields is asserted here.
No assertion identifies every
stable rank-three parameter with this family.
Under the two units with exponents $e$ and
$2e$, actual Type~II points and the specified
adjoint Type~I points are excluded; a
stable Type~I source outside this specified
normalized adjoint family is still a possible
remaining obstruction.

% END evidence/unitary_audit/reducible_comparison/adjoint_unit_exclusion.tex


% BEGIN evidence/lean_chain/determinant_normalization/scalar_normalization.tex
\section{Removing the determinant hypothesis from an odd adjoint origin}
\label{sec:kt-odd-adjoint-scalar-normalization}

The determinant normalization in
KT-ADJOINT-BRANCH-UNIT-EXCLUSION can be recovered by a global scalar
twist.  The input needed here is an odd rank-two origin over
$\mathbf Q$.  Descent of a self-dual automorphic rank-three system to
such an origin is a separate assertion.

\begin{theorem}[KT-ODD-ADJOINT-SCALAR-NORMALIZATION]
Let $p>3$, let $D<0$ be a fundamental discriminant with $p\nmid D$,
put $E=\mathbf Q(\sqrt D)$ and $\eta=\chi_D$, and let $e\ge2$ be
even with $(p-1)\nmid2e$.  Let $V$ be a continuous absolutely
irreducible two-dimensional representation of $G_{\mathbf Q}$ over
a finite extension of $\mathbf Q_p$.  Suppose $V$ is projectively
odd: complex conjugation acts with two distinct eigenvalues.
Put $A=\operatorname{Ad}^0(V)$.  Assume that
\begin{enumerate}
\item $A|_{G_E}$ is unramified at every finite place away from $p$;
\item $(\bar A|_{G_E})^{\mathrm{ss}}
       =1\oplus\omega^e\oplus\omega^{-e}$;
\item at a place $v\mid p$ of $E$, $A|_{G_{E_v}}$ has a full stable
      flag with ordered inertial characters
      $(\varepsilon_p^e,1,\varepsilon_p^{-e})$.
\end{enumerate}
Then, after a finite coefficient extension, there is a continuous
character $\theta:G_{\mathbf Q}\to\overline{\mathbf Q}_p^{\times}$,
with image in a finite coefficient field, such that
$V'=V\otimes\theta^{-1}$ satisfies
\[
 \det V'=\eta\varepsilon_p^e,\qquad
 (\bar V')^{\mathrm{ss}}=1\oplus\bar\eta\omega^e.       \tag{SN.1}
\]
Moreover, $V'|_{G_E}$ is unramified away from $p$, and $V'$ has a
$G_{\mathbf Q_p}$-stable line with unramified rank-one quotient.
Consequently $B_{1,\eta\omega^e}$ is not a $p$-adic unit.
There is no upper bound on $e$, and $p$ may split or be inert in $E$.
\end{theorem}

\begin{proof}
All lattices, eigenlines and characters below are defined after
finite coefficient extensions.  Their finiteness, including that
of the eventual global character, is justified in the construction.
Write $\psi=\omega^e$.
Since $p\nmid D$, $E\cap\mathbf Q(\mu_p)=\mathbf Q$ and the order
of $\psi|_{G_E}$ is greater than two.

\paragraph{The residual scalar character and its orientation.}
The residual-projective argument in the proof of
KT-ADJOINT-BRANCH-UNIT-EXCLUSION, up to (AJ.2), uses only
hypothesis~2 and this order condition.  It gives
\[
 \bar V^{\mathrm{ss}}=\lambda\oplus\mu,
 \qquad \lambda\mu^{-1}=\bar\eta^t\psi,
 \quad t\in\{0,1\}.                                    \tag{SN.2}
\]
For clarity, that argument first triangularizes the residual
adjoint on $G_E$.  Its faithful projective image is a finite group
with a normal unipotent $p$-subgroup and abelian prime-to-$p$
quotient.  A nontrivial such $p$-subgroup has a unique fixed point
on $\mathbf P^1$; if it is trivial, an element of order greater
than two forces the projective image into a torus.  Thus the two
residual characters exist on $G_E$.  Their ratio is conjugation
invariant and different from its inverse, so the global action
cannot exchange them.  This gives (SN.2).

Because complex conjugation is an involution and is not scalar
on $V$, its eigenvalues are $1,-1$.  Reduction preserves their
distinctness as $p>2$.  Since $e$ is even, $\psi(c)=1$ and
$\eta(c)=-1$, so (SN.2) has $t=1$.  Hence
\[
          \bar V^{\mathrm{ss}}=\mu\bar\eta\psi\oplus\mu.
                                                               \tag{SN.3}
\]
Here $\mu$ is an arbitrary finite residual character; it is not
assumed quadratic.

\paragraph{The quotient character at $p$.}
The highest line of the flag in hypothesis~3 is the unique inertia
eigenline of $A$ with character $\varepsilon_p^e$.
It is stable under $G_{\mathbf Q_p}$, since $E_v/\mathbf Q_p$ is
unramified or trivial and the cyclotomic character is invariant
under conjugation.  A nonzero element $N$ of that line is a
trace-zero matrix.  Conjugation preserves its determinant,
whereas inertia multiplies $N$ by the infinite-order character
$\varepsilon_p^e$.  Therefore $\det N=0$.
The line $\ker N=\operatorname{im}N$ in $V$ is
$G_{\mathbf Q_p}$-stable.  If $\delta_+$ and $\delta_-$ denote
its character and the quotient character, respectively, then
\[
            \delta_+/\delta_-\big|_{I_p}
                        =\varepsilon_p^e\big|_{I_p}.     \tag{SN.4}
\]
Intersecting this line with a stable lattice and taking the
saturated quotient identifies its reduced characters with the
two factors in (SN.3) on inertia.  Since $\eta$ is unramified
at $p$ and $\psi^2|_{I_p}\ne1$, the order in (SN.4) forces
\[
                   \bar\delta_-|_{I_p}=\mu|_{I_p}.        \tag{SN.5}
\]
Indeed the opposite order would give ratio $\psi^{-1}$.
We will prescribe $\theta|_{I_p}=\delta_-|_{I_p}$.

\paragraph{Compatible scalar choices at every other prime.}
Fix a rational prime $q\ne p$, and choose a decomposition group
$G_{\mathbf Q_q}$ and its inertia group $I_q$.
For $g\in I_q\cap G_E$, hypothesis~1 and the kernel of the
adjoint action give $V(g)=z_g I$.

If $q\nmid D$, this says that all of $V(I_q)$ is scalar.
A Frobenius lift has an eigenline after a finite coefficient
extension.  That line is stable under inertia and the Frobenius
lift, hence under the full decomposition group, whose quotient
by inertia is topologically procyclic.  Its character
$\alpha_q$ therefore satisfies
\[
               \alpha_q(g)=z_g\quad(g\in I_q),
               \qquad \bar\alpha_q|_{I_q}=\mu|_{I_q}.     \tag{SN.6}
\]

If $q\mid D$, the projective inertia image has order at most two.
It has order exactly two: (SN.3) has distinct inertial characters
$\mu\bar\eta$ and $\mu$, since $q\ne p$.
Choose $h\in I_q$ with $\eta(h)=-1$.  The matrix $V(h)$ has two
distinct eigenvalues with ratio $-1$.  They have distinct
reductions, so their spectral idempotents are integral on the
chosen stable lattice after a finite extension.  Every inertia
matrix is scalar or a scalar multiple of $V(h)$; thus there are
two inertia eigenlines, one with reduced character $\mu|_{I_q}$
and the other with $\mu\bar\eta|_{I_q}$.
The decomposition group preserves each line.  To see this,
conjugation by any of its elements sends the first character to
one whose reduction is still $\mu|_{I_q}$, because $\mu$ is a
global character.  The second line has a distinct reduction, so
the two lines cannot be exchanged.
Let $\alpha_q$ be the character of the first line.  Then
\[
 \bar\alpha_q|_{I_q}=\mu|_{I_q},\qquad
 \alpha_q(g)=z_g\quad(g\in I_q\cap G_E),                 \tag{SN.7}
\]
and the ratio of the other inertia character to $\alpha_q$
is exactly $\eta|_{I_q}$.  It is trivial on $I_q\cap G_E$
and equals $-1$ on the other coset.

\paragraph{Only finitely many of these inertia characters are nontrivial.}
We give this step explicitly; continuity alone is not being used
as a general finite-ramification assertion for arbitrary
$p$-adic representations.  For $q\ne p$, the restriction of a
continuous character of $G_{\mathbf Q_q}$ to inertia has finite
image in any finite extension of $\mathbf Q_p$.
Local class field theory identifies that image as a quotient of
$\mathbf Z_q^{\times}$.  An open pro-$q$ subgroup has trivial
image in a sufficiently small torsion-free pro-$p$ subgroup of
the target; the remaining quotient is finite.
In particular, when $q\nmid Dp$, the scalar image $V(I_q)$
consists of roots of unity in the original coefficient field $K$.

Choose $n$ such that $1+\varpi^n\mathcal O_K$ is torsion free.
The finite representation $V\bmod\varpi^n$ is unramified
outside a finite set of rational primes.  Outside that set and
the divisors of $Dp$, the scalar image $V(I_q)$ is both torsion
and contained in $1+\varpi^n\mathcal O_K$, so is trivial.
Together with (SN.6)--(SN.7), this proves that
$\alpha_q|_{I_q}=1$ for all but finitely many $q\ne p$.

\paragraph{A single global character with the prescribed inertias.}
For each relevant prime $q$, local class field theory for
$\mathbf Q_q$ says that the inertia restriction of a local
character factors continuously through the $q$-adic cyclotomic
character
\[
                 \varepsilon_q:I_q\longrightarrow
                                      \mathbf Z_q^{\times}.
\]
Indeed the maximal abelian extension is the compositum of the
unramified extension and $\mathbf Q_q(\mu_{q^{\infty}})$;
the latter is totally ramified.  Thus there is a continuous
$\beta_q:\mathbf Z_q^{\times}\to K'^{\times}$ such that
\[
 \beta_p\circ\varepsilon_p|_{I_p}=\delta_-|_{I_p},
 \qquad
 \beta_q\circ\varepsilon_q|_{I_q}=\alpha_q|_{I_q} (q\ne p).
                                                               \tag{SN.8}
\]
The field $K'$ can be chosen finite: only finitely many nontrivial
characters occur, and each was obtained from an eigenline over
a finite extension.  Set
\[
               \theta=\prod_q\bigl(\beta_q\circ
                         \varepsilon_q:G_{\mathbf Q}\to K'^{\times}\bigr),
                                                               \tag{SN.9}
\]
omitting the trivial factors.  This is a finite product of
continuous characters.  At $I_q$ every factor with a different
index is unramified, so (SN.9) has exactly the inertias prescribed
in (SN.8).  There is no separate choice at the real place.

Equations (SN.5)--(SN.7) show that $\bar\theta\mu^{-1}$ is
unramified at every finite prime of $\mathbf Q$.
Such a finite character is trivial: its fixed field would be
a finite extension of $\mathbf Q$ unramified at every finite
prime, which is trivial by Minkowski's discriminant bound.
Consequently $\bar\theta=\mu$, giving the residual statement
in (SN.1).

\paragraph{The determinant and the ordinary normalization.}
At $p$, twisting by (SN.9) makes the quotient character
unramified; (SN.4) makes the determinant on inertia equal to
$\varepsilon_p^e$.  At $q\nmid Dp$, all inertia matrices of
$V'$ are the identity.  At $q\mid D$, their two characters are
$1,\eta|_{I_q}$ by (SN.7).  It follows that
\[
                    \det V'\,(\eta\varepsilon_p^e)^{-1}
                                                               \tag{SN.10}
\]
is a continuous character unramified at all finite primes.
It too is trivial: every finite quotient of its compact image
would give a finite everywhere-unramified extension of
$\mathbf Q$.  This proves the determinant statement of (SN.1).
Equations (SN.6)--(SN.8) also show directly that
$V'|_{G_E}$ is unramified away from $p$.

Absolute irreducibility is unchanged by the scalar twist.
The residual representation, ordinary quotient and away-prime
unramifiedness now meet all hypotheses of
KT-ORDINARY-RANK2-UNIT-EXCLUSION with $a=e$.
That theorem gives the asserted nonunit conclusion.
\end{proof}

\paragraph{Source scope and imported class field theory.}
The local character factorization used in (SN.8) is the
$\mathbf Q_q$ specialization of J.~S.~Milne,
\emph{Class Field Theory}, version~4.03 (6 August 2020),
Chapter~I, \S1.10--1.11, printed pp.~23--24, especially
the displayed decomposition of $\mathbf Q_q^{\mathrm{ab}}$.
The local Kronecker--Weber statement is Theorem~4.8 and
Corollary~4.12, pp.~47--48.  The everywhere-unramified argument
is the one in Lemma~4.17 and its proof, pp.~49--50.
These are imported class field theory and discriminant inputs;
no new proof of them or Lean kernel check is claimed here.
The author PDF is preserved in the adjacent primary-source
directory with its retrieval hash.

An actual Type~I system with an identified odd $G_{\mathbf Q}$
adjoint origin therefore needs no additional hypothesis on the
origin's determinant, nebentypus or scalar twist.  This theorem
does not by itself construct such an origin from a
conjugate-self-dual rank-three representation over $E$.

% END evidence/lean_chain/determinant_normalization/scalar_normalization.tex


% BEGIN evidence/unitary_audit/reducible_comparison/selfdual_source_exclusion.tex
\section{Self-dual stable sources reduce to odd rank two}
\label{sec:kt-selfdual-source-exclusion}

The determinant normalization in the preceding adjoint theorem is
not automatic for an arbitrary proposed two-dimensional origin.
The following automorphic descent first produces an odd origin for
every self-dual stable source.  The separate scalar-normalization
lemma then gives the exact ordinary rank-two exclusion.  In this
section, self-dual means $R^\vee\simeq R$ on $G_E$, with no Tate
twist; it is stronger than conjugate self-duality.

\begin{theorem}[KT-SELFDUAL-SOURCE-ODD-DESCENT]
Let $p>3$, let $E=\mathbf Q(\sqrt D)$ be imaginary quadratic with
$p\nmid D$, put $\eta=\chi_D$, and let $e\ge2$ be even with
$(p-1)\nmid2e$.
Let $R$ be a continuous absolutely irreducible three-dimensional
representation of $G_E$ over a finite extension of $\mathbf Q_p$.
Suppose that its residual semisimplification is
\[
                  1\oplus\omega^e\oplus\omega^{-e}.       \tag{SD.1}
\]
Suppose also that $R$ is attached, in the weight-zero unitary
normalization used in KT-U3-GENERAL-COVERAGE, to a regular algebraic
conjugate-self-dual cuspidal representation $\Pi$ of
$\mathrm{GL}_3(\mathbf A_E)$ whose complex infinity parameter is
\[
                 (z/\bar z)^e\oplus1\oplus(z/\bar z)^{-e}.
                                                               \tag{SD.2}
\]
If $R^\vee\simeq R$, there is a holomorphic cuspidal eigenform $f$
over $\mathbf Q$ of weight $e+1$, and a finite-coefficient
two-dimensional Galois representation $V_f$, such that
\[
              R\simeq\operatorname{Ad}^0(V_f)|_{G_E}.     \tag{SD.3}
\]
The representation $V_f$ is absolutely irreducible and is
projectively odd: $V_f(c)$ is not scalar for a complex conjugation
$c$.  No determinant normalization or fixed classical level is
asserted for $V_f$ at this stage.
\end{theorem}

\begin{proof}
\paragraph{The determinant is exactly one.}
Self-duality gives $(\det R)^2=1$.  Thus every value of $\det R$
is a sign.  Formula (SD.1) says that its reduction is one on every
element of $G_E$.  Reduction is injective on $\{\pm1\}$ because
$p>2$, so
\[
                              \det R=1.                  \tag{SD.4}
\]
This uses both actual self-duality and the residual determinant;
it does not infer (SD.4) from a residual statement alone.

The unitary-normalized attachment identifies the full unramified
Satake polynomials with the Frobenius polynomials of $R$ at almost
all places of $E$.  Self-duality and (SD.4), followed by strong
multiplicity one for $\mathrm{GL}_3$ and $\mathrm{GL}_1$, therefore
give
\[
                  \Pi^\vee\simeq\Pi,
                  \qquad\omega_\Pi=1.                   \tag{SD.5}
\]
The source's initial algebraic attachment is to
$R\varepsilon_p^{-1}$; the common twist in the already checked
unitary normalization is essential here.  No additional Tate
factor is inserted in (SD.5).

\paragraph{Cuspidal quadratic descent and its central character.}
Conjugate self-duality and (SD.5) imply
$\Pi^c\simeq\Pi$.  Arthur--Clozel, Chapter~3,
Theorem~4.2(d), printed p.~202, supplies a \emph{cuspidal}
representation $\Sigma$ of $\mathrm{GL}_3(\mathbf A_{\mathbf Q})$
with base change $\Pi$.  Thus no noncuspidal descent is being
silently used.  Theorem~5.1, pp.~212--213, makes this base change
strong at every finite and infinite place.

Base change of central characters and (SD.5) give
$\omega_\Sigma\circ N_{E/\mathbf Q}=1$.
Global class field theory for the quadratic extension gives
\[
                              \omega_\Sigma\in\{1,\eta\}.
                                                               \tag{SD.6}
\]
Both $\Sigma^\vee$ and $\Sigma$ base change to $\Pi$.
The precise fiber theorem, Arthur--Clozel Chapter~3,
Theorem~3.1, p.~201, or the fiber clause in Theorem~4.2(d),
gives $\Sigma^\vee\simeq\Sigma\otimes\eta^j$ with
$j\in\{0,1\}$.  If $j=1$, the central characters would satisfy
\[
                 \omega_\Sigma^{-1}
                          =\omega_\Sigma\eta^3,
\]
contradicting $\omega_\Sigma^2=1$ and $\eta\ne1$.
Consequently $\Sigma$ is self-dual.  Set
\[
                          \Sigma_0=\Sigma\otimes\omega_\Sigma.
                                                               \tag{SD.7}
\]
It remains cuspidal and self-dual, has central character
$\omega_\Sigma^4=1$, and still base changes to $\Pi$.

\paragraph{The real parameter gives a holomorphic rank-two origin.}
Strong base change at infinity identifies the restriction of the
real Weil parameter of $\Sigma_0$ to $\mathbf C^\times$ with
(SD.2).  In particular $\Sigma_0$ is regular algebraic, and its
three distinct exponents are exactly $(e,0,-e)$.

Ramakrishnan's \emph{An exercise concerning the selfdual cusp forms
on $\mathrm{GL}(3)$}, Theorem~A, gives
\[
                            \Sigma_0\simeq\operatorname{Ad}(\pi)
                                                               \tag{SD.8}
\]
for a non-dihedral cuspidal representation $\pi$ of
$\mathrm{GL}_2(\mathbf A_{\mathbf Q})$; its central character
can be chosen of finite order.  The quadratic factor in that
theorem is the central character of $\Sigma_0$, so it is one.
Corollary~B gives regular algebraicity after the usual common
norm normalization.  The explicit archimedean compatibility
in Proposition~C(b) and formula (6b) also fixes the weight here:
the two real discrete-series characters have ratio
$(z/\bar z)^e$, so their discrete-series weight is $e+1$.
Since $e\ge2$, this is a genuine discrete series.  Thus $\pi$
is represented by a holomorphic cuspidal eigenform $f$ of that
weight (choosing its holomorphic vector), in the ordinary
unitary normalization of a classical form.

Let $V_f$ be its usual finite-coefficient Galois representation.
If $\xi$ is the finite nebentypus, then
\[
                 \det V_f=\xi\varepsilon_p^e,
                 \qquad \xi(-1)=(-1)^{e+1}=-1.           \tag{SD.9}
\]
Hence $\det V_f(c)=-1$, whereas $V_f(c)^2=1$; its two
eigenvalues are $1,-1$.  In particular it is projectively odd.

At almost every place of $E$, strong base change, (SD.8),
Gelbart--Jacquet's adjoint local compatibility, and the classical
Galois attachment give the same full characteristic polynomial
for $R$ and $\operatorname{Ad}^0(V_f)|_{G_E}$.  Chebotarev and
Brauer--Nesbitt identify their semisimplifications.  Since $R$
is irreducible of dimension three, the other representation
is itself irreducible and isomorphic to $R$, giving (SD.3).
It also follows that $V_f$ is absolutely irreducible: a
reducible two-dimensional representation has reducible adjoint.
Thus no separate all-primes big-image assertion is needed.
\end{proof}

\begin{corollary}[KT-SELFDUAL-SOURCE-UNIT-EXCLUSION]
Assume the base hypotheses on $p,E,e,R,\Pi$ and (SD.1)--(SD.2)
in KT-SELFDUAL-SOURCE-ODD-DESCENT, without assuming that $R$ is
self-dual.  Suppose in addition that $R$ is unramified at all finite places of $E$
away from $p$, and at a place above $p$ it has a stable full flag
with ordered inertial characters
$(\varepsilon_p^e,1,\varepsilon_p^{-e})$.
The unit condition
\[
                          B_{1,\eta\omega^e}\in\mathbf Z_p^\times
                                                               \tag{SD.10}
\]
implies $R^\vee\not\simeq R$.
\end{corollary}

\begin{proof}
If $R$ were self-dual, (SD.3) would give the projectively odd
origin $V_f$.  Its adjoint has precisely the residual,
away-prime and ordered ordinary properties in the hypotheses.
KT-ODD-ADJOINT-SCALAR-NORMALIZATION supplies a global scalar
twist $V'$ with the residual pair $1\oplus\eta\omega^e$,
an unramified ordinary quotient at $p$, and unramified
restriction to $G_E$ away from $p$.
KT-ORDINARY-RANK2-UNIT-EXCLUSION, with $a=e$, contradicts
(SD.10).
\end{proof}

\paragraph{Exact actual-source application and remaining gate.}
For the actual $N=2$, $\psi=1$ source and GD, under the carried
closed source convention KT-ONECHAR-CLOSED, the Type~I clause
in \path{FOneCharU.lean}, lines 47--54, supplies
\path{Adm3U_s}.  Its base predicate
\path{AfStmt.lean}, lines 29--32, supplies irreducibility,
the algebraic automorphic attachment and the local level.
For the original ordered flag, the exact definitions are
\path{AutomorphicForms/BlochKatoSelmerVanishes.lean},
\texttt{IsFlagOn} at lines 37--39, \texttt{KillsMu} at
lines 66--67, \texttt{LevelAt} at lines 102--107, and
\texttt{Level.above} at lines 111--114.
Every character satisfying \texttt{KillsMu 2} is trivial because
$\mu_2\subset\mathbf Q$.  Hence that level clause is an actual
decomposition-group stable upper flag with exactly the ordered
inertial characters $(\varepsilon_p^e,1,\varepsilon_p^{-e})$;
it is not inferred just from potential crystallinity.
\path{IsAutomorphicOverImagQuad.lean}, lines 35--36, is
explicitly conjugate-self-dual and cuspidal.
The checked common twist and \path{KeepLevel.lean},
lines 30--82 and 95--101, give the unitary-normalized attachment
and (SD.2).  The actual residual matrix from
\path{IsPtOf.lean}, lines 124--166, gives (SD.1), and
\path{PtInertiaAway.lean}, lines 54--69, gives the required
away-prime unramifiedness.  Finite coefficients follow from
\path{TanDet.lean}, lines 212--216, exactly as in
KT-AUXILIARY-UNIT-NO-TYPEII.

Accordingly, the exponent-$e$ auxiliary unit excludes
\emph{every self-dual actual Type~I source}, including any
adjoint origin over $E$ because an adjoint is self-dual.
The exponent-$2e$ unit already excludes every actual Type~II
source by KT-AUXILIARY-UNIT-NO-TYPEII.
Under these two units, any remaining actual source is a
non-self-dual stable rank-three source.  Conjugate self-duality
then says that its conjugate is its distinct dual.  This is a
strictly narrower source class; no assertion excludes that
remaining class or proves the original strict Selmer vanishing.

\paragraph{Primary source and review boundary.}
The new published input is D.~Ramakrishnan, \emph{Indian Journal
of Pure and Applied Mathematics} \textbf{45} (2014), 777--785,
\url{https://doi.org/10.1007/s13226-014-0088-1}.
The whole five-page author version at
\url{https://math.uchicago.edu/~fcale/papers/Ramakrishnan.pdf}
was read: Theorem~A and Corollary~B on author pp.~1--2,
Proposition~C(b) on p.~3, formula (6b) on p.~4, and the
complete Corollary~B proof on pp.~4--5.  The foundational
descent and transfer theorems it cites are imported through
Ramakrishnan's theorem, not freshly reconstructed here.
The author-version and published pagination are deliberately
distinguished.  Arthur--Clozel's precise quadratic descent,
fiber and strong-base-change statements are the same checked
primary inputs used earlier, reread here for the descent scope.
The scalar-normalization lemma and the descent and exclusion
proofs above have completed whole-proof and exact-source review.
No new Lean kernel build is claimed.

% END evidence/unitary_audit/reducible_comparison/selfdual_source_exclusion.tex


% BEGIN evidence/unitary_audit/reducible_comparison/essential_selfdual_reduction.tex
\subsection{The remaining source is not essentially self-dual}
\label{sec:kt-essential-selfdual-reduction}

\begin{lemma}[KT-ESSENTIALLY-SELFDUAL-REDUCTION]
Let $p>3$, let $E=\mathbf Q(\sqrt D)$ be imaginary quadratic with
$p\nmid D$, and let $e\ge2$ be even with $(p-1)\nmid2e$.
Let $R$ be a continuous three-dimensional representation of
$G_E$ over a finite extension of $\mathbf Q_p$, with residual
semisimplification $1\oplus\omega^e\oplus\omega^{-e}$.
Suppose it is unramified at every finite place away from $p$,
and at every place above $p$ it has an ordinary flag with
ordered inertial characters
$(\varepsilon_p^e,1,\varepsilon_p^{-e})$.
If $p\nmid h(E)$, then $\det R=1$.
Moreover, if there is a continuous character $\nu$ such that
\[
                          R^\vee\simeq R\otimes\nu,       \tag{ES.1}
\]
then $\nu=1$, and $R$ is self-dual.
\end{lemma}

\begin{proof}
The determinant is unramified at every finite place of $E$.
Away from $p$ this follows from the hypothesis on $R$; above
$p$ it follows from the exact product of its three inertial
characters, which is one.  Global class field theory therefore
makes $\det R$ a character of the finite ideal class group of
$E$.  Its reduction is one by the residual triple.
A finite subgroup of the principal units of a finite extension
of $\mathbf Q_p$ has $p$-power order: reduction is injective
on all roots of unity of order prime to $p$.
Thus the image of $\det R$ is a $p$-group quotient of
$\operatorname{Cl}(E)$.  The class-number hypothesis forces
\[
                              \det R=1.                  \tag{ES.2}
\]

Taking determinants in (ES.1) now gives $\nu^3=1$.
At a place $v\mid p$, semisimplifying the inertia flag in
(ES.1) gives equality of the character multisets
\[
 \{\varepsilon_p^{-e},1,\varepsilon_p^e\}
   =\{\nu\varepsilon_p^e,\nu,\nu\varepsilon_p^{-e}\}
       \quad\text{on }I_v.                               \tag{ES.3}
\]
In particular $\nu|_{I_v}$ is one of
$1,\varepsilon_p^e,\varepsilon_p^{-e}$.
The latter two characters have infinite order on inertia,
whereas $\nu^3=1$, so $\nu|_{I_v}=1$.

Reducing (ES.1) and semisimplifying globally on $G_E$ now
gives
\[
     \{\omega^{-e},1,\omega^e\}
          =\{\bar\nu\omega^e,\bar\nu,
                                      \bar\nu\omega^{-e}\}.
                                                               \tag{ES.4}
\]
Hence $\bar\nu$ is one of $1,\omega^e,\omega^{-e}$ as a
global character.  Since $p\nmid D$, cyclotomic inertia at
$v$ has full residual image, and $(p-1)\nmid2e$ in
particular gives $(p-1)\nmid e$.  The last two possibilities
are nontrivial on inertia, contrary to $\nu|_{I_v}=1$.
Thus $\bar\nu=1$.  Reduction is injective on cube roots
of unity because $p>3$, so $\nu=1$.
\end{proof}

\begin{corollary}[KT-NONESSENTIALLY-SELFDUAL-SOURCE-GATE]
In the actual $N=2$, $\psi=1$ source with GD and the carried
closed source convention, suppose
\[
 B_{1,\eta},\quad B_{1,\eta\omega^e},\quad
 B_{1,\eta\omega^{2e}}\ \in\mathbf Z_p^\times,
 \qquad \eta=\chi_D.                                    \tag{ES.5}
\]
Under KT-SELFDUAL-SOURCE-UNIT-EXCLUSION, every actual
\texttt{TwoPt} source is Type~I and its restricted rank-three
representation is not essentially self-dual.
Consequently, if the original strict Selmer length is nonzero,
there is an actual stable rank-three source of that more
restricted kind, with the same global residual triple.
\end{corollary}

\begin{proof}
The first unit gives $p\nmid h(E)$ by the imaginary quadratic
class-number formula, as checked in the $a=0$ case of
KT-AUXILIARY-UNIT-SELMER-BRIDGE.
The third unit excludes Type~II by
KT-AUXILIARY-UNIT-NO-TYPEII.  Every remaining actual point
is Type~I by the source's two-type classification, and its
residual and local conditions are those of the lemma.
If its representation were essentially self-dual, the lemma
would make it self-dual; the second unit excludes this by
KT-SELFDUAL-SOURCE-UNIT-EXCLUSION.
The final assertion follows from the actual-point existence
argument in K17-S17-JOINT, rechecked in KT-JOINT-RECHECK:
nonzero strict Selmer length forces an actual \texttt{TwoPt},
not merely a residual constituent of an ambient Hecke space.
\end{proof}

The three auxiliary units have therefore removed both the
reducible source branch and all essentially self-dual stable
sources.  The remaining class is genuinely smaller, but it
has not been shown empty.  In particular this corollary does
not replace the original Selmer problem by a proved uniform
vanishing statement.  Its use of actual source points is
essential; no converse from an ambient automorphic congruence
to a point of the source ring is assumed.

% END evidence/unitary_audit/reducible_comparison/essential_selfdual_reduction.tex


% BEGIN evidence/lean_chain/determinant_normalization/ordinary_imprimitive_exclusion.tex
\subsection{An ordinary source cannot be monomial}
\label{sec:kt-ordinary-imprimitive-exclusion}

This is a further source-scope reduction.  It gives no bound on
the strict Selmer group and is not needed for the determinant
normalization or self-dual exclusion above.

\begin{lemma}[KT-ORDINARY-IMPRIMITIVE-EXCLUSION]
Let $p>3$, let $E$ be a number field unramified at a place
$v\mid p$, and let $e>0$ with $(p-1)\nmid e$.
Let $R:G_E\to\mathrm{GL}_3(K)$ be a continuous absolutely
irreducible representation over a finite extension $K$ of
$\mathbf Q_p$.  Suppose
\[
              \bar R^{\mathrm{ss}}=1\oplus\omega^e
                                            \oplus\omega^{-e}
                                                               \tag{IM.1}
\]
and $R|_{G_{E_v}}$ has a stable full flag whose inertia
characters are $(\varepsilon_p^e,1,\varepsilon_p^{-e})$.
Then $R$ does not preserve a set of three lines with direct sum
equal to its underlying space, even after a finite coefficient
extension.  Its projective image is infinite.
\end{lemma}

\begin{proof}
Suppose first that three such lines $L_1,L_2,L_3$ exist.
Their permutation action is transitive, since a union of any
proper collection of orbits would span a proper nonzero
$G_E$-subrepresentation.  Let $P\subset S_3$ be the permutation
image and $H$ its kernel.

\paragraph{Reduction retains the permutation quotient.}
There is a stable monomial lattice.  Choose $0\ne v_1\in L_1$
and choose $g_i$ carrying $L_1$ to $L_i$; set $v_i=R(g_i)v_1$.
The stabilizer of $L_1$ acts on it by a continuous character
of a compact group, hence by units of the coefficient field.
For any $g\in G_E$, write $g g_i=g_j h$ with $h$ in that
stabilizer.  Then $R(g)v_i$ is a unit multiple of $v_j$.
Consequently $\bigoplus_i\mathcal O_K v_i$ is stable, after
one finite coefficient extension defining the lines.
Its reduction consists of monomial matrices with the same
permutation quotient $P$: the nonzero entries remain units.

By (IM.1), the finite residual image can also be triangularized.
The kernel of its diagonal-character map is a normal
$p$-subgroup, and the quotient is abelian.  Its image in $P$
is trivial because $|P|$ divides six and $p>3$.
Thus $P$ is abelian.  The only abelian transitive subgroup
of $S_3$ is $C_3$, so
\[
                          G_E/H\simeq C_3.                \tag{IM.2}
\]

\paragraph{The permutation action is trivial on $p$-adic inertia.}
Put $I=I_v$ and $I_0=I\cap H$.  This is a finite-index normal
subgroup of $I$.  The three lines are $I_0$-stable, and their
three characters are, in some order,
\[
                  \varepsilon_p^e|_{I_0},\quad1,
                              \varepsilon_p^{-e}|_{I_0}.  \tag{IM.3}
\]
Indeed these are the composition characters of the restricted
ordinary flag.  They are pairwise distinct: the cyclotomic
image on $I_0$ remains infinite, whereas $e>0$.
Each character in (IM.3) is invariant under conjugation by $I$.
Conjugation must therefore preserve each of the three distinct
character lines, so $I\subset H$.

\paragraph{The cubic self-twist contradicts the residual characters.}
Choose a nontrivial character
$\tau:G_E/H\to\mu_3$.  The monomial representation is induced
from the stabilizer $H$ of one line, and hence
\[
                              R\otimes\tau\simeq R.        \tag{IM.4}
\]
One can see this without an induction theorem: on the above
basis put $T v_i=\tau(g_i)v_i$.  The identity
$\tau(g)\tau(g_i)=\tau(g_j)$ gives
$T R(g)T^{-1}=\tau(g)R(g)$.
Since $I\subset H$, $\tau$ is unramified at $v$.
Reduction is injective on $\mu_3$ because $p>3$, so
$\bar\tau$ is still nontrivial.

Reducing (IM.4) and using (IM.1) gives equality of global
character multisets
\[
 \{1,\omega^e,\omega^{-e}\}
       =\{\bar\tau,\bar\tau\omega^e,
                                  \bar\tau\omega^{-e}\}. \tag{IM.5}
\]
In particular $\bar\tau=\omega^e$ or $\omega^{-e}$.
Both are nontrivial on $I_v$: $E_v/\mathbf Q_p$ is unramified
and $(p-1)\nmid e$.  This contradicts the unramifiedness of
$\tau$ at $v$.  This also covers the case in which
$\omega^e$ has order three.

Finally, if the projective image of $R$ were finite, a
finite-index subgroup $I'\subset I_v$ would act by scalar
matrices.  Its three composition characters would be equal.
The ordinary flag gives instead
$\varepsilon_p^e|_{I'},1,\varepsilon_p^{-e}|_{I'}$, which
are pairwise distinct as above.  This is impossible.
\end{proof}

For the actual source, the hypotheses on the field and exponent
follow from GD, and the ordered flag follows from
\texttt{LevelAt} at modulus two.  The lemma uses no auxiliary
unit condition and does not assume that a characteristic-zero
ordinary source can be recovered from an ambient residual
Hecke eigenvalue.  Combined with the three-unit exclusion,
it leaves primitive, nonessentially-self-dual actual stable
systems with infinite projective image.  No theorem comparing
their integral deformation rings with a Hecke algebra is
deduced here.

% END evidence/lean_chain/determinant_normalization/ordinary_imprimitive_exclusion.tex


% BEGIN evidence/unitary_audit/reducible_comparison/polarized_extension_graph.tex
\section{The auxiliary-unit extension graph does not close formally}
\label{sec:kt-polarized-extension-graph}

The three auxiliary-field units now have the precise
arithmetic consequence proved in
KT-AUXILIARY-UNIT-SELMER-BRIDGE: they kill the strict
groups with characters $\eta$, $\eta\chi^{-e}$ and
$\eta\chi^{-2e}$.  This section tests whether a
polarized generalized matrix algebra can turn that
information alone into exclusion of an irreducible
ordinary rank-three source.  The answer to that
\emph{formal implication} is negative.  The example
below is an abstract profinite group, not a number
field or an automorphic counterexample.

For the actual residual unitary action
$r_0=(\operatorname{diag}(\Psi,1,\Psi^{-1}),\eta)$,
$\Psi=\chi^e$, and
$\theta_{\rm Lie}(X)=-JX^tJ$, the root characters are
\[
\begin{array}{c|c}
\text{matrix line}&\text{global coefficient character}\\ \hline
E_{12}-E_{23}&\Psi\\
E_{12}+E_{23}&\eta\Psi\\
E_{21}-E_{32}&\Psi^{-1}\\
E_{21}+E_{32}&\eta\Psi^{-1}\\
E_{13}&\eta\Psi^2\\
E_{31}&\eta\Psi^{-2}.
\end{array}                                                \tag{EG.1}
\]
The negative rational line in this table is precisely
the original strict Selmer direction.  The positive
$\eta\Psi$ line is allowed by the upper ordinary flag;
it need not satisfy the strict local condition killed
by an auxiliary Bernoulli unit.  Even an additional
vanishing input for the rational positive-root group
does not remove these two remaining directions.

\begin{theorem}[KT-POLARIZED-EXTENSION-GRAPH-OBSTRUCTION]
Let $\ell>3$, let $e>0$ be even, and suppose that
the order $d$ of $\bar\chi^e$ in $\mathbf F_\ell^\times$
is greater than three.  There exist an abstract
profinite group $G$, characters $\chi_0,\eta$ on $G$,
and a continuous polarized representation
\[
 r:G\longrightarrow
          \mathrm{GL}_3(\mathbf Z_\ell)\rtimes\{\pm1\}
                                                               \tag{EG.2}
\]
with the following properties.  Put $\Psi=\chi_0^e$,
$G_E=\ker\eta$ and $k=\mathbf F_\ell$.

The reduction of $r$ is
$(\operatorname{diag}(\bar\Psi,1,\bar\Psi^{-1}),\eta)$,
and its characteristic-zero restriction to $G_E$ is
absolutely irreducible.  Among all one-dimensional
$k$-characters factoring through the finite quotient
constructed below, the only nonzero $H^1(G,k(\alpha))$
are
\[
 H^1(G,k(\eta\bar\Psi))\simeq k,
 \qquad H^1(G,k(\bar\Psi^{-1}))\simeq k.                  \tag{EG.3}
\]
In particular all of
\[
 H^1(G,k(\eta)),\quad
 H^1(G,k(\eta\bar\Psi^{-1})),\quad
 H^1(G,k(\eta\bar\Psi^{-2})),\quad
 H^1(G,k(\bar\Psi)),\quad
 H^1(G,k(\bar\Psi^2))                                    \tag{EG.4}
\]
vanish.  These five vanishings also hold with the corresponding
Teichm\"uller lifts of the characters and coefficients
$\mathbf Q_\ell/\mathbf Z_\ell$.
There is also a designated pair of
conjugate subgroups on which $r$ is upper triangular
in the prescribed order and the rational negative
class in (EG.3) restricts to zero.

Thus vanishing of these entire residual cohomology
groups, which is stronger than the corresponding
strict-group vanishing, is consistent with an
irreducible polarized representation and the indicated
ordinary flag.  No conclusion about actual Galois
groups or crystalline representations is asserted.
\end{theorem}

\begin{proof}
Choose a Teichm\"uller generator
$\zeta\in\mathbf Z_\ell^\times$ of order $\ell-1$
and put $q=\zeta^e$, of order $d>3$.
Let $F$ be the free pro-$\ell$ group on generators
$x,y$, and let
$H=\langle h\rangle\times\langle c\rangle
\simeq C_{\ell-1}\times C_2$ act by
\[
 hxh^{-1}=x^q,\quad hyh^{-1}=y^{q^{-1}},\qquad
 cxc=x^{-1},\quad cyc=y.                                 \tag{EG.5}
\]
Powers by $\mathbf Z_\ell$-units are defined in the
procyclic subgroups generated by $x,y$.
The indicated substitutions are automorphisms: their
inverse substitutions use $q^{-1},q$.  They commute,
and the substitution for $h$ has order dividing
$\ell-1$.  Hence $G=F\rtimes H$ is well defined.
Set $\chi_0(h)=\zeta$, $\chi_0(c)=-1$,
$\chi_0(F)=1$, and $\eta(c)=-1$, $\eta(h)=\eta(F)=1$.
Evenness of $e$ gives $\Psi(c)=1$.

Set
\[
 X=E_{12}+E_{23},\quad Y=E_{21}-E_{32},\quad
 T=\operatorname{diag}(q,1,q^{-1}).                       \tag{EG.6}
\]
Define
\[
 r(h)=(T,1),\quad r(c)=(1,-1),\quad
 r(x)=(\exp(\ell X),1),\quad
 r(y)=(\exp(\ell Y),1).                                  \tag{EG.7}
\]
The exponentials terminate in degree two and are
integral because two is a unit.  Their images lie in
the pro-$\ell$ group $1+\ell M_3(\mathbf Z_\ell)$,
so freeness gives the required continuous map on $F$.
Conjugation by $T$ scales $X,Y$ by $q,q^{-1}$.
Moreover $\theta_{\rm Lie}(X)=-X$ and
$\theta_{\rm Lie}(Y)=Y$.  Thus (EG.7) respects every
relation (EG.5) and defines (EG.2), including its
polarization and residual normalization.

For absolute irreducibility it suffices to work over
$\overline{\mathbf Q}_\ell$.  Since the eigenvalues
$q,1,q^{-1}$ of $T$ are distinct, a subspace invariant
under $T$ is spanned by a subset of the coordinate
vectors.  Invariance under $\exp(\ell X)$ and
$\exp(\ell Y)$ then forces every nonempty such subset
to contain all three vectors.  Equivalently, these
two logarithms generate $\mathfrak{sl}_3$:
\[
 [X,Y]=\operatorname{diag}(1,-2,1)=Z,\quad
 [Z,X]=3(E_{12}-E_{23}),\quad
 [Z,Y]=-3(E_{21}+E_{32}).                                 \tag{EG.8}
\]
These identities recover the four simple matrix
units, their brackets the two long units and the
two diagonal directions.

For (EG.3), $F$ acts trivially on each residual
coefficient module, and $\ell\nmid|H|$.
Inflation--restriction and averaging therefore give
\[
 H^1(G,k(\alpha))
 \simeq\operatorname{Hom}_H(F^{\rm ab}/\ell,
                                      k(\alpha)).          \tag{EG.9}
\]
The two generators of $F^{\rm ab}/\ell$ have exactly
the characters $\eta\bar\Psi$ and
$\bar\Psi^{-1}$ under $H$ by (EG.5).
This proves (EG.3).  None of the characters in
(EG.4) equals either of them when $d>3$.
For the divisible-coefficient assertion, each named
character is nontrivial modulo $\ell$, so its global
$H^0$ is zero.  The multiplication-by-$\ell$ exact
sequence identifies the $\ell$-torsion in its
divisible-coefficient $H^1$ with the corresponding
zero group in (EG.4).  Continuous cocycles from a
compact group into a discrete torsion module have
finite image, so this $H^1$ is $\ell$-primary torsion
and must itself be zero.

Finally let $D$ be the closed subgroup generated by
$h,x$.  Then $cDc=D$, so these give a designated
conjugate pair.  Its image under $r$ is upper
triangular with diagonal order $(\Psi,1,\Psi^{-1})$.
The cocycle corresponding to the second generator
of $F^{\rm ab}/\ell$ is zero on $D$, proving the
claimed strictness.  This is a group-theoretic local
analogue only; the theorem does not replace $D$ by
an actual local Galois group.
\end{proof}

The obstruction is the mixed pair of arrows
$E_{12}+E_{23}$ and $E_{21}-E_{32}$.  It connects all
three residual constituents.  Their commutators create
the missing odd negative and long-root directions as
higher products, as (EG.8) shows.  Vanishing of a
primitive extension quotient $A_{ij}/A'_{ij}$ can
therefore leave $A_{ij}=A'_{ij}$ nonzero.  Applying
Nakayama to the whole off-diagonal module without
first proving that the path products lie in its
coefficient maximal ideal would lose this distinction.

For the arithmetic problem, an additional argument
must control those mixed products, or relate an
integral extension module to an independently bounded
arithmetic quantity.  Requiring the remaining rational
negative strict group itself to vanish would be the
original Selmer hypothesis.  The three field units,
even together with rational positive-root vanishing,
do not establish that assertion by the formal GMA
argument tested here.

% END evidence/unitary_audit/reducible_comparison/polarized_extension_graph.tex


% BEGIN evidence/unitary_audit/reducible_comparison/local_mixed_paths.tex
\section{The mixed paths survive the actual local ordinary condition}
\label{sec:kt-local-mixed-paths}

The abstract group of KT-POLARIZED-EXTENSION-GRAPH-OBSTRUCTION
does not itself have local Galois subgroups.  The following
independent local calculation checks the two surviving
directions on $G_{\mathbf Q_\ell}$ itself.  It proves no
global existence statement and no smaller arithmetic gate.

\begin{lemma}[KT-LOCAL-MIXED-ROOT-PERSISTENCE]
Let $\ell>3$, let $e\ge2$ be even, and assume
$(\ell-1)\nmid2e$.  Let $\eta$ be the trivial or
unramified quadratic character of $G_{\mathbf Q_\ell}$.
Let $O$ be a finite coefficient valuation ring,
$K=\operatorname{Frac}O$, $k=O/(\varpi)$, and
$\chi=\varepsilon_\ell$.  Then
\[
 H^1_f(\mathbf Q_\ell,K(\eta\chi^e))
     =H^1(\mathbf Q_\ell,K(\eta\chi^e)),\qquad
 \dim_K H^1(\mathbf Q_\ell,K(\eta\chi^e))=1.             \tag{LM.1}
\]
Every class in $H^1(\mathbf Q_\ell,k(\bar\eta\bar\chi^e))$
lifts to an integral cocycle with coefficients
$O(\eta\chi^e)$.  After a finite coefficient extension
if necessary, there is an explicit two-parameter
local $GG_3$ family on the actual ordered ordinary
component whose tangent directions are the positive
twisted root and any prescribed local coboundary
in the negative rational root.  Its characteristic-zero
specializations are crystalline on $\ker\eta$ and have
ordered inertial characters $(\chi^e,1,\chi^{-e})$.
\end{lemma}

\begin{proof}
Put $V=K(\eta\chi^e)$.  It is crystalline, since
$\eta$ is unramified.  Its $\operatorname{Fil}^0$
de Rham subspace is zero, and $V^{G_{\mathbf Q_\ell}}=0$.
The Bloch--Kato dimension formula gives
$\dim_KH^1_f(\mathbf Q_\ell,V)=1$.
Local duality gives $H^2(\mathbf Q_\ell,V)=0$, since
$V^*(1)=K(\eta^{-1}\chi^{1-e})$ has no invariants.
The local Euler characteristic formula therefore
gives $\dim_KH^1(\mathbf Q_\ell,V)=1$, proving (LM.1).

For integral lifting, finite local duality gives
\[
 H^2(\mathbf Q_\ell,k(\bar\eta\bar\chi^e))=0.            \tag{LM.2}
\]
Indeed the dual coefficient has inertial character
$\bar\chi^{1-e}$; the exponent $1-e$ is odd and
cannot be zero modulo the even integer $\ell-1$.
Use the short exact sequences for successive
$O/(\varpi^n)$ coefficient modules.  Formula (LM.2)
makes the maps on $H^1$ surjective.  Choose compatible
classes and then compatible cocycle representatives;
the latter adjustment is by a lifted zero-cochain.
Their inverse limit is a continuous integral cocycle
$u:G_{\mathbf Q_\ell}\to O(\eta\chi^e)$ lifting the
given residual class.

Let $s$ be a formal parameter and define the exact
two-dimensional representation
\[
 \rho_s(g)=
 \begin{pmatrix}
  \eta(g)\chi(g)^e&s\,u(g)\\0&1
 \end{pmatrix}.                                           \tag{LM.3}
\]
The cocycle identity proves multiplicativity over
$O[[s]]$.  For every specialization $s\mapsto s_0$
in the maximal ideal of a finite coefficient ring,
its extension class lies in (LM.1); hence that
specialized representation is crystalline.

Enlarge coefficients to contain $\lambda$ with
$\lambda^2=-1/2$ and use the ordered basis
$(E,\lambda H,F)$ of $\mathfrak{sl}_2$ from (HS.3).
Write $A_s=\operatorname{Ad}^0\rho_s$ in this basis,
$D=\operatorname{diag}(-1,1,-1)$, and
$\eta(g)=(-1)^{j(g)}$.  The trace form and the
normalization identities in (HS.13)--(HS.14) give
\[
 r_s(g)=(A_s(g)D^{j(g)},\eta(g)).                          \tag{LM.4}
\]
This is an exact continuous homomorphism into $GG_3$.
It is upper triangular with diagonal
$(\chi^e,1,\chi^{-e})$ at every $g$, including the
nonsplit unramified quadratic case.  On $\ker\eta$
it is the adjoint of (LM.3), so all its
characteristic-zero specializations are crystalline.
The first-order upper-root matrix is a coefficient
unit times
\[
 X_+=E_{12}+E_{23},\qquad
 \operatorname{ad}(E)=-2\lambda X_+.                      \tag{LM.5}
\]
The semidirect adjoint action on this line is
$\eta\chi^e$, as required; its exact scalar
normalization does not change which residual
cohomology classes can occur.

Now suppose that a global negative rational cocycle
is strict at this place.  Its finite-coefficient
local class is zero by the local $H^0$ argument in
KT-ODD-CONORMAL-SELMER-EQUIVALENCE, so its local
cocycle is a coboundary $\operatorname{cob}(w)$
on the line $Y_-=E_{21}-E_{32}$.  Lift the scalar
$w\in k$ to $\widetilde w\in O$; the same argument
works for a prescribed finite $O/(b)$ chunk.
With a second formal parameter $t$, put
\[
 C_t=\exp(t\widetilde wY_-),\qquad
 r_{s,t}(g)=(C_t,1)r_s(g)(C_t^{-1},1).                    \tag{LM.6}
\]
The exponential and its inverse are degree-two
polynomials, because $Y_-^3=0$.  Formula (LM.6)
uses conjugation in $GG_3$, including its outer
action when $\eta(g)=-1$.  Choose the sign of
$\widetilde w$ according to the chosen coboundary
convention.  Its $t$-tangent is then precisely the
prescribed local negative rational coboundary.

Since $C_t\equiv1$ modulo the coefficient maximal
ideal $(\varpi,s,t)$, this is a framed change of
basis near the identity.  It carries the standard
ordered flag of (LM.4) to the ordered ordinary flag
of (LM.6), with exactly the same inertial diagonal.
Crystallinity is unchanged by conjugation.  Thus all
mixed $s,t$ terms in this explicit family are
compatible with the actual local condition, not
only with its linearization.  The integral
identification $U.R/I_t$ with the ordered ordinary
functor in KT-ORDINARY-INTEGRAL-COMPONENT places
(LM.6) on the source's chosen local component.
At a split auxiliary prime the same construction
applies at both conjugate places; at an inert prime
the semidirect formula already includes the
unramified quadratic descent.
\end{proof}

\paragraph{Exact imported local inputs.}
Laurent Berger, \emph{On $p$-adic Galois representations},
author PDF \texttt{barcelone.pdf}, \S5, Theorem~5.1,
author p.~4, states the local duality and Euler
characteristic formula for characteristic-zero
coefficients.  Section~16, author pp.~14--15,
identifies crystalline extension classes and derives
the displayed Bloch--Kato dimension formula from
the fundamental exact sequence.  The following
table on p.~15 explicitly treats positive Tate
twists; (LM.1) uses the same formula for the
unramified quadratic twist.  Finite local Tate
duality is the same imported theorem used in
Patrikis, \emph{Deformations of Galois
representations and exceptional monodromy},
Proposition~4.4, author p.~13, in the displayed
$H^2$--$H^0$ calculation.  The previously checked
integral ordinary-component theorem is used with
its arbitrary finite coefficient-DVR scope.

The construction is entirely local.  A global
negative rational class need not lift to a
characteristic-zero global representation merely
because its restrictions fit (LM.6).  Global cup
products, higher relations and integral saturation
remain separate questions.  This calculation
removes no global source and proves no additional
vanishing of the original strict Selmer group.

% END evidence/unitary_audit/reducible_comparison/local_mixed_paths.tex


% BEGIN evidence/unitary_audit/reducible_comparison/global_mixed_cup.tex
\section{The first mixed obstruction vanishes arithmetically}
\label{sec:kt-global-mixed-cup}

The source reductions above do not make the remaining rational
negative root disappear.  Here the first mixed obstruction is
computed over the actual rational Galois group, with its ordinary
local correction.  The class-number unit removes this obstruction;
it does not force the negative class to be zero.  The construction
stops at a characteristic-$p$ Artin algebra of nilpotence degree
three and supplies no characteristic-zero source point.

\begin{lemma}[KT-AUXILIARY-CLASS-UNIT-MIXED-CUP]
Let $p>3$, let $D<0$ be a fundamental discriminant prime to $p$,
put $E=\mathbf Q(\sqrt D)$ and $\eta=\chi_D$, and set
\[
                  S=\{\infty,p\}\cup\{q:q\mid D\}.
                                                               \tag{MC.1}
\]
Let $k$ be a finite extension of $\mathbf F_p$.
If $B_{1,\eta}$ is a $p$-adic unit, then
\[
 H^2(G_{\mathbf Q,S},k(\eta))=0,\qquad
 \dim_k H^1(G_{\mathbf Q,S},k(\eta))=1,                  \tag{MC.2}
\]
and localization induces an isomorphism
\[
 H^1(G_{\mathbf Q,S},k(\eta))
       \ \xrightarrow{\ \sim\ }
 H^1(\mathbf Q_p,k(\eta))/H^1_{\rm unr}(\mathbf Q_p,k(\eta)).
                                                               \tag{MC.3}
\]
The finite set in (MC.1) is part of the assertion; no claim is
made after allowing arbitrary extra ramified primes.
\end{lemma}

\begin{proof}
It suffices first to take $k=\mathbf F_p$, and then extend
scalars: all cochain spaces and local maps are finite direct
sums after a finite residue-field extension.
Write $L=E(\mu_p)$ and $\Delta=\operatorname{Gal}(L/\mathbf Q)$.
The $a=0$ case of KT-AUXILIARY-UNIT-SELMER-BRIDGE gives
\[
              p\nmid h(E),\qquad
              \operatorname{Cl}(L)\{p\}^{\eta\omega}=0.
                                                               \tag{MC.4}
\]
For $M=\mathbf F_p(\eta)$ the Cartier dual is
$M^D=\mathbf F_p(\eta\omega)$.
A class in $\Sha^1_S(\mathbf Q,M^D)$ is unramified outside
$S$ and zero on each decomposition group in $S$.
Its restriction to $G_L$ is therefore an everywhere-unramified
homomorphism.  Since $|\Delta|=2(p-1)$ is prime to $p$,
restriction is injective and class field theory embeds this
group into
\[
   \operatorname{Hom}_\Delta
       (\operatorname{Cl}(L)/p,\mathbf F_p(\eta\omega)).  \tag{MC.5}
\]
The character in (MC.5) is the displayed character, not its
inverse: equivariance is $f(\sigma x)=\eta(\sigma)
\omega(\sigma)f(x)$.  Thus (MC.4) makes (MC.5) zero.

Poitou--Tate duality pairs this zero group perfectly with
$\Sha^2_S(\mathbf Q,M)$.  At $p$, local duality gives
$H^2(\mathbf Q_p,M)=H^0(\mathbf Q_p,M^D)^\vee=0$,
because $\omega$ is nontrivial on inertia and $\eta$ is
unramified.  At $q\mid D$, the character $\eta\omega$
is nontrivial on inertia because $q\ne p$.  The local
$H^2$ is again zero.  At the real place it is zero by
averaging, since $p$ is odd.  Hence the global $H^2$ equals
$\Sha^2_S$ and is zero.  Global Euler characteristic for
the odd character $\eta$ now gives
$h^0-h^1+h^2=-1$.  Since $h^0=h^2=0$, this proves (MC.2).

Local Euler characteristic at $p$ gives $h^1=h^0+1$.
If $\eta(p)=-1$, the two dimensions $h^1$ and
$h^1_{\rm unr}$ are $1,0$; if $\eta(p)=1$, they are $2,1$.
Thus the target of (MC.3) has dimension one in either case.
Its kernel is zero.  Indeed a class in that kernel restricts
to an everywhere-unramified homomorphism on $G_E$: it is
unramified at $p$ by the kernel condition; outside $S$ this
is built into $G_{\mathbf Q,S}$; and at every $q\mid D$,
local duality and local Euler characteristic give
\[
               H^1(\mathbf Q_q,\mathbf F_p(\eta))=0
                                                               \tag{MC.6}
\]
because both local $H^0$ and $H^2$ vanish.
The homomorphism factors through $\operatorname{Cl}(E)/p=0$
by (MC.4).  Restriction to $G_E$ is injective because its
index is two.  This proves injectivity, and equality of
dimensions proves (MC.3).
\end{proof}

\begin{theorem}[KT-GLOBAL-QUADRATIC-MIXED-PERSISTENCE]
Retain the preceding hypotheses, let $e\ge2$ be even with
$(p-1)\nmid2e$, and put
\[
        \psi=\omega^e,\qquad
        \alpha=\eta\psi,\qquad \beta=\psi^{-1}.
                                                               \tag{MC.7}
\]
Let $u$ be a continuous cocycle in
$Z^1(G_{\mathbf Q,S},k(\alpha))$.
Let $v$ be a continuous cocycle in $Z^1(G_{\mathbf Q,S},k(\beta))$
whose cohomology class is zero at $\mathbf Q_p$ and unramified
on inertia at every finite prime other than $p$.
There is a continuous polarized residual deformation over
\[
                           B=k[s,t]/(s^2,t^2)             \tag{MC.8}
\]
of
\[
 r_0(g)=(\operatorname{diag}(\psi(g),1,\psi(g)^{-1}),\eta(g))
\]
in $\mathrm{GL}_3(B)\rtimes\{\pm1\}$, with first-order
terms
\[
       s\,u(g)(E_{12}+E_{23})+
                         t\,v(g)(E_{21}-E_{32}).          \tag{MC.9}
\]
Its restriction to $G_E$ is unramified away from $p$.
At $p$, it is conjugate by a matrix congruent to one to an
upper-triangular family with the prescribed specialized
inertial diagonal $(\psi,1,\psi^{-1})$.
Thus the two directions are compatible through their mixed
quadratic term, with the actual ordinary local condition.
\end{theorem}

\begin{proof}
Set
\[
 X=E_{12}+E_{23},\qquad Y=E_{21}-E_{32},\qquad
 Z=[X,Y]=\operatorname{diag}(1,-2,1).                     \tag{MC.10}
\]
The unitary outer automorphism is
$\theta(g)=J(g^t)^{-1}J$, with $J$ the antidiagonal identity.
The adjoint action of $r_0$ on $X,Y,Z$ has characters
$\alpha,\beta,\eta$, respectively.
As $(s,t)^3=0$, logarithms and exponentials terminate after
degree two; division by two is permitted.
For an unknown continuous cochain $w:G_{\mathbf Q,S}\to k$,
put
\[
 L(g)=s u(g)X+t v(g)Y+st w(g)Z,
 \qquad r(g)=\exp(L(g))r_0(g).                            \tag{MC.11}
\]
The Baker--Campbell--Hausdorff identity here is a finite
polynomial identity, not a convergence assertion.
It says that $r$ is a homomorphism exactly when
\[
 \begin{split}
 w(gh)-w(g)-\eta(g)w(h)
   =\tfrac12\bigl(&u(g)\beta(g)v(h)\\
                 &-v(g)\alpha(g)u(h)\bigr).
 \end{split}                                             \tag{MC.12}
\]
The right side is one half of the difference of the two
cup-product cochains.  It is a continuous $2$-cocycle in
$k(\eta)$ because $\alpha\beta=\eta$.
By (MC.2) a continuous $w$ solving (MC.12) exists.
This constructs the global semidirect-product homomorphism.
All three matrices in (MC.10) have trace zero, so its left
matrix has determinant one.  The right component is exactly
$\eta$, preserving the stated polarization.

\paragraph{The ordinary correction at $p$.}
Since $[v]|_{G_{\mathbf Q_p}}=0$, choose $d\in k$ with
$v(g)=\beta(g)d-d$ on this decomposition group.
Conjugate (MC.11) by $\exp(tdY)$.
The negative first-order term becomes zero, and the logarithm
of the remaining deformation factor is
\[
       s u(g)X+st w_p(g)Z,\qquad
       w_p(g)=w(g)-\tfrac12d(1+\beta(g))u(g).             \tag{MC.13}
\]
There are no other terms because $s^2=t^2=0$.
The homomorphism property shows that $w_p$ is now a
$k(\eta)$-valued cocycle on $G_{\mathbf Q_p}$.
By (MC.3), a global cocycle $h$ in $k(\eta)$ can be chosen
so that $[w_p+h]$ is unramified at $p$.
Replacing $w$ by $w+h$ preserves (MC.12).
Because $\eta$ is trivial on inertia at $p$, coboundaries
vanish on that inertia.  Unramifiedness of the class therefore
means that the actual cochain $w_p+h$ is zero on inertia.
Formula (MC.13) is upper triangular and its diagonal
correction on inertia is now zero.  This is precisely the
claimed ordered ordinary family.  Restriction to $G_E$
and conjugation transport the flag to either place above $p$.
The argument also works when $p$ is inert.

\paragraph{Away-prime inertia.}
There is no ramification outside $S$ by construction.
At $q\mid D$, both $\alpha$ and $\eta$ are nontrivial on
inertia, while the cyclotomic character is unramified.
Local duality and local Euler characteristic give
\[
 H^1(\mathbf Q_q,k(\alpha))=0,
 \qquad H^1(\mathbf Q_q,k(\eta))=0.                     \tag{MC.14}
\]
First conjugate by a matrix $\exp(saX)$ to remove the local
$X$-cocycle.  The remaining logarithm has the form
$t v(g)Y+st w_q(g)Z$, with $w_q$ an $\eta$-cocycle.
The second vanishing in (MC.14) lets a further conjugation
by $\exp(stbZ)$ remove that cocycle.
The coefficient of $Y$ is unramified by hypothesis on $v$.
On $I_q\cap G_E$ the resulting representation is therefore
the identity, as is $r_0$.  Conjugating back preserves this
identity.  This proves the asserted away-prime property.
\end{proof}

\paragraph{What the calculation decides.}
The odd character $\alpha=\eta\omega^e$ has
$h^1-h^2=1$ by global Euler characteristic, so a nonzero
positive class $[u]$ always exists at the set (MC.1).
A nonzero strict negative class $[v]$, in contrast, is an
instance of the original rational strict Selmer problem.
The theorem says that, \emph{if that class exists}, the
class-number unit supplies no contradiction from the first
mixed cup product or its ordinary local correction.
It does not assume that class zero, and it does not prove
that it exists.

The family (MC.8) has characteristic $p$ and nilpotents.
No lift through all higher orders, no characteristic-zero
irreducibility, no automorphy, and no point of the normalized
global source ring has been constructed.  At the local
ordinary functor it is a genuine Borel family followed by
near-identity conjugation, as in the already checked
KT-ORDINARY-INTEGRAL-COMPONENT.  Its noncritical residual-character
condition holds here: every nonzero difference among $e,0,-e$
is even, so none is congruent to $1$ modulo the even integer
$p-1$.  No characteristic-zero crystalline-point condition
on the Artin ring $B$ is being imposed.
This verifies the local condition but does not resolve the
higher global equations.
Accordingly this section is an exact obstruction audit,
not an additional arithmetic vanishing gate.

\paragraph{Imported duality and exact source scope.}
We use J.~S.~Milne, \emph{Arithmetic Duality Theorems},
second edition (2006), author PDF
\url{https://www.jmilne.org/math/Books/ADTnot.pdf}.
The definitions of $G_{\mathbf Q,S}$ and $\Sha^r_S$ are in
Chapter~I, pp.~48 and 56; the finite-module Poitou--Tate
pairing is Theorem~4.10(a), pp.~56--57.
Local duality is Corollary~2.3, p.~28, and local Euler
characteristic is Theorem~2.8, p.~31.
Global Euler characteristic is Theorem~5.1, p.~67, with
the finite-$S$ setup on p.~66 and its useful cohomological
form restated in Remark~5.2(c), p.~68.
These complete statements and surrounding definitions were
read for the present application.  Their foundational proofs
are imported as theorems; no new proof of Poitou--Tate duality
or arithmetic comparison with etale cohomology is claimed.

% END evidence/unitary_audit/reducible_comparison/global_mixed_cup.tex


% BEGIN evidence/lean_chain/cubic_cohomology.tex
\subsection{The two character spaces at the next mixed degree}
\label{sec:kt-cubic-cohomology}

Put $\psi=\omega^e$, where $p>3$, $e\ge2$ is even, and
$(p-1)\nmid2e$.  Let $D<0$ be fundamental and prime to $p$,
$\eta=\chi_D$, and
$S=\{\infty,p\}\cup\{q:q\mid D\}$.  All cohomology below uses
this exact finite set.  The coefficient field $k$ is a finite extension
of $\mathbf F_p$; the trace pairing identifies the Cartier dual of
$k(\rho)$ with $k(\rho^{-1}\omega)$.

\begin{lemma}[KT-CUBIC-TWISTED-COHOMOLOGY]
If $B_{1,\eta\omega^e}$ is a $p$-adic unit, then
\[
 H^2(G_{\mathbf Q,S},k(\eta\psi^{-1}))=0,
 \qquad
 H^1(G_{\mathbf Q,S},k(\eta\psi^{-1}))
   \xrightarrow{\;\sim\;}
 H^1(\mathbf Q_p,k(\eta\psi^{-1})).                \tag{CC.1}
\]
Both $H^1$ spaces have dimension one.  Thus an arbitrary local
$\eta\psi^{-1}$-valued cocycle class at $p$ can be corrected by a
global cocycle in that character, uniquely at the level of classes.
\end{lemma}
\begin{proof}
It suffices to use $k=\mathbf F_p$ and then extend scalars.
Set $L=\mathbf Q(\sqrt D,\mu_p)$ and
$\Delta=\operatorname{Gal}(L/\mathbf Q)$.
For $M=\mathbf F_p(\eta\psi^{-1})$ the Cartier dual is
$M^D=\mathbf F_p(\eta\omega^{e+1})$.
The even mirror in KT-AUXILIARY-UNIT-SELMER-BRIDGE gives
\[
          \operatorname{Cl}(L)\{p\}^{\eta\omega^{e+1}}=0.
                                                               \tag{CC.2}
\]
A class in $\Sha^1_S(\mathbf Q,M^D)$ restricts to an everywhere
unramified homomorphism on $G_L$.  Since $|\Delta|=2(p-1)$ is
prime to $p$, restriction and class field theory embed this group
in $\operatorname{Hom}_\Delta(\operatorname{Cl}(L)/p,M^D)$.
Its target is zero by (CC.2); the character seen by equivariance
is $\eta\omega^{e+1}$ itself, not its inverse.
Poitou--Tate gives $\Sha^2_S(\mathbf Q,M)=0$.

Every local $H^2$ in $S$ vanishes.  At $p$, the dual inertia
character is $\omega^{e+1}$, which is nontrivial because $e+1$
is odd and $p-1$ is even.  At $q\mid D$ the dual has the ramified
quadratic inertia character $\eta$.  At infinity the higher
cohomology vanishes because $p$ is odd.  Local duality and the
Poitou--Tate exact sequence now prove the first assertion in
(CC.1).

The character $\eta\psi^{-1}$ is odd and nontrivial.  Global
Euler characteristic gives $h^1=1$ after the just-proved $h^2=0$.
At $p$ the coefficient character itself is nontrivial on inertia
because $e\not\equiv0\pmod{p-1}$; its local $H^0$ and $H^2$
are zero, so local Euler characteristic gives $h^1=1$.
At every $q\mid D$, both local $H^0$ and $H^2$ vanish and
$q\ne p$, so local $H^1$ is zero.
Consequently the kernel of localization at $p$ is contained in
the residual strict group killed by the same auxiliary-unit
bridge.  Localization is injective, and the equal dimensions
make it an isomorphism.
\end{proof}

\begin{lemma}[KT-CUBIC-RATIONAL-OBSTRUCTION-SPACE]
Let $F=\mathbf Q(\mu_p)$, $\Delta_0=\operatorname{Gal}(F/\mathbf Q)$,
$A_F=\operatorname{Cl}(F)\{p\}$, and $\delta=\omega^{1-e}$.
Let $C_D$ be the subgroup of $A_F$ generated by the $p$-primary
projections of the ideal classes of all primes of $F$ over $q\mid D$.
There is an exact sequence of finite $k$-vector spaces
\[
\begin{split}
0\longrightarrow
 \operatorname{Hom}_{\Delta_0}
   (A_F/(pA_F+C_D),k(\delta))^\vee
\longrightarrow H^2(G_{\mathbf Q,S},k(\psi))
\longrightarrow
 \bigoplus_{\substack{q\mid D\\q^{e-1}\equiv1\ (p)}}k
\longrightarrow0 .
\end{split}                                               \tag{CC.3}
\]
The left term is $\Sha^2_S(\mathbf Q,k(\psi))$.
Here $\vee$ denotes the linear dual, using the trace pairing
to express the finite-module duality over $k$.
\end{lemma}
\begin{proof}
The Cartier dual of $k(\psi)$ is the rational odd character
$k(\delta)$.  Restriction to $G_F$ identifies
$\Sha^1_S(\mathbf Q,k(\delta))$ with the displayed equivariant
Hom space before taking its dual.  To check equality, a strict
class becomes an everywhere unramified homomorphism on $G_F$,
and its vanishing on every decomposition group in $S$ kills the
classes of the primes over $D$.  The unique prime over $p$ is
principal in $F$, so it imposes no further class-group quotient.
Conversely a homomorphism from that quotient extends uniquely
as a cohomology class over $\mathbf Q$ by prime-to-$p$
inflation--restriction.  Its restrictions to the decomposition
groups in $S$ vanish by the same quotient and the principality
at $p$.  All local restriction maps used here are injective
under the extension of degree dividing $p-1$.

Poitou--Tate identifies the dual of that $\Sha^1$ with
$\Sha^2_S(\mathbf Q,k(\psi))$.  Since $1-e$ is odd, the
character $\delta$ is nontrivial both globally and on inertia
at $p$.  Thus its global $H^0$ and the local $H^2$ at $p$
vanish.  At $q\mid D$, $\delta$ is unramified and its Frobenius
value is $q^{1-e}$, so local duality gives a one-dimensional
$H^2$ exactly when $q^{e-1}\equiv1\pmod p$, and zero otherwise.
The real place contributes zero.  The last part of the
Poitou--Tate sequence, whose final term is the zero global
$H^0$ of the dual, proves (CC.3).
\end{proof}

The three auxiliary units at $0,e,2e$ annihilate the specified
$\eta$-twisted class components.  They supply (CC.1), but they
do not supply a vanishing assertion for the rational character
$\delta$ in (CC.3).  In particular, a raw global $H^2$ calculation
must also retain the displayed local terms.  A particular mixed
obstruction known to extend over every prescribed local functor
maps to zero in those local terms and belongs to the left term;
this observation does not prove that its remaining class is
zero or nonzero.

\paragraph{A finite check of the raw-local distinction.}
For $p=7,e=2,D=-43$, the defining Bernoulli sums give
\[
 B_{1,\eta}=-1,\qquad
 B_{1,\eta\omega^2}\equiv5\pmod7,\qquad
 B_{1,\eta\omega^4}\equiv3\pmod7.                    \tag{CC.4}
\]
Thus all three auxiliary units hold, but $43\equiv1\pmod7$
gives a nonzero local quotient in (CC.3).  In particular the
raw $H^2(G_{\mathbf Q,S},\mathbf F_7(\omega^2))$ is nonzero.
The separate standard-library certificate
\path{lean_chain/cubic_cohomology_example.py} evaluates the
complete conductor sums modulo $7^2$ and independently checks
the two residues by exact rational Bernoulli polynomials of
indices $3,5$.  This example makes no claim that a specific
Massey obstruction or a strict Selmer group is nonzero.

\paragraph{Exact imported source scope.}
The unit-to-class and strict-cohomology bridge is the previously
whole-reviewed KT-AUXILIARY-UNIT-SELMER-BRIDGE.  The duality
inputs are Milne, \emph{Arithmetic Duality Theorems}, second
edition, Chapter I: the Cartier dual and $\Sha$ definitions
p.56, Theorem 4.10(a)--(b) pp.56--57 and its full exact sequence,
local Corollary 2.3 p.28, local Theorem 2.8 p.31, and global
Theorem 5.1 / Remark 5.2(c) pp.67--68.  Their complete statements
and exact finite-$S$ scope were independently read; the
foundational duality proofs remain imported.  These lemmas
describe the relevant cohomology spaces, and do not assert
that every abstract element is realized by the source's
higher ordinary deformation equations.

% END evidence/lean_chain/cubic_cohomology.tex


% BEGIN evidence/unitary_audit/reducible_comparison/cubic_massey_condition.tex
\subsection{The normalized ordinary cubic Massey condition}
\label{sec:kt-cubic-massey-condition}

The next condition is a concrete higher lifting condition, not a
vanishing theorem for the original Selmer group.  Its rational
class-group character is different from the character controlling
that original negative root.  The local matching conditions are
retained throughout; raw global $H^2$ is not substituted for the
ordinary obstruction space.

Let $p>3$, let $e\ge2$ be even with $(p-1)\nmid2e$, let
$D<0$ be fundamental and prime to $p$, and use
\[
 \eta=\chi_D,\quad \psi=\omega^e,\quad
 \alpha=\eta\psi,\quad\beta=\psi^{-1},\quad
 S=\{\infty,p\}\cup\{q:q\mid D\}.
                                                               \tag{CM.1}
\]
Let $k/\mathbf F_p$ be finite.  Assume the two auxiliary units
$B_{1,\eta}$ and $B_{1,\eta\omega^e}$.
The third unit at exponent $2e$ is not needed at this degree.
Fix continuous cocycles
\[
 u\in Z^1(G_{\mathbf Q,S},k(\alpha)),\qquad
 v\in Z^1(G_{\mathbf Q,S},k(\beta)),
                                                               \tag{CM.2}
\]
where $[v]$ is zero at $p$ and is unramified on all other finite
inertia groups.  These are the two directions in
KT-GLOBAL-QUADRATIC-MIXED-PERSISTENCE.

\paragraph{The specified graded equivariance.}
Put
\[
\begin{gathered}
 X=E_{12}+E_{23},\quad Y=E_{21}-E_{32},\quad
 Z=\operatorname{diag}(1,-2,1),\\
 U=E_{12}-E_{23},\qquad V=E_{21}+E_{32}.
\end{gathered}                                                 \tag{CM.3}
\]
For $T(a)=\operatorname{diag}(a,1,a^{-1})$, require the algebraic
$\mathbf G_m$-equivariance that sends $(s,t)$ to $(as,a^{-1}t)$
and acts on matrices by conjugation by $T(a)$.  This is an
integer grading, not merely equivariance for the finite group
$k^\times$.  Also require equivariance under
\[
 \kappa(g)=D_c\theta(g)D_c,\qquad
 D_c=\operatorname{diag}(-1,1,-1),\qquad
 \kappa(s,t)=(s,-t),                                    \tag{CM.4}
\]
where $\theta(g)=J(g^t)^{-1}J$ is the unitary outer automorphism.
It fixes the residual representation $r_0$.
On the five Lie vectors in (CM.3), the signs of $\kappa$ are
$+,-,-,-,+$.  Their Galois characters are respectively
$\alpha,\beta,\eta,\psi,\eta\psi^{-1}$.

Use $A_2=k[s,t]/(s,t)^3$ and $A_3=k[s,t]/(s,t)^4$.
The determinant-one, graded, $\kappa$-equivariant logarithms
through degree three necessarily have the form
\[
\begin{split}
 L(g)={}&s u(g)X+t v(g)Y+st w(g)Z\\
        &+s^2t\,m(g)U+st^2\,n(g)V.
\end{split}                                                   \tag{CM.5}
\]
Indeed the weight-two Lie line has $\kappa$ sign minus,
whereas $s^2$ has sign plus, and similarly for $t^2$.
At bidegree $(1,1)$ the trace-zero diagonal space of sign minus
is exactly $kZ$.  At degrees $(2,1)$ and $(1,2)$ the permitted
lines are $kU$ and $kV$.  There are no weight-three Lie lines.
Thus (CM.5) specifies the entire equivariant lifting problem,
rather than omitting other permitted pure-square corrections.

\begin{theorem}[KT-ORDINARY-CUBIC-MASSEY-CONDITION]
For the data (CM.1)--(CM.2), choose the ordinary quadratic
defining cochain $w$ supplied by
KT-GLOBAL-QUADRATIC-MIXED-PERSISTENCE.  There is a well-defined
ordinary cubic obstruction
\[
                  \mathfrak o_3(u,v)\in\mathcal O_\psi.          \tag{CM.6}
\]
Here $\mathcal O_\psi$ is the relative cochain obstruction space
defined below.  Put $F=\mathbf Q(\mu_p)$,
$\Delta_0=\operatorname{Gal}(F/\mathbf Q)$,
$A_F=\operatorname{Cl}(F)\{p\}$, and $\delta=\omega^{1-e}$.
Then
\[
 \dim_k\mathcal O_\psi
   =\dim_k\operatorname{Hom}_{\Delta_0}
                         (A_F/pA_F,k(\delta)).                 \tag{CM.7}
\]
The class (CM.6) is zero if and only if the given quadratic
family lifts over $A_3$ with the grading and involution
(CM.4), the exact ordered ordinary condition at $p$, and
unramified restriction to $G_E$ away from $p$.

Its image in raw global $H^2(k(\psi))$ is
$-[\Omega_+]$, where $\Omega_+$ is the explicit cochain
(CM.10) below.  It lies in the locally trivial subgroup.
If $b\in\{2,4,\ldots,p-3\}$ is congruent to $e$ modulo $p-1$
and $B_b/b$ is a $p$-adic unit, then
$\mathcal O_\psi=0$: every such quadratic family lifts to
this cubic order.  This extra rational unit removes the cubic
obstruction; it does not force $[v]=0$.
\end{theorem}

\begin{proof}
\paragraph{The quadratic defining cochain has no ordinary indeterminacy.}
The equation for $w$ is
\[
 w(gh)-w(g)-\eta(g)w(h)
  =\tfrac12\{u(g)\beta(g)v(h)-v(g)\alpha(g)u(h)\}.         \tag{CM.8}
\]
For a fixed local primitive $v(g)=\beta(g)d-d$ at $p$,
the ordinary condition requires the class of
$w_p=w-d(1+\beta)u/2$ to be unramified.
Two choices of $w$ satisfying these conditions differ by an
$\eta$-cocycle whose localization is unramified.  The
isomorphism (MC.3) makes this cocycle a global coboundary.
Conjugation by $\exp(st aZ)$ realizes that coboundary and
preserves all the indicated local conditions and equivariances.
Consequently the quadratic family is unique up to this
equivalence, once $u,v$ are fixed.  This is stronger than
merely saying that an unrestricted Massey product is defined.

\paragraph{The exact cubic equations.}
For a pair $g,h$, abbreviate
\[
\begin{array}{lll}
 A=u(g),&B=v(g),&C=w(g),\\
 a=\alpha(g)u(h),&b=\beta(g)v(h),&c=\eta(g)w(h),
\end{array}
\qquad F_{g,h}=Ab-Ba.                                   \tag{CM.9}
\]
Since
$[X,Y]=Z$, $[X,Z]=-3U$ and $[Y,Z]=3V$, the degree-three
Baker--Campbell--Hausdorff formula gives
\[
\begin{split}
 \Omega_+(g,h)&=-\tfrac32(Ac-Ca)
                         -\tfrac14(A-a)F_{g,h},\\
 \Omega_-(g,h)&=\tfrac32(Bc-Cb)
                         +\tfrac14(B-b)F_{g,h}.
\end{split}                                                   \tag{CM.10}
\]
The equations for the two remaining cochains are exactly
\[
\begin{split}
 m(gh)-m(g)-\psi(g)m(h)&=\Omega_+(g,h),\\
 n(gh)-n(g)-\eta(g)\psi(g)^{-1}n(h)&=\Omega_-(g,h).
\end{split}                                                   \tag{CM.11}
\]
The two right sides are continuous two-cocycles in the displayed
characters.  This follows either by substitution using (CM.8)
or by associativity: the degree-three kernel is an abelian
module, and the defect of a lift through a square-zero kernel
is a two-cocycle.  All exponential and BCH expressions here
are finite polynomial identities; the denominators divide
$12$, which is invertible.

\paragraph{The twisted cubic direction is removable.}
KT-CUBIC-TWISTED-COHOMOLOGY gives
\[
\begin{gathered}
 H^2(G_{\mathbf Q,S},k(\eta\psi^{-1}))=0,\\
 H^1(G_{\mathbf Q,S},k(\eta\psi^{-1}))
       \xrightarrow{\ \sim\ }H^1(\mathbf Q_p,k(\eta\psi^{-1})).
\end{gathered}                                                 \tag{CM.12}
\]
It also gives zero local $H^1$ in that character at every
$q\mid D$.  Thus the equation for $n$ has a global solution.
After the same local negative-root conjugation used at degree
two, its difference from an ordinary local lift at $p$ is a
cocycle in $k(\eta\psi^{-1})$.  Correct it by a global cocycle
using (CM.12).  At $q\mid D$ any difference from an
unramified-over-$E$ local lift is a coboundary.  Hence this
direction imposes no further global or local condition.
The two cubic directions are independent at this truncation,
because their interactions have degree at least four.

\paragraph{Local lifts of the rational direction exist.}
At $p$, $H^2(\mathbf Q_p,k(\psi))=0$, since its Cartier
dual has nontrivial inertia character $\omega^{1-e}$.
After removing the lower first-order term, the quadratic
family is upper triangular with fixed inertial diagonal.
Its cubic upper correction therefore exists, with arbitrary
local $H^1(k(\psi))$ difference between choices.

At $q\mid D$, local $H^1$ in both $\alpha$ and $\eta$ is
zero.  The degree-two argument therefore conjugates the
family to $\exp(t v(g)Y)r_0(g)$, with $v$ unramified.
This last formula is itself a homomorphism to every
truncation, because $Y$ spans an abelian Lie line and $v$
is a cocycle.  Extending the finite conjugating exponentials
gives an acceptable cubic local lift.  The character $\psi$
is unramified at $q$, and the difference between two such
local lifts, modulo local conjugation, is precisely an
unramified $H^1(k(\psi))$ class.  Indeed after the inertia
of $E$ is killed, the remaining inertia quotient has order
at most two; its positive-degree $p$-cohomology is zero.
There is no ramification outside $S$.

\paragraph{The relative obstruction space.}
Let $M=k(\psi)$, and set
\[
 L_p=H^1(\mathbf Q_p,M),\qquad
 L_q=H^1_{\rm unr}(\mathbf Q_q,M)\quad(q\mid D),
 \qquad L_\infty=0.                                    \tag{CM.13}
\]
Use the ordinary cochain differential
$df(g,h)=f(g)+g f(h)-f(gh)$.
Define $\mathcal O_\psi$ as the vector space of pairs
\[
 (c,(z_v)_{v\in S}),\qquad
 c\in Z^2(G_{\mathbf Q,S},M),\qquad dz_v=c|_{G_{\mathbf Q_v}},
                                                               \tag{CM.14}
\]
modulo pairs
$(df,(f|_{G_{\mathbf Q_v}}+\ell_v)_v)$, where $f$ is a
global continuous one-cochain and each $\ell_v$ is a local
one-cocycle with class in $L_v$.  Local coboundaries are
included in this equivalence.  Its defining exact sequence is
\[
\begin{split}
0\longrightarrow&
 \operatorname{coker}\!\left(
 H^1(G_{\mathbf Q,S},M)\longrightarrow
       \bigoplus_{v\in S}H^1(\mathbf Q_v,M)/L_v\right)\\
\longrightarrow&\mathcal O_\psi
\longrightarrow\Sha^2_S(\mathbf Q,M)\longrightarrow0 .
\end{split}                                                   \tag{CM.15}
\]
This follows directly: forgetting the local primitives gives
the last map; for an exact global cocycle, subtract a global
primitive and retain the local one-cocycle classes.

Take $c=-\Omega_+$ and choose the local primitives supplied
by the acceptable local cubic lifts.  Their classes modulo
the allowed local differences do not depend on the choices.
The resulting pair defines $\mathfrak o_3(u,v)$.
It vanishes precisely when a global $m$ solving (CM.11)
can match all those local lifts modulo (CM.13).
Together with the already removable $n$ direction, this
proves the asserted necessary and sufficient condition.
Changing $w$ by its permitted global coboundary conjugates
the lifting problem and acts trivially on the central
degree-three kernel, so does not change the obstruction.
Changing cocycle representatives of $u,v$ is likewise
realized by the first-order conjugations and does not
change its equivalence class.

\paragraph{Its exact rational arithmetic dimension.}
The Cartier dual is $M^D=k(\delta)$.
At $q\mid D$, local unramified orthogonality gives
$L_q^\perp=H^1_{\rm unr}(\mathbf Q_q,M^D)$.
At $p$ it gives $L_p^\perp=0$.
Put
\[
 Q=\ker\!\left(H^1(G_{\mathbf Q,S},M^D)
       \longrightarrow\bigoplus_{v\in S}
                   H^1(\mathbf Q_v,M^D)/L_v^\perp\right).
                                                               \tag{CM.16}
\]
The Poitou--Tate exact sequence and the perfect local
pairings give
\[
 \dim\operatorname{coker}(\mathrm{loc}/L)
       =\dim Q-\dim\Sha^1_S(\mathbf Q,M^D).
                                                               \tag{CM.17}
\]
For clarity, the annihilator of the image of global
$H^1(M)$ in the local $H^1(M)$ is the image of global
$H^1(M^D)$.  Intersecting this image with $L^\perp$
identifies the dual of the cokernel on the left with
$Q/\Sha^1_S(M^D)$.  Also
$\dim\Sha^2_S(M)=\dim\Sha^1_S(M^D)$ by Poitou--Tate.
Taking dimensions in (CM.15) therefore gives
$\dim\mathcal O_\psi=\dim Q$.

The group $Q$ is strict at $p$ and unramified at every
other finite prime.  Restriction to $G_F$ and class field
theory identify it with
\[
              Q=\operatorname{Hom}_{\Delta_0}
                          (A_F/pA_F,k(\omega^{1-e})).           \tag{CM.18}
\]
Restriction and descent are exact because $p\nmid|\Delta_0|$.
Conversely, an everywhere-unramified class-group homomorphism
is zero on the full decomposition group at the unique prime
over $p$: its Frobenius is the class of a principal prime.
Local restriction is injective, so the descended class is
strict at $\mathbf Q_p$.  No quotient by the prime classes
over $D$ occurs in (CM.18); only unramifiedness, rather
than zero localization, is prescribed there.
This proves (CM.7).  We assert this exact dimension and the
canonical cochain obstruction (CM.14); we do not obtain a
canonical class-group-valued functional merely by comparing
two short exact sequences.

Finally the carried Herbrand implication, explicitly (AB.12)
in KT-RATIONAL-POSITIVE-ROOT-UNIT, kills
$A_F^{\omega^{1-e}}$ when $B_b/b$ is a unit.
Thus (CM.7) makes $\mathcal O_\psi$ zero, proving the last
assertion.
\end{proof}

\paragraph{An explicit repeated Massey identity.}
The rational obstruction is not just assigned a character.
With pointwise product $uv$ define the two one-cochains
\[
 a_{12}=-w-\tfrac12uv,\qquad
 a_{23}=w-\tfrac12uv.
                                                               \tag{CM.19}
\]
Equation (CM.8) and the cocycle laws give
$da_{12}=u\smile v$ and $da_{23}=v\smile u$.
Thus
\[
 \Theta(u,v,u)=u\smile a_{23}+a_{12}\smile u
                      \in Z^2(G_{\mathbf Q,S},k(\psi))
                                                               \tag{CM.20}
\]
is a specified defining value of the repeated triple
Massey product.  Direct expansion gives the exact cochain
identity
\[
             \Omega_+=-\tfrac32\Theta(u,v,u)
                               +\tfrac14d(u^2v).              \tag{CM.21}
\]
Consequently the raw image of (CM.6) is
$\tfrac32[\Theta(u,v,u)]$.
The usual unrestricted defining-system indeterminacy is not
silently discarded: the ordinary quadratic normalization
above fixes the allowed $w$ modulo global coboundary.
For fixed $[u]$, the resulting relative obstruction is linear
in $[v]$, and scaling $[u]$ by $\lambda$ scales it by
$\lambda^2$.  This follows either from (CM.8)--(CM.10) by
choosing linear cochain splittings, whose choices do not
affect (CM.6), or from the two-parameter grading.

\paragraph{What would be needed for an exclusion.}
The first surviving arithmetic requirement is therefore
the vanishing of this particular relative cubic class,
whose dimension is the rational $\omega^{1-e}$ class
component.  It is more precise than the existence of a
raw $H^2$ group or an unrestricted Massey defining system.
The three auxiliary units do not evaluate this class.
Nor does a nonzero rational class component prove that
the value (CM.6) is nonzero.

An exclusion would require an actual arithmetic
nonvanishing or injectivity assertion for this map,
together with a proof that the relevant source is forced
to admit the specified two-parameter first-order jet.
Neither assertion is proved here.  In particular this
calculation is not a new upper bound for the original
strict Selmer group, and a characteristic-zero source
point is not inferred from a cubic lift.

\paragraph{Exact verification and imported scope.}
The separate certificate
\path{cubic_massey_check.py} verifies the finite BCH
identity, both two-cocycle identities, all five involution
signs, the determinant, both Massey defining equations,
and (CM.21) by exact symbolic polynomial arithmetic.
The arithmetic cohomology input (CM.12) is the separately
whole-reviewed KT-CUBIC-TWISTED-COHOMOLOGY.
For (CM.15)--(CM.18) the imported duality is Milne,
\emph{Arithmetic Duality Theorems}, second edition,
Chapter~I, Theorem~4.10(a)--(b) and its complete exact
sequence on pp.~56--57.  Local unramified orthogonality
is Theorem~2.6 on p.~30, with its complete proof continuing
to p.~31; at $q\ne p$ the finite dual takes values in
unramified units, so its $M^d$ equals $M^D$ here.
Those complete passages were read for this application.
The foundational duality theorems are imported.

% END evidence/unitary_audit/reducible_comparison/cubic_massey_condition.tex


% BEGIN evidence/unitary_audit/reducible_comparison/cubic_shuffle_comparison.tex
\subsection{An exact comparison with the consecutive repeated product}
\label{sec:kt-cubic-shuffle}

Retain all the data and local conditions of
KT-ORDINARY-CUBIC-MASSEY-CONDITION.  All cochains use
\(df(g,h)=f(g)+g f(h)-f(gh)\), and the cup convention is
\((x\smile y)(g,h)=x(g)\,g y(h)\).
Pointwise multiplication of one-cochains is denoted by
juxtaposition.  In particular it is not the cup product.

\begin{proposition}[KT-ORDINARY-CUBIC-SHUFFLE]
Define
\[
 h=-\frac12u^2,\qquad a=-w-\frac12uv,\qquad
 b=w-\frac12uv .                                      \tag{CS.1}
\]
Then \(dh=u\smile u\), \(da=u\smile v\), and
\(db=v\smile u\).  The cocycle
\[
 \Phi=u\smile a+h\smile v\in Z^2(G_{\mathbf Q,S},k(\psi))
                                                               \tag{CS.2}
\]
is a specified consecutive defining value for
\(\langle u,u,v\rangle\).  With
\(\Theta=u\smile b+a\smile u\) as in (CM.20), the exact
cochain identities are
\[
 \Theta+2\Phi=-d(ua),\qquad
 \Omega_+=3\Phi-dR,\qquad
 R=\frac32uw+\frac12u^2v.                              \tag{CS.3}
\]

If \((c,(z_v)_v)\), with \(c=-\Omega_+\), represents
the ordinary obstruction in (CM.14), put
\[
 z^{\mathrm{rep}}_v=\frac{R|_{G_{\mathbf Q_v}}-z_v}{3}.
                                                               \tag{CS.4}
\]
Then \(d z^{\mathrm{rep}}_v=\Phi|_{G_{\mathbf Q_v}}\),
and the relative class
\[
 \mathfrak m_{\mathrm{rep}}(u,u,v)
       =[(\Phi,(z^{\mathrm{rep}}_v)_v)]\in\mathcal O_\psi
                                                               \tag{CS.5}
\]
is well defined under the same permitted global and local
equivalences as the ordinary cubic obstruction.  Exactly
\[
             \mathfrak o_3(u,v)=-3\mathfrak m_{\mathrm{rep}}(u,u,v).
                                                               \tag{CS.6}
\]
Thus changing the order of the repeated entries introduces
no additional gap for this specified relative obstruction.
\end{proposition}

\begin{proof}
The cocycle identity for \(u\) gives
\(d(u^2)=-2u\smile u\).  The other two defining equations
are (CM.19).  The differential and cup rule show that the
two terms in \(d\Phi\) cancel, so \(\Phi\) is a cocycle.

For a direct sign check use the six variables of (CM.9).
They give
\[
\begin{split}
 \Phi(g,h)&=-Ac-\frac12Aab-\frac12A^2b,\\
 \Theta(g,h)&=Ac-Ca-\frac12Aab-\frac12ABa,\\
 d(uw)(g,h)&=-Ac-Ca-\frac12(A+a)(Ab-Ba),\\
 d(u^2v)(g,h)&=-A^2b-2ABa-2Aab-Ba^2.
\end{split}                                                   \tag{CS.7}
\]
Substitution proves the first identity in (CS.3), since
\(-ua=uw+u^2v/2\).  Substituting that identity in (CM.21)
gives the second identity, with precisely the displayed \(R\).
All denominators are units because \(p>3\).

For the relative assertion, (CS.3) gives
\(c=-3\Phi+dR\), and (CS.4) is therefore a local primitive
of \(\Phi\).  At the level of pairs one has the literal equality
\[
 (c,(z_v)_v)=-3(\Phi,(z^{\mathrm{rep}}_v)_v)
                              +(dR,(R|_v)_v).                 \tag{CS.8}
\]
The last pair is zero in (CM.14).  Since the allowed local
spaces \(L_v\) are vector subspaces and \(3\) is invertible,
any change by an allowed local cocycle transports to an
allowed local cocycle.  A change by a global primitive
transports to the corresponding global primitive.
The already proved independence of \(\mathfrak o_3\),
together with (CS.8), proves the asserted independence
of (CS.5).  This transports the actual ordinary local
choices; it does not replace them by arbitrary local
primitives for an unrestricted Massey product.
\end{proof}

\paragraph{The exact relation to Sharifi's conventions.}
Here is a further algebraic comparison, separate from the
proper norm-system requirement in Sharifi's theorem.
Restrict to \(K=\mathbf Q(\sqrt D,\mu_p)\), so the
coefficient characters become trivial, and first take
\(k=\mathbf F_p\).  A fixed primitive \(p\)-th root
identifies the coefficients with \(\mu_p\).
In the upper-unitriangular four-by-four group modulo its
central \((1,4)\) entry, the functions
\[
 \kappa_{12}=\kappa_{23}=u,\quad\kappa_{34}=v,\quad
 \kappa_{13}=-h=\frac12u^2,\quad\kappa_{24}=-a
                                                               \tag{CS.9}
\]
give a homomorphism.  Indeed Sharifi's defining equation
(his condition (c), printed p.~18) is
\(d\kappa_{ij}=-\sum_l\kappa_{il}\smile\kappa_{lj}\),
which is exactly (CS.1) with these signs.
Conjugation by \(I-E_{12}/2\) changes only the noncentral
entry \(\kappa_{13}\), replacing it by
\(\binom{u}{2}=(u^2-u)/2\).  Its values therefore lie
in the subgroup \(G_1^{(2)}/Z_1^{(2)}\) generated in
Sharifi's notation by the images of
\(X=I+E_{12}+E_{23}\) and \(Y=I+E_{34}\).
There is no remaining discrepancy in the repeated-entry
block or its binomial normalization.

With the convention (4.1) of that paper, the cocycle
of this last defining system is
\[
 \kappa_{12}\smile\kappa_{24}
       +\kappa_{13}\smile\kappa_{34}
       =-\Phi-\frac12d a .                              \tag{CS.10}
\]
Thus its raw class is \(-[\Phi]\), and the raw ordinary
cubic obstruction is three times this class.  The same
global-coboundary transport as (CS.8) adjusts the local
primitives.  For larger finite coefficient fields the
cochain identities remain valid over that field; a
literal application of a theorem with \(\mu_p\)
coefficients must still identify the requisite scalar
components rather than treating a general \(k\)-valued
character as one cyclic degree-\(p\) extension.

\begin{corollary}[KT-CUBIC-BOCKSTEIN-NORM-IDENTIFICATION]
Suppose \(k=\mathbf F_p\) and \([u]\ne0\).
Over \(K=\mathbf Q(\sqrt D,\mu_p)\), let \(S_K\) be
the primes above \(S\), set \(G=G_{K,S_K}\), and write
\(\chi=u|_G\).  This is a surjection onto \(H=\mathbf F_p\).
Let \(\mathcal R=\mathbf F_p[H]\), choose \(\tau\in H\)
corresponding to \(1\in\mathbf F_p\), and put \(x=[\tau]-1\),
\(I=(x)\).  The latter notation means the element of the
coimage corresponding to \(1\); it does not require a
finite-order lift in \(G\).
Then
\[
 f=v-a x\in Z^1(G,\mathcal R/I^2),\qquad
 \Psi^{(2)}([f])=-[\Phi|_G]x^2,                       \tag{CS.11}
\]
where \(G\) acts on \(\mathcal R\) by multiplication
by its image in \(H\), and \(\Psi^{(2)}\) is the
connecting map for the augmentation sequence through
\(I^3\).  The conjugated system (CS.9) is a binomial proper
system in the sense of Lam--Liu--Sharifi--Wake--Wang.
Its realization by a Sharifi norm system additionally
uses its zero absolute obstruction and
KT-CUBIC-MATCHED-NORM-SYSTEM below.

In particular, choose a primitive root \(\zeta\in\mu_p\)
and Kummer representatives \(A_0,B_0\in K^\times\)
for \(\chi,v|_G\) using \(\zeta\).  Put
\(L=K(A_0^{1/p})\), of degree \(p\), and let
\(\Omega_S\) be the maximal extension of \(K\)
unramified outside \(S_K\).  The matched-system
argument supplies a norm-system representative
\(y\in L^\times\) with
\[
 N_{L/K}y=B_0,\qquad D^{(1)}y\in\Omega_S^{\times p}.
                                                               \tag{CS.12}
\]
Sharifi's Theorem~5.4 applies to this system.  If
\(\mathfrak Y\) is the ideal appearing there and
\(P^{(1)}_{1,L/K}\) is its lower-product subgroup,
then the raw image \(\overline{\mathfrak o}_3\) of
the relative ordinary obstruction satisfies
\[
 \iota_\zeta\!\left(\operatorname{res}_{K/\mathbf Q}
                     \overline{\mathfrak o}_3\right)
     \equiv -3[N_{L/K}\mathfrak Y]
             \pmod{P^{(1)}_{1,L/K}}.                  \tag{CS.13}
\]
Here \(\iota_\zeta:\mathbf F_p\to\mu_p\) sends
\(1\) to \(\zeta\); equivalently one untwists
Sharifi's triple coefficient by \(\zeta^{\otimes2}\).
A nonzero right-hand coset in (CS.13) is sufficient
for a nonzero ordinary relative cubic obstruction.
No such nonzero coset is asserted here.
\end{corollary}

\begin{proof}
Restriction is injective in degree one because
\([K:\mathbf Q]\) is prime to \(p\).  The nonzero
\(\mathbf F_p\)-valued character \(\chi\) is therefore
surjective.  The action of \(g\in G\) on the group
ring is multiplication by \((1+x)^{\chi(g)}\).
The equations \(dv=0\), \(da=u\smile v\) show
directly that \(f\) is a cocycle modulo \(x^2\).
Use the same expression as a one-cochain modulo \(x^3\).
Its differential has coefficient of \(x^2\) equal to
\[
 u(g)(-a(h))+\binom{u(g)}2v(h),
                                                               \tag{CS.14}
\]
which is precisely (CS.10).  This proves (CS.11),
including its sign, without any comparison of
cohomology dimensions.

Definition~3.3.2 and the complete Lemma~3.3.3 of
Lam--Liu--Sharifi--Wake--Wang identify the binomial
block in (CS.9) as a proper partial-system extension.
In their Section~4.3 take \(n=2\), \(A=R=\mathbf F_p\):
the size condition is \(2p<p^2\), and \(p>3\)
ensures it.  Their Theorem~4.3.1 gives exactly the
same Bockstein class and fixes the binomial matrix
conventions.  That terminology alone does not
realize every specified last column by a norm
representative.  Here the additional hypothesis
needed by KT-CUBIC-MATCHED-NORM-SYSTEM holds:
the ordinary obstruction is locally zero everywhere
by (CM.14), including unramified places outside
\(S\).  After restriction to \(K\), its inflation
to \(H^2(G_K,\mu_p)=\operatorname{Br}(K)[p]\)
is zero by the Brauer local-global injection.
By (CS.10) its value is three times the Bockstein,
so the latter also has zero absolute inflation.
The matched-system lemma lifts the specified
truncated cocycle through the absolute regular
module, then uses Shapiro and the exact Kummer-dual
pairing to give (CS.12).  This argument supplies
the norm representative for the actual defining
class.  The lemma itself is independent of the
norm formula being applied here.

Apply Sharifi's Theorem~5.4 with his \(m=n=1\)
and repeated length \(k=2\).  Since \(2<p-1\),
its exceptional boundary term is absent.  Its
formula for the proper triple value is the
negative norm-ideal class.  Equations (CS.6)
and (CS.10) make the raw ordinary value three
times that triple value, proving (CS.13).
If its coset is nonzero then the restricted raw
class is nonzero, so the relative class cannot
vanish.  This establishes the final sufficient
condition without asserting that it is met.
\end{proof}

\paragraph{The remaining arithmetic requirements.}
The \(\mathbf F_p\)-valued ordering and norm-system
issues are discharged by (CS.3), (CS.9), the zero
absolute inflation, and the explicit matched-system
lemma.  The raw formula (CS.13) by itself does not
determine the local correction in (CM.15).
KT-CUBIC-EXCISION and KT-CUBIC-P-ONLY-NORM below
retain that correction in a class with ramification
only above \(p\), and give an exact ideal-class
formula when the norm representative is matched.
No norm-ideal class has been evaluated here;
neither nonvanishing nor injectivity in the
original strict class has been proved.
General finite \(k\)-valued input still
requires a scalar or cyclic reduction to use this
particular norm formula.  These are equivalences
and an available sufficient evaluation criterion,
not a new arithmetic exclusion or a source-to-jet
theorem.

The primary passages inspected for this comparison are
Sharifi, \emph{Massey products and ideal class groups},
author version dated 26 July 2006: Section~4, printed
pp.~17--22, especially condition (c), (4.1), the two
generators, Proposition~4.2 and the proper-system
definition; and Theorem~5.4 with its complete proof,
printed pp.~26--27.  The published article is
\emph{J.~reine angew.~Math.}~603 (2007), 1--33.
For the further imported properness identification
we read Lam--Liu--Sharifi--Wake--Wang,
\emph{Generalized Bockstein maps and Massey products},
\emph{Forum Math.~Sigma}~11 (2023), e5:
Definition~3.3.2, the complete proofs of
Lemma~3.3.3 and Theorem~3.3.4, and the complete
Section~4.3 (author-PDF pp.~24--27 and 30--32).
The identification paragraph and Theorem~4.3.1
were also checked in the published Cambridge
version.  Its Section~6 absolute-Galois vanishing
statements do not imply vanishing for this fixed
ramification set with the prescribed local
primitives.

% END evidence/unitary_audit/reducible_comparison/cubic_shuffle_comparison.tex


% BEGIN evidence/unitary_audit/reducible_comparison/cubic_excision_norm.tex
\subsection{Excision and the norm class with ramification only above p}
\label{sec:kt-cubic-excision-norm}

Retain the hypotheses, coefficients, and specified ordinary quadratic
family of KT-ORDINARY-CUBIC-MASSEY-CONDITION.  In particular \(p>3\),
\(e\) is even, \(\psi=\omega^e\), \(E=\mathbf Q(\sqrt D)\),
\(p\nmid D\), and
\(S=\{\infty,p\}\cup\{q:q\mid D\}\).
Write \(F=\mathbf Q(\mu_p)\), \(K=EF\), and
\(\delta=\omega^{1-e}\).  For a number field \(J\), the notation
\(G_{J,p}\) means its Galois group unramified outside the primes above
\(p\) and the infinite primes.  Let \(A_F\) denote the \(p\)-primary
ideal class group of \(F\).  This section proves an identification of
the actual relative obstruction, not just an equality of dimensions.

\begin{theorem}[KT-CUBIC-EXCISION]
For every finite coefficient field \(k/\mathbf F_p\) there is a
natural isomorphism
\[
 \mathcal O_\psi\ \simeq\ H^2(G_{\mathbf Q,p},k(\psi)).
                                                        \tag{CE.1}
\]
For \(k=\mathbf F_p\), fixing a primitive root \(\zeta\in\mu_p\)
then gives an isomorphism
\[
 \mathcal C_\zeta:\mathcal O_\psi
      \xrightarrow{\ \sim\ }(A_F/pA_F)^{\delta}.
                                                        \tag{CE.2}
\]
There is no quotient in (CE.2) by classes of primes above \(D\).
For larger finite \(k\), (CE.2) extends by scalars.
\end{theorem}

\begin{proof}
Put \(U=\operatorname{Spec}\mathbf Z[1/p]\),
\(V=\operatorname{Spec}\mathbf Z[1/pD]\), and \(M=k(\psi)\).
The corresponding finite locally constant sheaf extends to \(U\).
For \(q\mid D\), let \(U_q=\operatorname{Spec}(\mathbf Z_{(q)})^h\),
with fraction field \(\mathbf Q_q^h\) and residue field \(\mathbf F_q\).
Etale excision supplies the distinguished triangle
\[
 R\Gamma(U,M)\longrightarrow
 R\Gamma(V,M)\oplus\bigoplus_{q\mid D}R\Gamma(\mathbf F_q,M)
 \longrightarrow\bigoplus_{q\mid D}R\Gamma(\mathbf Q_q,M).
                                                        \tag{CE.3}
\]
The last arrow is the difference of restriction and unramified
inflation.  Here is the precise source of this local gluing statement.
For \(j:V\hookrightarrow U\), the complex form of Milne's
Lemma II.2.4 gives
\(R\Gamma(U,j_!M)\to R\Gamma(V,M)\to
\bigoplus R\Gamma(\mathbf Q_q^h,M)\).
The exact sequence \(0\to j_!M\to M\to i_*i^*M\to0\)
gives the second triangle with the same first term.
Comparing them gives (CE.3).
Milne II.1.1(b) identifies cohomology over \(U_q\) with that over
its residue field; the henselian fraction field and its completion
have equivalent finite separable extensions and the same absolute
Galois group, as recalled explicitly in the proof of I.2.14, p.~37.
Thus the final terms are exactly the completed local groups used
in (CM.14), with their actual unramified inflation maps.

One can read the degree-two group of the mapping fibre in (CE.3)
without discarding any local data.  A representative consists of
a global two-cocycle \(c\), unramified local two-cocycles \(b_q\),
and local one-cochains \(z_q\) with
\(dz_q=c|_q-b_q\).
Since \(H^2(\mathbf F_q,M)=0\), subtract primitives of the \(b_q\)
and set them to zero.  The remaining ambiguity in those primitives
is exactly unramified local \(H^1\), in addition to local
coboundaries.  The resulting quotient is therefore the pairs
\((c,(z_q)_{q\mid D})\) of (CM.14), with precisely its stated
equivalence and its allowed spaces \(L_q\).

Forgetting the \(p\)-primitive in (CM.14) is an isomorphism on
that quotient: local duality gives
\(H^2(\mathbf Q_p,M)=H^0(\mathbf Q_p,k(\delta))^\vee=0\),
because \(1-e\) is odd and is not zero modulo \(p-1\), while
\(L_p\) is all local \(H^1\).  The infinite place has no
positive-degree \(p\)-cohomology.  Finally Milne II.2.9 applies
on both affine schemes, since \(p\) is invertible, and identifies
their etale cohomology with the corresponding Galois cohomology.
This proves (CE.1) with its natural map, rather than by comparing
dimensions.

For (CE.2), the cover
\(U_F=\operatorname{Spec}\mathcal O_F[1/p]\to U\) is finite
etale of degree \(p-1\).  Exact averaging gives
\[
 H^2(U,\mathbf F_p(\psi))
     \simeq H^2(U_F,\mathbf F_p)^{\psi^{-1}}.
                                                        \tag{CE.4}
\]
The identification \(1\mapsto\zeta\) takes the right side to
\(H^2(U_F,\mu_p)^\delta\): it contributes the additional
cyclotomic character \(\omega\).  The Kummer exact sequence is
\[
 0\longrightarrow \operatorname{Pic}(U_F)/p
 \longrightarrow H^2(U_F,\mu_p)
 \longrightarrow \operatorname{Br}(U_F)[p]\longrightarrow0.
                                                        \tag{CE.5}
\]
There is one prime of \(F\) above \(p\), and \(F\) has no real
places.  The local invariant sequence for the Brauer group shows
\(\operatorname{Br}(U_F)[p]=0\): a single local invariant must
have sum zero.  The prime above \(p\) is principal, generated
by \(1-\zeta\), so \(\operatorname{Pic}(U_F)=\operatorname{Cl}(F)\).
Consequently (CE.4)--(CE.5) give exactly (CE.2).  Extending the
continuous cochain complexes by the finite field \(k\) proves
the scalar extension statement.
\end{proof}

\begin{lemma}[KT-CUBIC-MATCHED-NORM-SYSTEM]
Let \(J\) be a field of characteristic zero containing \(\mu_p\),
let \(\Gamma\) be a quotient of \(G_J\), and let
\(\chi:\Gamma\to H=\mathbf F_p\) be surjective.
Set \(\mathcal R=\mathbf F_p[H]=\mathbf F_p[x]/(x^p)\),
where \(x=[\tau]-1\) and \(\tau\) corresponds to \(1\in H\),
and write \(I=(x)\).
Suppose \(f_2\in Z^1(\Gamma,\mathcal R/I^2)\), with augmentation
\(v\), has a Bockstein class whose inflation to
\(H^2(G_J,\mathbf F_p)\) is zero.  Let \(L/J\) be cut out
by \(\chi\), and fix a Kummer representative \(B_0\in J^\times\)
for \(v\), using \(\zeta\).

There is a norm representative \(y\in L^\times\),
\(N_{L/J}y=B_0\), for which Sharifi's truncated norm system
represents the specified class \([f_2]\).  In particular its
proper triple value is the actual Bockstein of \(f_2\),
not merely its coset modulo lower products.  If
\(\Gamma=\operatorname{Gal}(\Omega/J)\), then
\(D^{(1)}y\in\Omega^{\times p}\).
\end{lemma}

\begin{proof}
All lifting in this paragraph takes place on the absolute
group \(G_J\).  The cyclic norm-residue sequence is exact at
its middle degree-two term:
\[
 H^1(G_J,\mathbf F_p)
 \xrightarrow{\chi\smile-}H^2(G_J,\mathbf F_p)
 \xrightarrow{\operatorname{res}}H^2(G_L,\mathbf F_p).
                                                        \tag{CE.9}
\]
This is the absolute Galois input stated in
Lam--Liu--Sharifi--Wake--Wang, Definition~6.1.2 and the
following paragraph, p.~48; their Proposition~6.1.5 uses
the same coefficient correction as below.  We do not
apply (CE.9) to a group with restricted ramification.

The assumed zero Bockstein first lifts the inflated
\([f_2]\) to \(H^1(G_J,\mathcal R/I^3)\).
For \(3\le n<p\), suppose \(f_n\) is a cocycle lifting it
through \(I^n\), and let \(\alpha_n\) be the obstruction
to lifting through \(I^{n+1}\).  On \(G_L\) the entire
regular module is trivial, so that coefficient sequence
splits there.  Thus \(\operatorname{res}\alpha_n=0\).
By (CE.9), write \(\alpha_n=\chi\smile\gamma_n\)
for a scalar one-cocycle \(\gamma_n\).
Replacing \(f_n\) by \(f_n-\gamma_n x^{n-1}\)
kills the obstruction and preserves its image modulo
\(I^2\).  The sign is explicit: for the differential used
throughout this report,
\[
 d(\gamma_n x^{n-1})(g,h)
       =\chi(g)\gamma_n(h)x^n\pmod{x^{n+1}}.
                                                       \tag{CE.10}
\]
The scalar coefficient is a cocycle and
\((1+x)^{\chi(g)}=1+\chi(g)x+O(x^2)\), which proves
(CE.10).  Iterating gives
\(F_0\in Z^1(G_J,\mathcal R)\) whose image in
\(H^1(G_J,\mathcal R/I^2)\) is the inflated \([f_2]\).
If the representative differs by a coboundary, lift
its zero-cochain to \(\mathcal R\) and subtract that
coboundary.  Its projection can therefore be taken
literally equal to the inflated \(f_2\).

Shapiro identifies \([F_0]\) with a Kummer class
\(y\in L^\times/L^{\times p}\); augmentation is
corestriction, so \(N_{L/J}y\equiv B_0\pmod{J^{\times p}}\).
We spell out the truncated-system identification.
Evaluate \(F_0|_{G_L}\) at the coefficient of the
identity of \(H\).  This is the Kummer character of
\(y\).  Equivariance makes its coefficient at
\(\sigma^j\) the Kummer character of \(\sigma^j y\).
Thus its image in the regular module is
exactly the map \(\phi(\lambda)=\sum_h\lambda(h)h\)
of Sharifi's Lemma~3.1, with its stated inversion
pairing, used in (4.2) and Lemma~4.1.
Projecting to \(\mathcal R/I^3\) and taking the action
\((1+x)^\chi\) gives the binomial four-by-four system;
its two last-column entries modulo the central entry
are precisely the two coefficients of \(f_2\).

One normalization in Lemma~4.1 causes no discrepancy.
For a lift \(\sigma\in G_J\) with \(\chi(\sigma)=1\),
\(F_0(\sigma)-v(\sigma)\) lies in \(I=x\mathcal R\).
Conjugation by a constant element of the additive
regular module removes that difference, since its
coboundary at \(\sigma\) is multiplication by \(x\).
This puts the image of \(\sigma\) into Sharifi's
\(X+v(\sigma)E_{34}\) form.  It changes \(f_2\)
only by a coboundary and preserves its Bockstein class.
The same conjugation descends on \(\Gamma\), since
the action \(\chi\) does.  This verifies matching of
the specified truncated class, rather than just
existence of some proper system with first entries
\(\chi,\chi,v\).

Choose an actual representative \(y\).  If its norm is
\(B_0c^p\), divide \(y\) by the base-field element \(c\).
Under Shapiro this subtracts a cocycle supported on
the norm element
\(\sum_{h\in H}[h]=x^{p-1}\), so it does not change
the projection modulo \(I^2\).  The norm is now
exactly \(B_0\).  Finally the matched system modulo
\(I^2\) factors through \(\Gamma\), so on \(G_\Omega\)
the full regular cocycle has values in \(I^2\).
By Sharifi's Lemma~3.3, \(D^{(1)}\) belongs to
\((I^2)^\perp\) for precisely the pairing just used.
The Kummer character of \(D^{(1)}y\) consequently
vanishes on \(G_\Omega\).  This proves
\(D^{(1)}y\in\Omega^{\times p}\) and the required
proper norm condition.
The full regular cocycle was constructed only on
\(G_J\), and its coefficients of degree at least two
may ramify outside the set allowed by \(\Gamma\).
Only its specified truncation is required to descend.
In particular \(D^{(2)}y\) need not have a \(p\)-th
root in \(\Omega\); the restricted-ramification
transgression being computed need not vanish.
\end{proof}

\begin{theorem}[KT-CUBIC-P-ONLY-NORM]
Assume \(k=\mathbf F_p\) and \([u]\ne0\).  Restrict the specified
ordinary quadratic family to \(K=E(\mu_p)\).
Its binomial defining system (CS.9), after the stated conjugation,
factors through \(G_{K,p}\).  Thus the proper norm construction
of KT-CUBIC-BOCKSTEIN-NORM-IDENTIFICATION can be made with
\(S_K\) consisting only of the primes above \(p\).

Let \(L/K\) be the degree-\(p\) extension cut out by \(u|_{G_{K,p}}\),
and use this smaller ramification set in its proper norm system.
Let \(\mathfrak Y\) and \(P^{(1)}_{1,L/K,p}\) be Sharifi's
ideal and lower-product subgroup for this system.  There is a
specified projection
\[
 \pi_\delta:H^2(G_{K,p},\mu_p)
           \longrightarrow (A_F/pA_F)^\delta
                                                        \tag{CE.6}
\]
which is inverse to extension on the \(E\)-positive
\(\delta\)-component.  The actual ordinary obstruction satisfies
\[
 \mathcal C_\zeta(\mathfrak o_3(u,v))
 \equiv -3\pi_\delta([N_{L/K}\mathfrak Y])
          \pmod{\pi_\delta(P^{(1)}_{1,L/K,p})}.
                                                        \tag{CE.7}
\]
In particular a nonzero right-hand coset proves
\(\mathfrak o_3(u,v)\ne0\).  All local correction information
in (CM.15) is retained by (CE.2) and this smaller ramification
set.  Moreover one may choose \(y\) to match the specified
ordinary system by KT-CUBIC-MATCHED-NORM-SYSTEM.  For that
matched choice, the ideal furnished by Sharifi's proof satisfies
the stronger equality
\[
 \mathcal C_\zeta(\mathfrak o_3(u,v))
            =-3\pi_\delta([N_{L/K}\mathfrak Y])
       \quad\hbox{in }(A_F/pA_F)^\delta .
                                                       \tag{CE.11}
\]
The quotient in (CE.7) is necessary when a proper norm
representative is used without retaining that matching.
\end{theorem}

\begin{proof}
At a prime \(q\mid D\), the actual ordinary quadratic family
is unramified after restriction to \(E\).  This was proved in
KT-GLOBAL-QUADRATIC-MIXED-PERSISTENCE: locally it is conjugate to
\(\exp(t v(g)Y)r_0(g)\), and both terms are one on inertia
over \(E\).  Identity on inertia is preserved by that fixed
conjugation.  Over \(K\), all residual characters are trivial.
Hence the restricted quadratic homomorphism is one on every
inertia group away from \(p\), and factors through \(G_{K,p}\).
Its logarithmic coordinates \(u,v,w\), and therefore the
polynomial binomial entries in (CS.9), factor through that
same quotient.  This statement uses the actual local family,
not just the vanishing of its inertia cohomology classes.

First restrict to \(G_{K,S_K^{\rm all}}\), where
\(S_K^{\rm all}\) contains the primes above \(pD\).
This is the finite-index subgroup of \(G_{\mathbf Q,S}\),
so restriction is injective in degree one by restriction
and corestriction, because \([K:\mathbf Q]\) is prime to \(p\).
Then descend through its quotient \(G_{K,p}\).
The descended \(\mathbf F_p\)-valued
character \(u\) remains nonzero and therefore cuts out
a degree-\(p\) extension unramified outside \(p\).
The binomial Bockstein calculation (CS.11)--(CS.14)
therefore applies on \(G_{K,p}\) itself.  The existence
of a matched norm representative will follow below;
it is not inferred solely from the terminology
``proper defining system.''

We next check that this is the image of the \emph{relative}
class in (CE.1).  The excision construction (CE.3) is functorial
for the finite map
\(U_K=\operatorname{Spec}\mathcal O_K[1/p]\to U\),
even though that map is ramified at primes over \(D\).
Use the corresponding local gluing triangles before and after
this base change.  Its pullback on a relative representative
restricts the global cocycle and all the local primitives.
At a prime above \(D\), the chosen acceptable cubic local
lift is unramified over \(K\).  Thus its central degree-three
coordinate is an unramified primitive of the pulled-back
\(-\Omega_+\).  Since the quadratic coordinates already
factor through \(G_{K,p}\), this is precisely the relative
representative of their raw cubic obstruction on \(G_{K,p}\).
Changing either primitive makes an unramified local cocycle,
which is killed by the equivalence in (CE.3).
There is no lost primitive at \(p\): it was forgotten on the
\(\mathbf Q\)-side because \(L_p\) was full.  On the
\(K\)-side this paragraph uses the gluing complex only at
primes over \(D\), so it does not incorrectly impose a
vanishing condition on \(H^2(K_w,\mathbf F_p)\) at \(w\mid p\).
It follows that the image of (CE.1) on \(U_K\) is exactly
the raw cubic class of the descended defining system.

For completeness that pullback does not lose the rational
class in (CE.2).  The fields \(E\) and \(F\) are disjoint,
since \(p\nmid D\), and
\(\operatorname{Gal}(K/\mathbf Q)
  =\operatorname{Gal}(E/\mathbf Q)\times\operatorname{Gal}(F/\mathbf Q)\).
Let \(e_{+,\delta}\) be the averaging projector onto the
component trivial on the first factor and equal to \(\delta\)
on the second.  It is defined over \(\mathbf F_p\), since
the group order \(2(p-1)\) is a unit.
The classes of primes of \(K\) above \(p\), and the local
Brauer invariant module at those primes, are permutation
modules on which the cyclotomic factor acts trivially.
Indeed \(F/\mathbf Q\) is totally ramified at \(p\), while
\(E/\mathbf Q\) is unramified there.  Since \(\delta\ne1\),
neither module has a \((+,\delta)\)-component.  Kummer and
the Brauer local invariant sequence therefore identify
\[
 e_{+,\delta} H^2(G_{K,p},\mu_p)
       \simeq (\operatorname{Cl}(K)/p)^{{+,\delta}}.
                                                        \tag{CE.8}
\]
Extension and norm on ideal classes have compositions
\(N_{K/F}\circ\mathrm{ext}=2\) and
\(\mathrm{ext}\circ N_{K/F}=1+c\), where
\(c\) generates \(\operatorname{Gal}(K/F)\).
They identify the right side of (CE.8) with
\((A_F/pA_F)^\delta\).  Define \(\pi_\delta\) by first
applying \(e_{+,\delta}\), then (CE.8), then
\(\tfrac12N_{K/F}\).  This proves (CE.6), including
injectivity of the pullback on the image of (CE.2).
There is no need to infer this injectivity from a putative
unramified Galois cover at the primes over \(D\).

The raw ordinary class on \(G_{K,p}\) lies in the
Picard subgroup by (CE.8).  Its inflation to the
absolute group \(G_K\) is therefore zero: the
restriction of a line bundle to \(\operatorname{Spec}K\)
is trivial, and the Kummer boundary is natural.
By (CS.10)--(CS.11), its value is three times the
Bockstein of the specified \(f_2=v-ax\).  Since
\(3\) is invertible, that Bockstein also inflates
to zero.  KT-CUBIC-MATCHED-NORM-SYSTEM now provides
a genuinely matched \(y\) over the absolute group,
whose truncated system descends to \(G_{K,p}\).

For this particular \(y\), Sharifi's Proposition~4.2
identifies the value of its specified defining
system exactly with \(\operatorname{Tra}[D^{(2)}y]\).
The proof of Theorem~5.4, equation (5.2), identifies
this exact transgression with \(-[N_{L/K}\mathfrak Y]\).
The possible common prime-to-\(p\) power taken in
that proof can be canceled in the \(\mathbf F_p\)-coefficient
group.  There is no boundary term because
\(2<p-1\).  The lower-product quotient is introduced
in the statement of the theorem to forget which
proper system was chosen; it is unnecessary for
these two identities attached to the same \(y\).
The exact cochain identities (CS.3)--(CS.10) say
that the raw ordinary obstruction is three times
this binomial-system value.  Applying \(\pi_\delta\)
therefore proves (CE.11).  If an arbitrary proper
norm representative is used, Sharifi's Theorem~5.4
instead gives its coset modulo lower products,
and proves (CE.7).
The notation \(\pi_\delta(P)\) denotes the image subgroup;
no stability assertion about \(P\) is needed for this
projection.  A nonzero resulting coset is a nonzero full
class, so (CE.2) proves the final implication.
\end{proof}

\paragraph{What this comparison changes, and what remains.}
The raw formula with \(S_K\) containing primes above \(D\)
could erase an ordinary obstruction supported on those
ideal classes.  Equations (CE.1)--(CE.11) remove that loss:
the specified obstruction is an element of the full
rational \(\omega^{1-e}\) class component, and the
proper norm system can be formed with ramification only
above \(p\).  A matched norm representative gives
the exact class (CE.11), but no such norm representative
or ideal class has been explicitly calculated here.
No nonzero value of that element or of its norm coset
has been established.  In particular this
comparison supplies neither an injection from the
original strict \(\omega^{-e}\) Selmer group nor a
source-to-independent-jet theorem.  The larger finite
coefficient-field norm calculation still requires its
own scalar or cyclic reduction.  An arithmetic exclusion
would require an evaluated nonzero norm coset (or a
matched evaluation before the lower-product quotient)
and the separate source-to-jet hypotheses.

\paragraph{Primary inputs and check scope.}
The etale statements used here are Milne,
\emph{Arithmetic Duality Theorems}, second edition (2006):
I.2.14 and its proof, pp.~36--37; II.0.8 and II.0.10,
pp.~145--148; the complete II.1.1, p.~149; II.2.1 and
Remark II.2.2(a), pp.~163--165; II.2.3--2.4 and the complex-level proof,
pp.~166--168; and the complete comparison proposition
II.2.9, pp.~170--172.  These passages supply the local
gluing maps, Kummer and Brauer computations, and the
comparison with the actual finite-ramification Galois
groups.  Milne's official addendum, revised 26 October 2025,
corrects the order of \(i^*\) and \(i_*\) at printed p.~168,
line~5; the sequence above uses the corrected \(i_*i^*\).
The correction is on p.~1 of
\url{https://www.jmilne.org/math/Books/add/ADT2006.pdf}.
The norm matching additionally uses the absolute cyclic
norm-residue exact sequence of Lam--Liu--Sharifi--Wake--Wang,
Definition~6.1.2 and its following paragraph, and the
coefficient correction in the complete Proposition~6.1.5,
pp.~48--49.  The exact Kummer-dual and norm-system
identifications are Sharifi's complete Lemma~3.1,
pp.~13--14, Lemma~4.1, pp.~19--21, Proposition~4.2,
p.~21, and equation (5.2), p.~27.
The proper norm theorem is Sharifi,
\emph{Massey products and ideal class groups},
Theorem~5.4, pp.~26--27 of the 26 July 2006 author
version, with the binomial matrix conventions of
Lam--Liu--Sharifi--Wake--Wang, Section~4.3.  Their
terminology alone is not used to infer a matched norm
representative: that requires the zero absolute
obstruction and KT-CUBIC-MATCHED-NORM-SYSTEM above.
The foundational theorems are imported.  This new proof
checks their scope and the relative comparison; it does
not reconstruct etale excision or the general norm
theorem, and it does not assert an arithmetic
nonvanishing computation.

% END evidence/unitary_audit/reducible_comparison/cubic_excision_norm.tex


% BEGIN evidence/unitary_audit/reducible_comparison/finite_coefficient_norm.tex
\subsection{The norm comparison over every finite coefficient field}
\label{sec:kt-finite-coefficient-norm}

The scalar restriction in the cyclic norm construction can be
removed by a finite polarization identity.  No extension with
additive Galois group \(k\) is treated as cyclic.

\begin{proposition}[KT-CUBIC-FINITE-COEFFICIENT-NORM]
Retain the hypotheses of KT-ORDINARY-CUBIC-MASSEY-CONDITION.
Over \(\mathbf F_p\), put
\[
\begin{split}
 U&=H^1(G_{\mathbf Q,S},\mathbf F_p(\eta\psi)),\\
 V&=\ker\left(H^1(G_{\mathbf Q,S},\mathbf F_p(\psi^{-1}))
       \longrightarrow H^1(\mathbf Q_p,\mathbf F_p(\psi^{-1}))
           \oplus\bigoplus_{q\mid D}H^1(I_q,\mathbf F_p(\psi^{-1}))
         \right).
\end{split}                                               \tag{FC.1}
\]
These are finite-dimensional spaces.  Choose bases
\((u_i)_i\) of \(U\) and \((v_j)_j\) of \(V\).
For any finite \(k/\mathbf F_p\), write
\(u=\sum_i c_i u_i\), \(v=\sum_j d_jv_j\)
under scalar extension.  The actual relative obstruction
\(q(u,v)=\mathfrak o_3(u,v)\) satisfies
\[
\begin{split}
q(u,v)={}&\sum_{i,j}c_i^2d_jq(u_i,v_j)\\
 &+\sum_{i<i',j}c_ic_{i'}d_j
     \bigl(q(u_i+u_{i'},v_j)-q(u_i,v_j)-q(u_{i'},v_j)\bigr).
\end{split}                                               \tag{FC.2}
\]
The right side belongs to
\(k\otimes_{\mathbf F_p}\mathcal O_{\psi,\mathbf F_p}\),
which is naturally \(\mathcal O_{\psi,k}\).
Every nonzero first input on the right is
\(\mathbf F_p\)-valued and cuts out a cyclic extension
of degree \(p\) after restriction to \(K=E(\mu_p)\).
Consequently KT-CUBIC-P-ONLY-NORM expresses every term
by its matched norm-ideal class.  Formula (FC.2)
therefore extends the exact norm comparison to all
finite coefficient fields.
\end{proposition}

\begin{proof}
The characters in (FC.1) take values in \(\mathbf F_p\).
Continuous cochains with finite \(k\)-coefficients
are obtained from their \(\mathbf F_p\)-versions
by scalar extension: a basis of \(k/\mathbf F_p\)
identifies them with a finite direct sum.
Cohomology, all local maps, their kernels, and the
relative pair quotient (CM.14) consequently commute
with that scalar extension.

Choose \(\mathbf F_p\)-linear cocycle representatives
of the two finite spaces in (FC.1).
For every pair \((u_i,v_j)\), choose an ordinary
quadratic defining cochain \(w_{ij}\) as in
KT-GLOBAL-QUADRATIC-MIXED-PERSISTENCE.  Define
\[
                 w(u,v)=\sum_{i,j}c_i d_j w_{ij}.
                                                        \tag{FC.3}
\]
Its differential has the required value because
the quadratic defining equation is bilinear.
The ordinary condition is bilinear as well.
Indeed the local primitive of the strict cocycle
\(v\) at \(p\) can be chosen linearly: it is unique
because \(H^0(\mathbf Q_p,\mathbf F_p(\psi^{-1}))=0\).
The correction \(d(v)(1+\beta)u/2\) in (CM.8)
is then bilinear, and (FC.3) retains the required
unramified class.  Away from \(p\), each chosen
quadratic family is literally trivial on inertia
over \(E\).  Its logarithmic coordinates vanish
there, so (FC.3) retains this condition too.
The uniqueness modulo global coboundary proved
before (CM.8) shows that this choice computes
the same intrinsic ordinary obstruction.

With (FC.3), the cocycle \(-\Omega_+\) in (CM.10)
is homogeneous of degree two in \(u\) and degree
one in \(v\): its terms are \(uw\) and \(u^2v\).
We must retain the same homogeneity in the local
primitives, rather than merely in its raw class.
At \(p\), every primitive is equivalent modulo
the allowed full local \(H^1\), and local \(H^2\)
vanishes.  Choose a linear primitive on the
finite span of the two-cocycles at issue.
At \(q\mid D\), the quadratic family factors
through \(G_{\mathbf Q_q}/I_{E,q}\).
The defect therefore factors through that group.
It has \(p\)-cohomological dimension one, so on
the same finite span choose a linear primitive
within its continuous cochain complex.
Such a primitive is precisely an acceptable
unramified-over-\(E\) cubic local lift in the
rational upper direction.  Its ambiguity is
the allowed unramified local \(H^1\).
Thus the relative pair, not just its raw
two-cocycle, can be chosen homogeneous of
bidegree \((2,1)\) in \((u,v)\).

Passing to the linear relative quotient proves
that \(q\) is quadratic in the first input and
linear in the second, over \(\mathbf F_p\),
with the same operation after scalar extension.
Expanding a quadratic map in a basis gives
exactly (FC.2); the off-diagonal coefficient
is its polarization
\(q(u_i+u_{i'},v_j)-q(u_i,v_j)-q(u_{i'},v_j)\).
No division by a field-dependent scalar is used.

Basis vectors \(u_i\) are nonzero, and
\(u_i+u_{i'}\ne0\) for distinct basis vectors.
Prime-to-\(p\) restriction to \(K\) preserves
their nonzero cohomology classes.  Their
\(\mathbf F_p\)-valued images are therefore
cyclic of order \(p\), exactly the scope of
the matched norm theorem.  If \(U\) or \(V\)
is zero, the formula is the empty sum and
the obstruction is zero.  Applying (CE.2)
and (CE.11) to the finitely many terms proves
the claimed norm comparison for arbitrary
finite \(k\).
\end{proof}

This proposition supplies the scalar reduction
needed for the finite-coefficient norm formula.
It evaluates none of its ideal classes, does
not prove the operation injective in \(v\),
and does not supply an actual source-to-jet map.

% END evidence/unitary_audit/reducible_comparison/finite_coefficient_norm.tex


% BEGIN evidence/unitary_audit/reducible_comparison/quartic_quintic_obstruction.tex
\subsection{The next two ordinary obstruction degrees}
\label{sec:kt-quartic-quintic-obstruction}

Retain (CM.1)--(CM.4), and now assume all three auxiliary units
at exponents \(0,e,2e\).  The hypotheses imply \(p\ge7\):
\(p=5\) and even \(e\) would give \((p-1)\mid2e\).
Thus exponential and logarithmic coordinates through degree five
are defined.  Write \(A_j=k[s,t]/(s,t)^{j+1}\),
\(F=\mathbf Q(\mu_p)\), \(K=EF\), and
\(A_J=\operatorname{Cl}(J)\{p\}\).
All assertions concern the specified equivariant ordinary Artin
families.  They do not assert that a characteristic-zero source
point supplies such a family.

Set
\[
 P=E_{13},\qquad H=\operatorname{diag}(1,0,-1),\qquad Q=E_{31}.
                                                        \tag{HQ.1}
\]
The permitted homogeneous terms after (CM.5) are exactly
\[
 \begin{array}{c|c|c|c}
 \text{degree}&\text{monomial}&\text{Lie vector}&\text{character}\\\hline
 4&s^3t&P&\eta\psi^2\\
 4&s^2t^2&H&1\\
 4&st^3&Q&\eta\psi^{-2}\\
 5&s^3t^2&X&\eta\psi\\
 5&s^2t^3&Y&\psi^{-1}
 \end{array}                                             \tag{HQ.2}
\]
Indeed the integer torus weights of \(\mathfrak{sl}_3\)
are \(-2,-1,0,1,2\).  The involution signs of
\(P,H,Q,X,Y\) are respectively \(-,+,-,+,-\).
These weights and signs give (HQ.2); there are no additional
permitted monomials.  In particular the positive long root
at degree four cannot be omitted.

\begin{theorem}[KT-ORDINARY-QUARTIC-CONDITION]
Fix an ordinary equivariant lift of the quadratic family to
\(A_3\).  It has a well-defined remaining quartic obstruction
\[
 \mathfrak o_4^+\in H^2(G_{\mathbf Q,S},k(\eta\psi^2)).
                                                        \tag{HQ.3}
\]
The cubic family lifts to \(A_4\), with the same ordinary
and away-prime conditions, if and only if
\(\mathfrak o_4^+=0\).  The other two quartic directions
have zero relative obstruction spaces.

The target of (HQ.3) has the exact description
\[
 H^2(G_{\mathbf Q,S},k(\eta\psi^2))
 \simeq
 \operatorname{Hom}_{\operatorname{Gal}(K/\mathbf Q)}
       (A_K/pA_K,k(\eta\omega^{1-2e}))^\vee.
                                                        \tag{HQ.4}
\]
Here the dual is over \(k\), using the trace pairing.
The extra unit
\[
                       B_{1,\eta\omega^{-2e}}\in\mathbf Z_p^\times
                                                        \tag{HQ.5}
\]
is sufficient to make this target zero.  It is not one of
the three assumed units in general.  This does not claim
that (HQ.5) is necessary for the specific class (HQ.3)
to vanish.
\end{theorem}

\begin{proof}
Use \(A,B,C,a,b,c\) as in (CM.9), and additionally put
\[
 M=m(g),\quad N=n(g),\quad
 m=\psi(g)m(h),\quad n=\eta(g)\psi(g)^{-1}n(h),
 \qquad f=Ab-Ba.
                                                        \tag{HQ.6}
\]
The repeated letters \(m,n\) on the right denote the
transported values at \(h\), only within this calculation.
The degree-four BCH defects, in the order \(P,H,Q\), are
\[
\begin{split}
 \Omega_{4,+}
   & =aM-Am+\tfrac12(A-a)(Ac-aC)-\tfrac14Aa f,\\
 \Omega_{4,0}
   & =\tfrac12(An-Bm+bM-aN)+\tfrac12(ABc+abC)\\
   &\hspace{12mm}-\tfrac14(C+c)(Ab+Ba)-\tfrac18(Ab+Ba)f,\\
 \Omega_{4,-}
   & =bN-Bn-\tfrac12(B-b)(Bc-bC)+\tfrac14Bb f.
\end{split}                                               \tag{HQ.7}
\]
If the new one-cochains are \(r_+,r_0,r_-\), their
equations in the convention of (CM.14) are
\(dr_i=-\Omega_{4,i}\).
One obtains (HQ.7) from the terms of total degree four in
\[
 A+B+\tfrac12[A,B]+\tfrac1{12}[A,[A,B]]
 +\tfrac1{12}[B,[B,A]]-\tfrac1{24}[B,[A,[A,B]]],
                                                        \tag{HQ.8}
\]
where in (HQ.8) the letters \(A,B\) denote the two
entire logarithmic matrices.  Alternatively the finite
matrix exponential, product, and logarithm give the same
three expressions directly.  The independent program
\path{quartic_bch_check.py} uses this latter computation
and checks all five degree-four monomials.
Associativity makes these defects two-cocycles: the
degree-four kernel is an additive central layer over
the fixed residual action.

We justify the actual local spaces before using global
cohomology.  At \(p\), the checked formally smooth
ordinary functor in KT-ORDINARY-INTEGRAL-COMPONENT gives
local lifts.  Its proof applies to these small extensions;
the even residual exponent differences never equal the
cyclotomic exponent one.  Fixed determinant is obtained
by removing the trace, since \(3\) is a unit.
The permitted local differences are all \(H^1\) in
an upper root, the unramified \(H^1\) in the diagonal
direction, and zero \(H^1\) in a lower root.  The last
assertion retains flag conjugation: a lower-root
difference can be produced by changing the ordinary
flag precisely when it is a local coboundary.

At \(q\mid D\), the quotient obtained by killing
inertia over \(E\) has an inertia subgroup of order
at most two and a procyclic unramified quotient.
It has \(p\)-cohomological dimension one.  Hence
acceptable local lifts exist across the square-zero
layer.  Its tangent \(H^1\) is zero for an
\(\eta\)-ramified coefficient and is ordinary
unramified \(H^1\) for a rational coefficient.
Projection to the required integer weight and
involution sign preserves these linear lifting
equations.  This gives equivariant local lifts as
well; it is not an inference from raw global
\(H^2\) alone.

For \(k(\eta\psi^{-2})\), the unit at \(2e\)
annihilates the even mirror
\(A_K^{\eta\omega^{2e+1}}\) and its strict
negative group by KT-AUXILIARY-UNIT-SELMER-BRIDGE.
Exactly the proof of (CC.1), using only
\(2e\not\equiv0\pmod{p-1}\), gives
\[
 H^2(G_{\mathbf Q,S},k(\eta\psi^{-2}))=0,
 \quad H^1(G_{\mathbf Q,S},k(\eta\psi^{-2}))
      \xrightarrow{\sim}H^1(\mathbf Q_p,k(\eta\psi^{-2})).
                                                        \tag{HQ.9}
\]
At every \(q\mid D\) its local \(H^1\) is zero.
Consequently a global solution can be corrected to
match the required lower-root local coboundary at
\(p\); the negative relative obstruction space is zero.

For the trivial diagonal coefficient, unramified
conditions away from \(p\) reduce the relative problem
to \(G_{\mathbf Q,p}\) by the same excision as (CE.3).
Ray class theory over \(\mathbf Q\), or Kronecker--Weber,
gives
\[
 \dim_k H^1(G_{\mathbf Q,p},k)=1,\qquad
 H^1(G_{\mathbf Q,p},k)
    \xrightarrow{\sim}
 H^1(\mathbf Q_p,k)/H^1_{\rm unr}(\mathbf Q_p,k).
                                                        \tag{HQ.10}
\]
The character of the degree-\(p\) subfield of
\(\mathbf Q(\mu_{p^2})\) generates the first group
and is ramified at \(p\).  The target also has
dimension one by local Euler characteristic.
Global Euler characteristic for the trivial even
module gives \(1-1+h^2=0\), so
\(H^2(G_{\mathbf Q,p},k)=0\).
The relative exact sequence and (HQ.10) show that
the diagonal obstruction space is zero.

For the positive coefficient \(M=k(\eta\psi^2)\),
its Cartier dual is \(k(\eta\omega^{1-2e})\).
Every local \(H^2\) in \(S\) is zero: the dual
has nontrivial inertia at \(p\), and quadratic
ramification at \(q\mid D\).  The local \(H^1\)
at such \(q\) is also zero.  Since the positive
local space at \(p\) is all \(H^1\), the relative
obstruction is exactly raw global \(H^2\).
Set \(\mathfrak o_4^+=-[\Omega_{4,+}]\).
The independent central directions and the preceding
two zero relative groups prove the lifting criterion.

Poitou--Tate duality identifies this global \(H^2\)
with the dual of strict \(H^1\) in the displayed
Cartier-dual character.  Restriction to \(K\),
prime-to-\(p\) averaging, and class field theory
identify that strict group with the Hom space in
(HQ.4).  There is no prime-class quotient: at a
prime over \(p\) the cyclotomic inertia exponent
\(1-2e\) is nontrivial, and at a prime over \(D\)
the quadratic inertia character is nontrivial.
The permutation modules of these primes therefore
have no component in that character.  These same
facts show that every class-group homomorphism in
(HQ.4) satisfies the required strict local conditions.
Finally apply the auxiliary-unit bridge with the
nonzero even residue \(a\equiv-2e\pmod{p-1}\).
Its even mirror is \(\eta\omega^{1-2e}\), proving
the sufficient condition (HQ.5).
\end{proof}

\begin{theorem}[KT-QUINTIC-RELATIVE-TARGET]
Suppose the fixed cubic family admits a quartic ordinary lift.
For any such lift the two possible quintic obstructions have
coefficients \(k(\eta\psi)\) and \(k(\psi^{-1})\).
The upper relative obstruction space is
\[
 H^2(G_{\mathbf Q,S},k(\eta\psi))
 \simeq\operatorname{Hom}_{\operatorname{Gal}(K/\mathbf Q)}
       (A_K/pA_K,k(\eta\omega^{1-e}))^\vee.
                                                        \tag{HQ.11}
\]
The additional unit \(B_{1,\eta\omega^{-e}}\) is
sufficient to kill (HQ.11).

Let \(\mathcal O^-_5\) denote the rational lower
relative obstruction space.  Its local condition at
\(p\) is zero, and its away-prime conditions are
unramified.  Put
\[
 d=\dim_k\operatorname{Hom}_{\operatorname{Gal}(F/\mathbf Q)}
                       (A_F/pA_F,k(\omega^{-e})).
                                                        \tag{HQ.12}
\]
There is an exact sequence
\[
\begin{split}
0\longrightarrow&
 \operatorname{coker}\bigl(H^1(G_{\mathbf Q,p},k(\psi^{-1}))
                  \to H^1(\mathbf Q_p,k(\psi^{-1}))\bigr)\\
\longrightarrow&\mathcal O^-_5
\longrightarrow H^2(G_{\mathbf Q,p},k(\psi^{-1}))
\longrightarrow0,
\end{split}                                               \tag{HQ.13}
\]
and precisely
\[
                         \dim_k\mathcal O^-_5=d+1.
                                                        \tag{HQ.14}
\]
The dimension \(d\) is that of the original residual
strict rational negative Selmer group.  The raw last
term in (HQ.13) has the odd mirror character
\(\omega^{e+1}\); it cannot replace the full relative
space.

For the fixed cubic family, its rational lower
quintic obstruction is independent of the choice
of acceptable quartic lift.  If the upper space
(HQ.11) is zero, the quartic family lifts through
degree five if and only if that rational relative
class is zero.  No nonzero value of this class is
asserted.
\end{theorem}

\begin{proof}
Table (HQ.2) lists the entire kernel of
\(A_5\to A_4\) allowed by the grading and involution.
It is an additive central layer, so the homomorphism
defect of any set-theoretic lift gives the two
two-cocycles, with the same relative construction
as (CM.14).  The local lifting and flag argument
in the preceding proof applies without change.
For the upper character the local space at \(p\)
is full and the spaces at \(q\mid D\) vanish.
The argument for (HQ.4), with \(2e\) replaced by
\(e\), proves (HQ.11) and its sufficient opposite
unit condition.

For the lower rational character, a permitted
local difference at \(p\) is a coboundary: it
comes from changing the ordered ordinary flag.
Thus \(L_p=0\), not full local \(H^1\).
At \(q\mid D\), killing inertia over \(E\)
gives exactly unramified \(H^1\) by the
order-at-most-two averaging already proved.
Apply etale excision at those primes while
retaining the \(p\)-primitive.  The result is
the mapping-fibre degree-two group for
\(R\Gamma(G_{\mathbf Q,p},k(\psi^{-1}))\to
R\Gamma(\mathbf Q_p,k(\psi^{-1}))\).
Local \(H^2\) is zero, because its dual inertia
character \(\omega^{e+1}\) is nontrivial.
The mapping-fibre sequence is exactly (HQ.13).

Write \(r=\dim_k H^2(G_{\mathbf Q,p},k(\psi^{-1}))\).
The Kummer and class-group argument in (CE.4)--(CE.5)
identifies this group with the scalar extension of
\((A_F/pA_F)^{\omega^{e+1}}\).
Since \(\psi^{-1}\) is even and nontrivial,
global Euler characteristic gives
\(\dim H^1(G_{\mathbf Q,p},k(\psi^{-1}))=r\).
Local \(H^1\) has dimension one: both its
local \(H^0\) and \(H^2\) vanish.

The kernel of that localization has dimension
\(d\).  To see the exact identification, restrict
to \(F\).  A strict class becomes an everywhere
unramified homomorphism.  Conversely an equivariant
homomorphism from \(A_F/pA_F\) in character
\(\omega^{-e}\) gives a global class unramified
outside \(p\); its restriction at the prime above
\(p\) is zero because that prime is principal.
The local and global restriction maps are
injective by averaging the degree \(p-1\).
This proves the claimed equality with the
original residual strict group.  Hence the
cokernel in (HQ.13) has dimension \(1-(r-d)\),
and adding \(r\) proves (HQ.14).

Finally compare two ordinary quartic lifts of
the same cubic family.  Their difference is
a one-cocycle in the three lines of degree four.
In the negative line its global strict class
is zero by the unit at \(2e\); in the diagonal
line its everywhere-unramified class is zero
by the class group of \(\mathbf Q\).
Modulo permitted global conjugations, only
the positive \(s^3tP\) difference remains.
Its brackets with the degree-one terms can
affect only \(s^3t^2X\), since
\([P,X]=0\) and \([Y,P]=X\).
They cannot affect \(s^2t^3Y\).
Thus the rational lower relative obstruction
does not change.  The two central degree-five
directions are independent at this order;
if (HQ.11) vanishes, solving the remaining
relative equation is therefore necessary
and sufficient for the degree-five lift.
\end{proof}

\paragraph{The exact remaining requirements.}
The three given auxiliary units remove the negative
quartic direction and the diagonal direction, but
do not in general remove the positive quartic
target (HQ.4).  An opposite auxiliary unit suffices
there; a second opposite unit suffices for the
upper quintic target.  These are sufficient group
vanishing inputs, not necessary conditions on
individual obstruction values.  The rational
quintic condition then sees a relative space of
dimension one plus the original target dimension.
No injectivity, nonzero value, or general lifting
contradiction follows from this dimension count.
The actual source-to-jet conditions remain
separate.  This section supplies exact later
lifting targets and a quartic cochain formula;
it does not settle an additional arithmetic gate.

\paragraph{Inputs and verification scope.}
The duality, Euler, ordinary-component, and auxiliary-unit
inputs are the same checked statements used in
(CM.12)--(CM.15), (CC.1)--(CC.3), and (AB.3)--(AB.5).
The diagonal calculation additionally uses the ray class
description over \(\mathbf Q\): Milne,
\emph{Class Field Theory}, Theorem~0.5 and the full
conductor paragraph, printed pp.~5--6, including
\(C_{m\infty}\simeq(\mathbf Z/m\mathbf Z)^\times\).
That exact statement and its surrounding conductor
scope were read; global class field theory is imported.
The independent rational matrix-exponential program
checks (HQ.7) and the six displayed bracket identities
used to identify the permitted directions.
No arithmetic norm or Massey value is computed.

% END evidence/unitary_audit/reducible_comparison/quartic_quintic_obstruction.tex


% BEGIN evidence/unitary_audit/reducible_comparison/source_to_cubic_jet.tex
\subsection{When a source point forces the cubic jet}
\label{sec:kt-source-to-cubic-jet}

A characteristic-zero point with nonzero opposite extension
entries does not, by itself, force the independent two-parameter
cubic jet of KT-ORDINARY-CUBIC-MASSEY-CONDITION.
The obstruction is present even in a flat domain with the exact
torus weights and involution signs.  A sufficient extraction
criterion is given afterwards, with its extra valuation and
coefficient-support hypotheses stated explicitly.

\begin{proposition}[KT-EQUIVARIANT-POINT-NOT-CUBIC-JET]
For every prime $p>3$, the complete local algebra
\[
 R=\mathbf Z_p[[x,y,z]]/(pz-x^2y)                         \tag{SJ.1}
\]
is a $\mathbf Z_p$-flat domain.  It has a continuous formal
$\mathbf G_m$ action induced by weights $(1,-1,1)$ on $(x,y,z)$,
algebraic on each quotient $R/\mathfrak m^n$,
an involution with signs $(+,-,-)$, and an augmentation to
$\mathbf Z_p$.  There is a continuous local characteristic-zero
specialization with $x,y$ both nonzero.

There is an equivariant quadratic jet
\[
 R\longrightarrow \mathbf F_p[s,t]/(s,t)^3,\qquad
             x\mapsto s,\quad y\mapsto t,\quad z\mapsto0.
                                                               \tag{SJ.2}
\]
There is no lift to $\mathbf F_p[s,t]/(s,t)^4$ with these
first-order values of $x,y$, even without requiring
equivariance of the lift.
\end{proposition}
\begin{proof}
The actions preserve the defining ideal: both terms of its
generator have weight one and involution sign minus.
The augmentation sets all three variables to zero.
The characteristic-zero point is
\[
                         x=p,\qquad y=p,\qquad z=p^2.
                                                               \tag{SJ.3}
\]
All three images are in the maximal ideal and satisfy (SJ.1),
so substitution defines a continuous local map to $\mathbf Z_p$.

For the domain assertion, filter the regular ambient ring
with positive weights
$\mathrm{wt}(p)=\mathrm{wt}(x)=\mathrm{wt}(y)=1$ and
$\mathrm{wt}(z)=2$.
This filtration is separated and cofinal with the
maximal-ideal filtration.  Its associated graded ring is
$\mathbf F_p[P,X,Y,Z]$.
Multiplicativity of initial forms gives
\[
 \operatorname{gr}R
      =\mathbf F_p[P,X,Y,Z]/(PZ-X^2Y).                  \tag{SJ.4}
\]
The polynomial on the right is primitive and degree one
in $Z$, hence irreducible by Gauss's lemma, since
$\gcd(P,X^2Y)=1$.  The polynomial ring is a unique
factorization domain, so this quotient is a domain.
The separated filtration then makes $R$ a domain.
The element $p$ survives, since its initial form $P$
is nonzero in (SJ.4).  Therefore $R$ is torsion-free,
and hence flat, over the discrete valuation ring
$\mathbf Z_p$.

Equation (SJ.2) is valid because its only relation has
degree three after reduction modulo $p$.
For a putative cubic lift, write
$x=s+O((s,t)^2)$ and $y=t+O((s,t)^2)$.
Its defining relation reduces to
\[
                      0=x^2y=s^2t\pmod{(s,t)^4},
                                                               \tag{SJ.5}
\]
which is false.  This proves the assertion.
\end{proof}

This example has no asserted Galois or automorphic realization.
Its role is exact: flatness, a characteristic-zero point,
the torus grading, the involution, and nonzero opposite
entries do not imply the desired independent cubic jet.
The term $pz$ can cancel $x^2y$ at the characteristic-zero
point and disappears in the characteristic-$p$ lifting problem.

\begin{proposition}[KT-SPARSE-VALUATIVE-CUBIC-EXTRACTION]
Retain the field, exponent and auxiliary-unit hypotheses
of KT-ORDINARY-CUBIC-MASSEY-CONDITION.
Let $\mathcal O$ be a finite coefficient DVR, with
uniformizer $\pi$, residue field $k$, and
$r_\pi=v_\pi(p)$.
Let $r:G_{\mathbf Q}\to
\mathrm{GL}_3(\mathcal O)\rtimes\{\pm1\}$ be continuous,
with right component $\eta$, determinant-one left component,
and residual representation $r_0$.
Assume its restriction to $G_E$ is unramified away from
$p$ and that it has the full ordered ordinary flag at $p$
with exact inertia characters $(\chi^e,1,\chi^{-e})$.
Here the flag is in the semidirect Borel sense: after
one local conjugation the left matrices for the full
rational decomposition group are upper triangular.

Suppose an integer $d\ge1$ satisfies
\[
                              r_\pi>3d,                       \tag{SJ.6}
\]
Assume $r(g)r_0(g)^{-1}\equiv1\pmod{\pi^d}$ for all $g$,
and, in the displayed global basis, the truncated logarithm of
$r(g)r_0(g)^{-1}$ modulo $\pi^{3d+1}$ has the sparse form
\[
\begin{split}
 L(g)={}&\pi^d\{u(g)X+v(g)Y\}
             +\pi^{2d}w(g)Z+\pi^{3d}W(g),\\
 &\hspace{30mm}W(g)\in\mathfrak{sl}_3(k).
\end{split}                                                   \tag{SJ.7}
\]
Here $r_0$ is lifted by the finite Teichm\"uller characters
before reduction, and (SJ.7) means there are no intermediate
coefficient powers below $\pi^{3d}$.
Then the cochains $u,v$ satisfy (CM.2), including the strict
and away-prime conditions on $v$, and
\[
                             \mathfrak o_3(u,v)=0.             \tag{SJ.8}
\]
In particular a point satisfying these explicit hypotheses
forces the graded, involutive two-parameter cubic jet, even
though it initially supplies only a one-parameter specialization.
\end{proposition}
\begin{proof}
The ring $\mathcal O/\pi^{3d+1}$ has characteristic $p$
by (SJ.6).  It contains a coefficient field $k$: the unique
lifts of residue elements satisfying $a^{|k|}=a$ form that
field, since the derivative is a unit and Frobenius is
additive in characteristic $p$.
Successive $\pi$-adic expansion therefore identifies it
with $k[\pi]/\pi^{3d+1}$.
The subalgebra generated by $\lambda=\pi^d$ is
$k[\lambda]/\lambda^4$.
The finite residual Teichm\"uller representation lies in
the constant coefficient field.
The truncated logarithm is the finite polynomial
$N-N^2/2+N^3/3$ for $N=r(g)r_0(g)^{-1}-1$.
The hypothesis modulo $\pi^d$ gives $N^4=0$ in this
quotient, so logarithms and exponentials terminate at
degree three.  No characteristic-zero analytic
logarithm or its convergence range is used.

Project the homomorphism identity first to the coefficient
of $\lambda$.  It gives the cocycle laws for $u$ and $v$.
Projection to $\lambda^2Z$ gives (CM.8).
At degree three put
\[
 m(g)=\tfrac12\{W(g)_{12}-W(g)_{23}\},\qquad
 n(g)=\tfrac12\{W(g)_{21}+W(g)_{32}\}.                  \tag{SJ.9}
\]
The $U,V$ lines are separately preserved by the residual
adjoint action.  The BCH calculation therefore gives
exactly (CM.11) for these $m,n$.
The other components of $W$ enter only linearly at this
order and can be discarded: none contributes to either
projected equation.  Substituting the independent monomials
$s,t,st,s^2t,st^2$ for the indicated five coefficient slots
now gives the required global homomorphism over $A_3$.
Its determinant, grading and involution follow from
the five trace-zero Lie vectors and (CM.3)--(CM.5).

It remains to verify the local conditions after this
separation of the parameters.
At every away-prime inertia subgroup $I_q\cap G_E$,
both $r$ and $r_0$ are the identity.  Every coefficient
in (SJ.7), including the projections (SJ.9), is therefore
zero there.  The extracted family is unramified over $E$.
For $v$, the larger rational inertia group has quotient
of order at most two after intersecting $G_E$, and its
coefficient character $\beta$ is unramified.
A homomorphism from that quotient into the additive
$p$-group $k$ is zero.  Hence $v$ is unramified on
the full rational inertia away from $p$ as required.

At $p$, choose an inertia element $\tau$ for which
$\psi(\tau),1,\psi(\tau)^{-1}$ are pairwise distinct.
This is possible because $(p-1)\nmid2e$.
The exact ordinary inertia characters reduce to these
constant characters in our characteristic-$p$ quotient.
The ordered eigenflag of $r(\tau)$ is therefore uniquely
defined in the neighborhood of the standard flag.
Its eigenprojectors are polynomials in $r(\tau)$ with
unit constant denominators.  Because (SJ.7) is defined
over $k[\lambda]/\lambda^4$, its normalized lower-unipotent
flag matrix is defined over the same subalgebra.
It is the flag furnished by ordinarity, so is stable
under the full decomposition group.

At first order the only negative Lie vector in (SJ.7)
is $Y$.  The flag equation consequently gives
$v(g)=\beta(g)a-a$ for one constant $a\in k$ on the full
decomposition group.  Conjugating by $\exp(\lambda aY)$
leaves
\[
 \lambda u(g)X+\lambda^2
       \{w(g)-\tfrac12a(1+\beta(g))u(g)\}Z
                       +\lambda^3W'(g).                       \tag{SJ.10}
\]
There is no remaining negative term through degree two.
Uniqueness of the normalized lower flag means that its
remaining correction starts at degree three.
The exact prescribed inertial diagonal implies that
the coefficient of $Z$ in (SJ.10) is zero on inertia;
thus $w$ has precisely the ordinary normalization used
in (CM.8).

The $V$ projection of $W'$ is a coboundary in
$k(\eta\psi^{-1})$: it is canceled by the degree-three
$V$ component of that lower flag matrix.
This projection is exactly the $st^2V$ coefficient of
the separated two-parameter family after its conjugation
by $\exp(t aY)$, since other components of $W$ cannot
interact before degree four.
A further degree-three $V$ conjugation makes the
extracted family upper triangular.
Its degree-three $U$ term is already upper triangular,
and its only diagonal correction is the $stZ$ term,
which is zero on inertia.  Thus its ordinary local
condition is exact.  This proves (SJ.8) by the cubic
criterion.
\end{proof}

\paragraph{What remains unproved for an arbitrary source.}
The actual source supplies a finite coefficient field,
the residual representation, and the local flags.
It does not currently supply (SJ.6), the sparse support
(SJ.7), or a theorem allowing other initial directions
to be removed while preserving the two specified
nonzero extension classes.
Passing to a ramified coefficient extension alone does
not establish (SJ.6): both $v_\pi(p)$ and the existing
entry valuations scale by the ramification degree.
The algebra (SJ.1) shows why these missing conditions
cannot be replaced just by flatness or by the existence
of a characteristic-zero point with nonzero edges.

Accordingly (SJ.8) is a proved conditional extraction
criterion, not an assertion that every remaining
primitive, nonessentially-self-dual TypeI point has
the independent cubic jet.

% END evidence/unitary_audit/reducible_comparison/source_to_cubic_jet.tex


% BEGIN evidence/lean_chain/source_point_cubic_scope.tex
\paragraph{KT-SOURCE-CUBIC-EXTRACTION-SCOPE.}
The exact archived source construction does not yet supply
the extra hypotheses (SJ.6)--(SJ.7).
In \path{GRepresents.lean:140--141}, \path{IsLocPt} requires
only that maximal-ideal elements have norm strictly less than
one.  The actual \path{IsPt} predicate at
\path{RhoKApply.lean:76--83} adds integral entries, residual
congruence, the marking, the away-prime condition, and the
local ordinary-component condition.  Its transport theorem
\path{IsPtOf.lean:124--166} preserves exactly those clauses.
\path{TwoPt}, at \path{ConjStable.lean:94--95}, adds failure
of the upper and lower alternatives; it does not specify a
first-order logarithmic support.

The complete normalization proof
\path{NormEx.lean:155--217} gives a congruence-one conjugation
fixing the complex-conjugation matrix and two marked entries.
The finite-coefficient argument
\path{TanDet.lean:110--138,182--216} is qualitative.
The torus point-extension theorem
\path{ExistsTorusTwistEqLocalPointOfWeightedPres.lean:95--137}
and the covering construction \path{ModelUpper.lean:103--250}
also conclude local points, without an entry-depth bound.

There are genuine first-order extension statements:
\path{ConjStable.lean:79--80} extends a specified point
derivation to characteristic-zero dual numbers, and
lines81--85 lift a compatible single-weight base derivation
to the ambient graded model.  Their complete implementations
in \path{ExistsDualNumberPointOfPointDerivationVanishing.lean}
and \path{ExistsBasePointDerivationOfWeightHomogeneous.lean}
do not assert its higher-order integrability.
The local power-series chart at \path{ConjStable.lean:25--37}
is over the fraction field and permits a nonzero clearing
factor; it is not an integral sparse global chart.

For a nonzero first-order coefficient at depth $d$, (SJ.6)
requires the normalized entry valuation
$d/v_\pi(p)<1/3$, whereas locality only gives positivity.
A ramified coefficient extension multiplies both numerator
and denominator by its ramification index.  Torus conjugation
shifts opposite-weight valuations by opposite amounts and
preserves their sum.  Neither operation supplies the missing
estimate or eliminates all intermediate coefficients.
This audit leaves open whether a particular arithmetic point
can satisfy the conditional extraction criterion; it supplies
no such point and no universal point-to-cubic-jet implication.

% END evidence/lean_chain/source_point_cubic_scope.tex


% BEGIN evidence/unitary_audit/reducible_comparison/actual_source_pairing_audit.tex
\subsection{The exact scope of a proposed nonzero source pairing}
\label{sec:kt-actual-source-pairing-audit}

This is an audit of a possible next arithmetic argument, not an
additional source exclusion.  Retain the hypotheses and notation
of (CM.1)--(CM.4) and the three auxiliary units.  Put
\(\alpha=\eta\psi\), \(\beta=\psi^{-1}\), and let
\(\mathfrak g\) be the residual trace-zero adjoint module of
the polarized rank-three datum.  The surviving first-order
classes have the form
\(u\in H^1(G_{\mathbf Q,S},k(\alpha))\) and a strict
\(v\in H^1(G_{\mathbf Q,S},k(\beta))\).

\begin{lemma}[KT-SOURCE-TATE-PAIRING-SCOPE]
For every positive integer \(m\),
\[
 \operatorname{Hom}_{G_{\mathbf Q}}
       (\mathfrak g^{\otimes m},k(\omega))=0.
 \tag{TP.1}
\]
In particular, the polarization or an adjoint trace contraction
does not itself provide a Tate-valued bilinear pairing of the
two displayed extension classes.  Their ordinary cup product
has coefficient \(k(\eta)\) and is zero in global cohomology.
\end{lemma}
\begin{proof}
Because \(p\nmid D\), the quadratic character is unramified
at \(p\).  On inertia there, the characters in
\(\mathfrak g\) are powers of \(\omega\) with exponents
\(0,\pm e,\pm2e\).  Choose an inertia element on which
\(\omega\) generates \(\mathbf F_p^\times\).
All these exponents are even.  Every tensor power therefore
has only even exponents, while the target exponent in (TP.1)
is one modulo the even integer \(p-1\).  The inertia
eigenspaces have disjoint characters, proving (TP.1).
This proof uses the actual local residual characters and
does not change upon imposing characteristic-zero
irreducibility on a source point.

Moreover \(\alpha\beta=\eta\), and the checked
quadratic mixed-cup theorem gives
\(H^2(G_{\mathbf Q,S},k(\eta))=0\).
Thus \(u\smile v=0\).  The Cartier dual of
\(k(\beta)\) is \(k(\omega\psi)\), rather than
\(k(\eta\psi)\).  No identification of these
characters is possible on inertia at \(p\).
\end{proof}

The lemma concerns equivariant contractions of the stated
coefficient modules.  It does not rule out a higher
Bockstein or Massey construction with an additional
arithmetic input.

\paragraph{The positive module that actually occurs in duality.}
Let \(U=\operatorname{Spec}\mathbf Z[1/p]\).
The lower relative quintic group of (HQ.13) is
\(H_c^2(U,k(\beta))\), with the usual modified real-place
term, which has no odd-primary contribution.  Artin--Verdier
duality gives a perfect pairing with
\[
 H^1(G_{\mathbf Q,\{p\}},k(\omega\psi)).
 \tag{TP.2}
\]
This assertion uses Milne, \emph{Arithmetic Duality Theorems},
second edition, Theorem II.3.1(b) and Corollary II.3.3,
printed pp.~176--178, together with the compact-support
definition already used in the excision proof.  Those
statements and the full short proof of the corollary were
read for this audit; the global duality theorem is imported.
The trace pairing of \(k/\mathbf F_p\) turns the resulting
finite-group duality into the stated \(k\)-duality.

Writing \(F=\mathbf Q(\mu_p)\), \(\Delta=\operatorname{Gal}(F/\mathbf Q)\),
and \(\mathcal E_F=\mathcal O_F[1/p]^\times\), Kummer
theory gives the exact sequence
\[
\begin{split}
0\longrightarrow&
 (k\otimes_{\mathbf F_p}\mathcal E_F/\mathcal E_F^p)^{(\omega^{-e})}
 \longrightarrow H^1(G_{\mathbf Q,\{p\}},k(\omega\psi))\\
 \longrightarrow&
 (k\otimes_{\mathbf F_p}A_F[p])^{(\omega^{-e})}
 \longrightarrow0.
\end{split}
 \tag{TP.3}
\]
Indeed, restriction to \(F\) and exact \(\Delta\)-invariants
are justified by \(p\nmid p-1\).  Identifying the Kummer
coefficient with \(\mu_p\) twists the character to
\(\omega^{-e}\).  The prime above \(p\) is principal,
so inverting it does not change the class group.
Its contribution to the group of \(p\)-units is a trivial
character and does not affect this nontrivial component.
The Dirichlet logarithmic unit representation contains
each nontrivial even character once.  Its integral
\(\omega^{-e}\)-part is free of rank one; the possible
\(p\)-power root-of-unity torsion has odd character
\(\omega\) and contributes nothing here.
Thus the left term in (TP.3) is a line, and the right
term has the original strict dimension \(d\).
This explains the previously proved dimension \(1+d\),
but supplies no value for the quintic obstruction.

Perfect duality says that a \emph{nonzero} relative
obstruction is detected by some element of (TP.2).
It does not say that irreducibility of a characteristic-zero
point makes a particular obstruction nonzero, nor does it
identify the auxiliary class \(u\) with an element of
the unit line in (TP.3).  The missing statement is an
arithmetic evaluation or nondegeneracy theorem for the
actual higher operation, followed by a verified
source-to-family comparison.

\paragraph{Two exact published nonvanishing tests and their boundaries.}
McCallum--Sharifi, \emph{A cup product in the Galois
cohomology of number fields}, author version dated
5 December 2003, pp.~10--14, explicitly assumes
Vandiver in Sections 4--5.  Lemma 4.1 uses its
cyclotomic-unit index consequence.  Proposition 4.2
expresses pairing with the cyclotomic root-of-unity
direction through the derivative of an odd Iwasawa
characteristic series; the following paragraph obtains
nonzero pairing when that component has lambda
invariant one.  The authors note that cyclicity suffices
for many parts.  The complete proposition proof and
the stated scope were read.  Neither its root-of-unity
input nor its integral hypotheses are supplied by
the auxiliary class \(u\) in our source problem.

Sharifi, \emph{Iwasawa theory and the Eisenstein ideal},
author pp.~6, 11 and 25--27, gives a different genuine
arithmetic test.  Under Hypotheses 1--2, including
\(p\mid B_{1,\theta\omega^{-1}}\) and
\(p\nmid B_{1,\omega\theta^{-1}}\), Theorem 5.6
equates the indicated cup-product surjectivity with
\(U_p-1\) generating the ordinary Eisenstein ideal.
Its entire proof was read, as were the exact
hypotheses and the subsequent finite-range computation.
Theorem 5.7 and Corollary 5.9 establish the required
generation and surjectivity for \(p<1000\).
They do not assert it at every prime.

To use that test on the surviving rank-three source,
one would still need an integral comparison from its
specified off-diagonal module to the indicated
Eisenstein ideal and the relevant generator or
Frobenius value.  The presently checked source
predicates and trace comparisons furnish neither.
A characteristic-zero point or a varying-weight
cuspidal family alone does not show that its
off-diagonal extension class lifts freely across
the reducibility quotient: products of other
off-diagonal entries can occur in its cocycle
equation.  No such cancellation is discarded here.

This last boundary is already visible in the actual matrix
identity.  Restrict a source representation to \(G_E\), put
\(D(g)=\operatorname{diag}(\chi^e(g),1,\chi^{-e}(g))\),
and write \(\rho(g)=(1+A_g)D(g)\) in its integral
split-residual lattice.  Then
\[
 A_{gh}=A_g+D(g)A_hD(g)^{-1}
                    +A_gD(g)A_hD(g)^{-1}.
 \tag{TP.4}
\]
The negative entry \(A_{21}\) therefore satisfies its
linear cocycle equation only after quotienting out the
three terms
\[
 A_{21}(g)A_{11}(h)
 +\chi^{-e}(g)A_{22}(g)A_{21}(h)
 +\chi^{-2e}(g)A_{23}(g)A_{31}(h).
 \tag{TP.5}
\]
They lie in the square of the residual maximal ideal,
but need not vanish in a higher coefficient quotient.
Absolute irreducibility does not set them to zero:
by the characteristic-zero Burnside theorem the linear
span of the matrices \(\rho(g)\) is the full matrix
algebra, so each off-diagonal entry is a nonzero
function of \(g\).  In particular the last product
in (TP.5), with independently varying \(g,h\), is
not identically zero.  This does not assert that
the sum in (TP.5) is a nonzero cohomology class.
It proves why one cannot infer vanishing of the
isolated extension's Bockstein merely from the
existence or irreducibility of the full lift.

This audit therefore proves no additional Selmer
upper bound, source exclusion, or vanishing degree.

% END evidence/unitary_audit/reducible_comparison/actual_source_pairing_audit.tex


% BEGIN evidence/unitary_audit/reducible_comparison/hsieh_unitary_scope.tex
\subsection{The ordinary Eisenstein theorem of Hsieh and the actual residual triple}
\label{sec:kt-hsieh-unitary-scope}

This is a direct applicability audit of a published-method input.
It supplies no new vanishing of a Selmer group.  We use the author
version of M.-L.~Hsieh, \emph{Eisenstein congruence on unitary groups
and Iwasawa main conjectures for CM fields}, dated
16 September 2025, retained at
\url{https://www.math.ntu.edu.tw/~mlhsieh/research/IMCCM.pdf}.

\begin{proposition}[KT-HSIEH-DIRECT-SCOPE]
Let $p>3$, let $E/\mathbf Q$ be imaginary quadratic with
$p\nmid\operatorname{disc}(E)$, and let $e$ be even.
The residual representation with semisimplification
\[
       \omega^e\oplus1\oplus\omega^{-e}\quad\hbox{on }G_E
       \tag{HU.1}
\]
does not match the residual Eisenstein branch used in Hsieh's
Theorems 7.18 and 8.14, even after a common scalar-character
twist and a permutation of the three constituents.
\end{proposition}
\begin{proof}
The source assumes that $E$ is $p$-ordinary: for $E/\mathbf Q$
this requires $p$ to split.  Hence there is no direct application
when $p$ is inert.

Suppose $p$ splits and write its chosen ordinary place as $v$,
with conjugate place $\bar v$.  Denote the source's finite branch
character by $\xi$, to distinguish it from our $\psi=\omega^e$.
Hypothesis 7.18(2), printed p.~91, requires that $\xi$ be
unramified at $\bar v$; it also imposes a condition at $v$,
which is not needed for this obstruction.  The exact normalized
three-character parameter in its equation (8.2), printed p.~95,
has residual constituents
\[
       \bar\xi^{-c}\omega,\qquad 1,\qquad
       \bar\xi\omega^{-1}.
       \tag{HU.2}
\]
On $I_v$ the ratio of the first two is therefore $\omega$,
because $\xi^c$ is unramified there.

By contrast, every ratio of two constituents in (HU.1)
restricts to a power of $\omega$ with exponent
$0,\pm e,\pm2e$.  All these exponents are even.
Since $\omega(I_v)=\mathbf F_p^\times$ has even order $p-1$,
none is the character $\omega$.  Common scalar twists cancel
in every ratio, and permutations leave the collection of ratios
unchanged.  Thus neither operation removes the mismatch.
The same argument works at whichever split place is selected
for the CM-type.
\end{proof}

There is a second, independent direction-of-inequality issue.
In the same source, Theorem 7.18 constructs a surjection to
$\Lambda/(L^S(\Psi))$ when $L^S(\Psi)$ is a nonunit.
Definition 7.19 and Corollary 7.20, printed pp.~91--92, give
\[
 \operatorname{Eis}(\Psi,S)
 \subset
 \bigl(L_p^S(-1,\Psi_+\tau_{E/\mathbf Q})
                         L_p^S(\Psi,\Sigma)\bigr).
 \tag{HU.3}
\]
The congruence component in that definition is constructed
under the nonunit hypothesis.  This is a production of
congruences from an $L$-value divisor, and does not state
that a unit $L$-value excludes every cuspidal congruence.
Even if one separately defined an appropriate congruence
ideal in the unit case, the displayed containment would
only put that ideal inside $\Lambda$.
The subsequent Selmer inequalities likewise have the
direction ``Selmer length at least congruence valuation
at least $L$-value valuation.''  A zero lower bound
does not force Selmer vanishing.

The full proof of Theorem 7.18, Definition 7.19,
Corollary 7.20, the exact parameter and GMA setup on
pp.~95--96, and the full proof of Theorem 8.14 on
pp.~101--102 were read for this audit.  The underlying
Hida theory, Fourier--Jacobi nonvanishing and Galois
attachment results are imported source inputs, not
newly reconstructed proofs.  The obstruction above
concerns this precise branch and these theorem hypotheses;
it is not a claim that every possible unitary Eisenstein
comparison has the same restriction.

Consequently this theorem does not provide the missing
integral comparison for our odd off-diagonal module,
nor does it supply a necessary three-unit condition for
the surviving primitive rank-three source.

% END evidence/unitary_audit/reducible_comparison/hsieh_unitary_scope.tex


% BEGIN evidence/main_gate/massey_literature_scope.tex
\subsection{Recent Massey theorems and the fixed ordinary problem}

\paragraph{KT-MASSEY-LITERATURE-SCOPE.}
Three further primary-source comparisons were made against the
specific relative obstruction in
KT-ORDINARY-CUBIC-MASSEY-CONDITION. None of the inspected statements
evaluates that obstruction.

Maire--Min\'a\v{c}--Ramakrishna--T\^an,
\emph{On the strong Massey vanishing property for number fields},
Mathematische Annalen 395 (2026), article102,
\url{https://doi.org/10.1007/s00208-026-03518-1},
Theorem1.1, concerns the absolute and maximal tame Galois groups
when $\zeta_p\notin K$. Definition1.1 and the following paragraph
distinguish the existence of one lifting defining system from the
lifting of every prescribed defining system. Remark1.1 explicitly
allows additional ramified primes in the construction; Theorem1.2
retains such primes for higher coefficient lifts. The end of the
introduction fixes trivial coefficients. Our problem has a fixed
finite set $S$, nontrivial character coefficients, and a specified
ordinary defining cochain. Trivializing the coefficient characters
does not supply descent with this fixed ramification set and
ordinary normalization.
Thus the inspected result cannot simply be applied to (CM.6).
The complete statements and these scope paragraphs were read;
the proof of the global embedding theorem was not re-proved.

Ray--Sujatha, \emph{Massey products and the Iwasawa theory of fine
Selmer groups}, author preprint arXiv:2508.17156v1,
\url{https://arxiv.org/pdf/2508.17156v1}, Proposition4.6 on p.14,
relates unobstructedness in an infinite extension to a fine-Selmer
$\mu$ condition under weak Leopoldt. Theorem7.1 on p.24 gives an
equivalence with spanning by products formed from repeated copies
of the extension character $\kappa$ and varying permitted defining
systems. Theorem6.10 and Corollary6.11 on p.23 specify that system.
These are neither evaluation of the fixed alternating-weight
product $(u,v,u)$ nor finite-level vanishing. In addition, the
ordinary application in Section5 assumes weight two and residual
irreducibility over $K(\mu_p)$, explicitly before Proposition5.2.
The residual three-character sum here does not satisfy those
hypotheses. The cited statements, their displayed proofs, and the
Section5 setup were read, without importing an unstated comparison
between the two ordinary deformation problems.

Ahlqvist--Carlson, \emph{Massey products in the \'etale cohomology
of number fields}, Journal f\"ur die reine und angewandte
Mathematik823 (2025),61--112,
\url{https://doi.org/10.1515/crelle-2025-0006},
Corollary4.10 on p.88 and Theorem4.12 on p.90 give explicit
arithmetic formulas on $X=\operatorname{Spec}\mathcal O_K$.
The latter uses an unramified degree-$p$ extension and ideal/norm
data to evaluate a repeated product, modulo its stated
indeterminacy. Its complete statement and short proof were read;
the preceding resolution calculations were not independently
re-proved. In our problem the positive direction may ramify
above $p$, and the local ordinary normalization must be retained.
Remark4.14 on p.93 already sketches the comparison for a punctured
arithmetic curve containing the required roots of unity, referring
to Sharifi's Theorem5.4. That observation must be retained rather
than treating every punctured-curve comparison as absent.
What is still required here is a comparison with the precise
relative ordinary defining system, followed by an actual evaluation
of its norm data. Neither follows merely from identifying a
class-group character.

\paragraph{Arithmetic-geometric location.}
The distinction is between extensions of finite coefficient
sheaves on a fixed punctured arithmetic curve, with local
conditions, and extensions after allowing additional punctures
or after passage to an infinite tower. The finite sheaf setting
is the one already specified at (FW.1)--(FW.2), with SGA4,
Expos\'eIX,3.2--3.3 as its Kummer-theoretic reference. This audit
does not change the remaining arithmetic Triple-NV gate in the
Teichm\"uller-tower crosswalk to the \emph{Esquisse}, Section2,
and the 1983 letter to Faltings, or prove a new integral Tate vanishing.

% END evidence/main_gate/massey_literature_scope.tex


% BEGIN evidence/main_gate/indivisibility_literature_scope.tex
\subsection{Further indivisibility inputs and the selected quadratic value}

\paragraph{KT-INDIVISIBILITY-LITERATURE-SCOPE.}
The missing modular-symbol assertion concerns the specified
$j=2$ value of the integral Eisenstein symbol, hence the same
negative fundamental discriminant in
$\beta_3(\chi_D)\beta_b(\chi_D)$.
It is a noncentral value: $L(E,\chi_D,3)$ with
$k=(p+9)/2\geq20$.  The following source checks record why
the particular imported results inspected do not prove that
this value is a unit.  They do not claim that no applicable
theorem can exist.

Heumann--Vatsal, \emph{Modular symbols, Eisenstein series,
and congruences}, J.~Th\'eor.~Nombres Bordeaux \textbf{26} (2014),
709--756, Theorem~5.1, pp.~747--748, gives nonvanishing with
characters of growing odd prime-power order and uses $j=0$.
The already checked Corollary~4.7 gives our quadratic
coefficient formula but no nonzero selected quadratic value.
Theorem~5.2 transfers congruences without adding that unit.
The full Theorem~5.1 proof was read in the separate detector
audit.

Bruinier, \emph{Nonvanishing modulo $\ell$ of Fourier
coefficients of half-integral weight modular forms},
Duke Math.~J.~\textbf{98} (1999), 595--611, Theorem~1,
requires the residue prime to be prime to the auxiliary
Legendre prime times one less than that prime.  It therefore
does not impose that condition at our chosen $q\equiv1\pmod p$.
Another auxiliary prime is not excluded.  More basically,
Theorems~2--4 concern one half-integral weight eigenform;
Theorem~6 and its ensuing odd-index remark concern one
Dirichlet value.  No eigenform with fundamental coefficients
equal to our product has been constructed here.  The exact
statements on author pp.~6 and 10--13 were read, together
with the BJKOSV survey, pp.~6--7 and 13--14; no foundational
whole-proof audit is claimed.

Ahlgren--Boylan, \emph{Central critical values of modular
$L$-functions and coefficients of half-integral weight modular
forms modulo $\ell$}, Amer.~J.~Math.~\textbf{129} (2007),
429--454, author version dated 29 June 2005,
Theorem~1.1, p.~2,
\url{https://people.math.sc.edu/boylan/reprints/THETAFINALNEW.pdf},
concerns central quadratic values of a cuspidal newform of
even weight and odd squarefree level.  Its hypothesis that
a minimum valuation is attained for two discriminants does
not say that this minimum is zero in our normalization.
Only the complete theorem statement and hypotheses were
retrieved; the later PDF download failed, so no local copy
or whole-proof reading is claimed.

Kayath--Lane--Neifeld--Ni--Xue, \emph{Subspaces spanned by
eigenforms with nonvanishing twisted central $L$-values},
Canadian J.~Math.,
\url{https://doi.org/10.4153/S0008414X25101697},
Theorem~1.2 and Proposition~2.4, pp.~7--8 of the author
version, give a complex spanning construction that divides
by selected nonzero complex central values.  The full
Proposition~2.4 proof was read.  Such division does not
establish saturation of an integral lattice at $p$;
the chosen values could be divisible by $p$.  Vatsal,
\emph{Canonical periods and congruence formulae},
Theorem~0.3, author p.~3, and its Eisenstein discussion
on p.~14, applies in its stated quadratic application to
ordinary elliptic curves with rational three-torsion and
squarefree conductor.  These inspected statements do not
give the arbitrary-prime, higher-weight, specified
noncentral coefficient required here.

\paragraph{KT-KIMURA-DISPLAYED-PROOF-SCOPE.}
I.~Kimura, \emph{Indivisibility of special values of zeta
functions associated to real quadratic fields},
RIMS K\^oky\^uroku Bessatsu B4 (2007), Theorem~3.1
and its full proof, printed pp.~11--15,
\url{https://www.kurims.kyoto-u.ac.jp/~kenkyubu/bessatsu/open/B4/pdf/B04_001.pdf},
was checked for a possible $p$-ramified fixed-weight input.
Two displayed steps need additional justification for
that application.  This is a gap in the displayed
argument as used here, not a counterexample to its theorem.

The form in the proof is first depleted at $p$.
Its surviving Fourier indices are prime to $p$.
The later sieves and $U_\ell,V_\ell$ operations with
$\ell\ne p$ preserve this property: multiplying or
dividing an index by $\ell$ cannot change its divisibility
by $p$.  Thus the final nonzero coefficient $H(r,n\ell)$
has $p\nmid n\ell$.  Writing $n\ell=Df^2$ then gives
$p\nmid D$.  This coefficient construction does not as
written yield the stated $p$-ramified discriminant
$D_\ell=d_\ell\ell p$.

There is also a separate coefficient-comparison issue.
Write $a=\ell^{r-1}$, $b=\ell^{2r-1}$ and $c=1-a+b$.
Equation~(3.6) is $U_\ell G_r=cV_\ell G_r$.
At the chosen index $D_0\ell^3$, the normalized
$U$ coefficient is $1+b+b^2-a-ab$ and the
$V$ coefficient is $c$.  Hence the comparison required is
\[
 (1+b+b^2-a-ab)-c^2
   =a(1-a+b)-b
   =a(1-\ell^r)(1-a),
 \tag{KS.1}
\]
where the last equality uses $b=a\ell^r$.
The printed subtraction instead compares the $U$
coefficient with $c$; it addresses $U\ne V$, not
$U\ne cV$.  Adding $\ell^{r-1}\ne1\pmod p$
repairs this coefficient comparison when an admissible
residue can be chosen.  It does not repair the distinct
support issue.  No result in this paper relies on the
unjustified $p$-ramified conclusion.

K.~Ono, \emph{Indivisibility of class numbers of real
quadratic fields}, Compositio Math.~\textbf{119} (1999),
1--11, Theorem~2, and D.~Byeon, \emph{Indivisibility of
special values of Dedekind zeta functions of real quadratic
fields}, Acta Arith.~\textbf{109.3} (2003),
Theorems~1--2, were also inspected for individual
quadratic-field indivisibility.  Only their indicated
statements and initial construction were read; neither
was subjected to a whole-proof audit in this attempt.
Their precise retrieval records, and the fuller bounded
quadratic-symbol search with its source-by-source scope,
are retained in
\path{main_gate/indivisibility_source_notes.md} and
\path{lean_chain/QUADRATIC_SYMBOL_NONVANISHING_SCOPE.md}.

These checks leave simultaneous auxiliary-unit existence,
the nonessentially-selfdual source exclusion, and the
original strict rational class component open.

% END evidence/main_gate/indivisibility_literature_scope.tex


% BEGIN evidence/soule_comparison.tex
\section{Soul\'e's 2003 bound and the explicit constants}
\label{sec:soule}

This section records a source comparison; Soul\'e's theorem
is cited and is not reproved. Let $F$ have degree
$r=r_1+2r_2$ and absolute discriminant $D$, and let
$A=\mathcal O_F$. For a finite abelian group $G$,
$\operatorname{card}_b(G)$ denotes the largest divisor of
$|G|$ whose prime factors are all strictly greater than $b$.
Thus it removes the entire primary parts at primes at most $b$.

Soul\'e, \emph{A bound for the torsion in the K-theory of
algebraic integers}, Documenta Mathematica, Extra Volume Kato
(2003), 761--788, Theorem~2, \S4.2, printed p.~779, states
\[
 \operatorname{card}_{\max(r+1,2m+2)}
       K_m(\mathcal O_F)_{\mathrm{tors}}\leq k(m).
 \tag{SO.1}
\]
The definitions in \S\S3.2, 4.1--4.2, pp.~777--778, are
\[
\begin{aligned}
 N&=2m+1,&
 \widetilde w&=r_1N(N+1)/2+r_2N^2-N,\\
 k(m)&=h(m,N),&
 h(q,N)&=a(N)^{(N+1)c(\widetilde w-q-1,N)},\\
 c(j,N)&=\binom{a(N)}{N+j},&
 a(N)&=2^{N(r+3)}
                  \left(\prod_{i=1}^N T_i\right)^{r/2}.
\end{aligned}
 \tag{SO.2}
\]
The constants $T_i$ are those in Proposition~2,
pp.~774--775; this is the exact recursive constant used by
the theorem. Its fully explicit simplified upper bounds
are Proposition~4, \S4.4, pp.~779--780:
\[
\begin{array}{ll}
\text{general number fields }(r\geq2):&
 \log\log k(m)\leq220\,m^4\log m\,
                   r^{4r}\sqrt D(\log D)^{r-1},\\[2pt]
\text{class number one }(r\geq2):&
 \log\log k(m)\leq210\,m^4\log m\,
                   r^{4r}\sqrt D(\log D)^{r-1},\\[2pt]
\text{imaginary quadratic}:&
 \log\log k(m)\leq1120\,m^4\log m\log D,\\[2pt]
\text{imaginary quadratic, class number one}:&
 \log\log k(m)\leq510\,m^4\log m\log D.
\end{array}
 \tag{SO.3}
\]
The general estimates are used for $m\geq2$ as in their proof,
pp.~781--784. These coefficients are the integers
$220,210,1120,510$; they are not powers of two, eleven, or five.
For $F=\mathbf Q$, Proposition~4(iv), printed p.~780, gives
\[
\begin{aligned}
 \log\log k(m)&\leq8m^4\log m\qquad(m\geq9),\\
 \log\log k(7)&\leq40545,\qquad
 \log\log k(8)\leq70130.
\end{aligned}
 \tag{SO.4}
\]
The pages containing (SO.1), (SO.3), and (SO.4) were checked
visually against the attached PDF.

\subsection{Comparison in the degrees of interest}

For $m=4k$, $k\geq3$, Soul\'e gives
\[
 \log\left(
   \operatorname{card}_{8k+2} K_{4k}(\mathbf Z)
 \right)
 \leq \exp\!\bigl(8(4k)^4\log(4k)\bigr).
 \tag{SO.5}
\]
The inherited VK-EXP7 result instead bounds the complete order
by
\[
       \log|K_{4k}(\mathbf Z)|\leq C_7(2k)^7.
 \tag{SO.6}
\]
There are two separate improvements: polynomial growth for
the logarithm in (SO.6), and inclusion of every primary part.
Soul\'e's excluded-prime cutoff grows with the degree.
The verified Vandiver range removes that cutoff only in
degrees where all its primes fall in the verified range;
it must not be silently removed for every $k$.
The fresh KT-C7-EXPLICIT derivation below makes one admissible
archive-based constant explicit. KT-QUARTIC-BOUND gives the
sharper, independently derived bound
$\log|K_{4k}(\mathbf Z)|<\frac25(2k)^4$.

\paragraph{Arithmetic-geometric meaning.}
Soul\'e's Theorem~3, \S4.3, printed p.~779, also states
\[
 \prod_{\substack{p\geq4n-1\\p\geq r+2}}
 \left|H^2_{\acute et}(\operatorname{Spec}\mathcal O_F[1/p],
                                      \mathbf Z_p(n))\right|
 \leq k(2n-2).
 \tag{SO.7}
\]
This is the direct scheme-theoretic comparison with the
integral Tate cohomology controlled by the Selmer certificates.
The mechanism of Soul\'e's proof is torsion homology of
arithmetic groups, a Hermitian reduction complex, and the
Hurewicz and Chern-class comparisons. It supplies no
vanishing theorem for the groups in (SO.7).

% END evidence/soule_comparison.tex


% BEGIN evidence/main_gate/soule_stabilizer_correction.tex
\subsection{KT-SOULE-STAB: the general-field stabilizer cutoff}
\label{sec:kt-soule-stab}

The transcription of Soul\'e's Theorem~2 and Proposition~4
in the comparison section agrees with the supplied original
PDF, pp.~779--780. In particular the coefficients are the
ordinary integers $220,210,1120,510$. There is nevertheless
a defect in the proof of its \emph{general-field} cutoff,
which does not affect the specialization to $\mathbf Q$.

\begin{proposition}[KT-SOULE-STAB]
Soul\'e's Lemma~9, pp.~777--778, asserting that a nontrivial
order-$p$ element of $\operatorname{GL}_N(\mathcal O_F)$
satisfies $p\leq1+\max([F:\mathbf Q],N)$, is false as
stated. A valid general replacement is
\[
                         p\leq N[F:\mathbf Q]+1.
                                                               \tag{SS.1}
\]
For $F=\mathbf Q$, the printed bound $p\leq N+1$ is valid.
\end{proposition}

\begin{proof}
Set $F=\mathbf Q(\sqrt5)$ and
$t=(\sqrt5-1)/2\in\mathcal O_F$, so $t^2+t-1=0$.
The matrix
\[
              M=\begin{pmatrix}0&-1\\1&t\end{pmatrix}
                \in\operatorname{SL}_2(\mathcal O_F)
\]
has characteristic polynomial $X^2-tX+1$.
Its roots are $\zeta_5,\zeta_5^{-1}$ after one embedding
of $F$; equivalently direct reduction using $t^2=1-t$
gives $M^5=I$ and $M\ne I$. It has order five, whereas
$1+\max(2,2)=3$. The equality $M^5=I$ is also checked
symbolically in the accompanying certificate.

For the replacement, regard $F^N$ as a rational vector space
of dimension $rN$, where $r=[F:\mathbf Q]$. A nonidentity
element of prime order $p$ has a primitive $p$th root of
unity as an eigenvalue; its rational minimal polynomial
therefore contains $\Phi_p$, of degree $p-1$.
Thus $p-1\leq rN$, proving (SS.1). For $r=1$ this is
the original cutoff. This argument does not require that
$\Phi_p(M)=0$ on every eigenspace, since a fixed subspace
may also be present.
\end{proof}

The precise invalid inference in the printed proof is that
$\Phi_p$ is irreducible over $F$ whenever $F$ does not
contain all $p$th roots of unity. A proper intersection
$F\cap\mathbf Q(\zeta_p)$ already makes it reducible.
The matrix example refutes that lemma, not the numerical
conclusion of Theorem~2, which might admit a separate repair.

Replacing just this step in the published homology argument
and carrying its other inputs gives the safe cutoffs
\[
 \operatorname{card}_{rN+1}
 H_q(\operatorname{GL}_N(\mathcal O_F),\mathbf Z)_{\rm tors}
       \leq h(q,N),\qquad
 \operatorname{card}_{r(2m+1)+1}
 K_m(\mathcal O_F)_{\rm tors}\leq k(m).
                                                               \tag{SS.2}
\]
Indeed every finite cell stabilizer has all prime divisors
at most $rN+1$ by Cauchy's theorem and (SS.1); the spectral
sequence and Hurewicz exclusions are then valid away from
those primes. The cell constants are unchanged. This is
a repair of the identified step conditional on the other
published cell and homology estimates, not a fresh full
proof of Soul\'e's paper. At $F=\mathbf Q$, (SS.2) gives
exactly the original cutoff $2m+2$, so the comparison used
for $K_{4k}(\mathbf Z)$ remains unchanged.


% END evidence/main_gate/soule_stabilizer_correction.tex


% BEGIN evidence/lean_chain/c7_section.tex
\section{An explicit choice of the constant in K17-VK-EXP7}
\label{sec:kt-c7-explicit}

\paragraph{Finding KT-C7-EXPLICIT.}
Under the stipulated archive theorems and the carried K17-VK-PT and
K17-VK-EXP7 argument, one explicit admissible choice is
\[
 C_7=10^{\,10^{3500}},\qquad
 \log|K_{4k}(\mathbf Z)|\le C_7(2k)^7\quad(k\ge1).
\]
This is a very coarse upper bound for a permissible constant.  It is
not the least constant and offers no numerical vanishing criterion.
Its purpose is to remove the unevaluated cohomology constants from the
source reconstruction requested in item 4.

All source locators below are relative to
\path{src/Development/AutomorphicArithmetic/AutomorphicVarieties/}
at commit \path{7cbf7a0c32fca1badf05931879a37c3f8a2b03b4}.
The theorem \path{cohBoundFree},
\path{AutomorphicForms/TemperedLocalComponents/CohBoundFree.lean}:93--99,
returns existential constants.  Its helper \path{cohBound_of}:69--91
obtains a base-level witness $c$ at line 75 and returns
$(c3^{dn^2},dn^2)$ at line 79, with $d=[F:\mathbf Q]$.
The witness ultimately comes from finite orbits and stabilizer
cohomology.  The following retention of the source's explicit lower
bounds replaces it with an integer expression.

\subsection{An explicit arithmetic-group cohomology bound}

Let $F$ be imaginary quadratic of absolute discriminant $D$ and let
$n\ge1$.  Define
\[
 I=D^n(2n)^{2n},\quad R=n^nI^2,\quad T=(2R+1)^{2n},
\]
\[
 A=I^{1+2n^2},\quad B=(T+1)^{2n^2},\quad H=AB^2T^n,
 \qquad b(D,n)=HAB^H.
\tag{KT-C7.1}
\]
\paragraph{Finding KT-COH-EXPLICIT.}
For every finite-dimensional complex representation $V$ of
$\mathrm{GL}_n(\mathcal O_F)$,
\[
 \sum_{q\ge0}\dim H^q(\mathrm{GL}_n(\mathcal O_F),V)
 \le b(D,n)\dim V.
\tag{KT-C7.2}
\]

\paragraph{Proof.}
\path{HermiteNormalForms/SmulStrictMono.lean}:17--19,244--249 defines
$\mathrm{dd}=dn$, $I_0=D^n(dn)^{dn}$, $RR=n^nI_0^2$, and
\[
 \operatorname{orbBound}(q)
 =I_0^{1+dn^2}\bigl((2RR+1)^{dn}+1\bigr)^{dn^2(q+1)}.
\]
For $d=2$ this is exactly $AB^{q+1}$.
\path{HermiteNormalForms/ColsMinSmul.lean}:34--75 bounds the number
of group orbits of ordered $q$-simplices of its Hermitian nerve by
this integer.  The augmented permutation resolution is exact by the
archive theorem \path{ConeExtensionAcyclic.sx_exact}, used at
\path{BuildingsResidueChars/BoundedComplexOfBoundedSimplicialComplex.lean}:18.

First bound a vertex stabilizer.  By
\path{HermiteNormalForms/Stabfin.lean}:24--26 choose a nonsingular
matrix $A_0$ whose $n$ columns are minimal vectors of the vertex $v$.
Every minimal vector of $v$ is in the set
$\operatorname{nearVec}(A_0)$ defined at
\path{HermiteNormalForms/SmulStrictMono.lean}:258--259.
By \path{HermiteNormalForms/Nearvec.lean}:9--38 this set has at most
$T$ elements.  A stabilizer element is determined by the images of
the $n$ columns: equality on those columns gives equality after
multiplication by $A_0$, and then equality because $\det A_0\ne0$.
This is the injection in \path{Stabfin.lean}:35--50, now with its
cardinality retained.  Hence
\[
 |\operatorname{Stab}(v)|\le T^n.
\tag{KT-C7.3}
\]

Choose one representative from each orbit of ordered $1$-simplices;
there are at most $AB^2$.  The injection in
\path{BuildingsResidueChars/AdelicRepDimensionBounded.lean}:22--68
maps the $q+1$ distinct vertices of an injective ordered $q$-simplex
into the disjoint union of the stabilizers of the first vertices of
these representatives.  This union has at most $AB^2T^n=H$ elements.
There is consequently no injective $q$-simplex when $q\ge H$.

The original resolution includes repeated tuples.  The normalized
alternating endomorphism is a chain map by
\path{BuildingsResidueChars/AltChainOfPermSplit.lean}:5--46.
It is the identity on the augmentation and zero in a degree in which
all tuples repeat a vertex, by
\path{BuildingsResidueChars/RepHomAlternatingSumEqZero.lean}:10--57.
The projective-resolution comparison in
\path{BuildingsResidueChars/ExtVanishingOfResolutionEndo.lean}:32--38
therefore gives $H^q=0$ for $q\ge H$.
Projectivity follows from finite stabilizers over $\mathbf C$, exactly
as in \path{AutomorphicForms/TemperedLocalComponents/BddArith.lean}:41--49.

In degree $q$, an equivariant homomorphism from the permutation term
to $V$ is determined by one vector for each orbit.  Its dimension is
at most $AB^{q+1}\dim V$; this is the source argument
\path{BuildingsResidueChars/PermutationModuleHomologyRank.lean}:4--34.
Cohomology is a subquotient of this Hom term.  Consequently
\[
 \sum_q\dim H^q
 \le\left(\sum_{q=0}^{H-1}AB^{q+1}\right)\dim V
 \le HAB^H\dim V.
\]
This proves KT-COH-EXPLICIT. \hfill$\square$

\subsection{The principal-level coefficient}

If $F$ has class number one, write $K(m)$ for the principal subgroup
of $\mathrm{GL}_n(\widehat{\mathcal O}_F)$ of level $m$.  We claim that
the same total adelic cohomology quantity used by \path{cohBoundFree}
satisfies
\[
 \operatorname{cohDim}(\mathrm{GL}_n(F),K(m),\xi)
 \le b(D,n)m^{2n^2}\dim\xi.
\tag{KT-C7.4}
\]

An adelic matrix $g$ determines a full $\mathcal O_F$-lattice $L$ in
$F^n$ whose completion is $g\widehat{\mathcal O}_F^n$.
For an explicit construction, choose an integer $a>0$ with
\[
 a\widehat{\mathcal O}_F^n\subseteq
 g\widehat{\mathcal O}_F^n\subseteq
 a^{-1}\widehat{\mathcal O}_F^n.
\]
The middle lattice gives a submodule of the finite quotient
$a^{-1}\mathcal O_F^n/a\mathcal O_F^n$ by Chinese remainders; take
its inverse image for $L$.  The same finite quotient shows that its
completion is the prescribed one.  Since $\mathcal O_F$ is a PID,
$L=\gamma\mathcal O_F^n$ for a matrix $\gamma\in\mathrm{GL}_n(F)$.
Thus $g=\gamma k$ with $k\in K(1)$.
There is one rational orbit at level one, with stabilizer
$\mathrm{GL}_n(\mathcal O_F)$.

The fibers at level $m$ have cardinality
$[K(1):K(m)]\le m^{2n^2}$, by
\path{BuildingsResidueChars/LevelSubgroupIndexBoundOfCardOhat.lean}:22--38.
The Shapiro fiber argument
\path{BuildingsResidueChars/OrbitBoundOfShapiroFunctorial.lean}:3--23,
with KT-COH-EXPLICIT, proves (KT-C7.4).
This is the one-orbit, base-level-one specialization of the helper
\path{gen} in \path{CohBoundFree.lean}:28--60.

\subsection{Substituting six explicit fields}

Take $F_0=\mathbf Q(\sqrt{-3})$ for the rank-three branch and take
the six rank-two fields with absolute discriminants
\[
 \mathcal D=\{3,7,11,19,43,67\}.
\]
They are pairwise coprime, so they satisfy the same six-field
requirement as K17's illustrative choices.
All have class number one.  To check this, Minkowski's bound says
that every ideal class has an integral representative of norm at
most $(2/\pi)\sqrt D$ (J.~S.~Milne,
\emph{Algebraic Number Theory}, version 3.08, Theorem 4.3, p.~70;
\url{https://www.jmilne.org/math/CourseNotes/ANT.pdf}).
For $D=3,7$ this is $<2$.  For $D=11,19$ it is $<3$, and $2$ is
inert.  For $D=43$ it is $<5$, and $2,3$ are inert.  For $D=67$ it
is $<6$, and $2,3,5$ are inert.  In the last two cases the only
possible nonunit norm is $4$, whose ideal is the principal ideal
$(2)$.  Thus all representatives are principal.  Inertness follows
from $-D\equiv5\pmod8$ at $2$ and the indicated quadratic
nonresidues at the odd primes; the accompanying exact script checks
each instance.

K17, \path{evidence/work/k_quantitative_refinement.md},
sections 4--5, permits constants $B_0$ for rank three at level $6$
and $B_D$ for rank two at level $18$.  Hence (KT-C7.4) allows
\[
 B_0=6^{18}b(3,3),\qquad B_D=18^8b(D,2).
\tag{KT-C7.5}
\]
The levels and the use of the full rational group are explicit in
\path{AutomorphicForms/TemperedLocalComponents/KlevThree.lean}:18--28
and \path{GaloisLibrary/CrystallineBaseChange/KlevThree.lean}:17--27:
both use \path{cohBoundFree ... n top} and replace a level $M$ by
$3M$.  The carried K17 exponent-seven proof therefore permits the
fully explicit integer
\[
 C_7^*=20\left(6^{18}b(3,3)+32\,18^8
                 \sum_{D\in\mathcal D}b(D,2)\right).
\tag{KT-C7.6}
\]
Every quantity in this expression is now given by integer additions,
multiplications, and powers in (KT-C7.1).

\subsection{Evaluation by elementary decimal bounds}

For $(D,n)=(3,3)$, $I=1{,}259{,}712<10^7$ and $2R+1<10^{14}$.
Thus
\[
 T<10^{84},\quad A<10^{133},\quad B\le10^{1512},
 \quad H<10^{3409},\quad \log_{10}b(3,3)<10^{3413}.
\]
For $n=2$ and $D\le67$, $I\le1{,}149{,}184<10^7$ and again
$2R+1<10^{14}$.  Hence
\[
 T<10^{56},\quad A<10^{63},\quad B\le10^{448},
 \quad H<10^{1071},\quad \log_{10}b(D,2)<10^{1074}.
\]
The $+1$ in $B$ causes no loss: an integer $T<10^u$ has
$T+1\le10^u$.
Also $6^{18}<10^{15}$, $18^8<10^{11}$, and
$20(10^{15}+192\cdot10^{11})<10^{18}$.
It follows that
\[
 C_7^*<10^{18+10^{3413}}<10^{10^{3414}}<10^{10^{3500}}.
\]
This proves KT-C7-EXPLICIT.  The exact-arithmetic replay is
\path{check_c7_arithmetic.py}; it verifies these inequalities without
trying to expand a double-exponential integer.

\paragraph{Scope.}
The source-standing result and the constant calculation are distinct.
The former is an explicit substitution in named archive theorems;
the latter is a new elementary quantitative refinement of the source's
finite-orbit proof, using those theorems as inputs.  Neither is a new
Lean kernel certificate.  The geometric constants arise from
Hermitian reduction complexes and finite congruence-level fibers.
Their interpretation as a bound for integral Tate pieces over
$\operatorname{Spec}\mathbf Z$ uses the carried K17-VK-PT bridge.
They do not establish vanishing or Vandiver's conjecture.

% END evidence/lean_chain/c7_section.tex

\section{Cote 149 and the geometric questions}

% BEGIN evidence/geometry_tower/finite_field_sections.tex
\subsection{KT-FF-SEC: finite-field section surjectivity is false}
\label{sec:kt-ff-sec}

Here the prime-to-$p$ arithmetic quotient means
\[
 \Pi_C^{(p')}=\pi_1(C)/
 \ker\bigl(\pi_1(C_{\overline k})\longrightarrow
                   \pi_1(C_{\overline k})^{(p')}\bigr);
\]
the full arithmetic group $G_k$ is retained. This is the convention
of AN\_GRH A05 and of the present question.

\begin{proposition}[KT-FF-SEC]
The map from rational points to conjugacy classes of sections of
$\Pi_C^{(p')}\to G_k$ is not surjective for all proper hyperbolic
curves over finite fields. The same failure occurs for the full
arithmetic fundamental group.
\end{proposition}

\begin{proof}
If $k=\mathbf F_q$, then $G_k\simeq\widehat{\mathbf Z}$.
Every surjection $E\twoheadrightarrow\widehat{\mathbf Z}$ of
profinite groups splits. Indeed, choose $g\in E$ mapping to $1$.
For every open normal subgroup $V\subset E$, the homomorphism
$\mathbf Z\to E/V$ sending $1$ to the image of $g$ factors through
a finite cyclic quotient of $\mathbf Z$. These homomorphisms give
a compatible continuous homomorphism
$s:\widehat{\mathbf Z}\to E$. Its composite with the given
surjection fixes the dense subgroup $\mathbf Z$, and hence is the
identity. This familiar finite-field observation is also the opening
paragraph of Stix, \emph{Rational Points and Arithmetic of Fundamental
Groups}, Chapter 15, p.~197.

For an explicit curve with no rational point, let $C/\mathbf F_2$
be the smooth projective curve with function field
\[
 \mathbf F_2(x,y),\qquad
 y^2+y=1+\frac{x^2+x}{x^3+x+1}.                         \tag{FF.1}
\]
The cubic $x^3+x+1$ has no root in $\mathbf F_2$ and is separable.
Its three roots over $\overline{\mathbf F}_2$ are simple poles of
the right-hand side: none is $0$ or $1$. The rational function has
no other poles. A function of the form $h^2+h$ has even pole order
at every pole, so (FF.1) defines a geometrically connected
Artin--Schreier cover of $\mathbf P^1$ of degree two. A simple pole
has different exponent two; the exact standard reference is
H.~Stichtenoth, \emph{Algebraic Function Fields and Codes}, second
edition, GTM~254 (2009), Proposition~3.7.8(c)--(d), p.~127.
Riemann--Hurwitz therefore gives
\[
 2g(C)-2=2(-2)+3\cdot2=2,
 \qquad g(C)=2.
\]
At each of the three rational points $x=0,1,\infty$ of
$\mathbf P^1(\mathbf F_2)$, the right-hand side is regular and
has value $1$. The corresponding fibre is
$y^2+y=1$, which has no $\mathbf F_2$-solution. Every rational
point of $C$ would map to one of these three rational points.
Thus $C(\mathbf F_2)=\varnothing$.

The fundamental exact sequence of SGA~1, Expos\'e IX,
Th\'eor\`eme 6.1, p.~195 in the 2003 recomposition, gives a
surjection $\pi_1(C)\to G_{\mathbf F_2}$. Its quotient
$\Pi_C^{(2')}\to G_{\mathbf F_2}$ is also surjective. The
preceding splitting argument applies to both, while their
rational-point domain is empty. This proves the claim. Since
$C$ is proper, no missing cuspidal or tangential section can
repair the assertion.
\end{proof}

\paragraph{Scope and attribution.}
The existence of sections over a finite field is standard, not a
new theorem. KT-FF-SEC supplies an explicit application answering
item~6(a) negatively. It does not conflict with the point-injectivity
statement AN\_GRH A05 at cote~149, scans~35--36. Injectivity and
surjectivity are distinct questions. It also does not refute
Grothendieck's section conjecture over finitely generated fields of
characteristic zero, whose base-field hypothesis is absent here.


% END evidence/geometry_tower/finite_field_sections.tex


% BEGIN evidence/geometry_tower/polycurve_counterexample.tex
\subsection{KT-POLY-FAIL: the maximal geometric prime-to-p quotient
can lose a whole fibre}
\label{sec:kt-poly-fail}

\begin{proposition}[KT-POLY-FAIL]
There is a smooth hyperbolic polycurve surface $X/\mathbf F_4$
and two distinct $\mathbf F_4$-points whose section classes in
$\Pi_X^{(2')}$ are equal. Thus the unrestricted extension of
AN\_GRH A05 from curves to hyperbolic polycurves, using the
maximal prime-to-$p$ quotient of the entire geometric fundamental
group, is false.
\end{proposition}

\begin{proof}
Write $k=\mathbf F_4=\mathbf F_2(\omega)$ with
$\omega^2+\omega+1=0$. Let $F/k$ be the smooth projective
genus-two curve
\[
 y^2+y=x^5,
\]
and write $\iota(x,y)=(x,y+1)$ for its hyperelliptic
involution. The unique point $q$ over $x=\infty$ is $k$-rational
and fixed by $\iota$. Set
\[
 B=\mathbf P^1_k\setminus\{0,1,\infty\},\qquad
 B'\longrightarrow B:\quad z^2+z=t+\omega.
                                                               \tag{PF.1}
\]
The second map is a geometrically connected finite \'etale
double cover, with deck involution $\tau(z)=z+1$.
Define
\[
 X=(F\times_k B')/\langle(\iota,\tau)\rangle.
                                                               \tag{PF.2}
\]
The diagonal action is free because its action on $B'$ is free.
The quotient exists (the product is quasi-projective), and finite
\'etale descent shows that $X\to B$ is smooth and proper with
geometric fibre $F$. Together with $B\to\operatorname{Spec}k$
this makes $X$ a hyperbolic polycurve. The fixed point $q$ descends
to a section of $X\to B$.

We work first over $\overline k$. Proper K\"unneth, SGA~1,
Expos\'e X, Corollaire 1.7, pp.~203--204, applied after pullback
to $B'$, and the section supplied by $q$, identify
\[
 \pi_1(X_{\overline k})\simeq
 D\rtimes H,\qquad
 D=\pi_1(F_{\overline k},q),\quad H=\pi_1(B_{\overline k}).
                                                               \tag{PF.3}
\]
The action of $H$ factors through the double-cover character
$H\twoheadrightarrow C_2$ and the induced pointed action
$\iota_*$ on $D$. To spell out the injection implicit in
(PF.3), the subgroup associated with $B'$ is
$D\times\pi_1(B'_{\overline k})$ by proper K\"unneth;
finite \'etale covers inject their fundamental groups. Its
quotient by $D$ is $\pi_1(B'_{\overline k})$, and descent along
$B'/B$ gives $H$. The descended fixed-point section splits this
extension and supplies exactly the stated action.

Let $\varphi:D\rtimes H\to Q$ be any finite quotient of odd
order. Pick $h\in H$ with nontrivial image in $C_2$.
On $\varphi(D)$, conjugation by $\varphi(h)$ is the involution
induced by $\iota_*$. Its order divides two. Its order also
divides the odd order of $\varphi(h)$, so it is the identity.
Hence the restriction $D\to\varphi(D)$ identifies $d$ and
$\iota_*(d)$ for every $d\in D$.

We claim that every finite odd quotient $D\twoheadrightarrow A$
on which $\iota_*$ is trivial is trivial. Let $Y\to F_{\overline k}$
be its connected \'etale Galois cover. The pointed invariance
lifts $\iota$ to an involution $j$ of $Y$ fixing the chosen point
over $q$ and commuting with $A$. Consequently
$Y\to\mathbf P^1_{\overline k}$ is Galois with group
$A\times\langle j\rangle$. Since $Y\to F_{\overline k}$ is
\'etale, every nontrivial inertia group of this composite has
order two. A direct product of an odd group with $C_2$ has a
unique element of order two, namely $j$. Thus the quotient
$Y/\langle j\rangle\to\mathbf P^1_{\overline k}$ is \'etale
with Galois group $A$. The geometric projective line has no
nontrivial finite \'etale connected cover (Riemann--Hurwitz),
so $A=1$. Applying this to $\varphi(D)$ proves
\[
 D\longrightarrow\pi_1(X_{\overline k})^{(2')}
 \quad\hbox{is trivial}.                                  \tag{PF.4}
\]
In fact (PF.3) now gives
$\pi_1(X_{\overline k})^{(2')}\simeq H^{(2')}$.

At $b=\omega\in B(k)$ the double cover (PF.1) has the two
rational points $z=0,1$. Choosing $z=0$ identifies the fibre
$X_b$ with $F$ over $k$. The points $(x,y)=(0,0)$ and $(0,1)$
of $F(k)$ give two distinct points $P,Q\in X_b(k)$.
The map $\pi_1(X_b)\to\Pi_X^{(2')}$ kills its geometric
subgroup by (PF.4). It therefore factors through the quotient
$G_k$, and both point sections of $\pi_1(X_b)\to G_k$ have
the same image. Their sections in $\Pi_X^{(2')}$ are equal
up to the usual geometric conjugacy. This proves the
counterexample.
\end{proof}

\paragraph{The precise obstruction.}
The missing operation is preservation of the fibre group when
passing from a hyperbolic fibration to the maximal prime-to-$p$
quotient of the entire geometric group. Here geometric monodromy
has order $p=2$. It disappears in every prime-to-$p$ quotient,
forcing the fibre's hyperelliptic involution to act trivially and
thereby killing the fibre. A fibrewise prime-to-$p$ completion
that retains the full geometric base group is a different group;
KT-POLY-FAIL makes no assertion against that alternative.


% END evidence/geometry_tower/polycurve_counterexample.tex


% BEGIN evidence/geometry_tower/artin_tower_counterexample.tex
\subsection{KT-ARTIN-FAIL: an affine tower of elementary fibrations}
\label{sec:kt-artin-fail}

\begin{proposition}[KT-ARTIN-FAIL]
The counterexample to prime-to-$p$ point separation can be taken
affine, with both steps in its polycurve tower elementary
fibrations in hyperbolic curves. It therefore applies to Artin
neighbourhoods understood to mean such towers, without the
additional product-embedding condition of a \emph{strongly}
hyperbolic Artin neighbourhood.
\end{proposition}

\begin{proof}
Keep $k$, $B$, $B'$, $\omega$ and $b=\omega$ from (PF.1).
Replace $F$ by the smooth projective curve
\[
 F:\quad y^2+y=x+\frac1x+\frac1{x+1}.
                                                               \tag{AF.1}
\]
Its three simple poles are $0,1,\infty$, so it has genus two
by Stichtenoth, Proposition~3.7.8(c)--(d). The corresponding
three ramification points $q_0,q_1,q_\infty$ are $k$-rational
and fixed by the involution $\iota:y\mapsto y+1$.
Put $F^\circ=F\setminus\{q_\infty\}$ and
\[
 X^\circ=(F^\circ\times B')/\langle(\iota,\tau)\rangle.
                                                               \tag{AF.2}
\]
This is a smooth affine surface. Its morphism to $B$ has a
smooth proper compactification obtained from $F\times B'$,
with complement the single section descended from $q_\infty$.
Both $X^\circ\to B$ and $B\to\operatorname{Spec}k$ are
therefore elementary fibrations with hyperbolic fibres.
The point $q_0$ gives a section of $X^\circ\to B$.

Fix compatible geometric base points using this section.
Let $\varphi:\pi_1(X^\circ_{\overline k})\twoheadrightarrow Q$
be a finite quotient of odd order. The image of the index-two
subgroup
\[
 \pi_1(F^\circ_{\overline k}\times B'_{\overline k})
 \subset\pi_1(X^\circ_{\overline k})
\]
has index dividing both two and $|Q|$, hence is all of $Q$.
Prime-to-two K\"unneth, SGA~1, Expos\'e XIII,
Proposition~4.6, pp.~310--311, applies to these smooth curves:
the required strong resolution is only in dimension at most
one. Consequently $Q$ is generated by the commuting images
$A$ of $\pi_1(F^\circ_{\overline k},q_0)$ and $H'$ of
$\pi_1(B'_{\overline k})$. In particular $A$ is normal in $Q$.

Take an element of $\pi_1(B_{\overline k})$ with nontrivial
double-cover character, and lift it using the $q_0$ section.
In this finite odd quotient, conjugation by its image on $A$
is induced by $\iota_*$. Indeed, the product trivialization
over $B'$ identifies the descent action on the pointed fibre
with $\iota$; any base-cover transport acts through $H'$, which
commutes with $A$ by prime-to-two K\"unneth. Only this relation
in $Q$ is used, not an exact full wild homotopy sequence for
the affine fibration. Consequently that conjugation on $A$
has order dividing two. It also has odd order, since $Q$ has
odd order. Thus $\iota_*$ acts trivially on $A$.

Let $Y\to F^\circ_{\overline k}$ be the connected \'etale
Galois cover with group $A$. The pointed trivial action of
$\iota_*$ gives a lift $j$ of $\iota$ of order two, fixing a
chosen point above $q_0$ and commuting with $A$.
The quotient map $F^\circ\to\mathbf A^1_x$ is finite, so
the normalization of $\mathbf A^1$ in the function field of
$Y$ is $Y$ itself. The composite $Y\to\mathbf A^1$ is
Galois with group $A\times\langle j\rangle$. Its inertia
at every point of $\mathbf A^1$ is either trivial or of
order two, because $Y\to F^\circ$ is \'etale. The only
order-two element in the direct product is $j$.
It follows that
\[
 Y/\langle j\rangle\longrightarrow\mathbf A^1_{\overline k}
\]
is \'etale and Galois of odd degree $|A|$.

There is no nontrivial finite odd Galois \'etale cover of the
geometric affine line in characteristic two. Indeed, its
smooth projective compactification can ramify only at infinity,
and ramification is tame because its Galois group has odd order.
If its degree is $d>1$ and the ramification index at infinity
is $e$, tame Riemann--Hurwitz reads
\[
 2g-2=-2d+\frac d e(e-1)=-d-\frac d e<-2,
\]
a contradiction. Hence $A=1$.

Since this holds for every finite odd quotient, the geometric
fibre group maps trivially to
$\pi_1(X^\circ_{\overline k})^{(2')}$.
The two distinct points of the fibre over $b$ corresponding
to $q_0$ and $q_1$ are rational over $k$ and have equal
arithmetic section classes in $\Pi_{X^\circ}^{(2')}$,
by the last argument of KT-POLY-FAIL.
\end{proof}

The result is compatible with KT-PRODUCT-SEP: an embedding
in a product of hyperbolic curves would contradict the
exhibited section collision, so this surface does not meet
that additional condition.


% END evidence/geometry_tower/artin_tower_counterexample.tex


% BEGIN evidence/geometry_tower/product_separation.tex
\subsection{KT-PRODUCT-SEP: the positive higher-dimensional consequence}
\label{sec:kt-product-sep}

\begin{proposition}[KT-PRODUCT-SEP]
Let $k$ be finitely generated over $\mathbf F_p$, and let $X/k$
be a geometrically connected variety. Suppose there are finitely
many $k$-morphisms $f_i:X\to C_i$ to hyperbolic curves such that
the map
\[
 X(k)\longrightarrow\prod_i C_i(k)
\]
is injective. Then the point-section map for the arithmetic
fundamental group with maximal prime-to-$p$ geometric quotient
is injective on $X(k)$.
In particular, this holds when $X$ embeds in a product of
hyperbolic curves, including any strongly hyperbolic Artin
neighbourhood in the sense of Schmidt--Stix, Definition~5.7.
\end{proposition}

\begin{proof}
Functoriality of the geometric prime-to-$p$ quotient gives
$\Pi_X^{(p')}\to\Pi_{C_i}^{(p')}$ over $G_k$ for every $i$.
If $P,Q\in X(k)$ have conjugate sections, their images under
every $f_i$ have conjugate sections. AN\_GRH A05 implies
$f_i(P)=f_i(Q)$ for every $i$. The assumed injectivity on
rational points gives $P=Q$.
\end{proof}

This is a direct corollary of the carried A05, not a new proof
of A05. No K\"unneth isomorphism and no exactness of a completed
fibration sequence is needed. The relevant definition is
A.~Schmidt and J.~Stix, \emph{Anabelian geometry with \'etale
homotopy types}, Annals of Mathematics 184 (2016), 817--868,
Definition~5.7 (author manuscript p.~17). Its additional
embedding condition must not be dropped: KT-POLY-FAIL concerns
general hyperbolic polycurves, where that condition is absent.


% END evidence/geometry_tower/product_separation.tex


% BEGIN evidence/geometry_tower/f149_review.tex
\subsection{KT-F149-REVIEW: the marked and unmarked statements}
\label{sec:kt-f149-review}

The complete proofs of K17 F149-GENERIC,
F149-MARKED-RECOGNITION, F149-BARE-EXTENSION and
F149-UNMARKED-CURVE have been read in this run. Their locators
are K17.tex, the subsections \emph{Generic recognition and the
necessary profinite input} and \emph{Why the generic marking
cannot simply be omitted}; the principal statements begin at
lines 13654, 13742 and 13908 of the supplied checkpoint-08
source. No fault was found in these arguments under their
stated hypotheses. This is a targeted proof review of those
statements, not an independent proof of the full anabelian
input theorems.

\paragraph{The positive statement and its exact scope.}
In the smooth characteristic-zero situation of F149-RELATIVE,
with normal integral base $S$, generic field $K$ finitely
generated over $\mathbf Q$, and normal line $V$, the checked
criterion is
\[
 s\text{ is induced by }V^\times(S)
 \quad\Longleftrightarrow\quad
 s_K\text{ is induced by }V_K^\times(K).
                                                               \tag{FR.1}
\]
The extension from the generic point is the key algebraic
step: the finite Kummer classes of a rational vector are
unramified at every codimension-one point, so every valuation
of its local coordinate is divisible by all positive integers
and hence zero. Normality extends the vector and its inverse.
The exact commutative-algebra input is EGA~IV$_2$,
Th\'eor\`eme~5.10.5 and its proof, pp.~115--116, as cited
in the supplied proof. Injectivity of inflation along
$G_K\twoheadrightarrow\pi_1(S)$ and injectivity of the full
Kummer map then identify the section and give uniqueness.
The argument does not use finite generation of $\operatorname{Pic}(S)$.

For the marked diagram
\[
 \pi_1(U)\longrightarrow G\longleftarrow G_K
                  \longrightarrow G_k
\]
with identified number field $k$ and regular finitely
generated $K/k$, Mochizuki's field and curve reconstruction
theorems recover $K$ and the generic marked curve.
The source statements are \emph{The local pro-$p$ anabelian
geometry of curves}, Theorems~A and~B (author manuscript
pp.~3--4), Theorem~16.5 together with its following full
profinite comparison (p.~86), and Theorem~17.1 (p.~87).
The statements, the full profinite comparison, and the proof
of Theorem~17.1 were checked directly in the supplied source.
Their regularity and base-Galois hypotheses are present here.
The full proof of Theorem~16.5 was not re-proved in this run.

Once the generic curve has been recovered, a nonzero normal
vector $b$ gives the membership test
\[
 [s_K]-[s_b]\in\kappa_K(K^\times)
       \subset H^1(G_K,\widehat{\mathbf Z}(1)).             \tag{FR.2}
\]
Changing $b$ translates by an element of the same Kummer
image, so the test is independent of that choice. F149-GENERIC
then makes it global. This is a functorial criterion on
realizable full marked diagrams. It is neither a finite
algorithm nor a recognition theorem for arbitrary profinite
groups, and it is not a section-surjectivity theorem.

\paragraph{What the tangential literature already supplies.}
Deligne, \emph{Le groupe fondamental de la droite projective
moins trois points}, \S15.19, p.~239, constructs relative
tangential specialization; \S15.26, p.~242, relates it to
formal coordinates. Stix, \emph{On cuspidal sections of
algebraic fundamental groups}, author version dated
24 February 2012, Proposition~16(2), p.~8, identifies the
cusp-inertia normalizers. His Theorem~17, p.~8, distinguishes
cuspidal packets, and Corollary~32, p.~12, recognizes a
cuspidal section by a cyclotomically normalized inertia
subgroup under the non-Tate hypothesis. These statements
recognize the cusp or its completed packet; they do not
assert that every completed Kummer class comes from an
actual nonzero tangent vector.

The origin of that group-theoretic cusp recognition is
Nakamura's anabelian weight filtration: \emph{Galois rigidity
of the \'etale fundamental groups of punctured projective
lines}, J. Reine Angew. Math.~411 (1990), \S3, and
\emph{Galois rigidity of pure sphere braid groups and
profinite calculus}, J. Math. Sci. Univ. Tokyo~1 (1994),
\S2.1, Theorem~2.1.1 and its proof, pp.~94--97. The latter
statement and complete displayed proof, including
Lemmas~2.1.2 and~2.1.4, were read in the original source;
the former locator is also explicitly given by Stix in \S5. The full
1990 proof was not retrieved and is not claimed checked.

Hoshi--Matsumoto--Nakayama, \emph{Tangential morphisms via
log arithmetic geometry}, RIMS-1995 (June 2026),
Propositions~1--3 and Remark~1, pp.~2--3 of the text
(PDF pp.~3--4), give a log-geometric formulation, coordinate
independence up to homotopy, and composition. In particular,
Remark~1 makes dependence on the normal cotangent basis
explicit. This supports the functoriality used by F149; it
does not remove the generic marking in (FR.2).

Sa\"idi's \emph{Good Sections of Arithmetic Fundamental
Groups}, arXiv:1010.1313v1, Theorems~2.6--2.7,
pp.~33--39, treats lifting uniformly good sections to
cuspidally abelian groups over slim, regular fields.
This is a different lifting problem. The published related
paper is \emph{The cuspidalisation of sections of arithmetic
fundamental groups}, Adv. Math.~230 (2012), 1931--1954;
its abstract was checked, but no identification of its
statement numbers with the preprint is asserted here.

\paragraph{The negative statement survives a full read.}
F149-UNMARKED-CURVE uses
$S=\mathbf G_{m,\mathbf Q}$ and
$U=S\times(\mathbf P^1\setminus\{0,1,\infty\})$.
Scaling the geometric base group
$\widehat{\mathbf Z}(1)$ by a profinite unit $u$, with
$u_2=3$ and $u_\ell=1$ for $\ell\ne2$, extends to the
full relative curve diagram and preserves its arithmetic
quotient and cusp data. It carries the tangent section of
$t b$ to slope $u^{-1}$ on base inertia. Actual global units
are $a t^m$ with $m\in\mathbf Z$, and their slopes are
integers; $u^{-1}$ has two-adic component $1/3$ and is not
an integer. The proof checks conjugacy as well: inertia
coboundaries vanish on geometric base inertia.

Thus a reference tangent vector and a rational base section
do not fix the missing marking when automorphisms may move
$G$. This does not conflict with Mochizuki's full-faithfulness
theorem, because that theorem retains the appropriate base
field Galois marking. No counterexample with $G$ fixed
pointwise, with its scheme realization fixed, or with $G_K$
retained is claimed.

\paragraph{Attribution and novelty.}
The anabelian reconstruction, tangent functors, cusp
normalizers and completed Kummer packets are established
literature. The K17 contribution reviewed here is their
assembly into the exact generic-to-global test (FR.1)--(FR.2)
and the explicit smaller-diagram automorphism obstruction.
No exhaustive priority claim for those formulations has
been established by this literature search. In Grothendieck's
terms, the result addresses the distinction between the
geometric subtorsor and the full profinite torsor on
cote~149 scan~85, with precisely the additional generic
marking of scans~83--84.


% END evidence/geometry_tower/f149_review.tex


% BEGIN evidence/geometry_tower/f149_normalizer_comparison.tex
\subsection{KT-F149-COMP: the two normalizer comparisons agree in scope}
\label{sec:kt-f149-comp}

GRH\_NIGHT AG-C149-NORMAL is an explicit isotrivial instance
of the general characteristic-zero comparison K17
F149-RELATIVE, and it also verifies the homotopy injection
which F149-RELATIVE retains as a hypothesis. Here is the
precise comparison, including the line-bundle convention.

Let $S/k$ be smooth, proper and geometrically connected in
characteristic zero, and choose
$\tau:\mathcal L^{\otimes n}\simeq\mathcal O_S$.
AG-C149-NORMAL defines
\[
 U=\{v\in\mathcal L^\times:
                         \tau(v^{\otimes n})\ne1\}.
\]
The zero-cusp normal line is $\mathcal L$; its conormal
line in F149-RELATIVE's convention is $\mathcal L^\vee$.
Thus their integral Chern classes differ by the stated
normal/conormal sign, not by an arithmetic discrepancy.

To verify the injection, pass to the finite \'etale torsor
$T=\{e:\tau(e^{\otimes n})=1\}$, and, after geometric
base change, choose a connected component. The pulled-back
family is the product of that proper component with
$\mathbf P^1\setminus(\{0,\infty\}\cup\mu_n)$.
SGA~1, Expos\'e X, Corollaire~1.7 injects the fibre
fundamental group into this product group. The finite
\'etale map from the product to $U$ injects its fundamental
group into $\pi_1(U)$. Right exactness is the relative
homotopy sequence of SGA~1, Expos\'e XIII, Proposition~4.1
and Examples~4.4. This proves exactly the sequence
(R149.1) required by F149-RELATIVE in this family.

F149-RELATIVE consequently identifies the zero-cusp
normalizer with $\pi_1(\mathcal L^\times)$.
AG-C149-NORMAL supplies an explicit map in the reverse
direction: the filling map $U\to\mathcal L^\times$
restricts to an isomorphism on the normalizer. The two
descriptions identify the same extension over $\pi_1(S)$,
with the same fibre inertia after the normal/conormal
convention is fixed. Tangential choices and conjugacies
must be transported together; an arbitrary automorphism
of an extension acting identically on kernel and quotient
need not itself be the identity.

AG-C149-FULLLIFT concerns pullback along each arithmetic
section $G_k\to\pi_1(S)$, not a splitting over all of
$\pi_1(S)$. Hence its coexistence with the nonsplit
normalizer in AG-C149-CUSP is consistent with
F149-RELATIVE. For the example
$S:y^2=x^3-x$ over $\mathbf Q$ and its nontrivial
two-torsion line, the global integral class is nonzero,
while its pullback along each arithmetic section vanishes.
Neither statement addresses the independent Ihara kernel
question LM6 A.45.

The complete supplied AG-C149-NORMAL and AG-C149-FULLLIFT
proofs were read, as was F149-RELATIVE. Their source
locators are GRH\_NIGHT.tex, lines 9063--9227, and
K17.tex, lines 13494--13644. The elliptic class computation
AG-C149-ELLIPTIC is carried from the intake, not
independently recomputed here. This comparison concerns
cote~149 scans~79--82; scan~85's geometric-subtorsor
recognition is the separate assertion reviewed above.


% END evidence/geometry_tower/f149_normalizer_comparison.tex


% BEGIN evidence/geometry_tower/route2_gate.tex
\subsection{KT-R2-POINT and the actual geometric gate}
\label{sec:kt-r2-point}

\begin{lemma}[KT-R2-POINT]
For the literal question in item~7, over
$\kappa=\overline{\mathbf F}_2$, the chosen point $P$ imposes
no extra condition on existence. A connected finite \'etale
cyclic cover $D\to E$ with primitive factor of two-rank zero
can be pointed above every $P\in E(\kappa)$.
\end{lemma}

\begin{proof}
A finite connected cover is surjective. Its fibre over $P$ is
a nonempty finite \'etale $\kappa$-scheme, hence a disjoint
union of copies of $\operatorname{Spec}\kappa$. Choosing one
of these points gives the pointing. Neither the cover nor its
primitive abelian factor is changed. All equations, maps,
$P$ and the chosen lift descend to some common finite field;
geometric two-rank is unchanged by this descent or by a
further extension of constants.
\end{proof}

This removes the parameter $P$ from the \emph{geometric}
existence question. It does not remove the prescribed-field
pointing condition in AN\_GRH A45--A52: those results measure
the rational two-primary group over a fixed finite field.

\paragraph{An exact genus-two theta-orbit formulation.}
Let $E/\kappa$ have genus two, let $L\in\operatorname{Pic}^0(E)$
have exact odd order $n>1$, and let $D\to E$ be its cyclic
cover. Write
\[
 0\longrightarrow\mathcal O_E\longrightarrow
 F_*\mathcal O_E\longrightarrow B_E\longrightarrow0
\]
for the locally exact differential sequence, using absolute
Frobenius and its semilinear action. Since the characteristic
is two, $B_E$ is a line bundle of degree one. Put
\[
 \Theta_{B_E}=\{M\in\operatorname{Pic}^0(E):
                         H^0(E,B_E\otimes M)\ne0\}.
                                                               \tag{R2.1}
\]
This is a translate of the Abel--Jacobi image of $E$.
The remaining existence question is exactly to find such
an odd-order $L$ for which every orbit of multiplication by
two on $(\mathbf Z/n\mathbf Z)^\times$ contains an $i$
with $L^i\in\Theta_{B_E}$.

Here is a proof of this particular form of the standard
Raynaud criterion. The character decomposition of
$H^1(D,\mathcal O_D)$ has summands $H^1(E,L^i)$.
For $i\ne0\pmod n$, Riemann--Roch and $H^0(E,L^i)=0$
give $\dim H^1(E,L^i)=1$. Tensoring the locally exact
differential sequence by $L^i$, then applying the projection
formula and coherent cohomology, identifies
\[
 \ker\bigl(F:H^1(E,L^i)\to H^1(E,L^{2i})\bigr)
       =H^0(E,B_E\otimes L^i).                              \tag{R2.2}
\]
The map is semilinear; each arrow between these one-dimensional
spaces is either zero or a semilinear isomorphism.
On a complete doubling orbit, stable Frobenius rank is zero
precisely when at least one arrow is zero. The primitive
factor selects the characters with $\gcd(i,n)=1$; because
$n$ is odd, its idempotent is integral at two and computes
this same stable part. Therefore, if
$h_n=\operatorname{ord}_n(2)$,
\[
 f_2(P_n)=h_n\cdot
 \#\{\text{doubling orbits of primitive }i:
                   L^i\notin\Theta_{B_E}\text{ for all }i
                   \text{ in the orbit}\}.                 \tag{R2.3}
\]
The identification of stable coherent Frobenius rank with
two-rank and the character version of the locally exact
differential criterion are the standard inputs; compare
Pop--Sa\"idi (2003), \S1.3, p.~112, carried in AN\_GRH A48,
and Sa\"idi, \emph{On complete families of curves with a
given fundamental group in positive characteristic},
Manuscripta Math.~118 (2005), \S\S1--2, pp.~428--431.
The displayed formulation is an explicit proof of the form
needed here, not a new general Raynaud theorem.

\paragraph{What was and was not advanced.}
The Rosati-fixed $\mathcal O_5$ multiplication remains an
extra restriction on the allowed genus-two curves. It is not
used in (R2.1)--(R2.3), and no argument here forces the required
orbit incidences on that restricted locus. The known failures
AN\_GRH A36--A44 and the sharper finite-field alternatives
A45--A52 are retained as supplied background. Searches of the
Raynaud--Sa\"idi--Tamagawa and Ozman--Pries literature did not
supply the missing all-curves incidence theorem. In particular,
Ozman--Pries, \emph{Ordinary and almost ordinary Prym varieties},
Theorem~1.1, concerns ordinary Pryms over a generic curve in a
fixed $p$-rank stratum; it is not an existence theorem for
two-rank-zero primitive factors over every prescribed
$\overline{\mathbf F}_2$-curve. No universal exponent or
annihilator-index bound in the actual primitive rational group
has been obtained, and no curve for which every odd cyclic
cover fails has been exhibited.


% END evidence/geometry_tower/route2_gate.tex


% BEGIN evidence/geometry_tower/rm_frobenius_avoidance.tex
\subsection{Arbitrary finite Frobenius data coexist with real-multiplication theta avoidance}
\label{sec:kt-rm-frobenius-avoidance}

The following is a strengthening of the \emph{carried method} in
AN\_GRH~A37, rather than a new claim about its published moduli
inputs.  A37 constructs avoidance with one specially selected
Frobenius class.  The additional argument below permits any
prescribed finite symplectic Frobenius class, including at the
primes to be avoided.  It identifies a limit of arguments using
finite torsion data; it does not address an exact Weil polynomial
or a fixed ground-field size.

Put $\mathcal O=\mathbf Z[T]/(T^2+T-1)=\mathcal O_5$.
All polarizations below are the canonical principal polarizations,
and the specified $\mathcal O$ action is fixed by Rosati.

\begin{theorem}[KT-RM-FROB-AVOID]
Let $m\geq1$ be odd, let $S$ be a finite set of odd primes, and
let
\[
                    \gamma\in\operatorname{SL}_2(\mathcal O/m\mathcal O).
                                                               \tag{RF.1}
\]
In every prescribed nonempty open of the ordinary genus-two
Hilbert--Blumenthal locus in characteristic two there is a curve
$C/\mathbf F_Q$,
with its specified integral Rosati-$\mathcal O$ action defined
over $\mathbf F_Q$, such that:
\begin{enumerate}
\item $Q\equiv1\pmod m$, and the action of $Q$-power Frobenius
      on the rank-two $\mathcal O/m$ module
      $J=\operatorname{Pic}^0(C^{(1)})[m]$ is conjugate to $\gamma$
      in an $\mathcal O$-linear symplectic frame;
\item the actual Raynaud theta support satisfies
\[
 \Theta_C(\overline{\mathbf F}_Q)\cap
 \bigl(J_{C^{(1)}}(\overline{\mathbf F}_Q)\{S\}\setminus\{0\}\bigr)
                          =\varnothing.                       \tag{RF.2}
\]
\end{enumerate}
The notation $\{S\}$ permits every exponent and every product
of primes in $S$.  No disjointness between $S$ and the prime
divisors of $m$ is assumed.  No bound for $Q$ is asserted.
\end{theorem}

\begin{proof}
We use two precise carried inputs.  AN\_GRH~A37, Steps~1--4,
gives an ordinary genus-two universal family on a geometrically
irreducible Hilbert--Blumenthal level scheme, full joint monodromy
$\operatorname{SL}_2(\mathcal O/n)$ at every odd level $n$, and
a nonempty open avoiding nonzero theta $n$-torsion for each fixed
$n$.  Its inverse-different pairing convention includes the
ramified prime five.  AN\_GRH~A19 gives the following actual
geometric statement: if $\Phi^d$ fixes $J[\ell]$ for every
$\ell\in S$, then every $S$-primary point on a translated genus-two
Abel curve defined over $\mathbf F_Q$ lies in
$J(\mathbf F_{Q^d})$.  It does not require semisimple Frobenius.

\paragraph{1. A lift of any finite class with bounded fixed torsion.}
At each $\ell\in S$, prescribe the reduction of $\gamma$ modulo
$\ell^{v_\ell(m)}$ when $\ell\mid m$; when $\ell\nmid m$,
prescribe the identity modulo $\ell$.  Let $D\geq1$ be a common
multiple of the orders of the resulting actions on the underlying
rank-four $\mathbf F_\ell$ modules.  Thus every lift $G_\ell$
of these prescribed classes satisfies $G_\ell^D\equiv1\pmod\ell$.

There is such a lift
\[
 G_\ell\in\operatorname{SL}_2(\mathcal O\otimes\mathbf Z_\ell)
 \quad\hbox{with}\quad
             \det_{\mathbf Z_\ell}(G_\ell^D-1)\ne0.          \tag{RF.3}
\]
Here the determinant is on the underlying rank-four module.
We give the existence argument explicitly.  Reduction maps for
$\operatorname{SL}_2$ over the local factors of
$\mathcal O\otimes\mathbf Z_\ell$ are surjective: lift a matrix
with determinant one and correct an entry using a unit cofactor,
then pass through the successive quotients.  This works at five
as well as at unramified primes.  Fix one lift $G_0$.

The regular function
$g\mapsto\det_{\mathbf Q_\ell}(g^D-1)$ on
$\operatorname{Res}_{F\otimes\mathbf Q_\ell/\mathbf Q_\ell}
\operatorname{SL}_2$, where $F=\mathbf Q(\sqrt5)$, is not zero.
For example $g=\operatorname{diag}(a,a^{-1})$, with a rational
$\ell$-adic unit $a$ of infinite order, has nonzero value.
Every congruence coset of this group is Zariski dense.  An
elementary verification uses the big cell
\[
 G_0\begin{pmatrix}
       1+x&y\\ z&(1+yz)/(1+x)
     \end{pmatrix},\qquad
 x,y,z\in\ell^r(\mathcal O\otimes\mathbf Z_\ell),             \tag{RF.4}
\]
for sufficiently large $r$.  Over $\mathbf Q_\ell$ this is an
open dense cell in an irreducible algebraic group; after choosing
a $\mathbf Z_\ell$ basis of $\mathcal O\otimes\mathbf Z_\ell$,
a nonzero numerator polynomial cannot vanish on the entire
product of six open $\ell$-adic balls.  The last assertion follows
by induction on the number of variables from the finiteness of
the roots of a nonzero one-variable polynomial.  Hence the
prescribed coset contains a point satisfying (RF.3).

Set
\[
 a_\ell=v_\ell\det_{\mathbf Z_\ell}(G_\ell^D-1),
 \qquad N=\prod_{\ell\in S}\ell^{a_\ell}.                  \tag{RF.5}
\]
These are finite integers.  Choose an odd integer $L$ divisible
by $m$ and by every $\ell^{a_\ell+1}$, and prescribe a class
$\delta\in\operatorname{SL}_2(\mathcal O/L)$ which reduces to
$\gamma$ modulo $m$ and to $G_\ell$ modulo
$\ell^{a_\ell+1}$ for every $\ell\in S$.  If
$v_\ell(m)>a_\ell+1$, retain the already prescribed higher
precision of $G_\ell$.  The Chinese remainder theorem gives
$\delta$; outside $S$, retain the prescribed components of
$\gamma$.

For any integral matrix $\Phi$ congruent to $\delta$ at this
level, polynomial continuity of powers and determinants gives
\[
 v_\ell\det(\Phi^D-1)=a_\ell\quad(\ell\in S).              \tag{RF.6}
\]
On a prime-to-two Tate module, the kernel of $\Phi^D-1$ on
$T_\ell J\otimes\mathbf Q_\ell/\mathbf Z_\ell$ is therefore
killed by $\ell^{a_\ell}$.  For completeness, its Smith normal
form has diagonal entries with total valuation $a_\ell$;
each diagonal valuation is at most their sum.  Equivalently,
the adjugate identity gives this annihilator directly.

\paragraph{2. Specialize inside the finite theta-avoidance open.}
Choose an auxiliary odd level $h\geq3$ prime to $L$.
Apply the generic avoidance in A37 to $N$, inside the prescribed
nonempty ordinary open, and then take a nonempty smooth affine
open $U$ on which the actual theta has no nonzero $N$-torsion.
After descending all data to a finite field
$k_0=\mathbf F_{q_0}$, enlarge it so that
$q_0\equiv1\pmod{hL}$.  This does not affect geometric
irreducibility or the theta-avoidance property.

The full level-$L$ torsor $V\to U$ is geometrically integral
with constant group $\operatorname{SL}_2(\mathcal O/L)$, by
the full joint monodromy input in A37.  Twist it by $\delta$
exactly as in A37, Step~5.  The affine descent datum is effective,
and the twist is geometrically integral.  Lang--Weil supplies a
point over $\mathbf F_{q_0^r}$ for every sufficiently large $r$.
Choose $r\equiv1\pmod{\operatorname{ord}(\delta)}$.
The resulting point of $U$ has arithmetic Frobenius class
$\delta$ in its level-$L$ torsor.  Let $C/\mathbf F_Q$ be
the universal curve at this point, with $Q=q_0^r$.
The actual Raynaud line bundle, its theta support, and the
specified real multiplication are defined over this exact field.
Thus item~1 holds, and $\Phi$ satisfies (RF.6).

\paragraph{3. Pass from a finite level to every supported order.}
Since $\Phi^D$ is the identity on each $J[\ell]$ for
$\ell\in S$, A19 puts every geometric $S$-primary theta point
in $J(\mathbf F_{Q^D})$.  By (RF.6) and the annihilator
calculation after it, that rational $S$-primary group is killed
by $N$.  The specialization was chosen to avoid all nonzero
$N$-torsion on theta, so no such point exists.  This proves
(RF.2).  If $S$ is empty, only the torsor specialization is
needed.  All constructions started inside the specified open,
which proves the final locus assertion.
\end{proof}

\begin{corollary}[KT-RM-FROB-NONDETERMINATION]
For every $b\geq1$ there are two ordinary genus-two curves over
finite fields of characteristic two, each with a specified
integral Rosati-$\mathcal O_5$ action, whose Frobenius actions
on $4159^b$-torsion are isomorphic as symplectic
$\mathcal O_5/4159^b$ modules, such that:
\begin{enumerate}
\item the first curve has an actual odd cyclic cover of degree
      $4159$ with primitive two-rank zero;
\item the second curve has no such cover of degree $4159^a$
      for \emph{any} $a\geq1$.
\end{enumerate}
Both field sizes can be required to be $1$ modulo $4159^b$.
No equality of the two field sizes or of their Weil polynomials
is asserted.  The common Frobenius class is chosen from the
actual positive curve; this corollary does not assert a positive
curve for an arbitrary prescribed class $\gamma$.
\end{corollary}
\begin{proof}
Use the actual A13 curve and its integral real multiplication,
as proved in the carried A33 supplement.  A13 provides a theta
point of order $4159$ and verifies
$\langle2,-1\rangle=(\mathbf Z/4159)^\times$; A11, or the
genus-two orbit criterion KT-R2-POINT, gives the first cover.
Extend its constants to contain $\mathbf F_{128}$ and so
that the field size is one modulo $4159^b$.  The torsion
point and the needed roots of unity are then rational, so
the cyclic cover is defined after this extension.  Record
the resulting symplectic Frobenius class $\gamma$.
Apply KT-RM-FROB-AVOID with $m=4159^b$, this $\gamma$, and
$S=\{4159\}$.  The resulting second curve has the same stated
finite torsion data and no nonzero $4159$-primary theta point.
Every primitive-character doubling orbit of a degree-$4159^a$
cyclic cover would need a theta point of that exact order, so
the absence proves item~2 for every exponent.
\end{proof}

\paragraph{Scope of the new obstruction.}
For witnesses whose orders have prime factors in a prescribed
finite set, this rules out a sufficient incidence criterion
depending only on a prescribed finite congruence precision of
Frobenius and the specified real-multiplication action, even
when arbitrarily high but finite precision is permitted in
advance.  It does
not rule out a criterion involving the exact field size, a
Weil polynomial, or primes selected adaptively from the
individual curve.  It does not exhibit a curve for which every
odd cyclic cover fails.  The universal existence question of
item~7 remains unchanged.

\paragraph{Source and proof-check scope.}
The carried inputs are AN\_GRH~A19, equations (FS.3)--(FS.10),
and A37, especially its Steps~1--5 and equations
(RM.2)--(RM.4), (RM.8)--(RM.13).  Their entire relevant
proofs were read in the supplied TeX.  Their named published
moduli, descent, and Lang--Weil inputs are cited as already
accepted results, rather than newly re-proved here.  The
additional argument is the arbitrary-congruence lift (RF.3),
the valuation freezing (RF.6), and their use in the same joint
torsor specialization.

% END evidence/geometry_tower/rm_frobenius_avoidance.tex


% BEGIN evidence/geometry_tower/route2_all_order_attempt/rm_theta_propagation_limit.tex
\subsection{A bounded all-order test of real-multiplication propagation}
\label{sec:kt-rm-theta-propagation-limit}

This note records an obstruction to one attempted incidence argument.
It does not produce a cyclic cover, a counterexample to the universal
cover question, or a new bound in A45--A52.

\begin{proposition}[KT-RM-THETA-STABILIZER]
Let $k$ be an algebraically closed field of characteristic two, let
$X/k$ be an ordinary smooth projective genus-two curve, and let
$A=\operatorname{Pic}^0(X)$.  Let $B$ be the actual locally exact
differential line bundle and
\[
 \Theta=\{\mathcal O_X(R)\otimes B^{-1}:R\in X(k)\}\subset A.
                                                               \tag{RT.1}
\]
Suppose $\mathcal O_5\hookrightarrow\operatorname{End}(A)$ is
a unital integral action fixed by the Rosati involution for the
canonical principal polarization.  If $0\ne u\in\mathcal O_5$
and $a\in A(k)$ satisfy
\[
                         u(\Theta)\subset\Theta+a,
                                                               \tag{RT.2}
\]
then $u=1$ or $u=-1$, and $a=0$.
Consequently, for every fixed $u\in\mathcal O_5\setminus\{0,1,-1\}$
and every $a\in A(k)$, the set
\[
                    \{z\in\Theta(k):u z\in\Theta+a\}
                                                               \tag{RT.3}
\]
is finite.  This conclusion concerns all torsion orders at once,
but does not exclude choosing exceptional points in that finite set.
\end{proposition}

\begin{proof}
Since $u\bar u=N_{\mathbf Q(\sqrt5)/\mathbf Q}(u)$ is a nonzero
integer, $u$ is an isogeny.  Thus (RT.2) is equality of reduced
curves, and restriction gives a finite surjective morphism of
genus-two curves of some degree $d\ge1$.  Write $D=\deg_A(u)$
for the degree of the isogeny, which need not equal $d$.

Let $h=[\Theta]\in\operatorname{NS}(A)\otimes\mathbf Q$.
Translation acts trivially on this group, and proper pushforward
therefore gives $u_*h=d h$.  The finite-flat identity
$u_*u^*=D$ shows that $u_*$ is invertible on this finite-dimensional
space.  Hence
\[
                       u^*h=(D/d)h.                          \tag{RT.4}
\]
The polarization map $[L]\mapsto\lambda_\Theta^{-1}\phi_L$
is injective on $\operatorname{NS}(A)\otimes\mathbf Q$.
The pullback identity
$\phi_{u^*L}=u^\vee\phi_Lu$ sends $u^*h$ to $u^\dagger u=u^2$.
It follows that
\[
                   u^2=c\in\mathbf Q_{>0},\qquad c=D/d.
                                                               \tag{RT.5}
\]
Since $u$ is an algebraic integer, $c$ is a positive integer.
Degrees on a surface give
$D^2=\deg(u^2)=\deg([c])=c^4$, so $D=c^2$ and $d=c$.
In particular $u^2=d$.

Factor the curve morphism as a power of relative Frobenius
followed by a separable map.  The Frobenius twist still has
genus two.  Riemann--Hurwitz for the separable factor gives
$2\ge2\deg(f_{\mathrm{sep}})$, so that factor has degree one
and is an isomorphism.  Therefore $d=2^r$ for some $r\ge0$,
and the original curve map is radicial.  The equation
$u^2=2^r$ in $\mathbf Q(\sqrt5)$ forces $r=2b$ and
$u=\pm2^b$: for odd $r$ it would put $\sqrt2$ in
$\mathbf Q(\sqrt5)$.

An ordinary genus-two curve in characteristic two has three
distinct Weierstrass points $W_0,W_1,W_2$.  Their classes
$z_i=\mathcal O_X(W_i)\otimes B^{-1}$ are three distinct
points of $\Theta\cap A[2]$, because $2W_i\sim K_X\sim2B$.
They are nonzero because $B$ is ineffective in the ordinary
case.  These are the carried A35, (T2F.3)--(T2F.4), facts for
the actual Raynaud translate.  If $b>0$, $u=\pm[2^b]$ sends
all three $z_i$ to zero, contradicting injectivity of a
radicial curve map on geometric points.  Thus $b=0$ and
$u=\pm1$.

The identity $2B\sim K_X$ makes $\Theta$ symmetric under
$[-1]$, by the hyperelliptic involution.  Hence (RT.2) now says
$\Theta+a=\Theta$.  Such an $a$ lies in the kernel of
$\phi_{\mathcal O_A(\Theta)}$, which is trivial because this
polarization is principal.  Thus $a=0$.

Finally, if (RT.3) were infinite, the closed subset
$\Theta\cap u^{-1}(\Theta+a)$ would be all of the irreducible
projective curve $\Theta$.  This would give (RT.2), already
excluded for the indicated $u$.
\end{proof}

\paragraph{The attempted arithmetic implication and its exact failure.}
Write $T^2+T-1=0$ for the chosen generator of $\mathcal O_5$.
The proposed propagation was to start with odd torsion on theta,
use powers of $T$ to produce more incidences, and combine these
with Frobenius to meet all primitive doubling orbits of one
cyclic subgroup.  Proposition KT-RM-THETA-STABILIZER rules out
the required global preservation of theta by every nontrivial
power of $T$, even allowing a translation.  For a finite collection
of such fixed pairs $(u,a)$, there are only finitely many exceptional
incidences of the indicated form.  An incidence at one specially
selected torsion point is still possible and has not been excluded.

Arithmetic Frobenius does preserve the actual theta curve over
its field of definition, but it need not preserve the cyclic
subgroup generated by a chosen theta point.  Even if a chosen
point $z$ has exact order $n$ and satisfies $\Phi z=vz$ for
$v\in(\mathbf Z/n)^\times$, the elementary sufficient condition
$\langle2,-1,v\rangle=(\mathbf Z/n)^\times$ still requires
that particular theta point.  The real-multiplication relation
does not supply it.  No existence assertion is inferred from
this conditional observation.

The converse attempt to avoid every odd theta order on one
finite-field curve is already excluded by the carried A42,
(PE.14)--(PE.15): every such theta curve has infinitely many
nonzero odd-order points.  Their full orders and their distribution
inside one cyclic subgroup remain uncontrolled.  The geometric
generic avoidance of A21/A37 cannot be specialized simultaneously
at all odd levels by the finite-level argument.  In particular,
the construction imposing $Q\equiv1\pmod m$ for every prescribed
finite $m$ supplies varying finite fields; one fixed finite
$Q$ cannot satisfy these congruences for all odd $m$.

These limitations coexist with the carried ordinary
$\mathcal O_5$ example A13/A33 that has a degree-$4159$
primitive two-rank-zero cover.  Thus they imply neither failure
of all odd cyclic covers nor their existence on an arbitrary
ordinary $\mathcal O_5$ curve.  This attempt is classified
\emph{nonmovement} for the universal Route~2 gate.

\paragraph{Exact inputs and proof-reading scope.}
For the elementary abelian-variety inputs see Milne,
\emph{Abelian Varieties}, version 2 January 2022:
Proposition~8.1 and Theorem~8.2, pp.~12--14;
Corollary~12.8, p.~22; Remark~12.13, p.~24;
Lemma~16.2 and its proof, p.~30; and Proposition~17.2 and
its proof, p.~35.  The finite-flat push--pull identity is
Stacks, Section~42.15, Lemma~42.15.2 (tag~\texttt{02RF});
proper pushforward and flat pullback preserve algebraic
equivalence by their compatibility with families in the same
section.  The curve factorization and genus preservation are
Stacks, Section~53.13, Proposition~53.13.7 and
Lemma~53.13.5 (tag~\texttt{0CCV}); Riemann--Hurwitz is
Section~53.12, Lemma~53.12.2 (tag~\texttt{0C1B}).
The statements and the displayed proofs of these local inputs
were read; foundational results cited inside them are imported.
The actual theta-characteristic and two-torsion facts are the
accepted A35 proof, and the odd-torsion abundance is accepted
A42, not newly derived here.  Retrieval metadata and exact
source hashes are retained with this note.

% END evidence/geometry_tower/route2_all_order_attempt/rm_theta_propagation_limit.tex


% BEGIN sections/framework.tex
\section{The arithmetic and geometric dictionary}
\label{sec:framework}

\subsection{Integral cohomology and rational mixed Tate motives}

Put $U_\ell=\operatorname{Spec}\mathbf Z[1/\ell]$.
For odd $\ell$, the exact comparison carried as K17-VK-PT is
\[
 K_{2e}(\mathbf Z)\otimes\mathbf Z_\ell
       \simeq H^2_{\acute et}(U_\ell,\mathbf Z_\ell(e+1)),
 \qquad e\ge2\text{ even}.
 \tag{FW.1}
\]
Its strict negative-twist dual is the Selmer group
\[
 \mathcal S_\ell(e)=\ker\left(
 H^1(U_\ell,\mathbf Q_\ell/\mathbf Z_\ell(-e))
 \longrightarrow
 H^1(\mathbf Q_\ell,\mathbf Q_\ell/\mathbf Z_\ell(-e))\right).
 \tag{FW.2}
\]
The relevant motivic weight is $e+1$: formally $2(e+1)-2=2e$.
Thus the completed new finite degrees concern twists $27$, $29$, $31$,
$33$, $35$, $37$, $39$ and $41$; all eight have checked finite vanishing theorems.
The realizations in (FW.1), not merely their rationalizations,
retain the finite groups whose orders the Hecke certificates
control. This use of the modern integral edge comparison is
not an assertion that a factorial-normalized Chern character
is invertible at every prime.

Deligne--Goncharov, \emph{Groupes fondamentaux motiviques de
Tate mixte}, author manuscript \path{Tate.pdf}, \S1.6,
(1.6.11)--(1.6.12), printed p.~14, gives
\[
 \operatorname{Ext}^1_{\mathrm{MT}(F)}(\mathbf Q(0),\mathbf Q(n))
       =K_{2n-1}(F)\otimes\mathbf Q,\qquad
 \operatorname{Ext}^2_{\mathrm{MT}(F)}(\mathbf Q(0),\mathbf Q(n))=0.
\]
Their \S1.7 defines the unramified category over $\mathcal O_S$;
Proposition~1.9, printed p.~16, retains higher-weight
$\operatorname{Ext}^1$ and vanishing $\operatorname{Ext}^2$.
Proposition~2.3, printed p.~19, gives the free graded Lie algebra,
while \S2.4 distinguishes its canonical abelianization from
noncanonical free lifts. These are rational categories.
Consequently their second Ext vanishing cannot prove the
vanishing of the torsion in (FW.1). Nor is an integral abelian
category of mixed Tate motives with identical Ext formulas
being postulated here.

The coefficient sheaves in (FW.1) are continuous inverse limits
of the finite prime-invertible Tate sheaves. SGA~4, Expos\'e~IX,
3.2--3.3 is the finite Kummer/Hilbert~90 locator used in K17.
The convention for weights is Deligne, \emph{Weil~II},
1.2.1--1.2.2, p.~153, and 1.2.5(iv), p.~154:
$\mathbf Q_\ell(1)$ has weight $-2$ for geometric Frobenius.
The rank-two modular systems have pure weight $2e$ before
normalization; their twists in the trace comparison have weight
zero. The three trace terms are those of
$\mathbf Q_\ell(-e)\oplus\mathbf Q_\ell\oplus\mathbf Q_\ell(e)$
in the corresponding convention. Purity of rational systems
does not eliminate their integral congruences.

\paragraph{$\lambda$-rings and $L$-functions.}
SGA~6, Expos\'e~V, Definition~2.1, Example~2.2(b) and
Lemma~2.2.1 locate the $\lambda$-operations on $K_0$ defined
by exterior powers. The determinant and the first
$\gamma$-filtration piece record line bundles. The normal-line
Chern class in cote~149 is in this framework; its Kummer
identification is SGA~5, Expos\'e~VII, Proposition~3.4(i).
The higher groups in (FW.1) use the separately cited higher
$K$-theory comparison. No $K_0$ identity alone proves their
vanishing.

The CM formula in KT-GL2-CM describes the Euler factors of
rank-one Hecke characters induced from $E$. The checked endoscopic
construction decomposes a rank-three parameter
into a character and a rank-two base change. KT-TWOFIELD-661
uses the exact special values
$L(-2e,\chi_D)=-B_{2e+1,\chi_D}/(2e+1)$.
The rational class component requires the separate ordinary
Bernoulli number $B_{\ell-1-e}$; the two $L$-value questions
are not interchangeable. No new crystalline, prismatic, or
divided-power comparison is asserted.

\paragraph{The integral fixed locus and its conormal.}
For the source deformation algebra, the involution defines
the closed fixed locus by $Q=(s-\kappa_S(s))$.
Its conormal is $Q/Q^2$: this is the language of EGA~IV$_4$,
Definition~16.1.2, whose locator and notation are also identified
in the Stacks Project, Tag~\texttt{01R1}.  The affine formula
is Tag~\texttt{01R3}, and the base-change surjection, with an
isomorphism under flat base change, is Tag~\texttt{0473}.
These standard conormal statements do not assert flatness
of a fixed locus.  KT-UNSQUARED-NORMAL-JACOBIAN identifies
the required odd cotangent with its pullback to the base
augmentation.  KT-ODD-CONORMAL-SELMER-EQUIVALENCE then
identifies its Matlis dual with the actual strict Tate-sheaf
extension group in (FW.2), after extending coefficients to $O$.
Thus the formal geometry preserves the integral arithmetic
obstruction exactly.  The scalar plus normalized fourth
symmetric-power normal representation at an adjoint point
is the coefficient system displayed in KT-ADJOINT-ODD-NORMAL-MODULE;
it is not the Tate dual of that displayed coefficient system.
Generic automorphic point detection alone does not determine
the integral conormal or remove its vertical torsion.

\subsection{The Teichm\"uller tower and the letter to Faltings}

Grothendieck's \emph{Esquisse d'un programme}, \S2,
printed pp.~9--12 in \emph{Geometric Galois Actions~I},
describes the compatible Galois action on the profinite
fundamental groupoids of the curve-moduli tower, with gluing
and forgetting markings. Printed pp.~9--10 single out the
tripod and the first nontrivial moduli level; pp.~11--12
discuss generators, relations, and carefully chosen base-point
systems. Ihara's question here concerns the fixed field of
the kernel after selecting one pro-$p$ realization of the
tripod action. A full-profinite faithfulness statement is
not a proof of the claimed maximal pro-$p$ image, and a
rational unipotent image is another realization again.

The letter to Faltings, 27 June 1983, in the same volume,
pp.~49--58, makes the base-Galois augmentation explicit in
(1)--(3), pp.~50--52; p.~51 discusses successive hyperbolic
curve fibrations and Artin neighbourhoods. Formula~(7),
p.~54, is the point-to-section map. The correction on
pp.~56--57 treats cuspidal packets, and the postscript on
p.~58 corrects the open-image/dominance conditions and gives
a proposed cyclotomic fixed-subgroup criterion. These
locators were read in the published transcription.
The finite-field counterexample below concerns an extension
beyond the characteristic-zero section conjecture formulated
there. The higher-dimensional counterexamples concern a
specific maximal geometric prime-to-$p$ quotient; they do
not refute the letter's full profinite formulation.

\subsection{Cote and run crosswalk}

The following table specifies the object and implication for
each family of findings. These are applications or limitations
of the named frameworks, not claims that Grothendieck's texts
already contain the later numerical theorems.

\small
\begin{longtable}{>{\raggedright\arraybackslash}p{.24\textwidth}
>{\raggedright\arraybackslash}p{.34\textwidth}
>{\raggedright\arraybackslash}p{.34\textwidth}}
\toprule
Finding family & Geometric or arithmetic object & Exact programme locator and change\\
\midrule\endhead
KT-ONECHAR-CLOSED; KT-TRACE-MULTIPLICITY; KT-QUARTIC-BOUND &
Strict extension classes of negative Tate sheaves, dual to
integral $H^2(U_\ell,\mathbf Z_\ell(e+1))$. &
Alp\"oge step~17; K17-S17-UNIT and VK-PT. Bridges Q06,
GRH\_NIGHT AG-Q18: a bound on every $K_{4k}$ is now explicit,
while the all-weight zero statement remains open.\\
KT-U3-WEIGHT, DIM, SECOND, SPECTRUM; KT-U3-ENDO-EXACT; KT-U3-SPECTRAL-SPLIT; KT-GL2-CM &
Algebraic automorphic coefficient spaces and Frobenius Euler
factors, with cohomological weight shift retained. &
K17-U3-HECKE, U3-GENERAL-COVERAGE, GL2-GENERAL; prompt
items1(a)--(b),2. These give uniform finite models, not a
universal theorem on their residual primes.\\
KT-K52-K56-K60-VANISH; KT-K64-VANISH; KT-K68-VANISH; KT-K72-VANISH;
KT-K76-VANISH; KT-COMPONENT-EXTENSION-19;
KT-K80-VANISH; KT-COMPONENT-EXTENSION-20 &
Twists27,29,31,33,35,37,39,41 of arithmetic $H^2$, the cyclotomic even components
$[-26],[-28],[-30]$, $[-32],[-34],[-36],[-38]$, and the corresponding odd ramified modules. &
K17 Q2-HIGHER4; K17-COMPONENT-CONTROL; LM6 A.44--A.45;
bridges Q06. The character interval grows to $[-16,20]$.
No complete remaining-prime kernel assertion follows.\\
KT-C7-EXPLICIT; KT-SOULE-STAB &
Torsion cohomology of arithmetic groups and integral higher
$K$-theory of rings of integers. &
K17 VK-EXP7; Soul\'e2003 \S4, Theorem2 and Proposition4;
Lemma9 pp777--778 requires the stated general-field correction.\\
KT-TWOFIELD-661; KT-SCHUR-PIVOT-SUPPORT &
Actual Eisenstein congruences, reflected rational characters,
and units in finite coefficient rings. &
K17-REFLECTOR (RB.1); prompt1(b)--(c). They identify
limitations of proposed uniform shortcuts; neither gives
nonzero $K_{4k}$ or a Vandiver counterexample.\\
KT-U3-SYM2-INJECT; KT-ORDINARY-SYM2-ADM3;
KT-ORDINARY-SYM2-ALL-LEVELS;
KT-STABLE-PROPAGATION; KT-ADJOINT-ACTUAL-HOST-POINT &
Weight-zero rank-three lisse sheaves with ordinary flags,
residual Tate constituents, cuspidal automorphic lifts, and
actual points of the normalized deformation ring. &
K17-S17-UNIT, actual Adm3U and CoversPt sources; the integral
Tate/Selmer dictionary (FW.1) is retained. Their congruences
can coexist with zero integral motivic cohomology. The strict
Selmer tangent line has the opposite sign from the adjoint
root; the new conormal comparison below retains that sign.\\
KT-ODD-CONORMAL-GATE; KT-ODD-CONORMAL-SELMER-EQUIVALENCE;
KT-FINITE-GRADED-NORMAL-MODULE &
An integral fixed-locus conormal, its weight and involution
summands, and the strict negative Tate extension group. &
EGA IV$_4$,16.1.2; SGA4 IX,3.2--3.3 for the finite Tate-sheaf
language; actual source ModelLower and ShaInjGen.
The conormal is exactly the Selmer dual. A unit minor remains
equivalent to the arithmetic vanishing problem; finite
generation does not supply flatness or a unit determinant.\\
KT-STABLE-HECKE-SELFDUAL-LOCALIZATION;
KT-TRACE-INVISIBLE-SELMER-TANGENT;
KT-RESIDUAL-TRACE-DISCRIMINANT &
Integral Hecke orders, Frobenius characteristic polynomials,
and extension data lost by the trace map at the reducible point. &
Alp\"oge step17, ModelLower and ShaInjGen; the conormal
framework EGA IV$_4$,16.1.2 and (FW.2).
An evaluated complementary factor gives an all-split-prime
Hecke identity, while the actual Selmer tangent is invisible
to traces and the split-residual matrix trace order has
nonunit discriminant. No unit normal comparison follows.\\
KT-FIELD-CHARACTER-REDUCTION; KT-EXPLICIT-TWO-UNIT-FAMILY;
KT-BERNOULLI-PREFIX-FORMULA; KT-THIRD-UNIT-CHARACTER-OBSTRUCTION;
KT-AUXILIARY-UNIT-SELMER-BRIDGE; KT-RATIONAL-POSITIVE-ROOT-UNIT;
KT-ODD-DISCRIMINANT-TWO-UNIT-FAMILY; KT-ODD-DISCRIMINANT-GD &
Special values of odd Dirichlet characters with finite
Teichm\"uller twists and a fixed negative discriminant. &
Prompt1(b)--(c), K17-REFLECTOR and the integral dictionary
(FW.1)--(FW.2). The exact same-discriminant three-unit
condition is not rational mixed-Tate Ext vanishing.
For the explicit weight progression (EU.2), two fields
supply its first two units and opposite local behavior
at the residue prime. The repaired odd-discriminant family
(OD.1) also satisfies every actual GD condition at $N=2$. The faithful cyclic-character projector
annihilates all units imported from multiquadratic fields; full
compositum $p$-rationality is stronger than the third-unit condition.
These are character decompositions of the finite Tate/Dirichlet
coefficient sheaves in (FW.1)--(FW.2), not vanishing of the
original rational strict extension group.\\
KT-GMA-REDUCIBILITY-DOES-NOT-BOUND-EXTENSIONS;
KT-RESIDUAL-CYCLICITY-SCOPE;
KT-POLARIZED-EXTENSION-GRAPH-OBSTRUCTION;
KT-AUXILIARY-CLASS-UNIT-MIXED-CUP;
KT-GLOBAL-QUADRATIC-MIXED-PERSISTENCE;
KT-CUBIC-TWISTED-COHOMOLOGY;
KT-CUBIC-RATIONAL-OBSTRUCTION-SPACE;
KT-ORDINARY-CUBIC-MASSEY-CONDITION;
KT-ORDINARY-CUBIC-SHUFFLE;
KT-CUBIC-BOCKSTEIN-NORM-IDENTIFICATION;
KT-CUBIC-EXCISION; KT-CUBIC-MATCHED-NORM-SYSTEM;
KT-CUBIC-P-ONLY-NORM;
KT-CUBIC-FINITE-COEFFICIENT-NORM;
KT-ORDINARY-QUARTIC-CONDITION; KT-QUINTIC-RELATIVE-TARGET;
KT-SOURCE-TATE-PAIRING-SCOPE;
KT-HSIEH-DIRECT-SCOPE;
KT-EQUIVARIANT-POINT-NOT-CUBIC-JET;
KT-SPARSE-VALUATIVE-CUBIC-EXTRACTION;
KT-SOURCE-CUBIC-EXTRACTION-SCOPE;
KT-MASSEY-LITERATURE-SCOPE &
Integral extension modules of a residual sum of three
rank-one sheaves; involution signs and mixed compositions
of extensions. &
The conormal framework EGA IV$_4$,16.1.2 and the exact
strict coefficient sheaf in (FW.2). Trace and reducibility
ideals alone do not bound the needed extension module.
The abstract graph retains irreducibility after five full
cohomology vanishings; it is not an arithmetic realization
or a counterexample to motivic vanishing. The next mixed
cohomology spaces distinguish an auxiliary quadratic Tate
sheaf killed by its unit from a rational Tate sheaf with
a strict class-group dual and explicitly retained local terms.
The specified ordinary cubic jet has a canonical relative
obstruction whose dimension is the full rational
$\omega^{1-e}$ class component. This is a finite-$S$
duality and lifting calculation, not a nonvanishing theorem
for an obstruction forced by an actual source. A flat formal
deformation ring can have a characteristic-zero point while
its specified cubic jet fails. The proved extraction criterion
requires additional valuation and coefficient-support hypotheses.
The exact shuffle preserves the local primitives, and for
$\mathbf F_p$ coefficients its raw value has the proved
Bockstein and norm-ideal interpretation. Etale excision on
$\operatorname{Spec}\mathbf Z[1/p]$ preserves the local primitives,
and Kummer identifies the result with the full rational Picard
component. An explicit absolute regular-module lift repairs the
proper-system matching step, and a matched norm representative
gives the exact ideal class. This uses Milne II.2.4 and II.2.9,
with the Kummer/Chern language already located at SGA5 VII.3.4(i).
The class has not been evaluated and source extraction is not proved.
The inertia exponents also prohibit a Tate-valued tensor contraction
of the residual adjoint module. Compact-support duality instead
gives the distinct unit-plus-class module, by Milne II.3.1 and II.3.3;
the actual matrix cocycle retains products of other extension entries.
Recent primary Massey theorems do not remove the fixed finite-set
and ordinary-normalization requirements.\\
KT-ORDINARY-RANK2-UNIT-EXCLUSION;
KT-AUXILIARY-UNIT-NO-TYPEII; KT-ADJOINT-BRANCH-UNIT-EXCLUSION;
KT-ODD-ADJOINT-SCALAR-NORMALIZATION;
KT-SELFDUAL-SOURCE-ODD-DESCENT;
KT-NONESSENTIALLY-SELFDUAL-SOURCE-GATE &
An ordinary rank-two sheaf, its opposite integral lattice,
and the resulting strict extension of a quadratic Tate sheaf. &
Alp\"oge step17, actual Type-II points; source
\path{TanDet.fin_quot} and \path{typeII_facts}.
A quadratic Bernoulli unit excludes this source branch
without a weight bound. The adjoint extension recovers and
removes a quadratic scalar sheaf, excluding the specified
normalized adjoint Type-I branch. Scalar normalization and
quadratic automorphic descent extend this to every self-dual
actual stable source. With the three units, only stable
sources that are not essentially self-dual remain. This is
an exclusion of specified weight-zero lisse sheaves, while
the remaining integral extension obstruction in (FW.2)
and the uniform auxiliary-field choice are still open.\\
KT-SMALL-ORDER-COMPONENT-REDUCTION;
KT-SEXTIC-CLASS-PRODUCT;
KT-SEXTIC-REGULATOR-RATIO-FOUR;
KT-ORDER-SIX-COMPONENT-REDUCTION;
KT-ORDER-FIVE-LENGTH-ONE;
KT-ORDER-FIVE-STRICT-CYCLIC;
KT-ORDER-FIVE-PDIV-WEIGHT;
KT-ODD-FAMILY-RATIONAL-MIRROR &
Character summands of the cyclotomic ideal class group and
their integral arithmetic cohomology realizations. &
K17-(K2), K17-REFLECTOR; (FW.1), SGA4 IX,3.2--3.3 for the
finite sheaf framework. Known low-degree class-number bounds
remove rational character orders one through four. The exact
sextic relative-regulator argument removes order six as well,
leaving order five or at least seven. The total length of the
four order-five ideal-class components is at most one. Exact
finite-coefficient transport and full VK-PT duality imply cyclic
strict and $K$ groups, with unknown exponent in general and
exponent at most $p$ when $p\mid e$; no full-length
class-group identification is used. The ordinary class-field
product gives the needed integral ratio; no rational-motive
argument is substituted for this torsion calculation. The
quadratic odd-mirror progression was already controlled by
reflection; its auxiliary construction does not settle a new
uncontrolled rational component.\\
KT-JACOBI-ODD-COEFFICIENT;
KT-JACOBI-TWO-UNIT-FIELD;
KT-JACOBI-TWO-UNIT-GD; KT-JACOBI-TYPEII-FIELD;
KT-JACOBI-THIRD-UNIT-NONIMPLICATION;
KT-JACOBI-FIXED-WEIGHT-REFORMULATION;
KT-KOHNEN-THETA-NONCONGRUENCE;
KT-EISENSTEIN-QUADRATIC-DETECTOR;
KT-SENARY-GENUS-TRACE-DETECTOR;
KT-ANALYTIC-RESIDUE-NONSATURATION;
KT-WITT-PULLBACK-MASS-OBSTRUCTION;
KT-GENUS-PROJECTOR-INTEGRALITY;
KT-NEIGHBOR-CONSTANT-LINE-LIMIT;
KT-SMALL-QUADRATIC-SPAN-LIMIT;
KT-EISENSTEIN-HECKE-EVALUATION-IDEAL;
KT-BETA3-SHORT-WINDOW;
KT-COHEN-SPECIAL-HEAT-FILTRATION;
KT-COHEN-WEIGHT-TWO-BOUND;
KT-COHEN-GENERAL-FIRST-UNIT-BOUND;
KT-SUBP-SEED-THIRD-UNIT-NONIMPLICATION;
KT-WEIGHT-TWO-KERNEL-NONFACTORIZATION;
KT-WEIGHTED-CHARACTER-ATTEMPT;
KT-HIGH-ORDER-TWIST-TARGET-SCOPE &
Modular coefficient sections, their reduction modulo $p$,
and two special-value units for an odd imaginary quadratic
field of discriminant less than $4p$. &
K17-S17-UNIT and actual Counts.GD; (FW.1)--(FW.2).
The heat lift, cusp integrality and Sturm bound provide an
actual field from the stated Bernoulli seed. Nonquadratic
mirror families are included. The simultaneous third unit
and the remaining rank-three sheaf exclusion are still open.
The selected field $p=31,D=-79$ has only the first two units;
this does not exclude another simultaneous choice. In the
automatic $e=(p+1)/2$ progression the two twisted units become
the weight-two value for the real field of discriminant $p|D|$
and the weight-three value for the imaginary field of
discriminant $D$. Their product is also a specified integral modular-symbol coefficient,
by Heumann--Vatsal Theorem3.23 and Corollary4.7. Quadratic
unit detection and simultaneous nonvanishing remain open.
The senary genus trace realizes the same product as a finite
weighted sum of integral coefficient sections. Its mass factor
can vanish modulo $p$ despite every individual coefficient being
a unit. Integral neighbor polynomials preserve that mass factor,
and a residual constant character can lift outside every spinor
character without transferring to the unitary source. The specified
weight-two residue kernel does not factor through the third value;
the direct weighted Burgess estimate misses the sufficient short
window. Higher order finite twists change the exact coefficient
sheaf and Selmer target in Counts.GD and Vlen.\\
KT-COHEN-UP-REFORMULATION;
KT-HECKE-SQUARE-CLASS-SCOPE &
Reduction of half-integral weight modular coefficient sections
and their fundamental square classes. &
Prompt1(b)--(c), K17-S17-UNIT and (FW.1)--(FW.2).
The complete coefficient identity $H_{(p+3)/2}=U_pH_2$ modulo $p$
includes the conductor-$p$ and constant terms. The coefficient
recurrences preserve each fundamental square class and do not
force a simultaneous unit. A formal eigenseries used to show
this limitation is not asserted to be modular. The missing
global modular constraint leaves the same arithmetic gate.\\
KT-INDIVISIBILITY-LITERATURE-SCOPE;
KT-KIMURA-DISPLAYED-PROOF-SCOPE &
Reduction of modular sections and specified Dirichlet special values,
with the same arithmetic coefficient sheaf retained. &
Prompt1(b)--(c), K17-S17-UNIT and (FW.1)--(FW.2).
The weight-two principal-part comparison sharpens one automatic
two-unit seed to $|D|\leq p-4$. A field inside this bound can still
fail the third unit. A sufficient $O(p^{2/5})$ window has no proved
seed-existence theorem, and the heat filtration stays linear.
The reviewed literature and displayed-proof gaps do not supply
simultaneous units or the original integral Tate extension vanishing.\\
KT-QUARTIC-QUARTER &
A refined finite-order bound for integral motivic cohomology
of Spec Z. &
K17 VK-PT and VK-EXP7; Deligne--Goncharov1.6(1.6.11)--(1.6.12)
records the rational comparison, which forgets this torsion.\\
KT-FF-SEC &
A point-free proper curve over $\mathbf F_2$ with a split
arithmetic fundamental-group extension. &
Cote149 scans35--36, AN\_GRH A05. Point separation does not
imply section surjectivity. SGA1 IX.6.1 supplies the
fundamental-group exact sequence.\\
KT-POLY-FAIL; KT-ARTIN-FAIL; KT-PRODUCT-SEP &
Wild isotrivial twisting lost on maximal geometric
prime-to-two completion; separation retained by maps to curves. &
Extension of cote149 A1 and AN\_GRH A05; Faltings letter p51.
SGA1 X.1.7, XIII.4.6 are the proper and prime-to-characteristic
Kunneth inputs in the respective proofs.\\
KT-F149-REVIEW &
Geometric tangent vectors inside a completed Kummer torsor,
with the generic Galois marking. &
Cote149 scans83--84, generic diagram; scan85, geometric
subtorsor. K17 F149-GENERIC, MARKED-RECOGNITION and
UNMARKED-CURVE. EGA IV$_2$,5.10.5 supplies normal extension;
the required anabelian reconstruction is separately cited.\\
KT-F149-COMP &
The normalizer of cusp inertia, the nonzero-vector line bundle,
and its integral first Chern class. &
Cote149 scans79--82; GRH\_NIGHT AG-C149-NORMAL, FULLLIFT,
CUSP and K17 F149-RELATIVE. Normal/conormal signs are kept
consistent. SGA5 VII.3.4(i) locates the Chern/Kummer language.\\
KT-R2-POINT and its remaining incidence question &
Odd cyclic etale covers of an ordinary genus-two curve,
primitive two-rank, and the Raynaud theta divisor. &
AN\_GRH A33 and A45--A52; original failed constructions
A36--A44 are retained. Over $\overline{\mathbf F}_2$ pointing
is automatic, but the RM incidence problem remains open.\\
KT-RM-FROB-AVOID; KT-RM-FROB-NONDETERMINATION &
The finite etale torsion local systems of the universal
Hilbert--Blumenthal Jacobian and the actual relative
Raynaud theta support. &
AN\_GRH A19 (FS.3)--(FS.10), A37 Steps1--5;
cote149 scans35--36 and its route-two substitute A33.
Arbitrary finite symplectic Frobenius data coexist with
avoidance at all exponents of any fixed odd prime support.
The all-odd-order cover question remains open.\\
KT-ORDER-FIVE-GREENBERG-UNIT-GATE;
KT-ORDER-FIVE-ABSTRACT-NORM-LIMIT &
Ideal classes and global units in the cyclotomic tower of
finite covers of $\operatorname{Spec}\mathbf Z[1/p]$;
their integral Iwasawa modules. &
K17-GREENBERG-KERNEL, Q3-IRREG, and Esquisse\S2.
The unit-norm index measures capitulation exactly in the stated
one-prime setting. The abstract module $\Lambda/(T+p)$
shows that a length-one base class and norm identities alone
do not give bounded growth. No arithmetic realization is asserted.
The exact unit-norm index and the remaining kernel gate stay open.\\
KT-SATURATED-RANKIN-COHEN-SCOPE &
Differential operators on modular coefficient sections and
reduction modulo $p$. &
K17-S17-UNIT, Prompt1(b)--(c), and (FW.1)--(FW.2).
The stated saturated order-$p-1$ symbol, even after addition
of the ordinary product, cannot universally select equal indices.
This is a coefficient-algebra limitation; no theorem that other
modular operators fail is asserted. The simultaneous-unit gate remains.\\
KT-GMA-MULTIGENERATOR-CONORMAL &
The conormal of the reducibility locus in an integral
three-character deformation space. &
EGA IV$_4$,16.1.2 is the conormal terminology already used above;
the arithmetic input is the separately cited Berger--Klosin setup.
A $d$-generated adjacent module contributes $d$ copies of the
adjacent Selmer length. The strengthened congruence inequality
and actual-source identification are not proved. This does not
eliminate torsion in the integral motivic realization (FW.1).\\
KT-QUINTIC-CLASS-PAIRING-SCOPE &
Cartier duals of finite Tate sheaves, unramified Selmer local
conditions, and the reciprocity pairing. &
The (FW.1)--(FW.2) integral cohomology dictionary and
Morgan--Smith2024, \S\S5.1--5.2.
The even order-five character has an odd Cartier dual.
The published symmetric reciprocity identity supplies no
perfect alternating self-pairing of the remaining class.\\
KT-ADJACENT-FACTOR-SCOPE;
KT-BK2026-NORMALIZED-DEPTH-REPAIR;
KT-BK2026-LONG-ROOT-SIGN-SCOPE &
Integral conormal lengths under coefficient extension,
and the sign of a block-algebra anti-involution. &
EGA IV$_4$,16.1.2 supplies conormal terminology;
the precise arithmetic hypotheses are those audited in
Berger--Klosin2026, \S\S6,8--10.
The numerical scalar-extension repair is conditional on
the retained arithmetic inputs. The block-sign countermodel
identifies an algebraic proof gap without asserting an
arithmetic counterexample. No extra congruence length is obtained.\\
KT-RM-THETA-STABILIZER &
The principally polarized Jacobian, its actual Raynaud theta
translate, Rosati-fixed real multiplication, and radicial
maps of curves. &
AN\_GRH A35(T2F.3--4), A42(PE.14--15), A45--A52;
cote149 scans35--36 and the route-two substitute A33.
Only $\pm1$ in $\mathcal O_5$ preserve theta up to translation.
This prevents automatic propagation by real multiplication;
exceptional torsion incidences and the universal cover gate remain.\\
Ihara gate &
Actual arithmetic pro-$p$ action and its plus Iwasawa modules,
including the cubic companion errors. &
Esquisse2; LM6 A.44--A.45; K17 Q3-IRREG and CUBIC-ODD;
GRH\_NIGHT AG-Q29. Rational Brown/Deligne--Goncharov
freeness does not supply arithmetic Triple-NV.\\
\bottomrule
\end{longtable}
\normalsize

The FGA representability programme provides the language of
families, Picard functors, and moduli. No FGA theorem is
invoked as a missing anabelian reconstruction or integral
$K$-theory comparison in this report, so no spurious FGA
statement number is assigned to one of these new assertions.

% END sections/framework.tex

\part{Unresolved questions and unproved proposals}

% BEGIN sections/open_gates.tex
\section{Unresolved questions and unproved proposals}

This part is not a proof of the main conjecture. It records the
remaining gates and the approaches that still require a
mathematical argument.

\paragraph{All weights.}
The closed source chain, complete source coverage, total
dimensions, and uniform exact spectrum algorithms do not
force the resulting determinants to be units. The all-weight
arithmetic gate is to eliminate the common residual support,
using ordinary reflection, other Frobenius elements, another
quadratic field, or a sharper admissible-source comparison.
KT-SMALL-ORDER-COMPONENT-REDUCTION applies the exact known
real quadratic, cubic and quartic class-number bounds to
remove every rational character of order at most four.
KT-ORDER-SIX-COMPONENT-REDUCTION additionally removes order six,
using the integral sextic class-number ratio and the exact general
relative-regulator and even-character bounds.
The remaining rational character therefore has order five or
order at least seven. KT-ORDER-FIVE-LENGTH-ONE bounds the sum
of the four order-five cyclotomic ideal-class lengths by one:
at most one order-$p$ component can remain.
KT-ORDER-FIVE-STRICT-CYCLIC then makes the strict rational
Selmer group and its $K$-group dual cyclic, with unknown exponent.
This does not remove the group or supply the total ordinary
$E$-Selmer cyclicity hypothesis.
A fourth-power upper bound for the logarithm of the order
cannot be promoted to order one. The simultaneous classical
congruence at661 shows why merely changing split primes in
the two classical levels need not suffice.

\paragraph{The refined spectral gate.}
The full stable/endoscopic decomposition is now proved for positive
even $e$. The rank-three branch uses only its stable factor, while
the classical branch remains. An exact four-operator test using
both split Hecke directions at seven and thirteen still leaves
one primitive common kernel for each small residual modulus in
degrees $26,28,30$. Thus adding the opposite split direction does
not settle those examples by itself. Their symmetric-square
origin is now proved by KT-U3-SYM2-EISENSTEIN, including the
residual congruence on the whole Galois group. The remaining
stable factors are units at both primes in these three degrees.
KT-ORDINARY-SYM2-ADM3 also verifies every literal admissibility
predicate. Along three explicit infinite progressions these
actual admissible systems have nonunit trace difference for
every group element, while the rational $K$-component vanishes.
KT-ADJOINT-ACTUAL-HOST-POINT now realizes these systems as
actual Type-I and TwoPt points on the chosen local component,
after an allowed finite coefficient extension. Consequently
every literal CoversPt cover contains their isomorphism class.
The untwisted strict-Selmer tangent line and the adjoint root
line have opposite complex-conjugation signs, as proved in
KT-HOST-TANGENT-SIGN. The new integral projector construction
does retain this sign through a conormal length comparison.
Its reverse map, however, proves that the odd weight-$-1$
conormal is exactly the Matlis dual of the original strict
Selmer group. A unit maximal minor is consequently an
equivalent formulation, not a smaller arithmetic gate.
The actual finite graded algebra and normal module have now
been constructed, with a surjection onto this conormal.
Flatness, the necessary support on every generic component,
and a unit determinant over this algebra remain unproved.
The conditional branchwise norm estimate cannot be applied
without those hypotheses; generic rigidity does not control
vertical contributions.

The precise auxiliary-field question is now (FV.8), simultaneous
nonvanishing of three degree-one generalized Bernoulli values
at one negative fundamental discriminant. The reviewed
individual-value theorems do not supply this assertion.
KT-EXPLICIT-TWO-UNIT-FAMILY now supplies its first two
units for two explicit discriminants in the progression
$e=(\ell-3)/2+(\ell-1)t$, when $\ell\equiv3\pmod4$;
the fields give opposite splitting behavior at $\ell$.
The third value in (EU.5) is still an arithmetic condition.
Those particular discriminants are even and do not satisfy
the odd-discriminant condition forced by actual GD with even
$N$; KT-EXPLICIT-TWO-UNIT-GD-SCOPE makes this limitation explicit.
KT-ODD-DISCRIMINANT-TWO-UNIT-FAMILY repairs that obstruction:
$D_s=-\operatorname{sf}(p+4s)$ supplies the first two units
and all actual GD conditions at $N=2,\psi=1$ in the same
progression. Its third unit remains open. The scope check
KT-ODD-FAMILY-RATIONAL-MIRROR proves that this particular
progression already has zero rational target component by
quadratic class-number reflection; it is not a new family of
previously uncontrolled rational components.
The new KT-JACOBI-TWO-UNIT-FIELD supplies a different, bounded
field construction. For an odd representative $3\le b\le p-2$
with $b\ne(p-1)/2$ and $B_{2b}/(2b)$ a unit, it produces
an odd negative fundamental $D$, prime to $p$, with $|D|<4p$
and both the class-number unit and $B_{1,\chi_D\omega^{b-1}}$
a unit. Under the stated genericity this field meets actual GD.
For $p\equiv3\pmod4$ and $e=(p+1)/2+(p-1)t$, the seed
is automatically a unit and the mirror character has full
order $p-1$. This family is not the quadratic-mirror family
already disposed of by reflection. Applying the same method
at $b\equiv2e+1$ gives a field excluding actual Type-II points
under its own seed condition, automatically for
$p\equiv7\pmod8$, $e=(p+1)/4+(p-1)t$.
These results do not supply both twisted units over one field;
units from two different discriminants cannot be combined into
the same-discriminant three-unit hypothesis.
KT-JACOBI-FIXED-WEIGHT-REFORMULATION identifies the two twisted
values in the first automatic progression with the weight-two
value of the real field of discriminant $p|D|$ and the
weight-three value of the imaginary field of discriminant $D$.
This exact congruence does not prove simultaneous nonvanishing.
KT-EISENSTEIN-QUADRATIC-DETECTOR gives a specified integral
modular-symbol coefficient equal to $-15/8$ times the same-field
product. The published full-symbol nonvanishing uses characters
of growing odd-prime-power order, so it supplies neither a unit
quadratic evaluation in this coefficient nor the class-number
condition for an unbounded discriminant.
The missing ordinary necessary direction has now been proved:
KT-AUXILIARY-UNIT-NO-TYPEII excludes actual Type-II source points
when the $a=2e$ auxiliary value is a unit, and the $a=e$ version
excludes the specified ordinary rank-two seeds.
KT-ADJOINT-BRANCH-UNIT-EXCLUSION extends the latter conclusion
to every normalized adjoint origin over $\mathbf Q$ satisfying
its stated determinant and adjoint hypotheses, by removing
the recovered quadratic scalar twist. There is no
prime-versus-weight restriction in these arguments.
KT-ODD-ADJOINT-SCALAR-NORMALIZATION removes the determinant
hypothesis for every projectively odd rank-two origin over
$\mathbf Q$. KT-SELFDUAL-SOURCE-ODD-DESCENT then supplies
such an origin for every self-dual actual stable source, using
quadratic descent and the self-dual $\mathrm{GL}_3$ adjoint
classification. Thus the $a=e$ unit excludes all self-dual
actual Type-I sources. Under the class-number unit, the exact
ordinary local conditions force determinant one and turn any
essential self-duality into actual self-duality, by
KT-ESSENTIALLY-SELFDUAL-REDUCTION.
The remaining field strategy now requires suitable fields
within actual GD and exclusion only of the stable rank-three
sources that are not essentially self-dual. This is a strict
conditional source reduction, not a uniform vanishing theorem.
The exact ordinary inertia and residual characters also rule
out monomial sources and finite projective image by
KT-ORDINARY-IMPRIMITIVE-EXCLUSION. This further specifies the
remaining source class without bounding its integral extension
modules or constructing a comparison theorem.
The precise bridge KT-AUXILIARY-UNIT-SELMER-BRIDGE identifies
the killed strict groups as the three quadratic-negative
characters, leaving the original rational negative-twist
group unchanged. KT-POLARIZED-EXTENSION-GRAPH-OBSTRUCTION
shows that those vanishings, even with rational positive-root
vanishing, cannot by a purely polarized matrix-algebra argument
exclude an irreducible ordinary source: mixed extension paths
can survive. This is an abstract scope obstruction, not an
arithmetic counterexample. The inspected comparison theorems
retain an extension upper bound or a cyclicity condition as an
input; the latter already forces the target vanishing in the
concrete split branches of KT-RESIDUAL-CYCLICITY-SCOPE.
The exact arithmetic calculation
KT-GLOBAL-QUADRATIC-MIXED-PERSISTENCE goes further:
the class-number unit kills the first global mixed obstruction
and permits its ordinary local correction. A hypothetical
strict negative class and a positive class therefore fit
together through the mixed quadratic term over
$k[s,t]/(s^2,t^2)$. This supplies neither a lift through all
higher orders nor a characteristic-zero source point, and
does not give a vanishing argument for the original class.
At the next mixed degree, KT-CUBIC-TWISTED-COHOMOLOGY kills
the $\eta\psi^{-1}$ obstruction space using the existing unit,
whereas KT-CUBIC-RATIONAL-OBSTRUCTION-SPACE retains the exact
$\omega^{1-e}$ class-group dual and possible auxiliary-prime
local quotients. No particular arithmetic higher obstruction
has been proved nonzero, and raw local $H^2$ nonvanishing
does not supply that missing assertion.
KT-ORDINARY-CUBIC-MASSEY-CONDITION further identifies the exact
relative obstruction, including local matching, for its specified
graded two-parameter jet. Its dimension is the full rational
$\omega^{1-e}$ class component. A rational $B_e$ unit removes
this cubic obstruction. Neither an arithmetic nonvanishing theorem
nor extraction of that jet from every relevant actual source is
proved, so the original vanishing gate is unchanged.
KT-CUBIC-EXCISION identifies this relative space naturally with
$H^2(G_{\mathbf Q,p},k(\psi))$ and, for scalar coefficients,
with the full rational class component. The explicit
KT-CUBIC-MATCHED-NORM-SYSTEM repairs the earlier unjustified
inference from proper-system terminology: zero absolute inflation
allows a lift preserving the specified truncated cocycle.
KT-CUBIC-P-ONLY-NORM then gives an exact matched norm-ideal
identity, retaining all local correction information. The norm
class has not been evaluated, and actual-source extraction remains
unproved. These comparisons do not assert a smaller arithmetic
vanishing gate.
KT-CUBIC-FINITE-COEFFICIENT-NORM extends the matched formula
to every finite coefficient field by polarization, without
calling an additive extension of that field cyclic.
KT-ORDINARY-QUARTIC-CONDITION retains a positive quartic
obstruction controlled sufficiently by the opposite $-2e$
auxiliary unit. KT-QUINTIC-RELATIVE-TARGET identifies the lower
quintic space with dimension $1+d$, where $d$ is the original
strict rational class dimension. Neither obstruction is evaluated.
KT-SENARY-GENUS-TRACE-DETECTOR gives a second exact same-field
product detector; its $p=31,D=-79$ example shows why individual
coefficient units do not imply a unit trace. An analytic residue
alone supplies no missing integral saturation.

The first-unit field choice now has a sharper special bound:
KT-COHEN-WEIGHT-TWO-BOUND gives $|D|\leq p-4$ when
$p\equiv3\pmod4$ and $b=(p+3)/2$.
The general principal-part argument gives $|D|\leq2b-7$
under its Bernoulli-seed hypotheses, with the explicitly
retained possibility $D=-p$ outside its prime-to-$p$ cases.
Neither result supplies the additional twisted unit.
The exact example $p=83,D=-43$ has the first two units
inside the sharper bound but $\beta_3=166$, divisible by $83$.
The sufficient short window $d^5\leq1600p^2$ guarantees
the third unit, but no first-unit seed is proved to exist there.
Finite quadratic orthogonality and the exact heat cycle do
not supply this existence. The canonical genus-mass projector
is integral precisely when the mass already is a unit.
The source-pairing and Hsieh-branch audits also leave the
original arithmetic extension and actual-source comparison open.

KT-EQUIVARIANT-POINT-NOT-CUBIC-JET exhibits a flat formal domain
with the required actions and a characteristic-zero point whose
specified quadratic jet does not lift cubically.
KT-SPARSE-VALUATIVE-CUBIC-EXTRACTION gives a positive extraction
criterion under explicit extra depth and sparse-support conditions;
these are not currently known for an arbitrary actual source.
Changing only the allowed level also retains the actual
ordinary symmetric-square systems, by
KT-ORDINARY-SYM2-ALL-LEVELS. No GD or universal comparison-set
claim for the changed level is made.

\paragraph{Higher finite degrees.}
The requested degrees $e=26,28,30$ have complete certificates and
checked vanishing proofs. The next degree $e=32$ is also complete:
KT-K64-VANISH uses two completed full initial-field replay methods,
both alternate-field unitary comparisons, and freshly replayed
classical data. KT-K68-VANISH now completes $e=34$ as well,
with two full initial-field replay methods and both alternate-field
unitary reports. KT-K72-VANISH completes $e=36$ using the full
sum-first replay and both complete alternate native reports.
KT-K76-VANISH completes $e=38$ with its full sum-first replay,
both alternate native reports and all classical comparisons;
the entire arithmetic proof and a fresh independent exact
audit have passed. Together with K17 this proves all $K_{4k}=0$
for $1\le k\le19$. KT-K80-VANISH now completes $e=40$:
the complete initial sum-first replay, both alternate native
reports and all classical comparisons pass. The whole proof
has author, root and independent peer PASS checks; the peer
also freshly recomputed both157-dimensional classical determinants.
Thus all $K_{4k}=0$ for $1\le k\le20$, with safe component
interval $[-16,20]$ and full-domain $\kappa_{41}$ onto.
No computation started in this continuation remains pending;
the all-weight arithmetic gate remains open.
Every actual complete report and failed launch is retained in the
evidence. No running process is recorded as a passed certificate.

\paragraph{Ihara and route two.}
The precise remaining Ihara and geometric cover conditions
are displayed in the following subsections. Neither is
settled by the new finite degrees. KT-RM-FROB-AVOID strengthens
the carried finite-support avoidance to arbitrary prescribed
finite Frobenius data. It rules out the proposed finite-data
incidence shortcut, but neither produces a counterexample to
all odd cyclic covers nor proves universal existence.
Any subsequent progress
must change one of these gates, not only restate it.

% END sections/open_gates.tex


% BEGIN evidence/main_gate/ordinary_eisenstein_ideal_attempt.tex
\subsection{Attempt HS1: the direction of an ordinary Eisenstein-ideal bound}

A possible all-weight route is to replace finite Hecke determinants
by an ordinary Eisenstein ideal and a special-value formula. The
specific result checked here does not supply the needed unit
exclusion.

Hsieh, \emph{Eisenstein congruence on unitary groups and Iwasawa
main conjectures for CM fields}, JAMS 27 (2014), 753--862,
Theorem~7.18, Definition~7.19 and Corollary~7.20
(author manuscript pp.~91--92), constructs an ordinary cuspidal
congruence from a nonunit product
\[
 L^S(\Psi)=L_p^S(-1,\Psi_+\tau_{K/F})L_p^S(\Psi,\Sigma).
\]
Besides primitivity away from $p$, the theorem imposes the
stated conductor, ordinary CM-type and $\mu$-invariant conditions.
The construction and its localized congruence ideal explicitly
assume that this product is not a unit. The corollary gives
\[
                \operatorname{Eis}(\Psi,S)\subseteq(L^S(\Psi)).
\]
The introduction, inequalities (L$\mid$E) and (E$\mid$S),
and Theorem~3 give the corresponding lower bounds on
congruence valuations and Selmer lengths. These exact
statements, and the complete short proof of Theorem~7.18,
were read; no proof audit of the entire article is claimed.

The direction matters. In a discrete valuation ring,
$I\subseteq(L)$ says $v(I)\geq v(L)$. If $L$ is a unit,
that inequality alone places no restriction on $I$:
$I=(\ell)\subseteq(1)$ is a concrete counterexample to
the proposed inference. An opposite inclusion, an equality,
or a separately proved upper bound on the appropriate
congruence module would be needed to infer a unit ideal
from a unit special value. Moreover, the actual theorem's
localized ideal is constructed in the nonunit case, so
its setup cannot simply be reused unchanged in the unit case.

No such opposite inclusion or equality for the precise
finite-level compact-unitary comparison of this report
has been established. Nor has the required specialization
from Hsieh's CM Selmer group to the rational strict negative
Tate Selmer group been proved. This approach therefore leaves
an actual arithmetic comparison to establish; it does not
turn the computed generalized-Bernoulli units into an
all-weight vanishing theorem.

The primary author-hosted manuscript is
\url{https://www.math.ntu.edu.tw/~mlhsieh/research/IMCCM.pdf};
its retrieval hash and local source audit are retained in the
evidence. The publication is cited as an imported theorem.

\paragraph{Attempt NV: simultaneous units and specialization.}
Ono--Skinner, \emph{Fourier coefficients of half-integral weight
modular forms modulo $\ell$}, Ann.\ of Math.\ 147 (1998),
453--470, main Theorem (author manuscript p.~2), treats one
good eigenform and all but finitely many primes. Its local
Corollary~2 (manuscript p.~3) still treats one form. These
statements do not give the same-discriminant product (FV.8),
nor a prime bound uniform in the varying weight. The primary
manuscript is
\url{https://uva.theopenscholar.com/files/ken-ono/files/029_8.pdf}.
Only the exact statements and their stated scopes are used here;
no new full-proof audit is claimed.

Castella--Liu--Wan's six-page \emph{Addendum to: Iwasawa--Greenberg
main conjecture for nonordinary modular forms and Eisenstein
congruences on GU(3,1)}, \S1, pp.~1--2, explains that the corrected
tower nonvanishing input gives infinitely many twists, insufficient
for their previous three-value selection. They use two auxiliary
families to recover a height-one assertion; \S5, Theorem~5.0.1,
p.~6 records the repaired conclusion. The primary source is
\url{https://web.math.ucsb.edu/~castella/CLW-addendum.pdf}.
The affected Hsieh special-value paper there is the 2014
\emph{Documenta} paper; it is distinct from the JAMS paper
cited in Attempt HI. We do not import an uncorrected
almost-all-twists assertion.

This suggests a two-family replacement for simultaneous units,
but KT-FIELD-SPECIALIZATION-LIMIT proves that height-one avoidance
alone cannot settle every specialized integral group here.
The needed additional argument must control integral specialization,
including higher-codimension contributions; merely substituting
two independent variables for (FV.8) does not prove it.

% END evidence/main_gate/ordinary_eisenstein_ideal_attempt.tex


% BEGIN evidence/main_gate/recent_arithmetic_route_audits.tex
\subsection{Two further literature routes and their precise limits}

\paragraph{KT-STABLE-CONGRUENCE-ROUTE-AUDIT.}
O.~Randal-Williams, \emph{Stable cohomology of congruence
subgroups}, Compositio Math.~162 (2026), 1274--1317,
\url{https://doi.org/10.1017/S0010437X26103315},
Theorem~A, p.~1275, factors congruence-subgroup cohomology
as an exterior algebra tensor completed cohomology in
degrees below $p-1$. Equation~(7.3), p.~1296, identifies
the latter with the cohomology of the fibre of the
completion map in $K$-theory. The remaining factor is
therefore arithmetic input, not an independently computed
vanishing. Section~7.3.2, p.~1297, explicitly retains the
unknown finite $K_{4k}$ terms. In \S7.5, p.~1301, those
terms occur in the fibre's exact sequence; the subsequent
examples assume Vandiver and additional low-degree
Bernoulli units. Thus applying the displayed factorization
does not eliminate the target group. This is a
statement-and-formula scope check of Theorem~A and
\S\S7.2--7.5, not a rereading of the paper's full proof.

\paragraph{KT-RANKIN-COHEN-NONVANISHING-SCOPE.}
R.~Takloo-Bighash, \emph{Simultaneous nonvanishing of
quadratic twists via Rankin--Cohen brackets},
\url{https://arxiv.org/abs/2609.19649}, Theorem~1.1 and
Corollary~1.2 in \S1, keeps a discriminant and a desired
number of forms fixed while the weight grows. Its
conclusion is nonzero complex central values for enough
level-one forms. It does not assert unit values at a
fixed residue prime, nor the three Dirichlet values at
one varying negative discriminant needed in (FV.8).
Consequently this theorem does not discharge that
simultaneous modular nonvanishing condition. Only the
stated theorem, definitions and scope were checked;
no full proof check or counterexample to its theorem
is claimed.

Both attempts leave the all-weight arithmetic gate
unchanged and are excluded from progress-review movement.

% END evidence/main_gate/recent_arithmetic_route_audits.tex


% BEGIN evidence/geometry_tower/ihara_gate.tex
\subsection{The remaining Ihara characters: no new arithmetic
nonvanishing theorem}
\label{sec:kt-ihara-gate}

The carried K17-GREENBERG-KERNEL states that finiteness of
$X^+$ implies $\Omega_p^*=\Omega_p$. Its source is
K17.tex, Theorem K17-GREENBERG-KERNEL, beginning at
line~9678. The supplied proof framework when odd torsion
persists is K17-CUBIC-ODD, beginning at line~10303.
Its conclusion requires, in addition to the structural
freeness and cyclicity assumptions, Pair-NV and Triple-NV
for the \emph{actual} arithmetic companion errors.
In the notation of that theorem, Triple-NV is
\[
 \left.
 \mathcal C_{o;a,b,c}(T_1,T_2,T_3)
 \right|_{(1+T_1)(1+T_2)(1+T_3)=\alpha}\ne0
                                                               \tag{IK.1}
\]
for at least one matching ordered odd triple for each root
$\alpha$ of the relevant odd companion polynomial.
The condition is (CU.1), K17.tex lines~10429--10440.

No arithmetic proof of (IK.1), and no positive-rank
replacement for it, has been obtained in this workstream.
The known finite-level $K$-group vanishings cannot be
silently replaced by all-weight control of the entire
plus Iwasawa module at an arbitrary prime. Brown's rational
mixed-Tate freeness likewise does not prove this integral
arithmetic nonvanishing condition. These are distinct
inputs in the supplied reduction, not alternative wordings
of the same theorem.

There is, however, an exact conditional dependence on the
main question. If all $K_{4k}(\mathbf Z)$ vanish, the carried
Kurihara equivalence gives Vandiver at every prime. The
carried K17-VANDIVER-KERNEL (K17.tex, beginning at line~9305)
then gives $\Omega_p^*=\Omega_p$ for all odd primes; K17-P2-MOMENT
already supplies $p=2$. Thus item~5 would follow from a
completed item~1, without a separate cubic calculation.
No converse and no new unconditional prime are asserted.

In the language of the \emph{Esquisse}, \S2, this remains
a question about the arithmetic action on the pro-$p$
level of the Teichm\"uller tower; the rational unipotent
mixed-Tate quotient does not by itself determine its
pro-$p$ kernel. The supplied LM6 A.44--A.45 and K17
Q3-IRREG rows therefore remain open precisely at the
arithmetic gate (IK.1), unless the main all-weight
vanishing theorem is supplied.


% END evidence/geometry_tower/ihara_gate.tex

\end{document}